\documentclass[11pt,openany]{book}
\usepackage[T1]{fontenc}
\usepackage{lmodern}
\usepackage[a4paper,margin=27mm,headheight=15pt]{geometry}
\usepackage{amsmath,amssymb,amsthm,mathtools,bm}
\usepackage{microtype,enumitem,longtable,booktabs,array}
\usepackage[dvipsnames]{xcolor}
\usepackage{fancyhdr}
\usepackage[unicode,colorlinks=true,linkcolor=MidnightBlue,citecolor=MidnightBlue,urlcolor=MidnightBlue]{hyperref}
\usepackage[nameinlink,noabbrev]{cleveref}
\numberwithin{equation}{chapter}
\newtheorem{theorem}{Theorem}[chapter]
\newtheorem{lemma}[theorem]{Lemma}
\newtheorem{proposition}[theorem]{Proposition}
\newtheorem{corollary}[theorem]{Corollary}
\theoremstyle{definition}
\newtheorem{definition}[theorem]{Definition}
\newtheorem{example}[theorem]{Example}
\theoremstyle{remark}
\newtheorem{remark}[theorem]{Remark}
\providecommand{\K}{\mathbb K}
\providecommand{\R}{\mathbb R}
\providecommand{\C}{\mathbb C}
\providecommand{\N}{\mathbb N}
\providecommand{\Z}{\mathbb Z}
\providecommand{\T}{\mathbb T}
\providecommand{\eps}{\varepsilon}
\providecommand{\id}{\mathrm{id}}
\providecommand{\dd}{\,\mathrm d}
\providecommand{\norm}[1]{\left\lVert#1\right\rVert}
\providecommand{\abs}[1]{\left\lvert#1\right\rvert}
\providecommand{\ip}[2]{\left\langle#1,#2\right\rangle}
\providecommand{\supp}{\operatorname{supp}}
\providecommand{\pvsupp}{\operatorname{supp}}
\providecommand{\dist}{\operatorname{dist}}
\providecommand{\spanop}{\operatorname{span}}
\providecommand{\Ran}{\operatorname{Ran}}
\providecommand{\Dom}{\operatorname{Dom}}
\providecommand{\spec}{\operatorname{spec}}
\providecommand{\esssup}{\operatorname*{ess\,sup}}
\providecommand{\conv}{\operatorname{conv}}
\providecommand{\wto}{\rightharpoonup}
\setlist{itemsep=.2em,topsep=.4em}
\setcounter{tocdepth}{1}
\emergencystretch=2em
\pagestyle{fancy}
\fancyhf{}
\fancyhead[LE,RO]{\thepage}
\fancyhead[LO]{\small Functional Analysis}
\fancyhead[RE]{\small\nouppercase{\leftmark}}
\renewcommand{\headrulewidth}{.3pt}
\title{Functional Analysis\\\large Unified Foundations and Graduate Theory\\\normalsize Expanded mathematical manuscript}
\author{MathGenome preparation notes}
\date{November 2026}
\begin{document}
\frontmatter
\maketitle
\chapter*{Purpose and conventions}
This manuscript organizes the important functional-analytic content of nine
specified textbooks into a common mathematical language. Sources are merged at
the level of precise assertions; a change of notation is not taken as a proof of
equivalence. The sources guide coverage rather than impose a single proof route.
Important genuinely different versions, estimates and counterexamples are kept.
The earlier mathematical manuscript and its formal code do not impose a ceiling
on the present scope or generality.

The scalar field $\K$ is either $\R$ or $\C$ unless specified otherwise.
Hilbert inner products are conjugate-linear in the first variable and linear in
the second. The continuous linear dual $X'$ is distinguished from the conjugate
dual. The Banach transpose $T'$ acts on linear functionals; the Hilbert adjoint
$T^*$ satisfies $\langle Tx,y\rangle=\langle x,T^*y\rangle$.
For an unbounded operator, its domain is part of its definition. Resolvents use
$R(\lambda,A)=(\lambda I-A)^{-1}$; unitary groups use $U(t)=e^{itA}$.
Any additional Fourier normalization is stated where it is used.

The prerequisites are ordinary set-theoretic mathematics (including choice),
linear algebra, general topology, measure and integration, and real and complex
analysis. A prerequisite is invoked in its precise mathematical form; any new
functional-analytic assertion requires its own proof or an explicitly identified
proved result. Specialist sequels are distinguished from material actually proved
here in the supplied source-disposition metadata.
\tableofcontents
\mainmatter
\part{Spaces, duality and Banach-space structure}

% BEGIN parts/part1.tex
% Part I: all labels carry the fa2:p1 prefix.

% BEGIN parts/part1-foundations.tex
\providecommand{\Q}{\mathbb Q}
\providecommand{\Lin}[2]{\mathcal L(#1,#2)}
\providecommand{\dual}[1]{#1'}
\providecommand{\ran}{\operatorname{ran}}
\providecommand{\cl}{\operatorname{cl}}
\providecommand{\clspan}{\overline{\operatorname{span}}}
\providecommand{\Spec}{\sigma}
\providecommand{\source}[1]{\par\smallskip{\footnotesize\textit{Source comparisons:} #1}\par\smallskip}


\chapter{Normed Spaces and Canonical Examples}\label{fa2:p1:c01}

Throughout these notes, $\K$ is $\R$ or $\C$. A subspace is a
$\K$-linear subspace unless a real underlying space is explicitly specified.
Completeness always refers to the displayed norm. We use elementary algebra,
completeness of the scalar field, and compactness of closed bounded subsets of
$\R^n$ (and hence $\C^n=\R^{2n}$).
\source{B05 Secs. 1.1--1.4 and Appendix C.}
\cite{B05} 


We use scalar completeness and the finite-dimensional Heine--Borel and extreme-value theorems in their ordinary real or complex coordinate forms.

% FAMILY D16.C01.01
\section{The metric carried by a norm}
\begin{definition}[Norm and seminorm]\label{fa2:p1:c01:norm}
A norm on a vector space $E$ is a function $x\mapsto\norm{x}$ with
$\norm{x}\geq0$, $\norm{x}=0$ exactly when $x=0$,
$\norm{a x}=\abs{a}\norm{x}$, and
$\norm{x+y}\leq\norm{x}+\norm{y}$. A seminorm satisfies the same conditions
except that nonzero vectors may have value zero. A normed space carries the
metric $d(x,y)=\norm{x-y}$ and its metric topology, the norm topology.
Write $B(x,r)=\{y:\norm{y-x}<r\}$ and
$\overline B(x,r)=\{y:\norm{y-x}\leq r\}$ for $r>0$.
\end{definition}

\begin{proposition}[Basic estimates and continuity]\label{fa2:p1:c01:operations}
The induced distance is a metric. The norm is $1$-Lipschitz, addition and
scalar multiplication are continuous, and translations and multiplication
by a nonzero scalar are homeomorphisms. The closure of a linear subspace is
a linear subspace. Every norm ball is convex, and a ball centered at zero
is balanced: $\abs{a}\leq1$ implies $aB\subset B$.
\end{proposition}
\begin{proof}
Positivity and symmetry of the metric follow from the norm axioms, and its
triangle inequality is the norm triangle inequality applied to
$x-z=(x-y)+(y-z)$. Applying the triangle inequality in both orders gives
$\abs{\norm{x}-\norm{y}}\leq\norm{x-y}$. Moreover,
\[
 \norm{(x+y)-(u+v)}\leq\norm{x-u}+\norm{y-v},\qquad
 \norm{a x-bu}\leq\abs{a}\norm{x-u}+\abs{a-b}\norm{u}.
\]
Near a fixed scalar $b$, $\abs{a}$ is bounded, so the latter estimate proves
joint continuity. The indicated bijections have inverses of the same form.
If $x_n,y_n$ lie in a subspace and converge to $x,y$, then
$a x_n+y_n\to ax+y$, proving the assertion about closure. Finally, for
$0\leq t\leq1$ the estimate
$\norm{t x+(1-t)y-z}\leq t\norm{x-z}+(1-t)\norm{y-z}$
proves convexity, and homogeneity proves balance.
\end{proof}

A seminorm gives a pseudometric rather than necessarily a metric. Its null
space $N=\{x:p(x)=0\}$ is linear, and $p(x+n)=p(x)$ for $n\in N$ by
the triangle inequality in both directions. Consequently $p$ defines a
genuine norm on the algebraic quotient $E/N$. This construction differs
from placing a distance-to-subspace seminorm on a previously normed space;
the latter is treated in Chapter~\ref{fa2:p1:c02}.

% FAMILY D16.C01.02
\section{Completeness and series}
\begin{definition}[Banach space and absolute summability]\label{fa2:p1:c01:banach}
A normed space is a Banach space if every Cauchy sequence converges in it.
A series $\sum_{n=1}^{\infty}x_n$ converges if its partial sums converge in
norm; it is absolutely summable if $\sum_n\norm{x_n}<\infty$.
\end{definition}

\begin{theorem}[Series criterion for completeness]\label{fa2:p1:c01:series-completeness}
A normed space $E$ is Banach if and only if every absolutely summable
series in $E$ converges. In a Banach space such a sum satisfies
$\norm{\sum_n x_n}\leq\sum_n\norm{x_n}$.
\end{theorem}
\begin{proof}
For $m>n$, the difference of partial sums has norm at most
$\sum_{j=n+1}^m\norm{x_j}$; thus partial sums are Cauchy and converge if
$E$ is complete. Passing to the limit in their finite triangle inequalities
gives the norm estimate. Conversely, let $(u_n)$ be Cauchy. Choose strictly
increasing indices $n_k$ so that
$\norm{u_m-u_n}<2^{-k}$ whenever $m,n\geq n_k$. The series
$\sum_{k\geq1}(u_{n_{k+1}}-u_{n_k})$ is absolutely summable, so its
partial sums $u_{n_{k+1}}-u_{n_1}$ converge. Hence $(u_n)$ has a convergent
subsequence. Given $\eps>0$, choose a Cauchy threshold and a subsequence
term beyond that threshold within $\eps/2$ of the subsequence limit; the
triangle inequality then puts every sufficiently late $u_n$ within
$\eps$ of that limit. Thus the whole Cauchy sequence converges.
\end{proof}

% FAMILY D16.C01.03
\section{Finite dimension}
\begin{theorem}[Finite-dimensional norm equivalence]\label{fa2:p1:c01:finite-dimension}
Let $E$ have finite dimension $n$ over $\K$. All norms on $E$ are
equivalent; $E$ is complete, its closed bounded sets are compact, and it
is locally compact. Every finite-dimensional subspace of any normed space
is closed. These statements include $n=0$.
\end{theorem}
\begin{proof}
When $n=0$, $E=\{0\}$, and all claims follow immediately. Otherwise choose
a basis $e_1,\ldots,e_n$ and write $A a=\sum_j a_j e_j$. Then
\[
 \norm{Aa}\leq\sum_j\abs{a_j}\norm{e_j}\leq C\abs{a}_2,
 \quad C=\Big(\sum_j\norm{e_j}^2\Big)^{1/2}.
\]
Thus $a\mapsto\norm{Aa}$ is continuous in the Euclidean topology. On the
compact Euclidean unit sphere it has a positive minimum $c$: if the
minimum were zero, its attaining point would contradict injectivity of
$A$. Homogeneity yields
$c\abs{a}_2\leq\norm{Aa}\leq C\abs{a}_2$ for all $a$. Repeating this
comparison for another norm proves equivalence. A Cauchy sequence in $E$
has Cauchy coordinates, whose scalar limits give its limit in $E$.
The inverse image of a closed bounded set is Euclidean closed and bounded,
thus compact; continuity of $A$ transfers compactness back. A closed unit
ball is consequently a compact neighborhood of zero, and translations
give local compactness. Finally, a convergent sequence in a finite-dimensional
subspace is Cauchy there and has a limit in that subspace; uniqueness of
the ambient metric limit proves closedness.
\end{proof}

The proof also shows that every algebraic linear map with finite-dimensional
domain is continuous: its coordinate formula is bounded above by a finite
sum and the coordinate norm is bounded by $c^{-1}\norm{x}$.

% FAMILY D16.C01.04
\section{Sequence spaces and continuous functions}
\begin{definition}[Canonical sequence spaces]\label{fa2:p1:c01:sequence-spaces}
For $1\leq p<\infty$, let
$\ell^p=\{x=(x_j)_{j\geq1}:\sum_j\abs{x_j}^p<\infty\}$ with
$\norm{x}_p=(\sum_j\abs{x_j}^p)^{1/p}$. Let $\ell^\infty$ be the bounded
scalar sequences with $\norm{x}_\infty=\sup_j\abs{x_j}$, and let $c_0$
be its subspace of sequences tending to zero, with the inherited norm.
The space $c_{00}$ consists of finitely supported sequences.
\end{definition}

\begin{lemma}[Discrete Holder and Minkowski inequalities]\label{fa2:p1:c01:holder}
If $1\leq p,q\leq\infty$ and $1/p+1/q=1$, with $1/\infty=0$, then
$\sum_j\abs{x_jy_j}\leq\norm{x}_p\norm{y}_q$ for every
$x\in\ell^p$ and $y\in\ell^q$.
For $1\leq p\leq\infty$ and every $x,y\in\ell^p$,
$\norm{x+y}_p\leq\norm{x}_p+\norm{y}_p$.
\end{lemma}
\begin{proof}
For $1<p<\infty$, Young's inequality $a b\leq a^p/p+b^q/q$
follows by minimizing $a^p/p-a b+b^q/q$ as a function of $a\geq0$;
the minimum is zero at $a=b^{q-1}$. Normalize nonzero finite vectors by
their respective norms and sum this inequality to obtain finite Holder.
The endpoint inequalities follow by taking out a supremum. Monotone
limits of finite sums give the infinite versions. For finite vectors and
$1<p<\infty$, apply Holder to
\[
 \sum_j\abs{x_j+y_j}^p
 \leq\sum_j(\abs{x_j}+\abs{y_j})\abs{x_j+y_j}^{p-1}
 \leq (\norm{x}_p+\norm{y}_p)\norm{x+y}_p^{p-1}.
\]
If the last norm is nonzero, divide by its $(p-1)$st power; otherwise the
claim is immediate. The cases $p=1,\infty$ follow directly from the scalar
triangle inequality. Applying finite Minkowski to initial segments and
taking increasing limits proves both membership of $x+y$ and the claimed
infinite inequality. Homogeneity and definiteness are immediate, so these
expressions really are norms.
\end{proof}

\begin{theorem}[Completeness of the canonical examples]\label{fa2:p1:c01:examples-complete}
The spaces $\ell^p$ for $1\leq p\leq\infty$ and $c_0$ are Banach.
If $K$ is compact Hausdorff, then $C(K,\K)$ with
$\norm{f}_\infty=\sup_{t\in K}\abs{f(t)}$ is Banach; for $K=\varnothing$
this is the zero space and the norm is zero.
\end{theorem}
\begin{proof}
If $(x^{(n)})$ is Cauchy in $\ell^p$, each coordinate is Cauchy, because
$\abs{x^{(n)}_j-x^{(m)}_j}\leq\norm{x^{(n)}-x^{(m)}}_p$.
Let its coordinate limit be $x_j$. For $p<\infty$, given $\eps>0$ and
$n,m\geq N$ with norm difference at most $\eps$, first pass $m\to\infty$
in any finite coordinate sum to obtain
$\sum_{j\leq k}\abs{x^{(n)}_j-x_j}^p\leq\eps^p$, then let
$k\to\infty$. Thus $x-x^{(N)}\in\ell^p$, hence $x\in\ell^p$, and
$\norm{x^{(n)}-x}_p\leq\eps$ for $n\geq N$. For $p=\infty$ the same
coordinate limit gives $\sup_j\abs{x^{(n)}_j-x_j}\leq\eps$ directly.
The space $c_0$ is closed in $\ell^\infty$: if $\norm{x-y}_\infty<\eps/2$
and $y_j\to0$, then $\abs{x_j}<\eps$ for all sufficiently large $j$.
This also proves its completeness by the preceding construction.

A uniformly Cauchy sequence $(f_n)$ has a scalar pointwise limit $f$,
and letting one index tend to infinity in its uniform Cauchy estimate shows
uniform convergence. To verify continuity at $t$, choose $n$ with
$\norm{f-f_n}_\infty<\eps/3$ and a neighborhood $U$ of $t$ on which
$\abs{f_n(s)-f_n(t)}<\eps/3$. Then
$\abs{f(s)-f(t)}<\eps$ on $U$. Thus $f\in C(K)$ and convergence occurs
in its norm. Compactness ensures that each continuous scalar function is
bounded, so the displayed norm is finite to begin with.
\end{proof}

Truncations $\sum_{j\leq n}x_j e_j$ converge to $x$ in $\ell^p$ for
$p<\infty$ and in $c_0$, by the tail formulas for their norms. Hence
$c_{00}$ is dense in each of these spaces. In $\ell^\infty$ its closure is
exactly $c_0$; the constant sequence $1$ has distance $1$ from $c_{00}$.

% FAMILY D16.C01.05
\section{The obstruction in infinite dimension}
\begin{lemma}[Riesz's lemma]\label{fa2:p1:c01:riesz-lemma}
If $M$ is a proper closed subspace of a normed space $E$ and
$0<\theta<1$, there exists $x\in E$ with $\norm{x}=1$ and
$\dist(x,M)>\theta$.
\end{lemma}
\begin{proof}
Choose $y\notin M$. Closedness gives $d=\dist(y,M)>0$. By the defining
infimum, choose $m\in M$ with $\norm{y-m}<d/\theta$, and set
$x=(y-m)/\norm{y-m}$. Translation by $m$ and scaling a linear subspace
show that $\dist(x,M)=d/\norm{y-m}>\theta$.
\end{proof}

\begin{theorem}[Compact balls characterize finite dimension]\label{fa2:p1:c01:compact-ball}
A normed space has compact closed unit ball if and only if it is
finite dimensional. It is locally compact if and only if it is finite
dimensional.
\end{theorem}
\begin{proof}
The finite-dimensional direction was proved above. In an infinite-dimensional
space choose $\norm{x_1}=1$, and repeatedly apply Riesz's lemma with
$\theta=1/2$ to the proper finite-dimensional, hence closed, span of the
previous vectors. Then $\norm{x_n-x_m}>1/2$ for $n>m$, so this bounded
sequence has no Cauchy, and therefore no convergent, subsequence. A compact
metric set has convergent subsequences: finite covers by balls of radii
$2^{-k}$ allow successive infinite subsequence selections, whose diagonal
is Cauchy, and compactness supplies its limit. Thus the unit ball is not
compact. Finally, a compact neighborhood of zero contains $B(0,r)$ for
some $r>0$, and its closed subset $\overline B(0,r/2)$ is compact. Scaling
would make the closed unit ball compact, a contradiction.
\end{proof}

\begin{example}[A visible separated sequence]\label{fa2:p1:c01:unit-vectors}
The coordinate vectors $(e_n)$ have norm $1$ in every $\ell^p$ and in
$c_0$. For $m\ne n$, their distance is $2^{1/p}$ when $p<\infty$ and
$1$ in the supremum norm. They supply a concrete bounded sequence with
no norm-convergent subsequence. The space $c_{00}$ with any of these norms
is incomplete: the finite truncations of $(2^{-j})_{j\geq1}$ are Cauchy
but have no limit in $c_{00}$, since norm convergence forces the stated
infinite-support coordinate limit.
\end{example}


\chapter{Subspaces, Quotients, Completions, and Sums}\label{fa2:p1:c02}
\source{B05 Secs. 1.7--1.8; B01 Sec. 11.2.}
\cite{B05} \cite{B01} 


% FAMILY D16.C02.01
\section{Inherited completeness}
\begin{proposition}[Complete subspaces]\label{fa2:p1:c02:complete-subspace}
A complete linear subspace $M$ of a normed space $E$ is closed in $E$.
If $E$ is Banach, a linear subspace $M$ is complete with its inherited
norm if and only if it is closed.
\end{proposition}
\begin{proof}
If $x\in\cl M$, choose $m_n\in M$ with $\norm{m_n-x}<1/n$.
This sequence is Cauchy and, when $M$ is complete, has a limit in $M$;
uniqueness of the ambient limit puts $x$ in $M$. Conversely, a Cauchy
sequence in a closed subspace of a Banach space converges in the ambient
space, and closedness keeps its limit in the subspace.
\end{proof}

% FAMILY D16.C02.02
\section{The quotient norm}
\begin{definition}[Distance quotient]\label{fa2:p1:c02:quotient-norm}
Let $M\subset E$ be a linear subspace of a normed space. On the algebraic
quotient $E/M$ put
\[
 \norm{x+M}_{E/M}=\inf_{m\in M}\norm{x-m}=\dist(x,M),
 \qquad q(x)=x+M.
\]
When $M$ is closed this is called the quotient norm.
\end{definition}

\begin{proposition}[Norm, topology, and factorization]\label{fa2:p1:c02:quotient-map}
The formula is a well-defined seminorm, and is a norm exactly when $M$
is closed. In that case $q$ is linear, surjective, contractive, and open;
more precisely $q(B_E(x,r))=B_{E/M}(q(x),r)$. A linear map $T:E\to F$
vanishing on $M$ factors uniquely as $\widetilde Tq$, and if $T$ is
bounded then $\widetilde T$ is bounded with $\norm{\widetilde T}=\norm{T}$.
\end{proposition}
\begin{proof}
Changing $x$ by an element of $M$ changes none of the candidate distances.
For two cosets choose representatives within $\eps$ of their infimum
norms; their sum represents the sum coset, giving the triangle inequality
after $\eps\downarrow0$. For a nonzero scalar, substitute $m=a n$ in
the infimum to get homogeneity; the zero scalar case is immediate.
The null cosets are exactly those represented by $x\in\cl M$, which
proves the assertion about definiteness. The inequality
$\norm{q(x)}\leq\norm{x}$ proves contractivity. If a coset $z$ has
$\norm{z-q(x)}<r$, the infimum formula gives a representative of
$z-q(x)$ with norm less than $r$, proving the displayed ball equality;
the reverse inclusion follows from contractivity. This equality proves
openness. Algebraic factorization is forced by
$\widetilde T(x+M)=Tx$, which is well-defined because $T(M)=0$.
For each $m\in M$,
$\norm{Tx}\leq\norm{T}\norm{x-m}$, so taking an infimum gives
$\norm{\widetilde T}\leq\norm{T}$. Composition with the contractive
$q$ gives the other inequality.
\end{proof}

For nonclosed $M$, the seminormed quotient has null space $\cl M/M$;
its Hausdorff quotient is isometric to $E/\cl M$. Thus a dense proper
subspace gives a nontrivial algebraic quotient with identically zero
seminorm, rather than a normed space.

% FAMILY D16.C02.03
\section{Banach quotients and completion}
\begin{theorem}[Completeness of a closed quotient]\label{fa2:p1:c02:quotient-complete}
If $E$ is Banach and $M$ is closed, then $E/M$ with the quotient norm is
Banach.
\end{theorem}
\begin{proof}
Let $(z_n)$ be Cauchy in $E/M$. Select an increasing subsequence
$(z_{n_k})$ such that $\norm{z_{n_{k+1}}-z_{n_k}}<2^{-k-1}$.
Choose a representative $x_1$ of $z_{n_1}$. At step $k$, the infimum
formula provides $u_k\in E$ with
$q(u_k)=z_{n_{k+1}}-z_{n_k}$ and $\norm{u_k}<2^{-k}$.
Define $x_{k+1}=x_k+u_k$. Induction gives $q(x_k)=z_{n_k}$, and
\[
 \norm{x_l-x_k}\leq\sum_{j=k}^{l-1}2^{-j}\quad(l>k).
\]
Hence $(x_k)$ converges to $x\in E$. Contractivity of $q$ implies
$z_{n_k}\to q(x)$. As a Cauchy sequence with a convergent subsequence,
$(z_n)$ converges to the same coset. Closedness of $M$ was used to have
a genuine norm; no choice of distance-minimizing representatives was
assumed.
\end{proof}

\begin{theorem}[Isometric completion and its universal property]\label{fa2:p1:c02:completion}
Every normed space $E$ has a Banach space $\widehat E$ and a linear
isometry $i:E\to\widehat E$ with dense image. If $F$ is Banach, every
bounded linear map $T:E\to F$ has a unique bounded extension
$\widehat T:\widehat E\to F$ satisfying $\widehat Ti=T$ and
$\norm{\widehat T}=\norm{T}$. Two completions are related by a unique
surjective linear isometry agreeing with their embeddings of $E$.
\end{theorem}
\begin{proof}
Take the vector space of Cauchy sequences $(x_n)$ in $E$ and identify
two sequences when $\norm{x_n-y_n}\to0$. Sums and scalar multiples
preserve Cauchy sequences and this equivalence. Put
$\norm{[(x_n)]}=\lim_n\norm{x_n}$; the real limit exists by the reverse
triangle inequality. The same inequality makes the value independent of
the representative, and the norm axioms pass to the limit. The value
zero is precisely the equivalence class of the zero sequence, so this is
a norm. Constant sequences define $i$, a linear isometry. If
$z=[(x_n)]$, then
$\norm{z-i(x_N)}=\lim_n\norm{x_n-x_N}\to0$ as $N\to\infty$;
thus $i(E)$ is dense.

To prove completeness without assuming it, let $(z_k)$ be Cauchy in the
new space. By density choose $a_k\in E$ with
$\norm{i(a_k)-z_k}<1/k$. The sequence $(a_k)$ is Cauchy in $E$ by
the triangle inequality. Set $z=[(a_k)]$. The density calculation just
given shows $i(a_k)\to z$, and then $z_k\to z$. This proves completeness.
For $z=[(x_n)]$, define $\widehat T z=\lim_n Tx_n$. The estimate
$\norm{Tx_n-Tx_m}\leq\norm{T}\norm{x_n-x_m}$ and completeness of $F$
give existence; the same estimate for equivalent representatives gives
independence. Passing to limits proves linearity and
$\norm{\widehat T z}\leq\norm{T}\norm{z}$. Restricting to $i(E)$ gives
equality of operator norms. Two continuous extensions agreeing on a dense
set agree everywhere. Applying this extension construction to the two
isometric dense embeddings produces maps in both directions; uniqueness
makes their compositions identities and the maps isometries.
\end{proof}

% FAMILY D16.C02.04
\section{Products and Hilbert sums}
\begin{proposition}[Finite products]\label{fa2:p1:c02:products}
For normed spaces $E_1,\ldots,E_n$, the product carries the norms
\[
 \norm{x}_p=\Big(\sum_{j=1}^n\norm{x_j}^p\Big)^{1/p}\quad(1\leq p<\infty),
 \qquad\norm{x}_\infty=\max_j\norm{x_j}.
\]
They give the same topology, coordinate projections and inclusions are
bounded, and the product is Banach exactly when every factor is Banach.
The empty product is the zero space.
\end{proposition}
\begin{proof}
Discrete Minkowski proves the triangle inequality. The estimates
$\norm{x}_\infty\leq\norm{x}_p\leq n^{1/p}\norm{x}_\infty$ prove
equivalence, and each coordinate map has norm at most $1$. A Cauchy
sequence is coordinatewise Cauchy and converges in the maximum norm when
the factors are complete. Conversely each factor is the closed coordinate
subspace of the product and inherits its norm, so is complete when the
product is complete.
\end{proof}

\begin{definition}[Hilbert direct sum over an arbitrary index set]\label{fa2:p1:c02:hilbert-sum}
Let $(H_j)_{j\in I}$ be Hilbert spaces over $\K$, with inner products
conjugate-linear in the first variable. Define
\[
 \bigoplus_{j\in I}^{\!2}H_j
 =\left\{(x_j):\sup_{A\subset I\text{ finite}}
                 \sum_{j\in A}\norm{x_j}^2<\infty\right\},\qquad
 \norm{x}^2=\sum_{j\in I}\norm{x_j}^2.
\]
Nonnegative sums here mean suprema over finite subsets. The inner product
is $\ip{x}{y}=\sum_{j\in I}\ip{x_j}{y_j}$.
\end{definition}

\begin{proposition}[Completeness and finite-support density]\label{fa2:p1:c02:hilbert-sum-complete}
This direct sum is Hilbert, each coordinate projection and inclusion is
contractive, and finite-support vectors are dense. Every vector has at
most countably many nonzero coordinates; the index set itself need not
be countable.
\end{proposition}
\begin{proof}
For each positive integer $m$, only finitely many coordinates can have
$\norm{x_j}\geq1/m$, since arbitrary finite sums of their squares are
bounded. Their countable union contains all nonzero coordinates. Moreover,
for any $\eps>0$ a finite set can be chosen whose squared sum is within
$\eps^2$ of the supremum; the remaining squared sum is at most
$\eps^2$. This proves finite-support density. Cauchy--Schwarz in each
$H_j$ and discrete Holder give
\[
 \sum_j\abs{\ip{x_j}{y_j}}\leq\sum_j\norm{x_j}\norm{y_j}
 \leq\norm{x}\norm{y}.
\]
Thus the scalar sum is absolutely convergent, independent of enumeration,
and defines an inner product with the stated norm and convention. Its
algebraic identities follow first on finite sums and then by absolute
convergence.

If $(x^{(n)})$ is Cauchy in the sum, its $j$th coordinate has a limit
$x_j\in H_j$. For every finite $A$, the inequalities
$\sum_{j\in A}\norm{x_j^{(n)}-x_j^{(m)}}^2\leq\eps^2$ for
$n,m\geq N$ pass to the limit $m\to\infty$ and then to the supremum
over $A$. Consequently $x-x^{(N)}$ has finite squared sum, $x$ belongs
to the direct sum, and $\norm{x^{(n)}-x}\leq\eps$ for $n\geq N$.
Coordinate contractivity and the isometry of each inclusion follow from
the norm formula. This proves all assertions, including the empty-index
case. The elementary inner-product Cauchy--Schwarz estimate used here
also follows by expanding $\norm{u-a v}^2\geq0$ and choosing
$a=\ip{v}{u}/\norm{v}^2$ for $v\ne0$.
\end{proof}

% FAMILY D16.C02.05
\section{Bounded projections and complements}
\begin{definition}[Topological direct sum]\label{fa2:p1:c02:topological-sum}
Subspaces $M,N\subset E$ give a topological direct sum $E=M\oplus N$
if each $x\in E$ has a unique decomposition $x=m+n$ and the coordinate
maps $x\mapsto m$, $x\mapsto n$ are bounded. A subspace with such a
complement is called complemented. A bounded projection is a bounded
linear map $P:E\to E$ satisfying $P^2=P$.
\end{definition}

\begin{proposition}[Projections and elementary complements]\label{fa2:p1:c02:projections}
A bounded projection gives the topological direct sum
$E=\ran P\oplus\ker P$, and both subspaces are closed. Conversely a
topological direct sum gives a bounded projection onto either summand.
Every finite-dimensional subspace of a normed space is complemented.
Every closed finite-codimensional subspace of a normed space is complemented.
\end{proposition}
\begin{proof}
The identity $x=Px+(x-Px)$ has its second term in $\ker P$, and the
intersection of range and kernel is zero. The range is
$\ker(\id-P)$, so both summands are closed. The coordinate maps are
$P$ and $\id-P$, hence bounded. The converse follows directly by taking
the first coordinate map.

For the finite-dimensional assertion we use the norm-preserving extension
theorem proved in Chapter~\ref{fa2:p1:c05}; none of the quotient or completion
arguments above uses this assertion. Choose a basis $e_1,\ldots,e_n$ of
$M$. Its coordinate functionals are bounded by the finite-dimensional
norm comparison. Extend each to $f_j\in E'$ by Hahn--Banach and put
$Px=\sum_j f_j(x)e_j$. This is bounded and equals the identity on $M$,
so is a projection onto $M$. If $M=\{0\}$ use $P=0$.
For a closed finite-codimensional $M$, $E/M$ is a finite-dimensional
normed space. Choose representatives of a basis to obtain a linear
section $s:E/M\to E$ of $q$. Finite-dimensional norm comparison makes
$s$ bounded, so $s q$ is a bounded projection with kernel $M$; its range
is the required finite-dimensional complement.
\end{proof}

Merely having an algebraic direct sum gives no estimate for its coordinate
maps. If $E$ is Banach, two closed complementary subspaces do give a
topological direct sum: the addition map $M\times N\to E$ is a bounded
bijection of Banach spaces, so Chapter~\ref{fa2:p1:c08} makes its inverse
bounded. On an incomplete space that implication requires a separate
argument and is not part of the definition.


\chapter{Bounded Operators and Perturbation of Inverses}\label{fa2:p1:c03}
\source{B05 Sec. 1.4; B06 Secs. 2.1 and 2.4.}
\cite{B05} \cite{B06} 


% FAMILY D16.C03.01
\section{From continuity to an estimate}
\begin{theorem}[Boundedness criterion for a linear map]\label{fa2:p1:c03:bounded-criterion}
For a linear map $T:E\to F$ between normed spaces, the following are
equivalent: continuity everywhere; continuity at zero; existence of
$C\geq0$ with $\norm{Tx}\leq C\norm{x}$ for every $x\in E$; and
boundedness of $T$ on the closed unit ball. They are also equivalent to
mapping bounded sets to bounded sets.
\end{theorem}
\begin{proof}
Continuity everywhere implies continuity at zero. Continuity at zero
gives $\delta>0$ such that $\norm{Tu}<1$ for $\norm{u}<\delta$.
For $x\ne0$, apply this to $u=\delta x/(2\norm{x})$ to get
$\norm{Tx}\leq2\norm{x}/\delta$; the zero case follows by linearity.
Such an estimate gives
$\norm{Tx-Ty}\leq C\norm{x-y}$, hence continuity everywhere, and
immediately sends bounded sets to bounded sets. Conversely the closed
unit ball is bounded, so bounded-set preservation bounds its image.
If that image has norms at most $C$, scale each nonzero vector to unit
norm to obtain the global estimate.
\end{proof}

% FAMILY D16.C03.02
\begin{definition}[Operator norm]\label{fa2:p1:c03:operator-norm}
Write $\Lin{E}{F}$ for the vector space of bounded linear maps $E\to F$,
with operator norm
\[
 \norm{T}=\sup_{\norm{x}\leq1}\norm{Tx}
 =\inf\{C\geq0:\norm{Tx}\leq C\norm{x}\text{ for all }x\}.
\]
The supremum is over a nonempty set even when $E=\{0\}$; then it is
zero. We write $\Lin{E}{E}=\mathcal L(E)$.
\end{definition}

\begin{proposition}[Norm estimates and finite-dimensional coordinates]\label{fa2:p1:c03:operator-estimates}
The operator norm is a norm, $\norm{Tx}\leq\norm{T}\norm{x}$, and
$\norm{ST}\leq\norm{S}\norm{T}$ whenever the composition is defined.
Restriction to a subspace cannot increase the norm. Every linear map on
a finite-dimensional domain is bounded; after bases are chosen its matrix
acts by the usual coordinate multiplication, and changing the norms or
bases gives equivalent operator norms in finite dimensions.
\end{proposition}
\begin{proof}
The bound on each vector follows by scaling, with $x=0$ handled
separately. It shows the equivalence of the supremum and infimum formulas.
Homogeneity and the triangle inequality follow by taking suprema of their
pointwise counterparts. If $\norm{T}=0$, every $Tx=0$, so the norm is
definite. For a composition, apply the vector estimate first to $S$ and
then to $T$ and take the supremum on the unit ball. Restriction uses a
smaller unit ball. If $x=\sum_j a_je_j$, the estimate
$\norm{Tx}\leq\sum_j\abs{a_j}\norm{Te_j}$ and Chapter~\ref{fa2:p1:c01}'s
lower coordinate bound prove boundedness. Matrix multiplication describes
the coefficients of this sum. Finally any two domain and codomain norms
are bounded above and below by fixed positive multiples of one another;
substitution in the defining supremum gives the analogous two operator
norm comparisons.
\end{proof}

In particular, $\norm{\id_E}=1$ when $E\ne\{0\}$ and is $0$ on the
zero space. A bounded functional need not attain the supremum in the
definition of its norm.

% FAMILY D16.C03.03
\section{Completeness of operator spaces}
\begin{theorem}[A complete target suffices]\label{fa2:p1:c03:operator-complete}
If $F$ is Banach and $E$ is any normed space, then $\Lin{E}{F}$ is
Banach in the operator norm.
\end{theorem}
\begin{proof}
Let $(T_n)$ be operator-norm Cauchy. For each $x$, the estimate
$\norm{T_nx-T_mx}\leq\norm{T_n-T_m}\norm{x}$ makes $(T_nx)$
Cauchy in $F$. Define $Tx=\lim_nT_nx$. Passing to the limit in sums
and scalar multiples proves linearity. The Cauchy sequence $(T_n)$ has
bounded norms, say at most $M$, so $\norm{Tx}\leq M\norm{x}$ and
$T$ is bounded. Given $\eps>0$, choose $N$ such that
$\norm{T_n-T_m}\leq\eps$ for $m,n\geq N$. Letting $m\to\infty$
in the pointwise estimate gives
$\norm{(T_n-T)x}\leq\eps\norm{x}$ for all $x$, and taking the
supremum gives $\norm{T_n-T}\leq\eps$ for $n\geq N$. This is
operator-norm convergence, with no completeness assumption on $E$.
\end{proof}

% FAMILY D16.C03.04
\section{Geometric inversion and the resolvent identity}
\begin{theorem}[Neumann inversion]\label{fa2:p1:c03:neumann}
Let $E$ be Banach and $T\in\mathcal L(E)$ satisfy $\norm{T}<1$.
Then $\id-T$ is invertible with bounded two-sided inverse
\[
 (\id-T)^{-1}=\sum_{n=0}^{\infty}T^n,\qquad
 \norm{(\id-T)^{-1}}\leq\frac1{1-\norm{T}}.
\]
For the partial sum through $n=N$, the norm error is at most
$\norm{T}^{N+1}/(1-\norm{T})$.
\end{theorem}
\begin{proof}
Submultiplicativity gives $\norm{T^n}\leq\norm{T}^n$ for $n\geq1$,
and $\norm{T^0}=\norm{\id}\leq1$. Completeness of $\mathcal L(E)$
and the series criterion give convergence to $S$ and the stated bounds.
The finite telescoping products satisfy
\[
 (\id-T)\sum_{n=0}^NT^n
 =\Big(\sum_{n=0}^NT^n\Big)(\id-T)=\id-T^{N+1}.
\]
Composition is continuous in the operator norm by its norm estimate;
letting $N\to\infty$ gives $(\id-T)S=S(\id-T)=\id$. The same proof
includes $E=\{0\}$, where its unique operator is also its identity.
\end{proof}

\begin{proposition}[Inverse difference and resolvent identities]\label{fa2:p1:c03:resolvent-identity}
For bounded invertible maps $A,B:E\to F$,
\[
 A^{-1}-B^{-1}=A^{-1}(B-A)B^{-1}.
\]
For $T\in\mathcal L(E)$ define $R(\lambda,T)=(\lambda\id-T)^{-1}$
when this bounded inverse exists. For any two such scalars,
\[
 R(\lambda,T)-R(\mu,T)
 =(\mu-\lambda)R(\lambda,T)R(\mu,T).
\]
The two resolvents commute.
\end{proposition}
\begin{proof}
Expanding the first right side gives $A^{-1}-B^{-1}$.
Apply it to $A=\lambda\id-T$ and $B=\mu\id-T$ for the second
formula. These two operators commute, and an invertible operator's inverse
commutes with every operator commuting with it: multiply $AC=CA$ on
both sides by $A^{-1}$. Applying this observation twice proves the final
claim. The sign is tied to the convention $\lambda\id-T$.
\end{proof}

% FAMILY D16.C03.05
\section{Stability of an inverse}
\begin{theorem}[Quantitative perturbation theorem]\label{fa2:p1:c03:perturbation}
Let $E,F$ be Banach, and let $A:E\to F$ be a bounded bijection whose
inverse is bounded. If $H:E\to F$ is bounded and
$\eta=\norm{A^{-1}H}<1$, then $B=A+H$ is boundedly invertible and
\[
 B^{-1}=\sum_{n=0}^{\infty}(-A^{-1}H)^nA^{-1},\qquad
 \norm{B^{-1}}\leq\frac{\norm{A^{-1}}}{1-\eta},\qquad
 \norm{B^{-1}-A^{-1}}\leq
 \frac{\eta\norm{A^{-1}}}{1-\eta}.
\]
Consequently the set of boundedly invertible maps in $\Lin{E}{F}$ is
open and inversion is norm-continuous there.
\end{theorem}
\begin{proof}
Factor $B=A(\id+A^{-1}H)$ and apply Neumann inversion on $E$.
This gives the series and the first estimate. Removing its term $n=0$
and summing the tail gives the second estimate. Since
$\eta\leq\norm{A^{-1}}\norm{H}$, an operator-norm ball around $A$
of sufficiently small positive radius satisfies the hypothesis. The final
estimate tends to zero as $\norm{H}\to0$, proving continuity.
\end{proof}

This theorem assumes the bounded inverse at its starting point. The
assertion that a bounded bijection of Banach spaces automatically has a
bounded inverse is a separate consequence of Baire category, proved in
Chapter~\ref{fa2:p1:c08}.


\chapter{Continuous Duals and Concrete Representations}\label{fa2:p1:c04}
\source{B05 Secs. 1.10--1.11 and Appendix C; B01 Sec. 4.3.}
\cite{B05} \cite{B01} 


% FAMILY D16.C04.01
\section{The continuous dual}
\begin{definition}[Dual and dual norm]\label{fa2:p1:c04:dual}
For a normed space $E$, its continuous linear dual is
$E'=\Lin{E}{\K}$, with
$\norm{f}=\sup_{\norm{x}\leq1}\abs{f(x)}$. Its algebraic dual
$E^*_{\mathrm{alg}}$ is the space of all linear scalar-valued maps, with
no continuity condition. The bidual is $E''=(E')'$.
\end{definition}

\begin{proposition}[Completeness and evaluation]\label{fa2:p1:c04:dual-complete}
The dual of every normed space is Banach, whether or not the original
space is complete. For $f\in E'$ and $x\in E$,
$\abs{f(x)}\leq\norm{f}\norm{x}$.
\end{proposition}
\begin{proof}
The scalar field is complete, so the operator-space completeness theorem
of Chapter~\ref{fa2:p1:c03} applies with target $\K$. The estimate is the
operator-norm estimate applied to a scalar-valued map.
\end{proof}

In infinite dimension the algebraic and continuous duals can differ.
For example choose linearly independent unit vectors $u_n$ inductively,
extend them to a Hamel basis by the maximality principle, and prescribe
$f(u_n)=n$, with arbitrary values, say zero, on the remaining basis
vectors. Here the maximality argument orders independent sets containing
all $u_n$ by inclusion: a chain union remains independent because every
finite relation occurs in one member, and a maximal such set spans the
space because an outside vector could otherwise be adjoined.
Extension by finite linear combinations defines an algebraic
functional which is unbounded on the unit sphere, hence discontinuous.
This algebraic use of a basis supplies no bounded extension theorem.

% FAMILY D16.C04.02
\begin{definition}[Canonical evaluation map]\label{fa2:p1:c04:canonical-map}
The canonical map $J_E:E\to E''$ is
$J_E(x)(f)=f(x)$ for $f\in E'$. A normed space is called reflexive if
$J_E$ is surjective.
\end{definition}

\begin{proposition}[Evaluation is initially contractive]\label{fa2:p1:c04:evaluation-contractive}
For each $x$, $J_E(x)$ is a continuous linear functional on $E'$.
The map $J_E$ is linear and $\norm{J_E x}\leq\norm{x}$.
\end{proposition}
\begin{proof}
Linearity in $f$ follows from pointwise addition and scalar multiplication
in $E'$, and the evaluation estimate gives boundedness with norm at most
$\norm{x}$. Linearity in $x$ follows from linearity of each $f$.
\end{proof}

No equality of norms or injectivity has been used here. Their proof in
Chapter~\ref{fa2:p1:c05} requires norming functionals. Thus the definition of
reflexivity does not silently build in an isometry theorem.

% FAMILY D16.C04.03
\section{Annihilators and transposes}
\begin{definition}[Annihilator and preannihilator]\label{fa2:p1:c04:annihilators}
For $A\subset E$ and $B\subset E'$ put
\[
 A^\perp=\{f\in E':f(a)=0\ (a\in A)\},\qquad
 {}^\perp B=\{x\in E:f(x)=0\ (f\in B)\}.
\]
The notation denotes a dual pairing, not a Hilbert inner product.
For $T\in\Lin{E}{F}$ its Banach transpose is
$T':F'\to E'$, $T'f=f\circ T$.
\end{definition}

\begin{proposition}[Basic dual maps]\label{fa2:p1:c04:dual-maps}
Annihilators and preannihilators are closed linear subspaces; they depend
only on the closed linear span of the set annihilated. Inclusion of sets
reverses inclusion of annihilators. Restriction $E'\to M'$ is contractive
for every subspace $M$. If $M$ is closed, the transpose of the quotient
map identifies $(E/M)'$ isometrically with $M^\perp$. The transpose is
bounded, $\norm{T'}\leq\norm{T}$, and
$(ST)'=T'S'$.
\end{proposition}
\begin{proof}
Each annihilator is an intersection of kernels of continuous linear
maps: evaluation is continuous on $E'$ by the preceding estimate, and
the members of $E'$ are continuous on $E$. Vanishing passes to linear
combinations and limits, proving the closed-span assertion and the
remaining elementary inclusions. The restriction estimate follows by
taking a supremum over the smaller unit ball. If $f\in M^\perp$, the
factorization theorem of Chapter~\ref{fa2:p1:c02} gives a unique
$\widetilde f\in(E/M)'$ with $f=\widetilde f q$ and
$\norm{\widetilde f}=\norm{f}$; this proves the quotient identification
without Hahn--Banach. Finally,
$\abs{f(Tx)}\leq\norm{f}\norm{T}\norm{x}$ proves the transpose bound,
and associativity of composition proves its product rule.
\end{proof}

Surjectivity of restriction, norm-preserving extensions from $M$, and
the equality $\norm{T'}=\norm{T}$ are established only after the
extension theorem in Chapter~\ref{fa2:p1:c05}.

\section{Concrete dual representations}
The sequence, Lebesgue and continuous-function representations are proved in Part IV, Theorems~\ref{fa2:p4:sequence-duals}, \ref{fa2:p4:lp-duality} and \ref{fa2:p4:riesz-markov}. Their pairings are linear in the vector acted on. Neither the algebraic dual nor a space of unquotiented pointwise representatives is substituted for the continuous dual.


\chapter{Hahn--Banach Extension and Norming Functionals}\label{fa2:p1:c05}
\source{B01 Secs. 1.1 and 11.4; B05 Secs. 1.9--1.11.}
\cite{B01} \cite{B05} 


% FAMILY D16.C05.01
\section{Real sublinear domination}
\begin{definition}[Real sublinear function]\label{fa2:p1:c05:sublinear}
On a real vector space, a finite real-valued function $p$ is sublinear
if $p(x+y)\leq p(x)+p(y)$ and $p(t x)=t p(x)$ for all $t\geq0$.
In particular $p(0)=0$ and $-p(-x)\leq p(x)$. Neither positivity nor
symmetry is required.
\end{definition}

\begin{theorem}[Real analytic Hahn--Banach]\label{fa2:p1:c05:real-hahn-banach}
Let $E$ be a real vector space, $M\subset E$ a linear subspace, and
$p:E\to\R$ sublinear. If $f:M\to\R$ is linear and $f(m)\leq p(m)$
for all $m\in M$, then there exists a real linear $F:E\to\R$ extending
$f$ and satisfying $F(x)\leq p(x)$ for all $x\in E$.
\end{theorem}
\begin{proof}
We use Zorn's maximality principle. First prove the one-dimensional
extension step. If $z\notin M$, prescribe
$h(m+t z)=f(m)+t a$. This is a well-defined linear extension because
the decomposition is unique. For it to be dominated by $p$, the
necessary and sufficient conditions on $a$ are
\[
 \sup_{m\in M}\{f(m)-p(m-z)\}\leq a
 \leq\inf_{n\in M}\{p(n+z)-f(n)\}.
\]
Indeed the upper inequality gives domination for $t>0$ after dividing
$m+t z$ by $t$; the lower inequality gives it for $t<0$ after dividing
by $-t$; $t=0$ is the original hypothesis. For every $m,n\in M$,
\[
 f(m)+f(n)=f(m+n)\leq p(m+n)
 \leq p(m-z)+p(n+z),
\]
so every lower candidate is at most every upper candidate. The lower
candidates include the finite number $-p(-z)$ and are bounded above by
$p(z)$ by taking $n=0$. The upper candidates are bounded below by
$-p(-z)$ and include $p(z)$. Their supremum and infimum are therefore
finite and ordered, so a real $a$ in this interval exists.

Order all dominated extensions of $f$ by extension of domain and values.
This set is nonempty. The union of a chain of domains is a subspace:
any two of its vectors occur together in one member of the chain. The
union functional is well-defined, linear, and dominated by $p$, hence
is an upper bound for that chain. Zorn's principle gives a maximal
extension. If its domain were proper, the one-dimensional construction
would extend it further, a contradiction. Its domain is consequently $E$.
\end{proof}

% FAMILY D16.C05.02
\section{Norm-preserving extension over both scalar fields}
\begin{theorem}[Real and complex norm-preserving Hahn--Banach]\label{fa2:p1:c05:norm-extension}
Let $E$ be a normed space over $\K$, $M\subset E$ any linear subspace,
and $f\in M'$. There exists $F\in E'$ with $F|_M=f$ and
$\norm{F}=\norm{f}$. Neither completeness nor closedness of $M$ is
required.
\end{theorem}
\begin{proof}
For real scalars apply the real theorem to $p(x)=\norm{f}\norm{x}$.
The domination for $x$ and $-x$ yields
$\abs{F(x)}\leq\norm{f}\norm{x}$. Thus $\norm{F}\leq\norm{f}$,
and restriction gives the reverse inequality.

For complex scalars regard $E$ and $M$ as real vector spaces and let
$u=\operatorname{Re}f$. It is real linear and bounded by
$\norm{f}\norm{x}$. Extend it by the real theorem to $h:E\to\R$
with that domination, and define
\[
                       F(x)=h(x)-i h(i x).
\]
It is real linear and
$F(i x)=h(i x)+i h(x)=iF(x)$, so is complex linear. For $m\in M$,
$h(m)=\operatorname{Re}f(m)$ and
$h(i m)=\operatorname{Re}(i f(m))=-\operatorname{Im}f(m)$,
so $F(m)=f(m)$. To bound its modulus, if $F(x)\ne0$ choose
$\zeta=\overline{F(x)}/\abs{F(x)}$, a scalar of modulus $1$.
Then
\[
 \abs{F(x)}=\operatorname{Re}F(\zeta x)=h(\zeta x)
 \leq\norm{f}\norm{\zeta x}=\norm{f}\norm{x}.
\]
The same bound is immediate when $F(x)=0$. Again restriction gives
equality of norms. The case $f=0$ causes no exception: the bound forces
the constructed extension to be zero.
\end{proof}

The passage through real parts changes scalar linearity twice; simply
using the real extension itself would not produce a complex linear
functional. The explicit reconstruction is what closes that bridge.

% FAMILY D16.C05.03
\section{Functionals which detect vectors and subspaces}
\begin{corollary}[Norming and distance-separating functionals]\label{fa2:p1:c05:norming}
For each nonzero $x\in E$ there exists $f\in E'$ with
$\norm{f}=1$ and $f(x)=\norm{x}$. For $x=0$ the functional $f=0$
suffices for the corresponding equality with norm at most $1$.
If $M$ is a subspace and $x\notin\cl M$, there exists $f\in E'$
with $f|_M=0$, $\norm{f}=1$, and $f(x)=\dist(x,M)>0$.
\end{corollary}
\begin{proof}
On the line $\K x$ define $g(a x)=a\norm{x}$. Its norm is $1$;
extend it by the preceding theorem. For the second assertion pass to
$E/\cl M$, whose quotient norm is a norm, and norm the nonzero coset
$q(x)$. Composing its norming functional with $q$ gives a functional
vanishing on $M$. The quotient-dual isometry in Chapter~\ref{fa2:p1:c04}
preserves its norm, and its value at $x$ is the quotient norm, which is
$\dist(x,\cl M)=\dist(x,M)$.
\end{proof}

% FAMILY D16.C05.04
\begin{theorem}[Norm formula and isometric bidual embedding]\label{fa2:p1:c05:bidual-isometry}
For every $x\in E$,
\[
 \norm{x}=\sup_{\norm{f}\leq1}\abs{f(x)}
          =\sup_{\norm{f}\leq1}\operatorname{Re}f(x).
\]
Thus $J_E$ is a linear isometry and $E'$ separates points of $E$.
Its image in the Banach space $E''$ is closed exactly when $E$ is
Banach. In particular reflexivity forces completeness. For every
$T\in\Lin{E}{F}$, $\norm{T'}=\norm{T}$.
\end{theorem}
\begin{proof}
The evaluation bound gives both supremum upper bounds, and a norming
functional gives equality for $x\ne0$; all values are zero when $x=0$.
By definition the first supremum is $\norm{J_E x}$, proving the
isometry. If $x\ne y$, a norming functional for $x-y$ distinguishes
them. An isometry transfers Cauchy sequences and limits, so $J_E(E)$
is complete exactly when $E$ is complete. A complete subspace is closed,
and a closed subspace of the Banach space $E''$ is complete, proving
the equivalence and the reflexivity assertion. Finally for every $x$,
\[
 \norm{Tx}=\sup_{\norm{f}\leq1}\abs{(T'f)(x)}
 \leq\norm{T'}\norm{x}.
\]
Taking the operator norm gives $\norm{T}\leq\norm{T'}$; the converse
bound was proved before Hahn--Banach. Zero spaces are covered by the
same estimates.
\end{proof}

% FAMILY D16.C05.05
\section{Restriction, quotient duality, and an endpoint counterexample}
\begin{theorem}[Restriction as an isometric quotient]\label{fa2:p1:c05:restriction}
For any subspace $M\subset E$, restriction $r:E'\to M'$ is a
surjective contraction with kernel $M^\perp$, and each $g\in M'$
has an extension of exactly the same norm. It induces a surjective
linear isometry
$E'/M^\perp\longrightarrow M'$. If $M$ is closed, then
$(E/M)'\cong M^\perp$ isometrically by composition with the quotient
map.
\end{theorem}
\begin{proof}
Surjectivity and the exact extension norm are Hahn--Banach; the kernel
is exactly the annihilator by definition. If $r(f)=g$, every other
extension is $f+h$ for $h\in M^\perp$, and contractivity gives
$\norm{g}\leq\norm{f+h}$. Taking the infimum proves that the quotient
norm is at least $\norm{g}$. A norm-preserving extension belongs to
this same coset, giving the reverse inequality. The quotient is a
genuine normed quotient because $M^\perp$ is closed. The final
identification is the factorization proved in Chapter~\ref{fa2:p1:c04};
it does not require another extension theorem.
\end{proof}

\begin{example}[A functional on $\ell^\infty$ outside $\ell^1$]\label{fa2:p1:c05:linfinity-counterexample}
Let $c\subset\ell^\infty$ be the subspace of convergent sequences.
The limit functional $\lambda(x)=\lim_jx_j$ is linear and has norm
$1$: its modulus is at most the supremum norm, and its value at the
constant sequence $\mathbf1$ is $1$. Hahn--Banach extends it to
$L\in(\ell^\infty)'$ with norm $1$. For every $j$, $L(e_j)=0$,
while $L(\mathbf1)=1$. If $L(x)=\sum_j a_jx_j$ for some
$a\in\ell^1$, testing $e_j$ would give $a_j=0$ for every $j$,
contradicting the value at $\mathbf1$. For counting measure on $\N$,
this is precisely a continuous functional on $L^\infty$ which no
$L^1$ density represents. The construction works over both scalar fields.
\end{example}


\chapter{Convex Separation, Polars, and Convex Closure}\label{fa2:p1:c06}
\source{B01 Secs. 1.2--1.4 and 3.3; B04 Chapter 3.}
\cite{B01} \cite{B04} 


Order assertions in this chapter concern real vector spaces. On a complex
space they concern real parts of continuous complex linear functionals.
Most statements are given for normed spaces; the stated locally convex
versions will also be needed in Chapter~\ref{fa2:p1:c24}. A locally convex
space here is a topological vector space with a base of convex balanced
neighborhoods of zero. Hausdorffness is imposed whenever stated.

% FAMILY D16.C06.01
\section{The gauge of a convex neighborhood}
\begin{definition}[Minkowski functional]\label{fa2:p1:c06:gauge}
Let $U$ be an absorbing convex subset of a real vector space with
$0\in U$. Its gauge is
$p_U(x)=\inf\{t>0:x\in tU\}$. Absorbing means that the set inside
this infimum is nonempty for every $x$.
\end{definition}

\begin{lemma}[Gauge and continuity]\label{fa2:p1:c06:gauge-properties}
The gauge is finite, nonnegative, positively homogeneous, and subadditive.
For the remaining assertions suppose that $U$ is an open convex
neighborhood of zero in a topological vector space. Then
$U=\{x:p_U(x)<1\}$, and every real linear functional $h$
dominated by $p_U$ is continuous. If in addition $U$ is balanced, its
gauge is a continuous seminorm. In a normed space with $B(0,r)\subset U$,
$p_U(x)\leq\norm{x}/r$.
\end{lemma}
\begin{proof}
Finiteness follows from absorption, and positivity from the infimum of
positive numbers. For $s>0$, substitution in the defining infimum gives
$p_U(sx)=s p_U(x)$, while $p_U(0)=0$. If $x\in aU$ and $y\in bU$
with $a,b>0$, then $(x+y)/(a+b)$ is the convex combination with weights
$a/(a+b)$ and $b/(a+b)$ of $x/a$ and $y/b$, so
$p_U(x+y)\leq a+b$. Choose $a,b$ arbitrarily close to the respective
infima to obtain subadditivity. This choice is possible because if
$x\in tU$, then $x\in sU$ for $s>t$ by convexity and $0\in U$.
If $p_U(x)<1$, some $t<1$ has $x\in tU\subset U$. Conversely if
$x\in U$ and $U$ is open, continuity of $s\mapsto sx$ gives an
$s>1$ with $sx\in U$, so $p_U(x)\leq1/s<1$.

For a dominated linear $h$, on the neighborhood $U\cap(-U)$ both
$h(x)\leq p_U(x)<1$ and $-h(x)=h(-x)<1$ hold. Thus $h$ is bounded
there and, by scaling that neighborhood, is continuous at zero and hence
everywhere. If $U$ is balanced, its gauge is absolutely homogeneous
(over $\C$ as well when $U$ is complex balanced); the same neighborhood
argument bounds $p_U$ near zero. The inequalities
$p_U(x)-p_U(y)\leq p_U(x-y)$ and their reversals prove continuity
everywhere. Finally $x/t\in B(0,r)\subset U$ whenever
$t>\norm{x}/r$, giving the norm bound by an infimum.
\end{proof}

% FAMILY D16.C06.02
\section{Which kind of separation is available}
\begin{theorem}[Open convex separation]\label{fa2:p1:c06:open-separation}
Let $A,B$ be nonempty disjoint convex subsets of a real topological vector
space, and suppose $A$ is open. There are a nonzero continuous real
linear functional $h$ and $\alpha\in\R$ with
\[
                  h(a)<\alpha\leq h(b)\qquad(a\in A, b\in B).
\]
For a complex locally convex space one may write $h=\operatorname{Re}f$
for a continuous complex linear $f$. In particular, these conclusions
hold in real or complex normed spaces.
\end{theorem}
\begin{proof}
First let $C$ be nonempty open convex and $x_0\notin C$. Choose
$c\in C$, set $U=C-c$ and $z=x_0-c$. The set $U$ is absorbing:
for a fixed $x$, continuity of $t\mapsto tx$ and $0\in U$ put
$tx\in U$ for all sufficiently small positive $t$. Its gauge therefore
exists, and $z\notin U$ implies $p_U(z)\geq1$.
On $\R z$ define $\ell(tz)=t p_U(z)$. For $t\geq0$ it equals
$p_U(tz)$; for $t<0$ it is nonpositive and therefore at most the
nonnegative $p_U(tz)$. Real Hahn--Banach gives an extension $h\leq p_U$.
It is continuous by the gauge lemma, nonzero since $h(z)=p_U(z)\geq1$,
and for $u\in U$,
$h(c+u)<h(c)+1\leq h(x_0)$.

Apply this conclusion with $C=A-B$, which is open as a union of
translates of $A$, convex, nonempty, and does not contain zero.
It yields $h(a)<h(b)$ for every $a,b$. Fixing one $b$ bounds
$\sup_A h$ above, and fixing one $a$ bounds $\inf_B h$ below; hence
$\sup_A h\leq\inf_B h$ are finite numbers. Choose $\alpha$ between
them. Since $A$ is open and $h\ne0$, no point of $A$ attains its
supremum: a sufficiently small displacement in a direction with positive
$h$ remains in $A$. Thus $h(a)<\alpha$ as required.
In the complex case perform the real argument on the underlying real
space and set $f(x)=h(x)-i h(i x)$. The reconstruction from
Chapter~\ref{fa2:p1:c05} proves complex linearity and
$\operatorname{Re}f=h$; continuity follows from continuity of $h$ and
multiplication by $i$.
\end{proof}

\begin{theorem}[Strong separation from closed convex sets]\label{fa2:p1:c06:strong-separation}
Let $C$ be a nonempty closed convex subset of a normed space $E$ and
$x_0\notin C$. There exists $f\in E'$, $f\ne0$, with
\[
           \sup_{c\in C}\operatorname{Re}f(c)
                      <\operatorname{Re}f(x_0).
\]
More precisely, for any $0<r<\dist(x_0,C)$ the functional may be chosen
so that the right side exceeds the supremum by at least $r\norm{f}$.
If $C$ is closed convex and $D$ is compact convex, both nonempty and
disjoint, there is a uniform positive gap between their functional
values. The point and compact-set conclusions also hold in a Hausdorff
locally convex space, without the norm-dependent quantitative formula.
\end{theorem}
\begin{proof}
In the normed case the open convex set $C+B(0,r)$ does not contain
$x_0$. Open separation gives
$\operatorname{Re}f(c+u)<\operatorname{Re}f(x_0)$ for all
$c\in C$, $\norm{u}<r$. The supremum of $\operatorname{Re}f$ on
this ball is $r\norm{f}$: for complex scalars a phase rotation changes
the value of $f(u)$ into its modulus without changing $\norm{u}$.
Taking suprema gives the quantitative gap, which is positive.

In a locally convex space choose an open convex balanced neighborhood
$V$ of zero with $(x_0+V)\cap C=\varnothing$. Balance makes it
symmetric, so $x_0\notin C+V$. Separate $C+V$ and $x_0$ as above.
Because the resulting nonzero $h$ is continuous and $V$ absorbs every
direction, choose $v\in V$ with $h(v)>0$. Then
$h(c)+h(v)\leq h(x_0)$ for all $c$, giving a uniform positive gap.

For the last assertion $C-D$ is convex, closed, and excludes zero.
To check closedness in this generality, if $z\notin C-D$, each
$d\in D$ has $z+d\notin C$. Continuity of addition supplies
neighborhoods $U_d$ of $z$ and $V_d$ of $d$ with
$(U_d+V_d)\cap C=\varnothing$. Finitely many $V_d$ cover compact
$D$, and the intersection of their $U_d$ is a neighborhood of $z$
disjoint from $C-D$. Apply the point-separation assertion to $C-D$
and zero and reverse the sign of the functional if necessary. Its
positive uniform gap gives the required separation of $C$ and $D$.
\end{proof}

Closedness alone for two disjoint convex sets does not supply the
compact-set conclusion. For example in $\R^2$ the closed convex sets
$\{(s,t):s>0, t\geq1/s\}$ and $\{(0,t):t\in\R\}$ are disjoint
and have distance zero. A functional bounded above or below on the
entire second set must have zero $t$ coefficient, and its values on the
first set approach its value on the second; hence there is no uniform
positive separating gap.

% FAMILY D16.C06.03
\section{Polars and double annihilators}
In this section $E$ is a normed space over $\K$.
\begin{definition}[Polars and pairing topologies]\label{fa2:p1:c06:polars}
For $A\subset E$, define
$A^\circ=\{f\in E':\abs{f(a)}\leq1\text{ for all }a\in A\}$.
For $B\subset E'$ define
$ {}^\circ B=\{x\in E:\abs{f(x)}\leq1\text{ for all }f\in B\}$.
The weak topology $\sigma(E,E')$ is the coarsest topology making all
$f\in E'$ continuous. The weak-star topology $\sigma(E',E)$ is the
coarsest topology making all evaluations $f\mapsto f(x)$ continuous.
Their full development is in Chapter~\ref{fa2:p1:c12}. An absolutely
convex set is convex and balanced over its scalar field.
\end{definition}

\begin{theorem}[Polar identities and the bipolar theorem]\label{fa2:p1:c06:bipolar}
Polars reverse inclusions, the polar of a union is the intersection of
the polars, and taking the closed absolutely convex hull does not change
a polar. For $r>0$, $(rA)^\circ=r^{-1}A^\circ$, and the polar of the
closed unit ball of $E$ is the closed unit ball of $E'$. Moreover,
\[
          {}^\circ(A^\circ)=\cl\operatorname{aconv}(A),
\]
where closure is in norm and the absolutely convex hull includes zero
(also when $A$ is empty).
\end{theorem}
\begin{proof}
Inclusion reversal, the union formula, and scaling follow by unwinding
the inequalities. A linear functional bounded in modulus by $1$ on $A$
has the same bound on sums $\sum_{j=1}^n a_jx_j$ with
$\sum_j\abs{a_j}\leq1$, and then on their norm closure by continuity.
The unit-ball statement is exactly the dual norm definition.
Write $C=\cl\operatorname{aconv}(A)$. The inclusion
$C\subset{}^\circ(A^\circ)$ follows from the preceding observation.
If $x_0\notin C$, strong separation gives an $f$ with
$\operatorname{Re}f(x_0)>s=\sup_C\operatorname{Re}f$.
Since $0\in C$ and $C$ is balanced,
$s\geq0$ and $s=\sup_C\abs{f}$: rotate a value by its scalar phase
and use balance. If $s>0$, $f/s\in A^\circ$ and
$\abs{f(x_0)/s}>1$. If $s=0$, a sufficiently large positive multiple
of $f$ belongs to $A^\circ$ and also has modulus greater than $1$
at $x_0$. Thus $x_0\notin{}^\circ(A^\circ)$, proving equality.
\end{proof}

\begin{theorem}[Double-annihilator closure in the stated pairing]\label{fa2:p1:c06:double-annihilator}
For a linear subspace $M\subset E$,
$ {}^\perp(M^\perp)=\cl M$.
For a linear subspace $N\subset E'$,
\[
             ({}^\perp N)^\perp
                    =\cl^{\sigma(E',E)}N.
\]
The closure on the right in the second formula is weak-star closure,
not in general norm closure.
\end{theorem}
\begin{proof}
Continuity gives the first inclusion $\cl M\subset{}^\perp(M^\perp)$.
If $x\notin\cl M$, the distance-separating functional of
Chapter~\ref{fa2:p1:c05} vanishes on $M$ but not at $x$, giving the reverse
inclusion. In the second formula, every evaluation by an element of
$ {}^\perp N$ vanishes on the weak-star closure of $N$, giving one
inclusion. If $f_0$ is outside that closure, a basic neighborhood supplies
finitely many $x_1,\ldots,x_n\in E$ for which
$(f_0(x_1),\ldots,f_0(x_n))$ is not in the closure of the image of
$N$ under this evaluation map. That image is a linear subspace of
$\K^n$ and hence is already closed. By extending a basis of this
subspace, choose a scalar-linear functional on $\K^n$ which vanishes
on it and does not vanish on the displayed vector. Write this functional
as $z\mapsto\sum_j a_jz_j$ and put $x=\sum_j a_jx_j$.
Then $g(x)=0$ for all $g\in N$ but $f_0(x)\ne0$, excluding $f_0$
from $( {}^\perp N)^\perp$.
\end{proof}

For example take $E=\ell^1$, so $E'=\ell^\infty$, and $N=c_0$.
Its preannihilator is zero by coordinate tests. Thus $c_0$ is weak-star
dense in $\ell^\infty$ although it is a proper norm-closed subspace.

% FAMILY D16.C06.04
\section{Convex closure in the weak topology}\label{fa2:p1:c06:convex-closure}
For every convex subset $C$ of a normed space over $\K$, the conclusion
to retain is $\cl^{\sigma(E,E')}C=\cl C$, with no completeness
assumption; in particular norm-closed convex sets are weakly closed.
Its canonical statement and proof are Theorem~\ref{fa2:p1:c12:convex-closure}
in Chapter~\ref{fa2:p1:c12}, after the weak topology is constructed in
detail. That proof uses exactly the strong separation proved here:
if $x$ lies outside the norm closure, a continuous real separating
functional (a real part over $\C$) gives a weakly open half-space
separating $x$ from $C$. Thus the dependence runs from geometric
Hahn--Banach to the weak-topology result, and does not require a weak
compactness theorem. Empty $C$ has empty closure in both topologies.

% FAMILY D16.C06.05
\section{Conjugate convex functions}
In the remainder of this chapter $E$ is a real normed space. An
extended-real function $\varphi:E\to(-\infty,+\infty]$ is proper if
its finite domain is nonempty; the value $-\infty$ is excluded. It is
lower semicontinuous if every real sublevel set is closed. Its epigraph
$\{(x,t)\in E\times\R:t\geq\varphi(x)\}$ is closed exactly when it
is lower semicontinuous: if $(x_n,t_n)$ in the epigraph converges to
$(x,t)$, lower semicontinuity gives
$\varphi(x)\leq\liminf\varphi(x_n)\leq t$; conversely, intersecting a
closed epigraph with a horizontal slice gives a closed sublevel set.
Convexity of the function is equivalent, by the defining inequality,
to convexity of its epigraph.

\begin{definition}[Conjugate and biconjugate]\label{fa2:p1:c06:conjugates}
For a proper $\varphi$ define
\[
 \varphi^*(f)=\sup_{x\in E}\{f(x)-\varphi(x)\}\quad(f\in E'),
 \qquad
 \varphi^{**}(x)=\sup_{f\in E'}\{f(x)-\varphi^*(f)\}.
\]
Terms with $\varphi(x)=+\infty$ or $\varphi^*(f)=+\infty$ contribute
$-\infty$ to the respective supremum.
\end{definition}

\begin{theorem}[Fenchel inequality and Fenchel--Moreau]\label{fa2:p1:c06:fenchel-moreau}
For any proper function, Fenchel's inequality
$f(x)\leq\varphi(x)+\varphi^*(f)$ holds whenever the sum has its
usual extended-real meaning. If $\varphi$ is in addition convex and
norm-lower-semicontinuous, then $\varphi^*$ is proper, convex and
weak-star lower semicontinuous, and
$\varphi^{**}(x)=\varphi(x)$ for every $x\in E$, including points
where $\varphi(x)=+\infty$.
\end{theorem}
\begin{proof}
The inequality is the defining supremum inequality. The conjugate is a
supremum of continuous affine functions of $f$ indexed by the finite
domain of $\varphi$, so it is convex and weak-star lower semicontinuous;
each of its values is greater than $-\infty$ since that domain is nonempty.

We first show that $\varphi$ has a continuous affine minorant. Pick
$x_1$ in its finite domain and $t_1<\varphi(x_1)$. Strongly separate
$(x_1,t_1)$ from its nonempty closed convex epigraph in the normed
product $E\times\R$. A continuous linear functional on this product
has the form $L(x,t)=\ell(x)+b t$, since its restrictions to the two
coordinate subspaces are continuous. Choose the separating orientation
so that
\[
 \ell(x_1)+b t_1<\alpha\leq\ell(y)+b s
           \quad((y,s)\in\operatorname{epi}\varphi).
\]
Sending $s$ to $+\infty$ at a fixed domain point gives $b\geq0$.
Taking $(y,s)=(x_1,\varphi(x_1))$ shows $b\ne0$. Thus $b>0$, and
$a_0(y)=(\alpha-\ell(y))/b$ is a continuous affine minorant.
If $a_0(y)=h_0(y)+c_0$, then $\varphi^*(h_0)\leq-c_0<\infty$,
so $\varphi^*$ is proper.

Set $\psi=\varphi-a_0\geq0$. Fix $x\in E$ and a real
$t<\psi(x)$, with arbitrary finite $t$ allowed when $\psi(x)=+\infty$.
Separate $(x,t)$ from $\operatorname{epi}\psi$ to obtain
\[
 \ell(x)+b t<\alpha\leq\ell(y)+b\psi(y)
                    \quad(y\in\operatorname{dom}\psi),\qquad b\geq0.
\]
If $b=0$, the separator may be vertical, so division by $b$ would be
invalid. Choose $\eps>0$ small enough that
$\ell(x)+(b+\eps)t<\alpha$. This is possible by the strict original
gap (and is immediate for every positive $\eps$ when $t\leq0$ and
$b=0$). Since $\psi(y)\geq0$,
$\ell(y)+(b+\eps)\psi(y)\geq\alpha$. This argument also works when
$b>0$. Therefore
$m(y)=(\alpha-\ell(y))/(b+\eps)$ is a continuous affine minorant of
$\psi$ with $m(x)>t$.

Adding $a_0$ proves that $\varphi$ is the supremum of its continuous
affine minorants, at every finite or infinite value. If $h+c\leq\varphi$,
then $\varphi^*(h)\leq-c$, and hence
$\varphi^{**}(x)\geq h(x)-\varphi^*(h)\geq h(x)+c$.
Taking the supremum over the minorants gives
$\varphi^{**}\geq\varphi$. Fenchel's inequality gives the reverse
inequality. This proves the equality without a Banach or finite-domain
assumption on $E$ or $\varphi$.
\end{proof}

\begin{example}[Norm and convex indicator]\label{fa2:p1:c06:conjugate-examples}
For $\varphi(x)=\norm{x}$, its conjugate is $0$ on the closed dual
unit ball and $+\infty$ outside it: the unit-ball case is the dual norm
inequality, while a vector with $f(x)>\norm{x}$ can be scaled to make
the defining supremum infinite. For a nonempty closed convex $C$, the
convex indicator $I_C$, equal to $0$ on $C$ and $+\infty$ elsewhere,
has conjugate $\sup_{x\in C}f(x)$, its support function. The identity
$I_C^{**}=I_C$ is a precise separation statement. These conclusions
are consequences of conjugation; no general optimization duality or
unproved minimax principle has been invoked.
\end{example}


\chapter{Baire Category and Uniform Boundedness}\label{fa2:p1:c07}
\source{B01 Secs. 2.1--2.2; B05 Secs. 1.5--1.6.}
\cite{B01} \cite{B05} 


% FAMILY D16.C07.01
\section{The category argument}
\begin{definition}[Nowhere dense and meagre sets]\label{fa2:p1:c07:category}
A subset of a topological space is nowhere dense if its closure has
empty interior. It is meagre if it is a countable union of nowhere
dense sets. A space is Baire if every countable intersection of open
dense subsets is dense.
\end{definition}

\begin{theorem}[Baire's complete-metric theorem]\label{fa2:p1:c07:baire}
Every complete metric space is Baire. In particular, a nonempty complete
metric space which is a countable union of closed sets has at least one
member with nonempty interior. No nonempty open subset of a complete
metric space is meagre in the ambient space.
\end{theorem}
\begin{proof}
Let $G_n$ be open dense sets and let $O$ be any nonempty open set.
Choose a closed ball $\overline B(x_1,r_1)$ with $r_1>0$,
$r_1<1/2$, contained in $O\cap G_1$. Inductively the nonempty open
ball $B(x_n,r_n)$ intersects dense $G_{n+1}$. At a point of that open
intersection choose a positive-radius closed ball
\[
 \overline B(x_{n+1},r_{n+1})\subset B(x_n,r_n)\cap G_{n+1},
 \qquad r_{n+1}<2^{-n-1}.
\]
Such closed balls exist because a sufficiently small closed ball around
a point of an open metric set stays inside it. The centers are Cauchy:
for $m\geq n$, $x_m\in\overline B(x_n,r_n)$, so their distances are
at most $r_n$. Their limit exists by completeness and belongs to every
closed ball by closedness. It therefore belongs to $O\cap\bigcap_nG_n$.
Since $O$ was arbitrary, the intersection is dense. Empty metric spaces
satisfy the density assertion vacuously.

If closed $F_n$ all have empty interior, their complements are open
dense and have dense intersection, so they cannot cover a nonempty
space. If a nonempty open $O$ were contained in a countable union of
nowhere dense sets, applying the just proved intersection argument to
the complements of their closures would give a point of $O$ outside
that union, a contradiction.
\end{proof}

Completeness is used in obtaining the common point of the nested balls.

% FAMILY D16.C07.02
\section{Banach--Steinhaus}
\begin{theorem}[Uniform boundedness principle]\label{fa2:p1:c07:uniform-boundedness}
Let $E$ be Banach, $F$ normed, and $(T_i)_{i\in I}$ any family in
$\Lin{E}{F}$. If for every $x\in E$,
$\sup_{i\in I}\norm{T_i x}<\infty$, then
$\sup_{i\in I}\norm{T_i}<\infty$. The target need not be complete
and the index set need not be countable. For an empty family the
nonnegative supremum is understood to be zero.
\end{theorem}
\begin{proof}
The empty-family assertion is immediate. Otherwise set
$A_n=\{x:\norm{T_i x}\leq n\text{ for all }i\in I\}$ for $n\geq1$.
Each is closed as an arbitrary intersection of closed inverse images,
and pointwise boundedness says $E=\bigcup_{n\geq1}A_n$.
Baire gives a ball $B(x_0,r)\subset A_N$ for some $N$ and $r>0$.
Both $x_0$ and $x_0+h$ belong to $A_N$ when $\norm{h}<r$, hence
\[
              \norm{T_i h}\leq\norm{T_i(x_0+h)}+\norm{T_ix_0}
                            \leq2N\qquad(i\in I).
\]
For a unit vector $u$, take $h=(r/2)u$ to obtain
$\norm{T_i u}\leq4N/r$. Scaling, including the zero vector, proves
$\norm{T_i}\leq4N/r$ for every $i$. The proof uses completeness only
in the domain through Baire.
\end{proof}

\begin{corollary}[Weak boundedness detects norm boundedness]\label{fa2:p1:c07:weak-boundedness}
For any normed space $E$, a set $A\subset E$ is norm bounded exactly
when $\{f(x):x\in A\}$ is bounded for every $f\in E'$.
\end{corollary}
\begin{proof}
The forward direction is the evaluation estimate. For the reverse,
consider the family $\{J_Ex:x\in A\}\subset(E')'$. It is pointwise
bounded by hypothesis on the Banach domain $E'$. Uniform boundedness
gives $\sup_{x\in A}\norm{J_Ex}<\infty$. The isometry of $J_E$ then
gives the desired norm bound. This argument does not require $E$ itself
to be Banach. Empty $A$ is bounded by convention.
\end{proof}

% FAMILY D16.C07.03
\section{Pointwise limits of operators}
\begin{theorem}[Continuity of an existing pointwise limit]\label{fa2:p1:c07:pointwise-limit}
Let $E$ be Banach, $F$ normed, and $T_n\in\Lin{E}{F}$. Suppose that
for every $x\in E$ the sequence $T_nx$ has a limit $Tx$ in $F$.
Then $T$ is bounded linear,
$\sup_n\norm{T_n}<\infty$, and
$\norm{T}\leq\liminf_n\norm{T_n}$.
\end{theorem}
\begin{proof}
Each convergent sequence $(T_nx)$ is bounded, so uniform boundedness
gives a common operator bound $M$. Continuity of the vector operations
lets limits pass through
$T_n(a x+b y)=aT_nx+bT_ny$, proving linearity of $T$.
Passing to the limit in $\norm{T_nx}\leq M\norm{x}$ gives its
boundedness. More precisely for each $x$,
$\norm{Tx}\leq(\liminf_n\norm{T_n})\norm{x}$, by passing to a
subsequence of operator norms realizing the lower limit when it is
finite; if the lower limit is infinite the assertion is immediate.
Taking the operator norm proves the stated inequality. Completeness
of $F$ is unnecessary because the pointwise limits were assumed to
exist in $F$.
\end{proof}

Pointwise convergence is often called strong operator convergence.
It does not imply norm convergence. On $\ell^2$, the rank-one maps
$P_nx=x_ne_n$ have norm $1$ and satisfy $\norm{P_nx}=\abs{x_n}\to0$
for each $x$, while $\norm{P_n-0}=1$ for all $n$.

% FAMILY D16.C07.04
\begin{theorem}[Extending convergence from a dense set]\label{fa2:p1:c07:dense-convergence}
Let $E$ be normed, $F$ Banach, $D\subset E$ dense, and
$T_n\in\Lin{E}{F}$ with $\sup_n\norm{T_n}\leq M<\infty$.
If $(T_nd)$ converges for every $d\in D$, then there is a bounded
linear $T:E\to F$ such that $T_nx\to Tx$ for every $x\in E$ and
$\norm{T}\leq M$. The domain need not be complete. If a bounded
$T:E\to F$ is already given and convergence to $Td$ holds on $D$,
then convergence to $Tx$ holds on all $E$ without completeness of $F$.
\end{theorem}
\begin{proof}
For $x\in E$ and $d\in D$,
\[
 \norm{T_nx-T_mx}\leq2M\norm{x-d}+\norm{T_nd-T_md}.
\]
Choose $d$ so that the first term is smaller than $\eps/2$ (if $M=0$
all maps vanish), then take $m,n$ large so that the second is smaller
than $\eps/2$. Thus $(T_nx)$ is Cauchy; completeness of $F$ supplies
its limit $Tx$. Passing through the linear identities and the estimate
$\norm{T_nx}\leq M\norm{x}$ proves linearity and the norm bound.
For the variant with an existing bounded $T$, use instead
\[
 \norm{T_nx-Tx}\leq(M+\norm{T})\norm{x-d}+\norm{T_nd-Td},
\]
and choose $d$ then $n$ in the same order. This proves convergence
directly without constructing any limit in an incomplete target.
\end{proof}

% FAMILY D16.C07.05
\section{Why the hypotheses matter}
\begin{example}[An incomplete domain]\label{fa2:p1:c07:incomplete-counterexample}
On $c_{00}$ with the supremum norm let $T_nx=n x_n$. Then
$\norm{T_n}=n$, by testing $e_n$, but for every finitely supported
$x$, $T_nx=0$ for all sufficiently large $n$. Thus pointwise boundedness
on an incomplete normed domain need not give uniform norm bounds.
Also $S_nx=\sum_{j\leq n}x_j$ has norm $n$ and converges pointwise
on this domain to $Sx=\sum_jx_j$. The limit is unbounded, since its
value on $\sum_{j=1}^n e_j$, a unit vector in the supremum norm, is $n$.
Thus even continuity of a pointwise limit can fail without completeness
of the domain.
\end{example}

\begin{example}[Boundedness on a dense set is insufficient]\label{fa2:p1:c07:dense-counterexample}
On the Banach space $c_0$, again take $T_nx=n x_n$, with norm $n$.
For each $x\in c_{00}$ the family is pointwise bounded, and $c_{00}$
is dense in $c_0$. Nevertheless for $x_j=j^{-1/2}$ one has
$x\in c_0$ and $T_nx=\sqrt n\to\infty$.
This identifies exactly the missing hypothesis: uniform boundedness
requires pointwise boundedness at every point of the Banach domain.
The dense-convergence theorem requires an independent uniform operator
bound and cannot obtain it from dense-set boundedness alone.
\end{example}


\chapter{Open Mapping, Bounded Inverses, and Closed Graphs}\label{fa2:p1:c08}
\source{B01 Secs. 2.3--2.5 and 2.7; B05 Sec. 1.6.}
\cite{B01} \cite{B05} 


% FAMILY D16.C08.01
\section{Openness by successive corrections}
\begin{theorem}[Open mapping theorem]\label{fa2:p1:c08:open-mapping}
Let $E,F$ be Banach spaces over the same scalar field, and let
$T:E\to F$ be bounded linear and surjective. Then $T$ maps open sets
to open sets. In particular, there is $c>0$ such that
\[
                       B_F(0,c)\subset T(B_E(0,1)).
\]
Consequently there is $C>0$ such that every $y\in F$ has at least one
preimage $x$ with $\norm{x}\leq C\norm{y}$. This selection is not
asserted to be linear or continuous.
\end{theorem}
\begin{proof}
Write $B=B_E(0,1)$ and $C_0=\cl T(B)$. Surjectivity implies
$F=\bigcup_{n\geq1}nC_0$: for any $y$ choose $Tx=y$ and an integer
$n>\norm{x}$. Each $nC_0$ is closed, so Baire on $F$ makes one of
them have nonempty interior. Scaling shows $C_0$ has nonempty interior.
It is convex and balanced, since $B$ and its linear image have those
properties and closure preserves them. Choose $y_0\in F$ and $a>0$
with $B_F(y_0,2a)\subset C_0$. As $y_0\in C_0$,
\[
 B_F(0,2a)\subset C_0-C_0\subset2C_0,
 \qquad B_F(0,a)\subset C_0.
\]
This step gives only a closure of an image; the next step removes that
closure using completeness of $E$.

Fix $y$ with $\norm{y}<a/2$. By scaling the inclusion just obtained,
$B_F(0,a/2)\subset\cl T(B_E(0,1/2))$, so choose
$x_1\in B_E(0,1/2)$ such that $\norm{y-Tx_1}<a/4$.
Inductively, if the current residual has norm less than $a2^{-n}$,
the inclusion
$B_F(0,a2^{-n})\subset\cl T(B_E(0,2^{-n}))$
provides $x_n$ with
\[
 \norm{x_n}<2^{-n},\qquad
 \norm{y-T(x_1+\cdots+x_n)}<a2^{-n-1}.
\]
The absolutely summable series $\sum_nx_n$ converges in $E$ to $x$.
Its norm is strictly less than $1$: its first term is strictly smaller
than $1/2$, and the remaining norm sum is at most $1/2$.
Continuity of $T$ and the residual estimate imply $Tx=y$. Hence the
displayed conclusion holds with $c=a/2$.

If $U\subset E$ is open and $y_0=Tx_0\in T(U)$, choose
$r>0$ with $B_E(x_0,r)\subset U$. Linearity and the ball inclusion
give $B_F(y_0,rc)\subset T(U)$, proving openness. For $y\ne0$ apply
the inclusion to $cy/(2\norm{y})$ and rescale to obtain a preimage
with norm at most $2\norm{y}/c$; for $y=0$ choose $x=0$.
Thus $C=2/c$ works. The proof includes zero spaces and does not divide
by an operator norm that might vanish.
\end{proof}

% FAMILY D16.C08.02
\begin{corollary}[Bounded inverse theorem]\label{fa2:p1:c08:bounded-inverse}
A bounded linear bijection between Banach spaces has bounded linear
inverse. If $B_F(0,c)\subset T(B_E(0,1))$, then
$\norm{T^{-1}}\leq1/c$.
\end{corollary}
\begin{proof}
The inverse is algebraically linear. For $y\ne0$ and any
$t>\norm{y}/c$, the ball inclusion gives a preimage of $y/t$ of norm
less than $1$. Injectivity identifies it with $T^{-1}y/t$, so
$\norm{T^{-1}y}<t$. Let $t\downarrow\norm{y}/c$. The zero case
is immediate. This is precisely boundedness with the stated norm bound.
\end{proof}

% FAMILY D16.C08.03
\section{Closed graphs}
\begin{theorem}[Closed graph theorem]\label{fa2:p1:c08:closed-graph}
Let $E,F$ be Banach spaces and let $T:E\to F$ be an everywhere-defined
linear map. Then $T$ is bounded if and only if its graph
$G(T)=\{(x,Tx):x\in E\}$ is closed in $E\times F$ with the product
norm $\norm{(x,y)}=\norm{x}+\norm{y}$. Equivalently, if
$x_n\to x$ and $Tx_n\to y$, then $y=Tx$.
\end{theorem}
\begin{proof}
If $T$ is bounded, the indicated convergences give $Tx_n\to Tx$,
so uniqueness of limits proves closedness. Conversely a closed graph
is a closed linear subspace of the product Banach space and hence is
Banach. The coordinate projection
$p:G(T)\to E$, $(x,Tx)\mapsto x$, is a bounded linear bijection;
surjectivity uses that $T$ is defined on all of $E$. Its inverse is
bounded by the bounded inverse theorem. The second product projection
is bounded, and its composition with $p^{-1}$ is $T$, hence $T$ is
bounded. In a metric product, closedness is equivalent to containing
all limits of convergent sequences, giving the sequential formulation.
\end{proof}

A closed linear map defined only on a proper domain need not be bounded
in the ambient norm. Even an everywhere-defined map between normed
spaces need not satisfy this theorem when one of the completeness
conditions is absent.

\begin{example}[Failure of inverse and graph conclusions without completeness]\label{fa2:p1:c08:incomplete-norms}
Take the vector space $\ell^1$ once with its usual Banach norm and once
with the supremum norm. The identity from the first space to the second
is a bounded bijection. Its inverse is unbounded, since a vector with
its first $n$ entries equal to $1$ has supremum norm $1$ and
$\ell^1$ norm $n$. The supremum-norm space is incomplete: truncations
of $(1/j)_{j\geq1}$ are Cauchy in that norm, but their coordinate limit
is not in $\ell^1$.
The unbounded inverse nevertheless has closed graph. Indeed, if
$x_n\to x$ in supremum norm and $x_n\to y$ in $\ell^1$ norm, the
second convergence also holds in supremum norm, so $x=y$. This gives
an everywhere-defined closed unbounded linear map from an incomplete
domain to a Banach target.
\end{example}

% FAMILY D16.C08.04
\section{Equivalent complete norms and closed-range estimates}
\begin{corollary}[Comparison of two complete norms]\label{fa2:p1:c08:equivalent-norms}
Let a vector space $E$ have two Banach norms $\norm{\cdot}_1$ and
$\norm{\cdot}_2$. If $\norm{x}_2\leq C\norm{x}_1$ for all $x$ and
some finite $C$, then also $\norm{x}_1\leq D\norm{x}_2$ for some
finite $D$; thus the norms are equivalent.
\end{corollary}
\begin{proof}
Apply the bounded inverse theorem to the identity from the first
Banach space to the second. Both completeness hypotheses are essential,
as the preceding example demonstrates.
\end{proof}

\begin{theorem}[Closed range and distance to the kernel]\label{fa2:p1:c08:closed-range}
For $T\in\Lin{E}{F}$ between Banach spaces, the range is closed if
and only if there is $C>0$ such that
\[
                      \dist(x,\ker T)\leq C\norm{Tx}
                                \qquad(x\in E).
\]
For an injective $T$ this reduces to boundedness below:
$\norm{x}\leq C\norm{Tx}$. In particular a bounded below map from
a Banach space into a normed space has closed range and a bounded
inverse on its range.
\end{theorem}
\begin{proof}
The kernel $N$ is closed, so $E/N$ is Banach by the explicit quotient
construction of Chapter~\ref{fa2:p1:c02}. The induced map
$\widetilde T:E/N\to\ran T$, $x+N\mapsto Tx$, is bounded and
bijective. If $\ran T$ is closed in Banach $F$, it is Banach, so the
bounded inverse theorem gives
$\norm{x+N}\leq C\norm{Tx}$, the claimed distance estimate.
If both quotient and range are zero any positive $C$ suffices.

Conversely suppose the estimate holds and $Tx_n\to y$ in $F$.
For differences the estimate gives
$\norm{(x_n+N)-(x_m+N)}\leq C\norm{Tx_n-Tx_m}$, so the cosets
are Cauchy. Completeness of $E/N$ supplies a limit coset $x+N$.
Continuity of $\widetilde T$ gives $Tx_n\to Tx$, hence $y=Tx$.
This proves closedness. When $N=\{0\}$ the estimate is the displayed
lower bound and directly bounds the inverse on the range. The same
argument for an injective map uses completeness only of $E$, so the
last assertion allows an incomplete target.
\end{proof}

The estimate concerns distance to the entire kernel. It does not choose
a bounded linear representative of each coset and therefore does not
by itself assert that the kernel has a complement.

% FAMILY D16.C08.05
\section{When one-sided inverses exist}
\begin{theorem}[Right and left inverse criteria]\label{fa2:p1:c08:splitting}
Let $E,F$ be Banach and $T\in\Lin{E}{F}$.
If $T$ is surjective, it has a bounded right inverse $S:F\to E$
with $TS=\id_F$ if and only if $\ker T$ is complemented in $E$.
If $T$ is injective, it has a bounded left inverse $S:F\to E$ with
$ST=\id_E$ if and only if its range is closed and complemented in $F$.
\end{theorem}
\begin{proof}
If $TS=\id_F$, then $P=ST$ is a bounded projection on $E$.
Its range is $S(F)$, since $ST(Sy)=Sy$, and its kernel is $\ker T$:
$Tx=0$ gives $Px=0$, while $Px=0$ gives $Tx=TPx=0$.
The bounded projection theorem therefore supplies the desired complement.
Conversely let $E=\ker T\oplus N$ be a topological direct sum.
Its summand $N$ is closed and Banach. The restriction $T|_N:N\to F$
is a bounded bijection, so its inverse is bounded. Compose that inverse
with the inclusion $N\hookrightarrow E$ to obtain the right inverse.

If $ST=\id_E$, then $Q=TS$ is a bounded projection on $F$ whose
range is exactly $T(E)$, since $Q(Tx)=Tx$. That range is closed and
complemented by the kernel of $Q$. Conversely, if $R=T(E)$ is closed
and $P:F\to R$ is its bounded projection, $T:E\to R$ is a bounded
bijection between Banach spaces, so $S=(T:E\to R)^{-1}P$ is bounded
and satisfies $ST=\id_E$. The identities and arguments remain valid
for zero spaces.
\end{proof}

In particular a surjection onto a finite-dimensional space has a bounded
right inverse, because its closed kernel has finite codimension. An
injection from a finite-dimensional space has a bounded left inverse,
because its finite-dimensional range is closed and complemented by
Chapter~\ref{fa2:p1:c02}. These special cases do not imply a complement
for an arbitrary closed subspace of an infinite-dimensional Banach space.

% END INPUT


% BEGIN parts/part1-tvs.tex
\chapter{Vector topologies and the general fundamental principles}
\label{fa2:p1:tvs}
The norm is one way of relating algebra to convergence. The topology of
uniform convergence on compact sets, or the topology determined by all
derivatives on compact sets, usually has no defining norm. We therefore
develop the topological statements themselves. Local convexity is imposed
only where the argument needs convex neighborhoods. Sources are
\cite[Chapters 1--3]{B04}, \cite[Chapter IV]{B02} and
\cite[Sections 2.1--2.4]{B05}.

\section{Neighborhoods, bounded sets and finite dimension}
\begin{definition}\label{fa2:p1:tvs-definitions}
A topological vector space, abbreviated TVS, is a vector space with jointly
continuous addition and scalar multiplication. We require it to be Hausdorff
unless explicitly stated otherwise. A set $A$ is balanced if $\lambda A\subset
A$ for $|\lambda|\leq1$, absorbing if every vector belongs to some positive
multiple of $A$, and bounded if for every zero neighborhood $U$ there is
$s>0$ with $A\subset tU$ whenever $t\geq s$. A locally convex space, or LCS,
has a zero-neighborhood base consisting of convex sets. An F-space is a
complete metrizable TVS; a Fr\'echet space is a locally convex F-space.
Completeness means convergence of every Cauchy net for the translation
uniformity: for every zero neighborhood $U$, the differences of sufficiently
late terms lie in $U$.
\end{definition}

\begin{proposition}\label{fa2:p1:balanced-bounded}
Every TVS has a base of balanced zero neighborhoods. An LCS has a base of
absolutely convex neighborhoods. Every zero neighborhood is absorbing.
Compact sets and convergent sequences are bounded. Closures, finite unions,
finite sums and scalar multiples of bounded sets are bounded. A linear map
is continuous if and only if it is continuous at zero; continuous linear
maps take bounded sets to bounded sets.
\end{proposition}
\begin{proof}
Continuity of $\lambda\mapsto\lambda x$ at zero proves absorption. Given
a zero neighborhood $U$, continuity of scalar multiplication at each
$(\lambda,0)$, $|\lambda|\leq1$, and compactness of the scalar unit disk give
a zero neighborhood $V$ with $\lambda V\subset U$ for all $|\lambda|\leq1$.
The union $\bigcup_{|\lambda|\leq1}\lambda V$ is balanced, contains $V$ and
lies in $U$; its interior is balanced and is a neighborhood. If $U$ is convex,
the absolutely convex hull of $V$ lies in $U$: a finite sum
$\sum a_jv_j$, $\sum|a_j|\leq1$, is a convex combination of points
$\lambda v_j\in U$ and zero. Taking the interior gives the LCS assertion.

For compact $K$ and an open balanced neighborhood $V$, the increasing open sets
$nV$ cover $K$ and a finite subcover gives $K\subset nV$. A convergent
sequence together with its limit is compact. Given $U$, choose a closed
neighborhood $V\subset U$; if $A\subset tV$, then
$\overline A\subset tV\subset tU$. Such closed neighborhoods exist because
one may choose $W$ with $W-W\subset U$, and then
$\overline W\subset W-W\subset U$. For sums choose $V+V\subset U$;
the remaining bounded-set assertions follow directly from the definition.
Translation of an increment proves the continuity-at-zero criterion. If
$T^{-1}(U)$ contains a zero neighborhood $V$ and $A\subset tV$, then
$T(A)\subset tU$, proving preservation of boundedness.
\end{proof}

\begin{theorem}\label{fa2:p1:finite-dimensional-tvs}
Every Hausdorff vector topology on a finite-dimensional vector space is its
usual Euclidean topology. Finite-dimensional subspaces of Hausdorff TVSs
are closed. A locally compact Hausdorff TVS is finite dimensional.
\end{theorem}
\begin{proof}
For a basis $e_1,\ldots,e_n$, the map
$A:\K^n\to E$, $A(a)=\sum a_je_j$, is continuous. Its image of the
Euclidean unit sphere is compact and does not contain zero. A balanced
zero neighborhood $V$ disjoint from that image satisfies
$A^{-1}(V)\subset\{a:|a|_2<1\}$: otherwise scale a point of $V$ onto the
sphere. Scaling $V$ proves continuity of $A^{-1}$. A net in a
finite-dimensional subspace which converges in the ambient space is
Cauchy there. The same uniformity argument and the inverse coordinate map
make it Euclidean Cauchy, so it converges in the subspace. Hausdorff
uniqueness makes its ambient limit a point of the subspace, proving
closedness.

For the last assertion choose a compact zero neighborhood $K$ and an open
neighborhood $U\subset K$. Finitely many translates
$x_j+\tfrac12 U$ cover $K$. Let $F=\operatorname{span}\{x_j\}$, a closed
finite-dimensional subspace. In the Hausdorff quotient $E/F$, with quotient
map $q$, the covering gives $q(K)\subset\tfrac12 q(K)$.
If $z\in q(K)$, induction gives $2^nz\in q(K)$ for every $n$. The set
$q(K)$ is bounded, so $2^{-n}(2^nz)$ tends to zero. Thus $z=0$.
Consequently $K\subset F$, and absorption of $K$ gives $E=F$.
The quotient facts used here are proved directly in
Proposition~\ref{fa2:p1:tvs-quotient}, independently of this finite-dimensional
conclusion.
\end{proof}

\section{Seminorms, metrization and quotients}
\begin{theorem}\label{fa2:p1:seminorm-model}
A separating family $(p_i)$ of seminorms determines a Hausdorff locally
convex topology with basic neighborhoods
$\{x:p_{i_j}(x)<\eps,\ 1\leq j\leq m\}$. Every Hausdorff LCS is obtained
in this way from its continuous seminorms. A linear $T:E\to F$ between LCSs
is continuous exactly when for each continuous target seminorm $q$ there
are continuous source seminorms $p_1,\ldots,p_m$ and $C>0$ such that
\[
 q(Tx)\leq C\max_jp_j(x)\quad(x\in E).
\]
\end{theorem}
\begin{proof}
The seminorm triangle inequality gives addition continuity; homogeneity
gives scalar continuity, including at $(0,x)$ by first bounding each of
the finitely many seminorms near $x$. Separation makes the topology
Hausdorff. Conversely, the gauges of open absolutely convex zero
neighborhoods are continuous seminorms and recover those neighborhoods,
by Lemma~\ref{fa2:p1:c06:gauge-properties}. If $T$ is continuous, some basic
source neighborhood is mapped into $\{q<1\}$. With
$p=\max p_j$, scaling $x$ by any $t>p(x)/\eps$ gives
$q(Tx)\leq p(x)/\eps$. If $p(x)=0$, arbitrary scaling gives $q(Tx)=0$.
The reverse implication follows by taking inverse images of the displayed
basic neighborhoods.
\end{proof}

\begin{theorem}\label{fa2:p1:tvs-metrization}
A Hausdorff TVS is metrizable if and only if it has a countable local base.
It then has a translation-invariant metric with balanced balls. For an
LCS the balls can also be chosen convex. A Hausdorff LCS is normable if
and only if it has a bounded zero neighborhood.
\end{theorem}
\begin{proof}
A metric gives a countable base of balls. In the converse direction take
balanced neighborhoods $U_0=E,U_1,U_2,\ldots$ forming a base and satisfying
$U_{n+1}+U_{n+1}+U_{n+1}\subset U_n$. Give a step in $U_n$ the cost
$2^{-n}$ and define $d(x,y)$ as the infimum of costs of finite chains
from $x$ to $y$. Concatenating chains proves the triangle inequality;
reversing or translating them proves symmetry and invariance. The
standard carrying argument is as follows. In a chain of cost less than
$2^{-n}$, replace two steps of the smallest equal cost by their sum,
which lies in the preceding neighborhood, and repeat. At most one step
of each cost remains; summing from the smallest costs upward, the
three-neighborhood inclusion keeps the partial sum in $U_n$.
It follows, with an immaterial shift of one index if the strict inequality
is replaced by a nonstrict one, that
\[
 \{d(0,x)<2^{-n}\}\subset U_n,
 \qquad U_n\subset\{d(0,x)\leq2^{-n}\}.
\]
These inclusions imply positive definiteness and equality of the two
topologies. Balance of all the $U_n$ implies $d(0,\lambda x)\leq d(0,x)$
for $|\lambda|\leq1$, so metric balls are balanced.

For an LCS choose a countable separating defining sequence $(p_n)$ and use
the simpler explicit metric
\[
 d(x,y)=\sup_{n\geq1}2^{-n}\min\{1,p_n(x-y)\}.
\]
Each summand is subadditive, hence so is the supremum. For a fixed positive
radius only finitely many seminorm constraints can restrict the open ball;
the remaining summands are uniformly smaller than that radius. The ball
is therefore an intersection of finitely many open seminorm balls and is
absolutely convex. The same observation proves compatibility.

If a bounded neighborhood exists, take an open absolutely convex
neighborhood $V$ inside it. Its gauge $p$ is a continuous seminorm. Since
$V$ is bounded, for every zero neighborhood $W$ there is $t$ with
$V\subset tW$. Thus the $p$-balls form a neighborhood base. If $p(x)=0$,
then $x\in sV$ for all $s>0$, and boundedness forces $x$ into every zero
neighborhood; hence $x=0$. Therefore $p$ is a defining norm. A norm unit
ball is plainly a bounded neighborhood, proving the converse.
\end{proof}

\begin{proposition}\label{fa2:p1:tvs-quotient}
For a closed subspace $M$ of a Hausdorff TVS $E$, the quotient topology
on $E/M$ is Hausdorff and makes the quotient map $q$ continuous, open and
linear. A map out of the quotient is continuous exactly when its
composition with $q$ is continuous. Local convexity, a countable local
base, and completeness in the metrizable case pass to this quotient.
Complete subspaces of Hausdorff TVSs are closed. In particular closed
subspaces and closed quotients of Fr\'echet spaces are Fr\'echet.
\end{proposition}
\begin{proof}
For open $U\subset E$, $q^{-1}q(U)=U+M$ is open. This proves openness and
the universal continuity statement. Addition and scalar multiplication
are continuous by lifting neighborhoods. Distinct cosets are separated
because a vector outside the closed $M$ has a neighborhood disjoint from
$M$; choose a smaller difference neighborhood to separate the cosets.
Images of a local base are a local base, and images of convex sets are
convex.

To prove completeness, let $(y_n)$ be a Cauchy sequence in a metrizable
quotient and fix balanced neighborhoods $U_n$ in $E$ so small that
$U_{n+1}+U_{n+1}\subset U_n$, $d_E(0,U_n)\leq2^{-n}$ in the sense of
metric diameter, and $(U_n)$ is a local base. Choose a subsequence
$(y_{n_k})$ with $y_{n_{k+1}}-y_{n_k}\in q(U_k)$. Lift its successive
differences to $u_k\in U_k$ and lift $y_{n_1}$ to $x_1$. The partial sums
$x_1+\sum_{k<m}u_k$ are Cauchy in $E$, so converge to $x$. Their quotient
images converge to $q(x)$. The original Cauchy sequence has the same
limit. In metric spaces sequence completeness and Cauchy-net
completeness agree: choose a Cauchy sequence of increasingly late net
terms and use the Cauchy property to transfer its limit to the net.

If $x$ is in the closure of a complete subspace, a net in that subspace
converging to $x$ is Cauchy and therefore has a limit in the subspace;
Hausdorff uniqueness proves closedness. A closed subspace of a complete
TVS is complete by taking ambient limits. These observations prove the
last assertions.
\end{proof}

\begin{example}\label{fa2:p1:nonlocally-convex}
On a nonatomic measure space and $0<p<1$, set
$d(f,g)=\int|f-g|^p$. This is a complete translation-invariant metric
on $L^p$. Its continuous linear dual is zero, so a nonzero such space
cannot be locally convex or normable.
\end{example}
\begin{proof}
The scalar inequality $|a+b|^p\leq|a|^p+|b|^p$ gives the metric triangle
inequality. From a Cauchy sequence choose a subsequence whose successive
metric differences have summable integrals. Tonelli makes the sum of
pointwise $p$th powers finite almost everywhere. The successive
differences are then absolutely summable almost everywhere (eventually
$|a|\leq |a|^p$); their sums define $f$, and the tail integral estimate
proves convergence in the metric. The Cauchy sequence itself converges.

If $\ell$ is continuous, homogeneity and its bound on a metric ball imply
$|\ell(f)|\leq C(\int|f|^p)^{1/p}$. For any $f$, nonatomicity partitions
the finite measure $|f|^p\,d\mu$ into $n$ equal pieces. Writing $f$ as
the sum of its restrictions gives
$|\ell(f)|\leq C n(\int|f|^p/n)^{1/p}$, which tends to zero since
$1-1/p<0$. Thus $\ell=0$. Hahn--Banach separates a nonzero point in
every Hausdorff LCS, contradicting this if the space were locally convex.
\end{proof}

\section{Equicontinuity, open mapping and separate continuity}
\begin{theorem}[Topological Banach--Steinhaus]\label{fa2:p1:tvs-ubp}
Let $\mathcal A$ be continuous linear maps from a TVS $E$ into a TVS $F$.
If the set of $x$ with bounded orbit $\{Tx:T\in\mathcal A\}$ is
nonmeagre in $E$, it is all of $E$, and $\mathcal A$ is equicontinuous:
for every zero neighborhood $V$ in $F$ there is a zero neighborhood $U$
in $E$ with $T(U)\subset V$ for all $T\in\mathcal A$. In particular
pointwise bounded families on an F-space are equicontinuous.
\end{theorem}
\begin{proof}
Choose a closed balanced zero neighborhood $W$ with $W-W\subset V$.
The closed sets $A_n=\bigcap_T T^{-1}(nW)$ cover the nonmeagre set in
the hypothesis. At least one $A_n$ has nonempty interior; otherwise their
union is meagre. Thus $x_0+U\subset A_n$ for a zero neighborhood $U$,
after reducing $U$ if necessary. Subtraction gives
$T(U)\subset n(W-W)$; hence $T(n^{-1}U)\subset V$. This is
equicontinuity. A singleton is bounded, and an equicontinuous family
maps a bounded set into one common bounded set: given $V$, first find
$U$, then use $B\subset tU$. Therefore every orbit is bounded. An
F-space is Baire by complete metrizability, which proves the last case.
\end{proof}

\begin{theorem}[Open mapping and closed graph for F-spaces]
\label{fa2:p1:fspace-open-graph}
A continuous linear surjection between F-spaces is open. Consequently a
linear map between F-spaces is continuous whenever its graph is closed.
Local convexity is unnecessary.
\end{theorem}
\begin{proof}
Let $T:E\to F$ be the surjection. For a zero neighborhood $U$ choose
$W-W\subset U$. Since $F=\bigcup_{n\geq1}T(nW)$, Baire gives interior
to some $\overline{T(nW)}$, so $\overline{T(U)}$ is a zero
neighborhood, by subtracting two points of that interior and scaling.

We remove the closure. Choose zero neighborhoods $U_0,U_1,\ldots$ in $E$
with $\overline{U_0}\subset U$, $U_{n+1}+U_{n+1}\subset U_n$, and
metric diameters tending to zero fast enough that any series
$\sum_{n\geq1}x_n$, $x_n\in U_n$, is Cauchy. For instance require the
diameter of $U_n$ at most $2^{-n}$ in a compatible invariant metric.
Choose zero neighborhoods $V_n$ in $F$, with diameter tending to zero,
such that $V_n\subset\overline{T(U_{n+1})}$. Starting with
$y\in V_0$, choose successively $x_{n+1}\in U_{n+1}$ so that
$y-\sum_{j\leq n+1}Tx_j\in V_{n+1}$. This is possible by the closure
inclusion and openness of $V_{n+1}$. The series converges to $x$ in
$E$; its partial sums lie in $U_0$, so $x\in U$. Continuity gives
$Tx=y$. Hence $V_0\subset T(U)$, proving openness.

For the graph theorem a closed graph $G\subset E\times F$ is an F-space.
The first coordinate projection $G\to E$ is a continuous linear
bijection and hence has continuous inverse by open mapping. Composing
that inverse with the second projection proves continuity of the map.
\end{proof}

\begin{proposition}\label{fa2:p1:separate-bilinear}
Let $B:E\times F\to G$ be separately continuous and bilinear, where
$E$ is an F-space and $F,G$ are TVSs. Then $B$ is sequentially jointly
continuous. If $F$ is metrizable, $B$ is jointly continuous.
\end{proposition}
\begin{proof}
For $y_n\to0$ the maps $S_n(x)=B(x,y_n)$ are pointwise convergent to
zero, hence pointwise bounded. Topological Banach--Steinhaus makes them
equicontinuous. Thus $x_n\to0$ implies $S_nx_n\to0$. At general
$(x,y)$ expand
$B(x_n,y_n)-B(x,y)=B(x_n-x,y_n-y)+B(x_n-x,y)+B(x,y_n-y)$.
If $F$ is metrizable, so is $E\times F$ and sequential continuity
implies continuity: a failure of continuity would choose a violating
point in every member of a countable local base.
\end{proof}

\section{General extension and duality}
\begin{corollary}[An equivalent algebra norm]
\label{fa2:p1:topological-algebra-norm}
If a Banach space has an associative separately continuous bilinear
multiplication, it admits an equivalent submultiplicative norm.  If it
has an identity and the algebra is nonzero, such a norm may be chosen
with $\|1\|=1$. In the zero algebra the identity is zero and has
norm zero; the submultiplicative renorming conclusion still holds.
\end{corollary}
\begin{proof}
Joint continuity, just proved, and scaling give
$\|xy\|\le C\|x\|\,\|y\|$ with $C>0$ (increase the bound if necessary).
The norm $C\|x\|$ is submultiplicative.
In the nonzero unital case use instead $\|x\|_a=\|L_x\|$, where $L_xy=xy$.
Associativity gives $\|xy\|_a\le\|x\|_a\|y\|_a$ and $\|1\|_a=1$;
the inequalities $\|x\|\le\|L_x\|\,\|1\|$ and
$\|L_x\|\le C\|x\|$ prove equivalence. Here $L_1$ is the identity
on a nonzero space, so its operator norm is one. In the zero algebra
every norm takes the value zero at its identity, which proves the
separate zero case.
\end{proof}

\begin{theorem}[Seminorm Hahn--Banach]\label{fa2:p1:hahn-banach}
Let $p$ be a seminorm on a real or complex vector space and let $f$ be a
linear functional on a subspace with $|f|\leq p$. It extends to a linear
$F$ on the whole space satisfying $|F|\leq p$. Every continuous linear
functional on a subspace of an LCS extends continuously to the whole
space. In a Hausdorff LCS the continuous dual separates points.
\end{theorem}
\begin{proof}
For real scalars analytic Hahn--Banach applied to $p$ gives $F\leq p$,
and applying it at $-x$ gives $|F(x)|\leq p(x)$. For complex scalars
extend $\operatorname{Re}f$ on the underlying real subspace and define
$F(x)=u(x)-iu(ix)$. This is complex linear and extends $f$. Choose
$|\lambda|=1$ so that $\lambda F(x)=|F(x)|$; then
$|F(x)|=u(\lambda x)\leq p(\lambda x)=p(x)$.
Continuity on a subspace gives a bound by the maximum of finitely many
continuous ambient seminorms, by the scaling argument in
Theorem~\ref{fa2:p1:seminorm-model}. Extend under that seminorm to get
continuity. If $x\ne0$, a continuous seminorm $p$ has $p(x)>0$;
the functional $\lambda x\mapsto\lambda p(x)$ extends and separates
$x$ from zero.
\end{proof}

\begin{theorem}[Compact--closed separation]\label{fa2:p1:lcs-strong-separation}
Disjoint nonempty convex sets in a TVS can be separated by a nonzero
continuous real linear functional if one is open. In an LCS, if $K$ is
compact convex and $C$ is closed convex, disjoint and nonempty, there is
a continuous real linear $f$ with
$\max_K f<\inf_C f$. For complex spaces use real parts of continuous
complex-linear functionals.
\end{theorem}
\begin{proof}
For the open-set assertion the gauge proof of
Theorem~\ref{fa2:p1:c06:open-separation} only uses a convex open
neighborhood and the continuity of the vector operations; it does not
use a base of convex neighborhoods. For the strict assertion choose
an absolutely convex zero neighborhood $U$ so small that
$(K+U)\cap(C+U)=\varnothing$. To obtain it, for each $k\in K$ choose
a neighborhood whose suitable difference translates miss $C$, and use
a finite subcover of $K$. Open-set separation gives
$\sup f(K+U)\leq\inf f(C+U)$. The nonzero continuous $f$ has
$\sup_U f>0$, so removing the two copies of $U$ leaves a strict gap
between $f(K)$ and $f(C)$. In the complex case every continuous
real-linear $u$ equals $\operatorname{Re}f$ for
$f(x)=u(x)-iu(ix)$.
\end{proof}

\begin{theorem}[Dual-pair bipolar theorem]\label{fa2:p1:general-bipolar}
For a separated dual pair $(E,F)$ define
$A^\circ=\{f\in F:|f(x)|\leq1\ (x\in A)\}$. Then
\[
 A^{\circ\circ}=\overline{\operatorname{aconv}(A)}^{\sigma(E,F)}.
\]
Double annihilators are the corresponding weak closures of linear spans.
Two locally convex topologies with the same continuous dual have the
same closed convex sets.
\end{theorem}
\begin{proof}
Each polar inequality defines a weakly closed absolutely convex set, so
the right side is contained in the left. If $x$ is outside the right
side, separation in the weak topology gives a continuous weak
functional $f\in F$ strictly separating $x$ from that hull. Absolute
convexity makes $\sup\operatorname{Re}f=\sup|f|$ on the hull;
rescaling $f$ gives $\sup_A|f|\leq1<|f(x)|$.
If the supremum is zero, scale by a sufficiently large positive
constant instead. Thus $x\notin A^{\circ\circ}$. Applying the same
separation argument to subspaces proves the annihilator assertion.
Finally a closed convex set is an intersection of closed halfspaces
determined by the continuous dual. The same halfspaces are available
for both compatible topologies, which proves the last statement.
\end{proof}

\begin{theorem}[Alaoglu for neighborhood polars]\label{fa2:p1:alaoglu}
If $U$ is a zero neighborhood in a TVS $E$, its polar $U^\circ\subset E'$
is compact for $\sigma(E',E)$. Every equicontinuous subset of $E'$ is
contained in such a polar. In particular the closed dual unit ball of
every normed space is weak-star compact, without requiring completeness
of its predual.
\end{theorem}
\begin{proof}
For each $x$, absorption gives a finite $M_x$ with
$|f(x)|\leq M_x$ for every $f\in U^\circ$. Embed the polar into
$\prod_{x\in E}\{z:|z|\leq M_x\}$. This product is compact by
Tychonoff. Pointwise limits preserve linear equations and
$|f(u)|\leq1$ for $u\in U$. A linear functional bounded on the
neighborhood $U$ is continuous, by scaling $U$; hence the image is
closed in the product. The product topology is precisely weak-star
topology. Equicontinuity supplies a common neighborhood on which all
the functionals have modulus at most one. For a normed predual take
$U$ to be the open unit ball; its polar is the closed dual unit ball.
\end{proof}

\section{Compact integration and holomorphic maps}
\begin{theorem}\label{fa2:p1:compact-hull-frechet}
The closed convex hull of a compact subset of a Fr\'echet space is compact.
\end{theorem}
\begin{proof}
For an absolutely convex neighborhood $U$, a compact set $K$ has a finite
$U$-net $F$. Then $\operatorname{conv}K\subset\operatorname{conv}F+U$.
The finite-dimensional compact simplex image $\operatorname{conv}F$
has a finite $U$-net. Thus the convex hull is totally bounded for every
zero neighborhood. Its closure remains totally bounded by first using
a smaller neighborhood whose closure lies in $U$. It is complete as
a closed subset of a Fr\'echet space. A complete totally bounded metric
space is compact: any sequence has a diagonal Cauchy subsequence from
successive finite covers, and completeness gives its limit.
\end{proof}

\begin{theorem}[Barycenters and continuous compact integrands]
\label{fa2:p1:barycenters}
Let $Q$ be compact Hausdorff, $\mu$ a regular probability measure, and
$f:Q\to E$ continuous into a TVS whose dual separates points. If
$C=\overline{\operatorname{conv}}f(Q)$ is compact, there is a unique
$b\in C$ satisfying $\ell(b)=\int_Q\ell(f)\,d\mu$ for every $\ell\in E'$.
Every point of the compact closed convex hull of a compact set $K$ is
the barycenter of a regular probability measure on $K$. In a Banach
space $\|b\|\leq\int\|f\|\,d\mu$.
\end{theorem}
\begin{proof}
For finitely many real and imaginary parts of functionals, apply their
joint map into a finite-dimensional real vector space. The integral
of the resulting compact-valued function lies in its compact convex
hull: a separating affine functional would otherwise have integral
strictly larger than its maximum on the range. Hence some finite
convex combination of $f(Q)$ matches the integrals of these finitely
many scalar tests. The corresponding closed subsets of the compact
$C$ have the finite-intersection property, giving $b$; separation gives
uniqueness. Conversely approximate a point of the closed convex hull
by finite convex combinations, represented by atomic probability
measures on $K$. The probability measures form a weak-star compact set
by Alaoglu and the regular-measure representation theorem in Part IV.
A subnet limit represents the desired point, by each scalar test.
For the norm estimate use
$|\ell(b)|\leq\|\ell\|\int\|f\|\,d\mu$ and the norming formula.
\end{proof}

\begin{theorem}[Weak holomorphy in Fr\'echet spaces]
\label{fa2:p1:frechet-holomorphy}
Let $E$ be a complex Fr\'echet space and $f:\Omega\to E$ a map on an
open subset of $\C$. If $\ell\circ f$ is holomorphic for every
$\ell\in E'$, then $f$ is continuous and holomorphic, has locally
convergent Taylor series in every defining seminorm, and satisfies
vector-valued Cauchy formulas. An entire weakly bounded map into any
complex TVS with separating dual is constant.
\end{theorem}
\begin{proof}
Fix a continuous seminorm $p$ and a compact disk $D\subset\Omega$.
The image $f(D)$ is weakly bounded. To see that it is $p$-bounded,
pass to the normed quotient $E/\ker p$ and use its continuous dual;
each of these functionals pulls back to $E'$. Weak boundedness equals
norm boundedness in a normed space, proved by applying uniform
boundedness to evaluation on its Banach dual. Thus
$\sup_Dp(f)<\infty$. Scalar Cauchy estimates on a slightly larger
disk, uniformly for functionals dominated by $p$, give
$p(f(z)-f(w))\leq C_p|z-w|$ on a smaller disk, by Hahn--Banach's
seminorm norming formula. This proves continuity.

Continuous parametrized contour integrands have Riemann integrals
in a Fr\'echet space: their Riemann sums are Cauchy in every defining
seminorm by uniform continuity; completeness supplies a common limit.
Apply scalar Cauchy formulas to these integrals. The separating dual
then gives
\[
 f(z)=\frac1{2\pi i}\int_{|\zeta-z_0|=r}
       \frac{f(\zeta)}{\zeta-z}\,d\zeta.
\]
Expanding the scalar kernel as a geometric series gives coefficients
$a_n=(2\pi i)^{-1}\int f(\zeta)(\zeta-z_0)^{-n-1}\,d\zeta$.
For each $p$, $p(a_n)\leq r^{-n}\sup p(f)$, so the series and its
derivatives converge uniformly in each smaller disk for every
seminorm. Completeness and the integral identity identify the sum
with $f$ and prove holomorphy. Finally, for the last assertion every
$\ell\circ f$ is bounded entire and hence constant by scalar
Liouville. Separation forces $f(z)=f(0)$.
\end{proof}

% END INPUT


% BEGIN parts/part1-foundation-12.tex
\chapter{Weak and Weak-Star Topologies}\label{fa2:p1:c12}

Weak topologies retain the scalar tests of a dual pairing. Their open
sets are controlled by finitely many tests, but convergence and
compactness need not be determined by sequences. Throughout, $E$ is
a real or complex normed space. We use the linear dual $E'$, not the
anti-dual.
\source{B01 Secs. 3.1--3.4}
\cite{B01} 


% FAMILY D16.C12.01
\section{An initial topology from a pairing}
\begin{definition}[Weak topology of a dual pair]\label{fa2:p1:c12:dual-pair}
Let $X,Y$ be vector spaces with a bilinear pairing
$\langle y,x\rangle$ separating points of $X$. The topology
$\sigma(X,Y)$ is the coarsest topology making all maps
$x\mapsto\langle y,x\rangle$, $y\in Y$, continuous. A basis of
neighborhoods of $x_0$ consists of
\[
 \{x:|\langle y_j,x-x_0\rangle|<\eps, 1\leq j\leq m\},
 \quad m<\infty,\quad\eps>0.
\]
Taking arbitrary unions of these finite intersections constructs the
topology. For $X=E$, $Y=E'$ and $\langle f,x\rangle=f(x)$ this is
the weak topology on $E$.
\end{definition}

\begin{proposition}[Hausdorffness and the continuous scalar tests]\label{fa2:p1:c12:initial-properties}
The topology $\sigma(X,Y)$ is Hausdorff and locally convex, and
addition and scalar multiplication are continuous. A linear functional
on $X$ is continuous for $\sigma(X,Y)$ if and only if it belongs to
the linear span of the pairing functionals. In particular the
continuous linear functionals on $(E,\sigma(E,E'))$ are exactly $E'$.
\end{proposition}
\begin{proof}
Distinct points have distinct values under one test; disjoint scalar
neighborhoods pull back to separate them. The basic neighborhoods
of zero are convex and balanced. The triangle inequality makes
addition continuous. For scalar multiplication near $(a,x)$, each
test satisfies
\[
 |\langle y,bz-ax\rangle|
 \leq |b|\,|\langle y,z-x\rangle|+|b-a|\,|\langle y,x\rangle|,
\]
and finitely many tests are controlled by taking $b$ bounded near
$a$ and shrinking the corresponding neighborhoods. This proves
continuity directly.

If a linear $g$ is continuous at zero, some finite-test neighborhood
is contained in $\{|g|<1\}$. If every test $y_1,\ldots,y_m$ vanishes
on $x$, every scalar multiple of $x$ lies in that neighborhood;
hence $g(x)=0$. Thus $g$ factors through the linear map
$V:X\to\K^m$, $Vx=(\langle y_j,x\rangle)_j$.
The induced linear functional on $V(X)$ extends algebraically to
$\K^m$ by extending a finite basis, so it has the form
$\sum_j c_j z_j$. Therefore $g(x)=\sum_jc_j\langle y_j,x\rangle$.
The converse follows from continuity of the tests. Hahn--Banach
makes $E'$ separate points of $E$, so this proof applies to $E$.
\end{proof}

% FAMILY D16.C12.02
\section{A topology belonging to a predual}
\begin{definition}[Weak-star topology]\label{fa2:p1:c12:weak-star}
On $E'$ the weak-star topology $\sigma(E',E)$ is generated by
$f\mapsto f(x)$ for $x\in E$. A basic neighborhood of $f_0$ is
\[
 \{f:|(f-f_0)(x_j)|<\eps, 1\leq j\leq m\}.
\]
The weak topology of the normed space $E'$ itself is
$\sigma(E',E'')$. Thus
\[
 \sigma(E',E)\ \subset\ \sigma(E',E'')\ \subset\
 \text{the norm topology on }E'.
\]
A specified predual means a normed space together with a specified
identification of its dual with the space under study. It supplies
the evaluation tests, and hence the particular weak-star topology.
\end{definition}

\begin{proposition}[When weak and weak-star agree]\label{fa2:p1:c12:weak-star-comparison}
The weak and weak-star topologies on $E'$ agree if and only if
$J_E:E\to E''$ is surjective. Every weak-star continuous linear
functional on $E'$ is evaluation at one $x\in E$.
\end{proposition}
\begin{proof}
Apply the preceding scalar-test characterization to $(E',E)$.
A finite linear combination of evaluations is evaluation at the
corresponding linear combination of vectors. If $J_E$ is onto,
the weak and weak-star test families coincide. Conversely, equality
of topologies makes every member of $E''$ weak-star continuous,
and therefore an evaluation. This proves surjectivity. In particular
on a nonreflexive Banach space the two topologies on its dual are
strictly different.
\end{proof}

% FAMILY D16.C12.03
\section{Nets, sequences and examples}
\begin{proposition}[Evaluation criterion for convergence]\label{fa2:p1:c12:convergence}
A net $(x_\alpha)$ converges to $x$ in $\sigma(E,E')$ if and only if
$f(x_\alpha)\to f(x)$ for every $f\in E'$. A net $(f_\alpha)$
converges weak-star to $f$ if and only if
$f_\alpha(x)\to f(x)$ for every $x\in E$. Norm convergence implies
the corresponding weak convergence. Bounded linear maps are
weak-to-weak continuous, and their transposes are weak-star to
weak-star continuous.
\end{proposition}
\begin{proof}
Convergence implies convergence of every continuous test. Conversely,
given a finite-test neighborhood, directedness supplies one index
beyond which all its finitely many scalar conditions hold.
For norm convergence use $|f(x_\alpha-x)|\leq\norm f\norm{x_\alpha-x}$.
For $T:E\to F$ bounded and $g\in F'$, the composite $gT=T'g$
belongs to $E'$, proving weak continuity. On transposes,
$(T'f)(x)=f(Tx)$ is precisely a weak-star test at $Tx$.
\end{proof}

\begin{example}[Three different notions of convergence]\label{fa2:p1:c12:examples}
The unit vectors $e_n$ of $\ell^2$ converge weakly to zero: Riesz
represents each functional by $y\in\ell^2$ and its values on $e_n$
are $\overline{y_n}\to0$. Their norms are all $1$.
In $\ell^\infty=(\ell^1)'$ let
$u_n(k)=1$ for $k\geq n$ and $0$ for $k<n$. Then
$u_n\to0$ in $\sigma(\ell^\infty,\ell^1)$ because
$|\sum_{k\geq n}a_k|\leq\sum_{k\geq n}|a_k|\to0$.
On the subspace of convergent scalar sequences the limit functional
$L(a)=\lim_k a_k$ is continuous of norm $1$. Hahn--Banach extends
it to $\ell^\infty$, and this extension has $L(u_n)=1$. Thus the
same sequence does not converge weakly to zero in $\ell^\infty$.
Finally in finite dimension all three topologies coincide: finitely
many coordinate functionals control the norm, and the canonical
bidual map is surjective by the finite-dimensional dimension formula.
\end{example}

% FAMILY D16.C12.04
\section{Boundedness and lower semicontinuity}
\begin{theorem}[Bounds for weakly convergent sequences]\label{fa2:p1:c12:bounded-sequences}
If a sequence $x_n$ converges weakly to $x$ in a normed space, then
$\sup_n\norm{x_n}<\infty$ and
$\norm x\leq\liminf_n\norm{x_n}$. If $E$ is Banach and
$f_n\to f$ in $\sigma(E',E)$, then
$\sup_n\norm{f_n}<\infty$ and
$\norm f\leq\liminf_n\norm{f_n}$. Every weakly compact subset
of a normed space is norm bounded.
\end{theorem}
\begin{proof}
The dual $E'$ is Banach even when $E$ is not. The bounded operators
$J_Ex_n:E'\to\K$ are pointwise bounded since each scalar sequence
$g(x_n)$ converges. Uniform boundedness, Chapter~\ref{fa2:p1:c07}, gives
$\sup_n\norm{J_Ex_n}<\infty$. The isometry of $J_E$ identifies
these norms with $\norm{x_n}$. For every unit $g\in E'$,
$|g(x)|=\lim_n|g(x_n)|\leq\liminf_n\norm{x_n}$; taking the
norming supremum gives the lower bound. For weak-star convergence
apply uniform boundedness directly to $f_n:E\to\K$, where the
Banach hypothesis on $E$ is now needed. The same scalar argument
and the definition of the dual norm give its lower bound.
If $K$ is weakly compact, each $g(K)$ is scalar compact and hence
bounded. Apply uniform boundedness to $\{J_Ex:x\in K\}$ on $E'$
and again use isometry.
\end{proof}

\begin{remark}[The completeness and net qualifications]\label{fa2:p1:c12:qualifications}
The Banach assumption for the weak-star sequence bound cannot simply
be omitted. On $c_{00}$ with the $\ell^1$ norm,
$f_n(x)=n x_n$ is continuous, has norm $n$, and converges pointwise
to zero, because each $x$ has finite support. Also a convergent
general net need not be bounded as a whole: its indices outside a
chosen final segment may contain arbitrarily large values. The
preceding sequence proof uses boundedness of every entire scalar
sequence; it is not a statement that arbitrary weakly convergent
nets are globally bounded.
\end{remark}

\begin{proposition}[Norm is weakly lower semicontinuous]\label{fa2:p1:c12:norm-lsc}
Each closed norm ball is weakly closed; each closed dual norm ball
is weak-star closed. More generally, a norm lower semicontinuous
convex function $\Phi:E\to(-\infty,+\infty]$ is weakly lower
semicontinuous.
\end{proposition}
\begin{proof}
The norming formula writes
$\{x:\norm x\leq r\}=\bigcap_{\norm f\leq1}\{x:|f(x)|\leq r\}$
for $r\geq0$. The dual ball is
$\bigcap_{x\in E}\{f:|f(x)|\leq r\norm x\}$.
These are intersections of closed scalar-test sets. For the last
claim, each sublevel set of $\Phi$ is norm closed and convex;
the next theorem makes it weakly closed, exactly the definition
of weak lower semicontinuity.
\end{proof}

% FAMILY D16.C12.05
\section{Convex closures and Mazur's lemma}
\begin{theorem}[Weak and norm closures of a convex set]\label{fa2:p1:c12:convex-closure}
For every convex subset $C$ of a normed space,
$\cl^{\sigma(E,E')}C=\cl^{\norm\cdot}C$.
In particular a norm closed linear subspace is weakly closed, and
the weak topology induced on a subspace $M$ equals
$\sigma(M,M')$ with its inherited norm.
\end{theorem}
\begin{proof}
The weak topology is coarser, so the norm closure is contained in
the weak closure. If $x$ is outside the norm closure and $C$ is
nonempty, strict separation of a point and a closed convex set
from Chapter~\ref{fa2:p1:c06} gives $f\in E'$ and a real number $a$
such that
$\operatorname{Re}f(x)>a\geq\sup_{z\in C}\operatorname{Re}f(z)$.
The weakly open half-space $\{y:\operatorname{Re}f(y)>a\}$
excludes $C$, so $x$ is outside the weak closure. The empty set
case is immediate. In the complex case this use of separation
is valid because every continuous real-linear functional $u$ is
$\operatorname{Re}f$ for the continuous complex-linear functional
$f(x)=u(x)-iu(ix)$.
For the subspace topology, restrictions of $E'$ belong to $M'$,
and Hahn--Banach extends every member of $M'$ to $E'$. Thus the
two families of scalar tests on $M$ coincide.
\end{proof}

\begin{corollary}[Mazur's lemma with tail combinations]\label{fa2:p1:c12:mazur}
If $x_n\to x$ weakly, there are integers $m_n\geq n$ and real
numbers $a_{n,j}\geq0$, $n\leq j\leq m_n$, with
$\sum_{j=n}^{m_n}a_{n,j}=1$, such that
$y_n=\sum_{j=n}^{m_n}a_{n,j}x_j$ converges to $x$ in norm.
\end{corollary}
\begin{proof}
For each $n$, the weak limit $x$ is in the weak closure of
$\{x_j:j\geq n\}$, hence in the weak closure of its convex hull.
By the theorem it is in that convex hull's norm closure. Choose
a finite convex combination at distance less than $1/n$; inserting
zero coefficients fills any skipped indices. This gives the stated
sequence and its norm convergence.
\end{proof}

% END INPUT


% BEGIN parts/part1-foundation-13.tex
\chapter{Banach--Alaoglu and Weak-Star Metrizability}\label{fa2:p1:c13}

The scalar field is $\R$ or $\C$, and $E$ is normed. No completeness
of $E$ is needed for the dual-ball compactness theorem. Compactness
here is topological compactness; subsequences enter only after a
metric has been justified.
\source{B01 Secs. 3.4, 3.6}
\cite{B01} 


% FAMILY D16.C13.01
\section{A compact set of pointwise data}
\begin{lemma}[Compact products of scalar disks]\label{fa2:p1:c13:scalar-product-compactness}
If $(D_i)_{i\in I}$ are closed bounded scalar disks centered at
zero, their product is compact in its product topology.
\end{lemma}
\begin{proof}
We include the needed product-compactness foundation. A proper
filter on a set is a family of subsets containing the whole set,
closed under finite intersections and supersets, and excluding
the empty set. Zorn's lemma extends every proper filter to a
maximal one, called an ultrafilter: the union of a chain of
proper filters is again proper. Maximality implies that for
each subset $A$ either $A$ or its complement belongs to the
ultrafilter. Indeed, if adjoining $A$ cannot generate a proper
filter, some existing filter member is disjoint from $A$, so
the complement already belongs to the filter.

Put $P=\prod_iD_i$, which is nonempty since all disks contain
zero. Given an ultrafilter $\mathcal U$ on $P$, its coordinate
image is the ultrafilter
$\mathcal U_i=\{A\subset D_i:\pi_i^{-1}(A)\in\mathcal U\}$.
Each $D_i$ is compact by finite-dimensional compactness
(Chapter~\ref{fa2:p1:c01}). The
closed sets $\cl A$, $A\in\mathcal U_i$, have the
finite-intersection property, so have a common point $z_i$.
Every open neighborhood $V$ of $z_i$ belongs to $\mathcal U_i$:
otherwise its closed complement would belong and have closure
containing $z_i$, a contradiction. Thus $\mathcal U_i$
converges to $z_i$. The limit is unique since $D_i$ is Hausdorff,
so these points specify $z=(z_i)\in P$. Every basic product
neighborhood of $z$ restricts only finitely many coordinates;
its inverse-image factors all belong to $\mathcal U$, hence so
does their intersection. Therefore $\mathcal U$ converges to $z$.

Finally any family of closed subsets of $P$ with the
finite-intersection property generates a proper filter. Extend
it to an ultrafilter and take its limit as above. That limit
lies in every one of the closed sets, since otherwise its open
complement would also belong to the ultrafilter. Thus every
such family has nonempty intersection, which, by taking
complements, is exactly compactness of $P$.
\end{proof}

\begin{theorem}[Banach--Alaoglu]\label{fa2:p1:c13:alaoglu}
The closed unit ball $B_{E'}=\{f\in E':\norm f\leq1\}$ is compact
in $\sigma(E',E)$. Every closed dual ball of finite radius is
weak-star compact.
\end{theorem}
\begin{proof}
For $x\in E$ let $D_x=\{z\in\K:|z|\leq\norm x\}$, a compact
real interval or complex disk, including $D_0=\{0\}$. The product
$P=\prod_{x\in E}D_x$ is compact by the preceding lemma. Define
\[
 L=\{\omega\in P:
 \omega_{x+y}-\omega_x-\omega_y=0\ (x,y\in E),\quad
 \omega_{ax}-a\omega_x=0\ (a\in\K,x\in E)\}.
\]
For each fixed $x,y,a$, the relevant expression is a continuous
map from $P$ to $\K$, being a finite linear combination of coordinate
projections. Its zero set is closed. The set $L$ is the intersection
of all these closed zero sets, including the constraints for every
complex scalar in the complex case. Thus $L$ is genuinely closed
in the compact product, and hence compact.

For $\omega\in L$ the function $f_\omega(x)=\omega_x$ is additive
and $\K$-homogeneous. The defining disks give
$|f_\omega(x)|\leq\norm x$, so it is continuous and has norm at
most $1$. Conversely every element of $B_{E'}$ satisfies precisely
these disk and linearity constraints. Consequently
$\iota:f\mapsto(f(x))_{x\in E}$ is a bijection from $B_{E'}$ to
$L$. Each coordinate of $\iota$ is a weak-star test, and every
basic weak-star neighborhood pulls back from finitely many product
coordinates. These two facts prove that $\iota$ and its inverse
are continuous. Thus $B_{E'}$ is homeomorphic to $L$, proving
compactness. Multiplication by a positive finite radius is a
weak-star homeomorphism; radius zero gives a singleton. Translation
gives balls with arbitrary centers.
\end{proof}

\begin{remark}[What product compactness supplies]\label{fa2:p1:c13:product-dependency}
The scalar-product lemma proves the precise product-compactness
case needed here from Zorn's lemma and finite-dimensional
compactness. Linearity and the norm bound have also been checked
here. The proof is valid for
incomplete $E$ because scalar boundedness already makes the
resulting linear function continuous.
\end{remark}

% FAMILY D16.C13.02
\section{A metric on a bounded part of a dual}
\begin{theorem}[Weak-star metrization for a separable predual]\label{fa2:p1:c13:metric}
If $E$ is separable, choose a sequence $(x_n)$ dense in its closed
unit ball. On any norm-bounded subset $A\subset E'$, the metric
\[
 d(f,g)=\sum_{n\geq1}2^{-n}\min\{1,|(f-g)(x_n)|\}
\]
induces the relative weak-star topology. When $E=\{0\}$ the dual
is a singleton, and the zero distance on that singleton is a metric.
\end{theorem}
\begin{proof}
The summands are bounded by $2^{-n}$; positivity, symmetry and the
triangle inequality follow from those of $\min(1,|s-t|)$ on
scalars. If $d(f,g)=0$, the functionals agree on the dense set
$(x_n)$; continuity and scaling give equality on all of $E$.

Each coordinate evaluation is continuous for $d$: if its difference
is to be less than $0<\eps<1$, require $d(f,g)<2^{-n}\eps$.
Suppose $\norm f,\norm g\leq R$ on $A$. For $x\ne0$ choose
$x_n$ close to $x/\norm x$. Then
\[
 |(f-g)(x)|\leq\norm x\bigl(2R\norm{x/\norm x-x_n}
 +|(f-g)(x_n)|\bigr).
\]
The first term is uniformly small on $A$, and a $d$-neighborhood
controls the second. Thus every evaluation at $x$ is $d$-continuous
on $A$ (the cases $R=0$ or $x=0$ are immediate), proving that
the relative weak-star topology is no finer than the metric one.
Conversely, to make $d(f,g)<\eps$, first choose $N$ with
$\sum_{n>N}2^{-n}<\eps/2$, then require finitely many evaluations
at $x_1,\ldots,x_N$ to differ by less than $\eps/2$.
This is a weak-star neighborhood contained in the metric ball,
proving the reverse inclusion. Boundedness was used precisely
to pass from the dense tests to all evaluations.
\end{proof}

% FAMILY D16.C13.03
\begin{corollary}[Weak-star subsequences]\label{fa2:p1:c13:subsequences}
If $E$ is separable and $(f_n)$ is norm bounded in $E'$, some
subsequence converges in $\sigma(E',E)$ to an element of $E'$.
\end{corollary}
\begin{proof}
All terms lie in a fixed closed dual ball. Alaoglu makes this ball
compact and the preceding theorem makes it metrizable. A compact
metric space is sequentially compact: recursively pass to infinitely
many terms in one member of a finite cover by balls of radii
$2^{-k}$, and diagonalize to obtain a Cauchy subsequence; a cluster
point supplied by compactness is its limit. Applied to the metric
above, this is exactly weak-star convergence.
\end{proof}

% FAMILY D16.C13.04
\section{Transferring topology through the bidual embedding}
\begin{proposition}[The weak topology as an induced weak-star topology]\label{fa2:p1:c13:embedding}
The canonical isometry $J_E:E\to E''$ is a homeomorphism from
$(E,\sigma(E,E'))$ onto its image with the topology induced by
$\sigma(E'',E')$. In particular $x_\alpha\to x$ weakly if and
only if $J_Ex_\alpha\to J_Ex$ weak-star in $E''$.
\end{proposition}
\begin{proof}
The coordinate at $f\in E'$ of $J_Ex$ is $J_Ex(f)=f(x)$.
Thus the two initial topologies have precisely the same scalar
tests and the same finite-test neighborhoods. Hahn--Banach makes
$J_E$ injective (indeed isometric), so this equality of topologies
is a homeomorphism onto its image.
\end{proof}

% FAMILY D16.C13.05
\section{Limits of the metric conclusion}
\begin{proposition}[The dual ball detects predual separability]\label{fa2:p1:c13:converse-metric}
For any normed space $E$, if $B_{E'}$ is weak-star metrizable, then
$E$ is separable. Hence $B_{E'}$ is weak-star metrizable if and
only if $E$ is separable.
\end{proposition}
\begin{proof}
Let $d$ be such a metric and, for each $n$, choose a relative
basic weak-star neighborhood $V_n$ of zero contained in the
$d$-ball of radius $1/n$. Let $D$ be the countable union of the
finite sets of vectors used by these neighborhoods. If $f\in E'$
vanishes on $D$, then, after normalizing $f$ to the unit ball if
necessary, it lies in every $V_n$, so its distance to zero is zero.
Thus only the zero continuous functional annihilates $D$.
Hahn--Banach separation implies $\clspan D=E$. Finite rational
linear combinations (complex rational combinations over $\C$)
of the countable set $D$ are countable and dense in $E$.
\end{proof}

\begin{example}[Nonmetrizable compactness and two separabilities]\label{fa2:p1:c13:nonmetrizable}
For an uncountable set $I$, the space $\ell^2(I)$ is nonseparable
by its uncountable orthonormal system. Its dual unit ball is
weak-star compact by Alaoglu and is not weak-star metrizable by
the proposition. Thus the compactness proof cannot in general
be read as a metric subsequence argument. Separability of $E$
and norm separability of $E'$ are different: $\ell^1$ is separable,
but $(\ell^1)'=\ell^\infty$ is nonseparable. Indeed its uncountable
family of $\{0,1\}$-valued sequences has pairwise distance $1$.
Nevertheless the closed unit ball of $\ell^\infty$ is metrizable
for its specified weak-star topology $\sigma(\ell^\infty,\ell^1)$.
\end{example}

\begin{remark}[Why boundedness cannot be dropped]\label{fa2:p1:c13:whole-dual}
If $E$ is an infinite-dimensional Banach space, the weak-star topology on all of
$E'$ is not first countable. To prove this, suppose a countable
local base at zero existed and refine each of its neighborhoods
to a finite-test neighborhood. Let $D$ be the resulting countable
collection of test vectors. For any $x\in E$, the neighborhood
$\{f:|f(x)|<1\}$ contains one of these finite-test neighborhoods.
Every functional annihilating that finite set must therefore
annihilate $x$, since all its scalar multiples stay in the
neighborhood. Hahn--Banach applied to its finite-dimensional
closed span implies $x$ belongs to that span. Consequently $E$
is the algebraic span of $D$. Baire contradicts this description:
an infinite-dimensional Banach space cannot be a countable union
of its finite-dimensional, closed, empty-interior subspaces.
Thus the topology is not metrizable. The Banach hypothesis in
this whole-space assertion matters: on $c_{00}$ the coordinate
tests already generate its dual's weak-star topology, which is
a metrizable product subspace topology. The bounded-set theorem
does not claim a metric for the entire dual of a Banach space.
\end{remark}

% END INPUT


% BEGIN parts/part1-foundation-14.tex
\chapter{Reflexivity, Weak Compactness and Uniform Convexity}\label{fa2:p1:c14}

The field is $\R$ or $\C$. Reflexivity refers to the canonical map
$J_Ex(f)=f(x)$, whose isometry was proved using Hahn--Banach in
Chapter~\ref{fa2:p1:c05}. It does not mean merely the existence of some
abstract isomorphism with the bidual.
\source{B01 Secs. 3.5--3.7, 4.3}
\cite{B01} 


% FAMILY D16.C14.01
\section{The unit ball and the bidual}
\begin{definition}[Reflexive space]\label{fa2:p1:c14:reflexive}
A normed space $E$ is reflexive when $J_E:E\to E''$ is surjective.
It is then automatically Banach, since it is isometric to the
Banach space $E''$. In the compactness characterizations below we
state the Banach hypothesis explicitly.
\end{definition}

\begin{theorem}[Goldstine's density theorem]\label{fa2:p1:c14:goldstine}\label{fa2:p1:goldstine}
For every normed space $E$, $J_E(B_E)$ is dense in $B_{E''}$ for
$\sigma(E'',E')$ relative to that ball.
\end{theorem}
\begin{proof}
Fix $\xi\in B_{E''}$ and finitely many $f_1,\ldots,f_m\in E'$.
Put $Vx=(f_1(x),\ldots,f_m(x))\in\K^m$ and
$a=(\xi(f_1),\ldots,\xi(f_m))$. We claim $a\in\cl V(B_E)$.
The latter set is closed, convex and bounded in finite-dimensional
real coordinates. If the claim fails, strict separation gives
scalars $c_j$ such that
\[
 \operatorname{Re}\sum_jc_j\xi(f_j)
 >\sup_{x\in B_E}\operatorname{Re}\sum_jc_jf_j(x)
 =\left\|\sum_jc_jf_j\right\|.
\]
Over $\C$, every real-linear functional on $\C^m$ has the displayed
real-part form; the equality on the right follows by rotating
$x$ by scalars of modulus one. But the left side is
$\operatorname{Re}\xi(\sum_jc_jf_j)\leq\norm\xi\norm{\sum_jc_jf_j}$,
a contradiction. Thus, for every $\eps>0$, some $x\in B_E$
satisfies $|f_j(x)-\xi(f_j)|<\eps$ for all $j$. These are exactly
the finite-test neighborhoods defining the asserted density.
No completeness or reflexivity was used.
\end{proof}

\begin{theorem}[Reflexivity and the weakly compact ball]\label{fa2:p1:c14:ball-characterization}
For a Banach space $E$, the following are equivalent: $E$ is
reflexive; $B_E$ is weakly compact; and every norm bounded weakly
closed subset of $E$ is weakly compact.
\end{theorem}
\begin{proof}
If $J_E$ is onto, its isometry identifies $B_E$ with $B_{E''}$.
The topology-transfer result of Chapter~\ref{fa2:p1:c13} identifies
the weak topology on the first with $\sigma(E'',E')$ on the
second. Alaoglu gives compactness. A bounded weakly closed set
lies as a closed subset in a suitable multiple of this compact
ball, proving the third statement. That statement implies the
second because $B_E$ is weakly closed. Conversely, if $B_E$ is
weakly compact, $J_E(B_E)$ is weak-star compact in the Hausdorff
space $E''$, hence weak-star closed. Goldstine says it is also
dense in $B_{E''}$. It therefore equals that ball; scaling gives
surjectivity of $J_E$.
\end{proof}

% FAMILY D16.C14.02
\section{Stability under subspaces, duals and quotients}
\begin{proposition}[Reflexivity is preserved by Banach isomorphisms]\label{fa2:p1:c14:isomorphism-invariance}
If $T:E\to F$ is a bounded bijection of Banach spaces with bounded
inverse, then $E$ is reflexive if and only if $F$ is. In particular
equivalent Banach norms have the same reflexivity property.
\end{proposition}
\begin{proof}
The transpose calculus gives $J_FT=T''J_E$, and $T''$ is invertible
with inverse $(T^{-1})''$. If $J_E$ is onto, this identity gives
$F''=T''(E'')=T''J_E(E)=J_FT(E)=J_F(F)$. Exchange $E,F$ for
the converse. Equivalent norms make the identity a Banach
isomorphism, so this special case follows.
\end{proof}

\begin{theorem}[Subspace, quotient and dual stability]\label{fa2:p1:c14:stability}
Let $E$ be Banach and $M\subset E$ a closed linear subspace.
If $E$ is reflexive, then $M$ and $E/M$ are reflexive. Also
$E$ is reflexive if and only if $E'$ is reflexive.
\end{theorem}
\begin{proof}
If $E$ is reflexive, $B_M=B_E\cap M$ is weakly closed in the
weakly compact ball: $M$ is norm closed and convex, hence weakly
closed. Hahn--Banach identifies its induced weak topology with
$\sigma(M,M')$. Thus $B_M$ is weakly compact, and $M$ is Banach
by closedness, so the ball characterization proves its reflexivity.

If $J_E$ is onto and $\Phi\in E'''$, define
$f(x)=\Phi(J_Ex)$. This belongs to $E'$. For every
$\xi=J_Ex\in E''$,
$\Phi(\xi)=f(x)=\xi(f)=J_{E'}f(\xi)$; surjectivity of $J_E$
therefore gives $\Phi=J_{E'}f$. Hence $E'$ is reflexive.
Conversely if $E'$ is reflexive, the implication just proved
makes $E''$ reflexive. The isometric image $J_E(E)$ is norm
closed because $E$ is Banach. The subspace argument makes it
reflexive; isomorphism invariance transfers that property to $E$.

Finally $(E/M)'$ is isometrically the annihilator $M^\perp\subset E'$
by the quotient duality theorem of Chapter~\ref{fa2:p1:c05}.
That annihilator is norm closed and thus reflexive when $E$ is
reflexive. The quotient $E/M$ is Banach by Chapter~\ref{fa2:p1:c02}.
Applying the dual equivalence to this quotient makes $E/M$
reflexive. This route does not assume in advance that the quotient
map sends the closed unit ball onto the closed quotient unit ball.
\end{proof}

% FAMILY D16.C14.03
\section{A Banach-space subsequence argument}
\begin{lemma}[A separable dual forces a separable space]\label{fa2:p1:c14:dual-separability}
If the continuous dual of a normed space $X$ is norm separable,
then $X$ is separable. Consequently a separable reflexive Banach
space has a norm-separable dual.
\end{lemma}
\begin{proof}
Choose a norm-dense sequence $f_n$ in $X'$. If $f_n\ne0$, choose
$x_n\in B_X$ with $|f_n(x_n)|\geq\norm{f_n}/2$; otherwise put
$x_n=0$. Let $M=\clspan\{x_n\}$. A functional $f$ annihilating
$M$ satisfies, whenever $\norm{f-f_n}<\eps$,
\[
 \tfrac12\norm{f_n}\leq|f_n(x_n)|
 =|(f_n-f)(x_n)|\leq\eps,
 \qquad\norm f\leq3\eps.
\]
Density permits arbitrarily small $\eps$, so $f=0$.
Hahn--Banach implies $M=X$. Rational linear combinations of the
$x_n$ give a countable dense set. For the consequence, if $X$
is reflexive and separable then $X''=J_X(X)$ is separable. Apply
the first assertion to the normed space $X'$, whose dual is $X''$.
\end{proof}

\begin{theorem}[Weak subsequences in a reflexive Banach space]\label{fa2:p1:c14:weak-subsequence}
Every norm-bounded sequence in a reflexive Banach space has a
weakly convergent subsequence. No separability assumption on the
original space is needed.
\end{theorem}
\begin{proof}
For a bounded sequence $(x_n)$ put
$M=\clspan\{x_n:n\in\N\}$. This is separable, since its finite
rational combinations are dense, and is a closed subspace of the
reflexive Banach space $E$. By stability it is reflexive, and the
lemma makes $M'$ norm separable. Alaoglu and the metrization
theorem of Chapter~\ref{fa2:p1:c13}, applied to predual $M'$, make
$B_{M''}$ weak-star compact and metrizable. The onto isometry
$J_M$ transfers this to weak compact metrizability of $B_M$.
The sequence lies in a finite multiple of that ball, so metric
sequential compactness supplies a subsequence converging weakly
in $M$ to some $x\in M$. For every $f\in E'$, its restriction
belongs to $M'$, hence $f(x_{n_k})\to f(x)$. This is weak
convergence in $E$. The proof uses separability of the sequence's
closed span, rather than any assertion that all weakly compact
topological spaces are sequentially compact. The full
Eberlein--Smulian equivalence is a further theorem and is not
needed here.
\end{proof}

% FAMILY D16.C14.04
\section{Uniform convexity without a circular compactness argument}
\begin{definition}[Uniformly convex norm]\label{fa2:p1:c14:uniform-convexity}
A normed space is uniformly convex if for each $0<\eps\leq2$
there is $\delta>0$ such that
\[
 \norm x,\norm y\leq1,\quad\norm{x-y}\geq\eps
 \quad\Longrightarrow\quad
 \left\|\frac{x+y}{2}\right\|\leq1-\delta.
\]
The condition is vacuous when the indicated pairs do not exist.
It implies strict convexity of the unit sphere.
\end{definition}

\begin{theorem}[Milman--Pettis]\label{fa2:p1:c14:milman-pettis}
Every uniformly convex Banach space is reflexive.
\end{theorem}
\begin{proof}
It suffices to put each $\xi\in E''$ of norm $1$ in $J_E(E)$;
zero and other norms follow by scaling. The image $J_E(E)$ is
norm closed because $E$ is complete. Fix $0<\eps\leq2$ and let
$\delta$ be the uniform convexity constant. By the definition of
$\norm\xi$, choose $f\in E'$ with $\norm f=1$ and, after a
scalar rotation if necessary,
$\operatorname{Re}\xi(f)>1-\delta/2$. The slice
\[
 S=\{x\in B_E:\operatorname{Re}f(x)>1-\delta\}
\]
is nonempty by Goldstine: every weak-star neighborhood of $\xi$
small enough in the $f$ coordinate meets $J_E(B_E)$. In fact
$\xi\in\cl^{\sigma(E'',E')}J_E(S)$, since any neighborhood of
$\xi$ may be intersected with the same open $f$-coordinate
condition and then Goldstine applied.

If $x,y\in S$, then
$\norm{(x+y)/2}\geq\operatorname{Re}f((x+y)/2)>1-\delta$,
so uniform convexity implies $\norm{x-y}<\eps$. Fix one $x\in S$.
It follows that $J_E(S)\subset J_Ex+\eps B_{E''}$.
This last ball is weak-star closed, being the intersection of
its scalar norm-bound constraints as in Chapter~\ref{fa2:p1:c12}.
Taking weak-star closures gives $\norm{\xi-J_Ex}\leq\eps$.
Since $\eps$ is arbitrary and $J_E(E)$ is norm closed, $\xi$
belongs to $J_E(E)$. Neither reflexivity nor existence of weak
subsequences was assumed in this proof: only Goldstine, uniform
convexity, and completeness of the norm-closed image were used.
\end{proof}

\begin{proposition}[Hilbert and $L^p$ uniform convexity]\label{fa2:p1:c14:uc-examples}
Every Hilbert space is uniformly convex, with
$\delta=1-\sqrt{1-\eps^2/4}$ for $0<\eps\leq2$.
For every positive measure $\mu$ and $1<p<\infty$,
$L^p(\mu,\K)$ is uniformly convex. The norms of $\ell^1$ and
$\ell^\infty$ are not uniformly convex.
\end{proposition}
\begin{proof}
The Hilbert assertion follows from
$\norm{(x+y)/2}^2=(\norm x^2+\norm y^2)/2-\norm{x-y}^2/4$.
We give an elementary argument for all interior $p$ at once.
The scalar function $a\mapsto|a|^p$ is strictly convex on
$\K$ viewed as a real vector space: the triangle inequality
and strict convexity of $t^p$ give strictness unless the two
scalars have the same modulus and phase, hence are equal.
For any $\eta>0$, compactness of the scalar set
\[
 \{(a,b): (|a|^p+|b|^p)/2=1,\ |a-b|\geq\eta\}
\]
therefore gives $\gamma_\eta>0$ such that
$|(a+b)/2|^p\leq1-\gamma_\eta$ on it. If it is empty,
choose any positive $\gamma_\eta\leq1$; otherwise the positive
gap has a positive minimum.

Take $u,v$ in the $L^p$ unit ball with $\norm{u-v}_p\geq\eps$.
Write $s=(|u|^p+|v|^p)/2$, $d=|u-v|^p$ and
$A=\{d\geq(\eps^p/2)s,\ s>0\}$. These are measurable objects
independent of the chosen a.e. representatives. Pointwise
$d\leq2^ps$, $\int s\leq1$, and
\[
 \int_A d\geq\int d-\frac{\eps^p}{2}\int_{A^c}s
 \geq\frac{\eps^p}{2},\qquad
 \int_A s\geq\frac{\eps^p}{2^{p+1}}.
\]
On $A$, divide $u,v$ by $s^{1/p}$ and apply the scalar estimate
with $\eta=\eps/2^{1/p}$. Off $A$ use ordinary convexity.
Integration gives
\[
 \left\|\frac{u+v}{2}\right\|_p^p
 \leq\int s-\gamma_\eta\int_A s
 \leq1-\gamma_\eta\frac{\eps^p}{2^{p+1}}<1.
\]
Taking a $p$th root supplies a uniform $\delta>0$ depending only
on $p,\eps$. In $\ell^1$, use $u=e_1,v=e_2$; their midpoint
has norm $1$ despite distance $2$. In $\ell^\infty$ use
$u=(1,1,0,\ldots)$ and $v=(1,-1,0,\ldots)$.
\end{proof}

\begin{proposition}[Weak convergence and convergence of norms]\label{fa2:p1:c14:radon-riesz}
In a uniformly convex normed space, if $x_n\to x$ weakly and
$\limsup_n\norm{x_n}\leq\norm x$, then $x_n\to x$ in norm.
\end{proposition}
\begin{proof}
Weak lower semicontinuity gives $\norm{x_n}\to\norm x$.
If $x=0$ this proves the assertion. Otherwise set
$r_n=\max(\norm{x_n},\norm x)$, $u_n=x_n/r_n$ and
$u=x/\norm x$. Then $r_n\to\norm x$, $\norm{u_n}\leq1$,
$\norm u=1$, and $u_n\to u$ weakly. Hence
$1\leq\liminf\norm{(u_n+u)/2}\leq1$.
Uniform convexity forces $\norm{u_n-u}\to0$: otherwise a
subsequence at distance at least some $\eps$ would have
midpoints of norm at most $1-\delta$. Rescaling gives the result.
\end{proof}

% FAMILY D16.C14.05
\section{Examples and the direct method of minimization}
\begin{corollary}[Reflexive and nonreflexive examples]\label{fa2:p1:c14:reflexive-examples}
Every Hilbert space and every $L^p(\mu,\K)$ with $1<p<\infty$
and arbitrary positive $\mu$ is reflexive. The spaces
$c_0$, $\ell^1$ and $\ell^\infty$ are not reflexive.
\end{corollary}
\begin{proof}
The positive assertions follow from completeness (for $L^p$, see
Theorem~\ref{fa2:p4:lp-completeness}), the uniform
convexity just proved, and Milman--Pettis. The canonical map in
the $L^p$ case can also be checked explicitly using
Theorem~\ref{fa2:p4:lp-duality}: write
$U_qg(f)=\int fg\,d\mu$, a linear isometric bijection
$L^q\to(L^p)'$. For $\xi\in(L^p)''$, the functional
$g\mapsto\xi(U_qg)$ on $L^q$ has, by the same theorem, a
density $h\in L^p$. Thus
$\xi(U_qg)=\int gh\,d\mu=(U_qg)(h)$ for every $g$.
Surjectivity of $U_q$ gives $\xi=J_{L^p}h$. This uses the
arbitrary-measure interior-exponent version twice and adds no
conjugation to its bilinear pairing.

For $c_0$, the sequence-duality theorem identifies its bidual with
$\ell^\infty$ and its canonical image with the sequences tending
to zero. The constant sequence $1$ is outside that image.
For $\ell^1$, a weakly convergent subsequence of the unit vectors
would have every coordinate of its limit zero, yet the continuous
functional $x\mapsto\sum_nx_n$ would have value $1$ along the
subsequence. This is impossible; the weak-subsequence theorem
therefore rules out reflexivity. Finally $c_0$ is a closed subspace
of $\ell^\infty$; subspace stability rules out reflexivity of
$\ell^\infty$. These endpoint examples concern counting measure;
one must not infer that every $L^1$ or $L^\infty$ space is
nonreflexive, since finite atomic examples are finite dimensional.
\end{proof}

\begin{theorem}[Attainment by weak compactness]\label{fa2:p1:c14:minimization}
Let $E$ be reflexive Banach, $C\subset E$ nonempty and weakly
closed, and $\Phi:C\to(-\infty,+\infty]$ proper and weakly
lower semicontinuous. Suppose its sublevel sets are norm bounded;
it suffices that $\Phi(x)\to+\infty$ as $x\in C$,
$\norm x\to\infty$, or that $C$ is bounded. Then $\Phi$
attains a finite minimum. In particular the conclusion applies
when $C$ is norm closed and convex and $\Phi$ is the restriction
of a norm lower semicontinuous convex function on $E$ with these
bounded-sublevel hypotheses.
\end{theorem}
\begin{proof}
Choose $a\in C$ with $\Phi(a)<\infty$. The set
$K=\{x\in C:\Phi(x)\leq\Phi(a)\}$ is nonempty, bounded and
weakly closed, hence weakly compact by the ball characterization.
The infimum on $K$ cannot be $-\infty$: otherwise the nonempty
closed sets $K\cap\{\Phi\leq-n\}$, for all sufficiently large
integers $n$, would be nested in a compact set and have a common
point, contradicting the exclusion of the value $-\infty$.
If $m=\inf_K\Phi$ is finite, the nonempty closed sets
$K\cap\{\Phi\leq m+1/n\}$ have the finite-intersection property;
compactness gives a common point with value $m$. Points outside
$K$ have larger value than $\Phi(a)$, so it is a global minimizer
on $C$. The last assertion follows from weak closedness of norm
closed convex sets and weak lower semicontinuity of norm-lower-
semicontinuous convex functions in Chapter~\ref{fa2:p1:c12}.
\end{proof}

\begin{corollary}[Norm attainment and best approximation]\label{fa2:p1:c14:attainment-examples}
Every continuous linear functional on a reflexive Banach space
attains its norm on the closed unit ball. In a uniformly convex
Banach space every nonempty closed convex set has a unique
nearest point to each specified vector.
\end{corollary}
\begin{proof}
The real part of a functional is weakly continuous on the weakly
compact ball; it attains a maximum, whose value is its norm by
the scalar rotation identity. For nearest points apply the
minimization theorem to $\Phi(z)=\norm{x-z}$; its sublevel sets
are bounded and it is convex and norm continuous. Uniform
convexity implies strict convexity, so two distinct minimizing
residuals of the same positive norm would have midpoint of
strictly smaller norm. If that norm is zero, the minimizer must
be $x$ itself. This proves uniqueness in all cases.
\end{proof}

% END INPUT


% BEGIN parts/part1-weak-advanced.tex
\chapter{Weak compactness, bounded dual topologies, and support}
\label{fa2:p1:weak-advanced}

Weak topologies are usually not metrizable.  In a Banach space they nevertheless
have a remarkably strong sequential compactness theory.  Completeness enters
both the convex-hull theorem and the bounded-dual theorem below; neither is a
formal fact about arbitrary locally convex spaces.  Our sources for these
results are \cite[Chapters V--VI]{B02}, \cite[Chapter 3]{B04}, and
\cite[Sections 2.7--2.9]{B05}.

\section{Sequences and weak compactness}

\begin{lemma}[A countable test family on a separable subspace]
\label{fa2:p1:separable-weak-compact}
Let $Y$ be a separable subspace of a normed space $E$.  There is a sequence
$(f_j)$ in $B_{E'}$ which separates points of $Y$.  On every weakly compact
subset of $Y$ the weak topology is metrizable.  A weakly convergent sequence
in $E$ is norm bounded.
\end{lemma}
\begin{proof}
Choose a dense sequence $(u_j)$ in the unit sphere of $Y$ and, by Hahn--Banach,
$f_j\in B_{E'}$ with $f_j(u_j)=1$.  If $y\ne0$, choose $u_j$ within $1/2$
of $y/\|y\|$; then $|f_j(y)|>\|y\|/2$.  Thus the countable family separates.
The map $y\mapsto(f_j(y))_j$ is weakly continuous and injective into a
countable product of scalar spaces.  On a compact subset it is a
homeomorphism onto its image, since a continuous bijection from a compact
space to a Hausdorff space has continuous inverse.  The image is metrizable,
for example by the metric $\sum_j2^{-j}\min(1,|s_j-t_j|)$.
Finally, if $x_n\to x$ weakly, the functionals $J_E x_n$ on the Banach
space $E'$ are pointwise bounded.  Uniform boundedness gives
$\sup_n\|x_n\|<\infty$.
\end{proof}

\begin{theorem}[Eberlein--\v Smulian]
\label{fa2:p1:eberlein-smulian}
For a subset $A$ of a normed space $E$, the following are equivalent:
\begin{enumerate}
\item its weak closure is weakly compact;
\item every sequence in $A$ has a subsequence weakly convergent in $E$;
\item every sequence in $A$ has a weak cluster point in $E$.
\end{enumerate}
In particular, a weakly closed set is weakly compact if and only if it is
weakly sequentially compact.  On a weakly closed set these conditions are
also equivalent to countable compactness for the weak topology.
Thus defining relative countable compactness by countable compactness
of the weak closure gives the same relative notion as condition 1.
\end{theorem}
\begin{proof}
If the weak closure $K$ of $A$ is compact, consider a sequence $(x_n)$ in
$A$ and put $Y=\overline{\operatorname{span}}\{x_n:n\ge1\}$.  Hahn--Banach
separation says that this norm closed linear space is weakly closed.
Consequently $K\cap Y$ is compact.  It is metrizable by the preceding
lemma, and a convergent subsequence follows.

Suppose next that every sequence in $A$ has a weakly convergent subsequence.
Then $A$ is bounded: otherwise choose $\|x_n\|\ge n$, contradicting the
boundedness of a weakly convergent subsequence.  In the bidual put
$K=\overline{J_E A}^{\,w^*}$.  This is weakly star compact by Alaoglu.
We show that $K\subset J_E E$.  Fix $\Phi\in K$ and construct $x_n\in A$
and finite sets $F_n\subset B_{E'}$ as follows.  A finite set of members
of $B_{E'}$ can norm the finite dimensional space
$\operatorname{span}\{\Phi,J_E x_1,\ldots,J_E x_n\}$ within a factor $2$:
for each point of its unit sphere choose a test functional giving absolute
value greater than $1/2$, and use a finite subcover.  Choose such a set
$F_n$.  Having chosen $F_1,\ldots,F_n$, use $\Phi\in K$ to choose
$x_{n+1}\in A$ such that
\[
 |f(x_{n+1})-\Phi(f)|<1/n
 \qquad(f\in F_1\cup\cdots\cup F_n).
\]
The initial $x_1$ is arbitrary; if $A$ is empty the assertion is trivial.
Take a subsequence $x_{n_k}\to x$ weakly.  Every $f$ in $F=\bigcup_nF_n$
then satisfies $f(x)=\Phi(f)$.  Also
$x\in\overline{\operatorname{span}}\{x_n:n\ge1\}$, since this space is
weakly closed.  The set $F$ norms, within a factor $2$, the norm closure
of $\operatorname{span}\{\Phi,J_E x_n:n\ge1\}$: the estimate first holds
on every finite initial span and passes to its norm closure.  It follows
that $\Phi-J_E x=0$.  Thus $K\subset J_E E$.  The weak star topology on
$J_E E$ is exactly the transported weak topology of $E$, so $K$ gives a
weakly compact set containing $A$.  Its weak closure is compact as well.

For clarity we also prove that the cluster-point formulation suffices.
If every sequence in $A$ has a weak cluster point, then for each $f\in E'$
the set $f(A)$ is bounded: a sequence with $|f(x_n)|\to\infty$ would have
no weak cluster point.  Uniform boundedness applied on $E'$ makes $A$
norm bounded.  For a sequence $(x_n)$ in $A$, let $Y$ be its separable
closed span and choose the countable separating family $(f_j)$ above.
Successive scalar subsequences and the diagonal procedure give a
subsequence $(y_n)$ for which all $f_j(y_n)$ converge.  It has a weak
cluster point $y$, and $y\in Y$ because $Y$ is weakly closed.  If $y_n$
does not converge weakly to $y$, some subsequence lies outside a weak
neighborhood $U$ of $y$.  It too has a weak cluster point $z$, lying in
the weakly closed set $Y\setminus U$.  But $f_j(z)=f_j(y)$ for every $j$,
so $z=y$, a contradiction.

Finally, sequential compactness implies countable compactness: a
countable open cover without a finite subcover permits a sequence whose
$n$th term avoids its first $n$ members, and no subsequence can converge.
In a countably compact $T_1$ space every sequence has a cluster point
(an infinitely repeated point suffices, and otherwise apply countable
compactness to the tails of its countable range).  The already proved
equivalence completes the last assertion.
\end{proof}

\begin{corollary}[Sequential density inside a weakly compact closure]
\label{fa2:p1:weak-compact-sequential-density}
Let $A$ be a relatively weakly compact subset of a normed space.
Every point of its weak closure is the weak limit of a sequence
of points of $A$. In particular a weakly compact set, with its
weak topology, has this closure-by-sequences property.
This does not assert it for arbitrary weakly star compact dual sets.
\end{corollary}
\begin{proof}
Fix $x\in\overline A^{\,w}$ and put $\Phi=J_Ex$.
Use exactly the finite-norming-set construction in the proof of
Eberlein--\v Smulian: choose $a_n\in A$ successively approximating
$\Phi$ on all previous finite sets, and let each new finite set
norm $\operatorname{span}\{\Phi,J_Ea_1,\ldots,J_Ea_n\}$ within
a factor two. Relative weak compactness supplies a weakly
convergent subsequence, say $a_{n_k}\rightharpoonup y$.
All the chosen tests take the same values on $J_Ey$ and $\Phi$.
Their union norms the closed span of those bidual vectors, and
$y$ is in the closed span of the $a_n$ because that subspace
is weakly closed. Thus $J_Ey=\Phi=J_Ex$, so $y=x$.
The extracted subsequence is the requested sequence from $A$.
\end{proof}

\begin{theorem}[Krein's compact convex-hull theorem]
\label{fa2:p1:weak-compact-hull}
If $K$ is weakly compact in a Banach space $E$, then
$\overline{\operatorname{conv}}K$ and
$\overline{\operatorname{aconv}}K$ are weakly compact.  Here the closures
may be taken in norm or weakly; the two closures of a convex set agree.
\end{theorem}
\begin{proof}
The empty case is immediate.  Put $C=\overline{\operatorname{conv}}K$ and
take $x_n\in C$.  Choose finite convex combinations $y_n$ of points of $K$
with $\|x_n-y_n\|<1/n$.  Let $Y$ be the norm closed linear span of the
countably many points used in these combinations.  Then $Y$ is a
separable Banach space and $L=K\cap Y$ is weakly compact and metrizable.
All the $y_n$ are barycenters of finitely supported probability measures
$\mu_n$ on $L$.

The probability measures on a compact metric space form a compact
metrizable set for convergence against continuous scalar functions.  To
see the sequential assertion needed here, choose a countable dense
subset of $C(L)$, take diagonal scalar subsequences, extend the resulting
bounded positive functional by uniform approximation, and use the
Riesz representation theorem (Theorem~\ref{fa2:p4:riesz-markov}).
Such a dense family exists: distances from a countable dense subset
of compact metric $L$ separate points, and the algebra they generate
with rational real coefficients (and rational imaginary coefficients)
is countable and dense by Stone--Weierstrass, proved below.
Thus, after passing to a subsequence, $\mu_n\to\mu$ against $C(L)$.

We spell out why this scalar limit has a vector barycenter in $Y$;
this is where completeness and separability are useful.  The identity
map from weak $L$ to norm $Y$ is Borel.  Indeed, norm closed balls are
weakly closed, open balls are countable unions of closed balls, and,
because $Y$ is separable, every norm open set is a countable union of
balls.  If $M=\sup_{z\in L}\|z\|$, cover $L$ by countably many norm balls
of radius $\eta$ centered in $L$.  Disjointify this Borel cover and keep
finitely many pieces with total measure greater than $1-\eta$.  Define
a finite simple function by their centers, using a fixed point of $L$
on the remaining part.  Its integral is a finite convex combination
of points of $L$, and its $L^1(\mu;Y)$ error from the identity is at most
$\eta+2M\eta$.  As $\eta\downarrow0$ these integrals are Cauchy.  Their
limit $b\in Y$ exists, belongs to $\overline{\operatorname{conv}}L$, and
satisfies $f(b)=\int_L f(z)\,d\mu(z)$ for every $f\in E'$.
Each restriction $f|_L$ is weakly continuous, hence
$f(y_n)\to f(b)$.  Therefore the selected $x_n$ also converge weakly
to $b\in C$.  The Eberlein--\v Smulian theorem makes $C$ weakly compact.
Apply this result to the weakly compact set
$\{\lambda x:x\in K,\ |\lambda|\le1\}$ to obtain the absolutely convex
version.  Equality of convex weak and norm closures follows directly
from Hahn--Banach separation.
\end{proof}

\begin{proposition}[Weak sequential completeness and the Schur property]
\label{fa2:p1:schur}
In $\ell^1$, weak convergence of sequences implies norm convergence.
The space $\ell^1$ is weakly sequentially complete.  The spaces $c_0$
and $C[0,1]$ are not weakly sequentially complete.
\end{proposition}

\begin{proof}
Suppose $x^n\to0$ weakly in $\ell^1$ but, along a subsequence,
$\|x^n\|_1\ge\delta>0$.  Each coordinate tends to zero.  Recursively
choose a subsequence $z^k$ and integers $0=m_0<m_1<\cdots$ so that
the mass of $z^k$ outside $(m_{k-1},m_k]$ is less than $\delta/8$.
The initial segment is made small by choosing $z^k$ sufficiently late,
and its remaining tail is made small by choosing $m_k$ sufficiently
large.  On this block choose the coordinates of $a\in\ell^\infty$ of
modulus at most one so that $a_jz^k_j=|z^k_j|$; the disjoint blocks
define one $a$.  Then $\operatorname{Re}\sum_j a_jz^k_j\ge3\delta/4$,
contradicting weak nullity.

If $(x^n)$ is weakly Cauchy in $\ell^1$, it is norm Cauchy.  Otherwise
choose $p_k,q_k\to\infty$ with $\|x^{p_k}-x^{q_k}\|_1\ge\delta$.
These differences are weakly null, contrary to the first assertion.
Norm completeness proves weak sequential completeness.  In $c_0$ the
partial sums $\sum_{j=1}^ne_j$ are weakly Cauchy, since $(c_0)'=\ell^1$,
but their only possible coordinatewise limit is the constant sequence
$1$, which is not in $c_0$.  In $C[0,1]$, the functions $t^n$ are weakly
Cauchy by dominated convergence for finite regular measures; their
pointwise limit is $0$ on $[0,1)$ and $1$ at $1$, so they have no weak
limit in $C[0,1]$.
\end{proof}

\begin{corollary}[Compact dual sets over a separable TVS]
\label{fa2:p1:separable-tvs-dual-compact}
If $E$ is a separable TVS, every weakly star compact subset $K$ of
its continuous dual is metrizable.
\end{corollary}
\begin{proof}
Evaluations on a countable dense subset separate continuous linear
functionals. They define a continuous injection of compact $K$
into a countable product of scalar spaces. It is a homeomorphism
onto its Hausdorff image, which is metrizable.
\end{proof}

\begin{proposition}[Metrizability of bounded weak topologies]
\label{fa2:p1:primal-weak-metrizable}
For a normed space $E$, the weak topology on $B_E$ is metrizable if and
only if $E'$ is norm separable.  The weak star topology on $B_{E'}$ is
metrizable if and only if $E$ is norm separable.
\end{proposition}
\begin{proof}
For sufficiency, a norm dense sequence of functionals in $B_{E'}$
generates the weak topology on a norm bounded set: uniform control of
the operator norm error makes each omitted test a uniform limit of
the chosen tests.  The displayed countable-product metric above applies.
The same argument on the dual ball uses a dense sequence in $E$.

Conversely, a countable neighborhood base at zero in $B_E$ can be
refined by weak neighborhoods involving finitely many members of $E'$.
Let $F$ be the norm closed span of all these countably many tests.
For $g\in E'$ and $\eps>0$, continuity of $g$ supplies finitely many
of the chosen tests $f_1,\ldots,f_m$ and $\delta>0$ such that
$x\in B_E$, $\max|f_j(x)|<\delta$ imply $|g(x)|<\eps$.
Then $\|g|_{\cap_j\ker f_j}\|\le\eps$.  Hahn--Banach extension from
this common kernel gives
$\operatorname{dist}(g,\operatorname{span}\{f_j\})\le\eps$:
extend $g$ restricted to the kernel with the same norm and subtract;
the difference factors through the finite dimensional test map.
Thus $F=E'$, which is separable.
For the dual ball, take the countably many points $x_j\in E$ occurring
in such a base.  Their closed span must be $E$, for otherwise
Hahn--Banach supplies a nonzero $f\in B_{E'}$ vanishing on all $x_j$;
it could not be distinguished from zero by the purported base.
\end{proof}

\begin{example}[An incomplete compact-hull obstruction]
\label{fa2:p1:incomplete-compact-hull}
In $c_{00}$ with the $\ell^1$ norm the set
$K=\{0\}\cup\{e_n/n:n\ge1\}$ is norm compact.  Its closed convex hull
in $c_{00}$ is not weakly compact.  Indeed the convex combinations
$y_m=\sum_{n=1}^m2^{-n}e_n/n$ belong to it (use the remaining weight
at zero).  Their coordinates converge to $(2^{-n}/n)_n$, which has
infinite support.  No subsequence can converge weakly in $c_{00}$,
since coordinate evaluations are continuous.  Eberlein--\v Smulian
therefore rules out weak compactness.  They are also a norm Cauchy
sequence without a limit in $c_{00}$, so the hull is not norm compact.
\end{example}

\section{The bounded weak star topology and Krein--\v Smulian}

\begin{theorem}[The bounded weak star topology]
\label{fa2:p1:bounded-weak-star}
Let $E$ be a normed space.  The seminorms
$p_{(x_n)}(f)=\sup_n|f(x_n)|$, where $x_n\to0$ in norm, define a
Hausdorff locally convex topology $\beta$ on $E'$.  It is the finest
topology which induces the weak star topology on every norm bounded
subset of $E'$.  A set $C\subset E'$ is $\beta$-closed if and only if
$C\cap rB_{E'}$ is weakly star closed for every $r>0$.
If $E$ is Banach, the continuous linear functionals for $\beta$ are
exactly the evaluations $f\mapsto f(x)$, $x\in E$.
\end{theorem}
\begin{proof}
The seminorms are finite, dominate each single evaluation by using the
sequence $(x,0,0,\ldots)$, and are norm continuous.  Their open unit
balls are exactly the sets $|f(x_n)|<1$ for all $n$: since $f(x_n)\to0$,
the supremum is a maximum unless it is zero.  Finite intersections of
such balls contain another of the same form by interlacing scaled
null sequences.  On a bounded subset the tail tests are uniformly
small, so these balls are relatively weakly star open.

Let $U$ be any set whose intersection with every bounded set is relatively
weakly star open and let $f_0\in U$.  We construct finite sets $F_n\subset E$
such that, for $n\ge0$,
\[
 \begin{aligned}
 &\{f:\|f-f_0\|\le n+1,\ |(f-f_0)(x)|\le1
                    \ (x\in F_0\cup\cdots\cup F_n)\}\subset U,
 \\
 &\|x\|\le1/n\quad(x\in F_n,n\ge1).
 \end{aligned}
\]
For $n=0$ use relative openness in $f_0+B_{E'}$ and double the finitely
many tests to replace strict inequalities by non-strict ones.  Suppose
$F_0,\ldots,F_{m-1}$ have been chosen.  If no suitable finite
$F_m\subset m^{-1}B_E$ exists, the sets
\[
 \{f:\|f-f_0\|\le m+1,\ |(f-f_0)(x)|\le1
       \ (x\in F_0\cup\cdots\cup F_{m-1}\cup F)\}\setminus U,
 \quad F\subset m^{-1}B_E\text{ finite},
\]
are nonempty weakly star compact sets with the finite intersection
property.  Their common point $h$ satisfies $\|h-f_0\|\le m$, since
all the tests in $m^{-1}B_E$ hold, and satisfies all the old tests.
The preceding induction step puts $h$ in $U$, a contradiction.
Thus the construction succeeds.  List the union of the finite $F_n$
(inserting zero if necessary) as a null sequence.  Its open polar
translated by $f_0$ lies in $U$: every $f$ has finite norm distance
from $f_0$, so one of the displayed bounds applies.  Hence $U$ is
$\beta$-open.  This proves the finest-topology assertion and its
closed-set version.

Suppose now that $E$ is Banach and $L$ is a $\beta$-continuous linear
functional on $E'$.  Interlacing and scaling gives a null sequence
$(x_n)$ with $|L(f)|\le\sup_n|f(x_n)|$.  The bounded linear map
$T:E'\to c_0$, $Tf=(f(x_n))$, therefore permits a well-defined bounded
functional $v(Tf)=L(f)$ on its range.  Extend $v$ to $c_0$ by
Hahn--Banach and represent the extension by $(a_n)\in\ell^1$.
Since $E$ is complete, $x=\sum_na_nx_n$ exists in $E$, and
$L(f)=\sum_na_nf(x_n)=f(x)$.  Conversely, each evaluation is manifestly
$\beta$-continuous.
\end{proof}

\begin{theorem}[Krein--\v Smulian]
\label{fa2:p1:krein-smulian}
Let $E$ be Banach and $C\subset E'$ convex.  Then $C$ is weakly star
closed if and only if $C\cap rB_{E'}$ is weakly star closed for every
$r>0$.  A linear subspace $M\subset E'$ is weakly star closed if and
only if $M\cap B_{E'}$ is weakly star closed.  If $E$ is separable, a
convex subset of $E'$ is weakly star closed if and only if it is
weakly star sequentially closed.
\end{theorem}
\begin{proof}
The forward assertion is immediate.  For the converse, the bounded
intersection condition makes $C$ closed for the locally convex topology
$\beta$.  A point outside $C$ is separated from it by a $\beta$-continuous
linear functional.  The preceding theorem makes this functional an
evaluation at a point of $E$, hence weakly star continuous.  Thus $C$
is also weakly star closed.  For a subspace use
$M\cap rB_{E'}=r(M\cap B_{E'})$.  For separable $E$, each dual ball is
weakly star metrizable.  Sequential closedness of $C$ therefore makes
each $C\cap rB_{E'}$ closed in that ball, and the first assertion applies.
\end{proof}

\begin{example}[Why completeness is indispensable]
\label{fa2:p1:krein-smulian-incomplete}
Let $E$ be an incomplete normed space, regard $E$ as a dense proper
subspace of its completion $\widehat E$, and choose $z\in\widehat E\setminus E$.
The hyperplane $M=\{f\in E':\widehat f(z)=0\}$ has weakly star closed
intersection with every dual ball but is not weakly star closed in $E'$.
Indeed, on each bounded ball evaluation at $z$ is the uniform limit of
evaluations at points of $E$.  If $M$ were weakly star closed, its
nonzero defining linear functional would be a scalar multiple of an
evaluation at some $x\in E$ (factor through a finite family of tests,
or apply separation to the hyperplane).  This would give $z\in E$
because $E'$ separates points of $\widehat E$, a contradiction.
\end{example}

\section{Variational support and norm-attaining functionals}

\begin{lemma}[Ekeland's variational principle]
\label{fa2:p1:ekeland}
Let $(X,d)$ be a complete metric space and $F:X\to(-\infty,+\infty]$
lower semicontinuous, bounded below, and not identically infinite.
If $F(x_0)\le\inf_XF+\alpha r$ with $\alpha,r>0$, there exists $\bar x$
such that $d(\bar x,x_0)\le r$, $F(\bar x)\le F(x_0)$, and
\[
 F(y)>F(\bar x)-\alpha d(y,\bar x)\qquad(y\ne\bar x).
\]
\end{lemma}
\begin{proof}
Put $S(x)=\{y:F(y)+\alpha d(y,x)\le F(x)\}$.  Starting at $x_0$,
choose $x_{n+1}\in S(x_n)$ whose $F$-value is within $2^{-n}$ of
$\inf_{S(x_n)}F$.  Triangle inequality gives
$S(x_{n+1})\subset S(x_n)$ and
$\alpha\sum_nd(x_{n+1},x_n)\le F(x_0)-\inf F\le\alpha r$.
Thus $x_n\to\bar x$, with the asserted distance bound.  Lower
semicontinuity gives $\bar x\in S(x_n)$ for every $n$.  Writing
$l=\lim F(x_n)$, the near-infimum choices imply
$\inf_{S(x_n)}F\to l$; since $\bar x\in S(x_n)$, we have $F(\bar x)\ge l$,
and lower semicontinuity gives the reverse inequality.  If
$y\in S(\bar x)$, then $y\in S(x_n)$ for every $n$, whence $F(y)\ge l$.
The defining inequality for $S(\bar x)$ forces $d(y,\bar x)=0$.
This proves the strict variational inequality.
\end{proof}

\begin{theorem}[Bishop--Phelps]
\label{fa2:p1:bishop-phelps}
Let $E$ be a real or complex Banach space and $C\subset E$ a nonempty
closed convex set.  If $f\in E'$ is bounded above in real part on $C$,
then for every $\alpha>0$ there are $x\in C$ and $g\in E'$ with
$\|g-f\|\le\alpha$ and
$\operatorname{Re}g(y)\le\operatorname{Re}g(x)$ for every $y\in C$.
The support points of $C$ are dense in its boundary.  In particular,
norm-attaining functionals are norm dense in $E'$.
\end{theorem}
\begin{proof}
First work over the real field.  Apply Ekeland to $F=-f$ on $C$.
It gives $x\in C$ with $f(y-x)\le\alpha\|y-x\|$ for all $y\in C$.
The convex set
$A=\{(y-x,f(y-x)):y\in C\}$ is disjoint from the open convex cone
$D=\{(z,t):t>\alpha\|z\|\}$.  Separate them by a nonzero continuous
linear functional $(h,s)$.  Since $0\in A$ and $D$ is a cone, orient
the separation so that $h(z)+st\ge0$ on $D$ and
$h(y-x)+sf(y-x)\le0$ on $A$.  Testing large positive $t$ gives $s\ge0$;
letting $t\downarrow\alpha\|z\|$ gives $\|h\|\le s\alpha$.
If $s=0$, this forces $h=0$, impossible.  Thus
$g=f+h/s$ supports $C$ at $x$ and $\|g-f\|\le\alpha$.
For a complex space apply the real argument to $\operatorname{Re}f$
and recover the complex functional from a real functional $h$ by
$h(z)-i h(iz)$; this correspondence preserves the norm.

We justify the boundary assertion separately, because an arbitrary
boundary point need not initially maximize any functional.  Let
$b\in\partial C$, take $z\notin C$ close to $b$, and apply Ekeland
on $C$ to $F(y)=\|y-z\|$, with $\alpha=1/2$ and
$r=2\|b-z\|$.  The resulting $x$ is within $2\|b-z\|$ of $b$ and
minimizes on $C$ the continuous convex function
$\|y-z\|+\tfrac12\|y-x\|$.
The convex subgradient sum rule yields functionals $u,v$ with
\[
 \|u\|=1,\quad u(x-z)=\|x-z\|,\quad \|v\|\le\tfrac12,
 \qquad -(u+v)\in N_C(x),
\]
where $N_C(x)=\{h:h(y-x)\le0\ (y\in C)\}$.  Here $x\ne z$ because
$z\notin C$.  To make this use of the sum rule explicit, its required
form follows from Fenchel--Rockafellar duality
(Theorem~\ref{fa2:p1:fenchel-rockafellar}): if a sum $\phi+\psi$
attains its minimum at $x$ and $\phi$ is continuous there, dual
attainment and the two Fenchel inequalities force equality in both,
so a functional $a$ satisfies $a\in\partial\phi(x)$ and
$-a\in\partial\psi(x)$.  Apply this first to the continuous sum of
the two norms and the indicator of $C$, then to the two continuous
norms.  The subgradients of a norm at $w\ne0$ are exactly its norm-one
norming functionals, by Hahn--Banach and the defining subgradient
inequality; at zero they are the dual unit ball.  This proves the
display.  Since $\|u+v\|\ge1/2$, the normal $-(u+v)$ is nonzero,
so $x$ is a support point.  Letting $z\to b$ proves density.

Finally take $C=B_E$.  For $f\ne0$ choose $0<\alpha<\|f\|$ in the
first assertion.  Then $g\ne0$ and its supporting point realizes
$\operatorname{Re}g(x)=\|g\|$, hence $|g(x)|=\|g\|$.
The zero functional is itself norm attaining.
\end{proof}

\begin{corollary}[Preserving a bounded separation]
\label{fa2:p1:bishop-phelps-separation}
Suppose $C$ is nonempty, closed and convex, $A$ is bounded, $\|f\|=1$, and
$\sup_C\operatorname{Re}f<\inf_A\operatorname{Re}f$.  For every
$\eps>0$ there are a support point $x\in C$ and $g\in E'$ such that
$\|g\|=1$, $\|g-f\|<\eps$, and
$\operatorname{Re}g(x)=\sup_C\operatorname{Re}g<\inf_A\operatorname{Re}g$.
\end{corollary}
\begin{proof}
Write $M=\sup_C\operatorname{Re}f$. Choose finite $d>0$ and $R\ge0$
such that $\operatorname{Re}f(a)\ge M+d$ and $\|a\|\le R$ for all
$a\in A$, taking $d=1$ and $R=0$ if $A=\varnothing$.
Fix $c\in C$ with defect $M-\operatorname{Re}f(c)<d/4$ and choose
$\alpha>0$ with $2\alpha<\eps$, $\alpha<1/2$, and
$\alpha(R+\|c\|+1)<d/4$. Apply Ekeland with starting point $c$,
parameter $\alpha$, and distance bound
$r=(M-\operatorname{Re}f(c)+\alpha)/\alpha$.
The support construction gives $h$ with $\|h-f\|\le\alpha$ and
$\|x-c\|\le r$. For every $a\in A$,
\[
 \operatorname{Re}h(a)-\operatorname{Re}h(x)
 \ge d-\alpha(R+\|c\|)-\alpha r>d/2.
\]
The gap is uniform on $A$, with $\inf\varnothing=+\infty$.
Normalize $h$ to $g=h/\|h\|$; then
$\|g-f\|\le2\alpha<\eps$, and the strict separation is preserved.
\end{proof}

% END INPUT


% BEGIN parts/part1-james.tex
\chapter{James's attainment criterion}
\label{fa2:p1:james-chapter}

Compactness implies attainment of continuous scalar extrema.  The converse
for weakly closed sets in a Banach space is a deep theorem, not a consequence
of Alaoglu alone: a weak star limit of points of $E$ may lie outside $E$.
We give the convex-block argument which detects precisely this obstruction.
The technical construction follows James's argument as organized in
\cite[Lemma 1.13.13 and Section 2.9]{B05}; the tail estimate is written out
below to make the strict inequality used at the end explicit.

\section{The convex-block lemma}

For a bounded sequence $(f_n)$ in the dual of a real normed space set
\[
 \begin{aligned}
 L(f_n)&=\{w\in E': w(x)\le\limsup_n f_n(x)\quad(x\in E)\},\\
 V(f_n)&=\{(g_n):g_n\in\operatorname{conv}\{f_j:j\ge n\}\}.
 \end{aligned}
\]
The set $L(f_n)$ is nonempty.  Indeed, $p(x)=\limsup_n f_n(x)$ is a
finite continuous sublinear function, so Hahn--Banach extends the zero
functional on $\{0\}$ to a linear $w\le p$; it is bounded because
$|w(x)|\le\sup_n\|f_n\|\,\|x\|$.  Applying the upper bound to $-x$
also gives
\[
 \liminf_n f_n(x)\le w(x)\le\limsup_n f_n(x).
\]
If $(g_n)\in V(f_n)$, then $V(g_n)\subset V(f_n)$ and
$L(g_n)\subset L(f_n)$.  A subsequence, indexed in increasing order,
also belongs to $V(f_n)$.

\begin{lemma}[James's convex-block lemma]
\label{fa2:p1:james-block-lemma}
Let $A$ be a nonempty balanced subset of $B_E$ in a real normed space.
Write $|h|_A=\sup_{x\in A}|h(x)|$.  Suppose $0<\theta<1$, $f_n\in B_{E'}$,
and
\[
 |f-w|_A\ge\theta
 \quad\bigl(f\in\operatorname{conv}\{f_n:n\ge1\},\ w\in L(f_n)\bigr).
\]
For any positive numbers $\beta_n$ with $\sum_n\beta_n=1$, there exist
$g_n\in B_{E'}$ and $\theta\le a\le2$ such that, for every $w\in L(g_n)$,
\[
 \left|\sum_{j=1}^\infty\beta_j(g_j-w)\right|_A=a,
 \qquad
 \left|\sum_{j=1}^n\beta_j(g_j-w)\right|_A
       <a\left(1-\theta\sum_{j>n}\beta_j\right)\quad(n\ge1).
\]
All infinite sums here converge in the norm of $E'$.
\end{lemma}
\begin{proof}
Put $t_n=\sum_{j\ge n}\beta_j$.  Choose positive $\eps_n\to0$ so small that
\[
 \sum_{n=1}^\infty \frac{\beta_n\eps_n}{t_nt_{n+1}}<1-\theta.
\]
This is possible, for example by making the $n$th summand less than
$(1-\theta)2^{-n-1}$ and also requiring $\eps_n<1/n$.
We recursively construct sequences $f^0=(f_j)$, $z^n$, $f^n$, points
$u_n\in A$, functionals $g_n$, and numbers $a_n$.  Given $f^{n-1}$ and
$g_1,\ldots,g_{n-1}$, let
\[
 a_n=\inf_{\substack{g\in\operatorname{conv}\{f_j^{n-1}:j\ge n\}\\
                         v\in V(f^{n-1})}}
       \ \sup_{w\in L(v)}
       \left|\sum_{j<n}\beta_jg_j+t_ng-w\right|_A.
\]
Every $L(v)$ is nonempty.  The expressions are at most $2$ and at
least $\theta$: the positive sum before $-w$ is a convex combination
of the original $f_j$, while $L(v)\subset L(f^0)$.  Thus
$\theta\le a_n\le2$.  Choose $g_n$ in the indicated tail convex hull
and $z^n\in V(f^{n-1})$ so that
\[
 a_n\le\sup_{w\in L(z^n)}
       \left|\sum_{j<n}\beta_jg_j+t_ng_n-w\right|_A
       <a_n(1+\eps_n).
\]
Choose $w_n\in L(z^n)$ and, using balance to choose the sign of a
nearly norming point, $u_n\in A$ such that
\[
 \left(\sum_{j<n}\beta_jg_j+t_ng_n-w_n\right)(u_n)
       >a_n(1-\eps_n).
\]
Finally choose a subsequence $f^n$ of $z^n$ for which
$f_j^n(u_n)\to\liminf_j z_j^n(u_n)$.  All its members have norm at
most one, and $f^n\in V(f^{n-1})$.

We verify the two nesting properties needed for this recursion.
First, $a_{n+1}\ge a_n$.  Indeed $V(f^n)\subset V(f^{n-1})$, and for
every admissible next choice $g$ the functional
\[
 (\beta_ng_n+t_{n+1}g)/t_n
\]
belongs to $\operatorname{conv}\{f_j^{n-1}:j\ge n\}$.
Substituting it in the infimum for $a_n$ gives exactly the expression
for the candidate defining $a_{n+1}$, with a possibly smaller class
of choices.  Second,
\[
 L(g_j)\subset\bigcap_{n\ge0}L(f^n)
             \subset\bigcap_{n\ge1}L(z^n).
\]
For the first inclusion, if $j>n$, repeatedly unravel the tail
convex combinations to obtain
$g_j\in\operatorname{conv}\{f_k^n:k\ge j\}$.  Consequently
$\limsup_jg_j(x)\le\limsup_jf_j^n(x)$ for each $x$.
The second inclusion uses that $f^n$ is a subsequence of $z^n$.

Fix $w\in L(g_j)$.  The chosen scalar subsequence gives
$w(u_n)=\lim_j f_j^n(u_n)=\liminf_jz_j^n(u_n)\le w_n(u_n)$.
Thus the lower estimate at $u_n$ remains true when $w_n$ is replaced
by $w$, while the preceding supremum gives the upper estimate:
\[
 a_n(1-\eps_n)
 <\left|\sum_{j<n}\beta_jg_j+t_ng_n-w\right|_A
 <a_n(1+\eps_n).
\]
The increasing bounded sequence $a_n$ has a limit $a\in[\theta,2]$.
Because $t_n\to0$, the middle functional converges in norm to
$\sum_j\beta_jg_j-w$.  This proves its $A$-seminorm is $a$.

It remains to prove the strict prefix bound; norm convergence alone
would not provide it.  Set
$G_n=\sum_{j\le n}\beta_j(g_j-w)$, $G_0=0$, and
$H_n=\sum_{j<n}\beta_jg_j+t_ng_n-w$.
The algebraic identity
\[
 G_n=\frac{\beta_n}{t_n}H_n+\frac{t_{n+1}}{t_n}G_{n-1}
\]
gives, on division by $t_{n+1}$ and iteration,
\[
 \frac{|G_n|_A}{t_{n+1}}
 \le\sum_{k\le n}\frac{\beta_k|H_k|_A}{t_kt_{k+1}}
 <a\sum_{k\le n}\frac{\beta_k(1+\eps_k)}{t_kt_{k+1}}.
\]
Since $\beta_k/(t_kt_{k+1})=1/t_{k+1}-1/t_k$ and $t_1=1$,
the right side is strictly smaller than
$a(1/t_{n+1}-1+1-\theta)=a(1/t_{n+1}-\theta)$.
Multiplying by $t_{n+1}$ proves the assertion.
\end{proof}

\section{The attainment theorem}

\begin{lemma}[A separable witness to failure of compactness]
\label{fa2:p1:james-witness}
Let $A\subset B_E$ be nonempty, norm separable and weakly closed in a
Banach space.  If $A$ is not weakly compact, there exist $0<\theta<1$
and $f_n\in B_{E'}$ such that $f_n(x)\to0$ for every $x\in A$ and
$|f|_A\ge\theta$ for every $f\in\operatorname{conv}\{f_n\}$.
This assertion holds over either scalar field.
\end{lemma}
\begin{proof}
Let $Y=\overline{\operatorname{span}}A$.  It is a separable Banach
space.  The weak star closure of $J_YA$ in $B_{Y''}$ contains a point
$\Phi\notin J_YY$: otherwise it would give a weakly compact set in
$Y$ whose closed subset $A$ is compact.  A point of this closure
which belongs to $J_YY$ must belong to $J_YA$, because $A$ is weakly
closed.  Since $J_YY$ is norm closed, $d=\operatorname{dist}(\Phi,J_YY)>0$.
Choose $0<\theta<\min(d,1)$ and a dense sequence $(a_j)$ in $A$.

For $M_n=\operatorname{span}\{a_1,\ldots,a_n\}$, the restriction of
$\Phi$ to $M_n^\perp\subset Y'$ has norm
$\operatorname{dist}(\Phi,J_YM_n)\ge d$.  Here is the distance formula:
extend this restriction to all of $Y'$ with its norm by Hahn--Banach;
the difference from $\Phi$ annihilates $M_n^\perp$, hence is a member
of $J_YM_n$.  The reverse inequality follows by restriction.
Choose $h_n\in B_{Y'}$ vanishing on $M_n$ with $\Phi(h_n)=\theta$,
first choosing a value larger than $\theta$ in absolute value and
then adjusting phase and scaling.  Extend $h_n$ to $f_n\in B_{E'}$.
Density of $(a_j)$ and the uniform norm bound give $f_n(x)\to0$ on $A$.
For any convex combination $f$ of the $f_n$, its restriction $h$ to
$Y$ satisfies $\Phi(h)=\theta$.  Since $\Phi\in\overline{J_YA}^{w^*}$,
$\theta\le\sup_{x\in A}|h(x)|=|f|_A$, as required.
\end{proof}

\begin{theorem}[James's weak compactness theorem]
\label{fa2:p1:james}
Let $A$ be a nonempty weakly closed subset of a real or complex Banach
space $E$.  The following are equivalent:
\begin{enumerate}
\item $A$ is weakly compact;
\item for every $f\in E'$, the finite supremum
      $\sup_{x\in A}|f(x)|$ is attained at a point of $A$;
\item for every continuous real linear functional $u$ on the underlying
      real space, $\sup_A|u|$ is finite and attained;
\item for every such $u$, $\sup_Au$ is finite and attained.
\end{enumerate}
Consequently $E$ is reflexive if and only if every member of $E'$
attains its norm on $B_E$.
\end{theorem}
\begin{proof}
Compactness gives all three attainment assertions immediately.
We first prove the converse for a real Banach space and a balanced
weakly closed $A\subset B_E$.  Suppose $A$ is not weakly compact.
By Eberlein--\v Smulian there is a sequence in $A$ with no weakly
convergent subsequence.  Its separable closed span $Y$ makes
$A_0=A\cap Y$ separable, weakly closed, and not weakly compact.
Apply the witness lemma to $A_0$.  It gives $f_n\in B_{E'}$ and
$0<\theta<1$ with $f_n\to0$ pointwise on $A_0$ and
$|f|_{A_0}\ge\theta$ for every convex combination $f$ of the $f_n$.
Every $w\in L(f_n)$ vanishes on $A_0$, so $|f-w|_A\ge\theta$.

Choose $0<q<\theta^2/2$ and use the block lemma with
$\beta_n=(1-q)q^{n-1}$.  Let $(g_n)$ and $a\ge\theta$ be its
conclusion, fix $w\in L(g_n)$, and put
$z=\sum_n\beta_n(g_n-w)$.  Then $|z|_A=a$.
For each $x\in A$, the inequality
$\liminf_j(g_j-w)(x)\le0$ supplies arbitrarily large $n$ with
\[
 (g_{n+1}-w)(x)<a\theta-2q,
\]
whose right side is positive because $a\theta\ge\theta^2>2q$.
Writing $t=\sum_{j>n}\beta_j=q^n$, so that
$\beta_{n+1}=(1-q)t$ and $\sum_{j>n+1}\beta_j=qt$, gives
\begin{align*}
 z(x)
 &\le\left|\sum_{j\le n}\beta_j(g_j-w)\right|_A
       +\beta_{n+1}(g_{n+1}-w)(x)+2\sum_{j>n+1}\beta_j\\
 &<a(1-\theta t)+(a\theta-2q)(1-q)t+2qt\\
 &=a-q(a\theta-2q)t<a.
\end{align*}
The factor $2$ bounds $\|g_j-w\|$, since $\|w\|\le1$.
Apply the same argument to $-x\in A$ to get $-z(x)<a$.
Thus $|z(x)|<a$ for every $x\in A$: the supremum is not attained.
This proves the contrapositive in the balanced real case.

Now assume assertion 2 for an arbitrary $A$.  Its scalar images are
bounded; uniform boundedness applied to $\{J_Ex:x\in A\}$ on $E'$
shows that $A$ is norm bounded.  Rescale to $A\subset B_E$ and set
\[
 B=\bigcap_{f\in E'}\{x\in E:|f(x)|\le\sup_{a\in A}|f(a)|\}.
\]
This is a balanced weakly closed subset of $B_E$ containing $A$.
For each $f$, the supremum of $|f|$ on $B$ equals that on $A$ and
is attained there.  In the real case the result just proved makes
$B$ weakly compact, hence its weakly closed subset $A$ compact.
In the complex case regard $E$ as real.  Every continuous real
linear functional has the form $\operatorname{Re}f$ for a unique
complex linear $f$, by $f(x)=u(x)-iu(ix)$.  Balance of $B$ gives
$\sup_B|u|=\sup_B|f|$, with attainment: rotate a maximizing point
of $|f|$ to make its value real.  The real balanced result again
makes $B$, and then $A$, compact.  The weak topologies of a complex
space and its realification agree, since $f=\operatorname{Re}f+
i\operatorname{Re}(-if)$.

Assertion 3 implies compactness by the real case applied directly
to $A$.  Assertion 4 implies assertion 3 because
$\sup_A|u|=\max(\sup_Au,\sup_A(-u))$, and both maxima are attained.
Finally the closed unit ball is weakly compact exactly when $E$
is reflexive, as proved in the reflexivity chapter.  Applying the
equivalence to this ball proves the last assertion.
\end{proof}

\begin{remark}[What the closedness hypothesis does]
\label{fa2:p1:james-closedness}
In a nonzero finite-dimensional normed space, every scalar linear
functional attains its absolute supremum on the unit sphere, but this does not turn an
arbitrary norm closed bounded nonconvex subset of an infinite
dimensional space into a weakly closed set.  The theorem explicitly
requires weak closedness; for convex norm closed sets Hahn--Banach
supplies it.  In particular its commonly used convex form says that
a nonempty norm closed bounded convex set is weakly compact if
every continuous real linear functional attains its supremum there.
The completeness of $E$ is used when $J_YY$ is declared norm closed
in the witness lemma.  Removing it destroys that step, and the
attainment criterion is not being asserted for incomplete spaces.
\end{remark}

% END INPUT


% BEGIN parts/part1-foundation-15.tex
\chapter{Schauder Bases and Complemented Subspaces}\label{fa2:p1:c15}

A basis can encode a vector by a finite algebraic expression, an
ordered convergent series, or an orthogonal expansion. These require
different hypotheses. In this chapter $E$ is a real or complex Banach
space unless explicitly stated otherwise.
\source{B05 Secs. 3.2, 4.1--4.2}
\cite{B05} 


% FAMILY D16.C15.01
\section{Coordinates whose continuity must be proved}
\begin{definition}[Schauder basis and its algebraic coordinate maps]\label{fa2:p1:c15:basis}
A sequence of nonzero vectors $(e_n)_{n\geq1}$ is a Schauder basis
of $E$ if for each $x\in E$ there is a unique scalar sequence
$(a_n(x))$ with
$x=\lim_{N\to\infty}\sum_{n=1}^Na_n(x)e_n$ in norm.
Write $e_n'(x)=a_n(x)$ and
$P_Nx=\sum_{n=1}^Ne_n'(x)e_n$, with $P_0=0$.
At this stage these maps are algebraic linear maps: uniqueness and
addition of convergent series prove linearity, without any continuity
assumption. A finite basis has the analogous finite expansion;
for a list of length $d$ put $P_N=P_d=I$ when $N\geq d$.
The empty basis is allowed for the zero space, with every $P_N=0$.
\end{definition}

\begin{theorem}[Continuity of Schauder coordinates]\label{fa2:p1:c15:coordinate-continuity}
If $(e_n)$ is a Schauder basis of a Banach space, all coordinate
functionals $e_n'$ and partial-sum maps $P_N$ are bounded.
The basis norm
\[
 \norm x_b=\sup_{N\geq0}\norm{P_Nx}
\]
is a complete norm equivalent to the original norm.
\end{theorem}
\begin{proof}
We give the completeness argument before applying any boundedness
theorem to $P_N$. Each supremum is finite since $P_Nx\to x$.
Linearity and uniqueness give the norm axioms, and taking the
limit of the partial sums gives $\norm x\leq\norm x_b$.
Let $(x^{(j)})$ be Cauchy in $\norm\cdot_b$, with coordinates
$a_n^{(j)}$. For every $n$,
\[
 |a_n^{(j)}-a_n^{(k)}|\norm{e_n}
 =\norm{(P_n-P_{n-1})(x^{(j)}-x^{(k)})}
 \leq2\norm{x^{(j)}-x^{(k)}}_b.
\]
Thus each coordinate converges to a scalar $a_n$. Given $\eps>0$,
choose $J$ with $\norm{x^{(j)}-x^{(k)}}_b<\eps$ for $j,k\geq J$.
For any fixed finite $m$, letting $k\to\infty$ in its partial-sum
inequality gives
\[
 \left\|\sum_{n=1}^m(a_n-a_n^{(j)})e_n\right\|\leq\eps
 \quad(j\geq J),
\]
uniformly in $m$. Fix such a $j$. For $m>l$ this yields
\[
 \left\|\sum_{n=l+1}^m a_ne_n\right\|
 \leq2\eps+\left\|\sum_{n=l+1}^m a_n^{(j)}e_n\right\|.
\]
The last term tends to zero as $l,m\to\infty$ because the series
for $x^{(j)}$ converges. Since $\eps$ is arbitrary, the partial
sums of $\sum a_ne_n$ are Cauchy in the original Banach space.
Let their limit be $x$. By uniqueness these $a_n$ are its basis
coordinates. The preceding uniform inequality now says
$\norm{x-x^{(j)}}_b\leq\eps$ for $j\geq J$. This proves
completeness of the basis norm.

The identity $(E,\norm\cdot_b)\to(E,\norm\cdot)$ is a bounded
bijection between Banach spaces. The bounded inverse theorem of
Chapter~\ref{fa2:p1:c08} gives $\norm x_b\leq C\norm x$ for some
$C<\infty$. Consequently $\norm{P_Nx}\leq C\norm x$ for
every $N$. Also
$|e_n'(x)|\norm{e_n}\leq2C\norm x$, proving continuity of each
coordinate. In finite dimension continuity of coordinates is the
finite-dimensional norm theorem, and the finite supremum norm is
equivalent and complete. The zero-space case is immediate.
\end{proof}

% FAMILY D16.C15.02
\section{The basis constant and a converse criterion}
\begin{proposition}[Partial-sum projections and basis bounds]\label{fa2:p1:c15:partial-sums}
For a Schauder basis,
$P_NP_M=P_{\min(N,M)}$, $P_Nx\to x$, and
$K=\sup_N\norm{P_N}<\infty$. If $E\ne0$, $K\geq1$.
For finitely supported scalars and $m\leq N$,
\[
 \left\|\sum_{n=1}^m a_ne_n\right\|
 \leq K\left\|\sum_{n=1}^N a_ne_n\right\|,
 \quad
 |e_n'(x)|\norm{e_n}\leq2K\norm x.
\]
The equivalent norm $\norm\cdot_b$ makes every $P_N$ contractive.
\end{proposition}
\begin{proof}
Uniqueness of coordinates proves the projection identities and
the displayed finite-sum bound. Uniform boundedness of the
operator norms was already obtained without circularity in the
coordinate theorem; it also follows now from uniform boundedness
since $P_Nx\to x$. Every nonzero initial span contains a fixed
vector of its projection, so its projection norm is at least $1$.
The difference $P_n-P_{n-1}$ gives the coordinate estimate, and
\[
 \norm{P_Nx}_b=\sup_M\norm{P_{\min(M,N)}x}\leq\norm x_b
\]
gives the final assertion.
\end{proof}

\begin{theorem}[A useful basis criterion]\label{fa2:p1:c15:basis-criterion}
Let $(e_n)$ be a sequence of nonzero vectors in a Banach space
whose linear span is dense. It is a Schauder basis if and only
if some $C<\infty$ satisfies
\[
 \left\|\sum_{n=1}^m a_ne_n\right\|
 \leq C\left\|\sum_{n=1}^N a_ne_n\right\|
 \quad(m\leq N)
\]
for every finite scalar family.
\end{theorem}
\begin{proof}
Necessity was just proved. For sufficiency, the inequality first
implies linear independence: in a finite zero combination its
first partial sum is zero, forcing its first coefficient zero;
continue inductively. Thus truncation defines $P_m$ on the
algebraic span. The estimate makes each $P_m$ bounded there
with norm at most $C$, so completeness of its finite-dimensional
range extends it to all of $E$. For $x\in E$ and finitely
supported $z$, the estimate
\[
 \norm{P_mx-x}\leq(C+1)\norm{x-z}+\norm{P_mz-z}
\]
shows $P_mx\to x$ by first approximating $x$ by $z$, then taking
$m$ beyond its support. The identities $P_mP_N=P_{\min(m,N)}$
extend from the dense span by continuity. Thus
$(P_n-P_{n-1})x=a_n e_n$ and
$P_mx=\sum_{n=1}^ma_n e_n$, giving an expansion. If another
series has the same limit, continuity of $P_m$ shows its first
$m$ terms equal $P_mx$; subtract consecutive terms to get the
same coefficients. This proves uniqueness.
\end{proof}

% FAMILY D16.C15.03
\section{Canonical examples and the separability boundary}
\begin{proposition}[The standard bases]\label{fa2:p1:c15:canonical}
The unit vectors form a Schauder basis of $c_0$ and of $\ell^p$
for $1\leq p<\infty$, with basis constant $1$. Their closed span
in $\ell^\infty$ is only $c_0$, and $\ell^\infty$ has no
countable Schauder basis at all.
\end{proposition}
\begin{proof}
In $\ell^p$ the norm of a truncation error is
$(\sum_{n>N}|x_n|^p)^{1/p}\to0$; in $c_0$ it is
$\sup_{n>N}|x_n|\to0$. The $n$th coordinate evaluation, bounded
by the norm, proves uniqueness. Truncations are contractions and
fix nonzero unit vectors, so their norms are $1$.
The uniform closure of finitely supported vectors is $c_0$:
truncations approximate every element of $c_0$, and a uniform
limit of sequences tending to zero still tends to zero.
Every space with a countable Schauder basis is separable by finite
rational combinations of its basis vectors. The uncountable family
of $\{0,1\}$-valued sequences has mutual $\ell^\infty$ distance
$1$, so $\ell^\infty$ is nonseparable and cannot have such a basis.
\end{proof}

\begin{remark}[Hamel, Schauder and orthonormal bases]\label{fa2:p1:c15:three-bases}
A Hamel expansion has only finitely many nonzero coefficients.
The vector $(1/n)_{n\geq1}\in\ell^2$ is not a finite linear
combination of the unit vectors, so that orthonormal Schauder
basis is not a Hamel basis. A countable complete orthonormal
system is an unconditional Schauder basis by
Theorem~\ref{fa2:p2:orthonormal-basis}; an uncountable Hilbert orthonormal basis
uses a finite-subset net and is not an ordered Schauder sequence.
The coordinate-continuity theorem above relies on completeness
of the original space and does not assert automatic continuity
for a basis defined on an arbitrary incomplete normed space.
\end{remark}

% FAMILY D16.C15.04
\section{Unconditional convergence at its basic level}
\begin{definition}[Unconditional series and basis]\label{fa2:p1:c15:unconditional}
A series $\sum_n v_n$ is unconditionally convergent if the net
$\sum_{n\in F}v_n$, over finite $F\subset\N$, converges in norm.
A Schauder basis is unconditional if each of its expansions is
unconditionally convergent. This is stronger than convergence in
the specified natural order.
\end{definition}

\begin{theorem}[Tail, rearrangement and subseries criteria]\label{fa2:p1:c15:unconditional-criteria}
For a series in a Banach space the following are equivalent:
unconditional convergence; the condition
\[
 \forall\eps>0\ \exists N\ \forall F\subset\{N+1,N+2,\ldots\}
 \text{ finite},\quad\left\|\sum_{n\in F}v_n\right\|<\eps;
\]
convergence of every subseries; and convergence of every
permuted series. All rearrangements then have the same sum.
Absolutely norm-summable series are unconditional; the converse
fails even in Hilbert space.
\end{theorem}
\begin{proof}
A convergent finite-subset net is Cauchy. Compare sums over
$F_0$ and $F_0\cup F$ for finite $F$ outside a sufficiently large
fixed $F_0$ to obtain the tail condition. Conversely the tail
condition makes the difference of two sufficiently advanced finite
sums small, by splitting their symmetric difference into two
tail sums. Completeness then gives a limit: the ordinary initial
sums are Cauchy and converge, and the same estimates compare any
finite sum containing $\{1,\ldots,N\}$ to them. The tail condition
also makes every subseries Cauchy. The prefixes of every
permutation eventually contain any prescribed finite set, so each
permuted series converges to the finite-subset limit.

If the tail condition fails, some $\eps>0$ permits, successively,
finite sets $F_k$ with $\max F_k<\min F_{k+1}$ and
$\norm{\sum_{n\in F_k}v_n}\geq\eps$. The subseries over
$\bigcup_kF_k$ has successive block differences of norm at least
$\eps$, so it is not Cauchy. Thus convergence of all subseries
forces the tail condition. The same failure constructs a bad
permutation: at stage $k$ first list the smallest unused integer,
then choose and list such a finite bad set beyond all indices
already listed. This lists every integer eventually and the bad
blocks prevent Cauchy convergence. Thus convergence of every
permuted series also forces the condition.
If $\sum\norm{v_n}<\infty$, the triangle inequality gives the
tail condition. In $\ell^2$, the series $\sum_n e_n/n$ is
unconditional by its square-summable orthogonal tails, although
$\sum_n\norm{e_n/n}=\sum_n1/n$ diverges.
\end{proof}

\begin{proposition}[Coordinate subspace projections for an unconditional basis]\label{fa2:p1:c15:subset-projections}
If $(e_n)$ is unconditional, for every $A\subset\N$ the map
$P_Ax=\sum_{n\in A}e_n'(x)e_n$ is a bounded projection onto
$\clspan\{e_n:n\in A\}$. There is one constant $C$ with
$\norm{P_A}\leq C$ for every $A$. The standard bases of $c_0$
and $\ell^p$, $p<\infty$, are unconditional, as are countable
complete Hilbert orthonormal systems.
\end{proposition}
\begin{proof}
Finite-subset coordinate maps $P_F$ are bounded by continuity of
each coordinate. For fixed $x$, their values are bounded: choose
a finite set after which all tail sums have norm at most $1$;
there are only finitely many sums within this fixed set, and
every other sum splits into one of these and a tail. Uniform
boundedness on the Banach space $E$ gives
$\sup_{F\text{ finite}}\norm{P_F}=C<\infty$. The subseries
criterion constructs $P_Ax$ and passing the bound to its limit
gives $\norm{P_Ax}\leq C\norm x$. Linearity follows by limits.
Coordinates, which are continuous, show that $P_A$ fixes exactly
the $A$ coordinates and kills the others. Its image lies in the
closed span of those vectors, and it fixes their finite span;
by continuity it fixes that closed span. This proves the range
and projection assertions. The canonical sequence-space tails
decrease when any coordinates are discarded, proving their
tail condition. Hilbert orthonormal tails were proved in
Theorem~\ref{fa2:p2:orthonormal-basis}.
\end{proof}

% FAMILY D16.C15.05
\section{Complementation and the scope of approximation}
\begin{proposition}[Bounded projections and topological sums]\label{fa2:p1:c15:complementation}
For a closed subspace $M$ of a Banach space $E$, being the range
of a bounded projection is equivalent to the existence of a
closed subspace $N$ such that $E=M\oplus N$ algebraically.
The projection is then the coordinate map $m+n\mapsto m$.
Every closed Hilbert subspace is complemented, and every
finite-dimensional or closed finite-codimensional Banach subspace
is complemented.
\end{proposition}
\begin{proof}
If $P^2=P$ is bounded, $\ran P=\ker(I-P)$ and $\ker P$ are
closed, and $x=Px+(I-P)x$ is the unique required split.
Conversely addition $M\times N\to E$ is a bounded bijection
between Banach spaces. Its inverse is bounded by the bounded
inverse theorem, making both coordinates bounded. The Hilbert
case is orthogonal decomposition. For a finite-dimensional $M$,
extend its finitely many coordinate functionals to $E$ by
Hahn--Banach; the sum of these extended coordinates times the
basis vectors is a bounded projection. If $E/M$ is finite
dimensional, choose representatives of a basis of $E/M$. Their
span $N$ maps bijectively onto $E/M$ with bounded inverse by
finite dimension, giving a bounded complementary coordinate
$E\to E/M\to N$ and the desired projection onto $M$.
\end{proof}

\begin{example}[Elementary failures of algebraic intuition]\label{fa2:p1:c15:complement-counterexamples}
The dense proper subspace $c_{00}\subset\ell^2$ cannot be the
range of a bounded projection, since that range would be closed.
Even a projection onto a closed one-dimensional range can be
unbounded if its complementary kernel is chosen only algebraically:
extend the independent vectors $e_n$ to a Hamel basis of $\ell^2$,
define an algebraic functional $f$ by $f(e_n)=n$ and arbitrary
(say zero) values on the added basis elements, and put
$Px=f(x)e_1$. Then $P^2=P$ because $f(e_1)=1$, but
$\norm{Pe_n}=n$ while $\norm{e_n}=1$.
Its range also has the usual bounded orthogonal projection;
complementation concerns the existence of a bounded split,
not boundedness of every algebraic split.
\end{example}

\begin{proposition}[Basis truncations approximate uniformly on compact sets]\label{fa2:p1:c15:compact-approximation}
For a Schauder basis, $P_N\to I$ uniformly on every norm-compact
subset of $E$. Consequently if $T:X\to E$ is compact, with
$X$ normed, then $P_NT$ has finite rank and
$\norm{P_NT-T}\to0$.
\end{proposition}
\begin{proof}
Let $K$ be compact and $K_0=\sup_N\norm{P_N}<\infty$.
Given $\eps>0$, choose finitely many points $z_j$ forming a
$\eps/(2(1+K_0))$-net for $K$. Since $P_Nz_j\to z_j$, for all
large $N$ the errors at all centers are less than $\eps/2$.
For $x\in K$ near $z_j$,
$\norm{(I-P_N)x}\leq(1+K_0)\norm{x-z_j}
+\norm{(I-P_N)z_j}<\eps$.
Apply this to the compact closure of $T(B_X)$ and take the
supremum over $B_X$. Each $P_NT$ has range in the finite initial
span. This proof depends on an actual Schauder basis of the
target. It supplies no unrestricted approximation-property
assertion for arbitrary Banach spaces.
\end{proof}

% END INPUT


% BEGIN parts/part1-bases-advanced.tex
\chapter{Basic sequences, dual bases and James's space}
\label{fa2:p1:bases-advanced}
The coordinates of a Schauder basis are continuous, but the existence of
an algebraic basis does not imply the existence of a Schauder basis.
This chapter develops the distinctions that become visible only in
infinite-dimensional Banach spaces. The principal source is
\cite[Chapters 3--4]{B05}.

\section{Perturbation and selection}
\begin{proposition}[Equivalent basic sequences and small perturbations]
\label{fa2:p1:equivalent-bases}
Basic sequences $(x_n)$ and $(y_n)$ are equivalent if and only if there
are $0<c\leq C<\infty$ such that
\[
 c\Big\|\sum a_nx_n\Big\|\leq\Big\|\sum a_ny_n\Big\|
 \leq C\Big\|\sum a_nx_n\Big\|
\]
for all finitely supported scalar sequences. If $(x_n)$ has basis
constant $K$ and the two sequences lie in one Banach space, the condition
\[
 2K\sum_n\frac{\|y_n-x_n\|}{\|x_n\|}<1
\]
suffices for equivalence. A Banach space with a basis has an equivalent
basis contained in any prescribed dense subset.
\end{proposition}
\begin{proof}
The two inequalities extend the correspondence $x_n\mapsto y_n$ and its
inverse by continuity to the closed spans. Conversely an isomorphism
gives the inequalities with its norm and inverse norm. If $x_n'$ are
the coordinate functionals, then
$x_n'(x)x_n=(P_n-P_{n-1})x$, so
$\|x_n'\|\leq2K/\|x_n\|$. Extend these functionals to the ambient
space by Hahn--Banach and set
$Ax=\sum_nx_n'(x)(y_n-x_n)$. The displayed sum of norms converges and
$\|A\|<1$. The Neumann series makes $I+A$ an automorphism of the
ambient space, and it maps the closed span of $(x_n)$ onto that of
$(y_n)$. Finally choose $y_n$ in the dense subset with weighted errors
whose sum is smaller than $1/(2K)$. If $(x_n)$ spans the whole space,
the automorphism is onto the whole space, so $(y_n)$ is a basis there.
\end{proof}

\begin{theorem}[Existence and weakly null selection]
\label{fa2:p1:basic-selection}
For every $M>1$, every infinite-dimensional Banach space contains a
normalized basic sequence with basis constant at most $M$. Every
weakly null sequence which does not tend to zero in norm has a basic
subsequence.
\end{theorem}
\begin{proof}
Choose positive $\eps_n$ with $\prod_n(1+\eps_n)<M$. On each
finite-dimensional span $E_n$ a finite subset $F_n$ of the dual unit
ball norms its unit sphere within the factor $1+\eps_n$; this follows
by choosing a norming functional at each sphere point and a finite
subcover. Having chosen $x_1,\ldots,x_n$, choose a unit $x_{n+1}$ in
the common kernels of all functionals so far selected. Their common
kernel has finite codimension and remains infinite dimensional. Since
$F_n$ norms $E_n$,
\[
 \|u\|\leq(1+\eps_n)\|u+tx_{n+1}\|
 \quad(u\in E_n,t\in\K).
\]
Iteration bounds every initial partial sum by $M$ times every longer
one. The basic-sequence criterion proved earlier makes $(x_n)$ a
basis for its closed span.

For the second statement pass to a subsequence with norms bounded
below by $a>0$; weak convergence gives a uniform upper norm bound.
The same construction can be performed with small corrections of
chosen sequence terms. Indeed for a finite family $f_1,\ldots,f_r$,
the map $Sx=(f_1(x),\ldots,f_r(x))$ has finite-dimensional range. Choose
a bounded linear right inverse $R$ on that range. Weak nullity gives
$RSx_n\to0$ in norm; hence $x_n-RSx_n$ lies in $\ker S$ and is as
close as desired to $x_n$. At stage $n$ choose this error smaller than
$a2^{-n}\eta$, with $\eta$ chosen so that the perturbation sum in the
preceding proposition is less than one, and also small enough to keep
all corrected norms above $a/2$. The corrected sequence has the
partial-sum inequalities just proved; the perturbation proposition
then transfers basicity to the selected original sequence.
\end{proof}

\begin{proposition}[Monotone and multiplier norms]
\label{fa2:p1:basis-renorming}
For a basis $(x_n)$ the equivalent norm
$\|x\|_b=\sup_n\|P_nx\|$ makes it monotone. For an unconditional
basis, $\|x\|_u=\sup_{|a_n|\leq1}\|\sum a_nx_n'(x)x_n\|$ is finite,
equivalent, and makes every coordinate multiplier of modulus at most
one contractive. On a real space this is a lattice norm for the
coordinate order. If $\|x_n\|=1$, coordinatewise multiplication makes
the space, with this latter norm, a commutative Banach algebra without
identity in infinite dimension.
\end{proposition}
\begin{proof}
The basis bound gives $\|x\|\leq\|x\|_b\leq K\|x\|$ and
$\|P_mx\|_b\leq\|x\|_b$. For an unconditional basis, the uniform
boundedness principle applied to finite coordinate multipliers gives
a uniform norm bound: unconditional convergence makes their values
bounded for each fixed $x$. Finite multipliers extend to arbitrary
bounded multipliers by the unconditional tail criterion. Composing
two multipliers proves contractivity in $\|\cdot\|_u$.
If $|a_n|\leq|b_n|$, multiplication by $a_n/b_n$ (zero where $b_n=0$)
proves lattice monotonicity. Coordinate projections are contractive
for the multiplier norm, and $\|x_n\|_u=\|x_n\|=1$ under the
normalization hypothesis. Thus $|x_n'(y)|\leq\|y\|_u$ for every $n$.
Multiplication of $x$ by the bounded sequence $(x_n'(y))$ therefore
defines $xy$ and gives $\|xy\|_u\leq\|x\|_u\|y\|_u$. Associativity
and commutativity hold coordinatewise. An identity would have all
coordinates one; its basis terms would fail to tend to zero because
$\|x_n\|=1$, a contradiction.
\end{proof}

\section{The sequence-space obstructions}
\begin{theorem}\label{fa2:p1:no-czero-wuc}
A Banach space contains no subspace isomorphic to $c_0$ if and only if
every weakly unconditionally Cauchy series is unconditionally norm
convergent. Here weakly unconditionally Cauchy means
$\sum_n|f(x_n)|<\infty$ for every $f\in E'$.
\end{theorem}
\begin{proof}
Uniform boundedness on the Banach space $E'$ shows that the finite
signed sums $\sum_{n\in F}a_nx_n$, $|a_n|\leq1$, have a common norm
bound $C$. If the series is not unconditionally convergent, the tail
Cauchy criterion gives disjoint successive finite sets $F_k$ with
$u_k=\sum_{n\in F_k}x_n$ of norm at least $\eps>0$. Absolute scalar
summability makes $(u_k)$ weakly null. Select a basic subsequence by
Theorem~\ref{fa2:p1:basic-selection}. It has a common upper bound on
all finite scalar sums with coefficients of modulus at most one,
inherited from the original series. Consequently
\[
 c\max|a_k|\leq\Big\|\sum a_ku_k\Big\|\leq C\max|a_k|.
\]
The lower bound follows from the uniform coordinate-functional bound
and $\inf\|u_k\|>0$. These inequalities identify its closed span
with $c_0$. Conversely the unit-vector series in $c_0$ is weakly
unconditionally Cauchy, since $c_0'=\ell^1$, and it does not converge
in norm. An isomorphic embedding transports this counterexample.
\end{proof}

\begin{theorem}[Orlicz--Pettis]\label{fa2:p1:orlicz-pettis}
If every subseries of $\sum x_n$ converges weakly in a Banach space,
then the series converges unconditionally in norm.
\end{theorem}
\begin{proof}
For every $f$, all subseries of the scalar series $\sum f(x_n)$
converge, so the scalar series is absolutely convergent. If norm
unconditional convergence failed, the preceding proof would produce
successive block sums $(u_k)$ equivalent to the unit basis of $c_0$.
The subseries consisting of all terms in those blocks has a weak
limit $u$ by hypothesis, and its block-end partial sums are
$\sum_{k\leq m}u_k$. Their weak limit belongs to the norm-closed
linear span of the $u_k$, which is weakly closed by Hahn--Banach.
Every coordinate functional of this basic sequence has value one
on $u$. Under the isomorphism with $c_0$ this would be the constant
sequence one, an impossibility. Thus the unconditional Cauchy
criterion holds, and completeness supplies the sum.
\end{proof}

\begin{theorem}[Shrinking and boundedly complete bases]
\label{fa2:p1:dual-bases}
For a basis $(x_n)$ of $E$ let $F=\overline{\operatorname{span}}\{x_n'\}$
in $E'$. The coordinate functionals form a basic sequence. The basis
is shrinking precisely when $F=E'$, equivalently
$\|f|_{\overline{\operatorname{span}}\{x_j:j>n\}}\|\to0$ for every
$f\in E'$. It is boundedly complete if every scalar sequence with
bounded initial sums $\sum_{j\leq n}a_jx_j$ has a convergent series.
The following hold:
\begin{enumerate}
\item $(x_n')$ is boundedly complete in $F$ if and only if $(x_n)$ is
shrinking;
\item $(x_n')$ is shrinking in $F$ if and only if $(x_n)$ is boundedly
complete;
\item $E$ is reflexive if and only if its basis is both shrinking and
boundedly complete.
\end{enumerate}
\end{theorem}
\begin{proof}
The adjoints $P_n'$ restrict to uniformly bounded initial projections
on the span of the $x_n'$, proving their basicity. The tail formula
$f-P_n'f=f(I-P_n)$ proves the shrinking criterion in both directions,
using the bound for $I-P_n$. For every $f\in E'$,
$P_n'f=\sum_{j\leq n}f(x_j)x_j'$ converges weak-star to $f$.

We give the two coefficient identifications underlying the assertions.
Each $\Phi\in F'$ gives coefficients $a_n=\Phi(x_n')$. The norming
property of $F$ (its finite coordinate projections norm $E$ up to
the basis constant) implies boundedness of
$s_n=\sum_{j\leq n}a_jx_j$. Conversely boundedness of these sums makes
$\sum a_j f(x_j)$ a bounded linear functional on the finite span of
the $x_j'$, extending uniquely to $F$. These constructions are inverse.
The canonical image of $E$ in $F'$ therefore consists exactly of
the coefficient sequences whose series $\sum a_jx_j$ converges.
It is onto exactly when the basis is boundedly complete. Its range
is also the closed span of the coordinate functionals of $(x_n')$,
because those are the restrictions of $J_Ex_n$ to $F$. Thus it is
onto exactly when $(x_n')$ is shrinking. This proves assertion 2.

To prove assertion 1, if $F=E'$ and the partial sums
$\sum_{j\leq n}b_jx_j'$ are bounded, they have a weak-star cluster
point $f\in E'$ by Alaoglu. Its coefficients are $f(x_j)=b_j$;
shrinking gives norm convergence to $f$. Conversely if the dual
basic sequence is boundedly complete, apply its criterion to the
bounded sequence $P_n'f$ for arbitrary $f\in E'$. Its norm limit
belongs to $F$ and its weak-star limit is $f$, so $f\in F$.

If $E$ is reflexive, bounded initial sums have a weak cluster point,
whose coordinates determine the desired basis sum; hence the basis
is boundedly complete. If $F\ne E'$, Hahn--Banach gives a nonzero
$\Phi\in E''$ annihilating $F$. Writing $\Phi=J_Ex$ would make every
coordinate of $x$ zero, a contradiction; hence the basis is shrinking.
Conversely, if both properties hold and $\Phi\in E''$, the bounded
sums $\sum_{j\leq n}\Phi(x_j')x_j$ converge to $x$. On each $x_j'$
the functionals $\Phi$ and $J_Ex$ agree, and shrinking makes their
span dense in $E'$. They agree on all of $E'$, so $J_E$ is onto.
\end{proof}

\begin{theorem}\label{fa2:p1:unconditional-reflexivity}
For a Banach space with an unconditional basis, that basis is shrinking
if and only if no $\ell^1$ embeds, and boundedly complete if and only
if no $c_0$ embeds. Consequently such a space is reflexive if and
only if it contains neither $c_0$ nor $\ell^1$.
\end{theorem}
\begin{proof}
If shrinking fails, the tail criterion produces a normalized block
sequence $(u_k)$ and $f\in E'$ with $|f(u_k)|\geq\delta>0$.
Adjust phases so these values are positive real. Unconditional
coordinate multipliers change phases of whole blocks with uniformly
bounded norm. Thus
\[
 c\sum|a_k|\leq\Big\|\sum a_ku_k\Big\|\leq\sum|a_k|,
\]
giving an $\ell^1$ subspace. If shrinking holds, $E'$ is separable.
An $\ell^1$ subspace would give a continuous surjection
$E'\to(\ell^1)'=\ell^\infty$ by Hahn--Banach, contradicting
nonseparability of $\ell^\infty$.

If bounded completeness fails, choose bounded initial sums with
non-Cauchy tails and choose successive block sums $(u_k)$ with norms
bounded below. Uniformly bounded coordinate multipliers and the
bound for the original initial sums show
$\|\sum a_ku_k\|\leq C\max|a_k|$. The coordinate estimate gives
the reverse bound, so their span is $c_0$. Conversely suppose a
$c_0$ subspace is present. Its unit basis is weakly null. A successive
subsequence can be approximated by successive blocks of the
ambient basis: at each stage its finitely many initial coordinates
tend to zero, and after choosing a term its sufficiently long basis
tail is small. Choose the errors summably small, so the perturbation
theorem makes the blocks equivalent to $c_0$. Their initial sums
are bounded, but they cannot converge since their norms are bounded
below. Bounded completeness of the ambient basis would force their
sum to converge: partial sums cut inside a block are bounded too,
since an ambient initial projection of a completed block sum has
norm at most the basis constant times that uniformly bounded sum.
This is a contradiction. The reflexivity conclusion now follows
from Theorem~\ref{fa2:p1:dual-bases}. In particular an unconditional
basic sequence spans a reflexive space or has a further subspace
isomorphic to $\ell^1$ or $c_0$.
\end{proof}

\section{Function-space bases}
\begin{theorem}\label{fa2:p1:schauder-haar}
The dyadic hat functions, together with $1,t$, form the classical
Schauder basis of $C[0,1]$. The Haar functions, ordered by successive
dyadic subdivisions and preceded by $1$, form a basis of $L^p[0,1]$
for $1\leq p<\infty$. The former basis and the Haar basis in $L^1$
are conditional. The Haar basis is unconditional when $1<p<\infty$.
\end{theorem}
\begin{proof}
For the hat functions the initial projections are piecewise linear
interpolants on the successive dyadic partitions. Their sup norm is
at most one, and uniform continuity gives convergence of the
interpolants. Each new coefficient is the value at the new midpoint
minus the average of the two old endpoint values; thus coefficients
are unique and continuous. The basis criterion proves the assertion.
To see conditionality, follow the nested dyadic intervals containing
$1/3$. Their successive midpoints $m_k$ alternate to the left and
right of $1/3$, and, after the first two levels, the endpoints of
the interval at level $k$ are $m_{k-1},m_{k-2}$. Choose a continuous
piecewise linear $f_N$, of sup norm one, with
$f_N(m_k)=(-1)^k$ for $k\leq N$, and zero endpoint values. At every
level $k\geq3$ the corresponding hat coefficient is $(-1)^k$,
since the two endpoint values cancel. The value of each of those
normalized hats at $1/3$ is $2/3$. Change the signs of these
coefficients to positive. The new value at $1/3$ is at least
$2(N-2)/3-O(1)$. Hence the finite sign multipliers are not uniformly
bounded, proving conditionality.

For Haar functions the initial projections are conditional
expectations onto finite dyadic interval partitions. Jensen gives
$\|P_nf\|_p\leq\|f\|_p$. Their union spans the dyadic step functions,
which are dense for finite $p$; coefficient uniqueness follows by
integration against the Haar functions. Thus they form a basis.
In $L^1$ take $f_N=2^N1_{[0,2^{-N}]}$. Its dyadic differences are
\[
 d_k=2^{k-1}(1_{[0,2^{-k}]}-1_{(2^{-k},2^{-(k-1)}]}),
 \qquad f_N=1+\sum_{k=1}^Nd_k.
\]
On the annulus $(2^{-j},2^{-(j-1)}]$ the alternating transform
$1+\sum(-1)^kd_k$ has modulus at least $2^j/3$ for $j\geq2$.
The annulus has measure $2^{-j}$; its $L^1$ norm is therefore at
least $(N-1)/3$, whereas $\|f_N\|_1=1$.

For completeness, a uniform $L^p$ bound for finite Haar sign
transforms follows from a short dyadic decomposition argument.
Their $L^2$ norm is at most one by orthogonality. For an integrable
$f$ and $\lambda>0$ take maximal dyadic intervals $I$ with
$|I|^{-1}\int_I|f|>\lambda$; their total length is at most
$\|f\|_1/\lambda$. Outside their union keep $g=f$, and on $I$
replace $f$ by its mean. Outside the exceptional union every
dyadic average of $|f|$ is at most $\lambda$. These averages
converge in $L^1$ to $|f|$ by density of dyadic step functions,
so a subsequence converges almost everywhere; hence
$|f|\le\lambda$ there. On a maximal proper interval the
parent average is at most $\lambda$, so its own average is at
most $2\lambda$. Thus $|g|\leq2\lambda$ and
$\|g\|_1\leq\|f\|_1$. Each $b_I=(f-f_I)1_I$ has mean zero;
its Haar transform is supported in $I$, since larger dyadic Haar
functions are constant on $I$. Outside the union only $Tg$ remains,
and Chebyshev plus the $L^2$ bound gives
\[
 |\{|Tf|>\lambda\}|\leq C\|f\|_1/\lambda.
\]
If the whole interval is exceptional, the same inequality follows
from $1\leq\|f\|_1/\lambda$; thus the argument covers every
$\lambda$. The constant is independent of the signs and truncation.
For $1<p<2$ split $f$ into its parts where $|f|>a\lambda$ and
$|f|\leq a\lambda$. Applying weak $(1,1)$ to the first and strong
$(2,2)$ to the second, integrate the distribution function against
$p\lambda^{p-1}d\lambda$. Tonelli gives the integrals
\[
 \int_0^{|f|/a}\lambda^{p-2}d\lambda
 =\frac{(|f|/a)^{p-1}}{p-1},\qquad
 \int_{|f|/a}^{\infty}\lambda^{p-3}d\lambda
 =\frac{(|f|/a)^{p-2}}{2-p},
\]
so $\|Tf\|_p\leq C_p\|f\|_p$. Duality gives the bound for
$p>2$, since the adjoint transform has conjugate signs and obeys
the same estimates. Uniform finite sign bounds together with basis
convergence imply unconditionality by the tail criterion.
\end{proof}

\section{A concrete noncomplementation theorem}
\begin{theorem}\label{fa2:p1:czero-uncomplemented}
The closed subspace $c_0$ of $\ell^\infty$ is not complemented.
Its canonical image is not complemented in $c_0''$, and $c_0$ is
not isomorphic to the dual of any normed space.
\end{theorem}
\begin{proof}
Call a Banach space countably separated if countably many continuous
linear functionals separate its points. This property passes to
subspaces and to isomorphic spaces. Coordinate evaluations countably
separate $\ell^\infty$. If a complement $N$ to $c_0$ existed,
$\ell^\infty/c_0\simeq N$ would therefore be countably separated.

Choose an uncountable almost disjoint family $(A_\alpha)$ of infinite
subsets of $\N$: identify $\N$ with the rationals in $(0,1)$ and
for each irrational $\alpha$ take a distinct rational sequence
converging to $\alpha$. Distinct such sets have finite intersection.
Put $z_\alpha=1_{A_\alpha}+c_0$. Distinct cosets are distinct.
For finitely many distinct indices and $|a_j|=1$,
$\|\sum a_jz_{\alpha_j}\|=1$: outside finitely many intersections
each coordinate has at most one nonzero summand, and infinitely
many coordinates have modulus one. If $f$ is a quotient functional
and $|f(z_{\alpha_j})|\geq1/m$, choose phases making all values
positive. Then $r/m\leq\|f\|$, where $r$ is the number of chosen
indices. Thus $f(z_\alpha)\ne0$ for only countably many $\alpha$.
Countably many $f$ cannot separate the uncountably many remaining
$z_\alpha$, a contradiction.

Under the explicit sequence-dual identifications $c_0''=\ell^\infty$
the canonical map is the usual inclusion, giving the bidual
assertion. Every dual $Y'$ is complemented in $Y'''$: the operator
$J_{Y'}J_Y':Y'''\to Y'''$ is a bounded projection onto $J_{Y'}Y'$
because $J_Y'J_{Y'}=I$. If $E\simeq Y'$, conjugating this projection
by the bidual isomorphism shows $J_EE$ complemented in $E''$.
This is impossible for $E=c_0$.
\end{proof}

\section{The cyclic-variation James space}
\begin{definition}\label{fa2:p1:james-space-definition}
For a real null sequence $a=(a_n)$ define
\[
 \|a\|_J=2^{-1/2}\sup_{m\geq2,\,p_1<\cdots<p_m}
 \left(\sum_{j=1}^{m-1}(a_{p_j}-a_{p_{j+1}})^2
                  +(a_{p_m}-a_{p_1})^2\right)^{1/2}.
\]
The real James space $J$ consists of the null sequences for which
this is finite. The cyclic last difference is part of this particular
norm; replacing it without analysis would lose the isometric statement
below.
\end{definition}

\begin{theorem}\label{fa2:p1:james-space}
$J$ is a real Banach space. Its unit vectors form a normalized monotone
shrinking basis. The canonical image $J_JJ$ has codimension one in
$J''$, but $J$ and $J''$ are isometrically isomorphic by a different
map. Consequently an abstract isometry onto the bidual does not imply
reflexivity.
\end{theorem}
\begin{proof}
Finite Euclidean triangle inequalities show that $\|\cdot\|_J$ is a
seminorm; for a null sequence, taking a second index tending to
infinity gives $|a_n|\leq\|a\|_J$, so it is a norm. A Cauchy
sequence converges uniformly in the coordinates to a null sequence
$a$. Passing to the limit in each finite cyclic expression shows
convergence in $J$ and finiteness of its norm, hence completeness.

Truncations are contractive. Indeed any finite cyclic list in the
truncated sequence either has no trailing zero, in which case it
is already a list in $a$, or has a trailing zero, which can be
approximated by $a_q$ with $q$ arbitrarily large. Repeated zeros
do not increase the expression. If the truncation tails did not
tend to zero in $J$, disjoint successive finite lists in arbitrarily
late tails would each have cyclic variation bounded below.
Since $a_n\to0$, their closing differences tend to zero. Joining
finitely many such lists would make the sum of their internal
squared differences arbitrarily large, contradicting $a\in J$.
Thus the unit vectors form a monotone basis, and each has norm one.

Let $\|\cdot\|_o$ be the same supremum without the cyclic closing
term, still multiplied by $2^{-1/2}$. Then
$\|a\|_o\leq\|a\|_J\leq2\|a\|_o$.
For normalized successive blocks $u_k$, and finitely many scalars
$b_k$, partition a finite chain of coordinate differences into
those within one block and those crossing between blocks.
The within-block sum is at most $2\sum b_k^2$, since each block
has $J$-norm one. For crossing differences use
$(b_j\alpha-b_k\beta)^2\leq2b_j^2+2b_k^2$ and
$|\alpha|,|\beta|\leq1$. A block enters at most once and leaves
at most once along an increasing chain, so these terms total at
most $4\sum b_k^2$. Hence
\[
 \Big\|\sum b_ku_k\Big\|_J\leq\sqrt{12}
                       \Big(\sum b_k^2\Big)^{1/2}.
\]
If the basis were not shrinking, the tail-functional criterion
would give normalized blocks with $f(u_k)\geq\delta>0$ after
changing their real signs. The displayed estimate shows
$\sum k^{-1}u_k$ converges in $J$, while its scalar $f$-series
dominates $\delta\sum k^{-1}$, a contradiction. Thus the basis
is shrinking.

For any monotone shrinking basis, its bidual is the space of
coefficient sequences with bounded initial sums, with norm
$\sup_n\|\sum_{j\leq n}a_je_j\|$. To verify this directly,
associate $a_j=\Phi(e_j')$ to $\Phi\in J''$; the partial-sum
norm is at most $\|\Phi\|$. Conversely bounded partial sums
define on finite linear combinations of $e_j'$ the bounded
functional $\sum b_je_j'\mapsto\sum a_jb_j$. It extends to $J'$
because the basis is shrinking. The two constructions are inverse
isometries.

Every sequence in this bidual model has a finite limit. If it were
not Cauchy in the scalar field, alternating indices
$p_1<q_1<p_2<q_2<\cdots$ with $|a_{p_j}-a_{q_j}|\geq\eps$
would force unbounded variation of its initial sums. Boundedness
of the partial sums also bounds every coordinate. If $L=\lim a_n$,
then $(a_n-L)$ is null and has finite cyclic variation: subtraction
of a constant does not change finite differences. It therefore
belongs to $J$. Conversely the constant sequence one has bounded
initial sums of norm one. Thus the bidual is $J\oplus\R\mathbf1$
as a vector space, and its canonical image is exactly the null
sequences, proving codimension one.

Finally define
\[
 T(a_1,a_2,\ldots)=(-L,a_1-L,a_2-L,\ldots),\qquad L=\lim a_n.
\]
The target is a null sequence in $J$. The bidual norm is the
supremum of cyclic lists of entries of $a$, with the option of
one extra trailing zero coming from truncation. A cyclic rotation
moves that zero to the beginning of the list. Subtracting $L$
from every entry then gives exactly the finite cyclic lists of
$Ta$. Hence $\|Ta\|_J=\|a\|_{J''}$.
Given $b\in J$, set $a_n=b_{n+1}-b_1$. Its limit is $-b_1$,
it belongs to the bidual model, and $Ta=b$. The map $T$ is
therefore a surjective real-linear isometry. It is not the
canonical inclusion, which has the nonzero one-dimensional
quotient just computed.
\end{proof}

% END INPUT


% BEGIN parts/part1-geometry.tex
\chapter{Convexity, smoothness and approximation geometry}
\label{fa2:p1:geometry}
Reflexivity concerns the canonical bidual map; strict and uniform
convexity concern the particular norm. An equivalent norm preserves the
former and need not preserve the latter. We keep these distinctions
throughout this chapter; compare \cite[Chapter 5]{B05}.

\section{Strict convexity and nearest points}
\begin{definition}\label{fa2:p1:rotund-definitions}
A norm is strictly convex, or rotund, if distinct unit vectors have a
midpoint of norm less than one. The norm is locally uniformly convex
at a unit vector $x$ if unit $y_n$ with $\|(x+y_n)/2\|\to1$ satisfy
$y_n\to x$ in norm. It is uniformly convex if
\[
 \delta_E(\eps)=\inf\{1-\|(x+y)/2\|:\|x\|,\|y\|\leq1,
                                      \ \|x-y\|\geq\eps\}>0
\]
for every $0<\eps\leq2$. For the zero space define
$\delta_E(\eps)=1$ when $\eps>0$, and always put $\delta_E(0)=0$;
the defining epsilon--delta implication is vacuous there. No real
infimum of an empty set is being used. The Radon--Riesz property means that weak
convergence together with convergence of norms implies norm convergence.
\end{definition}

\begin{theorem}\label{fa2:p1:strict-approximation}
Strict convexity is equivalent to each of the following: equality in
the triangle inequality for two nonzero vectors requires one to be a
positive real multiple of the other; every convex set has at most one
nearest point to any given vector; each supporting hyperplane to the
unit ball meets its sphere at at most one point. In a reflexive strictly
convex Banach space every nonempty closed convex set has a unique
nearest point.
\end{theorem}
\begin{proof}
If equality holds for $x,y\ne0$, their normalized vectors have a convex
combination of norm one with weights $\|x\|/(\|x\|+\|y\|)$ and its
complement. A nontrivial chord of the unit sphere would then contain
a norm-one midpoint, contrary to strict convexity. The converse is
immediate. Two distinct minimizing points in a convex set would have
a midpoint closer to the given vector, except when the common distance
is zero, in which case both points equal that vector. Conversely, if
distinct unit $x,y$ have a norm-one midpoint, the segment $[x,y]$ has
all its points at distance one from zero, by a norming functional at
the midpoint, contradicting uniqueness. The same norming-functional
argument proves the supporting-hyperplane characterization.
For existence choose a bounded minimizing sequence in the closed convex
set; reflexivity gives a weakly convergent subsequence, the set is weakly
closed, and the norm is weakly lower semicontinuous. Its limit is a
minimizer. Strict convexity supplies uniqueness.
\end{proof}

\begin{theorem}\label{fa2:p1:uc-stability}
Uniform convexity passes to subspaces and to quotients by closed
subspaces. Locally uniform convexity implies Radon--Riesz. Finite
$\ell^p$ sums, $1<p<\infty$, of uniformly convex spaces are uniformly
convex; the same is true for arbitrary sums if their moduli have a
common positive lower bound at every positive argument.
\end{theorem}
\begin{proof}
The subspace statement restricts the defining infimum. For a quotient,
lift unit-ball elements $u,v$ to $x,y$ of norms at most $1+\eta$.
If $\|u-v\|\geq\eps$, then $\|x-y\|\geq\eps$ and
\[
 \|(u+v)/2\|\leq(1+\eta)
      [1-\delta_E(\eps/(1+\eta))].
\]
Choose $\eta$ sufficiently small and use the fixed positive
$\delta_E(\eps/2)$ to get a uniform strict gap. To prove Radon--Riesz
under local uniform convexity, reduce by rescaling to $x_n\rightharpoonup x$
with $\|x_n\|\to\|x\|=1$. Weak lower semicontinuity gives
$\|(x+x_n/\|x_n\|)/2\|\to1$, and local uniform convexity at $x$
gives norm convergence. The zero-limit case is immediate.

We detail the sum argument to make its uniformity condition explicit.
Suppose unit vectors $x^{(n)},y^{(n)}$ in the sum have midpoint norms
tending to one. Set $a_i^{(n)}=\|x_i^{(n)}\|$ and
$b_i^{(n)}=\|y_i^{(n)}\|$. The scalar $\ell^p$ midpoint inequality
implies $\|a^{(n)}-b^{(n)}\|_p\to0$. The nonnegative deficits
\[
 d_i=(a_i+b_i)/2-\|(x_i+y_i)/2\|
\]
have weighted $p$-mass tending to zero. More precisely, put
$A_i=(a_i+b_i)/2$ and $c_i=\|(x_i+y_i)/2\|$. Then
$\sum_i(A_i^p-c_i^p)\to0$, and
$A_i^p-c_i^p\ge A_i^{p-1}(A_i-c_i)$, so
$\sum_iA_i^{p-1}d_i\to0$.
Fix $0<\eps<1$, let $\delta>0$ be a common lower bound of the
component moduli at $\eps$, and choose
$0<\eta\le\min(\eps,\delta,1)$.
The coordinates where $|a_i-b_i|>\eta A_i$ have total $A_i^p$-mass
at most $\eta^{-p}\|a-b\|_p^p$, which tends to zero.
On the remaining coordinates with $A_i>0$, write $u_i=x_i/a_i$,
$v_i=y_i/b_i$. Both denominators are positive. If
$\|x_i-y_i\|\ge2\eps A_i$, then
$\|u_i-v_i\|\ge2\eps-\eta\ge\eps$. Uniform convexity and
$\|(x_i+y_i)/2-A_i(u_i+v_i)/2\|\le\eta A_i/2$
give $c_i\le A_i(1-\delta+\eta/2)$, hence
$d_i\ge\delta A_i/2$. These coordinates consequently have total
$A_i^p$-mass at most $(2/\delta)\sum_iA_i^{p-1}d_i\to0$.
On both exceptional sets use $\|x_i-y_i\|\le2A_i$.
The remaining coordinates contribute at most
$(2\eps)^p\sum_iA_i^p$ to $\|x-y\|_p^p$. Letting
$\eps\downarrow0$ gives $\|x^{(n)}-y^{(n)}\|\to0$. This sequential
criterion is exactly uniform convexity: failure would choose one
pair at each midpoint deficit $1/n$ with fixed separation. Finite
sums have a common modulus by taking the minimum of finitely many
component moduli.
\end{proof}

\begin{theorem}[Clarkson inequalities]\label{fa2:p1:clarkson}
For any positive measure space and $f,g\in L^p$, the following hold.
For $2\leq p<\infty$,
\[
 \|f+g\|_p^p+\|f-g\|_p^p
       \leq2^{p-1}(\|f\|_p^p+\|g\|_p^p).
\]
For $1<p\leq2$, writing $q=p/(p-1)$,
\[
 \|f+g\|_p^q+\|f-g\|_p^q
       \leq2(\|f\|_p^p+\|g\|_p^p)^{q/p}.
\]
Thus $L^p$ is uniformly convex for $1<p<\infty$.
\end{theorem}
\begin{proof}
For the scalar two-by-two matrix $H(a,b)=(a+b,a-b)$, its norms
from $\ell^1_2$ to $\ell^\infty_2$, from $\ell^2_2$ to itself,
and from $\ell^\infty_2$ to itself are respectively $1,\sqrt2,2$.
The finite-matrix Riesz--Thorin theorem (proved in Part IV without
using reflexivity or uniform convexity) therefore gives
\[
 \|H\|_{\ell^p_2\to\ell^p_2}\leq2^{1-1/p}\ (p\geq2),\qquad
 \|H\|_{\ell^p_2\to\ell^q_2}\leq2^{1/q}\ (p\leq2).
\]
Integrating the first pointwise scalar inequality gives the first
formula. For the second use Minkowski for $q/p\geq1$:
\[
 \left(\|f+g\|_p^q+\|f-g\|_p^q\right)^{p/q}
 \leq\int\left(|f+g|^q+|f-g|^q\right)^{p/q}
 \leq2^{p/q}\int(|f|^p+|g|^p).
\]
Raising to $q/p$ proves the formula. For unit-ball $f,g$, dividing
by the appropriate power of two gives
$\|(f+g)/2\|^r+\|(f-g)/2\|^r\leq1$, with $r=p$ or $q$.
Therefore $\delta_{L^p}(\eps)\geq
1-[1-(\eps/2)^r]^{1/r}>0$. The same proof covers counting measure
and arbitrary index sets; the zero space is uniformly convex
vacuously, rather than by an infimum over a nonexistent sphere.
\end{proof}

\section{Smoothness and duality}
\begin{definition}\label{fa2:p1:smoothness-definitions}
At a unit vector $x$ put
$D(x)=\{f\in E':\|f\|=f(x)=1\}$, using the complex-linear dual
and real equality when the field is complex. The norm is smooth at
$x$ if $D(x)$ is a singleton. Define its modulus of smoothness by
\[
 \rho_E(t)=\max\left\{0,\sup_{\|x\|,\|y\|\leq1}
       \left(\frac{\|x+ty\|+\|x-ty\|}{2}-1\right)\right\}.
\]
The space is uniformly smooth if $\rho_E(t)/t\to0$ as $t\downarrow0$.
For nonzero spaces the same supremum can be taken over unit spheres,
by convexity in each vector and rescaling. The explicit maximum makes
$\rho_{\{0\}}=0$; the modulus identity in the next proof is used only
for nonzero spaces, with the zero-space equivalences checked directly.
\end{definition}

\begin{theorem}\label{fa2:p1:gateaux-smoothness}
The norm is smooth at $x\ne0$ if and only if it is real Gateaux
differentiable there. At a unit $x$ the derivative is
$h\mapsto\operatorname{Re}f_x(h)$ for the unique $f_x\in D(x)$.
A strictly convex dual makes the primal smooth; a smooth dual makes
the primal strictly convex. The converses hold in reflexive spaces.
\end{theorem}
\begin{proof}
Convexity gives the right directional derivative. For every
$f\in D(x)$,
$\|x+th\|\geq1+t\operatorname{Re}f(h)$ for $t>0$. Choose
$f_t\in D((x+th)/\|x+th\|)$. Then
\[
 \frac{\|x+th\|-1}{t}\leq\operatorname{Re}f_t(h).
\]
A weak-star subnet limit as $t\downarrow0$ belongs to $D(x)$,
because $\|f_t\|=1$ and $f_t(x)\to1$. Thus the right derivative
equals $\max_{D(x)}\operatorname{Re}f(h)$. The left derivative
equals the minimum. Uniqueness gives differentiability and its
formula; differentiability in every direction forces equality
of all real parts, hence equality of the complex functionals.

Two distinct members of $D(x)$ have a midpoint of norm one, so a
strictly convex dual forbids them. If $x,y$ are distinct unit
vectors with unit midpoint, a norming $f$ has $f(x)=f(y)=1$;
then $J_Ex,J_Ey$ are two distinct norming functionals at $f$ in
the dual space, contradicting smoothness of $E'$. If $E$ is
reflexive, apply these implications to $E'$ and identify $E''$
isometrically with $E$ to obtain both converses.
\end{proof}

\begin{theorem}[Uniform duality]\label{fa2:p1:uniform-smooth-duality}
A normed space is uniformly convex if and only if its dual is
uniformly smooth, and uniformly smooth if and only if its dual is
uniformly convex. Uniform smoothness is equivalent to uniform real
Fr\'echet differentiability of the norm on the unit sphere. It
implies reflexivity for Banach spaces. The norming map $x\mapsto f_x$
is then uniformly norm-to-norm continuous on the unit sphere.
\end{theorem}
\begin{proof}
For a nonzero space the exact dual modulus identity is
\[
 \rho_{E'}(t)=\sup_{0\leq\eps\leq2}
                      \{t\eps/2-\delta_E(\eps)\}.
\]
For the upper bound choose $f,g$ in the dual unit ball and unit
$x,y$ which nearly norm $f+tg$ and $f-tg$, with their phases adjusted
so the evaluations are real. Then
\[
 \|f+tg\|+\|f-tg\|
 \leq\|x+y\|+t\|x-y\|+o(1)
 \leq2-2\delta_E(\eps)+t\eps+o(1),
\]
where $\eps=\|x-y\|$. For the reverse bound take $x,y$ nearly
minimizing $\delta_E(\eps)$ and norming functionals for $x+y$
and $x-y$. Evaluating $f+tg$ at $x$ and $f-tg$ at $y$ gives
the reverse inequality, with $\|x-y\|\geq\eps$.
The usual approximation handles a zero vector being normed.

If all positive $\delta_E(\eps)$ are positive, split the
supremum into $\eps<\eta$ and $\eps\geq\eta$; for sufficiently
small $t$ the latter terms are nonpositive, while the former
are at most $t\eta/2$. Thus $\rho_{E'}(t)/t\to0$.
Conversely $\delta_E(\eps_0)=0$ would force
$\rho_{E'}(t)\geq t\eps_0/2$ for every $t$.
Moreover $\rho_{E''}=\rho_E$. One inequality follows from the
canonical embedding. For the other, nearly norm each of
$\Phi\pm t\Psi$ by two vectors of $E'$, and use Goldstine
simultaneously to approximate the evaluations of the two bidual
unit vectors by primal unit-ball vectors. The finite scalar
tests give the reverse inequality. Apply the first equivalence
to $E'$ to obtain the second.

If uniformly smooth, the one-sided derivative formula just proved
has a unique support functional, and convexity gives, for unit $x$,
\[
 0\leq\|x+h\|-1-\operatorname{Re}f_x(h)
                    \leq2\rho_E(\|h\|).
\]
This is the uniform Fr\'echet estimate. Conversely a uniform
Fr\'echet estimate applied to $h$ and $-h$ gives
$\rho_E(t)=o(t)$. For unit $x,y$ and unit $h$, comparison of the
two support inequalities with the displayed estimate yields
\[
 |(f_x-f_y)(h)|\leq2\rho_E(t)/t+\|x-y\|/t.
\]
In the complex case rotate $h$ to recover the modulus from its
real part. First choose $t$ small and then $\|x-y\|$ small;
this proves uniform continuity. Finally uniform smoothness makes
$E'$ uniformly convex and hence reflexive by Milman--Pettis.
Reflexivity of the dual implies reflexivity of a Banach predual,
as proved earlier.
\end{proof}

\begin{theorem}[Slices and minimizing sequences]
\label{fa2:p1:strong-rotundity}
The following are equivalent for a normed space:
\begin{enumerate}
\item for each $f\in E'$ with $\|f\|=1$, the diameters of the slices
$\{x:\|x\|\leq1,\operatorname{Re}f(x)>1-\delta\}$ tend to zero;
\item every sequence of unit vectors with
$\operatorname{Re}f(x_n)\to1$ for one norm-one $f$ is Cauchy;
\item every minimizing sequence in a nonempty convex set, for
distance to any fixed point, is Cauchy.
\end{enumerate}
This property, strong rotundity, implies strict convexity and
Radon--Riesz. In a Banach space it makes all such minimizing sequences
convergent. Reflexive strictly convex spaces with Radon--Riesz have
this property; in particular reflexive locally uniformly convex
spaces do.
\end{theorem}
\begin{proof}
Failure of the first condition chooses two points of slices at
$\delta=1/n$ with fixed positive separation; alternating and
normalizing them violates the second. The converse is immediate.
For condition 2 implying 3 translate the point to zero and let
$r=\inf_C\|x\|$. The case $r=0$ is immediate. Separate the open
ball $B(0,r)$ from $C$ to obtain a norm-one real part of a
functional $f$ with $\operatorname{Re}f(c)\geq r$. A minimizing
sequence $c_n$ has $\|c_n\|\to r$, so its normalized terms
satisfy condition 2 and hence it is Cauchy. Conversely, if
condition 2 fails, rescale the offending unit $x_n$ by the positive
numbers $\operatorname{Re}f(x_n)\to1$. The new $y_n$ satisfy
$\operatorname{Re}f(y_n)=1$, $\|y_n\|\to1$, and are not Cauchy.
They are a minimizing sequence for their convex hull, whose
distance from zero is one, contradicting condition 3.

Two distinct sphere points in one supporting hyperplane give a
constant alternating counterexample to condition 2, proving strict
convexity. If $x_n\rightharpoonup x$, $\|x_n\|\to\|x\|>0$,
a norming functional at $x$ puts the normalized $x_n$ in such
slices. They are norm Cauchy and weakly converge to $x/\|x\|$.
This implies norm convergence even without completeness: for fixed
$n$, weak lower semicontinuity gives
$\|z_n-x/\|x\|\|\le\liminf_m\|z_n-z_m\|$, whose right side tends
to zero by the Cauchy property. This proves Radon--Riesz. In a Banach space
Cauchy sequences converge. Finally, in a reflexive strictly
convex space a closed convex set has a unique nearest point.
Any bounded minimizing sequence has weakly convergent subsequences;
each weak limit is that point, and Radon--Riesz upgrades each
such subsequence to norm convergence. If the entire minimizing
sequence did not converge, a subsequence staying away from the
point would have a further convergent subsequence, a contradiction.
The same argument applies to the closure of an arbitrary convex
set. Local uniform convexity supplies Radon--Riesz by the preceding
theorem.
\end{proof}

% END INPUT


% BEGIN parts/part1-foundation-24.tex
\chapter{Compact convex geometry and fixed points}\label{fa2:p1:c24}

For the convex-geometric part, $E$ is a Hausdorff locally convex real
topological vector space; a complex space is regarded as real. All
closures in that part use the given locally convex topology. In the
fixed-point part we explicitly use a norm topology.
\source{B02 V.7, V.9; B05 Section 2.10.}
\cite{B02} \cite{B05} 


% FAMILY D16.C24.01
\section{Faces and extreme points}
\begin{definition}[Extreme points and faces]\label{fa2:p1:c24:faces}
For a convex set $K$, a point $x\in K$ is extreme if
$x=ty+(1-t)z$, $y,z\in K$, $0<t<1$, implies $y=z=x$.
The set of these points is $\operatorname{Ext}K$. A nonempty convex
subset $F\subset K$ is a face if every proper convex combination of
points of $K$ lying in $F$ has both endpoints in $F$.
A compact face is a face which is compact in the ambient topology.
\end{definition}

\begin{lemma}[Face operations and maximizing faces]\label{fa2:p1:c24:face-operations}
A singleton is a face exactly when its point is extreme. Nonempty
intersections of faces are faces, and a face of a face of $K$ is a
face of $K$. If $K$ is nonempty compact convex and $\ell$ a continuous real
linear functional, then
\[
 F_\ell=\{x\in K:\ell(x)=\max_K\ell\}
\]
is a nonempty compact face.
\end{lemma}
\begin{proof}
The singleton assertion is the definition. An intersection is convex,
and the endpoint implication holds in every intersectand. For transitivity,
a combination in the smaller face first puts its endpoints in the
larger face and then in the smaller one. A continuous real function
attains its maximum on a nonempty compact set, so $F_\ell$ is nonempty
and compact. If $tx+(1-t)y$ achieves that maximum, neither $\ell(x)$
nor $\ell(y)$ exceeds it; their strictly weighted average can equal it
only if both do. This proves the face assertion.
\end{proof}

\begin{lemma}[Separation in the locally convex topology]\label{fa2:p1:c24:lcs-separation}
Continuous real linear functionals separate distinct points of $E$.
If $C$ is a nonempty closed convex subset and $x_0\notin C$, there is
such a functional with $\ell(x_0)>\sup_{x\in C}\ell(x)$.
\end{lemma}
\begin{proof}
A locally convex neighborhood basis is given by finite intersections
of balls of continuous seminorms. Since $C$ is closed, choose a continuous
seminorm $p$ (the maximum of finitely many rescaled defining seminorms)
and $\eps>0$ such that $x_0+\{p<\eps\}$ misses $C$.
In the real normed quotient $E/\ker p$, with norm $p$, the image of
$x_0$ has distance at least $\eps$ from the image of $C$ and thus from
its closure. That closure is convex. The normed point--closed-convex
separation theorem of Chapter~\ref{fa2:p1:c06} gives a bounded real
functional strictly separating the point from that closure. Composing
with the quotient gives $\ell$; it is continuous since
$|\ell(x)|\leq Cp(x)$ for some constant $C$.
For distinct points, Hausdorffness gives a continuous seminorm with
$p(x-y)>0$, and a norming Hahn--Banach functional on the same quotient
separates them. Thus the use of separation here has been reduced to
the proved normed-space theorem with a concrete continuity bridge.
\end{proof}

\begin{proposition}[Extreme points in every compact face]\label{fa2:p1:c24:face-extreme}
Every nonempty compact convex set has an extreme point. More precisely,
every nonempty compact face of it contains an extreme point of the
original set.
\end{proposition}
\begin{proof}
Fix a nonempty compact face $F_0$ (the whole set is also a face).
Order the nonempty compact faces contained in it by reverse inclusion.
Any chain has an upper bound given by its intersection. To see
nonemptiness, the chain has the finite intersection property inside
the compact set $F_0$; each member is closed because the ambient space
is Hausdorff. Compactness gives a nonempty intersection, and the face
operations prove it is again a compact face.
Zorn's lemma therefore supplies an inclusion-minimal nonempty compact
face $F$. If it contained distinct $x,y$, a continuous real functional
would distinguish them by the separation lemma. Its maximizing face
inside $F$ would be a proper nonempty compact face, and by transitivity
a face of $F_0$ and of the original set, contradicting minimality.
Hence $F$ is a singleton, whose point is extreme in the original set.
\end{proof}

% FAMILY D16.C24.02
\begin{theorem}[Krein--Milman]\label{fa2:p1:c24:krein-milman}\label{fa2:p1:krein-milman}
If $K$ is a nonempty compact convex subset of a Hausdorff locally convex
space, then
\[
 K=\cl\operatorname{conv}(\operatorname{Ext}K).
\]
The closure is in the given locally convex topology, and the convex
combinations are finite real convex combinations.
\end{theorem}
\begin{proof}
Extreme points exist by the preceding proposition. Let
$C=\cl\operatorname{conv}(\operatorname{Ext}K)$. Since compact subsets
are closed in the Hausdorff ambient space and $K$ is convex, $C\subset K$;
it is nonempty closed convex and compact. Suppose $x_0\in K\setminus C$.
The separation lemma gives a continuous real functional with
$\ell(x_0)>\sup_C\ell$. Its maximizing face in $K$ is nonempty compact
and contains an extreme point $e$ of $K$. Then
$\ell(e)=\max_K\ell\geq\ell(x_0)>\sup_C\ell$, whereas $e\in C$.
This contradiction proves the equality. No metrizability or selection
of a convergent sequence from a general compact set is used: the
nonempty intersections above are genuine compactness arguments.
\end{proof}

\begin{example}[The topology of the closed convex hull matters]\label{fa2:p1:c24:topology-example}
The closed unit ball of an infinite-dimensional Hilbert space is weakly
compact by Chapter~\ref{fa2:p1:c14}; its extreme points are its unit sphere.
Indeed a vector of norm less than one is the midpoint of two distinct
ball vectors, while equality in the norm estimate for a unit vector
written as a proper convex combination forces both endpoints to be that
unit vector, by the strict convexity following from the parallelogram
identity. Krein--Milman here asserts a weak closed convex hull equality.
It does not assert that this ball is norm compact; Chapter~\ref{fa2:p1:c01}
shows that it is not. For a general locally convex compact set it is the
specified topology, rather than sequences in an unspecified norm,
which enters the theorem.
\end{example}

% FAMILY D16.C24.03
\section{Schauder and compact maps}
We use the finite-dimensional Brouwer fixed-point theorem: a continuous self-map of a nonempty compact convex set in a finite-dimensional real normed space has a fixed point, including sets with empty ambient interior. This is the ordinary finite-dimensional topological theorem; the compact approximation and limiting arguments needed in infinite dimension are proved below.

\begin{lemma}[Finite convex approximation of a compact set]\label{fa2:p1:c24:finite-approximation}
Let $Q$ be a nonempty compact subset of a normed space and $\eps>0$.
There are points $q_1,\ldots,q_m\in Q$ and a continuous map
$P_\eps:Q\to\operatorname{conv}\{q_1,\ldots,q_m\}$ such that
$\norm{P_\eps(x)-x}<\eps$ for every $x\in Q$.
\end{lemma}
\begin{proof}
Choose finitely many open $\eps$-balls with centers $q_i$ covering $Q$.
Put $w_i(x)=\max(0,\eps-\norm{x-q_i})$ and
\[
 P_\eps(x)=\frac{\sum_iw_i(x)q_i}{\sum_iw_i(x)}.
\]
At least one weight is positive; all weights and the positive
denominator are continuous, proving continuity. The normalized weights
are nonnegative with sum one. Every positive-weight center has distance
strictly less than $\eps$ from $x$, and the triangle inequality for the
finite convex combination proves the strict error estimate.
\end{proof}

\begin{theorem}[Schauder on a compact convex set]\label{fa2:p1:c24:schauder}
A continuous self-map $T:C\to C$ of a nonempty compact convex subset
of a real or complex normed space has a fixed point.
\end{theorem}
\begin{proof}
For $\eps>0$ apply the preceding lemma to $Q=C$ and let
$D_\eps=\operatorname{conv}\{q_1,\ldots,q_m\}\subset C$.
This is compact because it is the continuous image of the compact
finite-dimensional probability simplex. It is convex and lies in the
finite-dimensional real span of its vertices, endowed with the inherited
norm. The map $P_\eps\circ T$ restricted to $D_\eps$ is a continuous
self-map of $D_\eps$. The exact Brouwer theorem just stated gives
$x_\eps=P_\eps(Tx_\eps)$, so
$\norm{x_\eps-Tx_\eps}<\eps$.
Take $\eps=1/n$. A norm-compact set is a compact metric space, hence
a subsequence of $x_{1/n}$ converges to some $x\in C$.
Continuity of $T$ and the residual estimate give
$\norm{x-Tx}=0$, so $Tx=x$. For a complex ambient space all convex
combinations and the finite-dimensional span here are real; its
underlying real normed topology is identical, matching Brouwer exactly.
No uniqueness conclusion follows from this construction.
\end{proof}

% FAMILY D16.C24.04
\section{Relatively compact maps and the convex-hull step}
\begin{lemma}[Compactness of the closed convex hull in a Banach space]\label{fa2:p1:c24:compact-hull}
If $Q$ is a compact subset of a Banach space, then
$\cl\operatorname{conv}Q$ is compact in norm.
\end{lemma}
\begin{proof}
The empty case is immediate. Given $\eps>0$, choose an $\eps$-net
$q_1,\ldots,q_m$ for $Q$. Replace every entry of a finite convex
combination of points of $Q$ by a corresponding $q_i$. The resulting
point lies in the compact finite polytope $D=\operatorname{conv}\{q_i\}$
and differs from the original point by less than $\eps$ in norm,
because its coefficients have sum one. A finite $\eps$-net for $D$
is therefore a $2\eps$-net for $\operatorname{conv}Q$, and a slightly
larger-radius net covers its closure. Thus that closure is totally
bounded. It is complete because it is closed in a Banach space.
The complete-totally-bounded criterion gives compactness.
This identifies the use of Banach completeness: boundedness of the
convex hull alone would not establish compactness.
\end{proof}

\begin{theorem}[Compact-map fixed-point theorem]\label{fa2:p1:c24:compact-map}
Let $D$ be a nonempty closed convex subset of a Banach space and let
$T:D\to D$ be continuous with relatively compact image $T(D)$.
Then $T$ has a fixed point. In particular this holds for a continuous
compact self-map of a nonempty closed bounded convex set.
The first assertion also holds in an arbitrary normed space if the
compact closure of $T(D)$ is taken in that space.
\end{theorem}
\begin{proof}
In the Banach case let $Q=\cl T(D)$ and
$C=\cl\operatorname{conv}Q$. Since $D$ is closed and convex,
$Q\subset D$ and $C\subset D$. The hull lemma proves $C$ compact;
it is nonempty and convex. Moreover $T(C)\subset T(D)\subset Q\subset C$.
Thus $T|_C$ satisfies every hypothesis of the compact-convex Schauder
theorem, which gives a fixed point in $D$. If $D$ is bounded and $T$
is compact in the sense of taking bounded subsets to relatively compact
sets, the required relative compactness of the full image follows.

For the asserted normed-space extension one need not assume a complete
closed convex hull. Let the compact set $Q=\cl T(D)\subset D$ have
the finite approximation $P_{1/n}$ above, and put
$D_n=\operatorname{conv}\{q_1^{(n)},\ldots,q_{m_n}^{(n)}\}\subset D$.
Brouwer applied to $P_{1/n}\circ T$ on $D_n$ gives $x_n\in D_n$
with $\norm{x_n-Tx_n}<1/n$. Compactness of $Q$ supplies a convergent
subsequence of $Tx_n$; its residual estimate makes the corresponding
$x_n$ converge to the same limit $x$. Closedness keeps $x$ in $D$,
and continuity gives $Tx=x$. This proves the source-supported
normed-space generality while retaining the complete invariant
compact-convex reduction in the Banach formulation.
\end{proof}

\begin{example}[Why the hull lemma states completeness]\label{fa2:p1:c24:incomplete-hull}
In $c_{00}$ with the $\ell^1$ norm, the set
$Q=\{0\}\cup\{e_n/n:n\geq1\}$ is compact, since $e_n/n\to0$.
The vectors $x_N=\sum_{n=1}^N2^{-n}e_n/n$ are finite convex
combinations of points of $Q$ (put the unused coefficient at zero),
and they are Cauchy in $\ell^1$. Their completion limit has infinitely
many nonzero coordinates, so no subsequence converges in $c_{00}$.
Consequently the closed convex hull in that incomplete space is not
compact. This does not contradict the normed compact-map theorem,
whose direct proof uses compactness of the image itself.
\end{example}

% FAMILY D16.C24.05
\section{Existence versus contraction conclusions}
\begin{proposition}[What contraction adds]\label{fa2:p1:c24:contraction-comparison}
If $D$ is a nonempty closed subset of a Banach space and
$\norm{Tx-Ty}\leq q\norm{x-y}$ with $0\leq q<1$, then $T:D\to D$
has a unique fixed point and every iterate converges to it, with
\[
 \norm{T^nx-x_*}\leq\frac{q^n}{1-q}\norm{Tx-x}.
\]
Convexity and relative compactness of the image are not needed for this
assertion. Schauder supplies existence under its own hypotheses and
does not supply this uniqueness or iteration conclusion.
\end{proposition}
\begin{proof}
For completeness, successive distances satisfy
$\norm{T^{n+1}x-T^nx}\leq q^n\norm{Tx-x}$.
Summing the geometric tail shows that the iterates are Cauchy in the
closed complete set $D$, so they converge to $x_*\in D$ and give the
displayed estimate by letting the upper tail endpoint tend to infinity.
The Lipschitz bound makes $T$ continuous, hence $Tx_*=x_*$.
Two fixed points have distance at most $q$ times that same distance,
forcing equality. The geometric-series proof also includes $q=0$.
\end{proof}

\begin{example}[Schauder need not select a unique or attracting point]\label{fa2:p1:c24:comparison-example}
On the compact convex interval $[0,1]$, the identity map has every point
fixed and satisfies Schauder's hypotheses. The map $T(x)=1-x$ also
satisfies them and has the unique fixed point $1/2$, but its iterates
oscillate between $x$ and $1-x$ unless $x=1/2$. Its Lipschitz constant
is one, so the strict contraction hypothesis is absent. In contrast
$T(x)=x/2$ is a contraction with the unique fixed point zero and all
iterates converging geometrically. These examples separate existence,
uniqueness and convergence without identifying any of them implicitly.
\end{example}

% END INPUT


% BEGIN parts/part1-convex-advanced.tex
\chapter{Convex duality, representing measures and fixed points}
\label{fa2:p1:convex-advanced}
The compact convex theorems do more than provide existence of extreme
points. They give precise representation and invariance statements, whose
different topological hypotheses must be retained. We use
\cite[Chapters 3--5]{B04}, \cite[Chapter V]{B02} and
\cite[Chapters 3--4,13--14]{B08}.

\section{A qualification theorem in convex duality}
\begin{theorem}[Fenchel--Moreau in a locally convex space]
\label{fa2:p1:fenchel-moreau-lcs}
If $\varphi:E\to(-\infty,\infty]$ is proper, convex and lower
semicontinuous on a real Hausdorff LCS, then
\[
 \varphi(x)=\sup_{f\in E'}\{f(x)-\varphi^*(f)\},\qquad
 \varphi^*(f)=\sup_x\{f(x)-\varphi(x)\}.
\]
\end{theorem}
\begin{proof}
Fenchel's inequality gives the upper bound for the displayed supremum.
We show that continuous affine minorants separate every point below
the epigraph. If $r<\varphi(x_0)$, strictly separate $(x_0,r)$
from the closed convex epigraph by a functional of the form
$(x,t)\mapsto f(x)+st$, oriented so that
$f(x)+st\geq\alpha>f(x_0)+sr$ on the epigraph.
Its upward rays imply $s\geq0$. When $s>0$, the affine function
$(\alpha-f(x))/s$ is a minorant taking value greater than $r$
at $x_0$. At a point $\bar x$ where $\varphi$ is finite,
separation of $(\bar x,\varphi(\bar x)-1)$ necessarily has
$s>0$; hence at least one continuous affine minorant $m$ exists.
In the remaining case $s=0$, we have $f(x)\geq\alpha>f(x_0)$
on the effective domain. Each
$m_n(x)=m(x)+n(\alpha-f(x))$ is again a minorant, and
$m_n(x_0)\to\infty$. Thus a minorant exceeding $r$ always exists.
Writing a minorant as $h(x)-c$ gives
$c\geq\varphi^*(h)$, so it is bounded above by a member of the
displayed supremum. Let $r\uparrow\varphi(x_0)$, also when the
latter is infinite.
\end{proof}

\begin{theorem}[Fenchel--Rockafellar]
\label{fa2:p1:fenchel-rockafellar}
Let $\varphi,\psi:E\to(-\infty,\infty]$ be proper convex functions
on a real Hausdorff LCS. Suppose $x_0$ is in both effective domains
and $\varphi$ is continuous at $x_0$. Then
\[
 \inf_x(\varphi(x)+\psi(x))
 =\sup_{f\in E'}[-\varphi^*(f)-\psi^*(-f)].
\]
If the infimum is finite, the dual supremum is attained. No primal
attainment is asserted.
\end{theorem}
\begin{proof}
Put $a=\inf(\varphi+\psi)$ and
$p(u)=\inf_x[\varphi(x)+\psi(x-u)]$. This function is convex,
$p(0)=a$, and is bounded above in a neighborhood of zero because
$p(u)\leq\varphi(x_0+u)+\psi(x_0)$. If $a=-\infty$, weak
duality already forces the dual supremum to be $-\infty$.
Assume $a$ finite. Convexity and the local upper bound prevent
$p$ from being $-\infty$ anywhere: a point with that value,
combined with a suitable point near zero on the opposite ray,
would force $p(0)=-\infty$. On a balanced convex neighborhood
where $p\leq M$, the inequality
$2p(0)\leq p(u)+p(-u)$ gives a lower bound. Scaling within
this neighborhood proves continuity at zero: for $u\in tU$,
convexity gives $|p(u)-p(0)|\leq t(M-p(0))$, after replacing
$M$ by a larger bound if necessary.

The directional derivative
$q(v)=\lim_{t\downarrow0}(p(tv)-a)/t$ is finite and sublinear.
Finiteness and continuity bounds follow from the same neighborhood
estimate. Subadditivity follows by writing
$p(t(u+v))\leq[p(2tu)+p(2tv)]/2$ and taking the limit.
Hahn--Banach gives a continuous linear $f\leq q$; its continuity
follows from the local bound on $q$. Convexity gives
$q(u)\leq p(u)-a$, hence $p(u)\geq a+f(u)$.
For independent $x,y$ use $u=x-y$ to obtain
\[
 \varphi(x)-f(x)+\psi(y)+f(y)\geq a.
\]
Taking the two independent infima gives
$-\varphi^*(f)-\psi^*(-f)\geq a$. Weak duality gives the reverse
inequality for every $f$, proving equality and attainment.
\end{proof}

\section{Invariant extension and finitely additive representations}
\begin{theorem}[Agnew--Morse invariant Hahn--Banach]
\label{fa2:p1:invariant-hahn-banach}
Let $p$ be a finite real sublinear function on a real vector space,
and let $\mathcal S$ be a commuting semigroup of linear maps with
$p(Sx)\leq p(x)$. Let $Y$ be invariant under these maps and
$f:Y\to\R$ linear, $f\leq p$, $f(Sy)=f(y)$. There is an
extension $F\leq p$ to the whole space with $F(Sx)=F(x)$.
\end{theorem}
\begin{proof}
Adjoin the identity and let $\mathcal C$ be the convex hull of the
semigroup. Products of members of $\mathcal C$ remain there and
commute. Set $q(x)=\inf_{C\in\mathcal C}p(Cx)$.
Since $p(Cx)\leq p(x)$ and $p(Cx)\geq-p(-Cx)\geq-p(-x)$,
$q$ is finite. It is positively homogeneous. If $C$ nearly
minimizes at $x$ and $D$ at $y$, then
\[
 q(x+y)\leq p(CD(x+y))\leq p(Cx)+p(Dy),
\]
using commutation and subinvariance. Thus $q$ is sublinear.
For $y\in Y$, $f(y)=f(Cy)\leq p(Cy)$, so $f\leq q$ on $Y$.
Extend by ordinary Hahn--Banach under $q$.
For $S\in\mathcal S$, the average
$C_n=n^{-1}\sum_{j=0}^{n-1}S^j$ lies in $\mathcal C$, and
\[
 q((I-S)x)\leq p(C_n(I-S)x)
 \leq\frac{p(x)+p(-x)}n.
\]
Hence $q((I-S)x)\leq0$, and applying the same argument to $-x$
makes $F((I-S)x)=0$. Also $F\leq q\leq p$, as required.
\end{proof}

\begin{corollary}[Banach limits]\label{fa2:p1:banach-limits}
There is a positive norm-one real linear functional $L$ on
$\ell^\infty$ extending the usual sequence limit and invariant under
the left shift. It satisfies $\liminf x_n\leq L(x)\leq\limsup x_n$.
It can be chosen to agree with every convergent Ces\`aro mean.
It has no representation by an $\ell^1$ sequence.
\end{corollary}
\begin{proof}
Apply the invariant extension theorem to $p(x)=\limsup x_n$,
the left shift and the usual limit on convergent sequences.
Sublinearity and invariance of $p$ follow from the scalar limit
rules. Applying the bound to $x$ and $-x$ gives both inequalities,
positivity, $L(1)=1$ and $\|L\|=1$. For the additional choice
start instead on the linear space of sequences whose Ces\`aro
means converge. Their Ces\`aro limit is linear and invariant
under the shift, and lies between ordinary liminf and limsup;
the same extension therefore works. Since $L(e_n)=0$ for every
unit vector, any $\ell^1$ representation would have all
coefficients zero, contrary to $L(1)=1$.
\end{proof}

\begin{corollary}[Banach integration and extension of Lebesgue measure]
\label{fa2:p1:banach-measure}
The Lebesgue integral on bounded measurable functions on $\mathbb R/\mathbb Z$
extends to a positive translation-invariant norm-one functional on all
bounded functions on that circle.  There is an extended nonnegative
finitely additive function $\widehat m$ on every subset of $\mathbb R$,
translation invariant and agreeing with Lebesgue measure on every
Lebesgue measurable set.  Countable additivity cannot be required.
\end{corollary}
\begin{proof}
On the real vector space of all bounded circle functions take
$p(f)=\sup f$.  Circle translations commute and preserve $p$.
The subspace of bounded Lebesgue measurable functions is invariant,
and its ordinary integral $I_0$ is invariant and bounded above by
$p$.  The invariant Hahn--Banach theorem extends it to $I\le p$.
Applying this bound to $-f$ shows positivity, and $I(1)=1$ gives
its norm.  Complexify if desired.  Thus
$b_0(A)=I(1_A)$ is a finitely additive circle probability agreeing
with Lebesgue measure on measurable sets.

For a bounded subset $A$ of the real line, cut it into its finitely
many nonempty pieces $A\cap[n,n+1)$, project each piece to the
circle, and let $b(A)$ be the sum of their $b_0$ values.  This is
nonnegative and finitely additive.  It is translation invariant:
refine the cut simultaneously by the integer grid and its translate
by the proposed translation.  Each of the finitely many resulting
pieces lies in one unit interval before and after translation;
its two circle images differ by that translation, so have the same
$b_0$ value.  Summing gives the claim.  For arbitrary $A$ put
\[
 \widehat m(A)=\lim_{N\to\infty}b(A\cap[-N,N)).
\]
Monotonicity gives this limit in $[0,\infty]$.  For disjoint $A,B$,
finite additivity holds at each $N$ and passes to these nonnegative
limits.  Translating the exhaustion changes it to a cofinal bounded
interval exhaustion; monotonicity and translation invariance of
$b$ therefore prove translation invariance of $\widehat m$.
For measurable $A$, every bounded piece has its ordinary measure,
and monotone convergence proves $\widehat m(A)=m(A)$.

Finally choose one representative in $[0,1]$ from each equivalence
class modulo rational difference.  Its distinct rational translates
by rationals in $[-1,1]$ are disjoint, cover $[0,1]$, and lie in
$[-1,2]$.  Translation invariance and countable additivity would
give their representative set either positive mass, contradicting
the finite mass of $[-1,2]$, or zero mass, contradicting the mass
one of $[0,1]$.  This gives the precise assertions of
\cite[Chapter 1, Theorems 5.6--5.7]{B09}.
\end{proof}

\begin{proposition}\label{fa2:p1:ba-representation}
The continuous dual of real $\ell^\infty(S)$ is isometrically the
space of finitely additive signed set functions of bounded variation
on the power set of $S$. The correspondence is
$\mu(A)=L(1_A)$, and integration is first defined on finite simple
functions and then by uniform approximation. Positive functionals
correspond to positive finitely additive measures. The invariant
extension theorem gives a translation-invariant finitely additive
probability measure on the full power set of every abelian group.
\end{proposition}
\begin{proof}
Finite additivity follows from $1_{A\cup B}=1_A+1_B$ for disjoint
sets. For a finite partition $(A_j)$ choose signs so that
$\sum|\mu(A_j)|=L(\sum\eps_j1_{A_j})\leq\|L\|$.
Thus the variation is bounded. Conversely the variation bound
gives $|\int s\,d\mu|\leq\|\mu\|_{\mathrm{var}}\|s\|_\infty$
on simple functions. Every bounded function is uniformly
approximable by finite simple functions by partitioning its bounded
range into small intervals. The integral extends uniquely and
has exactly the variation norm, as the same signed partition
functions attain its supremum. The two constructions are inverse.
For a group use $p=\|\cdot\|_\infty$ and the constant functional
$c1\mapsto c$; the commuting translations preserve $p$.
The extension has norm one and value one at $1$, hence is positive:
for $0\leq h\leq1$, $L(1-h)\leq1$ implies $L(h)\geq0$.
Set $\mu(A)=L(1_A)$. Its invariance and normalization follow.
Countable additivity is not implied; for a Banach limit on $\N$
all singleton masses vanish although the whole space has mass one.
\end{proof}

\section{Milman's converse and Choquet representation}
\begin{theorem}[Milman's converse and extreme-point lifting]
\label{fa2:p1:milman-converse}
If $K=\overline{\operatorname{conv}}A$ is compact in a Hausdorff LCS,
then $\operatorname{Ext}K\subset\overline A$. If $T:K\to L$ is a
continuous affine surjection between compact convex sets, every
extreme point of $L$ is the image of an extreme point of $K$.
\end{theorem}
\begin{proof}
Replace $A$ by its compact closure in $K$. For a convex balanced
zero neighborhood $U$, finitely many $a_j+U$ cover $A$. Let
$K_j$ be the closed convex hull in $K$ of
$A\cap(a_j+U)$. Each is compact and contained in $a_j+\overline U$.
The convex hull of finitely many compact convex sets is compact:
it is the image of their product times a finite-dimensional simplex.
It contains $A$, so equals $K$. An extreme point of this convex
hull must belong to one of the $K_j$, and therefore to
$A+\overline U$. Intersecting over shrinking $U$ gives membership
in the closed $A$. For the second statement the fiber over an
extreme point is a nonempty compact face. Krein--Milman supplies
an extreme point of the fiber, which is extreme in $K$ by the
face property.
\end{proof}

\begin{theorem}[Two barycentric representation statements]
\label{fa2:p1:choquet}
Every point of a nonempty compact convex subset $K$ of a Hausdorff
LCS is the barycenter of a regular probability measure on
$\overline{\operatorname{Ext}K}$. If $K$ is metrizable, the
measure can be chosen with mass one on $\operatorname{Ext}K$
itself. The latter set need not be closed.
\end{theorem}
\begin{proof}
The first assertion follows from Krein--Milman and
Theorem~\ref{fa2:p1:barycenters}, applied to the compact closure
of the extreme points. We prove the stronger metrizable assertion.
Choose a sequence of bounded continuous real affine functions
$a_j$ separating points of $K$, with $|a_j|\leq1$; such a
sequence exists because the off-diagonal in the compact metric
$K\times K$ is Lindel\"of, and the separating dual supplies an
open cover of it. Then $q(x)=\sum2^{-j}a_j(x)^2$ is continuous
and strictly convex. For fixed $x$, the representing probability
measures form a nonempty weak-star compact set (the barycentric
conditions are closed affine scalar tests). Choose a measure
$\mu$ maximizing $\int q\,d\mu$ there.

Fix a compatible metric $d$ on $K$. Let $F_n$ be the compact
set of triples $(z,y,w)$ with $z=(y+w)/2$ and $d(y,w)\geq1/n$,
and let $A_n$ be its projection to $z$. These are compact and
$K\setminus\operatorname{Ext}K=\bigcup_n A_n$. The latter
identity follows because a nontrivial chord through a nonextreme
point can be shortened on its longer side to make that point
its midpoint. Thus the extreme set is Borel.

We need a Borel choice of a pair $(y(z),w(z))$ in each nonempty
fiber over $A_n$. Here is the elementary compact-metric selection
argument. For a compact closed graph $F\subset A_n\times K^2$,
the set of fibers meeting an open set is Borel: express the open
set as a countable union of closed balls contained in it, and
project the resulting compact intersections. Choose successively
the first closed ball of radius $2^{-j}$, from a countable
dense family, whose intersection with the previously chosen
balls and the fiber is nonempty. The choices are Borel by that
projection observation. The nested compact intersections have
diameter tending to zero; their unique limit is the required
Borel selection.

On the compact $F_n$ strict convexity gives a positive minimum
$c_n$ for $[q(y)+q(w)]/2-q(z)$. If $\mu(A_n)>0$, replace the
restriction of $\mu$ to $A_n$ by half its image under each of
the selected endpoint maps. The resulting probability measure
has the same barycenter, because the maps have midpoint $z$,
and its integral of $q$ increases by at least $c_n\mu(A_n)$.
This contradicts maximality. Thus every $A_n$ has measure zero,
and $\mu(\operatorname{Ext}K)=1$, as asserted.
\end{proof}

\begin{theorem}[Birkhoff--von Neumann, finite matrices]
\label{fa2:p1:birkhoff-von-neumann}
Every $n\times n$ doubly stochastic matrix is a convex combination
of permutation matrices. These are exactly the extreme points of
the compact convex set of doubly stochastic matrices.
\end{theorem}
\begin{proof}
The set is a closed bounded convex subset of $\R^{n^2}$. A
permutation matrix is extreme: any zero entry forces the same
zero in both matrices of a proper convex decomposition, and
the row sums then force the remaining entries to be one.
For the converse draw the bipartite graph of positive entries,
with rows and columns as vertices. An edge incident to a vertex
of degree one has weight one, by the corresponding row or column
sum. Its other endpoint also has degree one. Remove every such
isolated pair. If anything remains, every remaining vertex has
degree at least two, so a finite walk continuing without
backtracking contains a cycle. The graph is bipartite, hence
the cycle is even. Add and subtract a sufficiently small positive
$\eps$ on alternate edges of the cycle. Both resulting matrices
are nonnegative and have unchanged row and column sums, are
distinct, and average to the original. Thus a nonpermutation
matrix is not extreme. Krein--Milman in this finite-dimensional
compact convex set now gives its closed convex hull of permutation
matrices. There are finitely many permutations; their convex
hull is already compact and closed. This proves the asserted
finite convex decomposition, including $n=0$ by the empty-matrix
convention if that case is admitted.
\end{proof}

\section{A source endpoint that changes the representation}
\begin{theorem}[Bernstein representation with the endpoint explicit]
\label{fa2:p1:bernstein}
Let $f:[0,\infty)\to\R$ be finite and satisfy
$(-1)^m\Delta_{h_1}\cdots\Delta_{h_m}f(t)\geq0$ for
$t\geq0$, $h_j>0$, $m\geq0$. Then there is a unique finite
positive measure $\nu$ on the compactified interval $[0,\infty]$
with
\[
 f(t)=\int e^{-t\lambda}\,d\nu(\lambda),
\]
where at $\lambda=\infty$ the integrand is zero for $t>0$ and
one for $t=0$. The mass at infinity is
$f(0)-\lim_{t\downarrow0}f(t)$. Thus the classical representation
by a finite measure on $[0,\infty)$ holds at every $t\geq0$
exactly when $f$ is right-continuous at zero.
\end{theorem}
\begin{proof}
If $f(0)=0$, monotonicity and positivity make $f=0$. Otherwise
normalize $f(0)=1$. The set $K$ of all normalized functions
satisfying the inequalities is a closed convex subset of
$[0,1]^{[0,\infty)}$, hence compact in the product LCS topology.
Its functions are decreasing and convex; midpoint convexity
with the local boundedness $0\leq f\leq1$ gives ordinary convexity
and continuity on $(0,\infty)$.

If an extreme $e\in K$ has $0<e(a)<1$, both
\[
 g(t)=\frac{e(t)-e(t+a)}{1-e(a)},\qquad
 h(t)=\frac{e(t+a)}{e(a)}
\]
belong to $K$, by the finite-difference inequalities, and
$e=(1-e(a))g+e(a)h$. Extremality forces
$e(t+a)=e(t)e(a)$. Such a function cannot vanish at a positive
point: its values $e(na)=e(a)^n$ remain positive and monotonicity
would then be contradicted. It cannot equal one at a positive
point unless it is identically one, since a convex decreasing
function constant on an initial interval must remain constant.
Thus, unless it is constant one, all its positive values lie
strictly between zero and one, and the same argument for every
$a>0$ gives the multiplicative equation. Continuity gives
$e(t)=e^{-\lambda t}$ with $\lambda\geq0$.
If there is no positive point with value strictly between zero
and one, continuity on the connected interval $(0,\infty)$
makes $e$ either one there or zero there. The latter is the
function $e_\infty$ having value one only at zero.

The set $S=\{e_\lambda:0\leq\lambda\leq\infty\}$ is compact
in the product topology and contains all extreme points. Hence
$K=\overline{\operatorname{conv}}S$ by Krein--Milman. The
barycenter theorem produces the desired measure on $S$, identified
homeomorphically with $[0,\infty]$. For uniqueness set
$r=e^{-\lambda}\in[0,1]$. The values $f(n)$, $n=0,1,\ldots$,
are all moments of the pushed-forward measure. Polynomials are
dense in $C[0,1]$, so equality of all moments gives equality of
the measures. This also proves that each $e_\lambda$ is extreme.
Dominated convergence as $t\downarrow0$ shows that the right
limit is the mass of the finite-parameter part, proving the
endpoint assertion. Multiplication back by $f(0)$ completes the
general case.
\end{proof}

\begin{remark}\label{fa2:p1:bernstein-source-endpoint}
The finite-difference definition and the displayed extreme function
$e_\infty$ both occur in \cite[Section 14.3]{B08}. If its subsequent
integral notation is read as a measure only on finite parameters
and as an identity also at $t=0$, this extreme function is a
counterexample. The preceding theorem states the compactification
and continuity alternatives explicitly; it does not suppress
the endpoint or assume that a printed citation supplies it.
For $t>0$ differentiation under the integral is justified by
boundedness of $\lambda^me^{-t\lambda}$ away from zero in $t$,
so the represented function is smooth there with alternating
derivative signs.
\end{remark}

\begin{corollary}[The open-half-line Bernstein theorem]
\label{fa2:p1:bernstein-open-half-line}
For a finite function $f:(0,\infty)\to[0,\infty)$ the following are
equivalent: all alternating finite differences are nonnegative;
$f$ is smooth and $(-1)^mf^{(m)}\ge0$ for all $m$; and
\[
 f(t)=\int_{[0,\infty)}e^{-t\lambda}\,d\nu(\lambda)\quad(t>0)
\]
for a unique positive measure $\nu$, with the displayed integrals
finite. The measure may have infinite total mass but is finite on
bounded intervals. It is finite exactly when $f(0+)<\infty$, in
which case its mass is $f(0+)$.
\end{corollary}
\begin{proof}
The finite-difference assumptions imply monotonicity, midpoint
convexity and local boundedness, hence continuity in the interior.
For each $a>0$, apply the preceding theorem to $f_a(t)=f(t+a)$
on $[0,\infty)$. Its continuity at zero excludes an atom at infinity,
so it has a finite representing measure $\nu_a$ on finite parameters.
If $b>a$, uniqueness applied to $f(t+b)$ gives
$\nu_b=e^{-(b-a)\lambda}\nu_a$.
Put $\nu=e^\lambda\nu_1$. For every $a<1$, the preceding identity
with $b=1$ gives $\nu_a=e^{-a\lambda}\nu$; for $a\ge1$ it gives
the same formula directly. Consequently
$f(t+a)=\int e^{-(t+a)\lambda}\,d\nu$ for $t\ge0$.
Choosing $a<t$ proves the desired formula at every positive $t$.
The measure is locally finite because $e^\lambda$ is bounded on
bounded intervals. Any other representing measure has the same
finite weighted measure $e^{-a\lambda}\nu$ for each $a>0$, by the
finite theorem's uniqueness, so it agrees with $\nu$.
Differentiation under the integral is valid near any $t>0$ because
$\lambda^me^{-t\lambda}\le C_{m,t}e^{-t\lambda/2}$.
This proves the derivative signs. Conversely those signs imply
the finite-difference signs by integrating $f^{(m)}$ over the
rectangle $[0,h_1]\times\cdots\times[0,h_m]$.
Finally monotone convergence as $t\downarrow0$ gives
$f(0+)=\nu([0,\infty))$, including the infinite case.
\end{proof}

\section{Fixed points in locally convex spaces}
\begin{theorem}[Schauder--Tychonoff]\label{fa2:p1:schauder}
A continuous self-map of a nonempty compact convex subset of a
Hausdorff LCS has a fixed point.
\end{theorem}
\begin{proof}
For an open convex balanced zero neighborhood $U$, choose a finite
cover of $K$ by $(x_j+U)\cap K$. Compact Hausdorff spaces admit
a continuous partition of unity subordinate to a finite open
cover: shrink the cover using normality, choose finitely many
Urysohn functions with positive total, and divide by their sum.
Let $\phi_j$ be such a partition and
$P_Ux=\sum\phi_j(x)x_j$. Then $P_Ux-x\in U$ by convexity.
On the finite-dimensional compact convex hull of the $x_j$,
the map $P_U\circ T$ has a fixed point $x_U$ by Brouwer.
Thus $x_U-Tx_U\in U$. Direct the zero neighborhoods by reverse
inclusion. Compactness gives a convergent subnet of $x_U$,
with limit $x\in K$; the differences tend to zero and continuity
gives $x=Tx$.
\end{proof}

\begin{corollary}[Compact images and compact-homotopy continuation]
\label{fa2:p1:leray-schauder}
A continuous self-map $T$ of a nonempty closed convex subset $C$
of a Banach space has a fixed point if $\overline{T(C)}$ is compact.
More generally let $F:E\times[0,1]\to E$ be continuous, mapping
bounded sets to relatively compact sets, where $E$ is Banach.
Suppose $F(x,0)=0$ for every $x$, and that some $R>0$ satisfies
\[
 x=F(x,t),\quad 0\le t\le1\quad\Longrightarrow\quad\|x\|<R.
\]
Then $F(\cdot,1)$ has a fixed point.  In particular, a compact
continuous $T:E\to E$ has a fixed point whenever
$\{x:x=tT(x)\text{ for some }0\le t\le1\}$ is bounded.
\end{corollary}
\begin{proof}
For the first assertion the closed convex hull $K$ of
$\overline{T(C)}$ is compact by the complete-space compact-hull
theorem.  It lies in $C$, and $T(K)\subset K$; apply Schauder
to $K$.

For continuation choose $M>R$ and let $P_M:E\to\overline B_M$
be radial retraction, $P_My=y\min(1,M/\|y\|)$, with its value
at zero defined separately as zero.  On the nonempty closed
bounded convex subset $\overline B_M\times[0,1]$ of $E\times\mathbb R$
define
\[
 H(x,t)=\left(P_MF(x,t),
       \max\{0,\min\{1,t+R-\|x\|\}\}\right).
\]
This is continuous with relatively compact range: the first
coordinate comes from the continuous image of the compact closure
of $F(\overline B_M\times[0,1])$, and the second lies in a compact
interval.  The first assertion gives a fixed point $(x,t)$.
If $0<t<1$, the scalar equation forces $\|x\|=R<M$.
Since radial retraction can produce a vector of norm strictly
less than $M$ only without truncation, the other equation gives
$x=F(x,t)$, contradicting the strict bound $R$.  If $t=0$,
the first equation gives $x=0$, while the scalar equation would
require $R-\|x\|\le0$, again a contradiction.  Thus $t=1$.
The scalar equation now gives $\|x\|\le R<M$, so no truncation
occurs and $x=F(x,1)$, as required.  For the last assertion set
$F(x,t)=tT(x)$ and choose $R$ larger than the given bound.
This proves the compact-homotopy form in
\cite[Appendix A, Theorems A.3.6--A.3.7]{B07} without assuming
that $F$ itself maps a chosen ball into that ball.
\end{proof}

\begin{theorem}[Markov--Kakutani]\label{fa2:p1:markov-kakutani}
A commuting family of continuous affine self-maps of a nonempty
compact convex subset of an LCS has a common fixed point.
\end{theorem}
\begin{proof}
For one map $T$, the Ces\`aro means of any orbit lie in $K$ and
satisfy $Tz_n-z_n=(T^nx-x)/n\to0$, because compact sets are
bounded. A cluster point is fixed. The fixed set is compact
convex; commutation makes it invariant under every other map.
Induction gives common fixed points for every finite subfamily,
and compactness applied to these closed fixed sets gives a
point fixed by the whole family.
\end{proof}

\begin{lemma}[A small cap of a separable weakly compact set]
\label{fa2:p1:small-cap}
If $K$ is nonempty separable weakly compact convex in an LCS and
$p$ is a continuous seminorm, then for every $\eps>0$ there
is a proper weakly closed convex $C\subset K$ with
$p$-diameter of $K\setminus C$ less than $\eps$.
\end{lemma}
\begin{proof}
Let $D$ be the weak closure of $\operatorname{Ext}K$ and choose
a countable dense subset $A$ of $K$ in the original topology.
The weakly closed sets $a+\{p\leq\eps/8\}$, $a\in A$, cover
$D$. Its compact Hausdorff topology is Baire. Therefore some
such set contains a nonempty relatively weakly open $W\cap D$.
Let $K_1=\overline{\operatorname{conv}}^{w}(D\setminus W)$
and $K_2=\overline{\operatorname{conv}}^{w}(D\cap W)$.
If $K_1=\varnothing$, then $K=K_2$ has $p$-diameter at most
$\eps/4$, and $C=\varnothing$ proves the assertion. Hence assume
$K_1\ne\varnothing$ in the remainder of the proof.
Both are compact convex subsets of $K$, and $K_2$ has
$p$-diameter at most $\eps/4$. The compact convex hull of their
union equals $K$. Milman's converse gives $K_1\ne K$:
otherwise all extreme points, and then their weak closure $D$,
would lie in the closed $D\setminus W$, contradicting the
nonempty $W\cap D$.

For $0<r\leq1$ the set
$C_r=\{tx_1+(1-t)x_2:x_i\in K_i,\ r\leq t\leq1\}$ is
compact and convex. It cannot equal $K$, since an extreme point
in such a representation belongs to $K_1$, which would force
$K=K_1$. Every point outside $C_r$ has a representation with
$0\leq t<r$, and is therefore within $rd$ in seminorm of
$K_2$, where $d=p$-diameter of $K$ is finite. Finiteness follows
from weak boundedness and the normed quotient defined by $p$.
Thus $p$-diameter$(K\setminus C_r)\leq2rd+\eps/4$.
Choose $r$ small enough; the case $d=0$ is immediate by taking
any proper closed convex subset supplied by separation, or
$C=\varnothing$ when $K$ is a singleton.
\end{proof}

\begin{theorem}[Ryll--Nardzewski]\label{fa2:p1:ryll-nardzewski}
Let $K$ be nonempty weakly compact convex in an LCS. Let
$\mathcal S$ be a semigroup of weakly continuous affine self-maps
which is noncontracting in the original topology: for $x\ne y$,
$0\notin\overline{\{Sx-Sy:S\in\mathcal S\}}$.
Then the semigroup has a common fixed point.
\end{theorem}
\begin{proof}
For finitely many $T_1,\ldots,T_n$, their average $T_0$ has a
fixed point $x_0$ by the single-map affine fixed-point argument.
Discard any $T_j$ already fixing $x_0$; the average of the
remaining maps still fixes $x_0$. If any remain, noncontraction
gives one continuous seminorm $p$ and $\eta>0$ such that
$p(S(T_jx_0)-Sx_0)>\eta$ for all $S\in\mathcal S$ and all
remaining $j$. A maximum of finitely many seminorms supplies
one common $p$. The semigroup generated by these finitely many
maps is countable. Its orbit has a weakly closed convex hull
$Q\subset K$, weakly compact and separable in the original
topology: convex closures for the original and weak topologies
agree, so rational convex combinations of the countable orbit
are dense. Apply the cap lemma to get a proper closed convex
$C\subset Q$ with $p$-diameter of $Q\setminus C$ less than
$\eta$. Some orbit point $Sx_0$ lies outside $C$, since otherwise
its closed convex hull would be inside $C$. But
$Sx_0=n^{-1}\sum ST_jx_0$. Convexity forces at least one
$ST_jx_0$ outside $C$ as well, contradicting their separation
in seminorm. Thus every finite subfamily has a common fixed
point. Their fixed sets are weakly closed in the compact $K$
and have the finite-intersection property, proving the theorem.
\end{proof}

% END INPUT


% BEGIN parts/part1-supplements.tex
\chapter{Additional structural constructions and examples}
\label{fa2:p1:supplements}

These constructions complete variants from \cite{B01,B02,B03,B04,B05,B07,B08,B09}:
closed-sum estimates are more precise than algebraic splitting, a
separable extension theorem contrasts with the noncomplemented $c_0$
example, and the function-space examples retain their endpoint hypotheses.

\section{Sums, distances, and ordered extension}

\begin{proposition}[Classical differentiable and H\"older Banach spaces]
\label{fa2:p1:holder-banach}
On a compact interval $I=[a,b]$ with $a<b$, the space $C^m(I)$ is complete in
$\sum_{j=0}^m\|f^{(j)}\|_\infty$.  For $0<\alpha\le1$,
$C^{m,\alpha}(I)$ is complete in that norm plus
\[
 [f^{(m)}]_\alpha
 =\sup_{s\ne t\in I}\frac{|f^{(m)}(s)-f^{(m)}(t)|}{|s-t|^\alpha}.
\]
The same assertion for $m=0$ holds on an arbitrary metric space
when the uniform norm is required finite.  Pointwise multiplication
is continuous; the product of an $\alpha$-H\"older and a
$\beta$-H\"older function on $I$ is $\min(\alpha,\beta)$-H\"older.
If $m+\alpha<n+\beta$, with integer $m,n\ge0$ and
$0<\alpha,\beta\le1$, then
$C^{n,\beta}(I)\subset C^{m,\alpha}(I)$ continuously.
\end{proposition}
\begin{proof}
For a Cauchy sequence in $C^m$, every derivative converges uniformly
to a continuous function $g_j$.  Passing to the limit in
$f_k^{(j)}(t)-f_k^{(j)}(a)=\int_a^tf_k^{(j+1)}(s)ds$
proves $g_j'=g_{j+1}$, so $g_0\in C^m$ and convergence holds in
the stated norm.  If the sequence is Cauchy in the H\"older norm,
then for sufficiently large $k,l$ all its difference quotients
satisfy
$|(f_k^{(m)}-f_l^{(m)})(s)-(f_k^{(m)}-f_l^{(m)})(t)|
\le\varepsilon|s-t|^\alpha$.  Let $l\to\infty$ first, then
take the supremum in $s,t$; this proves convergence of the seminorm
as well.  For $m=0$ the same uniform-limit argument works on any
metric space.  Empty suprema for this metric-space version are
zero here.

The estimate
$[uv]_\alpha\le\|u\|_\infty[v]_\alpha+
\|v\|_\infty[u]_\alpha$ and the Leibniz formula give continuous
multiplication.  On a bounded interval an exponent can be lowered
using $|s-t|^{\beta}\le |I|^{\beta-\alpha}|s-t|^\alpha$
when $\alpha\le\beta$; bounded first derivatives also give a
Lipschitz bound by integration.  If $m=n$, the asserted inclusion
uses the former estimate; if $m<n$, it uses the latter for the
$m$th derivative and then lowers the exponent if necessary.
This proves all the claims in the exact interval and bounded-function
settings used in \cite[Appendix A.1]{B07} and
\cite[Chapter 1, Section 3.1]{B09}.
\end{proof}

\begin{theorem}[Metric completion]
\label{fa2:p1:metric-completion}
Every metric space $(X,d)$ has an isometric embedding $j$ into a
complete metric space $\widehat X$ with dense image.  Any two such
completions have a unique surjective isometry intertwining their
embeddings.  Every uniformly continuous map from $X$ to a complete
metric space extends uniquely to a uniformly continuous map on
$\widehat X$; its Lipschitz constant, if prescribed, is preserved.
Thus the linear completion above is also the metric completion;
the underlying metric assertion does not require a vector space.
\end{theorem}
\begin{proof}
Take all Cauchy sequences in $X$ and identify $(x_n)$ with $(y_n)$
when $d(x_n,y_n)\to0$.  The triangle inequality shows this is an
equivalence relation.  The real sequence $d(x_n,y_n)$ is Cauchy,
since its $m,n$ difference is at most
$d(x_m,x_n)+d(y_m,y_n)$.  Its limit is independent of the chosen
representatives and defines a metric on the equivalence classes.
Constant sequences give the isometry $j$.  If $z=[(x_n)]$, then
$j(x_n)\to z$, proving density.  To prove completeness, let $(z_m)$
be Cauchy in this metric and choose $x_m\in X$ with
$\widehat d(j(x_m),z_m)<1/m$.  The sequence $(x_m)$ is Cauchy and
defines $z\in\widehat X$.  Now $j(x_m)\to z$ and the displayed
approximation implies $z_m\to z$.

A uniformly continuous map $f:X\to Y$ sends Cauchy sequences to
Cauchy sequences: use its single uniform continuity radius for
each prescribed image accuracy.  Completeness of $Y$ therefore
defines $\widehat f([(x_n)])=\lim f(x_n)$.  Equivalent sequences
give the same limit.  To check uniform continuity of the extension,
choose $\delta>0$ so that $d(x,y)<\delta$ implies
$d_Y(f(x),f(y))<\varepsilon/2$.  If
$\widehat d(z,w)<\delta/2$, representative pairs eventually have
distance less than $\delta$; taking limits gives image distance at
most $\varepsilon/2<\varepsilon$.  Taking limits in a Lipschitz
inequality preserves its constant.  Density gives uniqueness of a
continuous extension.  Apply this construction to the isometric
identification between the dense images in two completions.  The
two extensions are inverse by density, so yield the asserted
unique surjective isometry.  The empty space is its own completion
and is handled separately throughout.  This proves the form in
\cite[Theorem 1.6-2]{B03} as well as its extension property.
\end{proof}

\begin{corollary}[The integral space associated with a positive functional]
\label{fa2:p1:positive-integral-completion}
Let $K$ be compact Hausdorff and $\ell$ a positive real functional
on $C(K)$.  Let $\mu$ be its finite regular representing measure.
The completion of $C(K)/\{f:\ell(|f|)=0\}$ in the norm
$\ell(|f|)$ is canonically $L^1(\mu)$, and the extension of
$\ell$ is integration.  In this completion the positive cone,
absolute value, and arbitrary Lipschitz scalar functional calculus
agree with their almost-everywhere function
versions.  An increasing sequence whose integrals are bounded
converges in this norm to its pointwise supremum; subtraction of
its first term reduces this assertion to a nonnegative sequence.
\end{corollary}
\begin{proof}
The Riesz representation theorem gives
$\ell(|f|)=\int|f|d\mu$, so the asserted quotient embeds
isometrically in $L^1(\mu)$.  This image is dense.  Indeed for a
measurable set $A$, regularity gives a compact $F$ and an open
$U$ with $F\subset A\subset U$ modulo null sets and
$\mu(U\setminus F)<\varepsilon$.  A continuous cutoff
$0\le\phi\le1$, equal to one on $F$ and zero outside $U$,
then satisfies $\|1_A-\phi\|_1<\varepsilon$.  Finite simple
functions are dense in $L^1$, proving density and hence the
completion assertion.  The same argument works for the completed
measure because subsets of null sets have arbitrarily small open
covers.  Integration is continuous and agrees with $\ell$ on
the dense image, so is its unique extension.

For a scalar Lipschitz map $q$, the bound
$|q(f)|\le\operatorname{Lip}(q)|f|+|q(0)|$ and finiteness of
$\mu$ give $q(f)\in L^1(\mu)$, and
$\|q(f)-q(g)\|_1\le\operatorname{Lip}(q)\|f-g\|_1$.
Thus its abstract extension coincides with the pointwise operation;
positive parts and absolute values are included.  For a
nonnegative increasing sequence $(g_n)$ with bounded integrals,
monotone convergence gives an integrable supremum $g$, and
$\|g-g_n\|_1=\int g-\int g_n\to0$.  This establishes the
abstract-completion and function-realization versions of
\cite[Appendix A.1--A.5]{B08}, without confusing finitely
additive set functions with countably additive measures.
\end{proof}

\begin{theorem}[Tight functionals on bounded continuous functions]
\label{fa2:p1:tight-cb-representation}
Let $X$ be any metric space and $L:C_b(X;\mathbb R)\to\mathbb R$
positive with $L(1)=1$.  Suppose that for every $\varepsilon>0$
there is a compact $K_\varepsilon\subset X$ such that
\[
 |L(f)|\le\sup_{K_\varepsilon}|f|+
              \varepsilon\|f\|_\infty\qquad(f\in C_b(X)).
\]
Then there is a unique finite positive Borel measure $\mu$ on $X$
such that $L(f)=\int f\,d\mu$.  It is a tight probability and is
inner regular by compact sets.  No completeness, separability,
or local compactness of $X$ is assumed.
\end{theorem}
\begin{proof}
A metric space is completely regular Hausdorff, so use its
Stone--\v Cech compactification $j:X\hookrightarrow\beta X$
from Theorem~\ref{fa2:p2:stone-cech}.  Every bounded continuous
real function on $X$ extends uniquely to $\beta X$, and restriction
is an isometric order isomorphism, since $j(X)$ is dense.
The functional induced by $L$ on $C(\beta X)$ therefore has a
regular probability measure $\nu$ by
Theorem~\ref{fa2:p4:riesz-markov}.

The image $j(K_\varepsilon)$ is compact, hence closed in $\beta X$.
For every continuous $0\le g\le1$ supported in its complement,
the stated bound gives $\int g\,d\nu\le\varepsilon$.  Regularity
and continuous cutoffs on the compact Hausdorff space imply
$\nu(\beta X\setminus j(K_\varepsilon))\le\varepsilon$.
Choose $K_n=K_{2^{-n}}$ and put $S=\bigcup_nj(K_n)$.  This is
Borel in $\beta X$ and has $\nu(S)=1$.  We do not assume that
$j(X)$ itself is Borel.  If $B$ is Borel in $X$, then
$j(B)\cap j(K_n)$ is Borel in the closed subspace $j(K_n)$,
hence Borel in $\beta X$.  Thus $j(B)\cap S$ is Borel, and
\[
 \mu(B)=\nu(j(B)\cap S)
\]
defines a countably additive probability on the Borel sets of $X$.
The integral identity follows first for simple functions on $S$,
then for the bounded function given by the continuous extension
of $f$.  The finite unions of the $K_n$ give compact sets with
$\mu$-masses tending to one.  To obtain compact inner regularity
inside an arbitrary Borel $B$, first restrict to such a finite
union and then use regularity of $\nu$ there; its compact subsets
pull back to compact subsets of $B$ in $X$.

Finally bounded continuous functions determine any finite positive
Borel measure on a metric space: for a closed nonempty $F$, the
functions $(1-n\,d(x,F))_+$ converge boundedly and pointwise to
$1_F$.  Dominated convergence determines the mass of every closed
set from the integrals.  The closed sets form an intersection-stable
class generating the Borel sets, so the finite-measure uniqueness
theorem gives equality on all Borel sets.  Empty $F$ is immediate.
This proves \cite[Chapter 1, Theorem 7.4]{B09}, including its
nonlocally compact metric-space scope.
\end{proof}

\begin{proposition}[The net tests used in duality]
\label{fa2:p1:nets}
A subset of a topological space is closed if and only if it contains
the limits of all its convergent nets.  A map is continuous if and
only if it preserves limits of all convergent nets.  A space is compact
if and only if every net has a convergent subnet.  A convergent net
in a TVS need not have bounded full range.
\end{proposition}
\begin{proof}
If $x\in\overline A$, direct its neighborhoods by reverse inclusion
and choose a point of $A$ in each neighborhood.  This net converges
to $x$, proving the closed-set criterion; its converse is immediate.
If a map fails continuity at $x$, choose in every neighborhood of
$x$ a point whose image avoids one fixed neighborhood of its image.
This gives a convergent net with nonconvergent image.  The forward
continuity implication follows from inverse images of neighborhoods.

In a compact space, the closures of all tails of a net have the
finite intersection property, so have a common point $x$.  It is
a cluster point: every neighborhood meets every tail.  To obtain
an actual subnet, use triples $(\alpha,\beta,U)$ with $\beta\ge\alpha$
and $x_\beta\in U$, ordered by increasing $\alpha,\beta$ and
decreasing neighborhoods $U$ of $x$.  They form a directed set:
first increase $\alpha$ beyond the finitely many old indices, then
use the cluster property to choose $\beta$ in the intersection
neighborhood.  The map to $\beta$ is order preserving and cofinal,
and the resulting subnet converges to $x$.
Conversely an open cover with no finite subcover gives a net
indexed by its finite subfamilies: choose a point outside each
finite union.  No subnet can converge, since a cover member
containing a proposed limit is eventually avoided.  This proves
the compactness criterion.  Finally the real net indexed by
$\mathbb N_{\ge1}\times\mathbb N_{\ge1}$ in the product order,
with value $m$ at $(m,1)$ and zero at $(m,n)$ for $n\ge2$, is eventually zero
but has unbounded range.
\end{proof}

\begin{example}[The weak topology is not determined by sequences]
\label{fa2:p1:weak-nonsequential}
In every infinite dimensional normed space choose nested subspaces
$E_n$ of dimension $n$ and put
$S=\bigcup_{n\ge1}\{x\in E_n:\|x\|=n\}$.
Then $S$ is weakly sequentially closed but is not weakly closed.
Indeed a weakly convergent sequence is norm bounded, so all its
terms lie in a finite union of the displayed finite dimensional
spheres. That union is norm compact, hence weakly compact and
weakly closed; the limit lies in $S$.
On the other hand any weak neighborhood of zero imposes only
$m$ scalar tests. If $n>m$, their common kernel in $E_n$ contains
a nonzero vector, which can be scaled to norm $n$. Thus the
neighborhood meets $S$, while $0\notin S$.
This does not contradict Eberlein--\v Smulian: that theorem
characterizes compactness, not all closed sets by sequences.
Compare \cite[Chapter 12, Theorem 1]{B08}.
\end{example}

\begin{proposition}[Global weak-to-norm continuity is finite rank]
\label{fa2:p1:weak-norm-finite-rank}
A bounded linear map $T:E\to F$ between normed spaces is continuous
from the weak topology of $E$ to the norm topology of $F$ if and
only if it has finite rank.
\end{proposition}
\begin{proof}
Continuity at zero gives finitely many $f_j\in E'$ and $\delta>0$
such that $|f_j(x)|<\delta$ for all $j$ implies $\|Tx\|<1$.
If all $f_j(x)=0$, every scalar multiple of $x$ satisfies the
same condition, so $Tx=0$. Thus $T$ factors through the finite
dimensional test map $x\mapsto(f_j(x))$, and has finite rank.
Conversely write a finite-rank bounded map in a basis of its
range as $Tx=\sum_{j=1}^mf_j(x)y_j$. The range coordinate
functionals are continuous in finite dimension, hence
$f_j\in E'$. Finitely many weak tests therefore control
$\|Tx\|\le\sum_j|f_j(x)|\,\|y_j\|$.
\end{proof}

\begin{remark}[Intrinsic completeness]
\label{fa2:p1:intrinsic-completeness}
A net is Cauchy in a TVS if for every zero neighborhood $U$ all
sufficiently late differences lie in $U$.  Any two compatible
translation-invariant metrics give this same condition, because
their balls give the same neighborhood filter at zero.  Thus
completeness for such metrics is intrinsic to the vector topology,
although completeness of an arbitrary compatible metric need not be.
For a metrizable TVS, sequence completeness for an invariant metric
is equivalent to this net completeness: choose from a Cauchy net
an increasingly late sequence whose tails have diameter at most
$2^{-n}$, take its limit, and use the Cauchy property to show the
whole net has that limit.  This justifies the completeness
conventions for F-spaces and Fr\'echet spaces used above.
\end{remark}

\begin{corollary}[Distance and the support function]
\label{fa2:p1:distance-support-duality}
Let $K$ be a nonempty closed convex subset of a real or complex
normed space $E$.  For $f\in E'$ put
$\sigma_K(f)=\sup_{k\in K}\operatorname{Re}f(k)\in(-\infty,\infty]$,
omitting $\operatorname{Re}$ in the real case.  Then
\[
 \operatorname{dist}(x,K)
 =\sup_{\|f\|\le1}
       \{\operatorname{Re}f(x)-\sigma_K(f)\}.
\]
A term with $\sigma_K(f)=\infty$ has value $-\infty$.
Neither a nearest point nor a maximizing functional is asserted.
\end{corollary}
\begin{proof}
The function $g(x)=\inf_{k\in K}\|x-k\|$ is finite and
$1$-Lipschitz.  It is convex: combine two approximating points
of $K$ using its convexity and then let their approximation
errors decrease to zero.  Its conjugate is computed by writing
$x=k+z$ with the two variables independent:
\[
 \begin{aligned}
 g^*(f)&=\sup_{k\in K,z\in E}
       \{\operatorname{Re}f(k+z)-\|z\|\}\\
 &=\sigma_K(f)+\begin{cases}0,&\|f\|\le1,\\
                         +\infty,&\|f\|>1.\end{cases}
 \end{aligned}
\]
The second supremum is zero in the first case by the dual norm
bound and the choice $z=0$.  In the other case some $z$ has
$\operatorname{Re}f(z)>\|z\|$, after a scalar phase change
if needed, and positive dilation makes the supremum infinite.
Fenchel--Moreau, applicable because $g$ is finite continuous
convex, gives the formula.  In a complex space apply that real
theorem to the underlying real space: its real continuous
functional $h$ is the real part of
$f(x)=h(x)-ih(ix)$, and $\|f\|=\|h\|$ by the scalar phase
test.  Thus the displayed supremum uses exactly the stated
complex linear dual.  This proves the source distance/support
formula \cite[Chapter 8, Theorem 18]{B08} also for unbounded $K$.
\end{proof}

\begin{corollary}[Nonmeagre images and finite algebraic codimension]
\label{fa2:p1:nonmeagre-range}\label{fa2:p1:finite-codimension-range}
Let $T:E\to F$ be bounded linear, with $E$ Banach and $F$ normed.
If the actual range $T(E)$ is nonmeagre in $F$, then $T$ is
surjective and open; in particular $F$ is Banach.
If both $E,F$ are Banach and the algebraic codimension of $T(E)$
in $F$ is finite, then $T(E)$ is closed.  Finite codimension
without the range-of-a-bounded-map hypothesis does not imply
that a subspace is closed.
\end{corollary}
\begin{proof}
Write $T(E)=\bigcup_{n\ge1}T(nB_E)$, using open unit balls.
If every $\overline{T(nB_E)}$ had empty interior, this would be
a meagre range.  Thus one closure contains a ball; taking
differences of two points in that ball and then rescaling gives
$\rho B_F\subset\overline{T(B_E)}$ for some $\rho>0$.
For $\|y\|<\rho/2$, successively choose
$\|x_n\|<2^{-n}$ so that
\[
 \left\|y-T\sum_{j=1}^nx_j\right\|<\rho\,2^{-n-1}.
\]
At each step this is possible by the rescaled closure inclusion
applied to the preceding residual.  The series converges in $E$
to an element of $B_E$, and continuity gives $T\sum x_n=y$.
Hence $(\rho/2)B_F\subset T(B_E)$, proving both surjectivity
and openness.  The induced bijection $E/\ker T\to F$ is now a
bounded linear isomorphism with bounded inverse.  Its domain is
Banach, so $F$ is complete.

For the second assertion choose a finite-dimensional algebraic
complement $Y$ of $T(E)$ in $F$.  The bounded map
$A:E\oplus Y\to F$, $A(x,y)=Tx+y$, is surjective between
Banach spaces.  Open mapping gives a constant $C$ such that
each $z\in F$ has a preimage with
$\|x\|+\|y\|\le C\|z\|$ (enlarge $C$ to avoid an infimum
attainment claim).  The component $y$ is unique because
$Y\cap T(E)=\{0\}$.  It therefore defines a linear bounded
projection $P:F\to Y$, with $\|Pz\|\le C\|z\|$.
Its kernel is exactly $T(E)$, which is closed.  For the last
warning, a discontinuous algebraic linear functional on an
infinite-dimensional Banach space has a nonclosed hyperplane
kernel: if that kernel were closed, its one-dimensional normed
quotient would make the functional continuous.  Such a functional
is obtained by giving unbounded values on a normalized independent
sequence and extending it to a Hamel basis.  The first assertion is
the nonmeagre-image theorem in \cite[Chapter 15, Theorem 16]{B08}.
Theorem 14 there literally assumes a finite-dimensional range,
which is automatically closed; the finite-codimension assertion
proved here is the stronger version obtained by its suggested
finite-dimensional extension of the domain.
\end{proof}

\begin{theorem}[Strong and weak operator topologies]
\label{fa2:p1:operator-topologies}
On $\mathcal L(E,F)$ the strong operator topology is defined by
$T\mapsto\|Tx\|$ for $x\in E$, and the weak operator topology by
$T\mapsto|g(Tx)|$ for $x\in E,g\in F'$.  Both are Hausdorff
locally convex, and
$\mathrm{WOT}\subset\mathrm{SOT}\subset\text{norm topology}$.
Their continuous linear duals coincide, consisting of finite sums
$T\mapsto\sum_jg_j(Tx_j)$; hence they have the same closed convex
sets.  Composition with a fixed bounded operator is continuous
for both topologies.  If $S_\alpha\to S$ and $T_\alpha\to T$
strongly and $\sup_\alpha\|T_\alpha\|<\infty$, then
$T_\alpha S_\alpha\to TS$ strongly.  For sequences the boundedness
condition follows from uniform boundedness when the intermediate
space is Banach.
\end{theorem}
\begin{proof}
The seminorm description proves the topological assertions and
the convergence tests.  A SOT-continuous linear functional is
bounded by $C\max_j\|Tx_j\|$ for finitely many $x_j$.  It factors
through $T\mapsto(Tx_j)_j$ into a finite product of $F$; extend
the resulting bounded functional by Hahn--Banach to that product
to obtain $\sum_jg_j(Tx_j)$.  These are WOT-continuous, and the
reverse inclusion is immediate.  Separation then proves equality
of convex closures.  For fixed composition the tests become
$Tx$ at another fixed vector or $g\circ T$ as another fixed
functional.  Finally, for each $x$,
\[
 \|T_\alpha S_\alpha x-TSx\|
 \le\|T_\alpha\|\,\|S_\alpha x-Sx\|
       +\|(T_\alpha-T)Sx\|\longrightarrow0.
\]
Uniform boundedness supplies the asserted sequence case.
\end{proof}

\begin{example}[Composition and adjoints: the exact boundaries]
\label{fa2:p1:operator-topology-boundaries}
On an infinite dimensional Hilbert space choose a unit vector $e$
and, for each finite dimensional subspace $M$, a unit vector
$u_M\perp M$.  Direct $(M,n)$ by inclusion and increasing positive
integers and put
\[
 S_{M,n}x=n^{-1}\langle e,x\rangle u_M,\qquad
 T_{M,n}x=n\langle u_M,x\rangle e.
\]
Then $S_{M,n}\to0$ in norm and $T_{M,n}\to0$ strongly, since every
fixed $x$ eventually belongs to $M$, but
$T_{M,n}S_{M,n}x=\langle e,x\rangle e$.
Thus joint SOT continuity of composition fails for general nets.
For the unilateral shift $V$ on $\ell^2$, both $V^n$ and $V^{*n}$
tend weakly to zero (first test finite support vectors and then
use their uniform norm bounds), but $V^{*n}V^n=I$.  Joint WOT
continuity fails even for sequences.  Also $V^{*n}\to0$ strongly
whereas $V^n$ does not, so adjunction is not SOT-continuous;
it is WOT-continuous because its matrix coefficients are the
conjugated transposed matrix coefficients of the original operator.
\end{example}

\begin{theorem}[Arbitrary Banach sums]
\label{fa2:p1:banach-sums}
Let $(E_i)_{i\in I}$ be Banach spaces.  For $1\le p<\infty$ the space
\[
 (\bigoplus_{i\in I}E_i)_p
 =\{x=(x_i):\sum_{i\in I}\|x_i\|^p<\infty\},\qquad
 \|x\|=(\sum_i\|x_i\|^p)^{1/p},
\]
is Banach.  An uncountable nonnegative sum means the supremum of its
finite subsums.  Also Banach are the bounded product with supremum
norm and its closed subspace
\[
 (\bigoplus_{i\in I}E_i)_{c_0}
 =\{x:\text{for every }\eps>0,\ \{i:\|x_i\|\ge\eps\}
                                           \text{ is finite}\}.
\]
Coordinate inclusions and projections are contractions.  Finite support
vectors are dense in the $p$-sum and the $c_0$-sum, but in general not
in the bounded product.
\end{theorem}
\begin{proof}
The triangle inequality in the $p$-sum follows by applying the finite
Minkowski inequality to each finite set of coordinates and taking a
supremum.  A Cauchy sequence $x^n$ has coordinate limits $x_i\in E_i$.
Given $\eps>0$, choose $N$ with $\|x^n-x^m\|\le\eps$ for $n,m\ge N$.
For every finite $F\subset I$, pass to the coordinate limits to get
$\sum_{i\in F}\|x_i^n-x_i\|^p\le\eps^p$.  Taking the supremum shows
$x-x^N$ belongs to the sum, hence so does $x$, and
$\|x^n-x\|\le\eps$ for $n\ge N$.  For the bounded product replace the
finite sums by coordinate supremums.  Uniform limits preserve the
displayed $c_0$ condition, so that subspace is closed and complete.
Truncating to a finite set makes the remainder arbitrarily small in
the first two spaces, by their definitions.  The constant scalar
family on an infinite index set cannot be uniformly approximated
by finite support families.  All coordinate norm estimates are
immediate from the defining norms.
\end{proof}

\begin{theorem}[Closed sums and quantitative decompositions]
\label{fa2:p1:closed-sums}
Let $G,L$ be closed subspaces of a Banach space $E$.  The following
conditions are equivalent:
\begin{enumerate}
\item $G+L$ is closed;
\item there is $C$ such that every $x\in G+L$ has a decomposition
      $x=g+l$ with $\|g\|+\|l\|\le C\|x\|$;
\item there is $C$ such that
      $\operatorname{dist}(g,G\cap L)\le C\operatorname{dist}(g,L)$
      for $g\in G$;
\item there is $C$ such that, for every $x\in E$,
\[
 \operatorname{dist}(x,G\cap L)
 \le C\bigl(\operatorname{dist}(x,G)+\operatorname{dist}(x,L)\bigr).
\]
\end{enumerate}
Moreover $G+L$ is closed if and only if $G^\perp+L^\perp$ is norm
closed in $E'$.  In that case
$(G\cap L)^\perp=G^\perp+L^\perp$.
\end{theorem}
\begin{proof}
If $G+L$ is closed, addition $G\oplus_1L\to G+L$ is a surjective
bounded map between Banach spaces.  Openness gives condition 2, with
a harmless enlargement of its constant to obtain actual preimages.
For $g\in G$ choose $l\in L$ nearly realizing its distance to $L$ and
apply condition 2 to $g-l=u+v$, $u\in G$, $v\in L$.  Then
$g-u=l+v\in G\cap L$, proving condition 3.  Conversely, if
$x_n=g_n+l_n\to x$, condition 3 permits replacing $g_n$ modulo
$G\cap L$ so that $\|g_n\|\le2C\operatorname{dist}(x_n,L)$ when
the distance is nonzero (the zero case has a representative zero).
This alone bounds rather than makes the representatives Cauchy.
Apply it instead to a subsequence of the differences with
$\sum_k\|x_{n_{k+1}}-x_{n_k}\|<\infty$.  Each difference has a
decomposition $u_k+v_k$ with
$\|u_k\|\le2C\|x_{n_{k+1}}-x_{n_k}\|$ and
$\|v_k\|\le(1+2C)\|x_{n_{k+1}}-x_{n_k}\|$.
The two absolutely convergent series in $G,L$ sum to a decomposition
of $x-x_{n_1}$.  Thus $G+L$ is closed.
Condition 4 implies 3 by taking $x\in G$.  For the reverse, choose
$g\in G$ and $l\in L$ with distances from $x$ arbitrarily close
to the corresponding infima.  Condition 3 and the triangle inequality
give
$\operatorname{dist}(x,G\cap L)\le\|x-g\|+C\|g-l\|$,
which implies condition 4 after adjusting the constant.

We prove the dual assertion without assuming the equality of the
annihilators.  If condition 4 holds and $f\in(G\cap L)^\perp$, then
\[
 |f(x)|\le C\|f\|\bigl(\operatorname{dist}(x,G)
                         +\operatorname{dist}(x,L)\bigr).
\]
Define $T:E/(G\cap L)\to E/G\oplus_1E/L$ by $Tx=(x+G,x+L)$.
The functional $f$ on its range is bounded by this inequality.
Extend it by Hahn--Banach to the product.  Its coordinate functionals
give $f=a+b$, $a\in G^\perp$, $b\in L^\perp$, with
$\max(\|a\|,\|b\|)\le C\|f\|$.  Hence the asserted sum equals
the closed annihilator.  Conversely, suppose $G^\perp+L^\perp$ is
norm closed.  Addition and the open mapping theorem give bounded
decompositions in that sum.  Its weak star closure is
$(G\cap L)^\perp$ by the annihilator theorem.  It is also weakly
star closed: its intersection with every dual ball is weakly star
closed, because bounded decompositions and Alaoglu produce weakly
star convergent subnet decompositions; Krein--\v Smulian then applies.
Thus it equals $(G\cap L)^\perp$.  The bounded decomposition and
the distance formula
$\operatorname{dist}(x,G\cap L)=\sup_{f\in B_{(G\cap L)^\perp}}|f(x)|$
now give condition 4.
\end{proof}

\begin{proposition}[Positive extension from a majorizing subspace]
\label{fa2:p1:positive-extension}
Let a real vector space $E$ be ordered by a convex cone, and let a
subspace $Y$ majorize $E$: every $x\in E$ has an upper bound in $Y$.
Every positive linear functional $f$ on $Y$ extends to a positive
linear functional on $E$.  In particular a positive functional on
a subspace of $C(K,\mathbb R)$ containing the constants extends
positively to $C(K,\mathbb R)$ and has norm $f(1)$.
\end{proposition}
\begin{proof}
Put $p(x)=\inf\{f(y):y\in Y,\ y\ge x\}$.  The set is nonempty by
majorization and bounded below: majorize $-x$ to find $y_0\in Y$
with $y_0\le x$; positivity gives $f(y)\ge f(y_0)$ for each admissible
$y$.  The function $p$ is sublinear and $p|_Y=f$, since $y'\ge y$
implies $f(y')\ge f(y)$.  Hahn--Banach gives $F|_Y=f$, $F\le p$.
For $c\ge0$ one has $p(-c)\le f(0)=0$, so $F(c)\ge0$.
In $C(K)$ constants majorize every continuous function, and
$-\|u\|_\infty1\le u\le\|u\|_\infty1$ proves
$|F(u)|\le f(1)\|u\|_\infty$, with equality at $1$ when $K\ne\varnothing$.
\end{proof}

\section{Isometries and function spaces}

\begin{example}[A concrete failure of norm attainment]
\label{fa2:p1:nonattainment-example}
On $c_0$ put $f(x)=\sum_{n\ge1}2^{-n}x_n$. Then $\|f\|=1$,
as shown by the vectors with their first $N$ coordinates one and
all other coordinates zero. But for $x\in B_{c_0}$ choose $N$
with $|x_n|\le1/2$ for $n\ge N$. Then
$|f(x)|\le\sum_{n<N}2^{-n}+\tfrac12\sum_{n\ge N}2^{-n}<1$.
Thus the norm is a supremum which no unit-ball point attains,
in both the real and the complex space.
\end{example}

\begin{proposition}[Finite dimensional best approximation]
\label{fa2:p1:finite-dimensional-approximation}
Every point of a normed space has a nearest point in every nonempty
closed subset of a finite dimensional subspace.  In particular it
has a nearest point in that subspace; strict convexity gives
uniqueness for convex subsets whenever a nearest point exists.
\end{proposition}
\begin{proof}
A minimizing sequence is bounded by triangle inequality.  Finite
dimensional compactness gives a convergent subsequence; closedness
and continuity of the norm give a minimizer.  Uniqueness for convex
sets in a strictly convex space was proved in
Theorem~\ref{fa2:p1:strict-approximation}.
\end{proof}

\begin{theorem}[Mazur--Ulam]
\label{fa2:p1:mazur-ulam}
A surjective isometry between real normed spaces is affine.
Consequently a surjective isometry taking zero to zero is real
linear.  Completeness is not required.
\end{theorem}
\begin{proof}
For an isometry $T:E\to F$ and fixed $x,y\in E$ define its midpoint
defect by
$d(T)=\|T((x+y)/2)-(Tx+Ty)/2\|$.
Triangle inequality bounds $d(T)\le\|x-y\|/2$.
If $T$ is onto, reflect $F$ about $(Tx+Ty)/2$, say
$Rz=Tx+Ty-z$, and set $S=T^{-1}RT$.  Then $S$ is a surjective
isometry of $E$ interchanging $x,y$, and
\[
 d(S)=\|S((x+y)/2)-(x+y)/2\|=2d(T).
\]
Indeed apply $T$ to the two points in the middle norm and use the
definition of the reflection.  Repeating this construction would
give arbitrarily large midpoint defects unless $d(T)=0$.
Thus $T$ preserves midpoints.  Subtract $T0$, use midpoint
preservation first for $0,2x$ and then for $x,y$ to obtain additivity,
then rational homogeneity, and finally real homogeneity by continuity.
\end{proof}

\begin{example}[Surjectivity and the scalar field]
The map $t\mapsto(t,\sin t)$ from $\mathbb R$ into
$(\mathbb R^2,\|\cdot\|_\infty)$ is an isometry which is not affine.
Complex conjugation on $\mathbb C$ is an onto real linear isometry
but is not complex linear.  Thus neither surjectivity nor the real
scalar conclusion can simply be omitted or strengthened.
\end{example}

\begin{theorem}[Sobczyk's separable extension theorem]
\label{fa2:p1:sobczyk}
Let $Y$ be a subspace of a separable normed space $E$.  Every bounded
map $T:Y\to c_0$ extends to a bounded map $S:E\to c_0$ with
$\|S\|\le2\|T\|$.  In particular every closed subspace isomorphic
to $c_0$ in a separable Banach space is complemented.  If the copy
is isometric, a projection of norm at most $2$ exists.
\end{theorem}
\begin{proof}
The zero operator is immediate.  Write $Ty=(f_n(y))$, where
$\|f_n\|\le R=\|T\|$ and $f_n(y)\to0$ for every $y\in Y$.
Choose Hahn--Banach extensions $g_n\in RB_{E'}$.  This ball is
weakly star compact and metrizable because $E$ is separable.
Every cluster point of $(g_n)$ belongs to the closed set
$F=RB_{E'}\cap Y^\perp$: evaluations on $Y$ converge to zero.
For a compatible metric $d$ on the ball, therefore
$\operatorname{dist}_d(g_n,F)\to0$.  Otherwise a subsequence with
distance bounded below would have a cluster point outside $F$.
Choose $h_n\in F$ with $d(g_n,h_n)\to0$.  Uniform continuity of
each evaluation on the compact ball gives
$(g_n-h_n)(x)\to0$ for every $x\in E$.
Consequently $Sx=((g_n-h_n)(x))_n$ lies in $c_0$, defines a linear
map of norm at most $2R$, and agrees with $T$ on $Y$.
If $i:c_0\to Y$ is an isomorphism, extend $i^{-1}$ and compose
the extension with $i$ to obtain a projection onto $Y$.
For an isometric $i$ its norm is at most $2$.
\end{proof}

\begin{proposition}[Sequential completeness in the weak topologies]
\label{fa2:p1:weak-sequential-completeness}
Reflexive Banach spaces are weakly sequentially complete.  The dual
of a Banach space is weakly star sequentially complete.
\end{proposition}
\begin{proof}
A weakly Cauchy sequence $(x_n)$ is norm bounded by uniform
boundedness on $E'$.  In a reflexive space its containing closed
ball is weakly compact, so it has a weakly convergent subsequence.
Every scalar sequence $f(x_n)$ is Cauchy and has the same limit
as that subsequence, proving weak convergence of the whole sequence.
If $(f_n)\subset E'$ is weakly star Cauchy and $E$ is Banach,
uniform boundedness gives $\sup\|f_n\|<\infty$; its pointwise limit
is then a bounded linear functional and is its weak star limit.
\end{proof}

\section{Polynomial approximation and convex function examples}

\begin{lemma}[Polynomial approximation on an interval]
\label{fa2:p1:weierstrass}
Every continuous scalar function on a compact real interval is a
uniform limit of polynomials.
\end{lemma}
\begin{proof}
Rescale to $[0,1]$ and use the Bernstein polynomials
$B_nf(x)=\sum_{k=0}^nf(k/n){n\choose k}x^k(1-x)^{n-k}$.
The weights are the binomial probabilities, with mean $x$ and
variance at most $1/(4n)$ for $k/n$.  Uniform continuity gives
$|f(s)-f(x)|<\eps$ when $|s-x|<\delta$.  Split the sum into this
set and its complement; the latter has total weight at most
$1/(4n\delta^2)$, by summing the squared deviations.  Hence
$\|B_nf-f\|_\infty\le\eps+2\|f\|_\infty/(4n\delta^2)\to\eps$.
Real and imaginary parts give the scalar assertion.
\end{proof}

\begin{theorem}[Stone--Weierstrass]
\label{fa2:p1:stone-weierstrass}
Let $K$ be compact Hausdorff.  A real subalgebra of $C(K,\mathbb R)$
which contains the constants and separates points is uniformly dense.
A complex subalgebra of $C(K,\mathbb C)$ with these properties is
uniformly dense if it is also closed under complex conjugation.
\end{theorem}
\begin{proof}
The uniform closure $B$ of the real algebra is again an algebra.
For $b\in B$, polynomial approximation to $|t|$ on the compact
range interval of $b$ shows $|b|\in B$.  Thus $B$ is closed under
finite maximum and minimum.  Given $f\in C(K,\mathbb R)$ and
$\eps>0$, for every $x,y$ point separation and affine rescaling
give $a_{x,y}$ in the algebra with
$a_{x,y}(x)=f(x)$ and $a_{x,y}(y)=f(y)$; for $x=y$ use a constant.
For fixed $x$, finitely many neighborhoods on which
$a_{x,y}>f-\eps$ cover $K$.  The maximum of the corresponding
functions, denoted $b_x\in B$, exceeds $f-\eps$ everywhere and
equals $f$ at $x$.  Finitely many neighborhoods where
$b_x<f+\eps$ again cover $K$.  The minimum of their $b_x$ is
within $\eps$ of $f$ everywhere.  Hence $f\in B$.
For a self-adjoint complex algebra, its real-valued members
contain real and imaginary parts and separate points, so the
real result applied to both parts proves the conclusion.
\end{proof}

\begin{example}[The conjugation hypothesis]
On the unit circle the polynomials in $z$ contain constants and
separate points, but cannot uniformly approximate $\bar z$.
Indeed the normalized circle integral of $zp(z)$ is zero for
every such polynomial, whereas that of $z\bar z$ is one.
Thus the complex conjugation hypothesis has mathematical content.
\end{example}

\begin{theorem}[Extreme points of normalized convex functions]
\label{fa2:p1:convex-function-extremes}
Let $\mathcal C$ consist of all nonnegative convex functions on
$[0,1]$ with $f(0)=0$ and $f(1)=1$, allowing an endpoint jump at
$1$.  It is a compact convex set for pointwise convergence.
Its extreme points are
\[
 e_r(x)=\frac{(x-r)_+}{1-r}\quad(0\le r<1),\qquad
 e_1(x)=\begin{cases}0&x<1,\\1&x=1.\end{cases}
\]
Every $f\in\mathcal C$ has a unique representation
$f(x)=\int_{[0,1]}e_r(x)\,d\mu(r)$ by a probability measure.
\end{theorem}
\begin{proof}
Convexity gives $0\le f(x)\le x$.  All defining inequalities are
closed conditions in $[0,1]^{[0,1]}$, so compactness follows from
the product theorem.  The right derivative of $f$ is increasing
on $(0,1)$, and $a=f'_+(0)=\lim_{x\downarrow0}f(x)/x\in[0,1]$.
Let $\nu$ be the positive locally finite Stieltjes measure on
$(0,1)$ associated with its right derivative.  For $x<1$,
\[
 f(x)=ax+\int_{(0,x)}(x-r)\,d\nu(r).
\]
For completeness, on each interior interval the left and right
derivatives bound the secant slopes.  Upper and lower Riemann
sums therefore give the integral of the monotone derivative as
the increment of $f$.  Letting the lower endpoint decrease to
zero and integrating the Stieltjes increments of that derivative
gives the displayed identity by Tonelli.  Letting $x\uparrow1$
then gives $f(1-)=a+\int_{(0,1)}(1-r)\,d\nu(r)\le1$.
Define $\mu$ to have atom $a$ at zero, restriction
$(1-r)d\nu(r)$ on $(0,1)$, and atom $1-f(1-)$ at one.
It is a probability measure, and the asserted representation
holds for $x<1$ by the formula and at $1$ because all $e_r(1)=1$.
Conversely any probability measure gives a member of $\mathcal C$.
Its mass at zero, its interior restriction, and its mass at one
are recovered by the right derivative and the endpoint jump,
so the representation is unique.

For every fixed $x$ the map $r\mapsto e_r(x)$ is continuous on
$[0,1]$, including $r=1$.  Thus the correspondence from probability
measures to $\mathcal C$ is affine and continuous for their compact
topologies, and is a bijection.  A probability measure is extreme
exactly when it is a Dirac measure: a non-Dirac regular measure
can be split over a Borel set with mass strictly between zero
and one, while a Dirac measure cannot be so split.  The affine
bijection therefore gives exactly the listed extreme functions.
\end{proof}

\section{Baire first class and generic nondifferentiability}

\begin{theorem}[Continuity points of a Baire first-class function]
\label{fa2:p1:baire-one}
Let $X$ be a complete metric space and $f_n:X\to\mathbb R$ continuous
with $f_n(x)\to f(x)$ at each point.  The continuity points of $f$
form a dense $G_\delta$ subset of $X$.
\end{theorem}
\begin{proof}
The oscillation $\omega_f(x)=\inf_{U\ni x}\sup_{y,z\in U}|f(y)-f(z)|$
has open strict sublevel sets.  We prove that
$\{\omega_f<\eps\}$ is dense for each $\eps>0$.
In a nonempty open set $U$, the relatively closed sets
\[
 E_N=\{x:|f_n(x)-f_m(x)|\le\eps/4\quad(n,m\ge N)\}
\]
cover $U$, by pointwise convergence.  An open subset of a complete
metric space is a Baire space: the usual nested closed-ball proof
starts in any ball whose closure lies in that subset.  Thus some
$E_N$ has a nonempty relatively open portion in $U$.
On it, passage to the limit gives $|f-f_N|\le\eps/4$.
Continuity of $f_N$ supplies a smaller nonempty open set on which
its oscillation is less than $\eps/4$.  The oscillation of $f$
there is less than $3\eps/4$.  Thus the desired sublevel is dense.
Baire applied to their intersection for $\eps=1/k$ proves a dense
$G_\delta$, and zero oscillation is exactly continuity.
\end{proof}

\begin{theorem}[A typical continuous function is nowhere differentiable]
\label{fa2:p1:generic-nondifferentiability}
In $C[0,1]$ with its supremum norm, the functions with no finite
derivative at any point of $(0,1)$, and no finite one-sided derivative
at the endpoints, contain a dense $G_\delta$.
\end{theorem}
\begin{proof}
For positive integers $n,m$ let $F_{n,m}$ be the set of functions
for which some $x\in[0,1]$ satisfies
\[
 |f(y)-f(x)|\le n|y-x|
 \quad\text{whenever }y\in[0,1],\ 0<|y-x|\le1/m.
\]
This set is closed.  If $f_k\to f$ uniformly and witnesses $x_k$
converge to $x$, for each $y$ with $0<|y-x|<1/m$ take nearby
$y_k$ within $1/m$ of $x_k$; passing to the limit proves the
inequality.  Boundary distances $|y-x|=1/m$ follow by continuity.

It has empty interior.  Approximate any continuous function by a
polygonal one and add an arbitrarily small, sufficiently fine
triangular wave so that every linear piece of the resulting
polygonal function $g$ has slope of absolute value greater than
$2n$.  This can be done by choosing the wave slope larger than
the maximum absolute slope of the initial polygonal function;
insert its finitely many vertices into the subdivision and avoid
coincident cancellation.  At each $x$ there is a point $y_x$ on
one adjacent linear piece, with $0<|y_x-x|<1/m$, such that
$|g(y_x)-g(x)|>n|y_x-x|$.  The inequality is strict and persists
for all $x'$ in a neighborhood of $x$, using the same $y_x$.
A finite such cover of $[0,1]$ gives a positive uniform margin,
so every sufficiently small uniform perturbation of $g$ still
lies outside $F_{n,m}$.  Hence $F_{n,m}$ is nowhere dense.
The complement of their countable union is a dense $G_\delta$.
A finite derivative at any point implies one of the local
Lipschitz inequalities in the display, for sufficiently large
$n,m$.  Thus every function in this complement has the stated
nondifferentiability property.
\end{proof}

\begin{theorem}[Universal embedding and Banach--Stone]
\label{fa2:p1:banach-stone}
Every normed space $E$ embeds linearly and isometrically into
$C(B_{E'},w^*)$ by $x\mapsto(f\mapsto f(x))$.
Let $K,L$ be compact Hausdorff spaces.  Every surjective linear
isometry $U:C(K,\mathbb K)\to C(L,\mathbb K)$ has the form
\[
 (Uf)(t)=h(t)f(\varphi(t)),
\]
where $\varphi:L\to K$ is a homeomorphism and $h\in C(L)$ satisfies
$|h(t)|=1$.  Conversely every such weighted composition is a
surjective linear isometry.  For real scalars the weight takes
values in $\{-1,1\}$.
\end{theorem}
\begin{proof}
The evaluation function is weakly star continuous, and
$\sup_{f\in B_{E'}}|f(x)|=\|x\|$ by Hahn--Banach.  Alaoglu supplies
the compact Hausdorff domain.

For the classification use the regular-measure representation of
$C(K)'$.  The extreme points of its unit ball are exactly
$\lambda\delta_s$, $s\in K$, $|\lambda|=1$.
Here are the details also for complex measures.  A measure with norm
less than one is the midpoint of two distinct measures in the ball,
unless the whole dual is zero.  A norm-one measure whose variation
is not concentrated at a single point has a Borel set $A$ with
$0<|\mu|(A)<1$.  It is the nontrivial convex combination of its
restrictions to $A$ and $A^c$, each normalized to norm one.
Such a set exists because regularity makes the support of a nonzero
measure nonempty, and two support points have disjoint neighborhoods
of positive variation.  If there is just one support point,
regularity puts all variation at that point.  Conversely, if
$\lambda\delta_s=(\nu_1+\nu_2)/2$ with $\|\nu_i\|\le1$, evaluating
the atomic mass at $s$ forces $|\nu_i(\{s\})|=1$ and equality in the
scalar triangle inequality.  Both measures have zero variation
off $s$ and equal $\lambda\delta_s$.  This proves extremality.

The dual isometry $U'$ maps extreme points bijectively.  Thus for
each $t\in L$,
$U'\delta_t=h(t)\delta_{\varphi(t)}$ with $|h(t)|=1$.
This is the desired formula, and $h=U1$ is continuous.
For every $f\in C(K)$ the function
$f\circ\varphi=(Uf)/h$ is continuous.  Continuous functions determine
the topology of a compact Hausdorff space (separate a point from a
closed set), so $\varphi$ is continuous.  Applying the same argument
to $U^{-1}$ shows that the corresponding point map is inverse to
$\varphi$, since continuous functions separate points.  Hence
$\varphi$ is a homeomorphism.  The converse follows by taking
supremums and composing with $\varphi^{-1}$.  Empty compact spaces
give zero spaces and are covered by the unique empty map.
\end{proof}

\begin{proposition}[Bounded Borel functions do not exhaust the bidual]
\label{fa2:p1:ck-bidual-borel}
For compact Hausdorff $K$, bounded Borel functions embed linearly
and isometrically in $C(K)''=M(K)'$ by
$h\mapsto[\mu\mapsto\int_Kh\,d\mu]$. This embedding need not be onto,
already for $K=[0,1]$.
\end{proposition}
\begin{proof}
The variation inequality bounds the functional norm by
$\|h\|_\infty$, and point masses detect the reverse inequality,
including when the supremum is not attained.
For a finite signed or complex measure $\mu$ on $[0,1]$, the set
of its nonzero point masses is countable and
$\sum_x|\mu(\{x\})|\le\|\mu\|$. Thus
\[
 \Phi(\mu)=\mu([0,1])-\sum_{x\in[0,1]}\mu(\{x\})
\]
is a bounded linear functional, of norm at most two. Linearity
follows by taking the countable union of the atoms of two measures
and using absolute convergence of their sums.
It vanishes on every Dirac measure but has value one on Lebesgue
measure. If $\Phi(\mu)=\int h\,d\mu$ for a bounded Borel $h$,
the Dirac tests would give $h(x)=0$ for every $x$, contradicting
the Lebesgue test. This proves non-surjectivity.
\end{proof}

\begin{remark}[A source bidual assertion requiring correction]
The identification of the full bidual with all bounded Borel
functions stated in \cite[Chapter 8, Theorem 14(ii)]{B08} does not
hold in the displayed generality. The preceding proposition gives
the valid embedding and a concrete obstruction to surjectivity.
This issue concerns that specific assertion; the regular-measure
representation of the first dual is the separate valid theorem.
\end{remark}

\section{Banach differential calculus}

\begin{definition}
\label{fa2:p1:frechet-derivative}
For open $\Omega\subset E$ and a map $F:\Omega\to F_0$ between
Banach spaces, the Fr\'echet derivative at $x$ is a bounded linear
map $DF(x)$ such that
$\|F(x+h)-F(x)-DF(x)h\|=o(\|h\|)$.
The derivative is unique, since the difference of two candidates
vanishes on each vector after replacing $h$ by $th$ and letting
$t\to0$.  The notation $C^1$ means that this derivative exists
everywhere and is continuous in operator norm.
\end{definition}

\begin{proposition}[Chain rule and the mean-value estimate]
\label{fa2:p1:banach-chain-rule}
If $F$ is differentiable at $x$ and $G$ at $F(x)$, then
$D(G\circ F)(x)=DG(F(x))DF(x)$.  If the segment $[x,y]$ is contained
in the domain and $\sup_{z\in[x,y]}\|DF(z)\|\le M$, then
$\|F(y)-F(x)\|\le M\|y-x\|$.
\end{proposition}
\begin{proof}
Write $F(x+h)-F(x)=Ah+r(h)$ with $r(h)=o(\|h\|)$, hence this
increment is $O(\|h\|)$.  Substitute it into the expansion of $G$;
its remainder is therefore $o(\|h\|)$, proving the chain rule.
For the estimate, apply any norm-one real functional on the target
to the path $t\mapsto F(x+t(y-x))$ and use the scalar mean-value
theorem.  Hahn--Banach then takes the supremum of these scalar
estimates.  Complex spaces are viewed as real for this estimate.
\end{proof}

\begin{theorem}[Inverse and implicit function theorems]
\label{fa2:p1:banach-implicit}
Let $F:\Omega\subset E\to F_0$ be $C^1$ between Banach spaces.
If $A=DF(x_0)$ is an isomorphism, $F$ has a $C^1$ local inverse
near $x_0$, with derivative
$D(F^{-1})(y)=DF(F^{-1}y)^{-1}$.
Consequently, if $H$ is $C^1$ near $(x_0,y_0)$ in $E\times Y$,
$H(x_0,y_0)=0$, and $D_yH(x_0,y_0):Y\to Z$ is an isomorphism,
the nearby solutions of $H(x,y)=0$ are exactly $y=\phi(x)$ for a
unique $C^1$ map, and
\[
 D\phi(x)=-D_yH(x,\phi(x))^{-1}D_xH(x,\phi(x)).
\]
If the initial map is $C^k$, so is the inverse or implicit map.
\end{theorem}
\begin{proof}
Translate the base points to zero.  On a sufficiently small closed
ball $\overline B(0,r)\subset\Omega$, continuity of $DF$ ensures
$\|I-A^{-1}DF(x)\|\le1/2$.  For $y$ close to $F(0)$ define
$T_y(x)=x+A^{-1}(y-F(x))$.  The mean-value estimate makes $T_y$
a contraction with constant $1/2$.  If
$\|A^{-1}(y-F(0))\|<r/2$, it maps the closed ball into itself.
The contraction theorem, proved by summing the geometric bounds
on its successive iterates, gives a unique solution $x=G(y)$.
For two such solutions,
$\|G(y)-G(z)\|\le2\|A^{-1}\|\,\|y-z\|$.
The same contraction estimate gives injectivity of $F$ on this
ball.  Restrict the target to a smaller open ball so that the
solutions stay in its interior.  Repeating the construction at
each point, or the identical estimate on a small ball centered
there, proves that the image of this interior is open.

At $x=G(y)$ the operator $DF(x)$ is invertible by the Neumann
lemma.  Put $k=G(y+h)-G(y)=O(\|h\|)$ and expand
$h=DF(x)k+o(\|k\|)$.  Thus
$k=DF(x)^{-1}h+o(\|h\|)$, proving the derivative formula.
Continuity follows from continuity of $DF$, $G$, and inversion.
For the implicit theorem apply this result to
$(x,y)\mapsto(x,H(x,y))$.  Its derivative is the invertible block
triangular map $(u,v)\mapsto(u,D_xH\,u+D_yH\,v)$.
The inverse preserves its first coordinate, giving the required
graph and uniqueness.  Differentiating $H(x,\phi(x))=0$ gives the
formula.  Finally inversion is $C^\infty$ on the invertible group:
the resolvent identity gives derivative $B\mapsto-A^{-1}BA^{-1}$,
and further derivatives follow by the product rule.  Induction in
the displayed derivative formulas proves the $C^k$ assertion.
\end{proof}

\begin{example}[Directional derivatives do not suffice]
The function $F(x,y)=x^3y/(x^6+y^2)$ off $(0,0)$, with $F(0,0)=0$,
has every directional derivative zero at the origin, but
$F(x,x^3)=1/2$ for $x\ne0$.  It is not even continuous there,
and therefore not Fr\'echet differentiable.  The Banach inverse
theorem above explicitly uses continuity in operator norm of the
Fr\'echet derivative.
\end{example}

% END INPUT


% BEGIN parts/part1-asplund.tex
\chapter{Asplund's arithmetic--infimal-convolution renorming}
\label{fa2:p1:asplund-chapter}

Strict convexity of a norm and of its dual norm are different
requirements. Two equivalent norms which supply the requirements
separately can be combined to supply both at once. We prove the
actual iteration, including its conjugate identities and geometric
error estimate, as in \cite[Problem 4, Parts A--C]{B01}.
The space in this chapter is real and normed; completeness is not
needed. Conjugates are taken against its continuous real dual.

\section{The general convex iteration}

\begin{theorem}[Arithmetic and infimal-convolution means]
\label{fa2:p1:asplund-convex-iteration}
Let $0\le\psi_0\le\varphi_0<\infty$ be convex functions on $E$
which vanish at zero. Define
\[
 \varphi_{n+1}(x)=\tfrac12(\varphi_n(x)+\psi_n(x)),\qquad
 \psi_{n+1}(x)=\tfrac12\inf_y
                    [\varphi_n(x+y)+\psi_n(x-y)].
\]
Then all iterates are finite convex functions vanishing at zero,
$0\le\psi_n\le\psi_{n+1}\le\varphi_{n+1}\le\varphi_n$, and
both sequences have the same convex limit $\theta$.
Their continuous-dual conjugates satisfy
\[
 \varphi_n^*\uparrow\theta^*,\qquad
 \psi_{n+1}^*(f)=\tfrac12[\varphi_n^*(f)+\psi_n^*(f)].
\]
At every $f$ for which $\psi_0^*(f)<\infty$, also
$\psi_n^*(f)\downarrow\theta^*(f)$; in particular this holds
on the whole dual when $D(\psi_0^*)=E'$.
If $\varphi_0$ is continuous at some $x_0$, all iterates are
continuous there, and
\[
 \varphi_{n+1}^*(f)=\tfrac12\inf_{g\in E'}
             [\varphi_n^*(f+g)+\psi_n^*(f-g)].
\]
When its value is finite the last infimum is attained.
\end{theorem}
\begin{proof}
Assume the order holds at stage $n$. Convexity and
$\varphi_n\ge\psi_n$ give
$[\varphi_n(x+y)+\psi_n(x-y)]/2\ge\psi_n(x)$.
Choosing $y=0$ gives the upper bound by
$\varphi_{n+1}(x)\le\varphi_n(x)$. Nonnegativity and the
choice $x=y=0$ give zero at the origin. Infimal convolution
is convex: choose near-minimizing splitting variables for two
points and take their convex combination; convexity of the
two summands bounds the value at the combined point, and
the errors tend to zero. This proves convexity of
$\psi_{n+1}$, while that of $\varphi_{n+1}$ is immediate.
Induction proves every assertion of order and finiteness.
Let $\varphi_n\downarrow a$ and $\psi_n\uparrow b$ pointwise.
Their arithmetic recurrence gives $a=(a+b)/2$, so $a=b=\theta$.
Passing convexity inequalities to the limit gives convexity.

For each $f$, $f(x)-\varphi_n(x)$ increases to
$f(x)-\theta(x)$. Interchanging two suprema gives
$\sup_n\varphi_n^*(f)=\theta^*(f)$.
For the other recurrence use $u=x+y$, $v=x-y$ to obtain
\begin{align*}
 \psi_{n+1}^*(f)
 &=\tfrac12\sup_{u,v}
       [f(u)+f(v)-\varphi_n(u)-\psi_n(v)]\\
 &=\tfrac12[\varphi_n^*(f)+\psi_n^*(f)].
\end{align*}
No infinity subtraction is involved: these conjugates are
nonnegative because the original functions vanish at zero.
They obey $\varphi_n^*\le\theta^*\le\psi_n^*$. At a point
where $\psi_0^*$ is finite, the decreasing sequence $\psi_n^*$
has a finite limit $b^*$, and the recurrence gives
$b^*=(\theta^*+b^*)/2$. Hence $b^*=\theta^*$.

Continuity of $\varphi_0$ supplies an upper bound $M$ on a
ball about $x_0$. Every iterate is convex, nonnegative and
bounded above there by $\varphi_0$. A finite convex function
bounded above and below on a ball is continuous at its
center: for an increment $h$ scale the segment from the
center to the sphere to get the upper error $O(\|h\|)$,
and use the opposite increment and convexity for the lower
error. Thus all iterates are continuous at $x_0$.
Fenchel--Rockafellar (Theorem~\ref{fa2:p1:fenchel-rockafellar})
applied to $\varphi_n+\psi_n-2f$ now gives
\[
 (\varphi_n+\psi_n)^*(2f)
 =\inf_{u+v=2f}[\varphi_n^*(u)+\psi_n^*(v)].
\]
The common finite domain point and continuity just checked
are exactly its qualification. A finite value has a minimizing
split by dual attainment. If the value is infinite, weak
duality and that same theorem give an infinite infimum, so
the equality still holds. Scaling by one half and writing
$u=f+g$, $v=f-g$ proves the last formula.
\end{proof}

\section{Quadratic homogeneity and a strictness estimate}

\begin{lemma}[Squares of seminorms]
\label{fa2:p1:quadratic-seminorm}
A finite nonnegative convex function $q$ with
$q(\lambda x)=\lambda^2q(x)$ for real $\lambda$ is the
square of a seminorm, and every squared seminorm has these
properties. For $0<t<1$,
\[
 q(x+y)\le q(x)/t+q(y)/(1-t),\qquad
 4q(x)\le q(x+y)/t+q(x-y)/(1-t).
\]
A norm is strictly convex in the unit-sphere sense if and
only if its square is a strictly convex function.
\end{lemma}
\begin{proof}
Convexity applied to
$x+y=t(x/t)+(1-t)(y/(1-t))$ gives the first inequality.
Minimize its right side in $t$, obtaining
$q(x+y)\le(\sqrt{q(x)}+\sqrt{q(y)})^2$.
When one value is zero take the appropriate limit of $t$;
when both vanish the inequality gives zero directly.
This is the triangle inequality for $\sqrt q$; its absolute
homogeneity follows from that of $q$. Conversely convexity
of a seminorm and scalar convexity of the square give
convexity of its square. Apply the first displayed inequality
to $(x+y)+(x-y)=2x$ to prove the second.
Finally
$p(tx+(1-t)y)^2\le(tp(x)+(1-t)p(y))^2
\le tp(x)^2+(1-t)p(y)^2$.
Equality in the second inequality requires $p(x)=p(y)$;
for distinct points of this common nonzero sphere strict
convexity makes the first inequality strict. The converse
is obtained by restricting strict convexity of $p^2$ to
two distinct unit vectors.
\end{proof}

\begin{theorem}[The factor four estimate and persistence of strictness]
\label{fa2:p1:asplund-rate-strictness}
Suppose additionally that $\varphi_0,\psi_0$ are quadratically
homogeneous and $\varphi_0\le(1+C)\psi_0$, $C\ge0$.
All iterates and $\theta$ have the same homogeneity, and
\[
 \varphi_n\le(1+C4^{-n})\psi_n.
\]
If either initial function is strictly convex, then $\theta$
is strictly convex.
\end{theorem}
\begin{proof}
Homogeneity passes through the arithmetic mean and through
the infimum by replacing $y$ by $\lambda y$; zero scalars
use the already proved value zero.
Suppose $\psi\le\varphi\le a\psi$, $a\ge1$, at one stage.
The preceding convex-order proof gives $\psi_{\rm new}\ge\psi$.
Use the second inequality of the lemma for $\varphi$ with
$t=1/(a+1)$ and then $\varphi\le a\psi$ to obtain
\[
 4\varphi(x)
 \le(a+1)\varphi(x+y)+(a+1)a^{-1}\varphi(x-y)
 \le(a+1)[\varphi(x+y)+\psi(x-y)].
\]
Taking the infimum gives
$\psi_{\rm new}\ge2\varphi/(a+1)$.
Consequently
\[
 \varphi_{\rm new}=\tfrac12(\varphi+\psi)
 \le\tfrac12[(a+1)/2+1]\psi_{\rm new}
 =\tfrac{a+3}{4}\psi_{\rm new}.
\]
Starting with $a=1+C$ proves the claimed factor $4^{-n}$.

Assume first that $\varphi_0$ is strictly convex.
Unrolling the arithmetic recurrence writes
$\varphi_n=v_n+2^{-n}\varphi_0$, where $v_n$ is convex
($v_0=0$, $v_{n+1}=(v_n+\psi_n)/2$).
Since $\psi_n\le\theta\le\varphi_n\le\varphi_0$ and
$\varphi_n\le(1+C4^{-n})\theta$, we have the pointwise bounds
\[
 v_n+(2^{-n}-C4^{-n})\varphi_0
              \le\theta\le v_n+2^{-n}\varphi_0.
\]
Fix distinct $x,y$ and $0<t<1$, put $z=tx+(1-t)y$, and
let $D=t\varphi_0(x)+(1-t)\varphi_0(y)-\varphi_0(z)>0$.
Convexity of $v_n$ gives
\[
 t\theta(x)+(1-t)\theta(y)-\theta(z)
 \ge2^{-n}D-C4^{-n}[t\varphi_0(x)+(1-t)\varphi_0(y)].
\]
For sufficiently large $n$ its right side is positive.
This proves strict convexity of $\theta$; pointwise limits
of strictly convex functions alone would not justify it.
If instead $\psi_0$ is strictly convex, their first arithmetic
mean $\varphi_1$ is strictly convex. Restarting the same
iteration at that stage proves the result, with its unchanged
limit and its already established comparability.
\end{proof}

\section{Simultaneous primal and dual strict convexity}

\begin{theorem}[Asplund's simultaneous renorming theorem]
\label{fa2:p1:asplund-renorming}
Suppose a real normed space has two equivalent norms $p_1,p_2$
such that $p_1$ and the dual norm $p_2^*$ are strictly convex.
The dual is taken for their common norm topology, which we use
on $E$ in this application.
It then has an equivalent norm $p$ for which both $p$ and $p^*$
are strictly convex. This is a conditional combination theorem;
it does not assert the existence of the two initial norms on
every Banach space.
\end{theorem}
\begin{proof}
Choose $\varepsilon>0$ so
$\varepsilon p_2^2\le p_1^2\le a\varepsilon p_2^2$ for some
$a\ge1$, and start the iteration with
$\varphi_0=p_1^2/2$, $\psi_0=\varepsilon p_2^2/2$.
They are continuous everywhere and are positive comparable
quadratic functions. Thus $p=\sqrt{2\theta}$ is an equivalent
norm and is strictly convex by the preceding theorem.

For any norm $q$, direct maximization in the radial variable
gives $(q^2/2)^*=(q^*)^2/2$. The upper bound follows from
$f(x)\le q^*(f)q(x)$, and the lower bound uses a sequence
of unit vectors nearly norming $f$ and their optimal scalar
multiples; actual norm attainment is not assumed.
Hence $\varphi_0^*=(p_1^*)^2/2$ and
$\psi_0^*=(p_2^*)^2/(2\varepsilon)$, finite everywhere.
Conjugation reverses order, and quadratic scaling gives
$\varphi_0^*\le\psi_0^*\le a\varphi_0^*$.

On the dual space put $\widetilde\varphi_n=\psi_n^*$,
$\widetilde\psi_n=\varphi_n^*$. The two exact conjugate
identities above show that these satisfy the very same
arithmetic/infimal-convolution iteration: the infimum's
two terms exchange by replacing $g$ with $-g$.
Their common limit is $\theta^*$, since $\psi_0^*$ is finite
everywhere. Its initial upper function
$\widetilde\varphi_0=\psi_0^*$ is strictly convex by the
hypothesis on $p_2^*$ and the squared-norm lemma.
The strictness theorem on the normed space $E'$ therefore
makes $\theta^*$ strictly convex. Finally
$\theta=p^2/2$ implies $\theta^*=(p^*)^2/2$, so the same
lemma makes $p^*$ strictly convex as well.
The proof uses neither reflexivity nor completeness, and
all continuity qualifications needed for conjugation were
verified for the actual iterates.
\end{proof}

\begin{corollary}[Complex norms]
\label{fa2:p1:asplund-complex}
The simultaneous renorming theorem also holds for two equivalent
complex norms with the stated primal and complex-dual strictness.
\end{corollary}
\begin{proof}
Apply the real construction to the underlying real space.
Every iterate is invariant under $x\mapsto e^{it}x$: the initial
functions are, arithmetic preserves this, and changing the
infimum variable to $e^{it}y$ preserves the other operation.
Thus its resulting norm is complex homogeneous. The map
$f\mapsto\operatorname{Re}f$ is an onto real linear isometry
from the complex continuous dual to the real continuous dual;
its inverse sends $u$ to $u(x)-iu(ix)$. It therefore preserves
strict convexity of both the initial and resulting dual norms.
\end{proof}

% END INPUT


% BEGIN parts/part1-infimal-convolution.tex
\chapter{Infimal convolution under a full-cone qualification}
\label{fa2:p1:infimal-convolution-chapter}

A continuity-at-a-common-point qualification is useful but does not
include every exact conjugate formula. The reflexive full-cone
condition below is a different sufficient hypothesis. We prove its
boundedness and weak-limit steps explicitly, as required in
\cite[Problem 12, Parts A--B]{B01}. The first two theorems use a
real reflexive Banach space, as in that source. The final theorem
gives the mature Banach-space form using Krein--\v Smulian.

\begin{theorem}[Closure of conjugate domains]
\label{fa2:p1:infconv-domain-closure}
Let $\varphi,\psi:E\to(-\infty,+\infty]$ be proper convex lower
semicontinuous functions with
$D(\varphi)\cap D(\psi)\ne\varnothing$. Set
\[
 h(f)=(\varphi^*\mathbin\square\psi^*)(f)
      =\inf_{g\in E'}[\varphi^*(f-g)+\psi^*(g)].
\]
Then $h$ is proper convex, never has value $-\infty$,
\[
 D(h)=D(\varphi^*)+D(\psi^*),\quad
 h^*=\varphi+\psi,
 \quad h^{**}=(\varphi+\psi)^*,
\]
where reflexivity identifies the dual of $E'$ with $E$.
Consequently
\[
 \overline{D((\varphi+\psi)^*)}
       =\overline{D(\varphi^*)+D(\psi^*)}.
\]
The closures are norm closures in $E'$; convexity makes them
also their weak closures. No exact equality of unclosed domains
is asserted without an additional hypothesis.
\end{theorem}
\begin{proof}
Choose $x_0$ in the common domain. Fenchel's inequality gives
\[
 \varphi^*(f-g)+\psi^*(g)
       \ge f(x_0)-\varphi(x_0)-\psi(x_0),
\]
so $h$ has a continuous affine lower bound and is never
$-\infty$. Both conjugates have nonempty effective domains:
the affine-minorant proof of Fenchel--Moreau supplies such
minorants for every proper lsc convex function. Choosing one
point in each domain makes $h$ finite somewhere. Near-minimizing
splits and convexity prove that $h$ is convex. A finite split
exists exactly when $f$ is the sum of two domain points; the
lower bound just proved makes its infimum finite. This proves
the domain identity.

For $x\in E=E''$, change variables $a=f-g$ to obtain
\begin{align*}
 h^*(x)
 &=\sup_{a,g}[a(x)-\varphi^*(a)+g(x)-\psi^*(g)]\\
 &=\varphi^{**}(x)+\psi^{**}(x)=\varphi(x)+\psi(x).
\end{align*}
The independent suprema cause no undefined infinity difference,
since these biconjugates are the proper original functions.
Conjugating once more gives the displayed biconjugate formula.

We justify the passage from a biconjugate to domain closure.
Take the norm closure of the epigraph of $h$ in $E'\times\mathbb R$.
It is a closed convex upward set, and its affine lower bound
excludes fibers extending to $-\infty$. Thus it is the epigraph
of a proper lsc convex function $\bar h\le h$.
A continuous affine minorant bounds $h$ if and only if it
bounds this closed epigraph, so $\bar h^*=h^*$ and
Fenchel--Moreau gives $\bar h=h^{**}$.
If $\bar h(f)<\infty$, a point over $f$ in this closed epigraph
is approached by points over $D(h)$; hence
$D(\bar h)\subset\overline{D(h)}$.
Conversely $\bar h\le h$ gives $D(h)\subset D(\bar h)$.
Taking closures and substituting $\bar h=(\varphi+\psi)^*$
proves the assertion. Equality of convex norm and weak closures
was proved by Hahn--Banach separation in Part I.
\end{proof}

\begin{theorem}[Exact convolution under the full-cone condition]
\label{fa2:p1:infconv-full-cone}
Under the preceding hypotheses assume
\[
 \bigcup_{\lambda\ge0}\lambda
                (D(\varphi)-D(\psi))=E.
\]
For every norm bounded set $B\subset E'$ and every real $\alpha$,
the set
\[
 \{g\in E':\text{some }f\in B\text{ satisfies }
                 \varphi^*(f-g)+\psi^*(g)\le\alpha\}
\]
is norm bounded. The infimal convolution $h$ is lsc, and
\[
 (\varphi+\psi)^*(f)
       =\min_{g\in E'}[\varphi^*(f-g)+\psi^*(g)]
\]
whenever the left side is finite; at infinite values the equality
holds with an infimum. In particular
\[
 D((\varphi+\psi)^*)=D(\varphi^*)+D(\psi^*).
\]
The corresponding primal--dual identity is
\[
 \inf_x[\varphi(x)+\psi(x)]
       =\sup_{g\in E'}[-\varphi^*(-g)-\psi^*(g)],
\]
with a finite maximizing value attained whenever the infimum
is finite. No primal minimizer is asserted. If the primal
infimum is $-\infty$, both sides are $-\infty$ in the
extended-real sense.
\end{theorem}
\begin{proof}
The empty $B$ case is trivial; otherwise write
$M=\sup_{f\in B}\|f\|$. If $g$ belongs to the displayed
set, then for any $x\in D(\varphi)$ and $y\in D(\psi)$,
the two Fenchel inequalities give
\[
 g(y-x)\le\alpha+\varphi(x)+\psi(y)-f(x)
          \le\alpha+\varphi(x)+\psi(y)+M\|x\|.
\]
For every fixed $v\in E$, the full-cone hypothesis represents
$v$ as a nonnegative multiple of some $y-x$: apply the
hypothesis to $-v$ and reverse its sign. It therefore gives
a uniform upper bound on $g(v)$. Applying it also to $-v$
gives a uniform lower bound. Thus this family of members of
$E'$ is pointwise bounded on the Banach space $E$.
Uniform boundedness gives the asserted common norm bound.
This proves both the fixed-$f$ sublevel bound and the stronger
uniform bound for varying bounded $f$.

For $h(f)<\infty$ take splits $g_n$ with costs tending to
$h(f)$. Their costs are bounded above; the preceding result
makes $g_n$ norm bounded. Reflexivity of $E'$ gives a weakly
convergent subsequence $g_n\rightharpoonup g$.
The conjugates are convex norm-lsc, hence weakly lsc.
Their sum at $(f-g_n,g_n)$ is weakly lsc as well: subtract
their affine lower bounds coming from the fixed common-domain
point $x_0$ to make both summands nonnegative before applying
the liminf inequality. It follows that
$\varphi^*(f-g)+\psi^*(g)\le h(f)$, proving attainment.

To prove lsc of $h$, let $f_n\to f$ in norm. Its affine lower
bound excludes $\liminf h(f_n)=-\infty$.
If that liminf is infinite there is nothing to prove.
Otherwise pass to a subsequence of finite values converging
to the finite liminf, and choose approximate splits whose
cost differs from $h(f_n)$ by at most $1/n$.
Now $f_n$ is bounded and the costs are bounded above, so the
uniform version of the first argument bounds $g_n$.
Take a weakly convergent subsequence. Then
$f_n-g_n\rightharpoonup f-g$, and the same nonnegative
liminf argument gives
$h(f)\le\varphi^*(f-g)+\psi^*(g)\le\liminf h(f_n)$.
Thus $h$ is norm-lsc. The preceding theorem and
Fenchel--Moreau now give
$h=h^{**}=(\varphi+\psi)^*$.
Together with finite-value attainment this proves the exact
formula and domain identity. At $f=0$, negate that formula
to obtain the stated primal--dual equality and attainment.
The common finite domain point rules out a primal value
$+\infty$; the case $-\infty$ has exactly the extended
meaning stated in the theorem.
\end{proof}

\begin{remark}[Why the two qualifications are kept separate]
The proof uses the full cone to obtain pointwise bounds on all
split functionals, completeness of $E$ for uniform boundedness,
and reflexivity for weakly convergent subsequences of those
splits. These are not consequences of an unstated continuity
point. The earlier Fenchel--Rockafellar theorem uses a continuity
qualification and works in a general LCS; this theorem supplies
the different source cone qualification and its stronger
sublevel-set information.
\end{remark}

\begin{theorem}[The Banach closed-subspace cone qualification]
\label{fa2:p1:infconv-banach-closed-cone}
Let $E$ be a real Banach space, and let $\varphi,\psi$ be proper
convex norm-lsc functions with a common finite domain point.
Suppose
\[
 L=\bigcup_{\lambda\ge0}\lambda(D(\varphi)-D(\psi))
\]
is a closed linear subspace of $E$. Then
\[
 (\varphi+\psi)^*(f)
       =\inf_{g\in E'}[\varphi^*(f-g)+\psi^*(g)],
\]
and every finite value has a minimizing split. Hence the exact
domain identity and the primal--dual equality of the preceding
theorem hold. No reflexivity hypothesis is needed in this form.
When $L=E$ the uniform bounded-split sublevel assertion holds
as well. For a proper subspace $L$, it concerns bounded
representatives after restriction to $L$, rather than all
ambient split functionals.
\end{theorem}
\begin{proof}
First let $L=E$. The proof of the bounded-split assertion used
only Fenchel's inequality, the full cone, and uniform boundedness
on Banach $E$: for bounded $f$ and bounded-above split cost it
gives a uniform upper bound on $g(y-x)$, then on $g(v)$ and
$g(-v)$ for every $v\in E$. Thus the same norm bound on $g$
holds here. The infimal convolution $h$ is proper convex and
has the affine lower bound from a common-domain point, without
using reflexivity.
For a finite value $h(f)$ choose near-minimizing $g_n$.
They are norm bounded, so Alaoglu supplies a weakly star
convergent subnet. Each conjugate is weakly star lsc, being
a supremum of evaluations minus constants. Subtracting its
affine minorant makes each summand nonnegative and justifies
the liminf of their sum. The limit split therefore attains
$h(f)$.

We need weak star lsc of the whole value function, not merely
norm lsc. Consider the convex epigraph of $h$ in
$E'\times\mathbb R=(E\times\mathbb R)'$.
For any net $(f_\alpha,t_\alpha)$ in this epigraph and in a
fixed norm ball which converges weakly star to $(f,t)$,
choose its minimizing splits $g_\alpha$; they exist because
$h(f_\alpha)\le t_\alpha<\infty$. The uniform split bound,
using bounded $f_\alpha$ and costs bounded above by the
same radius, puts $g_\alpha$ in one dual ball.
Take a weakly star convergent subnet. Lower semicontinuity
of the split cost gives
$h(f)\le\varphi^*(f-g)+\psi^*(g)\le t$.
The dual ball itself is weakly star closed, so the epigraph's
intersection with every norm ball is weakly star closed.
Krein--\v Smulian, applied to the Banach predual
$E\times\mathbb R$, makes the epigraph weakly star closed.

Apply Fenchel--Moreau on the LCS $(E',\sigma(E',E))$.
Its exact continuous dual is $E$, as proved for dual pairs
in Part I. Conjugating $h$ in this pairing gives
$h^*(x)=\varphi(x)+\psi(x)$ by the independent-suprema
calculation and Fenchel--Moreau for the original functions
on normed $E$. Weak star lsc now gives
$h=h^{**}=(\varphi+\psi)^*$, with the already proved
finite-value attainment. This proves the full-cone case.

For a closed linear $L$, fix $x_0$ in the common domain.
Both domains lie in $x_0+L$: subtract $x_0$ from a point
in either domain and use linearity of $L$.
On the Banach subspace $L$ set
$\varphi_L(z)=\varphi(x_0+z)$ and
$\psi_L(z)=\psi(x_0+z)$. They satisfy the full-cone
condition in $L$, so its case just proved applies.
For $f\in E'$ write $f_L=f|_L$. The translation identities are
\[
 \varphi^*(f)=f(x_0)+\varphi_L^*(f_L),\qquad
 (\varphi+\psi)^*(f)=f(x_0)+(\varphi_L+\psi_L)^*(f_L).
\]
Every member of $L'$ extends to $E'$ by Hahn--Banach.
Thus the infimum over ambient $g$ is exactly the infimum
over $g_L\in L'$, with the fixed additional term $f(x_0)$.
The formula and attainment follow by extending a minimizing
$g_L$. At a finite primal value take $f=0$ as before.
If $L\ne E$, adding a member of $L^\perp$ to a split
does not change its total cost, so a bound on every ambient
split would indeed be an unwarranted extra conclusion.
\end{proof}

% END INPUT


% END INPUT

\part{Hilbert operators, algebras and spectral theory}

% BEGIN parts/part2.tex
% Part II. Inner products are conjugate linear in the first variable.
\providecommand{\faip}[2]{\langle #1,#2\rangle}
\providecommand{\faB}{\mathcal B}
\providecommand{\faK}{\mathcal K}
\providecommand{\faD}{\mathcal D}
\providecommand{\faRan}{\operatorname{Ran}}
\providecommand{\faKer}{\operatorname{Ker}}
\providecommand{\faTr}{\operatorname{tr}}
\providecommand{\faind}{\operatorname{ind}}
\providecommand{\fasupp}{\operatorname{supp}}
\providecommand{\faess}{\operatorname{ess}}

% BEGIN parts/part2-hilbert.tex
\chapter{Hilbert geometry and variational equations}
We use real or complex Hilbert spaces. In the complex case
$\faip{x}{y}$ is conjugate linear in $x$ and linear in $y$.
The dual $H'$ consists of continuous \emph{linear} functionals; the antidual
consists of continuous conjugate linear functionals. These must not be confused
when an operator associated with a form is defined. The sources for this chapter
are \cite{B01,B02,B03}.

\begin{theorem}[Projection onto a closed convex set]\label{fa2:p2:projection}
If $C$ is a nonempty closed convex subset of $H$, each $x\in H$ has a unique
nearest point $p=P_Cx$. It is characterized by
\[
 \operatorname{Re}\faip{x-p}{v-p}\leq0\quad(v\in C).
\]
Moreover $\|P_Cx-P_Cy\|\leq\|x-y\|$. If $M$ is a closed linear subspace,
$H=M\oplus M^\perp$ and $P_M$ is linear, self-adjoint and contractive.
Every linear subspace $L$, closed or not, satisfies
$(L^\perp)^\perp=\overline L$.
\end{theorem}
\begin{proof}
Put $d=\inf_{v\in C}\|x-v\|$ and choose $v_n\in C$ with
$\|x-v_n\|\to d$. The parallelogram identity and convexity give
\[
 \|v_n-v_m\|^2
 =2\|x-v_n\|^2+2\|x-v_m\|^2
       -4\left\|x-\frac{v_n+v_m}{2}\right\|^2\longrightarrow0.
\]
Completeness and closedness give $v_n\to p\in C$, with $\|x-p\|=d$.
For $v\in C$ and $0<t\leq1$, compare $p$ with $p+t(v-p)$:
$0\leq-2t\operatorname{Re}\faip{x-p}{v-p}+t^2\|v-p\|^2$.
Divide by $t$ and let $t\downarrow0$. Conversely this inequality gives
$\|x-v\|^2=\|x-p\|^2+\|v-p\|^2
-2\operatorname{Re}\faip{x-p}{v-p}\geq\|x-p\|^2$.
If both $p,q$ are minimizers, the displayed parallelogram estimate, now with
constant sequences, gives $p=q$. Apply the variational inequality for $x,p$
with $v=q=P_Cy$, and for $y,q$ with $v=p$. Adding and rearranging gives
$\|p-q\|^2\leq\operatorname{Re}\faip{x-y}{p-q}$, proving nonexpansion.
For a subspace $M$, use $v=p+tz$ for real $t$, and also $iz$ in the complex
case, to obtain $x-p\perp M$. Uniqueness of this orthogonal decomposition
proves linearity. The identities $P_M^2=P_M$ and
$\faip{P_Mx}{y}=\faip{x}{P_My}$ follow by decomposing both vectors.
Finally apply this decomposition to $\overline L$; a vector in
$(L^\perp)^\perp$ has zero component in $(\overline L)^\perp$.
\end{proof}

\begin{proposition}[Projection lattice and monotone limits]\label{fa2:p2:projection-lattice}
Orthogonal projections form a complete lattice: the meet of a family
projects onto the intersection of their ranges, and its join onto
the closed span of their ranges. One has
$P\leq Q$ in operator order iff $\faRan P\subset\faRan Q$, iff
$PQ=QP=P$. Commuting projections have product equal to their meet.
If bounded self-adjoint $A_n$ satisfy $A_n\leq A_{n+1}\leq B$,
they converge strongly to a bounded self-adjoint supremum.
Commutation among the $A_n$ is not needed for this generalization
of the elementary monotone-operator theorem.
\end{proposition}
\begin{proof}
Intersections and closed spans are closed subspaces, so the projection
theorem gives their projections. Inclusion of ranges is exactly
the stated product identity; it makes $Q-P$ an orthogonal projection,
hence positive. Conversely if $Q-P\geq0$ and $Px=x$, then
$\faip{x}{Qx}\geq\|x\|^2$; contractivity of $Q$ forces $Qx=x$,
giving range inclusion. The meet and join are now the greatest
lower and least upper bounds by their subspace definitions.
If $PQ=QP$ the product is self-adjoint idempotent and its range is
the intersection, proving the meet formula.
For monotone limits, subtract $A_1$ to reduce to
$0\leq C_n\leq C$ with a common norm bound $M$.
The positive-form Cauchy--Schwarz inequality gives
$\|Dx\|^2\leq\|D\|\faip{x}{Dx}$ for every positive bounded $D$:
take the supremum over unit $y$ in
$|\faip{y}{Dx}|^2\leq\faip{y}{Dy}\faip{x}{Dx}$.
For $m\geq n$, apply it to $D=C_m-C_n$, whose norm is bounded
by $2M$. The scalar sequence $\faip{x}{C_nx}$ increases and is
bounded, so $C_nx$ is Cauchy. Its limit defines a bounded operator
by the uniform bound. Passing to inner products proves
self-adjointness and all order inequalities; any common upper
bound also dominates the limit. This proves the supremum assertion.
\end{proof}

\begin{theorem}[Riesz representation]\label{fa2:p2:riesz}
For every $f\in H'$ there is exactly one $h\in H$ such that
$f(x)=\faip{h}{x}$ for all $x$, and $\|f\|=\|h\|$.
Consequently every bounded sesquilinear form $a$ that is conjugate linear
in its first variable has a unique $A\in\faB(H)$ such that
$a(x,y)=\faip{x}{Ay}$, with $\|A\|=\|a\|$.
\end{theorem}
\begin{proof}
If $f\ne0$, the closed hyperplane $M=\ker f$ has an orthogonal complement.
Choose $z\notin M$ and put $u=z-P_Mz$. Then $f(u)=f(z)\ne0$ and
$u\perp M$. For any $x$, $x-f(x)u/f(u)\in M$, whence
$\faip{u}{x}=f(x)\|u\|^2/f(u)$. Thus
$h=\overline{f(u)}u/\|u\|^2$ represents $f$. The zero functional is
represented by zero. Testing a difference on itself proves uniqueness;
Cauchy--Schwarz and testing $x=h/\|h\|$ prove equality of norms.
For fixed $y$, apply the conjugate version of this representation to
$x\mapsto a(x,y)$. Uniqueness proves that the representing vector $Ay$
depends linearly on $y$, and its norm is at most $\|a\|\|y\|$.
Taking suprema in the form identity proves equality of norms.
\end{proof}

\begin{theorem}[Arbitrary orthonormal bases]\label{fa2:p2:orthonormal-basis}
Every Hilbert space has a maximal orthonormal family $(e_i)_{i\in I}$.
For such a family and every $x\in H$,
\[
 x=\sum_{i\in I}e_i\faip{e_i}{x},\qquad
 \|x\|^2=\sum_{i\in I}|\faip{e_i}{x}|^2.
\]
Sums mean the net over finite subsets of $I$. For each fixed $x$ only
countably many coefficients are nonzero. Two Hilbert spaces are unitarily
isomorphic exactly when their orthonormal bases have the same cardinality.
\end{theorem}
\begin{proof}
For a finite orthonormal family, subtracting its Fourier sum and expanding
the square gives Bessel's inequality. Hence for each positive integer $n$
only finitely many coefficients can have absolute value at least $1/n$;
otherwise their finite square sums would exceed $\|x\|^2$. Their union is
countable, and their square sum is finite. The corresponding finite-sum
net is Cauchy by Pythagoras and therefore has a limit $s$. The difference
$x-s$ is orthogonal to every $e_i$. Zorn's lemma supplies a maximal family
(the union of a chain is orthonormal); maximality forces $x=s$.
Pythagoras now gives Parseval. Conversely a total orthonormal family is
maximal by orthogonality.
For completeness we check the cardinal invariant. If $(e_i)_{i\in I}$ and
$(f_j)_{j\in J}$ are infinite bases, each $f_j$ has countable nonempty
support in the first basis, and each $e_i$ has nonzero coefficient against
at least one $f_j$. Thus $I$ is the union of $|J|$ countable sets, so
$|I|\leq |J|\aleph_0=|J|$; symmetry and the cardinal comparison theorem
give equality. If one basis is finite, linear independence and totality
give the ordinary finite-dimensional equality. A bijection of bases
defines an isometry on finite linear combinations. It extends by
completeness to a unitary, since its range is closed and contains a total
family. A unitary carries a basis to a basis, proving the converse.
\end{proof}

\begin{theorem}[Adjoint and reducing subspaces]\label{fa2:p2:adjoint}
For $T\in\faB(H,K)$ there is a unique $T^*\in\faB(K,H)$ with
$\faip{Tx}{y}=\faip{x}{T^*y}$. One has
\[
 \|T\|=\|T^*\|,\quad\|T^*T\|=\|T\|^2,\quad
 (ST)^*=T^*S^*,\quad T^{**}=T,
\]
and
$\overline{\faRan T}=(\ker T^*)^\perp$,
$\overline{\faRan T^*}=(\ker T)^\perp$.
A closed subspace $M$ is invariant under $T$ iff
$P_MTP_M=TP_M$; it reduces $T$ iff $P_MT=TP_M$, equivalently it is
invariant under both $T$ and $T^*$.
\end{theorem}
\begin{proof}
Riesz applied to $x\mapsto\faip{y}{Tx}$ constructs $T^*y$;
uniqueness proves linearity and $\|T^*y\|\leq\|T\|\|y\|$.
Taking the supremum of $|\faip{Tx}{y}|$ over both unit balls shows
$\|T\|=\|T^*\|$. Moreover
$\|Tx\|^2=\faip{x}{T^*Tx}\leq\|T^*T\|\|x\|^2$;
the reverse norm inequality is submultiplicativity. Testing against
arbitrary $x,y$ proves the product and double-adjoint identities.
The equivalence $y\perp\faRan T\iff T^*y=0$, followed by double
orthogonal complementation, proves the range assertions. The projection
identity for invariance says exactly that $T(M)\subset M$.
Applying it also to $M^\perp$ gives commutation. Finally
$T^*(M)\subset M$ is equivalent, by the adjoint identity, to
$T(M^\perp)\subset M^\perp$.
\end{proof}

\begin{theorem}[Closed range and minimum modulus]\label{fa2:p2:closed-range}
For $T\in\faB(H,K)$ its range is closed iff there is $c>0$ such that
$\|Tx\|\geq c\|x\|$ for $x\perp\ker T$. This is equivalent to closedness
of $\faRan T^*$. In this case the range equalities in the preceding theorem
hold without closures. The reduced minimum moduli
\[
 \gamma(T)=\inf\{\|Tx\|:x\perp\ker T,\ \|x\|=1\}
\]
and $\gamma(T^*)$ are equal (with the infimum of the empty set interpreted
as $+\infty$ for the zero operator on a zero complement).
\end{theorem}
\begin{proof}
The restriction $T_0:(\ker T)^\perp\to\faRan T$ is a bounded bijection.
If the range is closed, its inverse is bounded by the open mapping theorem;
this is the stated lower bound. Conversely the lower bound makes inverse
images of convergent range sequences Cauchy, proving closedness.
For $z\in(\ker T)^\perp$, define on $\faRan T$ the functional
$Tx\mapsto\faip{z}{x}$, taking $x\perp\ker T$. Its norm is at most
$\|z\|/c$. Extend by zero on $(\faRan T)^\perp$ and apply Riesz.
This represents $z=T^*y$ with $y\in\faRan T$ and
$\|y\|\leq\|z\|/c$. Thus $\faRan T^*=(\ker T)^\perp$ is closed.
Symmetry proves the equivalence. On the two nonzero complements,
$(T_0^*)^{-1}=(T_0^{-1})^*$ and hence
$\gamma(T^*)=\|(T_0^*)^{-1}\|^{-1}
=\|T_0^{-1}\|^{-1}=\gamma(T)$ when ranges are closed.
If they are not closed, both infima are zero by the first equivalence.
The zero case is read directly from the definition.
\end{proof}

\begin{theorem}[Lax--Milgram]\label{fa2:p2:lax-milgram}
Let $a:H\times H\to\mathbb C$ be bounded, conjugate linear in its first
variable, with $\operatorname{Re}a(v,v)\geq\alpha\|v\|^2$ for some
$\alpha>0$. For every continuous conjugate linear functional $F$ there is
exactly one $u\in H$ with $a(v,u)=F(v)$ for all $v$, and
$\|u\|\leq\|F\|/\alpha$. The same assertion holds over $\mathbb R$.
If $a$ is Hermitian, $u$ uniquely minimizes
$\frac12a(v,v)-\operatorname{Re}F(v)$.
\end{theorem}
\begin{proof}
Write $a(v,u)=\faip{v}{Au}$ and $F(v)=\faip{v}{f}$.
Coercivity implies $\|Au\|\geq\alpha\|u\|$, so $A$ is injective with
closed range. If $z\perp\faRan A$, then $A^*z=0$ and
$\operatorname{Re}\faip{z}{Az}=0$; coercivity gives $z=0$.
The range is therefore also dense, hence all of $H$. The lower bound
gives the estimate. If $a$ is Hermitian, insert $v=u+w$ into the energy:
its difference from the energy at $u$ is $a(w,w)/2$, because the cross
term is $\operatorname{Re}a(w,u)=\operatorname{Re}F(w)$.
\end{proof}

\begin{theorem}[Stampacchia inequality]\label{fa2:p2:stampacchia}
Under the same coercivity and boundedness assumptions, for a nonempty closed
convex $C\subset H$ and $F(v)=\faip{v}{f}$ there is a unique $u\in C$ with
\[
 \operatorname{Re}a(v-u,u)\geq\operatorname{Re}F(v-u)\quad(v\in C).
\]
If $a$ is Hermitian this is equivalent to minimizing the preceding energy
over $C$. For two right-hand sides $f,g$ the solutions satisfy
$\|u_f-u_g\|\leq\|f-g\|/\alpha$.
\end{theorem}
\begin{proof}
The projection characterization says that the displayed inequality is
equivalent to $u=P_C(u-\rho(Au-f))$ for any $\rho>0$.
If $H=\{0\}$ the assertion is immediate. Otherwise coercivity implies
$A\ne0$, so the following choice of $\rho$ is legitimate.
For $w=x-y$,
$\|(I-\rho A)w\|^2\leq(1-2\rho\alpha+\rho^2\|A\|^2)\|w\|^2$.
Choose $0<\rho<2\alpha/\|A\|^2$. Nonexpansion of $P_C$ makes this a
strict contraction of the complete metric space $C$ into itself.
The contraction theorem gives existence and uniqueness. Applying the
two inequalities to the two solutions and adding gives
$\alpha\|u_f-u_g\|^2\leq\operatorname{Re}\faip{u_f-u_g}{f-g}$,
proving the estimate. In the Hermitian case the energy difference is
$\operatorname{Re}(a(v-u,u)-F(v-u))+a(v-u,v-u)/2$;
this proves sufficiency. Necessity follows by taking the right derivative
along $u+t(v-u)$. Its sign is nonnegative at a minimum.
\end{proof}

\begin{proposition}[The parallelogram characterization]\label{fa2:p2:parallelogram}
A real or complex normed space is induced by an inner product iff
$\|x+y\|^2+\|x-y\|^2=2\|x\|^2+2\|y\|^2$.
\end{proposition}
\begin{proof}
Necessity follows by expansion. For sufficiency put
$b(x,y)=(\|x+y\|^2-\|x-y\|^2)/4$ on the underlying real space.
The parallelogram identity applied to $x+y,z$ and $x-y,z$ gives
$b(x+z,y)+b(x-z,y)=2b(x,y)$.
Setting $z=x$ gives $b(2x,y)=2b(x,y)$; replacing $x,z$ by
$(u+v)/2,(u-v)/2$ then gives additivity. Integer and rational homogeneity
follow, and continuity of the norm gives real homogeneity. Symmetry and
$b(x,x)=\|x\|^2$ give a real inner product. In the complex case define
$\faip{x}{y}=b(x,y)+i b(ix,y)$. The identity $b(ix,iy)=b(x,y)$ and
$b(ix,y)=-b(x,iy)$ verifies conjugate linearity in the first variable,
linearity in the second, conjugate symmetry and the required square norm.
\end{proof}

\begin{remark}\label{fa2:p2:geometry-open}
The norm-one two-dimensional projection characterization in dimension at
least three (Kakutani), and the nonexpansive radial retraction
characterization of de Figueiredo--Karlovitz, are separate obligations in
\cite{B01}, Theorems 5.13--5.14. The preceding polarization proof does
not prove them. Their independent geometric proofs are
\ref{fa2:p2:kakutani} and \ref{fa2:p2:radial-hilbert}.
\end{remark}

\begin{theorem}[Minty--Browder surjectivity]\label{fa2:p2:minty-browder}
Let $X$ be a real reflexive Banach space and $A:X\to X'$ be bounded on
bounded sets, hemicontinuous, and monotone:
$\langle Ax-Ay,x-y\rangle\geq0$.
Hemicontinuity means that $t\mapsto\langle A(x+ty),z\rangle$ is continuous
for all $x,y,z$. If
$\langle Ax,x\rangle/\|x\|\to+\infty$ as $\|x\|\to\infty$,
then $A$ is onto. If it is strictly monotone, each equation $Ax=f$ has
exactly one solution. In particular this applies to the continuous,
strictly monotone coercive variant in \cite{B01}, Theorem 5.16.
\end{theorem}
\begin{proof}
Fix $f$. Choose $R$ with $\langle Ax-f,x\rangle>0$ when $\|x\|=R$.
For each finite-dimensional subspace $E$ of $X$, the restricted map
$A_E:E\to E'$ is continuous. Indeed on a bounded set it is bounded;
if $x_n\to x$, any convergent subsequence $A_Ex_n\to g$ satisfies
$\langle g-A_Ey,x-y\rangle\geq0$. Put $y=x+tz$, divide by $t$ with
both signs, and use hemicontinuity to get $g=A_Ex$. Finite-dimensional
compactness now proves continuity of the whole sequence.
There is $x_E\in E$ with $\|x_E\|\leq R$ and
$\langle Ax_E-f,v\rangle=0$ for all $v\in E$.
To see existence, choose a Euclidean inner product on $E$. If the map
has no zero in the norm ball, the continuous map
$x\mapsto-R(A_E x-f_E)/\|A_E x-f_E\|$ maps that ball into itself.
Brouwer gives a boundary fixed point; its pairing with the residual is
negative, contrary to the outward inequality. A change of the arbitrary
Euclidean identification does not affect this contradiction.
Direct finite-dimensional spaces by inclusion. Reflexivity makes
$R B_X$ weakly compact, so a subnet $x_E$ converges weakly to some $x$.
For fixed $y\in X$ and every sufficiently large $E$ containing $y$,
monotonicity and the Galerkin equation give
$\langle Ay-f,y-x_E\rangle\geq0$. Passing to the subnet limit yields
$\langle Ay-f,y-x\rangle\geq0$ for every $y$.
Set $y=x+tz$ and use both signs and hemicontinuity; then
$\langle Ax-f,z\rangle=0$ for all $z$, so $Ax=f$.
Strict monotonicity applied to two solutions proves uniqueness.
\end{proof}

% END INPUT


% BEGIN parts/part2-prehilbert-variants.tex
\chapter{Pre-Hilbert and domain-sensitive variants in the sources}
\label{fa2:p2:prehilbert-variants-chapter}
Completeness of an ambient Hilbert space is stronger than completeness
of a particular set onto which we project. These source variants
retain that distinction. Inner products remain conjugate linear in
the first variable and linear in the second.

\begin{proposition}[Completion preserving the inner product]
\label{fa2:p2:prehilbert-completion}
Every real or complex inner-product space $E$ has a Hilbert
completion $H$ with a linear inner-product-preserving dense
embedding. It is unique up to the unique unitary intertwining
those embeddings. This is \cite[Theorem3.2-3]{B03}.
\end{proposition}
\begin{proof}
Take the normed completion constructed in Part I. If $x_n,y_n$
are Cauchy representatives of two of its vectors, set
$\langle[x_n],[y_n]\rangle=\lim_n\langle x_n,y_n\rangle$.
This scalar limit exists since
\[
 |\langle x_n,y_n\rangle-\langle x_m,y_m\rangle|
 \le\|x_n-x_m\|\|y_n\|+\|x_m\|\|y_n-y_m\|.
\]
The Cauchy sequences are bounded; the same inequality proves
independence of the representatives. Sesquilinearity and
conjugate symmetry pass to limits, and its square norm is
the already constructed completion norm squared. Definiteness
and completeness follow. The universal property of the normed
completion supplies the unique surjective linear isometry
between two such completions. Continuity of their inner
products, or polarization of their equal norms, makes this
isometry unitary.
\end{proof}

\begin{theorem}[Projection only requires the target to be complete]
\label{fa2:p2:prehilbert-projection}
Let $E$ be an inner-product space, with no ambient completeness
assumption, and let $C\subset E$ be nonempty, convex and complete
in the induced metric. Every $x\in E$ has a unique nearest
point $p\in C$, characterized by
\[
             \operatorname{Re}\langle x-p,z-p\rangle\le0
                         \quad(z\in C).
\]
The projection is nonexpansive. For a complete linear subspace
$Y$, this gives $E=Y\oplus Y^\perp$, with bounded orthogonal
projection onto $Y$. These are the exact weaker-ambient
versions of \cite[Theorem3.3-1 and Lemma3.3-2]{B03} and
\cite[Theorem1.25]{B06}.
\end{theorem}
\begin{proof}
Let $d=\inf_{z\in C}\|x-z\|$ and choose a minimizing sequence
$z_n$. The parallelogram identity and its midpoint in $C$ give
\[
 \|z_n-z_m\|^2\le2\|x-z_n\|^2+2\|x-z_m\|^2-4d^2\to0.
\]
Completeness of $C$, rather than of $E$, gives its limit $p$
in $C$, and continuity of the norm makes it a minimizer.
The same midpoint identity makes two minimizers identical.
For $0<t\le1$, compare $p$ with $p+t(z-p)$, expand the
square and divide by $t$; the limit as $t\downarrow0$ gives
the stated inequality. Conversely that inequality and expansion
of $\|x-z\|^2$ prove minimality.
For the projections $p,q$ of $x,y$, use $z=q$ in the first
inequality and $z=p$ in the second. Their sum gives
$\|p-q\|^2\le\operatorname{Re}\langle x-y,p-q\rangle
\le\|x-y\|\|p-q\|$, proving nonexpansiveness.
If $C=Y$ is linear, testing displacements $tv$ for both signs
and also $tiv$ over the complex field makes $x-p$ orthogonal
to $Y$. The resulting decomposition is unique because
$Y\cap Y^\perp=0$. Uniqueness proves linearity of its
coordinate projection, and Pythagoras gives its norm bound.
\end{proof}

\begin{proposition}[Orthonormal expansions without ambient completeness]
\label{fa2:p2:prehilbert-parseval}
For any orthonormal family $(e_i)_{i\in I}$ in an inner-product
space, finite-subset partial sums
$p_Fx=\sum_{i\in F}\langle e_i,x\rangle e_i$ satisfy
\[
 \|x-p_Fx\|^2=\|x\|^2-\sum_{i\in F}|\langle e_i,x\rangle|^2.
\]
Bessel's inequality and countable support of the coefficients
hold without completeness. For any particular $x$, the
following are equivalent: $x$ belongs to the closed linear
span of the family; $p_Fx\to x$; and Parseval equality holds
at $x$. Thus density of that span is equivalent to Parseval
for every $x$ in a pre-Hilbert space too.
Density implies that the orthogonal annihilator is zero.
The converse is valid in Hilbert spaces and can fail in
incomplete inner-product spaces.
\end{proposition}
\begin{proof}
Orthogonality gives the identity by finite expansion. Taking
the supremum proves Bessel; for every positive integer $m$
only finitely many coefficients can have modulus at least
$1/m$, which proves countable support. The identity gives
the equivalence of Parseval and convergence to $x$ directly.
If $x$ is in the closed span, approximate it by a finite
combination $z$ supported on $F_0$. Orthogonal projection
onto any finite span containing that support is best
approximation there, so $\|x-p_Fx\|\le\|x-z\|$ for
$F\supset F_0$. This proves convergence. Its converse
follows since all $p_Fx$ lie in the span.
An element orthogonal to a dense span is orthogonal to itself,
hence zero. In a Hilbert space, project any vector onto the
closed span; its residual is in the annihilator, proving
the converse there.

Here is an explicit failure in an incomplete space. In
$\mathbb K\oplus\ell^2$ put
$E=\{(\sum_j a_j,a):a\in c_{00}\}$ and
$M=\{(0,a):a\in c_{00},\ \sum_j a_j=0\}$.
The vectors with first coordinate one and $N$ tail coordinates
equal to $1/N$ converge in the ambient space to $(1,0)\notin E$,
so this inner-product space is incomplete.
The latter is a proper closed subspace of $E$, since it
is the kernel of the continuous first-coordinate functional.
Its closure in the ambient Hilbert space is
$\{0\}\oplus\ell^2$: approximate each tail unit vector by
itself minus the average of $N$ other tail unit vectors,
whose error norm is $1/\sqrt N$. Hence its ambient orthogonal
complement is $\mathbb K\oplus\{0\}$, which meets $E$
only at zero. So $M^\perp$ in $E$ is zero although $M$
is not dense in $E$. If an orthonormal-family example is
desired, apply the finite Gram--Schmidt recursion to
$(0,e_j-e_{j+1})$, $j\ge1$. These independent vectors span
$M$, and the recursion preserves every finite initial span.
Its orthonormal sequence therefore has the same closed
span and annihilator. This separates the two completeness
claims in \cite[Theorems3.6-2--3]{B03} and
\cite[Theorems1.26--1.28]{B06}.
\end{proof}

\begin{proposition}[Norm-total dual tests and a proper bounded domain]
\label{fa2:p2:source-convergence-domain-tests}
In any normed space $X$, a bounded sequence $x_n$ converges
weakly to a specified $x$ if and only if $f(x_n)\to f(x)$
for every $f$ in a set whose norm-closed linear span is $X'$.
Merely separating points is insufficient.
For a bounded linear $T:D\subset X\to Y$ between normed
spaces, closedness of $D$ implies closedness of its graph.
Conversely, a closed graph implies closedness of $D$ if
$Y$ is complete. Neither direction assumes completeness
of $X$. These are \cite[Lemmas4.8-7 and4.13-5]{B03}.
If $X,Y$ are Banach and $D$ is closed, a closed linear
operator on $D$ is bounded in the induced $X$ norm, the
proper-domain version of \cite[Theorem4.13-2]{B03}.
\end{proposition}
\begin{proof}
Tests on the given set extend by linearity to its span.
For general $g\in X'$, approximate it in operator norm
by such an $f$. The error in $g(x_n-x)$ is at most
$\|g-f\|(\sup_n\|x_n\|+\|x\|)$, uniformly in $n$.
First choose $f$, then let $n\to\infty$. The reverse
direction and boundedness of a weakly convergent sequence
were already proved. On $\ell^1$, the coordinate tests
separate points and $e_n$ is bounded and tends to zero
under each such test, yet the continuous sum functional
has value one. This disproves the weakened test condition.

If $D$ is closed and $x_n\to x$, $Tx_n\to y$, then
$x\in D$ and boundedness gives $y=Tx$. Conversely if
$Y$ is complete, for $x\in\overline D$ choose $x_n\in D$
converging to $x$. Boundedness makes $Tx_n$ Cauchy; its
limit in $Y$ and closedness of the graph put $x$ in $D$.
The target completeness cannot be omitted: take $X=\ell^2$,
$D=c_{00}$, and $Y=c_{00}$ with its inherited $\ell^2$
norm. The identity $D\to Y$ is bounded and has closed
graph in $X\times Y$, since any two limiting coordinates
must be the same vector of $Y$. Its domain is not closed
in $X$. For the final assertion, $D$ is itself Banach in
the induced norm, so apply the closed graph theorem to
$T:D\to Y$. Its graph is the same closed subset, now
viewed in $D\times Y$.
\end{proof}

\begin{proposition}[Bounded forms on two different Hilbert spaces]
\label{fa2:p2:bounded-form-representation}
If $H_1,H_2$ are real or complex Hilbert spaces and
$a:H_1\times H_2\to\mathbb K$ is bounded, conjugate linear
in its first and linear in its second argument, there is
a unique bounded linear $S:H_1\to H_2$ with
$a(x,y)=\langle Sx,y\rangle$, and $\|S\|=\|a\|$.
There is also the equivalent representation
$a(x,y)=\langle x,S^*y\rangle$.
This is \cite[Theorem3.8-4]{B03}; no coercivity or symmetry
is imposed and the two Hilbert spaces need not coincide.
\end{proposition}
\begin{proof}
For fixed $x$, Riesz representation of the bounded linear
functional $y\mapsto a(x,y)$ gives a unique vector $Sx$.
Its norm is that functional's norm, at most $\|a\|\|x\|$.
Conjugate linearity in $x$ and in the first inner-product
slot imply $S(\lambda x+\mu z)=\lambda Sx+\mu Sz$ by
uniqueness of the representing vector. Thus $S$ is bounded
and $\|S\|\le\|a\|$. Cauchy--Schwarz gives the reverse
bound for the form norm. The adjoint identity proves the
second representation. Uniqueness follows because vectors
whose inner products against every $y$ agree are equal.
All steps include a zero Hilbert space.
\end{proof}

\begin{proposition}[Elementary adjoint tests with the scalar qualification]
\label{fa2:p2:elementary-adjoint-tests}
For a bounded operator on a complex Hilbert space,
$T=T^*$ if and only if $\langle x,Tx\rangle$ is real
for every $x$. Over a real Hilbert space, real quadratic
values do not imply self-adjointness. For real or complex
Hilbert spaces, the product of two bounded self-adjoint
operators is self-adjoint exactly when they commute.
The self-adjoint operators form a norm-closed set.
These are \cite[Theorems3.10-3--5]{B03}, without a normality
assumption on the first test.
\end{proposition}
\begin{proof}
Self-adjointness makes $\langle x,Tx\rangle$ equal to its
conjugate. Conversely, if all those values are real, put
$A=T-T^*$ and $b(x,y)=\langle x,Ay\rangle$. Then
$b(z,z)=0$ for every $z$. For any complex sesquilinear
form, not only a Hermitian one, expansion gives
\[
 4b(x,y)=\sum_{k=0}^3i^{-k}b(x+i^ky,x+i^ky).
\]
Indeed the constant diagonal terms and the $b(y,x)$
terms cancel, leaving four copies of $b(x,y)$.
Thus $b=0$ and $A=0$. On $\mathbb R^2$ the nonzero
rotation $J(s,t)=(-t,s)$ has every quadratic value zero
but $J^*=-J$, disproving the real converse.
For self-adjoint $S,T$, $(ST)^*=TS$, proving the product
criterion. Finally $\|T_n^*-T^*\|=\|T_n-T\|\to0$,
so $T_n=T_n^*$ and $T_n\to T$ force $T=T^*$.
\end{proof}

\begin{theorem}[The absolute-diagonal Lax--Milgram variant]
\label{fa2:p2:absolute-lax-milgram}
Let $a$ be a bounded sesquilinear form on a real or complex
Hilbert space, conjugate linear in its first and linear in
its second argument. Suppose that, for some $b>0$,
\[
                   |a(v,v)|\ge b\|v\|^2\quad(v\in H).
\]
For every continuous conjugate-linear $F$ there is a unique
$u\in H$ with $a(v,u)=F(v)$ for all $v$, and
$\|u\|\le\|F\|/b$. The solution operator on the continuous
antidual is linear and has norm at most $1/b$.
This is the modulus condition actually stated in
\cite[Chapter6, Theorem6, condition(18)]{B08}, after converting
its opposite inner-product convention. Energy minimization
and Stampacchia retain their separately stated real-part
coercivity hypotheses.
\end{theorem}
\begin{proof}
Represent $a(v,u)=\langle v,Au\rangle$ by the bounded-form
representation above, using the adjoint of its first-slot
operator. The assumption and Cauchy--Schwarz give
$\|Av\|\ge b\|v\|$, so $A$ is injective and has closed
range. If $w\perp\operatorname{Ran}A$, then
$\langle w,Aw\rangle=0$, forcing $w=0$ by the same
modulus hypothesis. Thus the range is dense and closed,
hence all of $H$. Riesz representation of the antidual
writes $F(v)=\langle v,f\rangle$ with $\|f\|=\|F\|$.
The unique solution is $u=A^{-1}f$; the lower bound gives
the norm estimate. Additivity and scalar linearity follow
from uniqueness of the equation, including the zero space.
\end{proof}

\begin{corollary}[Compact self-adjoint operators before completion]
\label{fa2:p2:prehilbert-compact-spectral}
Let $E$ be a real or complex inner-product space and let
$T:E\to E$ be bounded and self-adjoint. Suppose that the
image of its unit ball is relatively compact in $E$:
its closure in $E$ is compact, rather than merely totally
bounded. Then there is a finite or countable orthonormal
family $(e_j)$ in $E$ and nonzero real numbers $\lambda_j$
such that
\[
 Te_j=\lambda_j e_j,\qquad
 Tx=\sum_j\lambda_j e_j\langle e_j,x\rangle\quad(x\in E).
\]
Every nonzero eigenspace has finite dimension; if there
are infinitely many listed eigenvalues, they tend to zero.
The displayed series converges in the norm of $E$ to $Tx$.
The family is empty when $T=0$. This conclusion, as in
\cite[Appendix A, Theorem A.2.28]{B07}, does not assert that
the incomplete space $E$ has an orthonormal basis.
\end{corollary}
\begin{proof}
Let $H$ be the Hilbert completion of $E$. Boundedness gives
a unique continuous extension $\overline T:H\to H$, and
the self-adjoint identity extends from pairs of vectors
in $E$ to pairs in $H$ by continuity. Write $K$ for the
closure in $E$ of $T$ applied to the unit ball of $E$.
Its image in $H$ is compact, hence closed. The unit ball
of $E$ is dense in the unit ball of $H$: approximate in
$E$ and, if necessary, divide each approximant by the
maximum of one and its norm. Thus
$\overline T(B_H)\subseteq K$, so $\overline T$ is compact.

In fact $\overline T(H)\subseteq E$. For $x\ne0$ choose
$x_n\in E$ with $x_n\to x$ and $\|x_n\|\le2\|x\|$.
The sequence $Tx_n$ has a subsequence converging in $E$
by the compactness of $2\|x\|K$. It also converges in
$H$ to $\overline T x$, so that subsequential limit equals
$\overline T x$ and belongs to $E$. The assertion for
$x=0$ is immediate.

For a complex space apply
Theorem~\ref{fa2:p2:compact-spectral} to $\overline T$.
For a real space first complexify $H$ and $\overline T$.
The complexified operator is compact because the real
and imaginary parts of the image of a bounded sequence
admit successively convergent subsequences; the adjoint
identity is preserved. Its eigenvalues are real, and
$u+iv$ is a $\lambda$-eigenvector exactly when both
$u$ and $v$ are real $\lambda$-eigenvectors. Thus each
complex eigenspace is the complexification of the real
one. Choose a real orthonormal basis there and restrict
the spectral expansion to real vectors. This gives the
same theorem for real self-adjoint operators.
Each of its nonzero eigenvectors
belongs to $E$, because
$e=\lambda^{-1}\overline T e$ when $\lambda\ne0$.
Choose orthonormal bases in these finite-dimensional
eigenspaces. The Hilbert-space theorem gives the stated
countability, eigenvalue properties, and spectral series.
For $x\in E$ every partial sum lies in $E$, and its
Hilbert-space limit is $\overline T x=Tx\in E$.
Since the norm on $E$ is the restricted Hilbert norm,
this is precisely convergence in $E$. No projection of
an arbitrary vector onto the full kernel, and no basis
of the ambient incomplete space, has been asserted.
\end{proof}

\begin{proposition}[Gram determinants and exact least-square distances]
\label{fa2:p2:gram-distance}
For vectors $y_1,\ldots,y_n$ of an inner-product space,
the Gram matrix $G=(\langle y_i,y_j\rangle)$ is positive
semidefinite and is positive definite exactly when they
are linearly independent. In that case, for
$Y=\operatorname{span}\{y_1,\ldots,y_n\}$ and
$b=(\langle y_i,x\rangle)_i$, the nearest point is
$p=\sum_j(G^{-1}b)_j y_j$ and
\[
 \operatorname{dist}(x,Y)^2
 =\|x\|^2-b^*G^{-1}b
 =\frac{\det\begin{pmatrix}G&b\\b^*&\|x\|^2\end{pmatrix}}
        {\det G}.
\]
For a dependent list $\det G=0$ and that quotient is not
defined; use any independent sublist spanning $Y$ instead.
For the empty list the empty determinant is one and the
distance is $\|x\|$. No completeness of the ambient space
is needed. This is \cite[Theorems6.5-1--2]{B03} with the
degenerate case made explicit.
\end{proposition}
\begin{proof}
For every coefficient vector $c$, $c^*Gc=\|\sum c_jy_j\|^2$.
This proves the positivity and independence criterion.
The finite-dimensional subspace $Y$ is complete, so the
pre-Hilbert projection theorem applies. Its orthogonality
equations are $Gc=b$, giving the coefficient formula
and, by expansion, the displayed squared distance.
Multiply the augmented matrix on the left by the block
triangular matrix
$\begin{pmatrix}I&0\\-b^*G^{-1}&1\end{pmatrix}$,
whose determinant is one. The lower right entry becomes
$\|x\|^2-b^*G^{-1}b$ and the lower left becomes zero,
proving the determinant ratio. Dependence and the empty
case are exactly the alternatives stated above, with no
division by a zero Gram determinant.
\end{proof}

% END INPUT


% BEGIN parts/part2-duality-regularization.tex
\chapter{Subgradients, duality maps, regularization, and asymptotic centers}
\label{fa2:p2:duality-regularization-chapter}
These are the Banach-space variants in \cite[Problems 2,13--15]{B01}.
The space $X$ is real throughout this chapter. Complex spaces can be
treated on their underlying real spaces when convex order is intended.
For a proper convex function $\phi$, put
$\partial\phi(x)=\{f:\phi(y)\ge\phi(x)+f(y-x)\ (y\in X)\}$
when $\phi(x)$ is finite, and put $\partial\phi(x)=\varnothing$
otherwise.

\begin{theorem}[Subgradient calculus with its qualification]
\label{fa2:p2:subgradient-calculus}
Subgradients are closed convex sets, the graph of $\partial\phi$
is monotone, and
\[
 f\in\partial\phi(x)\quad\Longleftrightarrow\quad
             \phi(x)+\phi^*(f)=f(x).
\]
For $q(x)=\|x\|^2/2$, $\partial q(x)$ is the normalized duality set
\[
 Jx=\{f\in X':\|f\|=\|x\|,\ f(x)=\|x\|^2\}.
\]
The subgradient of the norm at $x\ne0$ consists of its norm-one
norming functionals, and at zero is the closed dual ball. The
subgradient of $I_C$ is the normal cone
$N_C(x)=\{f:f(y-x)\le0\ (y\in C)\}$ at $x\in C$.
A differentiable convex function has just its derivative as subgradient.

If $\phi$ is bounded above on $\overline B(x,R)$, it is locally
Lipschitz on $B(x,R)$ and has a subgradient there. At the center,
if the upper bound is $M$, every subgradient satisfies
$\|f\|\le(M-\phi(x))/R$. Its directional derivative at the center is
\[
 \phi'(x;v)=\max_{f\in\partial\phi(x)}f(v).
\]
For a lower semicontinuous convex function on a Banach space,
$\operatorname{Int}D(\partial\phi)=\operatorname{Int}D(\phi)$.
If $\phi$ is continuous at one point of $D(\phi)\cap D(\psi)$,
then at every point
\[
 \partial(\phi+\psi)=\partial\phi+\partial\psi.
\]
Without this qualification only inclusion from right to left is assured.
\end{theorem}
\begin{proof}
The subgradient inequalities are closed affine inequalities in $f$.
Adding those at two graph points gives monotonicity. Taking the
supremum over $y$ gives exactly the Fenchel equality. The conjugate
of $q$ is $q^*(f)=\|f\|^2/2$: the upper bound is
$f(y)-\|y\|^2/2\le\|f\|^2/2$, and unit vectors on which $f$
approaches its norm, scaled by $\|f\|$, give the reverse bound.
Equality in $f(x)\le\|f\|\|x\|\le(\|f\|^2+\|x\|^2)/2$
proves the formula for $J$. Testing the norm subgradient along
scalar multiples of $x$ gives its value at $x$; testing arbitrarily
large multiples of any vector gives its norm bound. These conditions
also suffice by the norm inequality. Indicator functions give the
normal cone directly. For a differentiable convex function the
one-dimensional secant slopes give its global supporting derivative;
conversely every subgradient is squeezed by the two directional
derivatives and must equal that derivative.

The upper bound on a ball gives the lower bound
$2\phi(x)-M$ there by reflection in its center. On a smaller ball
$\overline B(x,r)$, extend the segment from $u$ toward $v$ a
distance $R-r$ beyond $u$. Convexity and the preceding upper and
lower bounds give
$|\phi(u)-\phi(v)|\le2(M-\phi(x))\|u-v\|/(R-r)$ when
$\|u-v\|\le R-r$. Divide a longer segment into shorter pieces
to obtain the same estimate for all $u,v$ in that smaller ball.
Thus the function is locally Lipschitz. At an interior point the
directional limit $p(v)=\lim_{t\downarrow0}
(\phi(x+tv)-\phi(x))/t$ is finite, continuous and sublinear.
Existence follows from increasing secant slopes; the Lipschitz
bound controls it, and convexity at the midpoint of $x+2tu$ and
$x+2tv$ proves subadditivity. Hahn--Banach supplies $f\le p$,
which is a subgradient because $p(v)\le\phi(x+v)-\phi(x)$.
For any specified $v$, first prescribe $f(v)=p(v)$ on its line;
$-p(v)\le p(-v)$ makes this domination valid for negative scalars
too. Extension proves the claimed maximum formula. The norm bound
at the center follows by testing $x\pm Rv$ for unit vectors $v$.

For the interior-domain assertion, place a closed ball inside the
finite domain. Its intersections with the closed sublevels cover
that complete ball. Baire supplies a smaller open ball on which
$\phi$ is bounded above. This local upper bound propagates to
any specified interior-domain point $z$: if its center is $u\ne z$,
choose $y=z+\eta(z-u)$ still in the finite domain, with $\eta>0$.
Then $z=(y+\eta u)/(1+\eta)$, and perturbations of $z$ are the
same convex combinations of $y$ and points in the bounded ball
about $u$. Convexity bounds $\phi$ above near $z$. The preceding
argument gives a subgradient near every such point. The reverse
interior inclusion follows from $D(\partial\phi)\subset D(\phi)$.

Inclusion in the sum rule follows by adding inequalities. For the
reverse let $f\in\partial(\phi+\psi)(x)$. The function
$\phi(y)-f(y)+\psi(y)$ attains its finite minimum at $x$.
Fenchel--Rockafellar duality, with the stated continuity
qualification, has an attaining dual functional. Equality of
the sum of its two Fenchel inequalities forces equality in each
at $x$. They yield $g\in\partial\phi(x)$ and
$f-g\in\partial\psi(x)$, as required. A strict-inclusion example
is $\phi(t)=-\sqrt t$ for $t\ge0$, infinite otherwise, and
$\psi=I_{\{0\}}$ on $\mathbb R$. Here
$\partial\phi(0)=\varnothing$, whereas $\phi+\psi=I_{\{0\}}$
has every scalar as subgradient at zero.
\end{proof}

\begin{proposition}[Global growth and a Slater normal cone]
\label{fa2:p2:subgradient-slater}
If a finite convex function satisfies $\phi(x)\le k\|x\|+C$
on $X$, it is globally $k$-Lipschitz and
$D(\phi^*)\subset kB_{X'}$. If $\phi$ is continuous convex,
$C_0=\{\phi\le0\}$, and $\phi(z)<0$ somewhere, then at $x\in C_0$
\[
 N_{C_0}(x)=\{\lambda f:f\in\partial\phi(x),\ \lambda\ge0,
                                  \lambda\phi(x)=0\}.
\]
\end{proposition}
\begin{proof}
For $t\ge1$, convexity bounds $\phi(y)-\phi(x)$ by
$(\phi(x+t(y-x))-\phi(x))/t$. Let $t\to\infty$ in the growth
bound, then exchange $x,y$. A functional of norm larger than $k$
has $f(v)>k\|v\|$ in some direction; scaling that direction makes
$f(tv)-\phi(tv)\to\infty$, proving the domain assertion.

Interior points of $C_0$ have zero normal. At a point with
$\phi(x)=0$ and $f\in N_{C_0}(x)$, the open convex subset of
$\mathbb R^2$
\[
 V=\{(r,s):r>\phi(y),\ s>-f(y-x)\text{ for some }y\}
\]
does not contain zero: otherwise a strictly feasible point would
have positive $f(y-x)$. Separate it from zero. Its upward rays
give a nonzero pair $(a,b)$ with $a,b\ge0$ and
$a\phi(y)-bf(y-x)\ge0$ for every $y$. Strict feasibility forces
$b>0$, because $b=0$ would contradict $\phi(z)<0$.
Put $\lambda=a/b$. If $\lambda=0$, the inequality on all of $X$
forces $f=0$; it has the required zero-multiplier representation
because $\partial\phi(x)$ is nonempty by continuity.
If $\lambda>0$, that inequality says
$f/\lambda\in\partial\phi(x)$. Conversely the subgradient
inequality and the multiplier conditions immediately give a normal.
When $\phi(x)<0$, continuity puts $x$ in the interior and the
multiplier must be zero, completing all cases.
\end{proof}

\begin{remark}[Finite-dimensional continuity needs no prior lower semicontinuity]
\label{fa2:p2:finite-convex-continuity}
A finite convex function on an open convex subset of
$\mathbb R^n$ is continuous. Around any point choose a small
closed coordinate cube contained in its domain. Convexity
bounds the function above there by the maximum of its finitely
many vertex values. The local boundedness and Lipschitz
argument in \ref{fa2:p2:subgradient-calculus} applies in its
interior. This supplies \cite[Problem 2.C2]{B01} without
assuming the continuity that is to be proved.
\end{remark}

\begin{theorem}[Duality maps and uniform smoothness]
\label{fa2:p2:duality-smoothness}
Let $X$ be a real Banach space. If $X'$ is strictly convex, $Jx$
is a singleton for every $x$, denoted $Jx$, and
$q(x)=\|x\|^2/2$ is G\^ateaux differentiable with derivative $Jx$.
The map $J:X_{\rm norm}\to X'_{w^*}$ is continuous, and
\[
 (Jx+Jy)(x+y)+(Jx-Jy)(x-y)=2(\|x\|^2+\|y\|^2).
\]
If also $X$ is reflexive and strictly convex, $J$ is bijective
and its inverse is the duality map of $X'$ under the canonical
identification $X''=X$.

The following are equivalent: $X'$ is uniformly convex; and $J$
is everywhere single-valued and uniformly norm continuous on
each bounded subset of $X$. Under these conditions $q$ is
Fr\'echet $C^1$ with derivative $J$.
\end{theorem}
\begin{proof}
Hahn--Banach makes $Jx$ nonempty. If $x\ne0$ and two distinct
norming functionals belonged to it, their midpoint would have
norm strictly less than $\|x\|$ but the same value $\|x\|^2$
at $x$, impossible. At zero $J0=\{0\}$. The directional maximum
formula of \ref{fa2:p2:subgradient-calculus}, applied in both
directions, proves G\^ateaux differentiability. If $x_\alpha\to x$
in norm, the functionals $Jx_\alpha$ are eventually in one fixed
dual ball. Every weak-star cluster point $f$ satisfies
$f(x)=\|x\|^2$ and $\|f\|\le\|x\|$, by the defining
equalities and uniform boundedness of that ball. Thus $f=Jx$.
Alaoglu and uniqueness of cluster points imply convergence of
the entire net to $Jx$: any subnet staying outside a neighborhood
would have a cluster point in its compact closed complement.
This proves continuity, including continuity along each line.
Expanding the displayed identity proves it exactly.

Reflexivity gives norm attainment for every $f\in X'$. A norming
unit vector scaled by $\|f\|$ is its preimage under $J$.
Strict convexity of $X$ gives uniqueness of that norming vector
by the midpoint test. The equations defining $J$ are symmetric
between $x$ and $f$ in the canonical bidual pairing, so this is
also the asserted inverse formula.

For unit $u,v$, the identity and the two norm inequalities give
\[
                 \|Ju+Jv\|+\|u-v\|\ge2.
\]
Suppose $X'$ is uniformly convex and uniform continuity on a
bounded set fails. There would be $x_n,y_n$ in that bounded set
with $\|x_n-y_n\|\to0$ but $\|Jx_n-Jy_n\|$ bounded below.
Pass to a subsequence with $\|x_n\|\to r$; then $\|y_n\|\to r$.
If $r=0$ the latter norm difference tends to zero immediately.
Otherwise normalize $u_n=x_n/\|x_n\|$, $v_n=y_n/\|y_n\|$.
Their distance tends to zero. The displayed inequality makes
$\|Ju_n+Jv_n\|\to2$, so uniform convexity of $X'$ gives
$\|Ju_n-Jv_n\|\to0$. Homogeneity $J(tx)=tJx$ for real $t$
and the convergent radii give the same conclusion before
normalization, a contradiction. The subgradient inequalities give
\[
 0\le q(x+h)-q(x)-Jx(h)
                 \le (J(x+h)-Jx)(h),
\]
proving Fr\'echet differentiability and $C^1$ continuity.

Conversely let $J$ have the stated uniform continuity. On the
ball of radius two let $\omega(t)\downarrow0$ bound
$\|J(x+h)-Jx\|$ for $\|x\|\le1$, $\|h\|\le t\le1$.
Then $q(x+ty)+q(x-ty)\le1+2t\omega(t)$ for $\|x\|,\|y\|\le1$.
For $f,g\in B_{X'}$ with $\|f-g\|\ge\varepsilon$, choose
$\|y\|\le1$ with $(f-g)(y)\ge\varepsilon/2$.
The two Fenchel inequalities for $q$ at $x+ty$ and $x-ty$ yield
\[
 1\ge(f+g)(x)+t(f-g)(y)-q(x+ty)-q(x-ty).
\]
Taking the supremum over the unit ball in $x$ gives
$\|f+g\|\le2-t\varepsilon/2+2t\omega(t)$.
Choose $t>0$ with $\omega(t)\le\varepsilon/8$.
It follows that $\|(f+g)/2\|\le1-t\varepsilon/8$,
which is uniform convexity of the dual ball.
\end{proof}

\begin{theorem}[Differentiable Banach Moreau envelopes]
\label{fa2:p2:banach-moreau}
Suppose $X'$ is uniformly convex, and let $\phi$ be proper,
lower semicontinuous and convex on $X$. For each integer $n\ge1$,
\[
 \phi_n(x)=\inf_y\{\phi(y)+n\|x-y\|^2/2\}
\]
is finite, convex and Fr\'echet $C^1$. Its derivative is the unique
maximizer $S_nx$ of
\[
 f(x)-\phi^*(f)-\|f\|^2/(2n),
\]
and $S_n:X\to X'$ is norm continuous. Moreover
$\phi_n(x)\uparrow\phi(x)$, including infinite values.
No uniqueness of the primal minimizing $y$ is needed.
\end{theorem}
\begin{proof}
A continuous affine minorant of $\phi$ bounds the primal
infimum below by a finite quadratic minimum; a fixed point of
$D(\phi)$ bounds it above. Convexity follows by taking convex
combinations of approximate minimizers. Fenchel--Rockafellar,
using the everywhere continuous quadratic, gives
\[
 \phi_n(x)=\max_f\{f(x)-\phi^*(f)-\|f\|^2/(2n)\}.
\]
One can also see attainment directly. For $y_0\in D(\phi)$,
$\phi^*(f)\ge f(y_0)-\phi(y_0)$, so the negative of the
maximized expression is coercive in $f$. The reflexive space
$X'$ has weakly compact bounded balls. Its quadratic and
conjugate terms are convex weakly lower semicontinuous, so the
minimum is attained. Strict convexity of the dual norm gives
uniqueness. This argument also bounds $S_nx$ uniformly for $x$
in a bounded set.

If $x_j\to x$, any weakly convergent subsequence of the bounded
$S_nx_j$ has limit minimizing the expression at $x$: test
against its unique minimizer and use weak lower semicontinuity.
Consequently $S_nx_j\rightharpoonup S_nx$, and their minimum
values converge by the same comparison. Lower semicontinuity
of $\phi^*$ then gives
$\limsup\|S_nx_j\|^2\le\|S_nx\|^2$; weak lower
semicontinuity gives the reverse bound. Radon--Riesz in the
uniformly convex dual proves norm convergence. Every subsequence
has this conclusion, so $S_n$ is norm continuous.

Comparing the two maximizing functionals at $x$ and $y$ gives
\[
 0\le\phi_n(y)-\phi_n(x)-S_nx(y-x)
                    \le(S_ny-S_nx)(y-x).
\]
This proves the derivative formula and its continuity. Clearly
$\phi_n\le\phi$ and $\phi_n$ increases with $n$. For every
$f\in D(\phi^*)$ the dual formula gives
$\phi_n(x)\ge f(x)-\phi^*(f)-\|f\|^2/(2n)$.
Pass to the limit and then take the supremum in $f$.
Fenchel--Moreau yields the required limit at all $x$.
\end{proof}

\begin{theorem}[Chebyshev and asymptotic centers in uniformly convex spaces]
\label{fa2:p2:asymptotic-center}
Let $X$ be uniformly convex Banach and $C\subset X$ nonempty,
closed and convex. Every nonempty bounded $A\subset C$ has a
unique center $c_A\in C$ minimizing $r_A(x)=\sup_{a\in A}\|x-a\|$.
If $\operatorname{diam}A>0$, its minimum is strictly smaller
than that diameter. For a bounded sequence $a_n\in C$, the function
$r(x)=\limsup_n\|x-a_n\|$ has a unique minimizer $c\in C$.
If $c_n$ is the center of the tail set $\{a_j:j\ge n\}$, then
\[
 c_n\to c\text{ in norm},\qquad
 r_n(c_n)\to r(c),\qquad r(c_n)\to r(c).
\]
Every minimizing sequence for either center problem converges
in norm to its center. If $a_n\to a$ strongly, $c=a$.
In a Hilbert space, weak convergence $a_n\rightharpoonup a$ gives
$r(x)^2=\|x-a\|^2+\limsup_n\|a_n-a\|^2$ and again $c=a$.
\end{theorem}
\begin{proof}
The functions $r_A,r$ are convex and $1$-Lipschitz by the
triangle inequality and passage to the limit of the decreasing
tail suprema. They are coercive since $A$ or the sequence is
bounded. Reflexivity and the direct method give a minimum on
$C$. We verify a useful stronger uniqueness statement. Let
$x_j$ be a minimizing sequence and let the minimum be $\rho$.
For $r$ the triangle inequality along the original tail gives
$\|x_j-x_k\|\le r(x_j)+r(x_k)$; the analogous bound holds for
$r_A$. If $\rho=0$, the minimizing sequence is therefore Cauchy.
If $\rho>0$ and $\|x_j-x_k\|\ge\varepsilon$ for arbitrarily
large $j,k$, fix a small $\eta>0$. For a common sufficiently
late original tail, both distances from $x_j,x_k$ to each $a_i$
are at most $\rho+\eta$. Uniform convexity then bounds their
midpoint distances by
$(\rho+\eta)(1-\delta)$, where one may use the fixed positive
modulus $\delta=\delta_X(\varepsilon/(\rho+1))$ and $\eta<1$.
If the parameter exceeds two no such pair is possible; otherwise
choose also $\eta<\rho\delta/2$. The resulting radius is
strictly less than $\rho$, a contradiction. For $r_A$ use all
points of $A$ instead of a tail. Thus every minimizing sequence
is Cauchy and converges in $C$ to a minimizer by continuity;
alternating two minimizers in such a sequence proves uniqueness.

Write $\rho_n=\min_C r_n$. These numbers decrease. For every
$x\in C$, $\rho_n\le r_n(x)\downarrow r(x)$, whereas
$\rho_n\ge\min_Cr=\rho$. Hence $\rho_n\downarrow\rho$.
Since $\rho\le r(c_n)\le\rho_n$, the tail centers form a
minimizing sequence for $r$. The preceding argument proves
their norm convergence and both radius limits.
To obtain $\min r_A<d=\operatorname{diam}A>0$, choose
$a,b\in A$ with $\|a-b\|>d/2$. For each $y\in A$ the
vectors $(a-y)/d,(b-y)/d$ lie in the unit ball and have
distance greater than $1/2$. Their midpoint therefore has
norm at most $1-\delta_X(1/2)$. The midpoint of $a,b$ belongs
to $C$ and has radius at most $(1-\delta_X(1/2))d<d$.
For strong convergence, $r(x)=\|x-a\|$ and $a\in C$.
For weak Hilbert convergence expand the squared norm about
$a$; its cross term tends to zero. Weak closedness of $C$
puts $a$ in $C$, and the displayed formula identifies its center.
\end{proof}

\begin{corollary}[A bounded orbit suffices for a nonexpansive fixed point]
\label{fa2:p2:bounded-orbit-fixed-point}
For a nonexpansive $T:C\to C$ in the preceding uniformly convex
setting, a bounded orbit $a_n=T^na$ has asymptotic center $c$
with $Tc=c$. Its fixed-point set is closed and convex.
\end{corollary}
\begin{proof}
Nonexpansiveness gives $r(Tx)\le r(x)$ by shifting the orbit
one place. Uniqueness of the center therefore forces $Tc=c$.
Continuity makes the fixed-point set closed. If $p,q$ are fixed
and $z=(1-t)p+tq$, the two inequalities
$\|Tz-p\|\le t\|p-q\|$ and
$\|Tz-q\|\le(1-t)\|p-q\|$ force equality in the triangle
inequality. Strict convexity places $Tz$ at exactly $z$ on that
segment (the case $p=q$ is immediate). This proves convexity
without a weak-continuity assumption on $T$.
\end{proof}

% END INPUT


% BEGIN parts/part2-hilbert-prox.tex
\chapter{Hilbert proximal maps and their parameter limits}
\label{fa2:p2:hilbert-prox-chapter}

In this chapter $H$ is a real Hilbert space, with scalar product
$(\cdot,\cdot)$ and norm $\|\cdot\|$.  Let
$\phi:H\to(-\infty,\infty]$ be proper, convex and lower
semicontinuous, and write $D(\phi)=\{x:\phi(x)<\infty\}$.
All closures below are norm closures.  These are the precise
Hilbert-space assertions of \cite[Problem 26, pp.~474--475]{B01};
the two domain closures in its questions 6 and 8 are essential.
The more general Banach duality-map regularization theorem does
not by itself give the Hilbert primal--dual decomposition here.

\begin{definition}[Proximal map, envelope, and Hilbert subgradient]
\label{fa2:p2:hilbert-prox-definition}
For $\lambda>0$ and $f\in H$, put
\[
 e_\lambda^\phi(f)=\inf_{v\in H}
       \left\{\phi(v)+\frac{\|f-v\|^2}{2\lambda}\right\}.
\]
Its minimizer, proved to exist uniquely below, is denoted
$P_\lambda^\phi f$ or $\operatorname{prox}_{\lambda\phi}f$.
For $u\in D(\phi)$ define
\[
 \partial\phi(u)=\{w\in H:
       \phi(v)\ge\phi(u)+(w,v-u)\text{ for every }v\in H\}.
\]
The Hilbert conjugate is
$\phi^*(w)=\sup_{v\in H}\{(w,v)-\phi(v)\}$.
\end{definition}

\begin{theorem}[Existence, the variational inequality, and firm continuity]
\label{fa2:p2:hilbert-prox}
For every $f\in H$ and $\lambda>0$ the envelope is finite and
has a unique minimizer $u=P_\lambda^\phi f\in D(\phi)$.
Equivalently,
\[
 (f-u,v-u)\le\lambda\bigl(\phi(v)-\phi(u)\bigr)
       \qquad(v\in D(\phi)),
 \qquad \frac{f-u}{\lambda}\in\partial\phi(u).
\]
For $u=P_\lambda^\phi f$ and $z=P_\lambda^\phi g$,
\[
 \|u-z\|^2\le(f-g,u-z),\qquad \|u-z\|\le\|f-g\|.
\]
If $\phi$ is the indicator of a nonempty closed convex set $K$,
then $P_\lambda^\phi=P_K$ for every $\lambda>0$.
\end{theorem}
\begin{proof}
By the continuous-affine-minorant consequence of Fenchel--Moreau
there are $a\in H$ and $b\in\mathbb R$ such that
$\phi(v)\ge(a,v)+b$.  Completing a square gives
\[
 J(v):=\tfrac12\|f-v\|^2+\lambda\phi(v)
 \ge\tfrac12\|v-(f-\lambda a)\|^2
       +\lambda(a,f)+\lambda b-\tfrac12\lambda^2\|a\|^2.
\]
Thus $m=\inf J> -\infty$, while evaluating at a point of
$D(\phi)$ gives $m<\infty$.  The parallelogram identity and
convexity yield, for $v,w\in D(\phi)$,
\[
 J\bigl((v+w)/2\bigr)
 \le\tfrac12J(v)+\tfrac12J(w)-\tfrac18\|v-w\|^2.
\]
A minimizing sequence $(v_n)$ therefore satisfies
$\|v_n-v_k\|^2\le4(J(v_n)+J(v_k)-2m)$ and is Cauchy.
Its norm limit $u$ exists in $H$.  Lower semicontinuity gives
$J(u)\le m$, so $J(u)=m$ and $u\in D(\phi)$.  The same
strict midpoint inequality proves uniqueness.

For $v\in D(\phi)$ and $0<t\le1$, compare $u$ with
$u+t(v-u)$.  Using the upper convexity bound for $\phi$ gives
\[
 0\le J(u+t(v-u))-J(u)
 \le t(u-f,v-u)+\tfrac12t^2\|v-u\|^2
                  +\lambda t(\phi(v)-\phi(u)).
\]
Divide by $t$ and let $t\downarrow0$.  This is the stated
variational inequality and, by definition, the subgradient
condition.  Conversely that inequality and the square expansion
give $J(v)-J(u)\ge\tfrac12\|v-u\|^2$ for every $v$ with
finite energy; competitors of infinite energy are automatic.
Hence it characterizes the minimizer.

Apply the inequality for $u$ with competitor $z$, and the one
for $z$ with competitor $u$.  Adding cancels the two function
values and gives
$\|u-z\|^2\le(f-g,u-z)$.  Cauchy--Schwarz gives the
Lipschitz bound, including the case $u=z$.  For an indicator,
finite-energy competitors are exactly the points of $K$ and
the minimization is exactly the metric projection problem.
\end{proof}

\begin{theorem}[Moreau's primal--dual decomposition and energy identity]
\label{fa2:p2:prox-moreau}
The conjugate $\phi^*$ is proper, convex and lower semicontinuous.
For $u=P_\lambda^\phi f$ and
$w=P_{1/\lambda}^{\phi^*}(f/\lambda)$ one has
\[
 u+\lambda w=f,\qquad
 \phi(u)+\phi^*(w)=(u,w),
\]
and
\[
 e_\lambda^\phi(f)
   +e_{1/\lambda}^{\phi^*}(f/\lambda)
       =\frac{\|f\|^2}{2\lambda}.
\]
In particular, at $\lambda=1$, the minimizers for $\phi$ and
$\phi^*$ at the same datum sum to $f$, and their minimum
energies sum to $\|f\|^2/2$.
\end{theorem}
\begin{proof}
The conjugate is a supremum of continuous affine functions, hence
convex and lower semicontinuous.  The affine minorant used above
gives $\phi^*(a)\le-b<\infty$, and a fixed point in $D(\phi)$
gives a finite lower bound for each $\phi^*(w)$.  Thus it is
proper, so its proximal minimizer also exists by the preceding
theorem.

Put $w=(f-u)/\lambda$.  The variational inequality says
$\phi(v)\ge\phi(u)+(w,v-u)$ for all $v$.  Taking the
supremum in the definition of the conjugate, with equality at
$v=u$, gives $\phi^*(w)=(w,u)-\phi(u)$.  Conversely the
definition of $\phi^*$ gives for every $z\in H$
\[
 \phi^*(z)\ge(z,u)-\phi(u)
             =\phi^*(w)+(u,z-w).
\]
Thus $u\in\partial\phi^*(w)$.  Since
$f/\lambda-w=u/\lambda$, the proximal characterization for
$\phi^*$ and parameter $1/\lambda$ identifies $w$ with the
claimed minimizer.  Finally substitute $f-u=\lambda w$ and
$f/\lambda-w=u/\lambda$ in the two envelopes.  Their sum is
\[
 \phi(u)+\phi^*(w)
       +\tfrac\lambda2\|w\|^2+\tfrac1{2\lambda}\|u\|^2
 =\frac{\|u+\lambda w\|^2}{2\lambda},
\]
which is the asserted identity.  All the displayed function
values are finite, so no subtraction of infinities occurs.
\end{proof}

\begin{theorem}[The small-parameter limit]
\label{fa2:p2:prox-small-parameter}
For every $f\in H$,
\[
 P_\lambda^\phi f\longrightarrow P_{\overline{D(\phi)}}f
       \quad\text{in norm as }\lambda\downarrow0.
\]
The domain need not be closed, and $f$ need not belong to its
closure.
\end{theorem}
\begin{proof}
Write $u_\lambda=P_\lambda^\phi f$ and fix $v_0\in D(\phi)$.
Comparison of the two energies, followed by the affine minorant,
gives
\[
 \|u_\lambda-f\|^2
 \le\|v_0-f\|^2+
       2\lambda\bigl(\phi(v_0)-(a,u_\lambda)-b\bigr).
\]
For $0<\lambda\le1$, writing
$u_\lambda=f+(u_\lambda-f)$ bounds its right side by
$C+2\|a\|\|u_\lambda-f\|$, with $C$ independent of
$\lambda$.  A quadratic inequality therefore bounds
$\|u_\lambda-f\|$, and hence $\|u_\lambda\|$, uniformly.
In particular the affine minorant bounds $\phi(u_\lambda)$
below by a constant independent of these $\lambda$.

Now compare with an arbitrary fixed $v\in D(\phi)$ and let
$\lambda\downarrow0$.  The same energy comparison gives
\[
 \limsup_{\lambda\downarrow0}\|u_\lambda-f\|^2
       \le\|v-f\|^2.
\]
Taking the infimum over such $v$ gives the squared distance
$d^2$ from $f$ to the nonempty closed convex set
$C=\overline{D(\phi)}$.  Let $p=P_Cf$.
Since $u_\lambda\in C$, its projection inequality gives
\[
 \|u_\lambda-p\|^2
       \le\|u_\lambda-f\|^2-\|p-f\|^2.
\]
The right side is nonnegative and has limsup at most zero,
proving the norm convergence directly, with no weak-to-strong
inference left implicit.
\end{proof}

\begin{theorem}[Large parameters and an empty minimizing set]
\label{fa2:p2:prox-large-parameter}
Let $M=\operatorname*{argmin}_H\phi$.
If $M\ne\varnothing$, it is closed and convex and
\[
 P_\lambda^\phi f\longrightarrow P_Mf
       \quad\text{in norm as }\lambda\to\infty.
\]
If $M=\varnothing$, then
$\|P_\lambda^\phi f\|\to\infty$.
\end{theorem}
\begin{proof}
For each fixed $v\in D(\phi)$, the energy comparison gives
\[
 \phi(u_\lambda)\le\phi(v)
              +\frac{\|f-v\|^2}{2\lambda}.
\]
Consequently $\phi(u_\lambda)\to m=\inf_H\phi$, allowing
$m=-\infty$.  If $M$ is nonempty, $m$ is finite and its
sublevel set $M=\{\phi\le m\}$ is closed and convex.
Put $p=P_Mf$.  The variational inequality at $p$ implies
\[
 (f-u_\lambda,p-u_\lambda)
       \le\lambda(m-\phi(u_\lambda))\le0.
\]
It follows both that $\|u_\lambda-f\|\le\|p-f\|$ and that
\[
 \|u_\lambda-p\|^2\le(f-p,u_\lambda-p).
\]
For clarity, the first bound follows by writing the left pairing
as $\|f-u_\lambda\|^2+(f-u_\lambda,p-f)$ and applying
Cauchy--Schwarz.  The second follows by writing
$f-u_\lambda=(f-p)-(u_\lambda-p)$.

Given any sequence $\lambda_n\to\infty$, boundedness and
Hilbert reflexivity give a weakly convergent subsequence with
limit $z$.  Convex lower semicontinuity implies weak lower
semicontinuity, so $\phi(z)\le m$ and $z\in M$.
The first bound and weak lower semicontinuity of the norm give
$\|z-f\|\le\|p-f\|$, whence $z=p$ by uniqueness of the
projection.  Along this subsequence the last displayed pairing
tends to zero, so $u_{\lambda_n}\to p$ in norm.
A sequence contradicting the asserted norm limit would have
such a convergent subsequence, a contradiction.  This proves
the full parameter limit.

If $M$ is empty and the norms do not tend to infinity, choose
$\lambda_n\to\infty$ with $(u_{\lambda_n})$ bounded and pass
to a weakly convergent subsequence.  Weak lower semicontinuity
would give $\phi(z)\le m$.  If $m=-\infty$ this contradicts
properness; if $m$ is finite it makes $z$ a minimizer.  Both
contradictions prove divergence of the norms.
\end{proof}

\begin{corollary}[The scaled limit and the conjugate domain]
\label{fa2:p2:prox-scaled-limit}
For every $f\in H$,
\[
 \frac{P_\lambda^\phi f}{\lambda}
       \longrightarrow- P_{\overline{D(\phi^*)}}0
       \quad\text{in norm as }\lambda\to\infty.
\]
\end{corollary}
\begin{proof}
Put $\tau=1/\lambda$.  Moreau's identity gives
\[
 \frac{P_\lambda^\phi f}{\lambda}
       =\tau f-P_\tau^{\phi^*}(\tau f).
\]
Firm continuity gives
$\|P_\tau^{\phi^*}(\tau f)-P_\tau^{\phi^*}0\|
\le\tau\|f\|$.  The small-parameter theorem for the proper
function $\phi^*$ gives
$P_\tau^{\phi^*}0\to P_{\overline{D(\phi^*)}}0$.
The triangle inequality proves the stated limit.
\end{proof}

\begin{remark}[Complex Hilbert spaces are realified here]
\label{fa2:p2:prox-realification}
For a complex Hilbert space the same results hold on its
underlying real Hilbert space with scalar product
$(x,y)_{\mathbb R}=\operatorname{Re}\langle x,y\rangle$.
Convexity is real convexity, the subgradients are real continuous
functionals identified by this real scalar product, and the
conjugate uses $\operatorname{Re}\langle w,v\rangle$.
This does not identify a nonlinear proximal map of a real-valued
convex function with a complex-linear operator or confuse the real conjugate
with a conjugate-linear Hilbert dual.
\end{remark}

% END INPUT


% BEGIN parts/part2-source-problems.tex
\chapter{Additional result-bearing problems in Banach geometry}
\label{fa2:p2:source-problems-chapter}
All order and monotonicity assertions in this chapter concern real
spaces. The complex versions of the sequence-ball and Clarkson
statements are specified separately.

\begin{lemma}[Extreme sets need not be assumed convex]
\label{fa2:p2:extreme-set-tests}
For a convex set $K$ and $a\in K$, the point $a$ is extreme
if and only if $K\setminus\{a\}$ is convex. Every finite convex
combination with strictly positive weights equal to an extreme
point has all its entries equal to that point.
If $K$ is compact convex in a Hausdorff locally convex space,
every nonempty closed extreme subset $A$ contains an extreme
point of $K$, even if $A$ itself is not assumed convex.
Here an extreme subset means that a proper convex combination
of points of $K$ in that subset has both endpoints in the subset.
\end{lemma}
\begin{proof}
The first equivalence says precisely that two points different
from $a$ cannot have a proper convex combination equal to $a$.
For a finite positive combination group one entry against the
convex combination of the others; extremality forces both to
equal $a$, and induction treats all entries.
Order the nonempty closed extreme subsets of $A$ by reverse
inclusion. Chains have nonempty intersections by compactness
of $K$, and the endpoint property survives intersection.
Choose an inclusion-minimal one by Zorn. If it contained two
distinct points, a continuous real linear functional would
distinguish them. Its maximizing set on that compact subset
would be nonempty, closed, proper and again extreme in $K$:
first the original extreme property puts both endpoints in
the subset, and equality of their functional average at its
maximum puts both in its maximizing set. This contradicts
minimality. The minimal subset is therefore a singleton,
whose point is extreme. This proves also the nonconvex
extreme-set version of \cite[Problem 1.3--7]{B01}.
\end{proof}

\begin{proposition}[Extreme points of the classical balls]
\label{fa2:p2:classical-ball-extremes}
For real scalars, the extreme points of the closed unit balls of
$\ell^\infty$ and $c$ are respectively the sequences with every
coordinate in $\{-1,1\}$, and those such sequences which converge.
The ball of $c_0$ has no extreme point. The extremes of the
$\ell^1$ ball are $\pm e_n$; those of $\ell^p$, $1<p<\infty$,
are all its unit vectors in the sense of norm one. The ball of
$L^1(\mathbb R)$ has no extreme points. For complex scalars replace
the signs by arbitrary modulus-one phases in the first two statements
and in the $\ell^1$ statement. These are the examples in
\cite[Problem 1.9]{B01}.
\end{proposition}
\begin{proof}
A scalar point of the closed disk is extreme exactly when its modulus
is one. Coordinatewise testing gives sufficiency in $\ell^\infty$
and $c$. If a coordinate has modulus less than one, perturb only
that coordinate in opposite directions inside its disk; both
perturbations remain in the given space. Every $c_0$ vector has
such a coordinate. In $\ell^1$ a norm-one vector supported at one
coordinate is extreme by equality in the triangle and norm bounds.
If two coordinates are nonzero, move a sufficiently small amount
of their magnitudes from one to the other, preserving their phases,
to obtain two distinct unit-ball vectors with the original midpoint.
A vector of norm less than one is also nonextreme, by a sufficiently
small coordinate perturbation. Strict convexity of $\ell^p$ gives
the interior-exponent assertion; interior ball points are nonextreme.
For nonzero $f\in L^1(\mathbb R)$ of norm one, choose a measurable
set $A$ with $0<\int_A|f|<1$. Such a set exists since
$t\mapsto\int_{(-\infty,t]}|f|$ is continuous and ranges from
zero to one. Its two normalized restrictions give a nontrivial
convex decomposition in the ball. Interior points are perturbed
by a small nonzero integrable function. This proves all claims.
\end{proof}

\begin{theorem}[Closed generating cones and automatic continuity]
\label{fa2:p2:positive-automatic-continuity}
Let $P$ be a convex cone in a real normed space $X$, and set $F=P-P$.
If every algebraic linear functional nonnegative on $P$ is continuous,
then $F$ is closed and has finite codimension. Conversely, if $X$ is
Banach, $P$ is closed, and $F$ is closed of finite codimension,
every such positive functional is continuous.
In particular a closed generating cone in a Banach space admits
bounded positive decompositions: some $C$ satisfies
$x=y-z$, $y,z\in P$, $\|y\|,\|z\|\le C\|x\|$ for all $x$.
This is \cite[Problem 5]{B01}.
\end{theorem}
\begin{proof}
An algebraic functional vanishing on $F$ is positive on $P$, hence
continuous by the first assumption. Algebraic separation of every
point outside $F$ shows that $F$ is the intersection of their
closed kernels, so is closed. All algebraic functionals on the
normed quotient $X/F$ are then continuous, by quotient factorization.
If that quotient were infinite dimensional, prescribe successively
the values $n$ on independent unit vectors and extend to a Hamel
basis. The resulting unbounded functional is a contradiction.

First prove the converse when $P-P=X$. Put
$K=(P\cap B_X)-(P\cap B_X)$. This set is balanced and convex,
and $X=\bigcup_{n\ge1}nK$. Baire applied to their closures,
followed by taking differences and scaling, gives
$B(0,r)\subset\overline K$ for some $r>0$.
For $\|x\|<r/2$, choose successively $y_j,z_j\in P$ with
$\|y_j\|,\|z_j\|\le2^{-j}$ and residual after $j$ steps
smaller than $r2^{-j-1}$. This is possible by scaling the
closure inclusion at each residual. Completeness and closedness
of the cone make $y=\sum y_j,z=\sum z_j$ belong to $P$;
their norms are at most one and $x=y-z$. Scaling gives the
stated bound, with an enlarged constant if necessary.

Let $f$ now be an algebraic positive functional. It is bounded
on $P\cap B_X$: otherwise choose $u_j\in P$ with
$\|u_j\|\le2^{-j}$ and $f(u_j)\ge1$. The convergent sum $u$
is in $P$, and $u-\sum_{j=1}^Nu_j\in P$ for every $N$.
Positivity would give the impossible inequality $f(u)\ge N$.
The positive decomposition bound now gives a bound for $f$ on
the whole unit ball, proving continuity without having assumed
continuity when applying it to the infinite sum.
Finally for closed finite-codimensional $F$, apply this argument
in the Banach space $F$ and extend the resulting bounded
restriction by Hahn--Banach. The difference from the original
functional factors through the finite-dimensional quotient, hence
is continuous too.
\end{proof}

\begin{remark}[Concrete cones and the completeness qualification]
\label{fa2:p2:positive-cone-examples}
For the nonnegative cone of $C[0,1]$, positive and negative parts
give $P-P=C[0,1]$, so every positive algebraic functional is
continuous; directly $|f(u)|\le\|u\|_\infty f(1)$.
If the cone additionally requires both endpoint values to be
zero, $P-P=\{u:u(0)=u(1)=0\}$, closed of codimension two,
and automatic continuity still follows. In
$X=\{u\in C^1[0,1]:u(0)=u(1)=0\}$ with its Banach $C^1$
norm, every $u$ satisfies $|u(t)|\le C t(1-t)$: use the
derivative bound separately from each endpoint and
$\min(t,1-t)\le2t(1-t)$. Thus
$u=(u+C t(1-t))-C t(1-t)$ is a positive decomposition and
again $P-P=X$. If instead this last space is given the
inherited supremum norm, the derivative functional $u\mapsto
u'(0)$ is positive but unbounded: positive functions vanishing
at zero have nonnegative right derivative, while
$u_n(t)=t(1-t)e^{-nt}$ have derivative one at zero and
supremum norm tending to zero. This explicitly distinguishes
the complete norm from the incomplete one in reading the
examples of \cite[Problem 5.C]{B01}.
\end{remark}

\begin{theorem}[Quadratic estimates for a closed positive operator]
\label{fa2:p2:positive-graph-estimate}
Let $A:D(A)\subset X\to X'$ be a closed real linear operator,
$X$ Banach, with $q(y)=(Ay)(y)\ge0$. The following are equivalent:
for every $x\in D(A)$, $\inf_{y\in D(A)}[q(y)-(Ay)(x)]> -\infty$;
and, for some $k\ge0$,
\[
 |(Ay)(x)|^2\le k(\|x\|+\|Ax\|)^2q(y)
                  \quad(x,y\in D(A)).
\]
If $D(A)=X$, then for any $\alpha>0$ the estimate
$\|Ay\|^2\le\alpha q(y)$ is equivalent to
\[
 q(y)-(Ay)(x)\ge-\alpha\|x\|^2/4\quad(x,y\in X),
\]
and holds if and only if the same estimate holds for the transpose
$A'|_X:X\to X'$ under canonical evaluation. This is the full
quadratic variant of \cite[Problem 6]{B01}.
\end{theorem}
\begin{proof}
The displayed estimate implies the lower bound by minimizing
$s^2-\sqrt{k}(\|x\|+\|Ax\|)s$ over $s\ge0$.
Conversely the function
$h(x)=\sup_y[(Ay)(x)-q(y)]$ is convex, lower semicontinuous
and finite on the Banach graph domain. The local boundedness
result of \ref{fa2:p2:subgradient-calculus} bounds it by $M\ge0$
on some graph ball of radius $R>0$ about zero. Substituting
$ty$ for $y$ gives
$t(Ay)(x)-t^2q(y)\le M$ for all real $t$ and graph-unit-scaled
$x$. A quadratic nonnegative for all $t$ has nonpositive
discriminant; hence $|(Ay)(x)|^2\le4Mq(y)$ on that ball.
Scaling $x$ gives $k=4M/R^2$.

When the domain is all of $X$, the closed graph theorem makes
$A$ bounded. For fixed $y$, the supremum over $x$ of
$(Ay)(x)-\alpha\|x\|^2/4$ is $\|Ay\|^2/\alpha$,
by the dual norm formula and scalar optimization. This proves
the asserted equivalence. Finally
\[
 q(x-y)-(A(x-y))(x)=q(y)-(Ax)(y).
\]
Thus replacing $y$ by $x-y$ turns the lower-bound inequality
for $A$ exactly into the one for $A'|_X$, with the same
$\alpha\|x\|^2/4$. Apply the equivalence again to obtain
the transpose estimate. No symmetry of $A$ was assumed.
\end{proof}

\begin{proposition}[Two further weak-convergence variants]
\label{fa2:p2:schur-cauchy-sphere}
Every weakly Cauchy sequence in $\ell^1$ converges in norm.
In $L^1(0,1)$ there is a weakly null sequence of norm one.
For a monotone map $A:B_X\to X'$ on the closed ball of a real
normed space with nonempty unit sphere,
\[
 A(B_X)\subset\overline{\operatorname{conv}A(S_X)}^{\,w^*}.
\]
These are \cite[Problems 8.5--6 and 9.B5]{B01}.
\end{proposition}
\begin{proof}
If a weakly Cauchy sequence were not norm Cauchy, two subsequences
would have differences bounded away from zero in norm. Their
differences are weakly null, contradicting the proved Schur theorem
\ref{fa2:p2:complete-continuity}. Completeness gives a norm limit.
For the $L^1$ example take the Rademacher functions
$r_n(t)=\operatorname{sign}(\sin(2^n\pi t))$. They have norm
one and form an orthonormal sequence in $L^2(0,1)$: on every
dyadic interval where the lower-frequency one is constant, the
higher-frequency one has zero integral. For $g\in L^\infty$
its $L^2$ coefficients therefore tend to zero by Bessel's
inequality. This is exactly weak nullity in $L^1$.

If $Ax$ were outside the displayed weak-star closed convex hull,
weak-star separation would give $v\in X$ with
$(Ax)(v)>\sup_{z\in S_X}(Az)(v)$. A boundary $x$ is already
in the image set. For an interior $x$, choose $t>0$ with
$z=x+tv\in S_X$ (the separator has $v\ne0$).
Monotonicity gives $t(Az-Ax)(v)\ge0$, contradicting separation.
\end{proof}

\begin{proposition}[Weak-star adjoint closures and elementary maximality]
\label{fa2:p2:weakstar-adjoint-closures}
For a closed densely defined $A:D(A)\subset X\to Y$ between
Banach spaces,
\[
 \overline{\operatorname{Ran}A'}^{w^*}=(\ker A)^\perp,
 \qquad \overline{D(A')}^{w^*}=Y'.
\]
If $X$ is reflexive the first closure may be norm closure;
if $Y$ is reflexive the second may be norm closure.
For a real linear monotone map $B:D(B)\subset X\to X'$,
maximal monotonicity by itself forces the graph to be closed and
the domain to be dense. An everywhere-defined linear monotone map
is maximal monotone and consequently bounded.
\end{proposition}
\begin{proof}
The graph annihilator in $X'\oplus Y'$ is exactly
$\{(-A'y',y'):y'\in D(A')\}$. A vector $x$ annihilating
$\operatorname{Ran}A'$ therefore makes $(x,0)$ annihilate
that graph annihilator, and the double-annihilator theorem
puts $(x,0)$ in the closed original graph. Thus $x\in\ker A$;
the converse is immediate. Similarly a $y$ annihilating
$D(A')$ puts $(0,y)$ in the graph, forcing $y=0$.
Weak-star double annihilators give the formulas. Under the
indicated reflexivity, weak-star is weak on the relevant dual,
and weak and norm closures of linear subspaces agree.

For the maximality assertions, a limit of graph points is
monotonically related to every graph point, so maximality
makes it a graph point. If $f$ annihilates $D(B)$, the pair
$(0,f)$ is monotonically related to the graph because
$(Bu-f)(u)=(Bu)(u)\ge0$. Maximality forces $f=B0=0$,
and Hahn--Banach gives density. Finally, for an everywhere
defined map and a related pair $(x,f)$, substitute
$u=x+tv$ for both signs of $t$ in
$(Bu-f)(u-x)\ge0$. Its linear coefficient must vanish:
$(Bx-f)(v)=0$ for every $v$. Hence $f=Bx$ and the map is
maximal. The closed graph theorem then proves boundedness.
These are the additional closure/domain assertions in
\cite[Problems 9.B3 and 16.A]{B01}.
\end{proof}

\begin{theorem}[The exact Clarkson inequalities]
\label{fa2:p2:clarkson}
Let $1<p<\infty$, $q=p/(p-1)$, and let $\mu$ be a positive
measure for which the usual $L^p$ duality holds, in particular
any sigma-finite measure. For $p\ge2$,
\begin{align*}
 \|f+g\|_p^p+\|f-g\|_p^p
 &\le2(\|f\|_p^q+\|g\|_p^q)^{p/q}\\
 &\le2^{p-1}(\|f\|_p^p+\|g\|_p^p).
\end{align*}
For $1<p\le2$,
\[
 \|f+g\|_p^q+\|f-g\|_p^q
                 \le2(\|f\|_p^p+\|g\|_p^p)^{q/p}.
\]
These inequalities hold for real or complex functions. In
particular $L^p$ is uniformly convex, with the valid modulus
$1-(1-(\varepsilon/2)^r)^{1/r}$ where $r=\max(p,q)$.
These retain the inequalities, in addition to the uniform
convexity consequence, of \cite[Problem 20]{B01}.
\end{theorem}
\begin{proof}
We first prove the scalar two-point inequality, rather than
assuming it. The matrix $H(a,b)=(a+b,a-b)$ has norm one
from $\ell^1_2$ to $\ell^\infty_2$, by the triangle inequality,
and norm $\sqrt2$ on $\ell^2_2$, by direct expansion of the
two squared moduli. The fully proved Riesz--Thorin theorem
\ref{fa2:p4:riesz-thorin}, with parameter $2/p$, gives
$\|H(a,b)\|_{\ell^p_2}\le2^{1/p}\|(a,b)\|_{\ell^q_2}$.
Raising to $p$ gives
$|a+b|^p+|a-b|^p\le2(|a|^q+|b|^q)^{p/q}$.
Apply it pointwise and use Minkowski in $L^{p/q}$, whose
exponent $p-1$ is at least one. The resulting bound is exactly
the first inequality. Convexity of $t^{p/q}$ gives the second.

For $p\le2$, apply the first mixed inequality at exponent
$q\ge2$ to $u,v\in L^q$:
\[
 (\|u+v\|_q^q+\|u-v\|_q^q)^{1/q}
            \le2^{1/q}(\|u\|_q^p+\|v\|_q^p)^{1/p}.
\]
Pair $(f+g,f-g)$ with $(u,v)$ and rewrite its integral as
the pairing of $(f,g)$ with $(u+v,u-v)$. Holder, first in
the measure variable and then for the two coordinates, bounds
the result by $2^{1/q}(\|f\|_p^p+\|g\|_p^p)^{1/p}$ when
$\|u\|_q^p+\|v\|_q^p\le1$. The dual norm of this
two-coordinate product is
$(\|f+g\|_p^q+\|f-g\|_p^q)^{1/q}$, as follows from
$L^p$ duality and the finite Holder equality. Taking the
supremum and raising to $q$ proves the second-exponent formula.
Finally, for $\|f\|_p,\|g\|_p\le1$, divide the applicable
inequality by $2^r$ to obtain
$\|(f+g)/2\|_p^r+\|(f-g)/2\|_p^r\le1$.
The asserted modulus follows immediately when
$\|f-g\|_p\ge\varepsilon$.
\end{proof}

\begin{corollary}[The endpoint-invariance interpolation bridge]
\label{fa2:p2:endpoint-interpolation-bridge}
On a finite measure space let $T:L^1\to L^1$ be bounded and
suppose only that $T(L^\infty)\subset L^\infty$. Then its
restriction to $L^\infty$ is automatically bounded. If the two
endpoint norms are $N_1,N_\infty$, then
\[
 \|Tf\|_p\le N_1^{1/p}N_\infty^{1-1/p}\|f\|_p
                   \quad(1<p<\infty).
\]
In particular this implies the weaker constant two requested
in \cite[Problem 22]{B01}.
\end{corollary}
\begin{proof}
The zero-measure case is immediate. For a sequence $f_n\to f$
in $L^\infty$ with $Tf_n\to g$ in $L^\infty$, finite measure
also gives both convergences in $L^1$. Boundedness of the original
operator forces $g=Tf$ there, hence as an a.e. class. Its
$L^\infty$ graph is closed and the closed graph theorem makes
the restriction bounded. If $T$ was initially real, complexify
it without increasing these two endpoint norms. For $L^1$,
on disjoint simple functions the triangle inequality gives
$\|\sum_jc_jT1_{A_j}\|_1\le N_1\sum_j|c_j|\mu(A_j)$;
density gives the bound on the complex space. For $L^\infty$,
take the supremum over a countable dense set of scalar phases
in $\operatorname{Re}(e^{i\theta}T_\mathbb C f)
=T(\operatorname{Re}(e^{i\theta}f))$ to obtain the same
bound $N_\infty$. Apply the fully proved complex Riesz--Thorin
theorem \ref{fa2:p4:riesz-thorin} to these two restrictions of
the same linear operator. Every $L^p$ is contained in $L^1$
on this finite measure space. Its interpolated extension agrees
with the original operator by bounded simple approximation
in both $L^p$ and $L^1$, so this proves the displayed assertion
for the original $T$, rather than for an unrelated extension.
\end{proof}

% END INPUT


% BEGIN parts/part2-geometry.tex
\chapter{Contractive projections characterize Hilbert geometry}
\label{fa2:p2:kakutani-chapter}

This chapter supplies the finite-dimensional geometric bridge in the
projection characterization. It is not obtained by assuming that the
norm is already smooth. The meridian argument is due to Phillips,
\emph{A characterization of Euclidean spaces}, Bull. Amer. Math. Soc.
46 (1940), 930--933, DOI 10.1090/S0002-9904-1940-07336-7.
We give a normal-cone argument for the preliminary regularity and use
a covariance integral for the final ellipse step. All spaces in the
geometric argument are real.

\section{The three-dimensional lemma}
Fix auxiliary Euclidean coordinates on a three-dimensional vector space.
Let $K$ be a centrally symmetric convex body with $0$ in its interior.
Suppose that every plane $L$ through $0$ is the range of a linear
projection $P_L$ with $P_LK\subset K$. Thus $P_LK=K\cap L$.
The operator norms of all these projections in the auxiliary Euclidean
norm have one common bound: choose $rB_2\subset K\subset RB_2$;
then $\|P_L\|_2\le R/r$.

For $p\in\partial K$ write
\[
 N(p)=\{f\in(\mathbb R^3)': f(x)\le f(p)\ (x\in K)\}.
\]
This is the closed cone of outward supporting functionals, including
zero. Every nonzero member satisfies $f(p)>0$. Supporting functionals
exist by separation. The same construction in a central section gives
its planar normal cone. Restriction of $N(p)$ to a central plane $L$
containing $p$ maps onto that planar normal cone: a planar supporting
functional $g$ extends as $gP_L$, and $gP_L\in N(p)$. This observation
will control possible corners.

\begin{lemma}[Regularity and uniqueness of the projection direction]
\label{fa2:p2:projection-regularity}
Under these hypotheses $\partial K$ has a unique supporting plane at
every point. The projection onto each central plane is unique, and its
kernel line depends continuously on that plane.
\end{lemma}
\begin{proof}
First suppose that the boundary curve of $K\cap L$ has a unique tangent
at $p$. Let $t$ denote its tangent direction. It is not parallel to $p$,
since $0$ is an interior point of the section. Every member of $N(p)$
vanishes on $t$, since its restriction to $L$ is a nonnegative multiple
of the unique planar normal. Thus $N(p)$ lies in a two-dimensional
linear space. If it is not a ray, it is a pointed two-dimensional
closed cone with two distinct extreme rays, generated by $f_1,f_2$.
Their common null line is $\mathbb Rt$, and their restrictions to $L$
are positive multiples of each other.

Let $M\ne L$ be another central plane containing $p$. Restriction from
$\operatorname{span}\{f_1,f_2\}$ to $M'$ is injective. Indeed a nonzero
linear combination vanishing on $M$ also vanishes on $p$, and the only
such line of combinations has kernel $L$: on $L$ both $f_j$ are
proportional and nonzero at $p$. Thus a kernel $M\ne L$ is impossible.
The normal cone of $K\cap M$ is the restriction of $N(p)$, as noted
above. Pulling its members back by $P_M$ and using this injectivity
shows that each $f\in N(p)$ equals $(f|_M)P_M$. Consequently
$f_1,f_2$ both annihilate $\ker P_M$, so $\ker P_M=\mathbb Rt$.
Choose planes $M_j\ne L$ containing $p$ and converging to $L$.
The uniformly bounded matrices $P_{M_j}$ have a convergent subsequence.
Its limit is a projection onto $L$ acting as the identity on $L$,
but it annihilates $t\in L$, a contradiction. Therefore $N(p)$ is a
ray whenever one central section is differentiable at $p$.

A planar convex curve is differentiable except at countably many
points. One proof parametrizes it locally as the graph of a convex
function; its one-sided derivatives are monotone and differ only at
their countably many jumps. Hence in every central section the points
just considered are dense. At such a point $p$, the unique supporting
plane must contain $\ker P_L$: a planar normal pulled back by $P_L$
is a supporting normal, and there is only one supporting direction.
If $P_L,Q_L$ had different kernels, their two kernel lines would span
a plane $D\ne L$. Every tangent to the section at its dense smooth
points would then be parallel to the fixed line $D\cap L$. This is
impossible for a convex curve bounding a planar body. More explicitly,
the curve is locally Lipschitz, and its a.e. derivative would be
parallel to one fixed line; absolute continuity would place the entire
closed curve in a line. Thus $P_L$ is unique.

If $L_j\to L$, every convergent subsequence of the uniformly bounded
matrices $P_{L_j}$ is a projection onto $L$ mapping $K$ into $K$;
uniqueness identifies the limit with $P_L$. Compactness of bounded
matrix sets now proves convergence of the whole sequence. Kernel
lines consequently vary continuously.

It remains to establish smoothness at an arbitrary boundary point
$p$. In a fixed central section through $p$, approach it along one
side through smooth section points $p_j$. Those points are smooth
for the body by the first paragraph. A subsequence of their supporting
plane directions converges to a plane direction $D$. For any central
plane $M$ containing $p$, choose $M_j$ containing $p_j$ with $M_j\to M$.
The supporting plane at $p_j$ contains $\ker P_{M_j}$, so $D$ contains
$\ker P_M$. These kernel lines, as $M$ ranges over the planes through
$p$, span at least a plane. If they all coincided with a line $\ell$,
the plane generated by $p,\ell$ would contain its own kernel (if
$\ell=\mathbb Rp$, every such plane would), which is impossible.
Thus the limiting direction $D$ is uniquely determined by these
kernel lines, independently of the chosen section and side.
The one-sided tangents at $p$ of every central section are limits of
tangents at its smooth points, and hence belong to $D$. If the two
tangent lines of a section in $L$ were different, they would span
$L$ and force $D=L$. But $D$ also contains $\ker P_L$, a line not
contained in $L$. This contradiction shows that every central section
is differentiable at $p$. The first paragraph now applies and says
that $N(p)$ is a ray. This proves the remaining claim.
For a finite-dimensional convex body uniqueness of the supporting
plane implies continuous variation of its normal: any subsequential
limit of unit supporting normals is a supporting normal at the limit
point, and uniqueness identifies it. Thus the boundary is $C^1$ in
its local graph coordinates, as needed below.
\end{proof}

\begin{lemma}[Meridians are linearly equivalent]
\label{fa2:p2:meridian}
Choose $p\in\partial K$ and let $E$ be the central plane parallel to
its supporting plane. The sections by any two planes containing
$\mathbb Rp$ are linearly equivalent by a map fixing $p$ and preserving
the coordinate along $p$ in the decomposition $\mathbb R^3=E\oplus\mathbb Rp$.
\end{lemma}
\begin{proof}
Write a point as $u+zp$, $u\in E$, so $p$ has height $1$ and the
opposite supporting plane has height $-1$. For $-1<z<1$ the point
$zp$ is interior to $K$, by convexity and the interior ball at $0$.
Choose Euclidean polar coordinates $u=r e_\theta$ in $E$. Its horizontal
section has a positive radial boundary $r=r_z(\theta)$, of class $C^1$:
a vertical supporting normal at an intermediate height is impossible,
and a radial tangent would place its interior center on a supporting
line. For the vertical plane $L_\theta=\operatorname{span}(p,e_\theta)$,
its kernel line lies in $E$ because its projection fixes $p$ and the
unique supporting plane at $p$ must contain its kernel. Choose a
continuous local vector $v_\theta$ on this line. It is not parallel to
$e_\theta$, since it is transverse to $L_\theta$.

At the boundary point $r_z(\theta)e_\theta+zp$ the same line is tangent
to $\partial K$, and is horizontal. It is consequently tangent to the
horizontal section there. Write
$v_\theta=a_\theta e_\theta+b_\theta e_\theta'$, where
$b_\theta\ne0$. Comparison with its boundary tangent gives
\[
 \frac{r_z'(\theta)}{r_z(\theta)}=
 \frac{a_\theta}{b_\theta}=F(\theta).
\]
The right side is continuous and independent of $z$; it is unaffected
by changing the sign or scale of $v_\theta$. Integrating this scalar
equation gives $r_z(\theta)=h(z)g(\theta)$ with $g>0$. At height zero
central symmetry gives $g(\theta+\pi)=g(\theta)$, and then symmetry
of $K$ gives $h(-z)=h(z)$. Thus a section in direction $e_\theta$ has
the form $|s|\le h(z)g(\theta)$. The linear map fixing $p$ and sending
$g(\theta)e_\theta$ to $g(\phi)e_\phi$ carries it precisely onto
the section in direction $e_\phi$. Endpoint sections at $z=\pm1$
follow by closedness (or by continuity of the meridian curves).
\end{proof}

\begin{theorem}[Three-dimensional ellipsoid lemma]
\label{fa2:p2:ellipsoid}
The body $K$ under consideration is an ellipsoid.
\end{theorem}
\begin{proof}
Fix a central plane $L$ and $p,q\in\partial K\cap L$.
Choose a supporting plane parallel to $L$ and a point $r$ where it
touches $K$; $r\notin L$. Applying the meridian lemma about the axis
$p$ carries the section in $L$ to that in $\operatorname{span}(p,r)$,
fixing $p$. Applying it about $r$ then carries this section to that in
$\operatorname{span}(q,r)$ and sends $p$ to $q$: relative to this axis
both points have height zero, and the horizontal plane is precisely
$L$. Finally use the lemma about $q$ to return to $L$, fixing $q$.
The composition is a linear automorphism of the planar body $K\cap L$
carrying $p$ to $q$. The cases of coincident axes are obtained by the
identity or central symmetry. Hence its linear automorphism group
is transitive on its boundary.

Here is a direct reason such a body is an ellipse. In Euclidean
coordinates on $L$ put $D=K\cap L$ and
$S=\int_D xx^{\mathsf T}\,dx$. This matrix is positive definite,
because $D$ contains a disk. Every automorphism $A$ of $D$ has
$|\det A|=1$ by the area change formula. Applying that formula to
the matrix integral yields $S=ASA^{\mathsf T}$, so $A$ preserves the
positive quadratic form $x^{\mathsf T}S^{-1}x$. Transitivity makes
this form constant on $\partial D$. Every ray meets that boundary
once; thus $D$ is the corresponding ellipse. Every central planar
section of $K$ is therefore an ellipse. Its Minkowski norm satisfies
the parallelogram identity on each plane, hence on all vector pairs.
Polarization, proved earlier, makes it a Euclidean norm. Its unit
body $K$ is an ellipsoid.
\end{proof}

\section{The normed-space conclusion}
\begin{theorem}[Kakutani's projection characterization]
\label{fa2:p2:kakutani}
Let $X$ be a real normed space of dimension at least three. Its norm
comes from an inner product if and only if every two-dimensional
linear subspace is the range of a projection of norm one. Completeness
is not required; when $X$ is complete the resulting inner-product
space is a Hilbert space.
\end{theorem}
\begin{proof}
In an inner-product space a finite-dimensional subspace has its
orthogonal projection, whose norm is one. Conversely restrict the
given projections to any three-dimensional subspace. They remain
contractive projections onto its planes. The preceding ellipsoid
lemma applies to its unit ball. Thus every two-dimensional subspace,
being contained in such a three-dimensional space, has an inner-product
norm. The parallelogram identity consequently holds for every pair
of vectors in $X$. Polarization constructs the required inner product
on the whole space. Completeness of the norm is exactly Hilbert
completeness. The dimension restriction matters: the identity is a
norm-one projection onto the sole two-dimensional subspace of any
two-dimensional normed space, whose unit ball need not be an ellipse.
\end{proof}

\section{The radial retraction and orthogonality}
For a real normed space put $R(x)=x/\max(1,\|x\|)$. The next result
is the radial characterization of de Figueiredo--Karlovitz,
Bull. Amer. Math. Soc. 73 (1967), 364--368,
DOI 10.1090/S0002-9904-1967-11753-1. The scalar convexity argument is
included because mere nondecrease on a ray does not locate all its
minimum points.

\begin{lemma}[Radial nonexpansion criterion]\label{fa2:p2:radial-criterion}
The map $R$ is nonexpansive if and only if the following property holds:
whenever $\|x\|=\|y\|>0$ and $\lambda_0$ minimizes
$\lambda\mapsto\|x-\lambda y\|$ on the real line, then
$|\lambda_0|\le1$.
\end{lemma}
\begin{proof}
For a convex coercive scalar function its minimizers form a compact
interval; it is strictly decreasing before that interval and strictly
increasing after it. This follows directly from monotonicity of
secant slopes. Thus the stated property is equivalent to strict
decrease on $(-\infty,-1]$ and strict increase on $[1,\infty)$ for
every equal-norm pair.

Assume first that $R$ is nonexpansive. For unit $x,y$ and $\lambda\ge1$,
comparison of $x,\lambda y$ gives
$\|x-y\|\le\|x-\lambda y\|$. Convexity makes this function
nondecreasing on $[1,\infty)$. It must in fact be strictly increasing.
Otherwise convexity gives an interval $[1,b]$, $b>1$, on which it is
constant, say $c$. Here $c>0$, since $x=y$ gives strict increase.
Put $z=\alpha x+(1-\alpha)y$ for a sufficiently small $\alpha>0$.
We have $\|z\|\le1$, $\|z\|>\alpha$, and
\[
 t_0=\frac{1-\alpha}{\|z\|-\alpha}\ge1,
 \qquad t_0\longrightarrow1\quad(\alpha\downarrow0).
\]
Choose $\alpha$ so $t_0<b$, then $\delta>0$ so small that
$\|z\|>(1+\delta)\alpha$ and
$t_1=(1+\delta)(1-\alpha)/(\|z\|-(1+\delta)\alpha)<b$.
Both $t_0,t_1$ lie in $[1,b]$. Directly factoring the two vectors gives
\[
 \left\|x-\frac z{\|z\|}\right\|
   =\left(1-\frac\alpha{\|z\|}\right)c,
 \quad
 \left\|x-\frac{(1+\delta)z}{\|z\|}\right\|
   =\left(1-\frac{(1+\delta)\alpha}{\|z\|}\right)c.
\]
The second is smaller, contradicting the already established
nondecrease for the unit pair $x,z/\|z\|$. Replacing $y$ by $-y$
proves the negative-ray assertion. Scaling gives the property for
all equal-norm pairs.

Conversely assume the property. If $\|x\|,\|y\|\le1$ there is nothing
to prove. If $0<\|x\|=s\le1<r=\|y\|$, apply the positive-ray
monotonicity to the equal-norm pair $x,s y/r$, at coefficients
$1/s$ and $r/s$, to get $\|x-y/r\|\le\|x-y\|$.
The case $x=0$ is immediate. If $r=\|x\|\ge s=\|y\|\ge1$,
apply it to the unit pair $y/s,x/r$, at coefficients $1$ and $r/s$:
\[
 \|x/r-y/s\|\le\|(r/s)x/r-y/s\|
                 =\|x-y\|/s\le\|x-y\|.
\]
These are all cases, up to interchanging $x,y$.
\end{proof}

\begin{theorem}[Radial characterization]\label{fa2:p2:radial-hilbert}
If a real normed space has dimension at least three, its radial
retraction onto the closed unit ball is nonexpansive exactly when
its norm comes from an inner product.
\end{theorem}
\begin{proof}
In an inner-product space $R$ is the metric projection onto the
closed unit ball, so the projection inequality proved earlier applies.
For the converse restrict to an arbitrary three-dimensional subspace
$E$. For any plane $L\subset E$, choose $x_0\notin L$ and its nearest
point $y_0\in L$, which exists by finite-dimensional compactness of
bounded minimizing sequences. Put $u=x_0-y_0$ and define the linear
projection $P(\alpha u+y)=y$. For $x=\alpha u+y$, with $x,y\ne0$,
minimality gives
$\|x-y\|\le\|x-\lambda y\|$ for all real $\lambda$ (the case
$\alpha=0$ is immediate). The pair
$x,(\|x\|/\|y\|)y$ has equal norm, and its minimizing coefficient
is $\|y\|/\|x\|$. The preceding criterion bounds it by one.
Thus $\|Px\|\le\|x\|$, also when $y=0$. Every plane in $E$ is
contractively complemented. Kakutani's theorem makes $E$ Euclidean.
The parallelogram identity follows for every pair in the whole space,
and polarization gives its inner product.
\end{proof}

\begin{proposition}[Why the two-dimensional exception is genuine]
\label{fa2:p2:radial-plane-exception}
Write $x\perp_B y$ if $\|x\|\le\|x+t y\|$ for all real $t$.
Radial nonexpansion is equivalent, in any real normed space, to
symmetry of this Birkhoff orthogonality. The norm
\[
 \|(s,t)\|=\max(|s|,|t|,|s+t|)
\]
on $\mathbb R^2$ has this property but is not an inner-product norm.
\end{proposition}
\begin{proof}
Use the preceding equal-norm minimum property. If $x\perp_B y$ but
$\|y+t x\|=m\|y\|<\|y\|$, then $0<m<1$; $m=0$ would make
nonzero $x,y$ dependent, contrary to $x\perp_B y$.
The vectors $y+t x$ and $m y$ have equal norm. Strict increase beyond
coefficient one gives
\[
 \|t x+(1-m)y\|<\|t x\|,
\]
contradicting $x\perp_B y$ after division by $|t|$.
Conversely, if $\|x\|=\|y\|>0$ and $x-\lambda_0y$ is a nearest
point to zero on this affine line, then
$x-\lambda_0y\perp_B y$. If this residual is zero then
$|\lambda_0|=1$. Otherwise symmetry, tested with coefficient
$1/\lambda_0$ when $\lambda_0\ne0$, gives
$\|y\|\le\|x\|/|\lambda_0|$. The zero coefficient already satisfies
the required bound. This proves the minimum property.

For the displayed planar norm let $[x,y]=x_1y_2-x_2y_1$ and define
its antinorm $\|y\|_a=\sup_{\|x\|\le1}|[x,y]|$. The six vertices
of its unit ball are $(1,0),(1,-1),(0,-1)$ and their negatives.
Maximizing the linear functional over these vertices gives
$\|(s,t)\|_a=\max(|s|,|t|,|s+t|)=\|(s,t)\|$.
Orthogonality involving a zero vector is immediate, so assume both
vectors nonzero in the following norming-functional argument.
Hahn--Banach on the quotient by $\mathbb R y$ says
$x\perp_B y$ iff a norming functional for $x$ annihilates $y$.
Every such planar functional is a scalar multiple of $[\,\cdot,y]$,
so this is equivalent to
$|[x,y]|=\|x\|\|y\|_a$. Equality of norm and antinorm makes this
condition symmetric in $x,y$. Thus radial nonexpansion holds.
Its hexagonal ball is not an ellipse, so its norm is not an
inner-product norm.
\end{proof}


% END INPUT


% BEGIN parts/part2-compact.tex
\chapter{Compactness, weak compactness, and adjoints}
Here $X,Y,Z$ are Banach spaces unless another assumption is stated.
The Banach adjoint is denoted by $T':Y'\to X'$; the canonical isometry
$J_X:X\to X''$ satisfies $(J_Xx)(x')=x'(x)$.
This chapter follows the distinctions in \cite{B01,B02,B05}.

\begin{proposition}[Banach adjoints and preadjoints]\label{fa2:p2:banach-adjoint}
For a bounded operator $T:X\to Y$, $T'y'=y'\circ T$ is bounded,
$\|T'\|=\|T\|$, and $T''J_X=J_YT$. For normed $X,Y$, a linear map
$S:Y'\to X'$ has the form $T'$ for a bounded $T:X\to Y$ exactly when
it is continuous from $\sigma(Y',Y)$ to $\sigma(X',X)$.
\end{proposition}
\begin{proof}
The upper norm bound follows from $|y'(Tx)|\leq\|y'\|\|T\|\|x\|$.
Hahn--Banach gives $\|Tx\|=\sup_{\|y'\|\leq1}|y'(Tx)|$, giving the
opposite bound. Evaluation against $y'$ proves bidual naturality.
If $S$ is weak-star continuous, for every $x\in X$ the scalar map
$y'\mapsto(Sy')(x)$ is weak-star continuous. Such a scalar functional is
a finite linear combination of evaluations on $Y$: continuity at zero
bounds it by finitely many defining seminorms, so it vanishes on their
common kernel and factors through their finite-dimensional image.
Thus there is a unique $Tx\in Y$ with $(Sy')(x)=y'(Tx)$ for all $y'$.
This defines a linear $T$. Its graph is closed: if $x_n\to x$ and
$Tx_n\to y$, then $y'(y)=\lim(Sy')(x_n)=(Sy')(x)=y'(Tx)$.
For Banach $X,Y$ the closed graph theorem gives boundedness. For general
normed spaces one can instead use uniform boundedness on the Banach
space $Y'$: the evaluations $y'\mapsto y'(Tx)$, $\|x\|\leq1$, are
pointwise bounded because $Sy'\in X'$, hence $\sup_{\|x\|\leq1}\|Tx\|<\infty$.
Conversely evaluations $x$ after $T'$ are evaluations $Tx$, proving
weak-star continuity.
\end{proof}

\begin{theorem}[Schauder and the compact ideal]\label{fa2:p2:schauder-adjoint}
The compact operators $\faK(X,Y)$ form a norm-closed linear space and a
two-sided operator ideal. Moreover $T$ is compact iff $T'$ is compact.
A compact operator sends every weakly convergent sequence to a norm
convergent sequence with the corresponding limit. This last assertion
is a statement about sequences, not continuity from the entire weak
topology to the norm topology.
\end{theorem}
\begin{proof}
The image of a bounded sequence under a sum of two compact operators has
a convergent subsequence by two successive extractions. Composing with
a bounded operator preserves the same property. If $T_n\to T$ in norm,
given $\varepsilon>0$ choose $n$ with $\|T-T_n\|<\varepsilon/3$.
A finite $\varepsilon/3$-net for $T_nB_X$ gives a finite
$\varepsilon$-net for $TB_X$. Thus $TB_X$ is totally bounded and its
closure is compact in the complete space $Y$.
Suppose $T$ compact. The restrictions of $B_{Y'}$ to the compact metric
space $K=\overline{TB_X}$ are uniformly bounded and uniformly Lipschitz.
Arzela--Ascoli gives compactness in $C(K)$; the identity
$\|T'y'-T'z'\|=\sup_{y\in K}|y'(y)-z'(y)|$ proves compactness of $T'$.
If $T'$ is compact, the forward direction makes $T''$ compact, and
$J_YT=T''J_X$ is compact. Since $J_Y$ is an isometry onto a closed
subspace, $T$ is compact.
Finally a weakly convergent sequence is norm bounded by uniform
boundedness. Every subsequence of its image has a norm-convergent
subsequence; its limit must equal the weak limit $Tx$, since continuous
functionals separate points. If the whole image did not converge in
norm, a subsequence at fixed positive distance from $Tx$ would contradict
this observation.
\end{proof}

\begin{definition}\label{fa2:p2:ap-definition}
$X$ has the approximation property (AP) if for each norm-compact $K\subset X$
and $\varepsilon>0$ there is a bounded finite-rank $F:X\to X$ with
$\sup_{x\in K}\|Fx-x\|<\varepsilon$. No uniform bound on $\|F\|$ is part
of AP. Such a bound is an additional bounded approximation property.
\end{definition}

\begin{lemma}[Compact sets and compact disks]\label{fa2:p2:compact-disk}
Every compact $K\subset X$ is contained in the closed absolutely convex
hull of a norm-null sequence, chosen so that $K$ is compact also for the
Minkowski norm below. The closed absolutely convex hull $C$ of
a norm-null sequence is compact. There is a Banach space $X_C$ and an
injective compact operator $i_C:X_C\to X$ whose closed unit ball has
image $C$.
\end{lemma}
\begin{proof}
Choose finite $2^{-2n}$-nets $F_n$ for $K$. For every $x\in K$ choose
$a_n(x)\in F_n$ with $\|x-a_n(x)\|\leq2^{-2n}$.
Then $x=a_1(x)+\sum_{n\geq1}(a_{n+1}(x)-a_n(x))$.
The finite sets of scaled increments $2^{n+1}(a_{n+1}-a_n)$ that occur
have norms tending to zero. Together with $2a_1$ they form a null
sequence $z_j$, and the coefficients $1/2,2^{-n-1}$ have sum one.
Thus $K\subset\overline{\operatorname{aconv}}\{z_j\}$.
For a null sequence, split an absolutely convex combination after its
first $N$ terms. Its tail has norm at most $\sup_{j>N}\|z_j\|$; its head
lies in a compact finite-dimensional absolutely convex hull. This proves
total boundedness; closure in $X$ proves compactness of $C$.
Let $X_C=\bigcup_{n\geq1}nC$ with Minkowski norm
$p_C(x)=\inf\{t>0:x\in tC\}$. Boundedness of $C$ makes this a norm,
and closedness gives its unit ball exactly $C$. A $p_C$-Cauchy sequence
is Cauchy in $X$, with a limit $x$; for each $\varepsilon>0$ its tails
lie in translates of $\varepsilon C$, which are closed in $X$.
Hence $x-x_n\in\varepsilon C$ eventually, proving completeness.
The inclusion has compact unit-ball image $C$ and is injective.
For the special disk constructed for $K$, the tails of its increment
expansions have $p_C$-norm at most $\sum_{n>N}2^{-n-1}$, uniformly on $K$.
There are only finitely many heads through a fixed $N$, so $K$ is
totally bounded for $p_C$. Completeness and the injective continuous
inclusion identify its limit points with points of the original $K$;
hence $K$ is $p_C$-compact.
\end{proof}

\begin{theorem}[The operator form of AP]\label{fa2:p2:approximation-property}
$X$ has AP iff for every Banach $Y$, every compact $T:Y\to X$ is a norm
limit of bounded finite-rank operators $Y\to X$. Hilbert spaces and
Banach spaces with a Schauder basis have AP; the latter have a common
bound on the norms of the approximants.
The existence of a Banach space without AP is the separate construction
at \ref{fa2:ap:counterexample}.
\end{theorem}
\begin{proof}
If $X$ has AP, approximate the identity uniformly on
$\overline{TB_Y}$ by $F$; then $FT$ is finite rank and
$\|FT-T\|<\varepsilon$.
Conversely let $K\subset C=i_C B_{X_C}$ be as in the lemma. By the
operator assumption choose a finite-rank $S:X_C\to X$ with
$\|S-i_C\|<\delta$. Write $Sx=\sum_{j=1}^m f_j(x)y_j$, $f_j\in X_C'$.
We need functionals from $X'$ in this formula. Each $f\in X_C'$ restricted
to the $p_C$-compact set $K$ is a uniform limit of restrictions from $X'$.
If this failed, Hahn--Banach in $C(K)$ and scalar Riesz representation
would give a finite complex measure $\mu$ annihilating all restrictions
from $X'$ but not $f|_K$. The Bochner integral
$z=\int_Kx\,d\mu(x)$ exists in $X_C$: the identity on the compact metric
set $K$ is continuous and has separable bounded range. Every $x'\in X'$
annihilates $i_Cz$, so Hahn--Banach gives $i_Cz=0$. Injectivity gives
$z=0$, contradicting $\int_K f(x)\,d\mu=f(z)\ne0$.
Choose $g_j\in X'$ with
$\sum_j\|y_j\|\sup_K|f_j-g_j|<\varepsilon/2$.
The bounded finite-rank map $Fx=\sum_jg_j(x)y_j$ then satisfies
$\sup_K\|Fx-x\|\leq\varepsilon/2+\delta\sup_Kp_C(x)<\varepsilon$
for a sufficiently small $\delta$. This proves the converse.
For a Hilbert space the projections onto finite subsets of an
orthonormal basis converge strongly to the identity and have norm at
most one. A finite-net argument makes this convergence uniform on every
compact set: approximate its points by a finite net, use convergence on
that net, and bound each remaining error by twice the net radius.
This proves AP directly for Hilbert spaces.
For a Schauder basis its bounded partial-sum maps converge strongly
to the identity; uniform boundedness gives a common operator norm
bound $M$. The same finite-net argument bounds a compact-set error
by the largest net-point error plus $(1+M)$ times the net radius.
Thus these partial sums give the bounded approximation property.
\end{proof}

\begin{theorem}[Gantmacher]\label{fa2:p2:gantmacher}
For $T\in\faB(X,Y)$ the following are equivalent:
$T$ is weakly compact; $T''X''\subset J_Y Y$;
$T':(Y',\sigma(Y',Y))\to(X',\sigma(X',X''))$ is continuous.
Also $T$ is weakly compact iff $T'$ is weakly compact. Weakly compact
operators form a norm-closed two-sided ideal.
\end{theorem}
\begin{proof}
Suppose $T$ is weakly compact and $x''\in B_{X''}$. Goldstine gives a net
$x_i\in B_X$ with $J_Xx_i\to x''$ weak-star. A subnet of $Tx_i$
converges weakly to $y\in Y$. Evaluating against $y'$ gives
$T''x''=J_Yy$. Conversely if this inclusion holds, the weak-star compact
set $T''B_{X''}$ lies in $J_YY$. The weak-star topology restricted there
is exactly the transported weak topology of $Y$. It is a compact set
containing $J_YTB_X$, proving weak compactness.
The scalar functions defining the asserted continuity of $T'$ are
$y'\mapsto x''(T'y')=(T''x'')(y')$. They are weak-star continuous iff
$T''x''$ is an evaluation at a vector of $Y$, proving equivalence.
If $T''X''\subset J_YY$ and $y'''\in Y'''$, put $y'=y'''J_Y\in Y'$.
For every $x''$,
$T'''y'''(x'')=y'''(T''x'')=x''(T'y')$.
Hence $T'''Y'''\subset J_{X'}X'$ and the same criterion makes $T'$
weakly compact. Applying this implication to $T'$ gives weak
compactness of $T''$, hence of its restriction to $J_XX$, namely $T$.
For norm closure, if $T_n\to T$ then $T_n''x''\to T''x''$ in norm;
$J_YY$ is norm closed, so the criterion applies. Naturality of canonical
embeddings proves the ideal property directly from the criterion.
\end{proof}

\begin{proposition}[Complete continuity]\label{fa2:p2:complete-continuity}
A bounded $T:X\to Y$ maps weakly compact sets to norm relatively compact
sets iff it maps weakly convergent sequences to norm convergent
sequences. Every compact operator has this property. If $X$ is reflexive,
the property is equivalent to compactness. Without reflexivity the
converse fails: the identity on $\ell^1$ is completely continuous and is
not compact.
\end{proposition}
\begin{proof}
The set consisting of a weakly convergent sequence and its limit is
weakly compact. Relative norm compactness of its image, together with
the unique weak limit, gives norm convergence. Conversely the
Eberlein--Smulian theorem gives a weakly convergent subsequence in every
sequence in a weakly compact set; the assumed property makes its image
norm convergent. Metric sequential relative compactness is relative
compactness. If $X$ is reflexive, its closed unit ball is weakly compact,
giving the converse to compactness.
For the example we prove Schur's property. If $x_n\to0$ weakly in
$\ell^1$ but $\|x_n\|_1\geq\varepsilon$, coordinate convergence and a
successive choice give a subsequence $y_j=x_{n_j}$ and integers
$0=m_0<m_1<\cdots$ such that
$\sum_{k\leq m_{j-1}}|y_j(k)|<\varepsilon/8$ and
$\sum_{k>m_j}|y_j(k)|<\varepsilon/8$.
On the block $(m_{j-1},m_j]$ choose $a_k$ of modulus one with
$a_ky_j(k)=|y_j(k)|$. This defines $a\in\ell^\infty$ and
$\operatorname{Re}\sum_k a_ky_j(k)\geq\varepsilon/2$, contradicting
weak convergence. Thus weak convergence in $\ell^1$ implies norm
convergence. The unit coordinate vectors have mutual distance two,
so its identity is not compact.
\end{proof}

\begin{theorem}[Dunford--Pettis testing criterion]\label{fa2:p2:dunford-pettis}
For a Banach space $X$ the following are equivalent:
every weakly compact map $X\to Y$, for every Banach $Y$, is completely
continuous; the same assertion for $Y=c_0$; and
\[
 x_n\rightharpoonup0\text{ in }X,\quad
 x_n'\rightharpoonup0\text{ in }X'
 \quad\Longrightarrow\quad x_n'(x_n)\longrightarrow0.
\]
\end{theorem}
\begin{proof}
Only the two nontrivial implications need proof. Given weakly null
$x_n'$, define $Tx=(x_n'(x))_n\in c_0$. Uniform boundedness gives
boundedness of $T$. For $a\in\ell^1=c_0'$, $T'a=\sum a_nx_n'$.
Its unit-ball image lies in the closed absolutely convex hull of the
weakly compact set $\{x_n':n\geq1\}\cup\{0\}$, which is weakly compact
by \ref{fa2:p1:weak-compact-hull}. Thus $T'$ and $T$ are weakly compact
by Gantmacher. Complete continuity of $T$ gives
$|x_n'(x_n)|\leq\|Tx_n\|_\infty\to0$.
Conversely let $T$ be weakly compact and $x_n\rightharpoonup0$.
If $\|Tx_n\|\not\to0$, pass to a subsequence and choose $y_n'\in B_{Y'}$
with $|y_n'(Tx_n)|\geq\varepsilon$. Weak compactness of $T'$ and
Eberlein--Smulian give $T'y_n'\rightharpoonup x'$ along a further
subsequence. Then the testing criterion applied to
$T'y_n'-x'$ and $x_n$ gives $y_n'(Tx_n)-x'(x_n)\to0$, a contradiction.
The compact-convex-hull step is supplied by the separate proof in
Part I, rather than replaced by a sequential compactness slogan.
\end{proof}

\begin{remark}\label{fa2:p2:dp-examples-open}
The classical assertions that $C(K)$ and $L^1(\mu)$ have the
Dunford--Pettis property are stronger than Schur's property for
$\ell^1$. Their source in \cite{B05}, 3.5.19, refers to outside proofs.
The separate proof at \ref{fa2:p2:dp-classical-examples} supplies
these examples, including its Vitali--Hahn--Saks bridge
\ref{fa2:p2:vitali-hahn-saks}.
\end{remark}

\begin{lemma}[Lomonosov's transitive-algebra lemma]\label{fa2:p2:lomonosov-lemma}
Let $\mathcal A\subset\faB(X)$ be a unital algebra on a complex Banach
space with no nontrivial common closed invariant subspace. For every
nonzero compact $K$ there is $A\in\mathcal A$ such that $AK$ has
eigenvalue one.
\end{lemma}
\begin{proof}
Rescale $K$ to have norm one and choose $x_0$ with $\|Kx_0\|>1$.
Put $S=\{x:\|x-x_0\|\leq1\}$. Then $0\notin\overline{KS}$, because
$\|Kx\|\geq\|Kx_0\|-1>0$ on $S$.
For each nonzero $y$, the closure of $\mathcal A y$ is a nonzero
invariant subspace, hence is $X$. Thus some $A_y\in\mathcal A$ has
$\|A_yy-x_0\|<1$. Compactness of $\overline{KS}$ gives finitely many
$A_1,\ldots,A_m$ whose corresponding open sets cover it.
Define $b_j(y)=\max(0,1-\|A_jy-x_0\|)$, and divide each by
$\sum b_j(y)>0$ to obtain a partition $c_j$ on $\overline{KS}$.
The continuous map $\Phi:S\to S$ given by
$\Phi(x)=\sum c_j(Kx)A_jKx$ has relatively compact image: it lies in
the convex hull of the finite union of the compact sets
$A_j\overline{KS}$, and a finite convex hull of compact sets is compact
(it is the continuous image of a compact simplex times their product).
The Schauder fixed-point theorem supplies $x\in S$ with $\Phi(x)=x$.
Since $0\notin S$, this is a nonzero eigenvector of
$AK$, where $A=\sum c_j(Kx)A_j\in\mathcal A$.
\end{proof}

\begin{theorem}[Lomonosov]\label{fa2:p2:lomonosov}
If a nonscalar $T\in\faB(X)$ on a complex Banach space commutes with a
nonzero compact operator, then it has a nontrivial closed hyperinvariant
subspace. In particular it has a nontrivial invariant subspace.
\end{theorem}
\begin{proof}
Suppose the commutant $\mathcal A=\{T\}'$ were transitive.
The lemma gives $A\in\mathcal A$ with $1\in\sigma_p(AK)$.
The eigenspace $E=\ker(AK-I)$ is nonzero and finite dimensional by the
compact-operator Riesz theorem proved in the Fredholm chapter.
Both $A$ and $K$ commute with $T$, so $E$ is $T$-invariant.
Over the complex field the finite-dimensional restriction has an
eigenvalue $\lambda$, hence $\ker(T-\lambda I)\ne0$.
This kernel is invariant under every member of $\{T\}'$.
Transitivity forces it to be $X$, which would make $T=\lambda I$.
This contradicts the nonscalar assumption. Thus the commutant has a
nonzero proper closed invariant subspace, which is precisely a
hyperinvariant subspace for $T$.
\end{proof}

% END INPUT


% BEGIN parts/part2-dp-examples.tex
\chapter{The classical Dunford--Pettis spaces}
\label{fa2:p2:dp-examples-chapter}

The weak compactness criterion for subsets of $L^1$ and the
Dunford--Pettis property of a Banach space are different assertions.
Here we prove the latter for $C(K)$ and $L^1$ by the testing criterion
of Theorem~\ref{fa2:p2:dunford-pettis}.  The needed uniform
absolute-continuity step is proved rather than inferred from pointwise
convergence of measures.  These are the examples recorded in
\cite[Theorem 3.5.19]{B05}.

\section{A uniform absolute-continuity theorem}

\begin{theorem}[Vitali--Hahn--Saks, finite control-measure form]
\label{fa2:p2:vitali-hahn-saks}
Let $\lambda$ be a finite positive measure on a sigma algebra
$\Sigma$, and let $\mu_n$ be finite signed or complex countably
additive measures absolutely continuous with respect to $\lambda$.
If $\mu_n(A)$ converges for every $A\in\Sigma$, then
\[
 \lim_{\delta\downarrow0}\sup_n
       \sup_{\lambda(A)<\delta}|\mu_n|(A)=0.
\]
The setwise limit is countably additive and absolutely continuous
with respect to $\lambda$.  In particular the assertion applies to
a weakly convergent sequence in the space of finite regular
measures on a compact Hausdorff space, after taking
$\lambda=\sum_n2^{-n}|\mu_n|/(1+\|\mu_n\|)$.
\end{theorem}
\begin{proof}
Identify measurable sets which differ by a $\lambda$-null set and
give the resulting Boolean algebra the metric
$d(A,B)=\lambda(A\mathbin\triangle B)$.  It is complete.  Indeed
from a Cauchy sequence choose a subsequence with summable successive
symmetric-difference measures.  Outside a null set the indicators
of this subsequence eventually stabilize, and their limiting
measurable set is its metric limit; the original Cauchy sequence
has the same limit.  Each map $A\mapsto\mu_n(A)$ is continuous
for this metric, since
$|\mu_n(A)-\mu_n(B)|\le|\mu_n|(A\mathbin\triangle B)$ and a finite
absolutely continuous measure has the epsilon-delta absolute
continuity property.  The latter follows directly by contradiction:
sets of measures tending to zero have a subsequence with summable
$\lambda$-measures; its limsup is null, and continuity from above
of the finite measure rules out variation bounded below on the
subsequence.

Fix $\eps>0$.  The closed sets
\[
 F_N=\{A: |(\mu_n-\mu_m)(A)|\le\eps
                           \quad(n,m\ge N)\}
\]
cover the complete measure algebra by setwise Cauchyness.  Baire
gives a metric ball about a set $A_0$, of radius $\delta>0$,
contained in some $F_N$.  If $\lambda(E)<\delta$ and $n,m\ge N$,
write $\nu=\mu_n-\mu_m$.  The three sets $A_0$, $A_0\cup E$,
and $A_0\setminus E$ belong to this ball.  Finite additivity gives
\[
 |\nu(E\setminus A_0)|\le2\eps,\qquad
 |\nu(E\cap A_0)|\le2\eps,
 \qquad |\nu(E)|\le4\eps.
\]
Now keep $m=N$ fixed.  By absolute continuity of $\mu_N$ and of
the finitely many initial measures, reduce $\delta$ so that
$\lambda(E)<\delta$ also implies $|\mu_j(E)|<\eps$ for $j\le N$.
It follows that $\sup_n|\mu_n(E)|\le5\eps$ on all such sets.
The same estimate applies to every measurable subset of $E$.
For a real signed measure, total variation on $E$ is at most twice
the supremum of the absolute values on subsets of $E$; apply
this to real and imaginary parts to get, for complex measures,
\[
 |\mu_n|(E)\le4\sup_{B\subset E}|\mu_n(B)|\le20\eps.
\]
This proves uniform absolute continuity of the variations.

The pointwise limit $\mu$ is finitely additive.  If
$E_j\downarrow\varnothing$, then $\lambda(E_j)\downarrow0$ and
the uniform estimate implies $\mu(E_j)\to0$.  A finitely additive
scalar set function with this continuity is countably additive:
for a disjoint union apply it to the tails after subtracting the
first finitely many pieces.  It vanishes on $\lambda$-null sets.
A finite countably additive signed or complex measure has finite
variation, the usual scalar-measure Jordan decomposition fact
used in Part IV.  In the application to a weakly convergent
sequence the variations are in addition uniformly bounded in
total mass by uniform boundedness.

For the last assertion the displayed $\lambda$ is finite and
regular (its positive series converges in variation), and dominates
every $\mu_n$.  Set evaluation $\mu\mapsto\mu(A)$ is a bounded
linear functional on the regular-measure Banach space for every
Borel $A$.  Thus weak convergence there gives exactly the setwise
convergence used above.  Notice that weak star convergence against
continuous functions alone would not justify this step.
\end{proof}

\section{The two examples}

\begin{lemma}[Weakly null bounded functions on finite-measure sets]
\label{fa2:p2:linfinity-weak-l1}
If $g_n\rightharpoonup0$ in the Banach space $L^\infty(\mu)$, then
$\int_A|g_n|\,d\mu\to0$ for every measurable $A$ of finite measure.
\end{lemma}
\begin{proof}
The zero measure algebra is trivial.  Otherwise complex
$L^\infty(\mu)$ is a commutative unital C*-algebra, and
Theorem~\ref{fa2:p2:commutative-cstar} identifies it isometrically
and as a *-algebra with $C(K)$ for a compact Hausdorff $K$.
Its Gelfand transforms $\widehat g_n$ are norm bounded and
converge pointwise to zero, since characters are bounded linear
functionals and $g_n$ converge weakly, not merely weakly star.
The positive bounded functional $g\mapsto\int_A g\,d\mu$ becomes
a finite positive regular measure $\nu_A$ on $K$ by Riesz--Markov.
The representation preserves absolute values: $|g|$ is the
positive square root of $g^*g$, so
$\widehat{|g|}=|\widehat g|$ by the continuous calculus.
Dominated convergence gives
$\int_A|g_n|=\int_K|\widehat g_n|\,d\nu_A\to0$.
Real functions give the same result by complexification: a
complex functional restricts to two continuous real functionals,
so real weak nullity implies weak nullity of that sequence in
the complex function space.  No Dunford--Pettis property is used
in this proof, only the independently proved commutative
representation and scalar dominated convergence.
\end{proof}

\begin{theorem}[The Dunford--Pettis property of $C(K)$ and $L^1$]
\label{fa2:p2:dp-classical-examples}
For every compact Hausdorff $K$, $C(K)$ has the Dunford--Pettis
property.  For every positive measure space, $L^1(\mu)$ has the
Dunford--Pettis property.  The statements hold over both scalar
fields, without separability hypotheses.
\end{theorem}
\begin{proof}
We verify the double weak-null sequence test.
Let $f_n\rightharpoonup0$ in $C(K)$ and
$\mu_n\rightharpoonup0$ in $C(K)'=M(K)$.
The norms of both sequences are bounded, say by $C$ and $M$,
and $f_n(t)\to0$ for every $t\in K$ by point evaluation.
Use the finite control measure in the preceding theorem.
Dominated convergence gives $\lambda\{|f_n|>\eta\}\to0$
for each $\eta>0$.  Uniform absolute continuity of the variations
then gives
\[
 |\mu_n(f_n)|
 \le\eta M+C|\mu_n|\{|f_n|>\eta\}\longrightarrow\text{at most }\eta M.
\]
Let $\eta\downarrow0$.  This proves the required pairing tends
to zero, and the testing criterion proves the first assertion.

Let now $x_n\rightharpoonup0$ in $L^1(\mu)$ and
$\ell_n\rightharpoonup0$ in its Banach dual.
There is a sigma-finite measurable set $S$ supporting every $x_n$:
take the union over $n,k$ of $\{|x_n|>1/k\}$, each of finite
measure by the integral bound.  Restrict $\ell_n$ to the closed
subspace $L^1(S)$ and represent these restrictions by
$g_n\in L^\infty(S)$.  The restriction and this dual
identification are bounded linear maps, so $g_n$ are weakly
null in $L^\infty(S)$ and uniformly bounded, say by $C$.
This does not assert the false unrestricted endpoint
identification $L^1(\mu)'=L^\infty(\mu)$ on every measure space.
Also $x_n$ are weakly null in $L^1(S)$, as can be seen by
extending each bounded scalar test on $S$ by zero.

The $L^1$ gliding-hump lemma of Part IV supplies uniform
integrability and finite-measure tail control.  Fix $\eps>0$;
choose finite-measure $A\subset S$ with
$\sup_n\int_{S\setminus A}|x_n|<\eps$.
Boundedness of $\|x_n\|_1$ and uniform absolute continuity then
give $R<\infty$ with
$\sup_n\int_{\{|x_n|>R\}}|x_n|<\eps$: the measures of these
sets are bounded by $\sup\|x_n\|_1/R$.
The preceding lemma yields $\int_A|g_n|\to0$, and hence
\[
 |\ell_n(x_n)|
 =\left|\int_S x_ng_n\,d\mu\right|
 \le2C\eps+R\int_A|g_n|\,d\mu\longrightarrow\text{at most }2C\eps.
\]
Letting $\eps\downarrow0$ proves the test and therefore the
Dunford--Pettis property for arbitrary $\mu$.
\end{proof}

% END INPUT


% BEGIN parts/part2-algebra.tex
\chapter{Banach algebras and holomorphic calculus}
All algebras in this chapter are complex. A unital Banach algebra here
is nonzero and has
$\|1\|=1$. For an operator our convention is
$R(\lambda,A)=(\lambda I-A)^{-1}$. Nonunital spectra are taken in the
unitization. We use the scalar Cauchy theorem and Runge approximation in
their usual precise forms from complex analysis, rather than treating
these as additional functional-analytic axioms. See \cite{B02,B04,B08}.

\begin{proposition}[Unitization and quotient]\label{fa2:p2:unitization}
If $\mathcal A$ is a Banach algebra, then
$\widetilde{\mathcal A}=\mathcal A\oplus\mathbb C$ with product
$(a,\lambda)(b,\mu)=(ab+\lambda b+\mu a,\lambda\mu)$ and norm
$\|a\|+|\lambda|$ is a unital Banach algebra. A closed two-sided ideal
$I$ has Banach quotient $\mathcal A/I$ with its quotient norm, and
multiplication is well defined and submultiplicative.
\end{proposition}
\begin{proof}
Associativity follows by expansion, and $(0,1)$ is an identity.
Submultiplicativity follows from the triangle inequality applied to
the four product terms. A Cauchy sequence is Cauchy in each component,
which proves completeness. Two representatives differing by elements
of $I$ have products differing by an element of $I$. For representatives
$a+i,b+j$, $\|(a+i)(b+j)\|\leq\|a+i\|\|b+j\|$; taking infima proves
the quotient norm inequality. Completeness of a quotient by a closed
subspace follows by choosing a subsequence with summable successive
quotient distances and successively correcting its representatives by
elements of $I$, obtaining a Cauchy sequence in $\mathcal A$.
\end{proof}

\begin{proposition}[Surjectivity is open]\label{fa2:p2:surjective-perturbation}
If $T:X\to Y$ is a bounded surjection between Banach spaces, every
sufficiently small norm perturbation is surjective. No complemented
kernel or bounded linear right inverse is assumed. A boundary point
of the spectrum of a bounded operator is never a surjective value.
\end{proposition}
\begin{proof}
Open mapping gives $C$ such that every $y$ has a lift $Tx=y$ with
$\|x\|\leq C\|y\|$. If $C\|B\|<1$, successively lift the
residuals: start $x_0=0$, take $Tu_n=y-(T+B)x_n$, and put
$x_{n+1}=x_n+u_n$. The next residual is $-Bu_n$, so it decreases
by the factor $C\|B\|$. The sum of the lifts converges, with norm
at most $C\|y\|/(1-C\|B\|)$, and its limit solves $(T+B)x=y$.
Thus the nearby surjections have a common lift bound. If
$\lambda\in\partial\sigma(A)$ and $\lambda-A$ were surjective,
nearby resolvent points would have inverses bounded by that lift
bound (their lift is unique). Resolvent norms must instead diverge
on approaching a spectral boundary, as proved above. This
contradiction proves the last assertion.
\end{proof}

\begin{proposition}[Characters in concrete function algebras]\label{fa2:p2:function-algebra-characters}\label{fa2:p2:stone-cech}
Characters of $C(K)$, $K$ compact Hausdorff, are point evaluations.
For completely regular Hausdorff $S$, the character space of $C_b(S)$
is its Gelfand compactification: point evaluations embed $S$ densely,
and each bounded continuous scalar function extends uniquely.
It has the Stone--Cech universal property for continuous maps into
compact Hausdorff spaces.
The characters of the disk algebra and of the polydisk algebra are
evaluations on the closed disk and closed polydisk respectively.
For empty $K$ or $S$, the corresponding function algebra is the zero
algebra and its character space is empty; the compactification assertions
then have their empty-space interpretation.
\end{proposition}
\begin{proof}
A maximal ideal $M\subset C(K)$ has a common zero. Otherwise finitely
many $f_j\in M$ would have no common zero by compactness; then
$\sum_j\bar f_jf_j\in M$ would be strictly positive and invertible,
impossible for a proper ideal. A common zero $x$ makes
$M\subset\{f:f(x)=0\}$, and maximality gives equality.
For $C_b(S)$, point evaluations are distinct and recover the topology
because $S$ is completely regular. They therefore give a topological
embedding in its compact character space. The commutative C*-theorem
identifies $C_b(S)$ isometrically with all continuous functions on
that space. If the evaluation image were not dense, a nonzero
continuous function vanishing on its closure would correspond to a
member of $C_b(S)$ vanishing at every point, a contradiction.
The Gelfand transform supplies each extension, and density gives
uniqueness. This is the compactification realization rather than an
assertion that the original $S$ was already compact.
For the universal property, a continuous $f:S\to K$ induces the
unital *-homomorphism $C(K)\to C_b(S)$, $g\mapsto g\circ f$.
Composing a character $\chi$ of $C_b(S)$ with this map gives a
character of $C(K)$, hence a unique point $F(\chi)\in K$ by the
first assertion. For every $g\in C(K)$,
$g(F(\chi))=\chi(g\circ f)$ is continuous in $\chi$.
Continuous scalar functions determine the compact Hausdorff topology,
so $F$ is continuous. On evaluation characters it equals $f$;
density makes this extension unique. This proves the universal
property with the precise completely regular Hausdorff hypothesis.
For the polydisk algebra, a character is determined by coordinate
values $\lambda_j$, which lie in the closed unit disks by its norm
one bound. Polynomials are uniformly dense: for $0<r<1$, $f(rz)$
has a Taylor series converging uniformly on the closed polydisk,
and tends uniformly to $f$ as $r\uparrow1$ by uniform continuity.
Thus the character is evaluation at $\lambda$. Every such evaluation
is a character. The one-dimensional disk case is included.
\end{proof}

\begin{proposition}[Shifts and adjoint spectra]\label{fa2:p2:shifts}
On $\ell^p(\mathbb N)$, $1\leq p<\infty$, the forward shift $S$
has spectrum the closed unit disk, empty point spectrum, and
approximate point spectrum the unit circle. The backward shift has
the same spectrum and point spectrum the open unit disk.
For a bounded complex Banach-space operator $T$,
$\sigma(T')=\sigma(T)$; for a Hilbert adjoint,
$\sigma(T^*)=\overline{\sigma(T)}$.
\end{proposition}
\begin{proof}
The shift norms are one, so the spectra lie in the closed disk by
the Neumann series. The forward eigenvector recurrence starts with
$0=\lambda x_1$ and forces every coordinate zero, also at $\lambda=0$
by injectivity. For $|\lambda|<1$,
$\|(S-\lambda)x\|\geq(1-|\lambda|)\|x\|$, so it is bounded below.
Its range is not dense: the nonzero bounded functional
$f(x)=\sum_{n\geq1}\lambda^{n-1}x_n$ satisfies $fS=\lambda f$.
Thus all interior points are spectral but not approximate points.
For $|\lambda|=1$, take the first $N$ entries
$1,\lambda^{-1},\ldots,\lambda^{-(N-1)}$ and normalize in $\ell^p$.
Only the two endpoint coordinates of $(S-\lambda)x$ remain, of norm
$2^{1/p}N^{-1/p}\to0$. This proves the assertions for $S$.
The backward recurrence gives eigenvectors
$(1,\lambda,\lambda^2,\ldots)$ exactly when $|\lambda|<1$.
The same finite truncations give approximate vectors on the circle,
so its full spectrum is the disk.
If $T$ is invertible so is $T'$, with inverse $(T^{-1})'$.
Conversely invertibility of $T'$ gives injectivity of $T$ with closed
range by the Banach closed-range theorem, and density of its range
because $\ker T'=0$. Thus $T$ is invertible. Apply this to
$T-\lambda I$ for the Banach spectrum identity. Under Hilbert Riesz
representation the scalar is conjugated in the adjoint, giving the
last formula.
\end{proof}

\begin{theorem}[Spectrum and spectral radius]\label{fa2:p2:spectrum-radius}
For $a$ in a nonzero unital $\mathcal A$, the spectrum $\sigma(a)$ is nonempty and compact,
contained in $\{|z|\leq\|a\|\}$. Its resolvent is analytic and
\[
 R(\lambda,a)-R(\mu,a)=(\mu-\lambda)R(\lambda,a)R(\mu,a),\qquad
 r(a)=\lim_{n\to\infty}\|a^n\|^{1/n}=\inf_n\|a^n\|^{1/n}.
\]
In particular complex Banach division algebras are $\mathbb C$.
\end{theorem}
\begin{proof}
If $\|b\|<1$, the norm-convergent series $\sum_{n\geq0}b^n$ is a
two-sided inverse for $1-b$. Hence invertibles form an open set, and
$R(\lambda,a)=\sum_{n\geq0}a^n/\lambda^{n+1}$ for $|\lambda|>\|a\|$.
The algebraic identity $u^{-1}-v^{-1}=u^{-1}(v-u)v^{-1}$ gives the
resolvent identity and its local power series, proving analyticity.
If the spectrum were empty, for each $f\in\mathcal A'$ the scalar
entire function $f(R(\lambda,a))$ would be bounded, tend to zero at
infinity, and hence vanish by Liouville. Hahn--Banach would make the
resolvent zero, an impossibility.
Submultiplicativity shows that $\lim\|a^n\|^{1/n}$ exists and equals
the infimum: writing $n=qk+s$ bounds its limsup by $\|a^k\|^{1/k}$.
Call this limit $r_0$. If $|\lambda|>r_0$ the same geometric series
converges by the root test, so $r(a)\leq r_0$.
For $R>r(a)$, Cauchy's coefficient formula for the resolvent gives
$a^n=(2\pi i)^{-1}\int_{|z|=R}z^nR(z,a)\,dz$; first verify it for
$R>\|a\|$ by the series and then deform the circle through the resolvent
set. Thus $\|a^n\|\leq R^{n+1}\max_{|z|=R}\|R(z,a)\|$ and
$r_0\leq R$. Let $R\downarrow r(a)$.
Finally in a division algebra choose $\lambda\in\sigma(a)$.
The noninvertible element $a-\lambda1$ must be zero, so every element
is scalar. The normalized identity makes this identification isometric.
\end{proof}

\begin{proposition}[Weak and strong analyticity]\label{fa2:p2:banach-analytic}
A function $f:U\to X$, $U\subset\mathbb C$ open and $X$ Banach, is
norm analytic iff $x'f$ is scalar analytic for every $x'\in X'$.
\end{proposition}
\begin{proof}
Only the converse needs proof. On a closed disk $D\Subset U$, every
$x'f$ is bounded; uniform boundedness applied to $\{J_Xf(z):z\in D\}$
gives $\sup_D\|f(z)\|<\infty$. For a smaller concentric disk, scalar
Cauchy estimates bound
$\|f(z)-f(w)\|\leq C|z-w|$ by Hahn--Banach. Thus $f$ is norm
continuous, so its contour integrals exist as Banach-valued Riemann
integrals. Scalar Cauchy's formula and separation give the vector
Cauchy formula. Expanding its kernel uniformly on a still smaller disk
gives a norm-convergent power series for $f$, proving norm analyticity.
\end{proof}

\begin{theorem}[Riesz--Dunford calculus]\label{fa2:p2:holomorphic-calculus}
Let $f$ be holomorphic on a neighborhood $U$ of $\sigma(a)$. Choose a
finite positively oriented contour cycle $\Gamma\subset U\setminus\sigma(a)$
whose winding number is one near the spectrum and zero off $U$. Then
\[
 f(a)=\frac1{2\pi i}\int_\Gamma f(z)R(z,a)\,dz
\]
is independent of the cycle and the chosen neighborhood, is a unital
algebra homomorphism in $f$, agrees with polynomials and resolvents,
and satisfies $\sigma(f(a))=f(\sigma(a))$. Uniform convergence on a
fixed contour implies norm convergence of the values. Every bounded
operator commuting with $a$ commutes with every $f(a)$.
\end{theorem}
\begin{proof}
A cycle can be obtained as the oriented boundary of finitely many
sufficiently small grid squares covering the compact spectrum, with
internal edges canceled. Independence follows by applying scalar
Cauchy's theorem to every functional of the analytic integrand between
the cycles and then separating points. For $f=1,z$ the value is $1,a$,
as follows from the resolvent series on a large circle; linearity and
independence allow contours appropriate to their domains.
For multiplicativity take nested cycles $\Gamma_1,\Gamma_2$ around
the spectrum and use
$R(z,a)R(w,a)=(R(z,a)-R(w,a))/(w-z)$.
Fubini is valid for continuous functions on the two compact curves.
In the double integral one Cauchy integral is $f(w)$ and the other is
zero, giving $f(a)g(a)=(fg)(a)$. This also proves polynomial agreement
and the inverse identity for $(\lambda-z)^{-1}$.
If $\mu\notin f(\sigma(a))$, shrink $U$ until $f-\mu$ has no zeros
there; applying the calculus to $1/(f-\mu)$ gives an inverse of
$f(a)-\mu$. Conversely if $\lambda\in\sigma(a)$ and $f(\lambda)=\mu$,
write $f(z)-\mu=(z-\lambda)g(z)$ on $U$, defining the removable value
at $\lambda$. If $f(a)-\mu$ were invertible, the two commuting factors
$a-\lambda$ and $g(a)$ would each be invertible: for commuting $b,c$,
an inverse of $bc$ commutes with $b,c$ and $c(bc)^{-1}$ is a two-sided
inverse of $b$. This contradicts $\lambda\in\sigma(a)$.
The norm convergence assertion follows by bounding the contour integral
by its length times the uniform difference and the maximum resolvent
norm. Commutation follows term by term under the integral.
\end{proof}
\pagebreak[4]
\begin{theorem}[Riesz projections and splitting]\label{fa2:p2:riesz-projection}
If $T\in\faB(X)$ and $\sigma(T)=K_1\mathbin{\dot\cup}K_2$ is a
disjoint union of compact sets, let $f=1$ near $K_1$ and $f=0$ near
$K_2$. Then $P=f(T)$ is a bounded projection commuting with $T$,
$X=PX\oplus(I-P)X$, and
$\sigma(T|_{PX})=K_1$, $\sigma(T|_{(I-P)X})=K_2$ (omit a zero
summand). For an isolated spectral point $\lambda$, take $K_1=\{\lambda\}$
in this split. Then $\lambda$ is a resolvent pole of order $m$ iff $T-\lambda$
on $PX$ is nilpotent of order $m$. On $(I-P)X$ it is invertible.
\end{theorem}
\begin{proof}
$f^2=f$ and multiplicativity give $P^2=P$; ranges and kernels of a
bounded projection are closed and give the decomposition. On $PX$
the calculus gives an inverse to $z-T$ for $z\notin K_1$: choose the
function equal to $(z-w)^{-1}$ near $K_1$ and zero near $K_2$.
Thus the restricted spectra lie in the corresponding $K_j$.
The spectrum of a diagonal two-block operator is the union of the
two block spectra (a bounded inverse exists exactly when both block
inverses exist). Since the union is $\sigma(T)$, equality follows.
For an isolated point use the split $K_1=\{\lambda\}$ and expand
$R(z,T)$ around $\lambda$. The Laurent coefficients are
$(T-\lambda)^kP$ for $(z-\lambda)^{-k-1}$, by the geometric expansion
first on a circle outside the restricted spectral radius, then by
uniqueness of Laurent coefficients. The series terminates exactly
when the restricted operator is nilpotent, with the stated order.
\end{proof}

\begin{proposition}[Permanence and polynomial hull]\label{fa2:p2:polynomial-hull}
Let $\mathcal B\subset\mathcal A$ be a closed unital subalgebra with the
same identity and $a\in\mathcal B$. Then
$\sigma_{\mathcal A}(a)\subset\sigma_{\mathcal B}(a)$,
$\partial\sigma_{\mathcal B}(a)\subset\partial\sigma_{\mathcal A}(a)$,
and their polynomial hulls agree. Each component of
$\mathbb C\setminus\sigma_{\mathcal A}(a)$ either lies entirely in
$\sigma_{\mathcal B}(a)$ or misses it entirely. In the closed algebra
generated by $1,a$ the spectrum is exactly the polynomial hull of the
ambient spectrum.
\end{proposition}
\begin{proof}
An inverse in $\mathcal B$ is an inverse in $\mathcal A$, proving the
inclusion. The set of $z$ in an ambient resolvent component for which
$R(z,a)\in\mathcal B$ is open by its local geometric series and closed
by norm continuity and closedness of $\mathcal B$. It is therefore all
or none of that component. It contains the unbounded component by the
large-$z$ series. If $z$ belongs to the boundary of the smaller-algebra
spectrum but is in the ambient resolvent, that same open-or-closed
argument would put a neighborhood wholly inside or outside the
smaller-algebra spectrum, a contradiction. In one complex variable
the polynomial hull fills precisely the bounded complementary
components; the preceding facts prove equality of hulls.
For the singly generated algebra put $K=\sigma_{\mathcal A}(a)$.
If $z\in\widehat K$, then for every polynomial
$|p(z)|\leq\max_K|p|=r(p(a))\leq\|p(a)\|$. Evaluation at $z$
therefore defines a well-defined bounded character on the polynomial
algebra in $a$, and extends to its closure; hence $z$ is in its spectrum.
If $z\notin\widehat K$, choose a polynomial separating $z$ from $K$;
the same strict separation holds for a sufficiently small compact
neighborhood $L$ of $K$, so $z\notin\widehat L$. Runge approximation
on $\widehat L$ approximates $(z-w)^{-1}$ uniformly by polynomials.
Choose a fixed contour for $K$ inside $L$. Uniform approximation there,
followed by the contour norm estimate, gives convergence of those
polynomials in $a$ to $R(z,a)$. Thus this inverse belongs to the
generated algebra. Both inclusions prove the assertion.
\end{proof}

\begin{theorem}[Characters and the Gelfand transform]\label{fa2:p2:gelfand}
Let $\mathcal A$ be a commutative unital Banach algebra. Its characters
are precisely the quotient maps onto $\mathbb C$ associated with
maximal ideals. Every character has norm one. The character space
$\Delta$ is compact Hausdorff in the weak-star topology, and
\[
 \sigma(a)=\{\chi(a):\chi\in\Delta\},\quad
 \widehat a(\chi)=\chi(a),\quad
 \|\widehat a\|_\infty=r(a).
\]
The Gelfand transform is a contractive unital homomorphism into
$C(\Delta)$ whose kernel is the intersection of maximal ideals.
\end{theorem}
\begin{proof}
A proper ideal cannot meet the open set of invertibles. Its closure
is still proper, since an element within distance one of the identity
is invertible. Thus maximal ideals, supplied by Zorn, are closed.
The quotient by one is a Banach division algebra and hence $\mathbb C$.
Conversely a nonzero scalar homomorphism is unital and has maximal
kernel. For a character $\chi$, $\chi(a)\in\sigma(a)$ since applying
$\chi$ to an inverse of $a-\chi(a)$ would give $1=0$.
Thus $|\chi(a)|\leq\|a\|$ and $\|\chi\|=1$.
The equations $\chi(1)=1$ and $\chi(ab)=\chi(a)\chi(b)$ define a
weak-star closed subset of the dual unit ball. Alaoglu proves
compactness; the weak-star topology is Hausdorff.
If $\lambda\in\sigma(a)$ the ideal generated by $a-\lambda$ is proper,
so a maximal ideal containing it gives a character with value
$\lambda$. This proves the spectral formula. Continuity of
$\widehat a$ is the definition of weak-star convergence; the remaining
algebra, norm and kernel assertions follow directly.
\end{proof}

\begin{theorem}[Gleason--Kahane--Zelazko]\label{fa2:p2:gkz}
If a linear $\varphi:\mathcal A\to\mathbb C$ satisfies $\varphi(1)=1$
and $\varphi(a)\ne0$ whenever $a$ is invertible, then $\varphi$ is a
continuous character. Continuity is a conclusion, not a hypothesis.
\end{theorem}
\begin{proof}
$a-\varphi(a)1$ is not invertible, so $|\varphi(a)|\leq\|a\|$.
For fixed $a$, $F(z)=\varphi(e^{za})$ is zero-free entire, $F(0)=1$,
and $|F(z)|\leq e^{\|a\||z|}$. Write $F=e^h$, $h(0)=0$, by integrating
$F'/F$ on the simply connected plane. The growth bound forces $h$ to
be affine. Here is the required analytic estimate: if $h$ is analytic
on $|z|\leq R$ and $M=\max\operatorname{Re}(h-h(0))$, the map
$(h-h(0))/(2M-h+h(0))$ has modulus at most one and vanishes at zero.
Schwarz's lemma gives
$\max_{|z|\leq r}|h(z)-h(0)|\leq2Mr/(R-r)$.
With $r=R/2$ and $M\leq\|a\|R+|h(0)|$, Cauchy's estimate gives
$|h''(0)|\leq C/R$. Apply the same argument centered at any point
and let $R\to\infty$. Thus $h''=0$.
Comparing the first two derivatives of $F$ at zero gives
$\varphi(a^2)=\varphi(a)^2$. Polarization yields
$\varphi(ab+ba)=2\varphi(a)\varphi(b)$.
If $\varphi(x)=0$, put $u=xy-yx$, $v=xy+yx$ and use
$u^2+v^2=2[x(yxy)+(yxy)x]$.
Applying the square and Jordan identities gives
$\varphi(u)^2+\varphi(v)^2=0$ and $\varphi(v)=0$; hence
$\varphi(u)=0$ and $\varphi(xy)=0$.
Take $x=a-\varphi(a)1$ to conclude
$\varphi(ab)=\varphi(a)\varphi(b)$.
\end{proof}

\begin{proposition}[Exponential components and logarithms]\label{fa2:p2:exponentials}
The identity component of the invertible group of a unital Banach
algebra consists exactly of finite products of exponentials. If
$\sigma(a)$ lies in a domain carrying a holomorphic logarithm, then
$a=e^b$ with $b=(\log)(a)$; holomorphic roots follow in the same way.
For a complex Hilbert space every invertible operator is a product of
two exponentials, and its invertible group is norm path connected.
No assertion of Kuiper contractibility is made.
\end{proposition}
\begin{proof}
$t\mapsto e^{tb}$ is a path of invertibles. Conversely subdivide a
path from $1$ to $a$ so finely that consecutive quotients differ from
$1$ by less than one in norm. The convergent logarithm series gives
an exponential for each quotient; multiplying gives $a$.
The calculus homomorphism and composition, verified by the nested
contour argument, give the holomorphic logarithm and roots.
On Hilbert space use polar decomposition $T=U|T|$.
For invertible $T$, $U$ is unitary and $|T|$ is positive invertible.
The bounded Borel calculus gives $U=e^{iB}$ for a bounded self-adjoint
$B$ using an argument function with values in $[-\pi,\pi]$;
the continuous calculus gives $|T|=e^C$, $C=\log|T|$ self-adjoint.
Then $t\mapsto e^{itB}e^{tC}$ is the required norm-continuous path.
\end{proof}

\begin{theorem}[Group-algebra characters and Wiener]\label{fa2:p2:wiener}
For a locally compact abelian group $G$ with Haar measure,
$L^1(G)$ is a commutative Banach algebra under convolution. Its
nonzero characters are exactly
$f\mapsto\int_G f(x)\overline{\gamma(x)}\,dx$ for continuous
homomorphisms $\gamma:G\to\mathbb T$.
In particular if $\sum_{n\in\mathbb Z}|a_n|<\infty$ and
$f(e^{it})=\sum a_ne^{int}$ never vanishes, then $1/f$ has an
absolutely convergent Fourier series.
\end{theorem}
\begin{proof}
Fubini and invariance of Haar measure give
$\|f*g\|_1\leq\|f\|_1\|g\|_1$ and associativity; commutativity
uses the abelian hypothesis. The integral against a character is
multiplicative by the same change of variables and has norm one.
Conversely let $\chi\ne0$ be an algebra character. It is bounded by
unitization and the spectral argument above. Write
$\tau_xf(t)=f(t-x)$ in additive notation, and choose $f_0$ with
$\chi(f_0)\ne0$. The identity
$(\tau_xf)*g=f*(\tau_xg)$ implies
$\chi(\tau_xf)=c(x)\chi(f)$, where
$c(x)=\chi(\tau_xf_0)/\chi(f_0)$.
Translation continuity in $L^1$ makes $c$ continuous;
the group law and invertibility of translations give
$c(x+y)=c(x)c(y)$ and $c(0)=1$.
Since all translations are isometric, $|c(x)|^n$ is bounded for
positive and negative integers $n$, so $|c(x)|=1$.
The Bochner identity $f*g=\int_G f(x)\tau_xg\,dx$ then gives
$\chi(f)\chi(g)=\chi(g)\int_G f(x)c(x)\,dx$.
Choose $g$ with $\chi(g)\ne0$ and put $\gamma=\overline c$.
For $G=\mathbb Z$ every character is $n\mapsto z^n$, $|z|=1$,
and $L^1(G)=\ell^1(\mathbb Z)$ is unital. Nonvanishing of all
character values makes $a$ invertible by the maximal-ideal criterion.
Its inverse $b\in\ell^1$ has Fourier series $1/f$, proving Wiener.
\end{proof}

% END INPUT


% BEGIN parts/part2-fredholm.tex
\chapter{Riesz theory, parametrices, and index}
The first theorem intentionally allows an incomplete normed space:
compactness means compact closure in that space, not merely in its
completion. Banach hypotheses are introduced when open mapping or
quotient completeness is used. See \cite{B02,B03,B06,B08}.

\begin{theorem}[Riesz--Schauder on normed spaces]\label{fa2:p2:riesz-schauder}
Let $K:X\to X$ be compact on a normed space and $A=I-K$. Then
$\ker A$ is finite dimensional, $\faRan A$ is closed, the ascending
kernels and descending ranges of $A^n$ stabilize, and for sufficiently
large $r$,
\[
 X=\ker A^r\oplus\faRan A^r.
\]
The restriction of $A$ to the second summand is a bounded isomorphism.
In particular $A$ is injective iff it is surjective; if bijective its
inverse is bounded. Over $\mathbb C$, every nonzero spectral value of
$K$ is an eigenvalue of finite algebraic multiplicity, and the only
possible accumulation point of its nonzero spectrum is zero.
\end{theorem}
\begin{proof}
On $N=\ker A$, $K$ is the identity, so its unit ball is relatively
compact. Riesz's lemma forces $N$ finite dimensional. There is $c>0$
with $\|Ax\|\geq c\operatorname{dist}(x,N)$. Otherwise choose
$\operatorname{dist}(x_n,N)=1$, $\|x_n\|\leq2$ and $Ax_n\to0$.
Compactness gives a convergent subsequence of $Kx_n$, hence of
$x_n=Ax_n+Kx_n$. Its limit lies in $N$, contradicting distance one.
If $Ax_n\to y$, subtract elements of $N$ to arrange
$\|x_n\|\leq2c^{-1}\|Ax_n\|+1/n$. Again a subsequence of $Kx_n$
converges, so $x_n$ converges in $X$, and continuity gives $y\in\faRan A$.
Each $A^n$ is $I$ minus a compact operator, so its range is closed.
If its kernels increase strictly forever, choose unit
$x_n\in\ker A^{n+1}$ at distance at least $1/2$ from $\ker A^n$.
For $m<n$, $Ax_n+Kx_m\in\ker A^n$; thus
$\|Kx_n-Kx_m\|\geq1/2$, a contradiction. If its ranges decrease
strictly forever, choose unit $x_n\in\faRan A^n$ at distance at least
$1/2$ from $\faRan A^{n+1}$. For $m>n$,
$Ax_n+Kx_m\in\faRan A^{n+1}$, giving the same contradiction.
Once equality occurs the corresponding chain stabilizes by applying
$A$ or taking an inverse image. Choose a common stabilization index $r$.
If $x=A^ry\in\ker A^r$, then $y\in\ker A^{2r}=\ker A^r$, so $x=0$.
For arbitrary $x$, $A^rx\in\faRan A^r=\faRan A^{2r}$, so
$A^rx=A^{2r}y$ for some $y$; hence
$x=(x-A^ry)+A^ry$ is the asserted direct sum.
The finite-dimensional first summand and closed second summand give
bounded projections: an unbounded projection would produce unit
vectors in the finite-dimensional summand approaching the closed
complement, contradicting their zero intersection. The restriction of
$A$ to the second summand is algebraically bijective, and the earlier
distance estimate, or its proof on this summand, bounds its inverse.
On the first summand $A$ is nilpotent; finite-dimensional rank-nullity
now proves injectivity iff surjectivity and index zero.
Apply this to $I-K/\lambda$, $\lambda\ne0$. Failure of invertibility
means a nonzero finite-dimensional generalized eigenspace. If distinct
nonzero eigenvalues $\lambda_n$ stayed away from zero, their eigenvectors
span a strictly increasing chain $E_n$. Choose unit $x_n\in E_n$ at
distance $1/2$ from $E_{n-1}$. Since
$Kx_n-\lambda_nx_n\in E_{n-1}$, the vectors $Kx_n/\lambda_n$ are
pairwise separated modulo earlier spaces; boundedness away from zero
and a convergent subsequence of $\lambda_n$ contradict compactness.
Thus there are only finitely many eigenvalues outside every disk
about zero. Stabilization supplies finite algebraic multiplicity.
\end{proof}

\begin{theorem}[Regularization in general dual systems]\label{fa2:p2:dual-regularization}
Let $(X_j,Y_j)$ be nondegenerate bilinear dual systems of normed
spaces, without an assumed bound on the pairings. Suppose bounded
$T:X_1\to X_2$, $T':Y_2\to Y_1$ are paired adjoints and bounded
$R:X_2\to X_1$, $R':Y_1\to Y_2$ are paired adjoints. Assume all
four defects $I-RT,I-TR,I-T'R',I-R'T'$ compact. Then the kernels
of $T,T'$ and their cokernels are finite dimensional and
\[
 \faRan T=\{f\in X_2:[f,y]=0\ (y\in\ker T')\},\quad
 \operatorname{codim}\faRan T=\dim\ker T'.
\]
The analogous assertion holds for $T'$.
For a left regularizer, the regularized equation has exactly the
original solutions for every right-hand side iff the regularizer
is injective. For a right regularizer, its solutions map onto all
original solutions for every right-hand side iff it is surjective.
\end{theorem}
\begin{proof}
$\ker T\subset\ker(RT)$ is finite dimensional by Riesz theory.
The normalized-distance compactness proof used in Atkinson shows
closedness of the range using only the compact defect and normed
space compactness, without open mapping or completeness.
Put $B=I-TR$, $B'=I-R'T'$, a compact adjoint pair on $X_2,Y_2$.
For a common stabilization exponent $r$, Riesz decompositions give
$X_2=E\dotplus M$, with $E=\ker(TR)^r$ finite dimensional and
$M=\faRan(TR)^r\subset\faRan T$, and similarly $Y_2=E'\dotplus M'$.
The pairings $E\times M'$ and $M\times E'$ vanish by the adjoint
identity; the finite-dimensional pairing $E\times E'$ is
nondegenerate, exactly as in the dual Fredholm proof.
Thus their projections are paired adjoints.
Write $\faRan T=M\dotplus W$ with $W=E\cap\faRan T$.
Every point outside this range has an $E$ component outside $W$,
so finite-dimensional duality supplies a $y\in E'$ annihilating
$W$ but not that component. It also annihilates $M$, hence all
$\faRan T$, which means $T'y=0$. This proves the annihilator
identity and, by finite-dimensional rank-nullity in $E,E'$,
the codimension formula. It also proves the finiteness assertions;
interchanging the two systems gives the other direction.
For left regularization, $RTx=Rf$ is equivalent to
$Tx-f\in\ker R$, so injectivity suffices. If $0\ne f\in\ker R$,
$x=0$ is a spurious solution of the regularized equation, proving
necessity. For right regularization $TRz=f$ sends $z$ to the
original solution $x=Rz$. Surjectivity of $R$ gives every original
solution in this way. Conversely any $x\notin\faRan R$ is an
original solution for $f=Tx$ but cannot be so represented.
This proves both exact equivalence criteria.
\end{proof}

\begin{proposition}[Algebraic pseudoinverses]\label{fa2:p2:algebraic-pseudoinverse}
For any linear map of vector spaces $T:X\to Y$, choices of algebraic
complements to its kernel and range give a linear $S:Y\to X$ with
$TST=T$ and $STS=S$. Its defects $I-ST,I-TS$ project onto the
kernel and a range complement. If these have finite dimensions,
their difference is the algebraic index, additive under products.
Equivalently a map has finite kernel and cokernel iff it has an
inverse modulo finite-rank defects (the finite-index convention for
``pseudoinverse'' in \cite{B08}). Finite-rank perturbations preserve
this algebraic index without any topology on the spaces.
No boundedness is asserted for these arbitrary algebraic complements.
\end{proposition}
\begin{proof}
Choose $X=\ker T\dotplus X_0$ and $Y=\faRan T\dotplus C$ by bases.
Set $S=(T|_{X_0})^{-1}$ on the range and zero on $C$.
Substitution proves all identities. When kernel and cokernel are
finite, the exact sequence in \ref{fa2:p2:index}, which is purely
algebraic, gives product additivity by alternating dimensions.
Conversely $ST=I+F$, $TS=I+G$ with finite-rank $F,G$ give
$\ker T\subset\ker(I+F)\subset\faRan F$ and
$\faRan T\supset\faRan(I+G)\supset\ker G$.
Thus both defects are finite dimensional, since
$Y/\ker G\cong\faRan G$ algebraically. The same parametrix
still works after a finite-rank perturbation $D$.
Put $E=\ker D$, of finite codimension. Restrictions of $T$ and
$T+D$ to $E$ agree. The inclusion $E\to X$ has index
$-\operatorname{codim}E$, so product additivity gives
$\faind(T|_E)=\faind T-\operatorname{codim}E$ and the identical
formula for $T+D$. Equality of the restrictions proves algebraic
finite-rank invariance.
\end{proof}

\begin{theorem}[Matrix Toeplitz symbols and a discontinuous counterexample]\label{fa2:p2:matrix-toeplitz}
For continuous $\Phi:\mathbb T\to GL_d(\mathbb C)$, its Toeplitz
operator on $H^2(\mathbb T;\mathbb C^d)$ is Fredholm and
$\faind T_\Phi=-\operatorname{wind}(\det\Phi,0)$.
Nonvanishing bounded measurable scalar symbols need not give Fredholm
operators: the unimodular boundary symbol of
$\theta(z)=\exp(-(1+z)/(1-z))$ is one explicit example.
\end{theorem}
\begin{proof}
The entrywise Hankel compactness proof gives
$T_\Phi T_\Psi-T_{\Phi\Psi}$ compact; $\Psi=\Phi^{-1}$ is a
parametrix. To compute its index we give the required loop homotopy
instead of merely citing its topological classification.
Polar decomposition deforms a continuous matrix loop to a unitary
loop. For $d\geq2$ its first column is a loop in $S^{2d-1}$.
Every loop in a sphere of dimension at least two contracts:
uniformly approximate it by a finite short-geodesic polygon, use
normalized straight interpolation for the close-loop homotopy, and
choose a point outside the finitely many great circles containing
the polygon. Stereographic projection from that point contracts
the polygon in Euclidean space. The final constant may be moved
to the first coordinate vector along a path.
This column homotopy lifts to a unitary matrix homotopy. Here is
an explicit local lift. For close unit vectors $v,w$ put
$P_v=v\otimes v$, $K=P_wP_v+(I-P_w)(I-P_v)$.
For sufficiently close $v,w$, $K$ is invertible and its polar
unitary sends $v$ to $w\langle w,v\rangle/|\langle w,v\rangle|$.
Multiplying on $\mathbb Cw$ by the inverse of that phase gives a
continuous local unitary $L(v,w)$ with $L(v,w)v=w$ and $L(v,v)=I$.
Uniform continuity on the compact homotopy square allows a finite
subdivision of its time variable; multiply these local lifts on
successive slabs. They preserve periodicity of the loop parameter.
At the end the first column is constant, so the loop has block form
$1\oplus U_{d-1}$. Inductively it deforms to
$1\oplus\cdots\oplus u(z)$ with scalar unitary $u$.
Its winding equals that of the original determinant, since determinant
is continuous throughout. Removing $z^k$, $k=\operatorname{wind}u$,
leaves a scalar loop with continuous logarithm, which contracts.
Thus $\Phi$ deforms through invertible symbols to a diagonal loop
with one entry $z^k$. Toeplitz operators along the path are Fredholm;
index constancy and the shift calculation give $-k$.
For the counterexample, $\theta$ is bounded analytic in the disk
since $\operatorname{Re}(1+z)/(1-z)>0$. Its boundary values have
modulus one except at $z=1$. Multiplication by it is therefore an
isometry of $H^2$ onto the closed subspace $\theta H^2$.
Suppose this range had finite codimension. Its orthogonal complement
would be finite dimensional and invariant under the backward shift.
A nonzero annihilating polynomial $q$ for that finite-dimensional
restriction would give a nonzero polynomial $p\in\theta H^2$:
if $S$ is the forward shift, take $p=\overline q(S)1$,
where the bar conjugates the coefficients. Indeed $q(S^*)h=0$
for every $h\perp\theta H^2$, so adjointness makes
$\faip{h}{p}=0$ for every such $h$.
Then $p/\theta\in H^2$, so the evaluation estimate
$|p(r)/\theta(r)|\leq C(1-r^2)^{-1/2}$ would hold for $0<r<1$.
But $1/\theta(r)=\exp((1+r)/(1-r))$, which grows faster than every
power of $(1-r)^{-1}$, whereas a nonzero polynomial has only a finite
order zero at $1$. This is impossible. The range has infinite
codimension, so the unimodular measurable symbol is not Fredholm.
Thus the continuous-symbol winding theorem cannot be extended to
arbitrary discontinuous symbols by dropping continuity.
\end{proof}

\begin{theorem}[Fredholm alternative in dual systems]\label{fa2:p2:dual-fredholm}
Let $X,Y$ be normed spaces with a nondegenerate bilinear pairing
$[\ ,\ ]:X\times Y\to\mathbb C$; boundedness of the pairing is not
required. Let compact $K:X\to X$ and $L:Y\to Y$ satisfy
$[Kx,y]=[x,Ly]$. Then either both $I-K,I-L$ are bijective, or their
kernels have the same positive finite dimension. In every case
\[
 \faRan(I-K)=\{f:[f,y]=0\text{ for }y\in\ker(I-L)\},
\]
and the analogous identity with $X,Y$ interchanged holds.
The sesquilinear version follows by conjugating one scalar structure.
\end{theorem}
\begin{proof}
Put $A=I-K$, $B=I-L$ and choose a common stabilization exponent $r$.
The pairing between $\ker A^r$ and $\ker B^r$ is nondegenerate:
an $x\in\ker A^r$ annihilating $\ker B^r$ also annihilates
$\faRan B^r$, because $[x,B^ry]=[A^rx,y]=0$.
The Riesz decomposition of $Y$ then makes it annihilate all of $Y$,
so $x=0$. The other direction is identical. Both spaces are finite
dimensional; nondegeneracy gives equal dimensions and identifies
their nilpotent restrictions as transposes. Hence
$\dim\ker A=\dim\ker B$.
The Riesz decomposition and finite-dimensional rank-nullity show
$\operatorname{codim}\faRan A=\dim\ker A$.
Every vector in $\faRan A$ annihilates $\ker B$. Conversely the
annihilator of the latter has codimension $\dim\ker B$: the map
$X\to(\ker B)'$ is onto, since otherwise a nonzero element of
$\ker B$ would annihilate all of $X$. Thus the contained subspaces
have equal finite codimension and are equal. Interchanging the spaces
proves the other assertion. The injective-surjective alternatives now
follow from the preceding theorem.
\end{proof}

\begin{proposition}[Kernels, Neumann series, and Volterra]\label{fa2:p2:kernel-compactness}
Let $Q$ be compact Hausdorff with a finite regular measure $\mu$.
A continuous kernel $k\in C(Q\times Q)$ defines a compact operator
$K:C(Q)\to C(Q)$ by $Kf(x)=\int k(x,y)f(y)\,d\mu(y)$.
If $\|T\|<1$ on a Banach space, $(I-T)^{-1}=\sum_{n\geq0}T^n$.
For a continuous Volterra kernel on $a\leq y\leq x\leq b$,
$Vf(x)=\int_a^xk(x,y)f(y)\,dy$ satisfies
$\|V^n\|\leq M^n(b-a)^n/n!$. Hence $I-\lambda V$ is invertible
for every $\lambda\in\mathbb C$ and $\sigma(V)=\{0\}$ on an
infinite-dimensional space.
\end{proposition}
\begin{proof}
The unit-ball image is bounded by $\mu(Q)\|k\|_\infty$.
Uniform continuity of the continuous map $x\mapsto k(x,\cdot)$ into
$C(Q)$ gives equicontinuity; Arzela--Ascoli gives compactness.
The geometric-series identity follows by multiplying partial sums
and letting the remainder tend to zero. Iterating the Volterra integral
integrates over $a\leq y_n\leq\cdots\leq y_1\leq x$, a simplex of
volume $(x-a)^n/n!$, proving the bound. The series
$\sum\lambda^nV^n$ converges for every $\lambda$ and is an inverse.
Zero is in the spectrum since a compact invertible operator on an
infinite-dimensional normed space would make the identity compact.
\end{proof}

\begin{definition}
A bounded $T:X\to Y$ between Banach spaces is Fredholm if its range
is closed and its kernel and cokernel $Y/\faRan T$ are finite
dimensional. Its index is
$\faind T=\dim\ker T-\dim(Y/\faRan T)$.
\end{definition}

\begin{theorem}[Atkinson and parametrices]\label{fa2:p2:atkinson}
$T$ is Fredholm iff there is $S:Y\to X$ bounded with
$ST-I_X$ and $TS-I_Y$ compact. It is equivalent to require finite-rank
defects. Thus for $X=Y$ it is equivalent to invertibility of the class
of $T$ in $\faB(X)/\faK(X)$. In Hilbert spaces left semi-Fredholmness
is equivalent to closed range and finite-dimensional kernel; right
semi-Fredholmness is the adjoint assertion.
\end{theorem}
\begin{proof}
For Fredholm $T$, choose bounded complements
$X=N\oplus X_0$, $Y=\faRan T\oplus C$ with $N=\ker T$ and $C$
finite dimensional. Finite-dimensional subspaces admit bounded
projections by Hahn--Banach; finite-codimensional closed subspaces
do also. The inverse of $T|_{X_0}$ is bounded by open mapping.
Define $S$ to be this inverse on the range and zero on $C$.
Then $I-ST=P_N$, $I-TS=P_C$ are finite rank.
Conversely $ST=I+K$ implies $\ker T\subset\ker(I+K)$, so the
kernel is finite dimensional. If no lower bound by
$\operatorname{dist}(x,\ker T)$ held, choose normalized representatives
$x_n$ with distance one, norm at most two and $Tx_n\to0$.
Then $x_n=STx_n-Kx_n$ has a convergent subsequence, contradicting
distance one. This lower bound and the compactness of $K$ prove
closed range by the same convergent-range argument used above.
From $TS=I+L$, adjunction and Schauder give finite-dimensional
$\ker T'$. The dual of $Y/\faRan T$ is precisely $\ker T'$;
Hahn--Banach then forces the quotient finite dimensional of the
same dimension. This proves Fredholmness.
On a Hilbert space the inverse of $T:(\ker T)^\perp\to\faRan T$
extended by zero gives a left inverse modulo the projection onto the
finite kernel; the converse was just proved using only the left
parametrix. Taking adjoints proves the right assertion.
\end{proof}

\begin{theorem}[Index calculus and stability]\label{fa2:p2:index}
Products and direct sums of Fredholm operators are Fredholm,
$\faind(ST)=\faind S+\faind T$ and
$\faind(S\oplus T)=\faind S+\faind T$.
The Fredholm set is norm open and its index is locally constant.
Compact perturbations preserve both Fredholmness and index.
For Hilbert adjoints $\faind T^*=-\faind T$.
On Hilbert spaces the semi-Fredholm index, allowing $+\infty$ and
$-\infty$, is also locally constant on the semi-Fredholm set.
\end{theorem}
\begin{proof}
Multiplying parametrices gives a parametrix for a product; direct sums
give one for a sum. For the product use the exact finite-dimensional
sequence
\[
0\to\ker T\to\ker ST\xrightarrow{T}\ker S
\to Y/\faRan T\xrightarrow{S}Z/\faRan ST
\to Z/\faRan S\to0.
\]
The map from $\ker S$ takes the class modulo $\faRan T$; the last
map is the natural quotient. At $\ker S$ exactness says a vector is
$Tx$ precisely when it lies in $\faRan T$; at the first cokernel it
says $Sy\in\faRan ST$ iff $y$ differs from a vector of $\faRan T$
by an element of $\ker S$. The other exactness identities follow
directly from the definitions. Alternating dimensions gives the
product formula. Kernels and cokernels split for direct sums.
For openness write $T$ in the complements used in Atkinson as
$\left[\begin{smallmatrix}T_0&0\\0&0\end{smallmatrix}\right]$,
where $T_0:X_0\to\faRan T$ is invertible. A small perturbation has
upper-left block $A$ still invertible. Multiplication by invertible
triangular block matrices reduces it to
$\operatorname{diag}(A,D-CA^{-1}B)$, whose second block is between
the finite-dimensional spaces $N,C$. Its index is always
$\dim N-\dim C$. This proves openness and local constancy.
The Calkin class is unchanged by a compact perturbation, so every
$T+tK$, $0\leq t\leq1$, is Fredholm. Local constancy on this connected
path gives index invariance. The closed-range identities identify the
Hilbert cokernel with $\ker T^*$, proving the adjoint formula.
For a left semi-Fredholm Hilbert operator use the same block
decomposition with finite kernel and possibly infinite cokernel.
The invertible range block stays invertible; its Schur complement
maps the finite kernel into the old cokernel. If the latter is
infinite, removing a finite-dimensional range leaves infinite
dimension, so the index remains $-\infty$. If finite, the earlier
calculation applies. Adjunction handles the right case.
\end{proof}

\begin{proposition}[Index components on Hilbert space]\label{fa2:p2:index-components}
On a fixed complex Hilbert space two Fredholm operators belong to the
same norm-path component iff their indices agree.
\end{proposition}
\begin{proof}
Necessity is local constancy. If $A$ has index zero, add a finite-rank
map carrying its kernel isomorphically onto its cokernel, zero on the
complement. Then $A+F$ is invertible and the path $A+tF$ is Fredholm.
Invertibles are path connected by the two-exponential construction.
For $A,B$ of equal index choose a parametrix $C$ for $B$, so
$CB=I+K$. Then $\faind C=-\faind B$ and $AC$ has index zero.
Choose a Fredholm path $\gamma$ from $I$ to $AC$. The path
$\gamma(t)B-tAK$ starts at $B$ and ends at
$ACB-AK=A$; every point is a compact perturbation of a product of
Fredholm operators. This proves sufficiency.
\end{proof}

\begin{theorem}[Essential and Weyl spectra]\label{fa2:p2:essential-spectrum}
For $T\in\faB(H)$ on an infinite-dimensional real or complex Hilbert
space, let spectra and Fredholm indices in the following formula use
that original scalar field, and set
$\sigma_e(T)=\{\lambda:T-\lambda I\text{ is not Fredholm}\}$ and
$\sigma_w(T)=\bigcap_{K\in\faK(H)}\sigma(T+K)$.
Then
\[
 \sigma_w(T)=\sigma_e(T)\cup
 \{\lambda:T-\lambda I\text{ Fredholm with nonzero index}\}.
\]
Both sets are invariant under compact perturbation. For normal $N$,
$\sigma_w(N)=\sigma_e(N)$. On a complex Hilbert space,
$\sigma_e(N)$ consists of all nonisolated spectral points together
with isolated eigenvalues of infinite multiplicity. On a real Hilbert
space, writing $N_{\mathbb C}$ for its canonical Hilbert
complexification, a real $\lambda$ belongs to $\sigma_e(N)$ exactly
when $\lambda\in\sigma(N_{\mathbb C})$ and either it is nonisolated
in that complex spectrum, or it is isolated there and
$\ker_{\mathbb R}(N-\lambda I)$ is infinite dimensional.
For real self-adjoint $N$ the complexified spectrum is real, so the
original isolated/nonisolated description also holds in $\sigma_{\mathbb R}(N)$.
\end{theorem}
\begin{proof}
The essential spectrum is the spectrum of the Calkin class by Atkinson,
so is compact and compact-perturbation invariant; over $\mathbb C$ it
is nonempty. Intrinsic real essential spectra may be empty.
A compact perturbation cannot make a non-Fredholm or nonzero-index
Fredholm operator invertible, proving one inclusion in the formula.
Over either scalar field, a Fredholm operator of index zero admits a
finite-rank correction to an invertible operator, by matching its
finite-dimensional kernel with a complement of its range. This proves
the other inclusion without assuming a complex perturbation descends
to a real one. Weyl-spectrum invariance also follows by translating
the set of all compact perturbations.
For complex normal $N$, $\ker(N-\lambda)=\ker(N-\lambda)^*$, so every
Fredholm index is zero. By its PVM, boundedness below off the kernel
is equivalent to a complex spectral gap around $\lambda$; the kernel
is $E\{\lambda\}H$. Thus Fredholmness is exactly the isolated
finite-multiplicity condition, or membership in the resolvent.
For a real bounded operator, complexification preserves closed range,
Fredholmness and index: its kernel and, for closed range, orthogonal
cokernel are the complexifications of the original real spaces, with
the same respective dimensions. Apply the complex normal conclusion
to $N_{\mathbb C}$ at each real $\lambda$ and use those identifications.
Normality also gives equal real kernel and adjoint-kernel dimensions,
so the real Weyl formula reduces to $\sigma_w(N)=\sigma_e(N)$.
For real self-adjoint $N$, the complexified spectrum is real and its
real-axis slice equals $\sigma_{\mathbb R}(N)$, giving the last assertion.
\end{proof}

\begin{proposition}[Approximate points and local alternatives]\label{fa2:p2:local-spectrum}
For a bounded operator on a Banach space,
$\lambda\in\sigma_{ap}(T)$ iff there are unit vectors $x_n$ with
$(T-\lambda)x_n\to0$; its complement means injectivity with closed
range. Always $\partial\sigma(T)\subset\sigma_{ap}(T)$.
For a Hilbert semi-Fredholm operator $A$ and $\|B\|<\gamma(A)$,
$A+B$ is semi-Fredholm and both nullity and deficiency cannot increase.
Near a semi-Fredholm spectral value $\lambda$, the nullity and
deficiency of $T-\mu$ are constant on a sufficiently small punctured
disk. Consequently a boundary spectral point is either isolated or
belongs to both the left and right essential spectra.
\end{proposition}
\begin{proof}
The first equivalence is the lower-bound and inverse-range criterion.
If $\lambda_n\in\rho(T)$ tends to a boundary point $\lambda$, the
resolvent norms tend to infinity; otherwise the resolvent identity
would make their inverses converge to an inverse at $\lambda$.
Choose unit $y_n$ nearly attaining these norms and normalize
$x_n=R(\lambda_n,T)y_n/\|R(\lambda_n,T)y_n\|$; then
$(T-\lambda)x_n\to0$.
For the perturbation assertion, on $\ker(A+B)$ one has
$\gamma(A)\operatorname{dist}(x,\ker A)\leq\|Ax\|
\leq\|B\|\|x\|$. A subspace of dimension greater than the finite
kernel contains a unit vector perpendicular to that kernel, giving
a contradiction. Apply the same argument to $A^*$ for deficiency.
Closedness of the perturbed range follows from a left parametrix and
a finite-dimensional correction, or from the block decomposition
with its bounded-below infinite block; this proves semi-Fredholmness.
For the punctured-disk assertion assume, by adjunction if necessary,
$\ker A$ finite and put $M=\bigcap_{n\geq1}\faRan A^n$.
All ranges are closed because powers are semi-Fredholm. The decreasing
spaces $\ker A\cap\faRan A^n$ stabilize by finite dimensionality.
This implies $A(M)=M$: for $x\in M$ choose inverse images in each
range, and correct one fixed inverse image by the stabilized kernel.
Thus $A|_M$ is surjective and semi-Fredholm. Small perturbations
$A|_M-\mu$ stay surjective with constant finite kernel dimension by
the block argument and index stability. For $\mu\ne0$,
$\ker(A-\mu)\subset M$, since $x=\mu^{-n}A^nx$.
Hence nullity is constant on a small punctured disk; index constancy
gives the same for deficiency. If $\lambda$ is also a boundary point,
the punctured disk meets the resolvent, so these constants are zero;
the whole punctured disk is therefore in the resolvent. This proves
the last alternative.
\end{proof}

\begin{theorem}[Toeplitz index for continuous scalar symbols]\label{fa2:p2:toeplitz}
On $H^2(\mathbb T)\subset L^2(\mathbb T)$ put
$T_\varphi=P_+M_\varphi|_{H^2}$, $\varphi\in C(\mathbb T)$.
If $\varphi$ never vanishes then $T_\varphi$ is Fredholm and
$\faind T_\varphi=-\operatorname{wind}(\varphi,0)$.
Its nullity is $\max(0,-\operatorname{wind}\varphi)$ and its
cokernel dimension is $\max(0,\operatorname{wind}\varphi)$.
The Hankel map $P_-M_\varphi|_{H^2}$ is compact.
\end{theorem}
\begin{proof}
$\|T_\varphi\|\leq\|\varphi\|_\infty$ by orthogonal projection.
For a trigonometric polynomial the Hankel map has finite rank, since
it can move only finitely many nonnegative Fourier modes into negative
ones. Uniform trigonometric approximation gives compactness for
continuous symbols. Consequently
$T_\varphi T_\psi-T_{\varphi\psi}
=-P_+M_\varphi P_-M_\psi|_{H^2}$ is compact.
Take $\psi=1/\varphi$ to obtain a parametrix.
Write $\varphi(z)=z^n e^{g(z)}$, $n=\operatorname{wind}\varphi$,
where a continuous logarithm $g$ exists after removing the winding.
The symbols $z^ne^{tg}$ give a Fredholm path to $T_{z^n}$.
For $n\geq0$ this is the shift by $n$, with zero kernel and
cokernel dimension $n$; for $n<0$ it is the backward shift with
kernel dimension $-n$ and is onto. Index invariance proves the formula.
To obtain the separate dimensions we prove Coburn's alternative.
If nonzero $f,g\in H^2$ lay in the kernels of $T_\varphi$ and
$T_{\bar\varphi}$ respectively, write
$\varphi f=\bar z\bar h$ and $\bar\varphi g=\bar z\bar k$ with
$h,k\in H^2$. Then
$\bar z\overline{hg}=zfk$ belongs simultaneously to the strictly
negative and strictly positive Fourier subspaces of $L^1$.
Uniqueness of Fourier coefficients makes it zero. A product of two
nonzero analytic functions is nonzero, so $k=0$, forcing $g=0$ since
$\varphi$ is nowhere zero. This contradiction says at least one kernel
is zero. Since $T_\varphi^*=T_{\bar\varphi}$, the index formula gives
the two dimensions as claimed.
\end{proof}

\begin{proposition}[Isolated points and finite Riesz rank]\label{fa2:p2:isolated-fredholm}
For an isolated $\lambda\in\sigma(T)$ on a Hilbert space, the
following are equivalent: $T-\lambda$ is semi-Fredholm; its Riesz
projection has finite rank; and $T-\lambda$ is Fredholm of index zero.
\end{proposition}
\begin{proof}
If it is semi-Fredholm, local index constancy and nearby resolvent
points give index zero, hence Fredholmness. Split by its Riesz
projection. On the complement $T-\lambda$ is invertible; on the
Riesz space it is quasinilpotent and Fredholm. If that space were
infinite dimensional, its class in the nonzero unital Calkin algebra
would be invertible yet have spectral radius zero, since quotient
maps decrease spectral radius. A nonempty spectrum of an invertible
element cannot be $\{0\}$, a contradiction. Thus the Riesz space
is finite dimensional. Conversely a finite Riesz space contributes
a finite-dimensional operator of index zero and the complement
is invertible, proving the reverse implication.
\end{proof}

\begin{proposition}[Index as a difference of traces]\label{fa2:p2:trace-index}
Suppose $T:H\to K$, $S:K\to H$ are bounded and both
$I_H-ST$ and $I_K-TS$ are trace class. Then $T$ is Fredholm and
\[
 \faind T=\faTr(I_H-ST)-\faTr(I_K-TS).
\]
More generally if their $m$th powers are trace class for an integer
$m\geq1$, the same formula holds with those powers.
\end{proposition}
\begin{proof}
For $m=1$ Atkinson gives Fredholmness. Let $G$ be its generalized
inverse, with $I-GT=P$, $I-TG=Q$ the finite-rank kernel and cokernel
projections. Since $(S-G)T$ is trace class,
$S-G=(S-G)TG+(S-G)Q$ is trace class. Cyclicity of the trace gives
$\faTr((S-G)T)=\faTr(T(S-G))$. Subtracting the two defect traces
therefore leaves $\faTr P-\faTr Q=\faind T$.
For general $m$, put $D=I-ST$, $E=I-TS$ and
$S_m=(I+D+\cdots+D^{m-1})S$.
Then $I-S_mT=D^m$, and $TD=ET$ gives
$I-TS_m=E^m$. These trace-class defects yield the same argument.
This is the current trace-index subcase of the mixed source family;
no unrestricted Banach nuclear determinant assertion is inferred.
\end{proof}

% END INPUT


% BEGIN parts/part2-ritz-semifredholm.tex
\chapter{Ritz approximation and one-sided Fredholm inverses}
\label{fa2:p2:ritz-semifredholm-chapter}

\begin{theorem}[Monotone Rayleigh--Ritz approximation]
\label{fa2:p2:ritz-approximation}
Let $T\ge0$ be compact on a Hilbert space, and let $\mu_k$ be
its eigenvalues in decreasing order, repeated with multiplicity
and padded with zero if necessary. For closed subspaces
$V_n\subset V_{n+1}$ with dense union, let $P_n=P_{V_n}$ and
$T_n=P_nT|_{V_n}$. Its ordered eigenvalues $\nu_k^{(n)}$ satisfy,
whenever $\dim V_n\ge k$,
\[
 0\le\nu_k^{(n)}\le\nu_k^{(n+1)}\le\mu_k,
 \qquad \nu_k^{(n)}\longrightarrow\mu_k.
\]
More precisely $|\nu_k^{(n)}-\mu_k|\le\|P_nTP_n-T\|\to0$,
with zero padding. For $R(x)=\langle x,Tx\rangle/\|x\|^2$,
and finite integers $1\le k\le\dim H$,
\[
 \mu_k=\max_{\dim L=k}\min_{x\in L\setminus0}R(x)
       =\min_{\dim M=k-1}\sup_{x\in M^\perp\setminus0}R(x).
\]
If $T$ is injective, the reciprocal forms are
\[
 \frac1{\mu_k}=\min_{\dim L=k}\sup_{x\in L\setminus0}\frac1{R(x)}
 =\sup_{\operatorname{codim}W=k-1}\inf_{x\in W\setminus0}\frac1{R(x)},
\]
where $W$ is closed. These include the monotone approximation
assertion of \cite[Problem 37.10]{B01}.
\end{theorem}
\begin{proof}
The compact spectral expansion proves the two min--max formulas:
the span $E_k$ of the first $k$ eigenvectors has minimum quotient
$\mu_k$, and every $k$-space meets $E_{k-1}^\perp$ nontrivially,
where the quotient is at most $\mu_k$. Dually every subspace
of codimension $k-1$ meets $E_k$, and $E_{k-1}^\perp$ gives
equality. Zero eigenvalues are handled by the same inequalities
and a subspace of the kernel when needed. Taking reciprocals
on these positive quotients reverses the inner extrema and
proves the displayed reciprocal formulas, including attainment
at those same spectral subspaces.
The quotient of the compression on $V_n$ is exactly $R|_{V_n}$.
Thus increasing its trial subspaces proves monotonicity and
the upper bound by $\mu_k$. The min--max formulas also show
that ordered eigenvalues of two positive compact operators
on the same space differ by at most their norm distance,
since each unit Rayleigh quotient differs by that bound.

Dense increasing $V_n$ give $P_nx\to x$: approximate $x$ by a
vector in one $V_N$ and use the minimizing property of the
orthogonal projection. This convergence is uniform on the
compact closure of $T(B_H)$ by finite nets and the common
bound $\|I-P_n\|\le1$. Hence $\|(I-P_n)T\|\to0$.
Adjunction gives $\|T(I-P_n)\|\to0$ too, and
$\|T-P_nTP_n\|\le2\|(I-P_n)T\|\to0$.
Apply the preceding eigenvalue bound to $T$ and the compression
extended by zero on $V_n^\perp$. This proves the error estimate
and convergence, without assuming finite-dimensional trial spaces.
\end{proof}

\begin{theorem}[One-sided parametrices on Banach spaces]
\label{fa2:p2:one-sided-fredholm}
For a bounded map $T:E\to F$ between Banach spaces, the following
three right-sided conditions are equivalent:
\begin{enumerate}
\item the range is closed of finite codimension and $\ker T$
      is complemented in $E$;
\item $TS=I_F+K$ for some bounded $S:F\to E$ and compact $K$ on $F$;
\item $TU=I_F-P$ for some bounded $U:F\to E$ and finite-rank
      projection $P$ on $F$.
\end{enumerate}
In the last form one may choose $\operatorname{rank}P=
\operatorname{codim}\operatorname{Ran}T$.
The following three left-sided conditions are likewise equivalent:
the range is closed and complemented and the kernel finite dimensional;
$ST=I_E+K$ with $K$ compact; and $UT=I_E-Q$ with $Q$ a
finite-rank projection. One may choose
$\operatorname{rank}Q=\dim\ker T$.
These are the general Banach-space statements of
\cite[Problem 38, A--B]{B01}; a one-sided inverse modulo compacts
does not remove the stated complement condition.
\end{theorem}
\begin{proof}
Under the first right-sided condition write
$E=\ker T\oplus E_0$ and $F=\operatorname{Ran}T\oplus F_0$.
Both splits are bounded: the first by hypothesis, the second
because the closed range has finite codimension. Invert
$T|_{E_0}:E_0\to\operatorname{Ran}T$ and compose with the
range projection. Its resulting $U$ satisfies $TU=I-P_{F_0}$,
with exactly the stated rank. The third condition implies the
second immediately. If $TS=I+K=D$, Riesz--Schauder gives a
finite-dimensional kernel and cokernel for $D$ and a bounded
regularizer $R$ with $DR=I-P$, where $P$ is a finite-rank
projection. This follows explicitly by inverting $D$ between
complements of its kernel and range. Thus $U=SR$ gives the
third condition.

Suppose now $TU=I-P$. The closed finite-codimensional space
$F_1=\ker P$ lies in $\operatorname{Ran}T$. Any linear space
containing $F_1$ is closed, since its image in $F/F_1$ is
finite dimensional. Thus the range is closed of finite codimension.
Decompose it as $F_1\oplus F_2$ with $F_2$ finite dimensional.
The restriction $U|_{F_1}$ is a bounded right inverse there.
Choose preimages of a basis of $F_2$ to obtain a bounded
right inverse on that finite-dimensional summand too.
Together they give a bounded $W:\operatorname{Ran}T\to E$
with $TW=I$. Then $I-WT$ is a bounded projection onto
$\ker T$, proving its complement condition.

For the left-sided statement, closed complemented range and
finite kernel permit the same inverse-between-complements
construction, giving $UT=I-Q_{\ker T}$.
From $ST=I+K=D$, take a left regularizer $R$ for the
Fredholm operator $D$, with $RD=I-Q$; then $U=RS$.
It remains to recover the first left-sided condition from
$UT=I-Q$. Its kernel is contained in $\operatorname{Ran}Q$,
hence finite dimensional. The range is closed: otherwise the
closed-range distance estimate would fail and there would be
$x_n$ with $\operatorname{dist}(x_n,\ker T)=1$,
$\|x_n\|\le2$, and $Tx_n\to0$. Finite rank of $Q$ gives
a subsequence for which $Qx_n$ converges; the equality
$x_n=UTx_n+Qx_n$ makes this subsequence converge to a kernel
point, contradicting distance one.

Put $E_1=\ker Q$. On $E_1$ one has $UT=I$, so $T(E_1)$
is closed. The bounded map $P_1=T(I-Q)U$ is a projection onto
$T(E_1)$, since it fixes $Tx$ for every $x\in E_1$.
Now $\operatorname{Ran}T=T(E_1)+T(\operatorname{Ran}Q)$
is a finite-dimensional extension of that complemented space.
In $\ker P_1$ its remaining component
$(I-P_1)T(\operatorname{Ran}Q)$ is finite dimensional and has
a bounded projection by Hahn--Banach. Combining that projection
with $P_1$ gives a bounded projection onto the entire range.
This proves the missing complement condition and completes
both equivalences and the exact finite-defect refinements.
\end{proof}

% END INPUT


% BEGIN parts/part2-classical-fredholm.tex
\chapter{The classical Fredholm series for continuous kernels}
\label{fa2:p2:classical-fredholm-chapter}
This is the continuous-kernel determinant, as in \cite[Chapter 24]{B08}.
It must be distinguished from the trace-class determinant: a continuous
kernel need not define a trace-class Hilbert operator. The two agree
when the trace-class hypothesis really holds.

Let $Q$ be compact metric, $\mu$ a finite positive Radon measure,
and $k\in C(Q\times Q;\mathbb C)$. Write
$Kf(x)=\int_Q k(x,y)f(y)\,d\mu(y)$ on $C(Q)$.
Put $\mu(Q)=m$ and $\|k\|_\infty=M$. Determinants of size zero
and their integrals have value one. Define
\begin{align*}
 D_k(z)&=\sum_{n=0}^\infty\frac{(-z)^n}{n!}
     \int_{Q^n}\det[k(t_i,t_j)]_{i,j=1}^n\,d\mu^n(t),\\
 B_k(x,y;z)&=\sum_{n=0}^\infty\frac{(-z)^n}{n!}
   \int_{Q^n}
   \det\begin{pmatrix}
     k(x,y)&k(x,t_1)&\cdots&k(x,t_n)\\
     k(t_1,y)&k(t_1,t_1)&\cdots&k(t_1,t_n)\\
     \vdots&\vdots&&\vdots\\
     k(t_n,y)&k(t_n,t_1)&\cdots&k(t_n,t_n)
   \end{pmatrix}\,d\mu^n(t).
\end{align*}
At $n=0$ the second integrand is $k(x,y)$.

\begin{theorem}[Fredholm series and resolvent kernel]
\label{fa2:p2:classical-fredholm}
The series define an entire scalar function $D_k$, with $D_k(0)=1$,
and an entire $C(Q\times Q)$-valued function $B_k$. They satisfy
\begin{align*}
 B_k(x,y;z)&=D_k(z)k(x,y)
          +z\int_Q k(x,t)B_k(t,y;z)\,d\mu(t),\\
 B_k(x,y;z)&=D_k(z)k(x,y)
          +z\int_Q B_k(x,t;z)k(t,y)\,d\mu(t),\\
 D_k'(z)&=-\int_Q B_k(t,t;z)\,d\mu(t).
\end{align*}
The operator $I-zK$ is invertible exactly when $D_k(z)\ne0$.
At such a parameter its inverse is
\[
 (I-zK)^{-1}f(x)=f(x)+
       \frac z{D_k(z)}\int_Q B_k(x,y;z)f(y)\,d\mu(y).
\]
In particular the exceptional parameters form a discrete set and
are precisely the reciprocals of the nonzero eigenvalues of $K$.
\end{theorem}
\begin{proof}
Hadamard's inequality bounds the $n$th determinant by
$n^{n/2}M^n$ and the augmented determinant by
$(n+1)^{(n+1)/2}M^{n+1}$. After division by $n!$ and multiplication
by $m^n|z|^n$ both majorant series converge for every bounded set of
parameters: the ratio estimate, or $n!\ge(n/e)^n$, makes the $n$th
root tend to zero. The latter factorial bound follows by summing
$\log j\ge\int_{j-1}^j\log t\,dt$.
Thus the series converge locally uniformly, uniformly also in $x,y$
for $B_k$, and termwise differentiation is legitimate.

Let $d_n$ and $b_n(x,y)$ be their coefficients of $z^n$.
Expansion of the augmented determinant along its first row gives
$b_n=d_n k+K b_{n-1}$ for $n\ge1$, with $b_0=k$.
To see the coefficient, the first-column term is $d_nk$. In each
of the other $n$ terms rename its chosen integration variable as
$t$ and cyclically move its remaining row to the first position.
The permutation sign together with the cofactor sign is minus.
The common factor $n(-1)^n/n!$ is therefore
$(-1)^{n-1}/(n-1)!$, exactly the coefficient of $K b_{n-1}$.
The remaining determinant is the augmented $(n-1)$-minor at $(t,y)$.
Column expansion gives the corresponding right composition.
Finally $n d_n=-\int b_{n-1}(t,t)\,d\mu(t)$ directly from the
definitions. Summing proves all three identities.

If $D_k(z)\ne0$, multiply the proposed inverse on either side by
$I-zK$. The first two identities cancel its integral kernel and
leave the identity operator, proving two-sided invertibility.
For the converse suppose $I-z_0K$ is invertible. In a disk about
$z_0$ the inverse is analytic by the Neumann perturbation formula.
Its resolvent kernel is the continuous, norm-analytic function
\[
 r_z(x,y)=\big((I-zK)^{-1}k(\,\cdot,y)\big)(x).
\]
Here $y\mapsto k(\cdot,y)$ is norm continuous and the inverse is
uniformly bounded on smaller disks. The first identity, with the
inverse applied to its column, yields $B_k=D_k r_z$ throughout the
disk. Therefore
$D_k'=-D_k\int_Qr_z(t,t)\,d\mu(t)$ there.
An analytic scalar solution of $D'=-hD$ that is zero at $z_0$ is
identically zero on that disk: multiply by the exponential of a
local primitive of $h$. This would make the entire $D_k$ zero,
contrary to $D_k(0)=1$. Hence $D_k(z_0)\ne0$.
Compactness of the continuous-kernel operator, proved in Part V,
and Riesz--Schauder identify the exceptional parameters with the
stated eigenvalue reciprocals. Alternatively that identification
follows from the Fredholm equivalence already proved in Part II.
\end{proof}

\begin{theorem}[Multiplication and diagonal logarithmic derivative]
\label{fa2:p2:classical-fredholm-product}
The determinant is continuous under uniform convergence of continuous
kernels, locally uniformly in its parameter. If $L$ has continuous
kernel $\ell$ on the same space and
$h(x,y)=k(x,y)+\ell(x,y)-\int_Q k(x,t)\ell(t,y)\,d\mu(t)$, then
\[
 D_h(1)=D_k(1)D_\ell(1).
\]
More generally substitute any scalar multiples of the kernels in
this formula. Write $k_j$ for the continuous kernel of $K^j$ and
$c_j=\int_Q k_j(t,t)\,d\mu(t)$. Near zero,
\[
 \log D_k(z)=-\sum_{j=1}^\infty\frac{z^j}{j}c_j.
\]
If the operator on $L^2(\mu)$ is trace class, $c_j=\operatorname{tr}K^j$
for every $j\ge1$ and $D_k(z)=\det(I-zK)$ in the trace-class sense.
\end{theorem}
\begin{proof}
On a uniformly bounded set of kernels the preceding Hadamard
majorants are uniform. Each coefficient is a finite polynomial
integral, hence uniformly continuous under uniform convergence.
Dominated convergence of the series proves the continuity claim.
For a finite-rank continuous kernel
$k(x,y)=\sum_{j=1}^N a_j(x)b_j(y)$, set
$M_{ij}=\int_Q b_i(t)a_j(t)\,d\mu(t)$.
Cauchy--Binet expands its $n$th coefficient as a sum of products
of two $n$ by $n$ minors. Expanding each into its permutations and
integrating separately shows that the integral is $n!$ times the
sum of the $n$th principal minors of $M$. Hence
$D_k(z)=\det(I_N-zM)$, with all coefficients for $n>N$ zero.
This also follows by the same finite permutation expansion and
does not presuppose an infinite-dimensional determinant theorem.
For two finite-rank kernels choose the common finite-dimensional
span of all their range functions. The ordinary determinant on
that space proves the multiplication identity. Finite sums of
products $a(x)b(y)$ uniformly approximate any continuous kernel,
by Stone--Weierstrass. Their compositions converge uniformly as
well, so the continuity just proved passes the identity to $k,\ell$.

For $|z|mM<1$, the resolvent kernel has the uniformly convergent
series $r_z=\sum_{j\ge1}z^{j-1}k_j$, because
$\|k_j\|_\infty\le M^j m^{j-1}$. Divide the derivative identity
by $D_k$ and integrate its power series, using $D_k(0)=1$.
This proves the logarithmic formula.
If $K$ is trace class on $L^2$, its powers have continuous kernels
and remain trace class. The diagonal trace identity proved with
finite measurable averaging partitions in
\ref{fa2:p2:mercer} gives $c_j=\operatorname{tr}K^j$ even without
positivity. The trace-class determinant has the same logarithmic
series near zero. Both determinants are entire, so they agree
everywhere by the identity theorem.
\end{proof}

\begin{corollary}[The general eigenvalue product is regularized]
\label{fa2:p2:classical-fredholm-eigenvalues}
Let $\lambda_j$ be the nonzero eigenvalues of the continuous-kernel
operator on $L^2(\mu)$, with algebraic multiplicity. Then
\[
 D_k(z)=\exp\!\left(-z\int_Q k(t,t)\,d\mu(t)\right)
          \prod_j(1-z\lambda_j)e^{z\lambda_j}.
\]
The product converges locally uniformly. Each nonzero root has the
algebraic multiplicity of the corresponding eigenvalue. The nonzero
eigenvalues and their algebraic multiplicities on $C(Q)$ and $L^2$
agree. The unregularized product is not being asserted in the
absence of trace-class hypotheses.
\end{corollary}
\begin{proof}
The kernel is Hilbert--Schmidt, hence $\sum_j|\lambda_j|^2<\infty$
by the product-to-convex-sum argument in the Weyl inequalities,
now with $F(t)=e^{2t}$. For bounded $z$ the tails have
$|z\lambda_j|<1/2$ and
$|\log(1-z\lambda_j)+z\lambda_j|\le2|z\lambda_j|^2$.
Thus the product is entire and its zeros have exactly the indicated
multiplicities. For $j\ge2$, $K^j$ is trace class, and its continuous
kernel gives $c_j=\operatorname{tr}K^j=\sum_i\lambda_i^j$ by
Lidskii and compact spectral mapping. Absolute summation identifies
the logarithm of the product with the preceding logarithmic formula
near zero. Analytic continuation gives the displayed identity.

Finally let $J:C(Q)\to L^2(\mu)$ be the natural bounded map and
$A:L^2(\mu)\to C(Q)$ the kernel map, bounded by Cauchy--Schwarz
and continuity of $x\mapsto k(x,\cdot)$ in $L^2$.
The two operators are $AJ$ and $JA$. Their nonzero spectra agree:
for $\lambda\ne0$, an inverse of $\lambda-AJ$ gives one of
$\lambda-JA$ by
$\lambda^{-1}[I+J(\lambda-AJ)^{-1}A]$, and conversely.
The maps $J,A$ intertwine their generalized eigenspaces; on either
such finite-dimensional space the corresponding operator is
invertible. Thus $J$ and $A$ are injections in opposite directions
between those spaces and their dimensions are equal. This proves
the full algebraic-multiplicity assertion, even when $J$ is not
injective on all of $C(Q)$ because $\mu$ lacks full support.
\end{proof}

% END INPUT


% BEGIN parts/part2-fredholm-holder.tex
\begin{proposition}[The improved Fredholm bound for a H\"older kernel]
\label{fa2:p2:fredholm-holder}
Let $Q=[a,b]$, $L=b-a>0$, and let $\mu$ be a finite positive
Radon measure of mass $m$. Suppose that the continuous kernel $k$
satisfies $|k(x,y)|\le M$ and
$|k(x,y)-k(x,z)|\le H|y-z|^\alpha$, where $0<\alpha\le1$.
For $n\ge2$ and arbitrary $x_i,y_j\in Q$,
\[
 \left|\det[k(x_i,y_j)]_{i,j=1}^n\right|
 \le n^{n/2}M H^{n-1}
       \left(\frac{L}{n-1}\right)^{\alpha(n-1)}.
\]
The same bound holds for H\"older regularity in the first variable.
Consequently the coefficient $d_n$ of $z^n$ in $D_k(z)$ satisfies
\[
 |d_n|\le\frac{m^n}{n!}n^{n/2}M H^{n-1}
       \left(\frac{L}{n-1}\right)^{\alpha(n-1)},
\]
and $D_k$ has entire-function order at most $(\alpha+1/2)^{-1}$.
Here order means
\[
 \limsup_{r\to\infty}
 \frac{\log\log\max\{e,\max_{|z|\le r}|D_k(z)|\}}{\log r}.
\]
\end{proposition}
\begin{proof}
Permute the columns so that $y_1\le\cdots\le y_n$; this changes
only the determinant's sign. Replace column $j$ by the difference
between columns $j$ and $j+1$ for $1\le j<n$, leaving the last
column unchanged. This transformation has determinant one.
Its first $n-1$ columns have Euclidean norms at most
$\sqrt n H(y_{j+1}-y_j)^\alpha$; the last has norm at most
$\sqrt n M$. Hadamard's inequality bounds the determinant by
$n^{n/2}M H^{n-1}\prod_{j<n}(y_{j+1}-y_j)^\alpha$.
The $n-1$ nonnegative gaps have sum at most $L$. Their product is
therefore at most $(L/(n-1))^{n-1}$ by the arithmetic--geometric
mean inequality, proving the claim, including coincident points.
Transposition proves the first-variable version. Inserting
$x_i=y_i=t_i$ and integrating gives the coefficient estimate.

Put $\beta=\alpha+1/2$. If $M=0$ the determinant is identically
one. If $H=0$, all columns coincide and $D_k$ is a polynomial
of degree at most one. Otherwise $n!\ge(n/e)^n$ and
$(n-1)^{n-1}\ge n^{n-1}e^{-1}$ show, after enlarging constants,
that $|d_n|\le C A^n n^{-\beta n+\alpha}$.
For any $0<\beta'<\beta$, absorb $n^\alpha$ and the excess
factor $n^{-(\beta-\beta')n}$ into a constant, obtaining
$|d_n|\le C' A'^n n^{-\beta'n}$.
For $t=A'r$, set $N=\lceil(2t)^{1/\beta'}\rceil$.
The logarithm of every summand $t^n n^{-\beta'n}$ is at most
$C_{\beta'}t^{1/\beta'}$, by maximizing
$s\log t-\beta's\log s$ over positive real $s$.
The first $N$ terms thus have sum at most
$N\exp(C_{\beta'}t^{1/\beta'})$; for $n>N$ the summand is at
most $2^{-n}$. Hence $\log\max_{|z|\le r}|D_k(z)|
\le C''(1+r^{1/\beta'}+\log(2+r))$.
The order is at most $1/\beta'$, and letting $\beta'\uparrow\beta$
proves the assertion.
\end{proof}

\begin{remark}[The spacing factor in the source]
This supplies the H\"older refinement discussed in
\cite[pp.~262--263]{B08}. With $n-1$ successive gaps the valid
arithmetic--geometric mean factor is $(L/(n-1))^{n-1}$.
The printed factor $n^{-n}$ on a unit interval is not valid for
these gaps: for $n=2$, $y_1=0,y_2=1$, their product is one.
The corrected estimate retains the improved entire order and
does not alter the Fredholm resolvent identities.
\end{remark}

% END INPUT


% BEGIN parts/part2-trace.tex
\chapter{Compact Hilbert operators, singular values, and traces}
Throughout this chapter $H,K$ may be nonseparable. Every compact
operator has separable range closure; its nonzero singular part is
therefore countable. Arbitrary-dimensional kernels are retained.
The trace assertions here require their stated ideal hypotheses.
A continuous kernel alone does not imply trace class. See
\cite{B01,B02,B06,B08}.

\begin{theorem}[Compact normal spectral expansion]\label{fa2:p2:compact-spectral}
Let $N$ be compact and normal on a complex Hilbert space $H$,
which need not be separable. Then $H$ is the orthogonal Hilbert sum
of $\ker N$ and its finite-dimensional nonzero eigenspaces.
Distinct eigenspaces are orthogonal; there are at most countably
many nonzero eigenvalues, with only finitely many of modulus at
least $\varepsilon$ for each $\varepsilon>0$. Thus
$Nx=\sum_{\lambda\ne0}\lambda P_\lambda x$, in norm, and finite
spectral sums ordered by decreasing modulus converge to $N$ in
operator norm. The norm of an omitted tail is the supremum of
the omitted eigenvalue moduli, with supremum zero for an empty tail.
Two complex compact normal operators are unitarily equivalent
exactly when all eigenvalue multiplicities agree, including the
possibly arbitrary cardinal dimension of their kernels.

For a compact normal operator on a real Hilbert space, the correct
orthogonal decomposition instead consists of its kernel, real
nonzero eigenvalue lines, and two-dimensional reducing planes on
which the matrix is
\[
 B(a,b)=\begin{pmatrix}a&-b\\ b&a\end{pmatrix},\qquad b>0.
\]
There are at most countably many nonzero blocks, with only
finitely many having modulus at least $\varepsilon$ for each
$\varepsilon>0$; their moduli are $|\lambda|$ for real lines and
$\sqrt{a^2+b^2}$ for planes. Finite block truncations converge in
operator norm, and their tail norms are the supremum of omitted
block moduli. Orthogonal equivalence is characterized by the
multiplicity of each real eigenvalue, each pair $(a,b)$ with
$b>0$, and the real Hilbert dimension of the kernel.
In particular a compact self-adjoint operator on either scalar
field has the eigenvector expansion above with real eigenvalues;
in the real case no nonreal rotation blocks occur. Zero spaces,
the zero operator and finite nonzero spectra are included.
\end{theorem}
\begin{proof}
First work over $\mathbb C$. For normal $N-\lambda$, equality of
its norm and adjoint norm on each vector gives
$\ker(N-\lambda)=\ker(N^*-\bar\lambda)$.
Thus the eigenspaces reduce $N$, and the adjoint identity makes
distinct eigenspaces orthogonal. Compact Riesz theory makes each
nonzero eigenspace finite-dimensional and leaves only finitely
many eigenvalues outside any disk about zero. The orthogonal
complement of their closed span is reducing. Its restriction is
compact normal and has no nonzero eigenvalue, hence no nonzero
spectral value. Normality identifies norm and spectral radius,
so the restriction is zero. This proves the complete decomposition.
Pythagoras gives pointwise convergence and the exact tail norm;
the finite-outside-each-disk property gives operator-norm convergence.
A unitary intertwiner preserves each eigenspace. Conversely take
the Hilbert sum of chosen unitaries between corresponding
eigenspaces, including an arbitrary-dimensional kernel.

For real $H$, use its Hilbert complexification
$H_{\mathbb C}=H\oplus iH$, with
$\|x+iy\|^2=\|x\|^2+\|y\|^2$, conjugation $C(x+iy)=x-iy$, and
$N_{\mathbb C}(x+iy)=Nx+iNy$. Direct inner-product calculation gives
$(N_{\mathbb C})^*=(N^*)_{\mathbb C}$ and
$\|N_{\mathbb C}\|=\|N\|$. The complexification is compact:
from a bounded sequence $x_j+iy_j$, extract a subsequence on which
$Nx_j$ converges and then one on which $Ny_j$ converges.
Normality therefore passes to $N_{\mathbb C}$, to which the complex
case applies. Since $CN_{\mathbb C}=N_{\mathbb C}C$, conjugation
interchanges the eigenspaces for $\lambda$ and $\bar\lambda$.
For real $\lambda$, its eigenspace is the complexification of its
fixed-real subspace; a real orthonormal basis of that subspace is
also a complex orthonormal basis of the eigenspace.

For $\lambda=a+ib$, $b>0$, choose an orthonormal basis $(z_j)$ of
its eigenspace. The vectors $Cz_j$ are an orthonormal basis of the
orthogonal conjugate eigenspace. Set
\[
 e_j=\frac{z_j+Cz_j}{\sqrt2},\qquad
 f_j=\frac{i(z_j-Cz_j)}{\sqrt2}.
\]
These vectors are fixed by $C$, are real orthonormal, and span the
fixed-real part of the paired eigenspaces. Direct calculation gives
$Ne_j=ae_j+bf_j$ and $Nf_j=-be_j+af_j$, yielding $B(a,b)$.
Take these planes, the real-eigenvalue lines and the fixed-real
kernel. The complex decomposition and conjugation show that their
real Hilbert sum is all of $H$. Their norms and finite truncation
tails give the real convergence assertions. A real orthogonal
intertwiner complexifies to a unitary intertwiner, so it preserves
the listed multiplicities. Conversely identify the corresponding
real lines, standard planes and kernels and take their orthogonal
Hilbert sum. This proves both directions of the real classification.

If $N=N^*$, every eigenvalue of the complexification is real,
since its quadratic form is real. Hence the nonreal paired
eigenspaces are absent, proving the real self-adjoint specialization.
The distinction is essential: on $\mathbb R^2$, $N(x,y)=(-y,x)$
satisfies $N^*=-N$, $N^*N=NN^*=I$ and $N^2=-I$, but it has no
nonzero real eigenvector. It is one plane $B(0,1)$, not a sum of
real eigenvalue lines.
\end{proof}

\begin{theorem}[Courant--Fischer and singular systems]\label{fa2:p2:singular-system}
For compact $T:H\to K$ there are positive singular values
$s_1\geq s_2\geq\cdots\downarrow0$ and orthonormal families
$(v_j)$, $(u_j)$ with
\[
 Tx=\sum_js_j u_j\faip{v_j}{x},\qquad T^*u_j=s_jv_j,
 \quad (\ker T)^\perp=\overline{\operatorname{span}}\{v_j\}.
\]
Finite rank is included by terminating the list and then putting
$s_j=0$. One has
\[
 s_n(T)=\inf_{\operatorname{rank}F<n}\|T-F\|,
 \quad s_n(ATB)\leq\|A\|\|B\|s_n(T).
\]
For compact positive $C$ with eigenvalues $\lambda_n$ in decreasing
order, counted with multiplicity,
\[
 \lambda_n=\inf_{\dim L=n-1}\ 
 \sup_{\substack{x\perp L\\\|x\|=1}}\faip{x}{Cx}
 =\sup_{\dim L=n}\ \inf_{\substack{x\in L\\\|x\|=1}}
 \faip{x}{Cx}.
\]
These formulas apply below the dimension of the space, padding by
zeros when appropriate; positive and negative eigenvalues of compact
self-adjoint operators are treated separately.
\end{theorem}
\begin{proof}
Apply the compact positive spectral theorem to $T^*T$ and take square
roots of its positive eigenvalues. Set $u_j=Tv_j/s_j$.
The adjoint identity proves orthonormality and $T^*u_j=s_jv_j$.
On the kernel $T$ vanishes, so the expansion follows.
Truncating after $n-1$ terms gives approximation error $s_n$.
Conversely any rank less than $n$ map has a nonzero kernel vector in
$\operatorname{span}\{v_1,\ldots,v_n\}$; on a unit such vector its
error from $T$ has norm at least $s_n$. This proves the approximation
formula. Composing a rank-$<n$ approximant with $A,B$ proves the
ideal inequality. For the minmax upper bound take the first $n-1$
eigenvectors. For its lower bound any $(n-1)$-dimensional space has
a nonzero perpendicular vector in the span of the first $n$
eigenvectors; its Rayleigh quotient is at least $\lambda_n$.
For the maxmin formula reverse the argument: the first $n$
eigenvectors attain the lower bound, and every $n$-dimensional space
has a nonzero vector perpendicular to the first $n-1$ eigenvectors,
whose quotient is at most $\lambda_n$.
\end{proof}

\begin{theorem}[Hilbert--Schmidt operators and kernels]\label{fa2:p2:hilbert-schmidt}
The number $\|T\|_2^2=\sum_i\|Te_i\|^2$ is independent of the
orthonormal basis. Finiteness defines the Hilbert--Schmidt ideal $S_2$;
its operators are compact and $\|T\|_2^2=\sum_js_j(T)^2$.
It is a Hilbert space with inner product $\faTr(S^*T)$.
For sigma-finite measure spaces, a kernel $k\in L^2(\nu\times\mu)$
defines $T_k:L^2(\mu)\to L^2(\nu)$ by integration and
$\|T_k\|_2=\|k\|_2$. Conversely every Hilbert--Schmidt map between
these $L^2$ spaces has such a kernel, unique almost everywhere.
\end{theorem}
\begin{proof}
For bases $(e_i)$ and $(f_j)$, Parseval and nonnegative summation give
$\sum_i\|Te_i\|^2=\sum_{i,j}|\faip{f_j}{Te_i}|^2
=\sum_j\|T^*f_j\|^2$. This proves independence and adjoint equality.
Finite-column truncations converge in the $S_2$ norm because the
series is summable, and their operator norm errors are no greater
than their $S_2$ errors by Cauchy--Schwarz. Thus $T$ is compact.
Use a singular basis to obtain the singular-value formula.
The matrix coefficients identify $S_2(H,K)$ isometrically with
square-summable matrices; conversely such a matrix defines an operator
by Cauchy--Schwarz, proving completeness and the Hilbert identity.
For a kernel, Cauchy--Schwarz defines the operator for almost all
target points. Product $L^2$ functions are limits of countable sums
of simple rectangles; their source-$L^2$ sections therefore have
essentially separable range. Only countably many source basis
coordinates are active outside one common null set, so the following
Parseval--Tonelli step also applies when the full $L^2$ space is
nonseparable. For fixed target point Parseval in the source variable,
then Tonelli, gives $\sum_i\|T_ke_i\|^2=\int|k|^2$.
Conversely the $L^2$-convergent sum
$k(y,x)=\sum_i(Te_i)(y)\overline{e_i(x)}$ has this norm and reproduces
$T$ first on finite linear combinations and then by density.
The sum is countable on its nonzero coefficients. Uniqueness follows
by the same Parseval identity; density of simple rectangles in the
product $L^2$ space justifies the tensor identifications used here.
\end{proof}

\begin{theorem}[Trace class and cyclicity]\label{fa2:p2:trace-class}
A compact $T$ is trace class if $\|T\|_1=\sum_js_j(T)<\infty$.
Trace class is a complete two-sided ideal,
$\|ATB\|_1\leq\|A\|\|T\|_1\|B\|$, and
\[
 \faTr T=\sum_i\faip{e_i}{Te_i}
\]
is absolutely convergent and basis independent, with
$|\faTr T|\leq\|T\|_1$ and $\faTr(AT)=\faTr(TA)$.
It consists precisely of products of two Hilbert--Schmidt operators
(with intermediate Hilbert space allowed), and
$\|AB\|_1\leq\|A\|_2\|B\|_2$.
\end{theorem}
\begin{proof}
The singular expansion is a norm-absolutely summable series
$T=\sum_js_j(u_j\otimes v_j)$, where
$(u\otimes v)x=u\faip{v}{x}$. For every basis,
\[
 \sum_i|\faip{e_i}{Te_i}|
 \leq\sum_js_j\sum_i|\faip{e_i}{u_j}\faip{v_j}{e_i}|
 \leq\sum_js_j.
\]
Absolute summation and Parseval give
$\faTr T=\sum_js_j\faip{v_j}{u_j}$, independent of the basis.
In particular $\faTr(u\otimes v)=\faip{v}{u}$.
For finite orthonormal families $(a_i),(b_i)$,
$\sum_i|\faip{a_i}{Tb_i}|\leq\sum_js_j$ by the same two Bessel
inequalities. Taking the first singular pairs gives equality in the
supremum over all such finite families. This characterization proves
the triangle inequality for $\|\cdot\|_1$. The singular approximation
inequality proves its ideal bound.
A trace-norm Cauchy sequence converges in operator norm to a compact
$T$. For each finite pair of orthonormal families pass to the limit
in the characterization; taking the supremum gives
$\|T-T_n\|_1\leq\liminf_m\|T_m-T_n\|_1$. This proves completeness.
The rank-one trace formula gives cyclicity for each term of the
absolutely summable singular expansion, and the ideal bound permits
summation, proving cyclicity for every bounded factor.
If $A,B$ are Hilbert--Schmidt, an intermediate basis gives
$AB=\sum_i(Ae_i)\otimes(B^*e_i)$ and
$\sum_i\|Ae_i\|\|B^*e_i\|\leq\|A\|_2\|B\|_2$.
The finite-family characterization shows that any such nuclear
series is trace class with norm no greater than this sum.
Conversely $T=U|T|^{1/2}|T|^{1/2}$ factors into two Hilbert--Schmidt
operators, each with squared $S_2$ norm $\sum s_j$.
\end{proof}

\begin{lemma}[Weyl inequalities]\label{fa2:p2:weyl-inequalities}
Let $\lambda_j(T)$ be the nonzero eigenvalues of a compact operator,
counted with algebraic multiplicity and ordered by decreasing modulus.
Then
\[
 \prod_{j=1}^n|\lambda_j(T)|\leq\prod_{j=1}^ns_j(T),\qquad
 \sum_j|\lambda_j(T)|\leq\sum_js_j(T),
\]
and for $r\geq0$,
$\prod_j(1+r|\lambda_j|)\leq\prod_j(1+rs_j)$ whenever the right
side converges.
\end{lemma}
\begin{proof}
Choose an invariant $n$-dimensional subspace from the generalized
eigenspaces, selecting an initial Jordan flag if needed. The modulus
of its determinant is $\prod_{j\leq n}|\lambda_j|$.
The induced map on the $n$th exterior power has norm
$\prod_{j\leq n}s_j$: this follows by applying the singular system to
the orthonormal exterior basis. Its norm bounds the modulus of that
determinant, proving the product inequality.
For completeness, the passage from products to sums uses the
following finite majorization fact. If decreasing real lists satisfy
$\sum_{j\leq k}a_j\leq\sum_{j\leq k}b_j$ for every $k$, then
$\sum F(a_j)\leq\sum F(b_j)$ for every increasing convex $F$ with
$F(-\infty)=0$. Indeed it holds for $F(t)=(t-c)_+$ since its sum is
$\max_k(\sum_{j\leq k}a_j-kc)$; piecewise-linear convex interpolation
and monotone approximation give the general assertion. Apply this
to logarithms of eigenvalue moduli and singular values, padding zeros
by limits at $-\infty$. Taking $F(t)=e^t$ and
$F(t)=\log(1+re^t)$ gives the stated inequalities on finite lists.
Pass to the limit by monotone convergence.
\end{proof}

\begin{theorem}[Lidskii]\label{fa2:p2:lidskii}
For every trace-class operator on a complex Hilbert space,
\[
 \faTr T=\sum_j\lambda_j(T),
\]
with algebraic multiplicity and absolute convergence. Normality is
not required.
\end{theorem}
\begin{proof}
We first prove that a trace-class quasinilpotent operator $Q$ has
trace zero. Its active space $\overline{\faRan Q+\faRan Q^*}$ is
separable, so choose finite-rank orthogonal projections $P_n\to I$
there. Then $Q_n=P_nQP_n\to Q$ in trace norm: verify this on finite
rank operators and use their trace-norm density and $\|P_n\|\leq1$.
Thus $\faTr Q_n\to\faTr Q$ and $Q_n\to Q$ in operator norm.
The maximum modulus $\rho_n$ of eigenvalues of $Q_n$ tends to zero.
For a given $\delta>0$, the resolvent of $Q$ is bounded on
$\delta\leq|z|\leq2\|Q\|+1$; sufficiently small norm perturbations
have no spectrum there by the geometric inverse formula.
Put $D_n(z)=\det(I-zQ_n)$. On compact $z$-sets, for large $n$,
\[
 \frac{D_n'(z)}{D_n(z)}
 =-\sum_{k\geq1}z^{k-1}\faTr Q_n^k
 \longrightarrow-\faTr Q.
\]
Indeed for $k\geq2$, finite-dimensional trace and Weyl give
$|\faTr Q_n^k|\leq\rho_n^{k-1}\sum_j|\lambda_j(Q_n)|
\leq\rho_n^{k-1}\|Q\|_1$.
Integrating the logarithmic derivative with $D_n(0)=1$ yields
$D_n(z)\to e^{-z\faTr Q}$ locally uniformly.
Weyl and compression contraction give
$|D_n(z)|\leq\prod_j(1+|z|s_j(Q))$.
For every $M$ this product is at most
$p_M(|z|)\exp(|z|\sum_{j>M}s_j(Q))$, with $p_M$ a fixed polynomial.
If $\faTr Q\ne0$, take $z$ on the ray for which
$-z\faTr Q=|z||\faTr Q|$. Let $|z|\to\infty$ and then $M\to\infty$;
the resulting inequality $|\faTr Q|\leq\sum_{j>M}s_j(Q)$ is impossible.
Thus $\faTr Q=0$.
For general $T$, let $M$ be the closed span of all generalized
eigenvectors for its nonzero eigenvalues, and $K=M^\perp$.
Then $M$ is invariant for $T$ and $K$ for $T^*$. The compact
restriction $T^*|_K$ is quasinilpotent. To see this, if it had
nonzero eigenvalue $\bar\lambda$ and vector $u$, the adjoint Riesz
projection $P_\lambda^*$ would fix $u$ by its contour formula on that
eigenvector. But $P_\lambda^*u=0$ since $u$ annihilates
$\faRan P_\lambda$, the generalized eigenspace contained in $M$.
This is a contradiction. Both diagonal restrictions are trace class
by the ideal property. The first part gives zero trace on $K$.
Order the finite-dimensional generalized eigenspaces successively;
orthonormalizing their invariant flags gives an orthonormal basis of
$M$ in which the diagonal entries are exactly the eigenvalues with
their algebraic multiplicities. Absolute convergence of the diagonal
trace gives $\faTr(T|_M)=\sum_j\lambda_j(T)$.
Combining the two diagonal traces proves the theorem. Absolute
convergence was already established by Weyl's inequality.
\end{proof}

\begin{theorem}[Fredholm determinant in trace class]\label{fa2:p2:fredholm-determinant}
For $T\in S_1$,
\[
 \det(I+zT)=\sum_{n=0}^\infty z^n\faTr(\mathop{\bigwedge}\nolimits^n T)
 =\prod_j(1+z\lambda_j(T))
\]
is entire. It is continuous under trace-norm convergence locally
uniformly in $z$, vanishes exactly when $I+zT$ is not invertible,
and satisfies
$\det((I+S)(I+T))=\det(I+S)\det(I+T)$ for $S,T\in S_1$.
\end{theorem}
\begin{proof}
The singular values of $\bigwedge^nT$ are the products
$s_{j_1}\cdots s_{j_n}$, $j_1<\cdots<j_n$.
Thus the sum of their trace norms is bounded by
$\prod_j(1+|z|s_j)\leq e^{|z|\|T\|_1}$, proving entire convergence.
Here is a precise continuity estimate. For a finite sum of rank-one
operators $F=\sum_j u_j\otimes v_j$, expansion by principal minors gives
\[
 \det(I+F)=\sum_J\det[\faip{v_i}{u_j}]_{i,j\in J}.
\]
The exterior Gram inequality bounds each minor by
$\prod_{j\in J}\|u_j\|\|v_j\|$ (it is the inner product of the
two exterior products). Expand $\det(I+F+\varepsilon D)$ using
singular rank-one decompositions of the two finite-rank operators.
The terms containing exactly one factor from $\varepsilon D$ have
total absolute coefficient at most
$\|D\|_1\prod_j(1+s_j(F))\leq\|D\|_1e^{\|F\|_1}$;
the remaining terms have degree at least two in $\varepsilon$.
Thus along $F_t=(1-t)S+tT$,
$|\frac d{dt}\det(I+F_t)|\leq\|T-S\|_1e^{\|F_t\|_1}$.
Integration gives
\[
 |\det(I+T)-\det(I+S)|
 \leq\|T-S\|_1e^{\max(\|T\|_1,\|S\|_1)}
\]
for finite rank. This passes to arbitrary trace class as follows.
For each fixed $n$, antisymmetrization is a bounded projection on
the $n$-fold Hilbert tensor power, and telescoping tensor products
shows trace-norm continuity of $\bigwedge^nT$ under trace-norm
approximation. Its trace norm is at most $\|T\|_1^n/n!$.
Hence finite-rank approximations converge in each coefficient and
the factorial tails converge uniformly on bounded $z$-sets. They
therefore converge to the exterior-series definition of the determinant.
Passing to the limit proves the displayed continuity bound, locally
uniformly after replacing $S,T$ by $zS,zT$. The finite-dimensional
multiplicative identity passes to the limit too, since products of
the trace-class perturbations converge in trace norm.
For sufficiently small $z$, termwise differentiation, first for
finite rank and then by trace-norm approximation, gives
\[
 \log\det(I+zT)=\sum_{k\geq1}\frac{(-1)^{k+1}z^k}{k}\faTr T^k.
\]
Lidskii applied to $T^k$ and absolute summation identify this with
$\sum_j\log(1+z\lambda_j(T))$. The product is entire by
$\sum|\lambda_j|<\infty$, so the identity theorem proves the product
formula everywhere. Its zeros are exactly $-1/\lambda_j$, the same
as failure of invertibility by Riesz theory. This proves the remaining
assertions, without treating a merely continuous kernel as trace class.
\end{proof}

\begin{theorem}[Mercer and the diagonal trace]\label{fa2:p2:mercer}
Let $Q$ be compact metric, $\mu$ a finite positive measure with full
support, and $k$ continuous Hermitian on $Q\times Q$. If its integral
operator on $L^2(\mu)$ is positive, then
\[
 k(x,y)=\sum_j\lambda_j e_j(x)\overline{e_j(y)}
\]
absolutely and uniformly on $Q\times Q$, where $\lambda_j>0$ are
its nonzero eigenvalues and $e_j$ are normalized continuous
eigenfunctions. It is trace class and
$\faTr K=\int_Qk(x,x)\,d\mu(x)$.
More generally the same diagonal trace identity holds for any
trace-class integral operator with continuous kernel on this space,
without positivity or Hermitian symmetry.
\end{theorem}
\begin{proof}
Positivity of the operator implies positive definiteness of the
kernel matrices: approximate point masses at finitely many points
by nonnegative normalized continuous bumps of shrinking supports;
full support makes their integrals positive and uniform continuity
passes the quadratic integral inequality to the finite matrix.
Form the Hilbert space completion of the span of symbols $k_x$ with
$\faip{k_x}{k_y}=k(x,y)$, quotienting null vectors.
The identity
$\|k_x-k_y\|^2=k(x,x)+k(y,y)-2\operatorname{Re}k(x,y)$ shows that
$x\mapsto k_x$ is continuous. Vectors give continuous functions
$h(x)=\faip{k_x}{h}$. The map
$Sf=\int_Q k_x f(x)\,d\mu(x)$ is bounded, and $S^*h=h$ as an
$L^2$ function. Thus $S^*S=K$.
If $S^*h=0$, continuity and full support give $h(x)=0$ for all $x$,
so $h=0$ by density of the kernel vectors. Hence $\faRan S$ is dense.
The vectors $u_j=Se_j/\sqrt{\lambda_j}$ are orthonormal and total
in the kernel Hilbert space: the orthogonal complement is killed
by $S^*$, using the compact spectral expansion of $K$ and
$\ker S=\ker K$. Parseval therefore gives
$k(x,x)=\sum_j\lambda_j|e_j(x)|^2$ and the off-diagonal expansion.
The diagonal partial sums increase continuously to the continuous
function $k(x,x)$; Dini's theorem gives uniform convergence there.
Cauchy--Schwarz bounds every off-diagonal absolute tail by the square
root of the product of its two diagonal tails, proving uniform
absolute convergence. Integrating the diagonal gives
$\sum_j\lambda_j=\int k(x,x)<\infty$, and Lidskii or the singular
trace formula proves the trace identity.
For the nonpositive trace-class case choose nested finite measurable
partitions with mesh tending to zero, discarding null cells, and let
$P_n$ be conditional averaging onto their cell-constant functions.
Uniform continuity and density of continuous functions give
$P_n\to I$ strongly. Trace-norm approximation by finite rank gives
$\faTr(P_nKP_n)\to\faTr K$. The finite-dimensional trace is
$\sum_{A\in\mathcal P_n}\mu(A)^{-1}\int_{A\times A}k(x,y)\,d\mu(x)d\mu(y)$.
Its difference from $\int k(x,x)d\mu$ is bounded by $\mu(Q)$ times
the modulus of continuity of $k$ at the mesh size. This tends to zero,
proving the general diagonal identity.
\end{proof}

\begin{theorem}[Strong order positivity on $C(Q)$]\label{fa2:p2:perron-cq}
Let $Q$ be nonempty compact Hausdorff and $K:C(Q)\to C(Q)$ compact, preserving
real functions and satisfying $Kf>0$ everywhere for every nonzero
$f\geq0$. Then $r=r(K)>0$ is an algebraically simple eigenvalue
with a strictly positive eigenfunction. Every other eigenvalue has
modulus strictly less than $r$. There is a positive eigenmeasure
$\varphi$ for $K'$ and
$r^{-n}K^n\to P$ in operator norm, where
$Pf=u\varphi(f)/\varphi(u)$ and $Ku=ru$.
If $K1=1$, then $r=1$, $u=1$, and the normalized $\varphi$ is
the unique invariant probability measure.
\end{theorem}
\begin{proof}
$K1$ has a positive minimum $c$, so $K^n1\geq c^n1$ and $r\geq c>0$.
Riesz theory gives an eigenvector $z$ with eigenvalue $\lambda$ of
modulus $r$. Positivity implies $|Kz|\leq K|z|$: evaluate at each
point after rotating the complex phase and use
$\operatorname{Re}(\alpha z)\leq|z|$.
Put $y=|z|$ and $h=Ky-ry\geq0$. If $h\ne0$, then $Kh\geq\delta1$
for some $\delta>0$. Hence
$K(Ky)=rKy+Kh\geq(r+\delta/\|Ky\|_\infty)Ky$.
Iterating this inequality on the nonzero nonnegative vector $Ky$
would give a spectral radius larger than $r$, a contradiction.
Thus $Ky=ry$; strict positivity gives $y>0$.
For any peripheral eigenvector $z$ and a point $q$, choose a unit
scalar $\alpha$ so $\alpha(Kz)(q)=|(Kz)(q)|$.
The nonnegative function $|z|-\operatorname{Re}(\alpha z)$ has
$K$-value zero at $q$, since $K|z|=r|z|$ and $|\lambda|=r$.
Strict positivity forces that function identically zero.
Thus $\alpha z$ is real nonnegative everywhere, actually strictly
positive, and its eigenvalue must be the positive real $r$.
If $u,v>0$ are two eigenfunctions for $r$, put
$a=\min_Q u/v>0$. The nonnegative $u-av$ has a zero somewhere.
If nonzero its $K$-image is everywhere positive, contradicting
$K(u-av)=r(u-av)$. Hence the eigenspace is one-dimensional.
To obtain an eigenmeasure, choose a point evaluation $\delta_q$.
For $s>r$, $R(s,K)'\delta_q$ is positive by its Neumann series;
its norm tends to infinity, since its value on $u$ is
$u(q)/(s-r)$. Normalize to positive norm-one measures and take a
weak-star convergent subnet as $s\downarrow r$. Positivity makes
their norms their values on $1$, so the limit is still norm one.
The resolvent identity gives $K'\varphi=r\varphi$.
Since $u>0$, $\varphi(u)>0$.
If a generalized eigenvector satisfied $(K-r)w=u$, applying
$\varphi$ would give $0=\varphi(u)>0$.
The finite-dimensional Jordan structure therefore makes $r$
algebraically simple. Compactness and absence of other peripheral
eigenvalues give a spectral gap; the Riesz decomposition and the
radius formula make the complementary powers tend to zero after
division by $r^n$. Its one-dimensional spectral projection has the
stated formula. If $K1=1$, positivity gives
$\|K\|=\|K1\|_\infty=1$, so $r=1$ and $u=1$.
An invariant probability annihilates the complementary range and
has value one at $1$, so it must be $\varphi$.
\end{proof}

\begin{theorem}[Quadratic perturbations of an orthonormal basis]\label{fa2:p2:riesz-basis}
Let $(e_n)$ be an orthonormal basis and
$\sum_n\|y_n-e_n\|^2<\infty$. If the bounded map
$B(\sum c_ne_n)=\sum c_ny_n$ has zero kernel, then it is invertible
and $(y_n)$ is a Riesz basis. In particular the conclusion holds
without a separate independence assumption if
$\sum_n\|y_n-e_n\|^2<1$.
The zero-kernel condition means no nonzero square-summable infinite
linear relation; mere independence of finite subsets is insufficient.
\end{theorem}
\begin{proof}
Cauchy--Schwarz gives
$\|(B-I)\sum c_ne_n\|\leq(\sum|c_n|^2)^{1/2}
(\sum\|y_n-e_n\|^2)^{1/2}$.
The same bound on the tail shows $B-I$ is a norm limit of finite-rank
maps, hence compact. Riesz--Schauder makes injective $B=I+(B-I)$
surjective with bounded inverse. This is exactly the image-of-a-basis
definition of a Riesz basis, and gives the unique convergent expansion
of every vector with square-summable coefficients. In the small-sum
case $\|B-I\|<1$, so the Neumann series gives invertibility directly.
This is the precise almost-orthogonal basis variant in \cite{B08},
Chapter 22; it does not assert arbitrary perturbation stability.
\end{proof}

\begin{theorem}[Strong Krein--Rutman on a solid cone]\label{fa2:p2:krein-rutman}
Let $X$ be a real Banach space and $P\subset X$ a closed convex cone
with nonempty interior, $P\ne X$. If a compact linear $T$ satisfies
$T(P\setminus\{0\})\subset\operatorname{int}P$, then it has
$Tu_0=\mu_0u_0$ with $u_0\in\operatorname{int}P$, $\mu_0>0$.
Every eigenvector in $P\setminus\{0\}$ is a positive multiple of
$u_0$, with eigenvalue $\mu_0$. This eigenvalue is algebraically
simple, and every other nonzero eigenvalue of the complexification
has modulus strictly less than $\mu_0$. In particular $\mu_0$ is
the complex spectral radius. No normality of the cone is assumed.
\end{theorem}
\begin{proof}
First $\operatorname{int}P\cap(-P)=\varnothing$: otherwise adding
the negative of an interior point would put a ball about zero in
$P$, forcing $P=X$. Strong positivity then also forces pointedness:
if $x,-x\in P$ with $x\ne0$, then $Tx\in\operatorname{int}P$ and
$-Tx\in P$, a contradiction.
Fix $u\in\operatorname{int}P$. Closedness and $-u\notin P$ give
$\inf_{x\in P}\|x+u\|>0$. Also $Tu-ru\in P$ for some $r>0$,
since $Tu$ is interior. Write $x\geq y$ when $x-y\in P$.
If $T(x+\varepsilon u)=\lambda x$ for $x\in P$, then
$\lambda>0$ and $\lambda x\geq\varepsilon r u$.
Applying $T$ successively and using $Tx\leq\lambda x$ gives
$(\lambda/r)^n x\geq\varepsilon u$ for every $n\geq1$.
If $\lambda<r$, passage to the limit would put $-u\in P$.
Thus $\lambda\geq r$.
For $\varepsilon>0$ the map
$F_\varepsilon(x)=T((x+\varepsilon u)/\|x+\varepsilon u\|)$
is continuous on $P$ and has relatively compact image in
$P\cap\|T\|B_X$. Its denominator is bounded away from zero there
(indeed on all of $P$). Schauder on this closed bounded convex set
gives $x_\varepsilon=F_\varepsilon(x_\varepsilon)$.
Hence $T(x_\varepsilon+\varepsilon u)=\lambda_\varepsilon x_\varepsilon$,
$\lambda_\varepsilon=\|x_\varepsilon+\varepsilon u\|\geq r$ and
$\|x_\varepsilon\|\leq\|T\|$.
As $\varepsilon\downarrow0$, compactness and these bounds give a
subsequence with both $\lambda_\varepsilon$ and
$x_\varepsilon=(Tx_\varepsilon+\varepsilon Tu)/\lambda_\varepsilon$
convergent. Its limit satisfies
$Tu_0=\mu_0u_0$, $\mu_0=\|u_0\|\geq r$, so $u_0\ne0$ and
strong positivity makes $u_0$ interior.
If $Tx=\nu x$ with $0\ne x\in P$, then $\nu>0$ and $x$ is
interior. If $x$ is not proportional to $u_0$, let
$s=\max\{t\geq0:u_0-tx\in P\}$. It is finite and positive,
and $y=u_0-sx$ is a nonzero boundary point of $P$.
If $\nu\geq\mu_0$, then
$Ty=\mu_0y+s(\mu_0-\nu)x\leq\mu_0y$.
But $Ty$ is interior, and adding a positive vector keeps a point
interior; this would make $y$ interior, a contradiction.
Thus $\nu<\mu_0$. Reversing $x,u_0$ gives the opposite strict
inequality, again impossible. This proves uniqueness in the cone.
We prove the full spectral comparison without using a possibly
unavailable equivalent order norm on all of $X$.
On any finite-dimensional invariant real subspace $F$ containing
$u_0$, define
$M(x)=\inf\{c:x\leq cu_0\}$ and
$m(x)=\sup\{c:cu_0\leq x\}$.
They are finite and continuous because $u_0$ is interior; their
extreme inequalities are attained because $P$ is closed.
The oscillation $M(x)-m(x)$ induces a norm on $F/\mathbb Ru_0$:
pointedness makes its kernel exactly $\mathbb Ru_0$.
For a nonconstant class both $M(x)u_0-x$ and $x-m(x)u_0$ are
nonzero cone points. Strong positivity sends them to interior,
so the normalized $T/\mu_0$ strictly decreases oscillation.
On the compact unit sphere of this finite-dimensional quotient,
continuity gives a common contraction factor less than one.
For a real eigenvector outside $\mathbb Ru_0$ take its span with
$u_0$ as $F$; for a nonreal eigenvalue take also the real and
imaginary parts of an eigenvector of the complexification.
The induced quotient has that eigenvalue divided by $\mu_0$,
so its modulus is less than one. This proves the strict spectral
comparison and also $\ker(T-\mu_0)=\mathbb Ru_0$.
Finally if $(T-\mu_0)w=u_0$, then
$T^nw=\mu_0^nw+n\mu_0^{n-1}u_0$.
Interiorness gives $-cu_0\leq w\leq cu_0$ for some $c$.
Apply $T^n$ and divide by $\mu_0^n$; the upper inequality would
give $(c-n/\mu_0)u_0-w\in P$ for all $n$.
Divide by $n$ and close the cone to obtain $-u_0/\mu_0\in P$,
a contradiction. The finite-dimensional generalized eigenspace
from Riesz theory therefore has no nontrivial Jordan chain.
It is one-dimensional. All nonzero spectral points of the compact
complexification are eigenvalues, so the spectral-radius assertion
follows as well.
\end{proof}

\begin{remark}
Order positivity on $C(Q)$ means preservation of nonnegative real
functions. It differs from Hilbert quadratic-form positivity.
The strong-cone theorem and the preceding $C(Q)$ proof use their
order hypotheses explicitly; neither is inferred from a
nonselfadjoint kernel by calling it a positive Hilbert operator.
\end{remark}

% END INPUT


% BEGIN parts/part2-kernel-trace-criteria.tex
\chapter{Regularity criteria and counterexamples for integral traces}
\label{fa2:p2:kernel-trace-criteria-chapter}
The smooth-kernel sufficient condition in \cite[Chapter30, Theorem13]{B08}
is different from a trace formula which assumes trace class. The
H\"older version also underlies \cite[Chapter24, Theorem8]{B08}.
We prove that sufficient condition and give concrete obstructions
to replacing it by continuity alone. Lebesgue measure is used below.

\begin{theorem}[One-variable H\"older regularity implies trace class]
\label{fa2:p2:holder-trace-class}
Let $k\in C([0,1]^2)$ and suppose, for some $1/2<\alpha\le1$,
\[
 |k(x,y)-k(x',y)|\le H|x-x'|^\alpha\quad(x,x',y\in[0,1]).
\]
Then its $L^2$ integral operator $K$ is trace class, with
\[
 \|K\|_1\le\|k\|_{L^2([0,1]^2)}+
                          \frac{H}{1-2^{1/2-\alpha}}.
\]
The same conclusion holds for regularity in the second variable.
In particular a continuously differentiable kernel, and hence
a smooth kernel, is trace class. For its nonzero eigenvalues
$\lambda_j$ with algebraic multiplicity,
\[
 \int_0^1 k(t,t)dt=\operatorname{tr}K=\sum_j\lambda_j,
 \qquad D_k(z)=\prod_j(1-z\lambda_j).
\]
Both the eigenvalue sum and the product converge absolutely,
the latter locally uniformly in $z$. Affine change of variables
gives the corresponding statement on every compact interval.
\end{theorem}
\begin{proof}
Let $P_n$ be orthogonal projection onto the step functions
constant on the $n$ equal subintervals. The kernel of $P_nK$
is the average of $k$ in the first variable on its cell.
It differs pointwise from $k$ by at most $Hn^{-\alpha}$.
The Hilbert--Schmidt kernel identity gives
\[
                  \|(I-P_n)K\|_2\le Hn^{-\alpha}.
\]
If $s_j$ are the decreasing singular values, then for any
orthogonal projection $P$ of rank at most $n$,
$\sum_{j>n}s_j^2\le\|(I-P)K\|_2^2$.
Indeed in a singular system with left vectors $u_j$,
the latter squared norm is
$\sum_js_j^2(1-\|Pu_j\|^2)$. The numbers $\|Pu_j\|^2$
lie in $[0,1]$ and their sum is at most $n$ by Bessel.
For decreasing weights their weighted sum is at most
$\sum_{j\le n}s_j^2$: move any weight beyond $n$ to an
available place among the first $n$, first on finite sums
and then by increasing limits. This proves the tail bound.
Consequently, on the dyadic block $2^m<j\le2^{m+1}$,
\[
 \sum_{j=2^m+1}^{2^{m+1}}s_j
 \le2^{m/2}\Big(\sum_{j>2^m}s_j^2\Big)^{1/2}
 \le H2^{m(1/2-\alpha)}.
\]
Summing the geometric series and using $s_1\le\|K\|_2$
proves the asserted trace norm bound. Regularity in the second
variable applies this argument to $K^*$, whose trace norm is
the same. Continuous differentiability bounds the first partial
derivative on the square and gives the estimate with $\alpha=1$.
If an exponent larger than one is used, subdividing a segment
into $n$ pieces bounds its total variation by a constant times
$n^{1-\alpha}\to0$; the kernel is then constant in that
variable and the operator has rank at most one.

Now the nonpositive trace-class diagonal theorem
\ref{fa2:p2:mercer} and Lidskii give the sum formula and its
absolute convergence. The general classical product
\ref{fa2:p2:classical-fredholm-eigenvalues} is regularized by
$\exp(-z\int k(t,t)dt)$. The sum formula makes that exponential
cancel $\prod_je^{z\lambda_j}$: on finite products the remaining
factor is the exponential of the absolutely convergent tail.
It tends uniformly to one on bounded $z$ sets. Thus the product
is unregularized here, for a proved sufficient condition rather
than an unproved trace-class assumption.
\end{proof}

\begin{example}[Volterra is Hilbert--Schmidt but not trace class]
\label{fa2:p2:volterra-not-trace-class}
On real or complex $L^2(0,1)$, $Vf(x)=\int_0^x f(t)dt$ is
compact, Hilbert--Schmidt and quasinilpotent, but not trace class.
Its singular values are $((n+1/2)\pi)^{-1}$, $n\ge0$.
The ordinary trigonometric orthonormal basis has diagonal sum $1/2$
in either field, even though $V$ has no nonzero eigenvalue. Over
$\C$ the exponential basis has no unordered diagonal sum.
This retains the counterexample of \cite[Section30.5]{B08}.
\end{example}
\begin{proof}
Its kernel $1_{t<x}$ has squared $L^2$ norm $1/2$, proving
the first two properties. Fubini gives
$(V^*g)(t)=\int_t^1g(x)dx$. If $V^*Vf=\lambda f$,
$\lambda>0$, two weak derivatives give
$\lambda f''=-f$, $f'(0)=0$, $f(1)=0$.
The elementary ODE and these traces give exactly
$f(t)=c\cos((n+1/2)\pi t)$ and
$\lambda=((n+1/2)\pi)^{-2}$; substitution in the integral
equation verifies every one. The operator $V$ is injective
by its weak derivative, so $V^*V$ has zero kernel. Its compact
positive spectral theorem supplies no additional eigenvalues
or missing orthogonal part. Taking square roots proves the
singular values, whose sum diverges.

An eigen-equation $Vf=\lambda f$ with $\lambda\ne0$ gives
$\lambda f'=f$ and $f(0)=0$, hence $f=0$. Compact spectral
theory therefore gives $\sigma(V)=\{0\}$.
Finally the complete orthonormal system
$1,\sqrt2\cos(2\pi nt),\sqrt2\sin(2\pi nt)$ has first
diagonal entry $\int_0^1x\,dx=1/2$. Every other diagonal
entry is zero: integrate the products of cosine with its
sine primitive, and of sine with $1-\cos(2\pi nt)$,
over one period. This gives $1/2$, with no trace-class interpretation.
In the complex basis $e_n(t)=e^{2\pi int}$, for $n\ne0$ we have
$Ve_n=(e_n-1)/(2\pi in)$, so the diagonal entry is $(2\pi in)^{-1}$.
The absolute sum diverges, hence the unordered diagonal sum does not exist.
\end{proof}

\begin{example}[A continuous symmetric kernel which is not trace class]
\label{fa2:p2:continuous-kernel-not-trace-class}
There is a real continuous even function $f$ on the circle
such that the continuous symmetric kernel $k(x,y)=f(x-y)$
defines a compact self-adjoint operator on $L^2(\mathbb T)$
which is not trace class. Thus continuity and symmetry do
not replace either positivity in Mercer or regularity in
\ref{fa2:p2:holder-trace-class}.
\end{example}
\begin{proof}
Define polynomials $P_0=Q_0=1$ and
\[
 P_{r+1}(z)=P_r(z)+z^{2^r}Q_r(z),\qquad
 Q_{r+1}(z)=P_r(z)-z^{2^r}Q_r(z).
\]
Induction shows that each has degree $2^r-1$ and exactly
$2^r$ coefficients, all $\pm1$. On $|z|=1$,
$|P_{r+1}|^2+|Q_{r+1}|^2=2(|P_r|^2+|Q_r|^2)$,
so $\|P_r\|_\infty\le2^{(r+1)/2}$.
Choose $N_m=1+\sum_{j=1}^{m-1}2^{2j}$ and put
\[
 F(z)=\sum_{m\ge1}\frac{2^{-m}}{m^2}z^{N_m}P_{2m}(z),
 \qquad f(t)=\operatorname{Re}F(e^{2\pi it}).
\]
The supremum norm of the $m$th term is at most
$\sqrt2/m^2$, so this series converges uniformly and $f$
is real continuous. Its disjoint positive frequency blocks
have coefficients $\pm2^{-m}/m^2$ for $F$; taking real
parts divides these by two and gives the same real
coefficients at the opposite frequencies. In particular
$f$ is even. Uniform convergence permits termwise Fourier
coefficient integration. Hence
\[
 \sum_{n\in\mathbb Z}|\widehat f(n)|
       =\sum_{m\ge1}2^{2m}\frac{2^{-m}}{m^2}=\infty.
\]
Convolution with $f$ is diagonal on the complete Fourier
orthonormal basis, with real eigenvalues $\widehat f(n)$.
Its continuous square-integrable kernel makes it compact
and Hilbert--Schmidt. For this diagonal self-adjoint operator
the singular values are exactly the absolute Fourier
coefficients, with multiplicities. Their divergent sum proves
that it is not trace class. This gives the explicit continuous
counterexample behind the warning in \cite[Section30.6]{B08}.
\end{proof}

% END INPUT


% BEGIN parts/part2-spectral.tex
\chapter{C*-algebras and general spectral measures}
The Hilbert spaces here need not be separable unless explicitly stated.
The involution in an abstract C*-algebra is written $a\mapsto a^*$;
in a representation it becomes the Hilbert adjoint. Scalar measure
representation refers to \ref{fa2:p4:riesz-markov}. See \cite{B02,B04,B08}.

\begin{proposition}[The C*-unitization norm]\label{fa2:p2:cstar-unitization}
The algebraic unitization of a nonunital nonzero C*-algebra becomes a
C*-algebra with
$\|(a,\lambda)\|=\|L_a+\lambda I\|_{\faB(\mathcal A)}$, where
$L_ax=ax$. This is not in general the Banach-algebra sum norm.
It contains $\mathcal A$ isometrically as a closed ideal.
\end{proposition}
\begin{proof}
$\|L_a\|=\|a\|$: the reverse to the product bound follows by testing
$x=a^*/\|a\|$, using $\|aa^*\|=\|a\|^2$.
The representation is injective. Otherwise a nonzero scalar multiple
of $a$ would be a left identity; its adjoint is a right identity,
and their product makes the two equal, giving an identity in
$\mathcal A$, a contradiction.
Write $b=a+\lambda1$ algebraically. For $x\in\mathcal A$,
$\|bx\|^2=\|x^*b^*bx\|\leq\|x\|\|b^*bx\|$, so
$\|L_b\|^2\leq\|L_{b^*b}\|$.
Apply this also to $b^*$ and use submultiplicativity to get
$\|L_{b^*}\|^2\leq\|L_b\|\|L_{b^*}\|$, hence
$\|L_{b^*}\|\leq\|L_b\|$; reversing gives equality. Thus both
inequalities combine to the C*-identity.
Also $|\lambda|\leq\|L_a+\lambda I\|$. If this failed, $L_a$
would be invertible as a bounded operator by a Neumann series.
Solving $ae=a$ then gives $a(ey-y)=0$ for every $y$, so $e$ is a
left identity, the same contradiction. Therefore
$\|a\|+|\lambda|\leq3\|(a,\lambda)\|$, and the opposite norm
bound is the triangle inequality. Equivalence with the complete
sum norm proves completeness. For the zero algebra use $\mathbb C$;
if an algebra already has an identity and a larger unitization is
desired, use $\mathcal A\oplus\mathbb C$ with its direct-sum maximum
C*-norm under $(a,\lambda)\mapsto(a+\lambda1,\lambda)$.
\end{proof}

\begin{theorem}[Separable abelian von Neumann algebras]\label{fa2:p2:abelian-von-neumann}
Suppose $H$ is separable, over $\mathbb R$ or $\mathbb C$.
Every maximal abelian unital WOT-closed *-algebra has a cyclic
vector, and every abelian von Neumann algebra has a separating vector.
For complex $H$, every such maximal algebra is unitarily the scalar
multiplication algebra $L^\infty(\mu)$ on $L^2(\mu)$. For a bounded
complex normal $N$, there is a finite scalar control measure $\mu$
on $\sigma(N)$ such that
\[
 W^*(N)=\{f(N):f\text{ bounded Borel}\}\cong L^\infty(\mu).
\]
Under this isomorphism $\sigma(L^\infty,L^1)$ equals WOT on the
whole algebra.

For real $H$, use the canonical Hilbert complexification $H_{\mathbb C}$,
its conjugation $J$, and $\mathcal A_{\mathbb C}=\mathcal A+i\mathcal A$.
If $\tau$ is a measure-preserving involution, write
\[
 L^p_\tau(\mu)=\{g\in L^p(\mu;\mathbb C):
                   g(\tau x)=\overline{g(x)}\text{ a.e.}\},
 \qquad p=1,2,\infty,
\]
with their inherited real structures. Every real maximal abelian
algebra has the orthogonal scalar model of all multipliers from
$L^\infty_\tau(\mu)$ on $L^2_\tau(\mu)$, for a finite regular
Borel measure on a compact metrizable space and an involutive
measure-preserving homeomorphism $\tau$.
For a bounded real normal $N$, take
$X=\sigma_{\mathbb C}(N_{\mathbb C})$ and $\tau z=\bar z$.
There is a finite $\tau$-invariant scalar control measure $\mu$
for the complex spectral PVM $E_{\mathbb C}$, meaning
$E_{\mathbb C}(B)=0$ iff $\mu(B)=0$ for every Borel $B\subset X$, and
\[
 W^*_{\mathbb R}(N)=
 \{f(N_{\mathbb C})|_H:f\text{ bounded Borel},
                         f(\bar z)=\overline{f(z)}\}
       \cong L^\infty_\tau(\mu).
\]
Here the restrictions are real bounded operators. The real duality
$L^\infty_\tau(\mu)=(L^1_\tau(\mu))'$ uses
$(f,g)\mapsto\int fg\,d\mu\in\mathbb R$, and its weak-star topology
equals real WOT on the whole algebra. For real self-adjoint $N$,
$\tau$ is the identity on its spectrum and this is the usual real
scalar calculus. All these finite-control-measure assertions are
restricted to separable $H$; zero spaces and zero measures are included.
\end{theorem}
\begin{proof}
We first prove the scalar multiplication assertions over $\mathbb C$.
The commutant of all scalar multiplications on $L^2(\mu)$,
for finite $\mu$, is that same multiplication algebra. If $T$
commutes with them, put $h=T1$. Then $T1_B=1_Bh$, and the norm
bound on these vectors gives $|h|\leq\|T\|$ almost everywhere.
It follows on simple functions, then by density, that $T=M_h$.
For a sigma-finite measure partition into finite pieces and apply
the same proof on the reducing pieces.
Let $\mathcal M$ be maximal abelian. It equals its commutant: any
self-adjoint element of the commutant could otherwise be adjoined
to form a larger abelian *-algebra, and real and imaginary parts
handle every commutant element. Choose a maximal orthogonal family
of cyclic reducing spaces $H_j$ for $\mathcal M$.
Separability makes it countable, and maximality makes their sum $H$.
Their projections belong to $\mathcal M'=\mathcal M$.
Choose unit cyclic $e_j\in H_j$ and put $e=\sum_j2^{-j}e_j$.
Since $P_je=2^{-j}e_j$, $e$ is cyclic. Its scalar measure
representation identifies $\mathcal M$ with an algebra containing
continuous multiplications and their spectral projections, hence
all $L^\infty$ multiplications; the preceding commutant calculation
gives equality. Equivalently apply the cyclic PVM construction to
the commutative representation.
Extend any abelian von Neumann algebra $\mathcal A$ to a maximal
abelian one by Zorn. Its cyclic vector $e$ is separating for
$\mathcal A$: if $Ae=0$, commutation makes $A$ vanish on the dense
maximal-algebra orbit of $e$, so $A=0$.
Apply this to $W^*(N)$ and take $\mu=\mu_e$.
Then $E(B)=0$ iff $\mu(B)=0$, by the separating property.
Let $H_e=\overline{C^*(N)e}$ and $T\in W^*(N)={N,N^*}''$.
Its projection reduces $T$, and on $H_e\cong L^2(\mu)$ the
commutant calculation gives $T=M_f$ for some $f\in L^\infty(\mu)$.
For any $h\in H$, $\mu_h\ll\mu$. The isometry
$L^2(\mu_h)\to L^2(\mu)$, $g\mapsto g\sqrt{d\mu_h/d\mu}$,
intertwines coordinate multiplication. Transport it to the cyclic
spaces $H_h,H_e$ and extend by zero on $H_h^\perp$.
This partial isometry commutes with $N,N^*$, so $T$ commutes with it.
Therefore $T=f(N)$ on every $H_h$, hence on $H$.
A bounded completed-measure function has a bounded Borel representative
on the compact metric spectrum, so this is the asserted Borel calculus.
Every coefficient measure $\mu_{xy}$ has an $L^1(\mu)$ density,
giving weak-star continuity of every WOT coefficient. Conversely any
$g\in L^1(\mu)$ is a coefficient on $H_e$: take
$u=\sqrt{|g|}$, $v=\operatorname{sgn}(g)\sqrt{|g|}$, so
$\faip{u}{M_fv}=\int fg\,d\mu$. Thus the two families of defining
linear functionals coincide, proving equality of topologies.

For real $H$, a maximal orthogonal family of cyclic reducing spaces
for an abelian algebra is countable. A convergent weighted sum of
unit cyclic vectors is separating: its components are killed
separately, and commutation then makes each operator vanish on
every cyclic space. For a real maximal abelian algebra
$\mathcal A$, its commutant equals $\mathcal A$: split a commutant
operator into its symmetric and skew-adjoint parts, each of which
can be adjoined while preserving commutativity and the involution.
The projection onto $\overline{\mathcal A e}$ therefore belongs to
$\mathcal A$, so a separating $e$ is cyclic. The zero space is immediate.
Every complex operator on $H_{\mathbb C}$ splits uniquely into the
complexifications of two real operators. This identifies commutants
and both operator topologies under complexification; in particular
$\mathcal A_{\mathbb C}$ is a complex maximal abelian von Neumann
algebra, and $W^*_{\mathbb R}(N)_{\mathbb C}=W^*(N_{\mathbb C})$.
Choose a countable strongly dense subset of $\mathcal A_{\mathbb C}$'s unit ball and
norm-generate a commutative unital C*-algebra stable under $T\mapsto JTJ$.
Its compact metric character space has the involution induced by this
conjugation. The vector-state measure at the real cyclic vector $e$
is invariant under that involution by Riesz uniqueness. The cyclic
unitary sending $e$ to $1$ intertwines $J$ with
$g\mapsto\overline{g\circ\tau}$. Strong closure gives all complex
multipliers; taking fixed points gives exactly the asserted real model.

For real normal $N$, choose a real separating $e$ for $W^*_{\mathbb R}(N)$.
Its complexification is separating for $W^*(N_{\mathbb C})$.
Uniqueness of the complex normal PVM gives
$JE_{\mathbb C}(B)J=E_{\mathbb C}(\bar B)$.
Thus $\mu(B)=\|E_{\mathbb C}(B)e\|^2$ is finite, invariant under
conjugation, and has exactly the null sets of $E_{\mathbb C}$.
The complex Borel-image theorem identifies the fixed algebra with
$L^\infty_\tau(\mu)$; bounded Borel representatives may be
symmetrized by $f\mapsto(f+\overline{f\circ\tau})/2$.
The real $L^1$ predual follows by extending real functionals on its
fixed subspace complex-linearly and applying complex $L^1$ duality;
the fixed-space phase tests give the exact norm.
Every coefficient density for real vectors lies in $L^1_\tau(\mu)$,
by the PVM covariance and Radon--Nikodym uniqueness.
Conversely the cyclic subspace at $e$ has real model $L^2_\tau(\mu)$.
For $g\in L^1_\tau(\mu)$ the functions
$u=\sqrt{|g|}$ and $v=\operatorname{sgn}(g)\sqrt{|g|}$ both lie in
that real form and satisfy $\langle u,M_fv\rangle=\int fg\,d\mu$.
These are actual real vectors after transport. Equality of the
families of real coefficient functionals proves the whole-algebra
weak-star/WOT equality, including the zero-measure case.
\end{proof}

\begin{theorem}[Separable spectral multiplicity and the commutant]\label{fa2:p2:spectral-multiplicity}
A bounded normal operator on a separable complex Hilbert space is classified,
up to unitary equivalence, by a finite measure class $[\mu]$ on
$\mathbb C$ and a measurable multiplicity function
$m:\mathbb C\to\{1,2,\ldots,\infty\}$, modulo $\mu$-null sets.
It has the model $M_z$ on
$\int^\oplus\mathbb C^{m(z)}d\mu(z)$, supported on its spectrum.
Its commutant is exactly the essentially bounded measurable fields
$z\mapsto B(z)\in\faB(\mathbb C^{m(z)})$ acting fiberwise.
This classification statement is separable; the arbitrary-Hilbert
cyclic-sum theorem remains valid without this restriction.
\end{theorem}
\begin{proof}
Choose a countable cyclic reducing sum with scalar finite measures
$\mu_j$ and put
$\mu=\sum_j2^{-j}\mu_j/(1+\mu_j(\mathbb C))$.
Write $r_j=d\mu_j/d\mu$ and $A_j=\{r_j>0\}$.
The map $f_j\mapsto\sqrt{r_j}f_j$ identifies its summand with
$L^2(A_j,\mu)$ and preserves multiplication by $z$.
Thus the full space is the measurable subspace of
$L^2(\mu;\ell^2)$ whose $j$th coordinate vanishes off $A_j$.
At each $z$, enumerate the active indices $j$ in increasing order.
The $k$th active index is a measurable integer-valued function,
since the event that it equals $j$ is described by finitely many
membership conditions among $A_1,\ldots,A_j$.
Reordering these coordinates is therefore a measurable fiberwise
unitary onto the stated model, with $m(z)=\sum_j1_{A_j}(z)$.
For uniqueness, an intertwining unitary intertwines all spectral
projections by Fuglede--Putnam and Borel calculus. Therefore the
two scalar control measures have the same null sets; transport one
to the other by the square root of its Radon--Nikodym derivative.
We may now use the same finite $\mu$ on both sides.
Embed the two variable-multiplicity models into $L^2(\mu;\ell^2)$,
with coordinate supports $S_k=\{m\geq k\}$ and $T_j=\{n\geq j\}$.
Each matrix entry of the intertwiner, extended by zero to scalar
$L^2(\mu)$, commutes with every scalar multiplication. The preceding
commutant calculation makes it multiplication by a measurable
function $u_{jk}(z)$. Bounds on the operator imply bounds on the
matrix field: test each finite rational-coordinate vector times
$1_B$ for arbitrary measurable $B$. The integral norm inequality
forces the corresponding pointwise inequality almost everywhere;
there are only countably many such tests, so one common null set
works. Thus $u(z)$ is a contraction between the fibers.
The same applies to the adjoint. Testing the identities $U^*U=I$
and $UU^*=I$ on coordinate vectors, and passing to the countable
orthogonal sums of rows, gives
$u(z)^*u(z)=I_{m(z)}$, $u(z)u(z)^*=I_{n(z)}$ almost everywhere.
Hence it is a fiber unitary and $m=n$ almost everywhere.
Conversely equal measure classes and multiplicities give a unitary
by Radon--Nikodym transport and the identity on fibers.
The identical matrix-entry proof, without imposing unitarity,
describes all operators commuting with scalar multiplications as
essentially bounded measurable operator fields, their norm being
the essential supremum of the fiber norms. An operator commuting
with $M_z$ also commutes with $M_{\bar z}$ by Fuglede and hence
with all scalar Borel multiplications. Conversely every such field
commutes with $M_z$. This proves the commutant assertion.
\end{proof}

\begin{theorem}[Commutative C*-representation and permanence]\label{fa2:p2:commutative-cstar}
The Gelfand transform of a commutative unital C*-algebra is an isometric
*-isomorphism onto $C(\Delta)$. A unital C*-subalgebra has the same
spectra as its ambient C*-algebra. For normal $a$, continuous functions
on $\sigma(a)$ have a unique isometric unital *-calculus, generated by
$z\mapsto a$, with full spectral mapping.
\end{theorem}
\begin{proof}
For self-adjoint $b$, $e^{itb}$ is unitary, so its norm is one.
Holomorphic spectral mapping gives $|e^{it\lambda}|\leq1$ for every
$\lambda\in\sigma(b)$ and every real $t$, forcing $\lambda$ real.
For normal $a$, the C*-identity and commutation give
$\|a^2\|^2=\|(a^*a)^2\|=\|a^*a\|^2=\|a\|^4$.
Iteration and the radius formula give $r(a)=\|a\|$.
In the commutative case all elements are normal, so the Gelfand
transform is isometric. Characters take real values on self-adjoint
elements, hence preserve involution. Its range is closed, unital,
self-conjugate and separates the characters; Stone--Weierstrass makes
it all of $C(\Delta)$.
For a self-adjoint $b$ in a unital C*-subalgebra, its subalgebra
spectrum is real and so consists entirely of boundary points.
The boundary inclusion from spectral permanence for Banach subalgebras
then forces equality of spectra. If general $a$ is ambient-invertible,
$a^*a$ has no zero in its ambient spectrum, hence none in its
subalgebra spectrum; $(a^*a)^{-1}a^*$ is an inverse in the subalgebra.
For normal $a$, the algebra generated by $1,a,a^*$ is commutative.
Its characters are determined by their value at $a$, and its maximal
ideal space is homeomorphic to $\sigma(a)$: surjectivity is the
character spectral formula and permanence, and injectivity follows
from generation. The preceding representation is exactly the desired
calculus. The spectrum of a continuous function in $C(\sigma(a))$ is
its range; permanence proves full spectral mapping.
\end{proof}

\begin{theorem}[Positivity and order]\label{fa2:p2:cstar-positive}
For $a$ in a C*-algebra, the following are equivalent:
$a=a^*$ with $\sigma(a)\subset[0,\infty)$; $a=b^2$ with $b=b^*$;
$a=x^*x$ for some $x$. Positive elements form a norm-closed convex cone,
each has a unique positive square root, and
$0\leq a\leq b$ implies $\|a\|\leq\|b\|$.
For bounded complex-linear operators $A$ on a complex Hilbert space $H$,
positivity is equivalent to $\faip{x}{Ax}$ being a nonnegative real number for every $x\in H$.
\end{theorem}
\begin{proof}
Continuous calculus gives square roots and positive and negative parts
of every self-adjoint element. For self-adjoint $a$,
$a\geq0$ is equivalent to $\|a-t1\|\leq t$ for some
$t\geq\|a\|/2$. Indeed its real spectrum lies in $[0,2t]$ exactly
when this norm bound holds. This criterion shows closedness and
addition: if $\|a-t\|\leq t$ and $\|b-s\|\leq s$, then
$\|(a+b)-(t+s)\|\leq t+s$. Positive scalar multiplication is immediate.
It remains to prove that $x^*x$ is positive without assuming this
characterization in advance. Write $x^*x=u-v$ with $u,v\geq0$ and
$uv=0$, and put $y=xv^{1/2}$. Then $y^*y=-v^2\leq0$.
The nonzero spectra of $y^*y$ and $yy^*$ agree: $1-ab$ invertible
iff $1-ba$ invertible, with inverse
$1+b(1-ab)^{-1}a$ for the latter. Hence $yy^*\leq0$ too.
Write $y=b+ic$ with $b,c$ self-adjoint. Then
$b^2+c^2+i(bc-cb)=-v^2\leq0$, so
$i(bc-cb)\leq0$, because squares of self-adjoints are positive by
their continuous calculus. But
$b^2+c^2-i(bc-cb)=yy^*\leq0$ gives the opposite inequality.
The cone is proper by its spectral definition, so the commutator
is zero, and $v^2=0$. The C*-identity gives $v=0$.
This proves positivity of $x^*x$, and the three characterizations.
If $c\geq0$ and $c^2=a$, it commutes with $a$ and with its positive
square root $b$. The commutative calculus for $b,c$ gives pointwise
nonnegative functions with equal squares, so $b=c$.
For operators on a complex Hilbert space, $A\geq0$ gives nonnegative
quadratic forms. Conversely, real-valuedness of $\faip{x}{Ax}$ for every
$x$ implies self-adjointness by complex polarization. A negative spectral
point would give, by continuous calculus, a nonzero spectral vector on which the form is negative.
Thus the operator is positive. For abstract order norm monotonicity
one may use the faithful representation constructed below; alternatively
the norm criterion gives $\|a\|\leq\|b\|$ by spectral comparison in
the positive cone. In a faithful representation the direct proof is
$\|a\|=\sup_{\|x\|=1}\faip{x}{ax}
\leq\sup_{\|x\|=1}\faip{x}{bx}=\|b\|$.
The construction of that representation below uses only positivity
and the norm criterion, not this monotonicity assertion.
\end{proof}

\begin{proposition}[Closed ideals and automatic norm control]\label{fa2:p2:cstar-quotient}
Every closed two-sided ideal $I$ in a C*-algebra is *-invariant and
has an approximate identity of positive contractions. The quotient
with its quotient norm is a C*-algebra. Every *-homomorphism of
C*-algebras is contractive, isometric when injective, and has closed
range.
\end{proposition}
\begin{proof}
For $a\in I$, the positive $a^*a$ lies in $I$ even before *-invariance
is known. The functions $t/(t+\varepsilon)$ vanish at zero, so their
calculus values $e_\varepsilon$ lie in $I$. Moreover
$e_\varepsilon a^*\to a^*$: the square of the error norm is bounded
by $\sup_{t\geq0}t\varepsilon^2/(t+\varepsilon)^2\leq\varepsilon/4$.
Thus $a^*\in I$.
For a finite set $F\subset I$ put
$b=\sum_{a\in F}(a^*a+aa^*)$ and $e=b(b+\varepsilon)^{-1}$.
It is a positive contraction in $I$. For $a\in F$,
$\|a(1-e)\|^2\leq\|(1-e)b(1-e)\|\leq\varepsilon/4$;
the left estimate is the same. The norm comparison uses order
monotonicity obtained from the faithful GNS representation in the
next theorem. That construction uses neither ideals nor quotient
norms, so this forward reference is not circular.
Directing finite sets by inclusion and
$\varepsilon\downarrow0$ gives an approximate identity.
For it, $\|a+I\|=\lim\|a(1-e)\|$: the lower bound is immediate;
for $i\in I$, $a(1-e)=(a+i)(1-e)-i(1-e)$ gives the upper bound
after an infimum over $i$. Hence
\[
 \|a+I\|^2=\lim\|(1-e)a^*a(1-e)\|=\|a^*a+I\|,
\]
where the upper inequality uses any representative $a^*a+i$ and
$\|(1-e)i(1-e)\|\to0$, and the lower inequality is the quotient
norm bound. This proves the C*-identity in the quotient.
For a unital *-homomorphism $\pi$,
$\|\pi(a)\|^2=r(\pi(a^*a))\leq r(a^*a)=\|a\|^2$.
Unitization gives the nonunital case. If $\pi$ is injective and a
self-adjoint spectrum shrank strictly, choose a nonzero continuous
function supported in the missing part. Its value at $a$ would be
nonzero but map to zero, since polynomial approximation intertwines
the calculi. Thus self-adjoint spectra, hence all norms, are preserved.
Factor through the C*-quotient by the closed kernel; the resulting
injection is isometric, so its complete image is closed.
\end{proof}

\begin{theorem}[States and GNS]\label{fa2:p2:gns}
Every positive functional $\varphi$ on a unital C*-algebra satisfies
$|\varphi(b^*a)|^2\leq\varphi(a^*a)\varphi(b^*b)$ and
$\|\varphi\|=\varphi(1)$. A state is a positive functional with value
one at the identity. Each state has a cyclic representation
$\pi_\varphi:\mathcal A\to\faB(H_\varphi)$ with unit cyclic vector
$\xi_\varphi$ and
$\varphi(a)=\faip{\xi_\varphi}{\pi_\varphi(a)\xi_\varphi}$,
unique up to a unitary carrying the cyclic vector to the cyclic vector.
The direct sum of the GNS representations of all states is faithful.
\end{theorem}
\begin{proof}
Apply positivity to $(a+tb)^*(a+tb)$ for arbitrary complex $t$;
minimizing its scalar quadratic polynomial proves Schwarz, including
the case $\varphi(b^*b)=0$ by letting $|t|\to\infty$.
For $0\leq a\leq\|a\|1$, positivity gives
$\varphi(a)\leq\|a\|\varphi(1)$; alternatively this inequality
follows immediately because $\|a\|1-a$ is positive by calculus.
Schwarz with $b=1$ and $a^*a\leq\|a\|^21$ gives the norm formula.
On $\mathcal A/N$, $N=\{a:\varphi(a^*a)=0\}$, set
$\faip{[a]}{[b]}=\varphi(a^*b)$.
Schwarz makes this well defined. Also
$\varphi(a^*c^*ca)\leq\|c\|^2\varphi(a^*a)$, since
$a^*(\|c\|^21-c^*c)a$ is positive. Thus $N$ is a left ideal and
$[a]\mapsto[ca]$ is bounded by $\|c\|$.
After completion these operators define a unital *-representation;
the adjoint identity follows from
$\varphi((ca)^*b)=\varphi(a^*c^*b)$. The class of $1$ is cyclic.
Any other cyclic realization has the isometric dense map
$\pi_\varphi(a)\xi_\varphi\mapsto\pi(a)\xi$, which extends uniquely
to the required unitary.
For existence of sufficiently many states, choose $a\ne0$ and an
evaluation state on the commutative algebra generated by $a^*a,1$
at its largest spectral value $\|a\|^2$. Hahn--Banach extends it to
a functional of norm one on $\mathcal A$ still taking value one at $1$.
Such a functional is positive: for self-adjoint $b$,
$|1+it\varphi(b)|\leq\|1+itb\|
\leq(1+t^2\|b\|^2)^{1/2}$ forces $\varphi(b)$ real by using both
signs and letting $t\to0$. If $0\leq b\leq1$,
$|1-\varphi(b)|\leq\|1-b\|\leq1$, so $\varphi(b)\geq0$.
Rescaling gives positivity for every positive $b$.
The chosen state has $\|\pi_\varphi(a)\xi_\varphi\|^2
=\varphi(a^*a)=\|a\|^2$, so the sum representation is faithful
(and in fact isometric). This also completes any order norm comparison
deferred to a faithful representation above.
\end{proof}

\begin{definition}[Projection-valued measure]\label{fa2:p2:pvm-definition}
A PVM on a measurable space $(S,\Sigma)$ is a map
$E:\Sigma\to\faB(H)$ with $E(S)=I$, each $E(B)$ an orthogonal
projection, $E(B)E(C)=E(B\cap C)$, and
$E(\bigcup_n B_n)x=\sum_nE(B_n)x$ for disjoint $B_n$, with norm
convergence for every $x$. Put
$\mu_{xy}(B)=\faip{x}{E(B)y}$ and $\mu_x=\mu_{xx}$.
For bounded measurable $f$, uniformly approximating it by simple
functions defines $f(E)$, a contractive *-homomorphism with
\[
 \faip{x}{f(E)y}=\int f\,d\mu_{xy},\qquad
 \|f(E)x\|^2=\int|f|^2\,d\mu_x.
\]
\end{definition}
\begin{proof}
For disjoint simple functions, orthogonality gives
$\|\sum c_jE(B_j)x\|^2=\sum|c_j|^2\mu_x(B_j)$.
This proves the norm bound, independence of the simple representation
and the norm-square identity. Common refinements verify multiplication
and involution for simple functions, and uniform approximation
preserves them. Polarization gives the scalar integral formula.
The scalar complex measure has variation at most $\|x\|\|y\|$:
for a finite partition use
$\sum|\faip{E(B_j)x}{E(B_j)y}|
\leq(\sum\|E(B_j)x\|^2)^{1/2}
(\sum\|E(B_j)y\|^2)^{1/2}$.
\end{proof}

\begin{theorem}[Representations give PVMs]\label{fa2:p2:representation-pvm}
If $S$ is compact Hausdorff and $\pi:C(S)\to\faB(H)$ is a unital
*-representation, there is a unique regular Borel PVM $E$ such that
$\pi(f)=\int f\,dE$ for continuous $f$. The representation extends
to all bounded Borel functions as a *-homomorphism.
\end{theorem}
\begin{proof}
Scalar Riesz representation gives regular measures $\mu_{xy}$ for
$f\mapsto\faip{x}{\pi(f)y}$, with total variation at most
$\|x\|\|y\|$. For bounded Borel $g$, the forms
$\int g\,d\mu_{xy}$ are bounded sesquilinear forms, hence define
operators $\widetilde\pi(g)$.
To prove multiplicativity, regard $\mu\mapsto\int g\,d\mu$ as an
element of $C(S)''$ of norm at most $\|g\|_\infty$.
Goldstine gives a net of continuous $g_i$, uniformly bounded by this
norm (or an arbitrarily small excess), converging against every
regular measure to $g$. Thus $\pi(g_i)\to\widetilde\pi(g)$ weakly.
For fixed continuous $f$, convergence also holds against $f\mu_{xy}$,
so $\widetilde\pi(gf)=\widetilde\pi(g)\pi(f)
=\pi(f)\widetilde\pi(g)$. For bounded Borel $f$, approximate it in
the same way and test against $g\mu_{xy}$; this is again a finite
regular Borel measure. A second weak limit proves the product formula
for both bounded Borel factors. Conjugation of the scalar measures
proves the adjoint identity.
Set $E(B)=\widetilde\pi(1_B)$. The algebra and adjoint identities make
these orthogonal projections. For disjoint $B_n$, scalar measure
additivity and orthogonality give
$\|(E(\bigcup B_n)-\sum_{n\leq N}E(B_n))x\|^2
=\mu_x(\bigcup_{n>N}B_n)\to0$.
This proves strong countable additivity. Uniqueness follows from
uniqueness of each scalar regular measure, followed by polarization.
\end{proof}

\begin{theorem}[Bounded normal spectral theorem]\label{fa2:p2:normal-spectral}
Every bounded normal $N$ on an arbitrary complex Hilbert space has a
unique regular Borel PVM on $\sigma(N)$ with
$N=\int z\,dE(z)$. Its bounded Borel calculus is a *-homomorphism,
$\|f(N)\|$ is the $E$-essential supremum of $|f|$, and
$\sigma(f(N))$ is its $E$-essential range. If $f_n\to f$ pointwise
$E$-almost everywhere and $\sup_n\|f_n\|_\infty<\infty$, then
$f_n(N)\to f(N)$ strongly. A bounded operator commuting with both
$N,N^*$ commutes with all its spectral projections.
\end{theorem}
\begin{proof}
Apply the commutative C*-calculus and the preceding representation
theorem to the inclusion of $C^*(I,N)$ into $\faB(H)$.
Uniqueness follows since any proposed PVM agrees on polynomials in
$z,\bar z$, hence on continuous functions by Stone--Weierstrass,
then on all scalar measures by Riesz uniqueness.
The norm-square identity proves the norm bound by essential supremum.
For its reverse choose a nonzero vector in a spectral set where
$|f|$ exceeds a given smaller constant. For the spectral formula,
if $|f-\lambda|$ is bounded away from zero almost everywhere its
reciprocal gives an inverse. Otherwise choose unit vectors in
$E\{|f-\lambda|<1/n\}H$, obtaining approximate eigenvectors, so
$\lambda$ is spectral. Dominated convergence in each finite measure
$\mu_x$ proves strong convergence. Commutation first holds for
continuous functions by polynomial approximation; the bounded
continuous nets in the PVM construction pass it to Borel functions.
\end{proof}

\begin{theorem}[Cyclic models, including nonseparable spaces]\label{fa2:p2:cyclic-model}
If $e$ is cyclic for $C^*(I,N)$, then $N$ is unitarily equivalent to
$M_z$ on $L^2(\mu_e)$, with $e$ corresponding to $1$.
In general $H$ is an orthogonal direct sum of cyclic reducing spaces,
with no separability assumption; equivalently $N$ is multiplication
by a bounded function on a disjoint union of finite-measure spaces.
That union need not be sigma-finite and there need not be one finite
scalar control measure on $\sigma(N)$ for the whole representation.
\end{theorem}
\begin{proof}
The map $f\mapsto f(N)e$ on continuous functions has squared norm
$\int|f|^2\,d\mu_e$. Continuous functions are dense in $L^2(\mu_e)$
by regularity and simple-function approximation. Thus it extends to
an isometry onto the closed cyclic space and intertwines multiplication
by $z$ with $N$. That space is invariant under both $N,N^*$, so
is reducing. A maximal orthogonal collection of nonzero cyclic
reducing spaces exists by Zorn. Its orthogonal complement is reducing;
if nonzero it contains another cyclic space, contradicting maximality.
Their direct sum therefore equals $H$. On the disjoint union assign
the sum measure: the measure of a set is the supremum of the finite
partial sums of its component measures. Its $L^2$ space is exactly
the Hilbert direct sum (each square-integrable vector has at most
countably many nonzero components). The component coordinate
functions give the multiplication symbol. This proves the stated
model. To see the two possible obstructions, use $H=\ell^2(\mathbb R)$
with its unit vectors $e_t$ and the bounded self-adjoint diagonal operator
$Ne_t=(\arctan t)e_t$. Its spectral projection on each singleton
$\{\arctan t\}$ is nonzero. A finite scalar measure can give positive
mass to only countably many singletons (for each positive integer $n$
there are only finitely many of mass at least $1/n$). Hence no finite
scalar measure controls this spectral measure. The coordinate
representation uses counting measure on an uncountable set and is
not sigma-finite. In fact a sigma-finite multiplication representation
would give a finite control measure: partition the multiplication space
into measurable finite-measure pieces $A_n$ and weight its measure on
$A_n$ by $2^{-n}/(1+\mu(A_n))$ for $n\ge1$. The resulting finite measure
has the same null sets. Its pushforward under the multiplication
symbol controls the projections, since the projection on a Borel
set is multiplication by the indicator of its inverse image. This
contradicts the singleton projections.
\end{proof}

\begin{theorem}[Polar decomposition]\label{fa2:p2:polar}
For $T\in\faB(H,K)$ let $|T|=(T^*T)^{1/2}$. There is a unique
partial isometry $U:H\to K$ with initial space $(\ker T)^\perp$,
final space $\overline{\faRan T}$ and $T=U|T|$.
For normal $T$ it commutes with $|T|$ and $U$; for invertible $T$,
$U$ is unitary. For a general operator it need not be unitary.
\end{theorem}
\begin{proof}
Define $U(|T|x)=Tx$. Equality $\||T|x\|=\|Tx\|$ makes this
well defined and isometric on $\faRan|T|$. Extend it to its closure
and set it zero on $\ker|T|=\ker T$. The range and uniqueness follow
from this construction. If $T$ is normal its PVM gives
$|T|=|z|(T)$ and $U=u(T)$ with $u(z)=z/|z|$ off zero and $u(0)=0$,
which proves commutation. If $T$ is invertible both initial and
final spaces are the whole spaces.
\end{proof}

\begin{theorem}[Fuglede--Putnam]\label{fa2:p2:fuglede-putnam}
For bounded normal $N,M$ and bounded $T$ with $NT=TM$, one has
$N^*T=TM^*$. Consequently $T$ intertwines all corresponding Borel
spectral functions. Similar bounded normal operators are unitarily
equivalent.
\end{theorem}
\begin{proof}
The operator-valued entire function
$F(z)=e^{zN^*}Te^{-zM^*}$ can, using $NT=TM$, be written
$e^{zN^*-\bar zN}T e^{\bar zM-zM^*}$.
The two outside factors are unitary, since their exponents are
skew-adjoint and the normal factors commute. Thus $\|F(z)\|\leq\|T\|$.
Scalar Liouville and separation make $F$ constant. Differentiation
at zero gives $N^*T=TM^*$. Polynomial approximation followed by the
PVM bounded-net construction gives intertwining for Borel functions.
If $T$ is invertible and intertwines $M,N$, then $T^*T$ commutes with
$M,M^*$, and its positive invertible square root does too. In
$T=U|T|$, the unitary $U$ therefore still intertwines $M,N$.
\end{proof}

\begin{theorem}[Operator topologies and double commutant]\label{fa2:p2:double-commutant}
The strong operator topology (SOT) is pointwise norm convergence and
the weak operator topology (WOT) is convergence of all matrix
coefficients. On norm-bounded sets adjunction is WOT continuous and
multiplication is separately SOT and WOT continuous. A unital
*-subalgebra $\mathcal A\subset\faB(H)$ satisfies
$\overline{\mathcal A}^{\mathrm{SOT}}
=\overline{\mathcal A}^{\mathrm{WOT}}=\mathcal A''$.
No separability assumption is required.
\end{theorem}
\begin{proof}
The continuity assertions follow by testing a vector or a coefficient;
for SOT multiplication by a fixed factor, use its norm bound.
Commutants are WOT closed, so both closures are contained in
$\mathcal A''$. If $T\in\mathcal A''$ and $x\in H$, the projection
onto $\overline{\mathcal A x}$ commutes with $\mathcal A$, because
the subspace reduces a *-algebra. Hence it commutes with $T$, and
since $I\in\mathcal A$, $Tx\in\overline{\mathcal A x}$.
For simultaneous approximation of $x_1,\ldots,x_n$, use the diagonal
algebra on $H^n$. An operator commuting with all diagonal $A$ has
matrix entries in $\mathcal A'$. Thus diagonal $T$ commutes with
that entire commutant. Applying the one-vector argument to
$(x_1,\ldots,x_n)$ produces $A\in\mathcal A$ with
$\sum_j\|(T-A)x_j\|^2<\varepsilon^2$. These are exactly the SOT
neighborhood conditions. Thus $T$ lies in the strong closure.
The WOT closure lies between strong closure and $\mathcal A''$, so
all three agree.
\end{proof}

\begin{proposition}[Mean ergodic theorem]\label{fa2:p2:mean-ergodic}
If $T$ is a contraction on a Hilbert space, then
$N^{-1}\sum_{k=0}^{N-1}T^k$ converges strongly to the orthogonal
projection onto $\ker(I-T)$. For a strongly continuous contraction
semigroup the same holds for $t^{-1}\int_0^tT(s)\,ds$ onto the common
fixed-vector space.
\end{proposition}
\begin{proof}
If $Tx=x$, contractivity and $\faip{x}{T^*x}=\|x\|^2$ force
$T^*x=x$ by equality in Cauchy--Schwarz; the reverse is identical.
Thus $\ker(I-T)=\ker(I-T^*)=\overline{\faRan(I-T)}^\perp$.
The averages fix this kernel and send $(I-T)y$ to
$(y-T^Ny)/N\to0$. Their norms are at most one, so convergence
extends from the range to its closure and hence to all of $H$.
For the semigroup, the orthogonal complement of the common fixed
space is the closure of the linear span of $(I-T(s))H$, by the same
adjoint fixed-space identity for each $s$. Its average on such a
vector is a difference of two intervals of length $s$, with norm at
most $2s\|y\|/t$. Uniform boundedness of the averages again extends
convergence from their span to its closure.
\end{proof}

\begin{corollary}[Spectral families and normal corollaries]\label{fa2:p2:spectral-corollaries}
For bounded normal $N$ on a complex Hilbert space, self-adjointness is equivalent to a real
spectrum, and unitarity to spectrum contained in $\mathbb T$.
Also $E(\{\lambda\})H=\ker(N-\lambda)$.
For bounded normal $N$ on a real Hilbert space, $N$ is self-adjoint
exactly when $\sigma(N_{\mathbb C})\subset\mathbb R$, and is orthogonal
exactly when $\sigma(N_{\mathbb C})\subset\mathbb T$.
For real $\lambda$, the real eigenspace identifies with the fixed-real
part of $E_{\mathbb C}(\{\lambda\})H_{\mathbb C}$.
For self-adjoint $A$ over either scalar field the spectral family
$E_t=E(({-}\infty,t])$ is increasing and strongly right continuous,
with limits zero and $I$ at the two infinities. If its spectrum lies
in $[m,M]$, Riemann--Stieltjes sums for
$\int_{m-0}^M t\,dE_t$ converge in operator norm to $A$;
the corresponding scalar matrix coefficients are ordinary
Riemann--Stieltjes integrals. The lower endpoint includes its atom.
\end{corollary}
\begin{proof}
In the PVM model $N^*$ has symbol $\bar z$, so a real spectrum gives
$N=N^*$ and a unit-circle spectrum gives $N^*N=I=NN^*$.
The converse spectral containments follow from real self-adjoint
spectra and the inverse formula for unitaries, or the same norm-square
identity. A vector is killed by $N-\lambda$ exactly when its spectral
measure is concentrated at $\lambda$, proving the eigenprojection
formula. For real normal operators apply these complex tests to the
canonical Hilbert complexification; it preserves adjoints and products,
and real eigenvectors identify with the fixed-real vectors in the
complex eigenspace. For a real self-adjoint operator the unique
complex spectral PVM preserves that fixed-real subspace on real-axis
Borel sets, giving its real PVM and bounded calculus.
Monotonicity and right continuity of $E_t$ follow from
intersection of the decreasing measurable intervals and strong
countable additivity. For a partition the Stieltjes sum is the
calculus of a step function differing from $t$ by at most its mesh
on the spectrum. The operator norm error is therefore at most that
mesh. The scalar assertion follows by applying a coefficient
functional to these sums. The value at $m-0$ ensures the interval
containing $m$ includes $E\{m\}$.
\end{proof}

% END INPUT


% END INPUT


% BEGIN parts/part2-ap.tex
\chapter{A Banach space without the approximation property}
\label{fa2:ap:chapter}

The statement that the approximation property is not automatic is a structural
counterexample, not an exercise following from the definition. We give the
construction needed here. Its finite Walsh blocks follow Enflo's original
construction \cite{Enflo1973}. We use an absolutely convergent rank-one series
to finish the argument; in particular, no unproved equivalence between AP and
BAP on reflexive spaces is needed.

\section{An obstruction expressed by an absolutely convergent series}

Recall that a Banach space $X$ has the approximation property (AP) if, for every
compact $K\subset X$ and $\eps>0$, there is a bounded finite-rank $T:X\to X$
with $\sup_{x\in K}\|Tx-x\|<\eps$. No uniform bound on $\|T\|$ is imposed.

\begin{lemma}[Trace obstruction]\label{fa2:ap:trace-obstruction}
Suppose $f_j\in X'$, $x_j\in X$, and
\[
 \sum_j\|f_j\|\|x_j\|<\infty,\qquad
 \sum_j f_j(x)x_j=0\quad(x\in X),\qquad
 \sum_j f_j(x_j)\ne0.
\]
Then $X$ does not have AP.
\end{lemma}
\begin{proof}
Discard zero terms. Put $b_j=\|f_j\|\|x_j\|$. Choose numbers $a_j\to\infty$
with $\sum a_jb_j<\infty$: choose increasing integers $m_r$ with
$\sum_{j\ge m_r}b_j<2^{-r}/r^2$, and take $a_j=r$ on successive blocks,
enlarging the initial finite block if necessary. Set
$y_j=x_j/(a_j\|x_j\|)$ and $g_j=a_j\|x_j\|f_j$. Then $y_j\to0$,
$K=\{0,y_1,y_2,\ldots\}$ is compact, and $\sum\|g_j\|<\infty$.
For a finite-rank operator write $T=\sum_{r=1}^q h_r(\cdot)z_r$ with
$h_r\in X'$. Absolute convergence gives
\[
 \sum_jg_j(Ty_j)=
 \sum_{r=1}^q h_r\left(\sum_jg_j(z_r)y_j\right)=0.
\]
If AP held, choose $T$ approximating the identity within $\eps$ on $K$.
It follows that
$|\sum_j f_j(x_j)|=|\sum_jg_j(y_j-Ty_j)|
 \le\eps\sum_j\|g_j\|$.
Letting $\eps\downarrow0$ contradicts the last hypothesis.
\end{proof}

We shall also use the following elementary way to handle rank-one series.
For a finite formal tensor $u=\sum f_j\otimes x_j$ set
\[
 \nu(u)=\inf\sum_j\|f_j\|\|x_j\|,
 \qquad \operatorname{tr}u=\sum_jf_j(x_j),
 \qquad \widehat u(x)=\sum_jf_j(x)x_j.
\]
The infimum runs over finite algebraic representations of the tensor, rather
than representations merely of its associated operator. The distinction is
essential: injectivity of the completed tensor-to-operator map is precisely
what must not be assumed in this argument. Both the trace and the associated
operator have norm at most $\nu(u)$. If
$u_1+\sum_{m\ge1}v_m$ has $\sum_m\nu(v_m)<\infty$, choose finite
representations within $2^{-m}\eta$ of the respective infima. Flattening the
series gives an absolutely summable rank-one series. Its trace and operator
are the limits of those of its finite partial sums. This assertion follows
directly from the two preceding norm bounds and needs no tensor-product theorem.

\section{Two adjacent Walsh levels}

Fix $n\ge2$, $H_n=\{0,1\}^{2n}$, and
$\chi_S(a)=(-1)^{\sum_{j\in S}a_j}$ for $S\subset\{1,\ldots,2n\}$.
The characters are orthonormal for the uniform probability on $H_n$.
Let $W_{n,r}$ be the set with $|S|=r$, and put
$F_{n,r}=\sum_{|S|=r}\chi_S$. The two levels $r=n-1,n+1$ have the same
cardinality $t(n)=\binom{2n}{n-1}$.

\begin{lemma}[A small Walsh difference]\label{fa2:ap:walsh}
For
\[
 h_n=t(n)^{-1}(F_{n,n-1}-F_{n,n+1})
\]
one has $\|h_n\|_\infty=2/n$.
\end{lemma}
\begin{proof}
The coefficient identity
$\sum_r F_{n,r}(a)z^r=(1-z)^q(1+z)^{2n-q}$, where $q=\sum a_j$,
shows that $F_{n,r}(0)=\binom{2n}{r}$ and
$F_{n,r}(a)=(1-r/n)\binom{2n}{r}$ if $q=1$.
Thus $|h_n(a)|=2/n$ at $q=1$. Replacing $a$ by its complementary vector
multiplies both odd/even levels by the same sign, so also handles $q=2n-1$;
at $q=0,2n$ the difference is zero.

For $2\le q\le2n-2$, the coefficient integral on $|z|=1$ bounds
$|F_{n,n-1}(a)|$ by
\[
 \frac{2^{2n}}{2\pi}\int_0^{2\pi}
 |\sin(\theta/2)|^q|\cos(\theta/2)|^{2n-q}\,d\theta.
\]
The integrand for an intermediate $q$ is a geometric interpolation of the
integrands for $q=2$ and $q=2n-2$. H\"older's inequality bounds its integral
by their common value. The resulting bound is
$\binom{2n}{n}/(2n-1)\le t(n)/n$ (the equality before the inequality can
also be checked by the beta integral or its integration-by-parts recurrence).
Complementing a subset $S$ gives
$F_{n,n+1}(a)=(-1)^q F_{n,n-1}(a)$.
Consequently $|F_{n,n-1}-F_{n,n+1}|\le2t(n)/n$ in this remaining range.
\end{proof}

Let $E_n\subset C(H_n;\R)$ be the span of these two levels. The finite group
$G_n$ generated by translations of $H_n$ and permutations of its $2n$
coordinates acts isometrically on $E_n$. For an endomorphism $S$ of $E_n$,
write $\tau_-(S),\tau_+(S)$ for the average of its diagonal coefficients on
the two Walsh levels. Averaging $U^{-1}SU$ over $G_n$ kills every off-diagonal
coefficient: distinct characters have distinct translation signs, whose
average product is zero. Coordinate permutations make the remaining diagonal
coefficient constant on each level. Evaluation at $0\in H_n$ therefore gives
the exact finite-dimensional identity
\begin{equation}\label{fa2:ap:finite-twirl}
 \tau_-(S)-\tau_+(S)
 =\frac1{|G_n|}\sum_{U\in G_n}(U^{-1}SUh_n)(0).
\end{equation}
This identity, not just an operator-norm inequality, will control the rank-one
representations below.

\section{The incidence sets}

Fix $a=6/5$ and take $n_m=\lfloor a^{m+m_0}\rfloor$ with $m_0$ sufficiently
large. Put $t_m=t(n_m)$ and $k_m=\lfloor t_m^2\rfloor=t_m^2$.
Increasing $m_0$ is permitted in every estimate below.
Let $M_m$ be a set of $k_mt_m$ symbols arranged in $k_m$ rows of $t_m$
symbols. There exist $k_{m+1}$ subsets $N_{m,j}\subset M_m$, each of size
$t_{m+1}$, with
\begin{align}
 |N_{m,i}\cap N_{m,j}|&\le t_{m+1}/(2n_{m+1})\quad(i\ne j),
 \label{fa2:ap:intersection}\\
 |N_{m,j}\cap\text{any row of }M_m|&\le1,
 \label{fa2:ap:row}\\
 \sum_{e\in M_m}\left|\frac1{k_mt_m}
 -\frac{r_m(e)}{k_{m+1}t_{m+1}}\right|&\le\frac1{n_{m+1}},
 \quad r_m(e)=|\{j:e\in N_{m,j}\}|.
 \label{fa2:ap:balance}
\end{align}
Here is an explicit construction and the estimates establishing existence.
Write $t=t_m$, $s=t_{m+1}$, $k=k_m$, $q=k_{m+1}$,
$L=\lfloor k/s\rfloor$, and $\ell=Ls$. Identify the first $\ell$ rows with
$\{1,\ldots,\ell\}\times\Z/t\Z$. For each integer $\rho\ge0$, list
\[
 (r,r\rho+\kappa\pmod t),\qquad
 \kappa=0,\ldots,t-1,\quad r=1,\ldots,\ell,
\]
in that order. Cut this list into successive sets of $s$ elements, and use
the first $q$ such sets over successive $\rho$. A set has one point per row,
proving \eqref{fa2:ap:row}. Put $v=\lfloor q/(Lt)\rfloor+1$; all used
values of $\rho$ are less than $v$.

Within one value of $\rho$ the sets are disjoint. For two different values,
common rows satisfy
$r(\rho_1-\rho_2)\equiv\kappa_2-\kappa_1\pmod t$.
Two such rows differ by at least $t/|\rho_1-\rho_2|\ge t/v$.
The row interval of either set has length $s$, so the intersection has at most
$1+sv/t$ elements. Stirling's estimate
$\log t(n)=n\log4+O(\log n)$ gives $s=t^{a+o(1)}$,
$L=t^{2-a+o(1)}$, $q=t^{2a+o(1)}$ and
$v=t^{3a-3+o(1)}$. Since $4-3a>0$, both $v/t=o(1/n_{m+1})$ and
$1/s=o(1/n_{m+1})$. Thus
$1+sv/t\le s/(2n_{m+1})$ for all $m$, after taking $m_0$ large.

Each full $\rho$-list covers the active $\ell t$ points once. In a partial
list, a point occurs either once or zero times. Hence
$|r_m(e)-qs/(\ell t)|\le1$ on the active points and $r_m(e)=0$ elsewhere.
The left side of \eqref{fa2:ap:balance} is at most
\[
 \frac{2(k-\ell)}k+\frac{\ell t}{qs}
 \le\frac{2s}{k}+\frac{Lt}{q}.
\]
The first term is $t^{a-2+o(1)}$ and the second
$t^{3-3a+o(1)}$. Both are $o(1/n_{m+1})$, establishing the balance bound.
All inequalities are uniform after removal of a finite initial segment.

\section{The Banach space and the rank-one series}

Let $K_m$ be the disjoint union of $k_m$ copies of $H_{n_m}$, one for each
row of $M_m$. In the real reflexive Banach space
\[
 Y=\left(\bigoplus_{m\ge1}C(K_m;\R)\right)_{\ell^2},
\]
define one vector $e$ for each symbol in $M_m$. On its own row in $K_m$ it
is its distinct character from $W_{n_m,n_m+1}$. On row $j$ of $K_{m+1}$ it
is zero if $e\notin N_{m,j}$ and otherwise its distinct character from
$W_{n_{m+1},n_{m+1}-1}$, using any bijection for that row. It vanishes on
all other blocks. Let $X$ be the closed span of these countably many vectors
in $Y$. Each summand of $Y$ is finite dimensional. Here is the dual
identification used to infer reflexivity. If $F\in Y'$, let $F_m$ be its
restriction to the $m$th summand. Testing finitely many coordinates on
vectors nearly norming $F_m$, with scalar lengths proportional to
$\|F_m\|$, gives $\sum_{m\in A}\|F_m\|^2\le\|F\|^2$ for every
finite set $A$. Conversely Cauchy--Schwarz makes every sequence
$(F_m)$ with square-summable norms a bounded functional by
$F(y)=\sum_m F_m(y_m)$, with the matching norm; finite-support vectors
are dense, so these descriptions are inverse. Applying the same
description to the dual sum identifies $Y''$ with the sum of the
finite-dimensional double duals, and its canonical embedding with
the componentwise onto embeddings. Thus $Y$ is reflexive. A closed
subspace of a reflexive space is reflexive by weak compactness of
its unit ball, and the countable generating family makes $X$
separable. Thus $X$ is separable and reflexive.

Within any row all the participating characters are distinct: the upper
degree comes from its own $M_m$ row, and the lower degree from $N_{m-1,j}$.
Orthogonality in the uniform $L^2$ norm therefore proves, for any finite sum,
\begin{equation}\label{fa2:ap:coordinate}
 \left\|\sum_e c_e e\right\|_X\ge |c_e|\|e\|_X
 \quad\text{for each one of its symbols }e.
\end{equation}
Indeed, in each block in which $e$ is nonzero, one incident row has supremum
norm at least $|c_e|$. The block norm is a maximum over its rows, not a sum;
$e$ itself has block norm one. Squaring and summing over its one or two
nonzero blocks proves the inequality. This also proves linear independence.
The coordinate functionals
$f_e(\sum c_d d)=c_e$ extend continuously to $X$ with
$\|f_e\|\|e\|=1$. Define the finite algebraic tensors
\[
 u_m=\frac1{k_mt_m}\sum_{e\in M_m}f_e\otimes e,
 \qquad \operatorname{tr}u_m=1,\qquad \nu(u_m)\le1.
\]

\begin{lemma}[Summable adjacent differences]\label{fa2:ap:differences}
One has $\nu(u_{m+1}-u_m)\le5/n_{m+1}$.
\end{lemma}
\begin{proof}
For a finite set $A$ of symbols put
$u_A=|A|^{-1}\sum_{e\in A}f_e\otimes e$.
The balance estimate and \eqref{fa2:ap:coordinate} give
\[
 \nu\left(u_m-\frac1{k_{m+1}}\sum_j u_{N_{m,j}}\right)
 \le1/n_{m+1}.
\]
Also $u_{m+1}$ is the average of $u_{M_{m+1,j}}$ over its rows.
It is therefore enough to prove
$\nu(u_{N_{m,j}}-u_{M_{m+1,j}})\le4/n_{m+1}$.

Put $n=n_{m+1}$, $s=t_{m+1}$ and
$E=\operatorname{span}(N_{m,j}\cup M_{m+1,j})$.
Restriction $R:X\to C(H_n)$ to row $j$ of block $K_{m+1}$ has norm at
most one and range exactly $E_n$. Its restriction $R_E$ is an isomorphism
from $E$ onto $E_n$, identifying its two groups of symbols with the two
Walsh levels. For $U\in G_n$ set
\[
 x_U=R_E^{-1}Uh_n,\qquad \lambda_U(z)=(U^{-1}Rz)(0).
\]
Then $\|\lambda_U\|\le1$ and $\|x_U\|\le4/n$.
To verify the latter bound explicitly, the coefficients of $x_U$ are signed
$1/s$ on each of its two groups. In its distinguished row its norm is
$2/n$. In another row of the same block its norm is bounded by
$|N_{m,j}\cap N_{m,i}|/s\le1/(2n)$.
In the preceding block the row estimate gives a bound $1/s$, and in the
following block it gives $1/s$ again. There are no other nonzero blocks.
Thus the squared $Y$ norm is at most $(2/n)^2+2/s^2\le(4/n)^2$.

Expanding the finite Walsh identity \eqref{fa2:ap:finite-twirl} in its
matrix coefficients now gives the algebraic tensor identity
\[
 u_{N_{m,j}}-u_{M_{m+1,j}}
 =\frac1{|G_n|}\sum_{U\in G_n}\lambda_U\otimes x_U.
\]
For clarity, every $\lambda_U$ is a linear combination of the finitely many
$f_e$ with $e\in N_{m,j}\cup M_{m+1,j}$, because only those symbols have
nonzero restriction to that row. Translation averaging kills the mixed
character coefficients and permutation averaging gives the normalized
diagonals $1/s$ and $-1/s$. Thus this is a finite algebraic identity, not an
inference that the tensor-to-operator map is injective. Taking the displayed
rank-one representation proves its $4/n$ bound and the lemma.
\end{proof}

\begin{theorem}[Nonautomatic approximation property]\label{fa2:ap:counterexample}
There exists a separable real reflexive Banach space without AP, and consequently
without a Schauder basis. There also exists a complex Banach space without AP.
\end{theorem}
\begin{proof}
Since $\sum 1/n_m<\infty$, the series
$u_1+\sum_m(u_{m+1}-u_m)$ admits an absolutely summable rank-one
representation as above. Call the resulting formal series $u$. Its trace is
$\lim_m\operatorname{tr}u_m=1$.

For any bounded operator $T$ define
$\tau_m(T)=(k_mt_m)^{-1}\sum_{e\in M_m}f_e(Te)$.
Its modulus is at most $\|T\|$. If $T$ is finite rank, approximate its
finitely many range vectors by finite combinations of the generating symbols.
This gives finite-rank $T_\eps$ with $\|T-T_\eps\|<\eps$ and range in
one finite span of symbols. For all sufficiently large $m$,
$\tau_m(T_\eps)=0$, hence $\tau_m(T)\to0$.
Take the rank-one $T(z)=f(z)x$. Then
$\tau_m(T)=f(\widehat u_m(x))\to0$ for every $f\in X'$ and $x\in X$.
The rank-one series converges in operator norm, so
$f(\widehat u(x))=0$ for all $f,x$. Hahn--Banach implies $\widehat u=0$.
The trace obstruction applies and proves that this $X$ lacks AP.

If $(e_j)$ were a Schauder basis, its partial-sum projections $P_n$ would
converge pointwise to the identity and be uniformly bounded by uniform
boundedness. A finite-net argument would give uniform convergence on every
compact set, proving AP. Thus no such basis exists.

Finally complexify $X$ with, for example,
$\|x+iy\|=\sup_\theta\|x\cos\theta-y\sin\theta\|$.
This is a complete complex norm equivalent to the product norm. Its real
part contains $X$ as a complemented real subspace. AP passes to a complemented
subspace: compose a compact-set approximant with its bounded projection and
inclusion. Complex finite-rank approximants are also real finite rank.
Hence AP of this complexification would imply AP of $X$, a contradiction.
\end{proof}

The conclusion is an existence theorem. It does not claim that all subspaces
of a space with AP inherit AP, nor that a particular compact operator is
approximable merely because its domain and range have familiar notation.

% END INPUT


% BEGIN parts/part2-extension-geometry.tex
\chapter{Extension, almost orthogonality, and nonlinear Hilbert geometry}
These statements retain important result-bearing problems from
\cite{B01}, rather than treating a chapter heading as coverage of its
exercises. The variational statements use real Banach spaces; the
Hilbert assertions also apply to complex spaces by real restriction.

\begin{theorem}[Br\o ndsted--Rockafellar approximation]
\label{fa2:p2:brondsted-rockafellar}
Let $X$ be a real Banach space and $\phi$ a proper lower
semicontinuous convex function. Put
\[
 \partial_\varepsilon\phi(x)=
 \{f\in X':\phi(y)\geq\phi(x)+f(y-x)-\varepsilon
                                      \text{ for all }y\}.
\]
If $f_0\in\partial_\varepsilon\phi(x_0)$, $\varepsilon>0$,
then for every $\lambda>0$ there are $x_1$ and
$f_1\in\partial\phi(x_1)$ with
\[
 \|x_1-x_0\|\leq\varepsilon/\lambda,
 \qquad \|f_1-f_0\|\leq\lambda.
\]
Moreover $\partial_\varepsilon\phi(x)$ is nonempty for every
$x\in D(\phi)$ and $\varepsilon>0$, and
\[
 \overline{D(\partial\phi)}=\overline{D(\phi)},\qquad
 \overline{\operatorname{Ran}\partial\phi}
                         =\overline{D(\phi^*)}
\]
in the respective norm topologies.
\end{theorem}
\begin{proof}
Set $\psi=\phi-f_0$. The epsilon-subgradient condition says
$\psi(x_0)\leq\inf\psi+\varepsilon$, so $\psi$ is bounded
below. The quantitative Ekeland principle
(\ref{fa2:p1:ekeland}) gives $x_1$ at distance at most
$\varepsilon/\lambda$ for which
$\psi(y)+\lambda\|y-x_1\|\geq\psi(x_1)$ for all $y$.
The subgradient sum rule from Fenchel--Rockafellar applies because
the norm term is continuous everywhere. Thus at $x_1$,
$0=f_1-f_0+g$ with $f_1\in\partial\phi(x_1)$ and
$g\in\lambda\partial\|\cdot-x_1\|(x_1)=\lambda B_{X'}$.
The last identity follows directly by testing the defining
subgradient inequality on all directions and their negatives.
It proves the desired norm estimate, including the constants.

Fenchel--Moreau says
$\phi(x)=\sup_f(f(x)-\phi^*(f))$. At a finite value choose
a functional within $\varepsilon$ of this supremum; it belongs
to $\partial_\varepsilon\phi(x)$ by Fenchel's inequality.
For any $x\in D(\phi)$, applying the first part with
$\varepsilon\downarrow0$ and, for example,
$\lambda=\sqrt\varepsilon$ proves the first closure identity.
For $f_0\in D(\phi^*)$, choose $x_0$ within $\varepsilon$
of the defining conjugate supremum. Then
$f_0\in\partial_\varepsilon\phi(x_0)$ and the same choice of
$\lambda$ approximates $f_0$ in norm by an actual subgradient.
The reverse inclusions follow from
$\phi(x)+\phi^*(f)=f(x)$ at every subgradient pair.
\end{proof}

\begin{theorem}[Kirszbraun--Valentine]
\label{fa2:p2:kirszbraun}
Let $H_1,H_2$ be real or complex Hilbert spaces, $S\subset H_1$,
and $f:S\to H_2$ have Lipschitz constant $L\geq0$. It extends
to all of $H_1$ with that same constant. No dimension,
separability, boundedness or closedness assumption on $S$ is needed.
\end{theorem}
\begin{proof}
The empty domain has the constant zero extension. If $L=0$ and
$S$ is nonempty the map is constant, so again the conclusion is
immediate. Divide values by $L$ to reduce the remaining case to
$L=1$. To adjoin one point $x$, it suffices to intersect the balls
$\overline B(f(s),\|s-x\|)$ for $s\in S$.

First take finitely many points $s_1,\ldots,s_n$, writing
$y_i=f(s_i)$ and $r_i=\|s_i-x\|$. Minimize
$F(y)=\max_i(\|y-y_i\|^2-r_i^2)$.
Orthogonal projection onto the finite-dimensional linear span
of the $y_i$ cannot increase this function. On that span it is
continuous and coercive, so it attains a minimum at $y$.
For the active indices, $0$ lies in the convex hull of $y-y_i$.
Otherwise finite-dimensional separation gives a direction whose
inner product with every active gradient is strictly negative;
a small move lowers all active expressions, and the finitely
many inactive expressions remain strictly below the old maximum.
This contradicts minimality. Hence there are $t_i\geq0$,
$\sum t_i=1$, supported on the active indices, with
$y=\sum t_i y_i$. The variance identity gives
\begin{align*}
 F(y)&=\sum_i t_i\|y_i-y\|^2-\sum_i t_i\|s_i-x\|^2\\
 &=\frac12\sum_{i,j}t_it_j
       (\|y_i-y_j\|^2-\|s_i-s_j\|^2)
                   -\left\|\sum_i t_i s_i-x\right\|^2\leq0.
\end{align*}
The Lipschitz hypothesis proves the inequality. Thus every finite
set of the balls intersects.

If there are any balls, intersect them inside one fixed ball.
That ball is weakly compact by Hilbert reflexivity and all the
balls are weakly closed by convexity. The finite intersection
property proves intersection of the entire family, including
uncountably many centers. This gives a one-point extension.
Partially order all extensions of $f$ by inclusion. A chain has
its union as an upper bound with the same Lipschitz estimate.
Zorn gives a maximal extension; a missing point could be adjoined
by the proved intersection result, a contradiction. This completes
the extension on all of $H_1$. The variance argument uses real
parts in the complex case and is unchanged.
\end{proof}

\begin{theorem}[Cotlar--Stein almost orthogonality]
\label{fa2:p2:cotlar}
For a finite family of bounded operators $T_i$ on a Hilbert space put
\[
 A=\sup_i\sum_j\|T_i^*T_j\|^{1/2},\qquad
 B=\sup_i\sum_j\|T_iT_j^*\|^{1/2}.
\]
Then $\|\sum_iT_i\|\leq\sqrt{AB}$.
For an infinite family the same bound holds for every finite partial
sum if these row sums are uniformly bounded. If those sums converge
strongly on a dense subspace, they therefore extend to a strong
sum on the entire Hilbert space with that bound.
\end{theorem}
\begin{proof}
Write $a_{ij}=\|T_i^*T_j\|^{1/2}$ and
$b_{ij}=\|T_iT_j^*\|^{1/2}$. Both matrices are symmetric by
adjunction. Let $C=\max_i\|T_i\|$ and let $n$ be the number
of terms. If either row bound is zero all operators vanish.
Otherwise expand $(S^*S)^m$, $S=\sum_iT_i$.
The term with indices $i_1,\ldots,i_{2m}$ has norm at most
$\prod_{k=1}^m a_{i_{2k-1},i_{2k}}^2$ by grouping adjacent
star/nonstar pairs. Grouping the other way, with two end factors,
bounds its norm by
$C^2\prod_{k=1}^{m-1} b_{i_{2k},i_{2k+1}}^2$.
The geometric mean of these two bounds is
\[
 C a_{i_1i_2}b_{i_2i_3}a_{i_3i_4}\cdots a_{i_{2m-1}i_{2m}}.
\]
Summing successively from the last index bounds the full sum by
$nC A^mB^{m-1}$. Since $S^*S$ is positive,
$\|S\|^{2m}=\|(S^*S)^m\|\leq nC A^mB^{m-1}$.
Taking $2m$th roots and letting $m\to\infty$ gives the assertion.
The same estimate applies to every finite subfamily of the infinite
family. A uniform operator bound and convergence on a dense set
make the partial sums Cauchy on every vector by a three-error
approximation; completeness gives the strong sum. No unconditional
strong convergence is inferred from row bounds alone.
\end{proof}

\begin{theorem}[Normal structure and nonexpansive fixed points]
\label{fa2:p2:kirk-fixed-point}
A bounded convex set has normal structure if every convex subset
of positive diameter contains a point whose supremum distance from
that subset is strictly smaller than its diameter. Every nonempty
weakly compact convex set with normal structure has a fixed point
for every nonexpansive self-map. In particular every nonempty
closed bounded convex subset of a uniformly convex Banach space
has this property (Browder--G\"ohde--Kirk).
\end{theorem}
\begin{proof}
Among nonempty weakly compact convex invariant subsets, a decreasing
chain has nonempty invariant intersection by compactness. Zorn
therefore gives a minimal one $K$. Also
$K=\overline{\operatorname{conv}}T(K)$: the latter is a nonempty
closed convex subset of $K$ and is invariant because its points
still lie in $K$, so their images lie in $T(K)$.
If $d=\operatorname{diam}K>0$, normal structure gives $z\in K$
and $r<d$ with $\sup_{y\in K}\|z-y\|\leq r$.
The set $K_r=\{z\in K:\sup_{y\in K}\|z-y\|\leq r\}$ is
nonempty, weakly compact and convex. For $z\in K_r$,
$\|Tz-Ty\|\leq\|z-y\|\leq r$ for every $y\in K$;
the convex-hull identity extends this bound to every point of $K$.
Thus $K_r$ is invariant. It is proper, since otherwise $d\leq r$,
contradicting minimality. Hence $d=0$ and $K$ is a fixed point.

For a bounded convex subset $D$ of a uniformly convex space, of
diameter $d>0$, choose $x_0,x_1\in D$ with distance greater than
$d/2$. For every $y\in D$, the two vectors
$(x_i-y)/d$ are in the unit ball and separated by more than $1/2$.
The unit-ball modulus of uniform convexity gives
$\|(x_0+x_1)/2-y\|\leq(1-\delta(1/2))d$.
Their midpoint is in $D$, proving normal structure. Uniform
convexity implies reflexivity by the earlier Milman--Pettis theorem;
the specified closed bounded convex set is therefore weakly compact.
Apply the first assertion. This proof does not assume weak
continuity of the nonlinear map.
\end{proof}

\begin{theorem}[Alternating convex projections and the subspace improvement]
\label{fa2:p2:alternating-projections}
If $C_1,C_2$ are nonempty closed convex subsets of a Hilbert space
with nonempty intersection, the alternating projection iterates
converge weakly to a point of the intersection. If both are closed
linear subspaces, the iterates converge in norm to the orthogonal
projection of the initial vector onto their intersection.
\end{theorem}
\begin{proof}
Put $x_{2n+1}=P_{C_1}x_{2n}$ and
$x_{2n+2}=P_{C_2}x_{2n+1}$. For each $p\in C_1\cap C_2$,
the variational projection inequality gives
\[
 \|x_{n+1}-p\|^2+\|x_{n+1}-x_n\|^2\leq\|x_n-p\|^2.
\]
Thus the iterates are bounded, every distance to $p$ has a limit,
and successive differences tend to zero by summation. Both
distances to $C_i$ tend to zero. Every weak subsequential limit
lies in both sets: their distance functions are convex norm
continuous, hence weakly lower semicontinuous by separation.
Such limits exist by reflexivity and Eberlein--Smulian.
If $p,q$ are two weak limits, the convergent scalar sequence
$\|x_n-p\|^2-\|x_n-q\|^2$ has respective subsequence limits
$-\|p-q\|^2$ and $\|p-q\|^2$, so $p=q$. If the sequence
failed to converge weakly to that point, a subsequence outside
a fixed weak neighborhood would have a further weakly convergent
subsequence, the same point, a contradiction.

For subspaces $M,N$, restrict the positive self-adjoint contraction
$Q=P_NP_M$ to $N$. Its powers converge strongly by its PVM:
$t^n\to1_{\{1\}}(t)$ on $[0,1]$, with bounded domination.
Its fixed space is $M\cap N$, since equality in the two projection
norm contractions forces membership in both. The corresponding
limit is $P_{M\cap N}$ on $N$. Apply it to the first iterate
that lies in $N$; its projection onto the intersection equals
that of the original vector because each previous orthogonal
projection preserves that component. The intervening $M$
iterates have the same norm limit by continuity. This proves the
strong subspace assertion without claiming norm convergence for
arbitrary convex pairs.
\end{proof}

% END INPUT


% BEGIN parts/part2-projection-displacement.tex
\chapter{Alternating projections without an intersection}
\label{fa2:p2:projection-displacement-chapter}
The missing disjoint-set variants of \cite[Problem 27(A),(B)]{B01}
are separated from the previously proved intersection theorem.
All inner-product inequalities below use their real parts in a
complex Hilbert space.

\begin{theorem}[The minimum-displacement vector]
\label{fa2:p2:projection-displacement}
Let $C_1,C_2$ be nonempty closed convex subsets of a real or
complex Hilbert space $H$. Starting from any $u_0\in H$, set
$u_{2n-1}=P_{C_1}u_{2n-2}$ and $u_{2n}=P_{C_2}u_{2n-1}$.
Write
\[
 D=\overline{C_2-C_1},\qquad
 p=P_D0,\qquad \delta=\operatorname{dist}(C_1,C_2).
\]
Then $u_{2n}-u_{2n-1}\to p$ in norm. Neither a nonempty
intersection, a bounded orbit, nor attainment of the distance
is required. The closure in the definition of $D$ is essential.
\end{theorem}
\begin{proof}
For projection $P_C$, its variational inequality applied to
$x,y$ gives
\[
 \|P_Cx-P_Cy\|^2
 \le \operatorname{Re}\langle P_Cx-P_Cy,x-y\rangle.
\]
Indeed the two inequalities for $x-P_Cx$ against $P_Cy-P_Cx$
and for $y-P_Cy$ against $P_Cx-P_Cy$ add to this one.
Expanding the square consequently gives the firm estimate
\[
 \|P_Cx-P_Cy\|^2+
 \|(x-P_Cx)-(y-P_Cy)\|^2\le\|x-y\|^2.
\]
Compare two alternating orbits $(u_j),(v_j)$ and put
$a_j=u_j-v_j$. Applying this estimate at each step and
summing yields
\[
 \sum_{j=1}^N\|a_j-a_{j-1}\|^2
       \le\|a_0\|^2-\|a_N\|^2\le\|a_0\|^2.
\]
In particular, the differences of their even gaps tend to zero:
\[
 (u_{2n}-u_{2n-1})-(v_{2n}-v_{2n-1})
                  =a_{2n}-a_{2n-1}\longrightarrow0.
\]

Given $\varepsilon>0$, choose $v_0\in C_2$ with
$\operatorname{dist}(v_0,C_1)\le\delta+\varepsilon$.
Such a point exists by the definition of infimum. Its first
gap has norm $\|v_1-v_0\|\le\delta+\varepsilon$.
Every following gap has no larger norm: $v_{j-1}$ belongs
to the set onto which $v_j$ is projected when choosing
$v_{j+1}$, so
$\|v_{j+1}-v_j\|\le\|v_j-v_{j-1}\|$.
Thus every even gap of this comparison orbit is at most
$\delta+\varepsilon$. The preceding vanishing difference
implies
\[
 \delta\le\liminf_n\|u_{2n}-u_{2n-1}\|
 \le\limsup_n\|u_{2n}-u_{2n-1}\|\le\delta+\varepsilon.
\]
The first inequality follows from membership in $C_2-C_1$.
Let $\varepsilon\downarrow0$ to obtain convergence of the
gap norms to $\delta$.

The set $D$ is nonempty closed convex and its distance from
zero is $\delta$, so $\|p\|=\delta$. Its projection inequality
gives $\operatorname{Re}\langle p,w-p\rangle\ge0$ for
$w\in D$. Consequently
\[
 \|w-p\|^2\le\|w\|^2-\delta^2.
\]
Apply this to $w=u_{2n}-u_{2n-1}$; convergence of its norm
proves the claimed vector convergence.
\end{proof}

\begin{corollary}[Weak convergence to an attained nearest pair]
\label{fa2:p2:projection-nearest-pair}
In the preceding theorem, suppose $p\in C_2-C_1$. There
are $a\in C_1,b\in C_2$ with $b-a=p$ such that
$u_{2n-1}\rightharpoonup a$ and $u_{2n}\rightharpoonup b$.
The pair may depend on $u_0$. In particular the assumption
holds when one of the two sets is bounded.
\end{corollary}
\begin{proof}
The nearest-pair set
\[
 M=\{(a,b)\in C_1\times C_2:b-a=p\}\subset H\oplus H
\]
is nonempty closed convex. Each $(a,b)\in M$ satisfies
$P_{C_1}b=a$ and $P_{C_2}a=b$: these points attain the
global minimum distance $\delta$. Nonexpansiveness gives
\[
 \|u_{2n}-b\|\le\|u_{2n-1}-a\|
                  \le\|u_{2n-2}-b\|.
\]
Both even and odd distances to the corresponding component
are nonincreasing. Thus $x_n=(u_{2n-1},u_{2n})$ is bounded
and $\|x_n-z\|$ has a limit for every $z\in M$.
Every weak subsequential limit of $x_n$ lies in $C_1\times C_2$
because closed convex sets are weakly closed; the norm convergence
of its component difference to $p$ puts that limit in $M$.
Weak subsequential limits exist by reflexivity and
Eberlein--Smulian. There cannot be two different limits $z,w$:
the convergent scalar sequence
$\|x_n-z\|^2-\|x_n-w\|^2$ has respective subsequence limits
$-\|z-w\|^2$ and $\|z-w\|^2$, obtained by expanding the
squares and using weak convergence. Therefore $z=w$.
If weak convergence to this point failed, a subsequence
outside one fixed weak neighborhood would have a further
weakly convergent subsequence, giving the same point and
a contradiction. This proves both component limits.

Finally, if for example $C_1$ is bounded, choose pairs
$(a_n,b_n)\in C_1\times C_2$ with
$\|b_n-a_n\|\to\delta$. Both component sequences are bounded,
so a subsequence converges weakly in $H\oplus H$ to $(a,b)$
in the weakly closed product. Weak lower semicontinuity gives
$\|b-a\|\le\delta$, hence equality. The displacement $b-a$
belongs to $D$ and has its minimum norm, so by uniqueness
of $P_D0$ it equals $p$. The case of bounded $C_2$ is the same.
\end{proof}

% END INPUT


% BEGIN parts/part2-minimax.tex
\chapter{Minimax and noncoercive variational inequalities}
These are the abstract results of \cite[Problems 30--31]{B01}.
They separate compactness-based existence from coercive uniqueness.

\begin{theorem}[Convex--concave minimax]
\label{fa2:p2:minimax}
Let $K,L$ be nonempty compact convex subsets of Hausdorff
topological vector spaces. Suppose $F:K\times L\to\mathbb R$
is convex and lower semicontinuous in its first variable and
concave and upper semicontinuous in its second. Then
\[
 \min_{x\in K}\max_{y\in L}F(x,y)
            =\max_{y\in L}\min_{x\in K}F(x,y).
\]
All four extrema are attained, and a minimizing $x_0$ and
maximizing $y_0$ give a saddle pair
$F(x_0,y)\leq F(x_0,y_0)\leq F(x,y_0)$.
In particular this applies to closed bounded convex sets of
reflexive Banach spaces with norm-lower/upper-semicontinuous
convex/concave sections, using their weak topologies. It includes
the continuous bilinear von Neumann theorem.
\end{theorem}
\begin{proof}
Let $\alpha=\sup_y\inf_xF(x,y)$. It is finite: a fixed
$y_*$ gives a finite lower bound by lower semicontinuity and
compactness of $K$, and a fixed $x_*$ gives a finite upper
bound by upper semicontinuity and compactness of $L$.
For $a>\alpha$, the closed convex subsets
$K_y=\{x:F(x,y)\leq a\}$ have the finite intersection property.
Indeed fix $y_1,\ldots,y_m$ and form the upward set
\[
 V=\{v\in\mathbb R^m:\exists x\in K,
                              F(x,y_i)\leq v_i\ (1\leq i\leq m)\}.
\]
It is convex by convexity of the sections. It is closed: for
a converging net of its vectors, compactness supplies a subnet
of the witnesses, and lower semicontinuity preserves each
inequality. If $a(1,\ldots,1)\notin V$, strict separation of
this point from the closed convex set gives a nonzero vector
$t=(t_i)$, with $t_i\geq0$ because $V$ is upward, such that
\[
 \inf_{x\in K}\sum_i t_iF(x,y_i)>a\sum_i t_i.
\]
Normalize $\sum t_i=1$. Concavity at $y=\sum t_i y_i\in L$
gives $F(x,y)\geq\sum t_iF(x,y_i)$ for every $x$, so its
infimum is greater than $a>\alpha$, a contradiction.
Thus the finite intersection is nonempty. Compactness of $K$
gives $x_a$ satisfying $F(x_a,y)\leq a$ for every $y$.
It follows that $\inf_x\sup_yF(x,y)\leq\alpha$; the opposite
inequality holds for every function by comparing each pair.

The function $x\mapsto\max_yF(x,y)$ is lower semicontinuous,
as a supremum of lower semicontinuous functions, and is finite.
It therefore attains its minimum on $K$. Similarly the
upper-semicontinuous function $y\mapsto\min_xF(x,y)$ attains
its maximum. Comparing their common value at $(x_0,y_0)$
gives equality there and the saddle inequalities. The section
extrema themselves exist by compactness. Finally convex norm
lower semicontinuity on a Banach space is weak lower
semicontinuity by separation; apply this also to $-F$.
Reflexivity supplies the stated weak compactness.
\end{proof}

\begin{theorem}[A bounded-section replacement for coercivity]
\label{fa2:p2:noncoercive-vi}
Let $A$ be a bounded linear monotone operator on a real Hilbert
space, $K$ a nonempty closed convex set, and $f\in H$. If
there is $v_0\in K$ for which
\[
 \{u\in K:\langle Au-f,v_0-u\rangle\geq0\}
\]
is bounded, then there is $u\in K$ with
$\langle Au-f,v-u\rangle\geq0$ for every $v\in K$.
This condition holds when $K$ is bounded, or when $A$ is
coercive. Without strict monotonicity uniqueness is not asserted.
\end{theorem}
\begin{proof}
Apply Stampacchia to $A+\varepsilon I$ and
$f+\varepsilon v_0$. Its solution $u_\varepsilon$ obeys,
on testing $v_0$,
\[
 \langle Au_\varepsilon-f,v_0-u_\varepsilon\rangle
                 \geq\varepsilon\|u_\varepsilon-v_0\|^2\geq0.
\]
The hypothesis bounds these vectors. Take a weakly convergent
subsequence as $\varepsilon\downarrow0$, with limit $u\in K$.
Monotonicity rewrites their inequalities in the form
\[
 \langle Av-f,v-u_\varepsilon\rangle
       \geq-\varepsilon\langle u_\varepsilon-v_0,
                                             v-u_\varepsilon\rangle
       \quad(v\in K).
\]
The right side tends to zero. Pass to the weak limit; then
put $v=u+t(w-u)$ for any $w\in K$, divide by $t>0$ and
let $t\downarrow0$. Boundedness of $A$ gives the required
inequality at $u$. In the coercive case the section inequality
implies
$\alpha\|u\|^2\leq\|A\|\|v_0\|\|u\|+
\|f\|(\|u\|+\|v_0\|)$, a quadratic bound, proving the
last sufficient condition. For $A=0$ and bounded $K$, it says
exactly that a continuous linear functional attains its maximum
on a closed bounded convex Hilbert set.
\end{proof}

\begin{theorem}[A coercive form plus a convex potential]
\label{fa2:p2:potential-vi}
For bounded linear coercive $A$ on a real Hilbert space and a proper
lower semicontinuous convex $\phi$, there is a unique
$u\in D(\phi)$ with
\[
 \langle Au,v-u\rangle+\phi(v)-\phi(u)\geq0
                                      \quad(v\in D(\phi)).
\]
Taking $\phi=I_K-\langle f,\cdot\rangle$ recovers
Stampacchia with a forcing term.
\end{theorem}
\begin{proof}
The proximal map $P_\rho z$ uniquely minimizes
$\phi(v)+\|v-z\|^2/(2\rho)$ and is nonexpansive.
Existence follows from an affine lower bound for $\phi$,
coercivity of the quadratic, weak lower semicontinuity and
Hilbert reflexivity; strict convexity gives uniqueness. Its
optimality inequality, obtained on line segments, is
$(z-P_\rho z)/\rho\in\partial\phi(P_\rho z)$.
Adding the two subgradient inequalities gives
$\|P_\rho z-P_\rho w\|^2\leq
\langle z-w,P_\rho z-P_\rho w\rangle$, proving nonexpansion.
For $0<\rho<2\alpha/\|A\|^2$ the map
$u\mapsto P_\rho(u-\rho Au)$ is a strict contraction, since
the norm-square estimate for $I-\rho A$ is
$1-2\rho\alpha+\rho^2\|A\|^2<1$.
The zero Hilbert space is immediate and avoids division by
$\|A\|=0$. The fixed-point theorem gives $u$, and its
proximal condition is $-Au\in\partial\phi(u)$, exactly the
displayed inequality. Adding two such inequalities and using
coercivity proves uniqueness as well.
\end{proof}

\begin{theorem}[The weak Minty inequality on an arbitrary monotone domain]
\label{fa2:p2:minty-compact-domain}
Let $X$ be a real reflexive Banach space, $K\subset X$ nonempty closed bounded
convex, and $A:D\subset X\to X'$ monotone, with no assumption
that $D$ is all of $X$ or that $K\subset D$. There is $u\in K$
such that $\langle Av,u-v\rangle\leq0$ for every $v\in D\cap K$.
If $K\subset D$ and $A$ is hemicontinuous along its segments,
this is equivalent to
$\langle Au,v-u\rangle\geq0$ for all $v\in K$.
\end{theorem}
\begin{proof}
For each $v\in D\cap K$ its proposed inequality defines a
weakly closed subset of the weakly compact set $K$.
For finitely many $v_i$ set
$F(s,t)=\sum_{i,j}s_it_j\langle Av_j,v_i-v_j\rangle$
on the two finite probability simplexes. For $s_i\geq0$,
$\sum s_i=1$, monotonicity gives
\[
 F(s,s)=-\tfrac12\sum_{i,j}s_is_j
                  \langle Av_i-Av_j,v_i-v_j\rangle\leq0.
\]
Thus $\max_t\min_sF(s,t)\leq0$, by taking $s=t$.
The finite bilinear minimax theorem just proved gives some
$s$ with $F(s,t)\leq0$ for every $t$.
The vector $u=\sum s_iv_i\in K$ satisfies each of the
finitely many inequalities by testing the coordinate vertices.
Compactness and the finite intersection property give the
simultaneous solution; if $D\cap K$ is empty every point of
$K$ is a solution. For the final assertion use monotonicity
in one direction, and in the other substitute
$v=u+t(w-u)$, divide by $t$, and pass to $Au$ by
hemicontinuity. Only probability coefficients were used in
the quadratic identity; its sign is not claimed for arbitrary
signed coefficients.
\end{proof}

\begin{remark}[Two sign qualifications in the source problem]
\label{fa2:p2:minty-source-signs}
In \cite[Problem 31(A), printed p.~483]{B01}, the stated goal
prints $\langle Av,u-v\rangle\ge0$, whereas its subsequent
sets $K_v$ use $\langle Av,v-u\rangle\ge0$.
The latter, equivalently the inequality proved above, is the
correct version. The former fails for $A(v)=v$ and $K=[-1,1]$:
the tests $v=1,-1$ would require both $u\ge1$ and $u\le-1$.
Likewise the source's quadratic sign for all signed coefficients
is not needed and is false: with $v_1=0,v_2=1$, $A(v)=v$
and $(s_1,s_2)=(1,-1)$, its displayed $F(s,s)$ is $1$.
The probability-simplex identity in the proof has the valid sign.
These corrections preserve the actual monotone existence theorem,
rather than turn the two inconsistent printed formulas into claims.
\end{remark}

% END INPUT

\part{Unbounded operators, groups and semigroups}

% BEGIN parts/part3.tex
% Part III. Every unbounded operator carries its domain as part of its data.

% BEGIN parts/part3-domains.tex
\chapter{Domains, adjoints, and self-adjoint extensions}
An operator is the pair $(D(A),A)$, not just a differential expression.
Products have $D(AB)=\{x\in D(B):Bx\in D(A)\}$; sums have
$D(A+B)=D(A)\cap D(B)$. Equality includes equality of domains.
The resolvent convention remains $R(\lambda,A)=(\lambda I-A)^{-1}$.
See \cite{B01,B02,B03,B04,B08}.

\begin{theorem}[Graphs and Hilbert adjoints]\label{fa2:p3:closed-adjoint}
Let $A:D(A)\subset H\to K$ be densely defined. Its adjoint has domain
\[
 D(A^*)=\{y\in K:\text{there is }z\in H\text{ with }
              \faip{Ax}{y}=\faip{x}{z}\ (x\in D(A))\},
 \qquad A^*y=z.
\]
It is closed. The operator $A$ is closable iff $D(A^*)$ is dense;
then $\overline A=A^{**}$, whose graph is the closure of the graph
of $A$. For a closed densely defined $A$,
$(\faRan A)^\perp=\ker A^*$ and
$(\faRan A^*)^\perp=\ker A$.
\end{theorem}
\begin{proof}
Density makes $z$ unique. Put $J(x,y)=(-y,x)$ on $H\oplus K$.
The defining identity says $G(A^*)=(JG(A))^\perp$ after exchanging
the two Hilbert factors, so the adjoint graph is closed.
A linear subspace of $H\oplus K$ is a graph precisely when it contains
no nonzero vertical vector $(0,y)$. Therefore $A$ is closable iff
$\overline{G(A)}$ has no such vector. By the graph orthogonality
identity a vertical vector in that closure is exactly a vector
orthogonal to $D(A^*)$. This proves the density criterion.
Taking orthogonal complements twice and applying $J^{-1}$ gives
$G(A^{**})=\overline{G(A)}$.
Finally $y\perp\faRan A$ means the defining identity with $z=0$,
so exactly $A^*y=0$. Apply this to $A^*$ and use $A^{**}=A$ for
the other assertion.
\end{proof}

\begin{proposition}[A typical Schrödinger potential]\label{fa2:p3:schrodinger-potential}
If $q=q_2+q_\infty$ is real with $q_2\in L^2(\mathbb R)$ and
$q_\infty\in L^\infty(\mathbb R)$, then
$-d^2/dx^2+q$ on exactly $H^2(\mathbb R)$ is self-adjoint and
bounded below, and $C_c^\infty(\mathbb R)$ is a core.
\end{proposition}
\begin{proof}
For compactly supported smooth $u$, integration of $(|u|^2)'$
gives $\|u\|_\infty^2\leq2\|u\|_2\|u'\|_2$.
Also $\|u'\|_2^2\leq\|u\|_2\|u''\|_2$ by integration by parts.
Density extends both inequalities to their Sobolev domains. Hence
$\|q_2u\|_2\leq\sqrt2\|q_2\|_2
\|u\|_2^{3/4}\|u''\|_2^{1/4}$.
Young's inequality makes this at most
$\varepsilon\|u''\|_2+C_\varepsilon\|u\|_2$ for every
$\varepsilon>0$. The bounded part adds only
$\|q_\infty\|_\infty\|u\|_2$.
Thus the real multiplication potential is symmetric and has
arbitrarily small relative bound with respect to the free
self-adjoint operator on $H^2$. Kato--Rellich proves self-adjointness
on that domain; equivalence of graph norms transfers the standard
cutoff-and-mollification core $C_c^\infty$ from the free operator.
For the lower bound,
$|\int q_2|u|^2|\leq\|q_2\|_2\|u\|_4^2
\leq\|q_2\|_2\|u\|_\infty\|u\|_2$.
The first Sobolev inequality and Young again bound this by
$\varepsilon\|u'\|_2^2+C_\varepsilon\|u\|_2^2$.
Choosing $\varepsilon<1$ and including the bounded potential gives
a finite lower bound on the quadratic form, hence on the operator.
\end{proof}

\begin{example}[One expression, different adjoints]\label{fa2:p3:derivative-domains}
On $L^2(0,1)$ let $S=-i\,d/dx$ with
$D(S)=H_0^1(0,1)$, where $H^1$ functions use their absolutely
continuous representatives. Then $D(S^*)=H^1(0,1)$ and $S^*=-i\,d/dx$.
The self-adjoint extensions are precisely
$S_\theta=-i\,d/dx$ on
$\{u\in H^1(0,1):u(1)=e^{i\theta}u(0)\}$,
$0\leq\theta<2\pi$.
\end{example}
\begin{proof}
Integration by parts gives
$\faip{Su}{v}-\faip{u}{S^*v}=i[\bar u v]_0^1$.
Testing compactly supported smooth $u$ in the adjoint definition
shows the distributional derivative of $v$ belongs to $L^2$;
the one-dimensional Sobolev characterization then gives $v\in H^1$.
Conversely integration by parts for such $v$ has zero boundary term
on $H_0^1$, proving the adjoint domain.
The boundary form on trace pairs $(u(0),u(1))\in\mathbb C^2$ is
$i(\overline{u(1)}v(1)-\overline{u(0)}v(0))$.
Its maximal isotropic subspaces are exactly the lines
$\{(a,e^{i\theta}a):a\in\mathbb C\}$: isotropy of a nonzero vector
requires equal moduli of its coordinates, and no two-dimensional
isotropic subspace exists. Vanishing of the boundary form against
that line imposes the same line on the adjoint trace. Thus those
domains are self-adjoint and exhaust the extensions.
\end{proof}

\begin{theorem}[Banach unbounded adjoints and closed range]\label{fa2:p3:banach-closed-range}
Let $A:D(A)\subset X\to Y$ be densely defined and closed between
Banach spaces. Define $D(A')\subset Y'$ by boundedness of
$x\mapsto y'(Ax)$ in the $X$ norm, and $A'y'$ by its unique
continuous extension. Then $A'$ is closed, and
\[
 \faRan A\text{ closed}\iff\faRan A'\text{ norm closed}
 \iff\faRan A'=(\ker A)^\perp
 \iff\faRan A={}^{\perp}(\ker A').
\]
Closedness of $\faRan A$ is also equivalent to
$\operatorname{dist}(x,\ker A)\leq C\|Ax\|$ on $D(A)$.
Surjectivity of $A$ is equivalent to
$\|y'\|\leq C\|A'y'\|$ on $D(A')$;
surjectivity of $A'$ is equivalent to
$\|x\|\leq C\|Ax\|$ on $D(A)$.
\end{theorem}
\begin{proof}
Passing to limits in $y'(Ax)=(A'y')(x)$ proves closedness of $A'$.
The graph norm makes $D(A)$ Banach. If its range is closed, apply
open mapping to $D(A)/\ker A\to\faRan A$; this bounds the quotient
graph norm, hence the displayed distance. Conversely that estimate
allows representatives of a convergent range sequence to be corrected
by kernel elements into an $X$-Cauchy sequence; closedness of $A$
then gives a preimage of the limit.
If the range is closed, for $f\in(\ker A)^\perp$ define
$g(Ax)=f(x)$. The distance estimate bounds this functional on the
range by $C\|f\|$. Hahn--Banach extends it to $y'\in Y'$, and
$A'y'=f$. Hence $\faRan A'=(\ker A)^\perp$.
For the converse assume $\faRan A'$ norm closed. Open mapping on
the graph-norm domain modulo the kernel supplies preimages $y'$
with $\|y'\|\leq C\|f\|$ for $f=A'y'$.
Each bounded piece $\faRan A'\cap rB_{X'}$ is weak-star closed:
choose these bounded preimages, use Alaoglu to take a weak-star
subnet in $Y'$, and pass to the defining identity for each
$x\in D(A)$. Krein--Smulian
(\ref{fa2:p1:krein-smulian}) now makes the whole range weak-star
closed. Its preannihilator is $\ker A$, by separation of the closed
graph (or its annihilator identity). Thus it equals $(\ker A)^\perp$.
Using the bounded preimages and the Hahn--Banach distance formula,
\[
 \operatorname{dist}(x,\ker A)
 =\sup_{\substack{f\perp\ker A\\\|f\|\leq1}}|f(x)|
 \leq C\|Ax\|,
\]
which proves closed range. The identity involving $\ker A'$ follows
by Hahn--Banach from $\overline{\faRan A}={}^{\perp}\ker A'$.
If the latter equality holds without closure the range is closed,
so all equivalences follow. The surjectivity criteria combine these
equivalences with the lower-bound closed-range criterion for the
adjoint, or use the same bounded-preimage argument and uniform
boundedness on the graph domain. For example $A$ onto makes
$\{y':\|A'y'\|\leq1\}$ pointwise bounded on every $Ax\in Y$;
uniform boundedness gives its norm boundedness. The converse lower
bound makes $A'$ injective with closed range, hence $A$ onto.
For $A'$ onto, uniform boundedness applied to the evaluations at
$x\in D(A)$ proves the second lower bound; conversely the estimate
extends $Ax\mapsto f(x)$ by Hahn--Banach for every $f\in X'$.
\end{proof}

\begin{theorem}[Symmetry, deficiency spaces, and Cayley]\label{fa2:p3:cayley}
A densely defined $S$ is symmetric if $S\subset S^*$; it is
self-adjoint if equality holds, including domains. It is essentially
self-adjoint if its closure is self-adjoint. For closed symmetric $S$,
\[
 \|(S-z)x\|^2=\|(S-\operatorname{Re}z)x\|^2
                    +|\operatorname{Im}z|^2\|x\|^2.
\]
It is self-adjoint iff $\faRan(S+i)=\faRan(S-i)=H$, equivalently
$N_+=\ker(S^*-i)=0=N_-=\ker(S^*+i)$.
Its Cayley isometry $V=(S-i)(S+i)^{-1}$ maps the closed space
$\faRan(S+i)=N_+^\perp$ onto $\faRan(S-i)=N_-^\perp$.
Self-adjoint operators correspond bijectively to unitary operators
$U$ with $\ker(I-U)=0$, by
\[
 U=(S-i)(S+i)^{-1},\quad
 S=i(I+U)(I-U)^{-1},\quad D(S)=\faRan(I-U).
\]
\end{theorem}
\begin{proof}
Symmetry kills the cross term in the norm expansion, proving the
identity and the lower bound off the real axis. Closedness then
implies closedness of those ranges. Their orthogonal complements
are the displayed deficiency spaces by the adjoint definition.
For self-adjoint $S$, nonreal eigenvalues of $S^*=S$ are impossible,
so both ranges are dense and closed. Conversely if both deficiencies
vanish, for $x\in D(S^*)$ solve
$(S+i)y=(S^*+i)x$; the zero kernel of $S^*+i$ gives $x=y\in D(S)$.
The norm identity makes the Cayley map isometric. On its initial
space $(I-V)(S+i)x=2ix$, so its difference range is dense.
For unitary $U$ without fixed vectors, $\faRan(I-U)$ is dense
because $\ker(I-U^*)=0$. The inverse formula defines a symmetric
operator there: expanding the two inner products and using
$\faip{Ux}{Uy}=\faip{x}{y}$ verifies symmetry. Its graph is closed
because $(I-U)x\mapsto i(I+U)x$ has the bounded inverse coordinate
$x=((I-U)x-i[i(I+U)x])/2$. Its two nonreal ranges are all of $H$,
as $(S+i)(I-U)x=2ix$ and
$(S-i)(I-U)x=2iUx$. It is therefore self-adjoint. Substitution checks
that the two constructions are inverse.
\end{proof}

\begin{theorem}[All closed symmetric extensions]\label{fa2:p3:deficiency-extensions}
For a closed symmetric $S$ one has the graph-orthogonal decomposition
$D(S^*)=D(S)\dotplus N_+\dotplus N_-$.
Every closed symmetric extension is, uniquely, a restriction $S_W$
of $S^*$ with
\[
 D(S_W)=D(S)\dotplus\{u+Wu:u\in L_+\},\quad
 S_W(x+u+Wu)=Sx+iu-iWu,
\]
where $W:L_+\to N_-$ is an isometry on a closed subspace
$L_+\subset N_+$. It is self-adjoint exactly when
$L_+=N_+$ and $W$ is onto $N_-$. In particular self-adjoint
extensions exist iff the two deficiency indices are equal, including
infinite cardinalities. Deficiency dimensions are constant in each
open half-plane.
\end{theorem}
\begin{proof}
In the graph inner product
$\faip{x}{y}_{S^*}=\faip{x}{y}+\faip{S^*x}{S^*y}$,
the three displayed subspaces are mutually orthogonal. If a graph
vector is orthogonal to all three, orthogonality to $D(S)$ says
$S^*x\in D(S^*)$ and $(S^*)^2x=-x$. Thus
$(S^*+i)x\in N_+$ and $(S^*-i)x\in N_-$.
Orthogonality to those deficiency graphs then forces both vectors
zero, hence $x=0$. Since all three graph subspaces are closed,
this proves the decomposition.
Every symmetric extension lies between $S$ and $S^*$.
Modulo $D(S)$ its domain is a subspace $L\subset N_+\oplus N_-$
on which the boundary form vanishes. For $u+v\in L$, applying that
form to itself gives $\|u\|=\|v\|$; applying it to sums and complex
scalar multiples gives $\faip{u}{u'}=\faip{v}{v'}$.
Thus $v=Wu$ for an isometry on a subspace $L_+$.
Graph closedness is equivalent to closedness of $L_+$, since the
deficiency graph norm is a fixed multiple of $\|u\|$.
Conversely these identities make each displayed restriction symmetric
and closed. Its deficiency spaces have dimensions
$\dim(N_+\ominus L_+)$ and
$\dim(N_-\ominus WL_+)$: compute
$\faRan(S_W+i)=\faRan(S+i)\oplus L_+$ and similarly for the other
sign. The self-adjoint criterion above proves the characterization.
For constancy at nonreal $z$, if $|w-z|<|\operatorname{Im}z|$ and
$u\in\ker(S^*-w)\cap\ker(S^*-z)^\perp$, write
$u=(S-\bar z)x$ and pair $(S^*-w)u=0$ with $x$.
The lower bound gives $\|x\|\leq\|u\|/|\operatorname{Im}z|$,
while the pairing gives
$\|u\|^2\leq|w-z|\|u\|\|x\|$, so $u=0$.
The orthogonal projection between the two deficiency spaces is
therefore injective. An injective bounded map of Hilbert spaces gives
the corresponding cardinal dimension inequality: its adjoint has
dense range, and the images of a basis span a dense space, with each
vector having countable orthonormal support. Reversing $w,z$ in a sufficiently small disk
gives equality of their Hilbert dimensions. The dimensions are locally
constant and hence constant on each connected half-plane.
\end{proof}

\begin{theorem}[Friedrichs extension]\label{fa2:p3:friedrichs}
If a densely defined symmetric $S$ satisfies
$\faip{x}{Sx}\geq c\|x\|^2$, it has a self-adjoint extension with
the same lower bound, canonically associated with the closure of its
quadratic form. More generally a densely defined closed positive
sesquilinear form $q$ on $V\subset H$, with
$\|u\|_V^2=q(u,u)+\|u\|^2$, determines a unique nonnegative
self-adjoint $A$ by
\[
 D(A)=\{u\in V:\exists f\in H, q(v,u)=\faip{v}{f}\ (v\in V)\},
 \qquad Au=f.
\]
\end{theorem}
\begin{proof}
First shift $S$ so that $S\geq I$. On $D(S)$ put
$q(v,u)=\faip{v}{Su}$. Complete it in the norm $q(u,u)^{1/2}$.
This completion embeds injectively into $H$. Indeed if $u_n$ is
form-Cauchy and $u_n\to0$ in $H$, then for fixed $v\in D(S)$,
$q(v,u_n)=\faip{Sv}{u_n}\to0$; density of $D(S)$ in the form
completion makes its form limit zero. This proves closability and
the injective embedding of the completed form domain.
For a general closed nonnegative form, Lax--Milgram on $V$ solves
$q(v,u)+\faip{v}{u}=\faip{v}{f}$ uniquely for every $f\in H$.
Thus the operator in the statement has $\faRan(I+A)=H$.
Its inverse $B$ is a bounded positive self-adjoint contraction:
the form identity gives $\|Bf\|\leq\|f\|$, positivity, and symmetry
by testing the two solutions against one another. It is injective,
and its range is dense because $\ker B^*=0$.
Its inverse on $\faRan B$ is self-adjoint: the adjoint identity for
an injective bounded self-adjoint operator with dense range shows
$(B^{-1})^*=B^{-1}$ directly. Hence $A=B^{-1}-I$ is self-adjoint,
and its form identity gives nonnegativity. Uniqueness follows since
the form equation uniquely fixes $(I+A)^{-1}$.
For the completed form of $S$, each original $u\in D(S)$ satisfies
$q(v,u)=\faip{v}{Su}$ by form-density and continuity, so the
associated operator extends $S$. Undo the initial scalar shift to
obtain the claimed lower bound.
\end{proof}

\begin{theorem}[Relative symmetric perturbations]\label{fa2:p3:kato-rellich}
If $A=A^*$ and $B$ is symmetric on a domain containing $D(A)$, and
$\|Bx\|\leq a\|Ax\|+b\|x\|$ on $D(A)$ with $a<1$, then $A+B$
on exactly $D(A)$ is self-adjoint. If $A$ is essentially self-adjoint
on a core $D_0$ and the same estimate holds there, the closure of
$A+B$ on $D_0$ is self-adjoint on $D(\overline A)$.
\end{theorem}
\begin{proof}
The spectral theorem gives
$\|(A\pm i\lambda)^{-1}\|\leq1/\lambda$ and
$\|A(A\pm i\lambda)^{-1}\|\leq1$.
Thus $\|B(A\pm i\lambda)^{-1}\|\leq a+b/\lambda<1$ for large
$\lambda$. Factoring $A+B\pm i\lambda$ and using the geometric
series makes both ranges all of $H$. The relative bound also gives
equivalence of the graph norms of $A$ and $A+B$, so the latter is
closed. It is symmetric and has both nonreal ranges onto, hence
self-adjoint by the criterion above. On a core, the bound extends
$B$ continuously for the graph norm to $D(\overline A)$, preserving
symmetry. Equivalent graph norms make $D_0$ also a core for the sum,
so the preceding argument proves the second assertion.
\end{proof}

\begin{theorem}[Unbounded measurable spectral calculus]\label{fa2:p3:unbounded-calculus}
For a PVM $E$ and finite-valued measurable $f$ define
\[
 D(f(E))=\{x:\int|f|^2\,d\mu_x<\infty\},\quad
 f(E)x=\lim_n(f1_{\{|f|\leq n\}})(E)x.
\]
It is a closed densely defined normal operator, its adjoint is
$\bar f(E)$, and
\[
 D(f(E)g(E))=D(g(E))\cap D((fg)(E)),\qquad
 \overline{f(E)g(E)}=(fg)(E).
\]
Likewise the closure of $f(E)+g(E)$ on the intersection domain is
$(f+g)(E)$. A real-valued $f$ gives a self-adjoint operator.
Even when $f$ is bounded the domain of the product $f(E)g(E)$ need
not be the full domain of $(fg)(E)$.
\end{theorem}
\begin{proof}
The truncated spectral ranges increase densely to $H$, giving density
and a well-defined Cauchy limit by the norm-square identity. Split
the space into disjoint spectral sets on which $f$ is bounded.
On their orthogonal sum the operator has domain consisting exactly
of component vectors whose image squared norms are summable.
Such a direct sum is closed: graph convergence gives componentwise
limits, whose squared image norms sum to the norm of the limiting
image. Its adjoint is the direct sum of the component adjoints, with
the same square-summability condition. This proves the stated
adjoint and normality.
If $x\in D(g(E))$, the spectral measure of $g(E)x$ is
$|g|^2\mu_x$, first for bounded truncations and then by norm
convergence. Hence it lies in $D(f(E))$ exactly when
$\int|fg|^2d\mu_x<\infty$, proving the product domain formula and
its agreement with $(fg)(E)$ there. Truncating simultaneously where
$|f|,|g|\leq n$ gives a graph core for $(fg)(E)$ and for
$(f+g)(E)$, proving both closure statements. Taking $f=0$ and an
unbounded $g$ exhibits the last warning: the product domain is
$D(g(E))$, while $(fg)(E)=0$ is defined on all of $H$.
\end{proof}

\begin{proposition}[Commutation in the unbounded calculus]\label{fa2:p3:unbounded-commutation}
Let $N$ be closed densely defined normal and $B$ bounded with
$BN\subset NB$ and $BN^*\subset N^*B$. Then $B$ commutes with
every spectral projection of $N$. For every finite-valued Borel $f$,
$B D(f(N))\subset D(f(N))$ and $Bf(N)\subset f(N)B$.
\end{proposition}
\begin{proof}
The assumptions imply that $B$ preserves $D(N^*N)$ and commutes
with $N^*N$ there. Solving its shifted equation gives commutation
with $(I+N^*N)^{-1}$. Bounded Fuglede and spectral calculus give
commutation with all spectral pieces of this positive inverse.
On each piece bounded away from zero, $N$ is bounded normal and
commutes with the restriction of $B$; its PVM therefore commutes
too. The construction in the next theorem assembles these bounded
PVMs into the full one, proving the first assertion. For every
Borel set $S$, $\mu_{Bx}(S)=\|E(S)Bx\|^2
\leq\|B\|^2\mu_x(S)$. Integrating $|f|^2$ proves domain preservation.
The bounded truncations of $f$ commute with $B$; taking their norm
limits on $D(f(N))$ gives the stated operator inclusion, with its
exact domain meaning.
\end{proof}

\begin{theorem}[Unbounded normal spectral theorem]\label{fa2:p3:normal-spectral}
Every closed densely defined normal $N$, meaning $N^*N=NN^*$
including product domains, has a unique Borel PVM on $\mathbb C$
with $N=\int z\,dE(z)$ on exactly its square-integrability domain.
It is supported on $\sigma(N)$ and every open set meeting the spectrum
has nonzero spectral projection. No assumption that the resolvent
is nonempty is needed. In particular a self-adjoint operator has a
PVM on $\mathbb R$; a general normal operator is a direct sum of
bounded normal pieces and hence a multiplication operator, including
on nonseparable Hilbert spaces.
\end{theorem}
\begin{proof}
For any closed densely defined $A$, orthogonal decomposition of
$H\oplus H$ into the rotated graph of $A$ and the graph of $A^*$
shows $I+A^*A$ is onto: writing $(0,h)=(-Ax,x)+(y,A^*y)$ gives
$y=Ax$ and $h=x+A^*Ax$. Also
$\|(I+A^*A)x\|\geq\|x\|$, and its inverse $B$ is positive,
self-adjoint, injective and contractive. The graph of $A$ restricted
to $D(A^*A)$ is dense in $G(A)$: a graph vector orthogonal to all
such vectors would satisfy $\faip{x}{(I+A^*A)y}=0$ for every such
$y$, hence $x=0$. This proves $A^*A$ self-adjoint and the graph-core
assertion without using the desired unbounded spectral theorem.
For normal $N$, this core and equality $N^*N=NN^*$ give
$D(N)=D(N^*)$ and $\|Nx\|=\|N^*x\|$: prove it on the common
product domain, then close the two graphs.
Put $B=(I+N^*N)^{-1}$ and $C=NB$, both bounded. The commutation
has a domain-sensitive justification. For any closed densely defined
$A$, put $B_A=(I+A^*A)^{-1}$ and $C_A=(I+AA^*)^{-1}$.
For $x\in D(A)$, $A^*AB_Ax=x-B_Ax\in D(A)$, so
$AB_Ax\in D(AA^*)$ and $(I+AA^*)AB_Ax=Ax$.
Thus $AB_Ax=C_AAx$. For normal $N$, $C_N=B_N=B$, hence
$BN\subset NB$. Apply this on $Bx\in D(N^*N)$ to obtain
$BC=CB$; the same argument applies with $N^*$.
Let $P$ be the bounded positive spectral measure of $B$.
On $H_\delta=P([\delta,1])H$, $B$ is boundedly invertible and
$H_\delta\subset\faRan B\subset D(N)$.
Since $C$ commutes with $B$ and its spectral projections,
$N|_{H_\delta}=CB^{-1}|_{H_\delta}$ is bounded; the analogous
identity for $N^*$ makes this a reducing normal restriction, with
norm at most $((1-\delta)/\delta)^{1/2}$.
Partition $(0,1]$ into $(1/(n+1),1/n]$. Their spectral spaces are
orthogonal and sum to $H$, since $B$ is injective. Apply the bounded
normal theorem on each and take the direct sum of their PVMs.
The direct-sum domain equals $D(N)$: one inclusion follows from
closedness of $N$, and the other first on the graph core $D(N^*N)$
using $NP_n=P_nN$, then on its graph closure.
Thus $N=\int z,dE$. Its spectrum is the closed union of the bounded
piece spectra: away from that union the piece resolvents are bounded
by the reciprocal distance and their direct sum is an inverse;
points in the union are spectral by approximate vectors. This proves
the support assertions. Uniqueness follows because the bounded
operator $B$ and its spectral partition are determined by $N$, and
uniqueness holds on each bounded piece. The multiplication model
follows by taking the cyclic models of these pieces.
\end{proof}

% END INPUT


% BEGIN parts/part3-lp-multiplication.tex
\chapter{Multiplication operators on general $L^p$ spaces}
\label{fa2:p3:lp-multiplication-chapter}

The Hilbert spectral theorem is not a substitute for this model on
$L^p$ when $p\ne2$. Throughout, $1<p<\infty$ and $q=p/(p-1)$.
For Banach duality we use the scalar-linear, bilinear pairing
$[f,g]=\int fg\,d\mu$. Thus a Banach transpose does not conjugate
the multiplier. The conclusions below include all ten questions of
\cite[Problem33]{B01}, with the measure and zero-space qualifications
made explicit. The domain, transpose, bounds, and compactness results
are proved over either scalar field, not inferred for real $L^p$ from
a theorem stated only on a complex space.

\begin{theorem}[Domains, transpose, kernels, and range]
\label{fa2:p3:lp-multiplication}
Let $(\Omega,\Sigma,\mu)$ be any positive measure space, and let
$a:\Omega\to\mathbb K$, where $\mathbb K=\mathbb R$ or $\mathbb C$,
be measurable and finite almost everywhere.
On $X=L^p(\mu;\mathbb K)$ put
\[
 D(M_a)=\{f\in X:af\in X\},\qquad M_af=af,
 \qquad Z=\{a=0\}.
\]
Then $M_a$ is densely defined and closed. Under the isometric
identification $X'=L^q(\mu)$ of
Theorem~\ref{fa2:p4:lp-duality}, its Banach transpose has exactly
\[
 D(M_a')=\{g\in L^q:ag\in L^q\},\qquad M_a'g=ag.
\]
Write $L^r(E)=\{f\in L^r:f=0\text{ almost everywhere on }\Omega\setminus E\}$.
Then
\[
 \begin{aligned}
 \ker M_a&=L^p(Z),& (\ker M_a)^\perp&=L^q(\Omega\setminus Z),\\
 \ker M_a'&=L^q(Z),& (\ker M_a')_\perp&=L^p(\Omega\setminus Z),\\
 \overline{\operatorname{Ran}M_a}&=L^p(\Omega\setminus Z).
 \end{aligned}
\]
The unclosed range is exactly
\[
 \operatorname{Ran}M_a
 =\left\{h\in L^p:h1_Z=0,\quad
             1_{\Omega\setminus Z}\frac{h}{a}\in L^p\right\}.
\]
These statements require no semifiniteness or $\sigma$-finiteness.
On complex $L^2$, the Hilbert adjoint instead is $M_{\overline a}$
with its maximal domain; in particular real $a$ gives a self-adjoint
operator.
\end{theorem}
\begin{proof}
Changing $a$ on a null set does not affect the operator, so take it
finite everywhere. For $f\in L^p$, set $f_n=f1_{\{|a|\le n\}}$.
Then $f_n\in D(M_a)$ and $f_n\to f$ by dominated convergence,
proving density.

Suppose $f_n\to f$ and $af_n\to h$ in $L^p$. Choose a subsequence
along which the sums of both $p$th-power errors are finite. Integrating
the two nonnegative sums shows that each sum is finite almost
everywhere; thus $f_n\to f$ and $af_n\to h$ pointwise along that
subsequence outside a null set. Finiteness of $a$ gives $h=af$ there.
Consequently $f\in D(M_a)$ and $M_af=h$, which is closedness.

If $g,ag\in L^q$, H\"older gives
$[M_af,g]=[f,ag]$ on $D(M_a)$, so $g\in D(M_a')$ and
$M_a'g=ag$. Conversely, for $g\in D(M_a')$ there is $h\in L^q$
with $[M_af,g]=[f,h]$ for every $f\in D(M_a)$. On
$E_n=\{|a|\le n\}$, every $L^p$ function supported in $E_n$
belongs to that domain, and $1_{E_n}ag\in L^q$. Uniqueness in
$L^p$ duality therefore gives $1_{E_n}h=1_{E_n}ag$. Since
$E_n\uparrow\Omega$, this proves $h=ag\in L^q$ and the exact
transpose domain. With the Hilbert pairing
$\langle g,f\rangle=\int\overline g f$, the same equality reads
$\langle g,M_af\rangle=\langle M_{\overline a}g,f\rangle$;
the preceding domain argument proves the asserted Hilbert adjoint.

The kernel formulas are the pointwise equation $af=0$. Their
annihilators follow either from duality on the range of the bounded
projection $f\mapsto1_Zf$, or directly by testing the norming
function $\overline g|g|^{q-2}1_Z$ against $g\in L^q$ (with
value zero where $g=0$), and by the analogous test with $p,q$
interchanged. A functional annihilates $\operatorname{Ran}M_a$
exactly when it belongs to $\ker M_a'$. Hahn--Banach separation
then identifies its norm closure with the stated preannihilator.
Finally, solving $af=h$ off $Z$ gives the exact range formula:
necessity follows from $f=h/a$ there; for sufficiency choose
$f=1_{\Omega\setminus Z}h/a$, so $af=h\in L^p$ and $f$ lies
in the required domain.
\end{proof}

\begin{theorem}[The ordinary essential-bound tests]
\label{fa2:p3:lp-multiplication-bounds}
Assume now that $\mu$ is semifinite: every measurable set of positive
measure contains a measurable subset of finite positive measure.
Then
\[
 D(M_a)=L^p\quad\Longleftrightarrow\quad a\in L^\infty.
\]
When this holds, $\|M_a\|=\|a\|_\infty$. Moreover,
$\operatorname{Ran}M_a$ is closed exactly when, for some $c>0$,
$|a|\ge c$ almost everywhere on $\Omega\setminus Z$. The operator
is surjective exactly when
\[
                    |a|\ge c\quad\hbox{almost everywhere on }\Omega
                    \quad\hbox{for some }c>0.
\]
In that case its inverse is the bounded, everywhere-defined
multiplier $M_{1/a}$ into $D(M_a)$, and $M_a$ is injective.
\end{theorem}
\begin{proof}
An essentially bounded multiplier is bounded on all of $L^p$, with
norm at most its essential supremum. Conversely, if the domain is
all of $L^p$, the closed graph theorem gives
$\|af\|_p\le C\|f\|_p$. Were $\{|a|>C+\varepsilon\}$ of
positive measure, semifiniteness would provide a subset $F$ of
finite positive measure. Its indicator belongs to the full domain
and contradicts that estimate. Thus $|a|\le C$ almost everywhere.
The same indicator test, now for any number below
$\|a\|_\infty$, gives equality of the operator norm. If $L^p=0$,
semifiniteness implies every measurable set is null; both norms
are zero and the assertions retain their usual vacuous meaning.

Let $Y=L^p(\Omega\setminus Z)$. If $|a|\ge c$ there, each
$h\in Y$ has $h/a\in L^p$, so the exact range formula gives
$\operatorname{Ran}M_a=Y$, which is closed. Conversely, if the
range is closed it equals $Y$ by the first theorem. The inverse
of $M_a:D(M_a)\cap Y\to Y$, regarded as a map $B:Y\to Y$,
has closed graph: if $h_n\to h$ and $Bh_n\to f$, closedness
of $M_a$ gives $af=h$, and its unique solution in $Y$ is $Bh$.
The closed graph theorem supplies a constant $C>0$ such that
\[
              \|f\|_p\le C\|af\|_p
              \qquad(f\in D(M_a)\cap Y).
\]
If $\{0<|a|<1/(2C)\}$ had positive measure, a finite positive
subset $F$ would have $1_F\in D(M_a)\cap Y$ and violate this
estimate. Thus $|a|\ge1/(2C)$ off $Z$, almost everywhere.
This argument also covers $Y=0$, where the conclusion is vacuous.

Surjectivity forces $Y=L^p$ and the range closed. Semifiniteness
makes $L^p(Z)=0$ equivalent to $\mu(Z)=0$, by the same finite-set
indicator test. Combining this with the closed-range criterion
proves the necessary uniform lower bound. Conversely such a bound
makes $1/a$ bounded, and $f=h/a$ solves $af=h$ for every $h\in L^p$.
It belongs to the maximal domain because $af=h\in L^p$.
The zero level is null, so the solution is unique.
\end{proof}

\begin{theorem}[Spectrum, eigenvalues, and continuous multipliers]
\label{fa2:p3:lp-multiplication-spectrum}
Let $\mu$ be semifinite and use complex $L^p$ for the spectrum.
For real $a$ on real $L^p$, ``spectrum'' means the spectrum of its
natural complexification on complex $L^p$: its domain is
$D(M_a;\mathbb R)+iD(M_a;\mathbb R)$, since multiplication by
real $a$ preserves real and imaginary parts and their integrability.
The bounded or unbounded maximal multiplier has
\[
 \sigma(M_a)=\operatorname{essran}_\mu a
 =\{\lambda\in\mathbb C:
      \mu(\{|a-\lambda|<\varepsilon\})>0
                         \text{ for every }\varepsilon>0\}.
\]
For $\lambda$ outside that closed set,
\[
 R(\lambda,M_a)f=\frac{f}{\lambda-a},\qquad
 \|R(\lambda,M_a)\|
       =\operatorname*{ess\,sup}\frac1{|\lambda-a|}.
\]
The eigenspace at $\lambda$ is $L^p(\{a=\lambda\})$;
thus $\lambda$ is an eigenvalue exactly when that level set has
positive measure. If $a$ is real, the spectrum is real. If also
$a\in L^\infty$ and $L^p\ne0$, its essential infimum and
essential supremum both belong to the spectrum, the spectrum lies
between them, and $\sigma(M_a)=\{0\}$ exactly when $a=0$ almost
everywhere.

For a nonempty open $\Omega\subset\mathbb R^n$ with Lebesgue
measure and continuous finite $a$, one has the exact formula
\[
                    \sigma(M_a)=\overline{a(\Omega)}.
\]
In particular the closure is necessary even for bounded continuous
multipliers.
\end{theorem}
\begin{proof}
If $\lambda$ is outside the displayed essential range, there is
$c>0$ with $|\lambda-a|\ge c$ almost everywhere. Multiplication
by $b=(\lambda-a)^{-1}$ is bounded. For $h\in L^p$,
$a(bh)=\lambda bh-h\in L^p$, so $bh\in D(M_a)$ and
$(\lambda I-M_a)bh=h$. Conversely
$b(\lambda I-M_a)f=f$ for every $f\in D(M_a)$. This proves
the resolvent assertion, including its exact norm by the preceding
bounded-multiplier norm formula.

If $\lambda$ is in that essential range, for each $\varepsilon>0$
choose a finite positive-measure subset
$F\subset\{|a-\lambda|<\varepsilon\}$. The normalized indicator
$f=1_F/\mu(F)^{1/p}$ has norm one, belongs to $D(M_a)$ because
$|a|\le|\lambda|+\varepsilon$ on $F$, and satisfies
$\|(\lambda I-M_a)f\|_p\le\varepsilon$.
A bounded inverse would give
$1\le\|R(\lambda,M_a)\|\varepsilon$ for all $\varepsilon$,
which is impossible. This proves the spectrum formula for the
unbounded operator as well as the bounded one. The eigenvalue
statement is the kernel formula for $M_{a-\lambda}$ together with
semifiniteness. For real $a$, a nonreal $\lambda$ is separated
from its values by $|\operatorname{Im}\lambda|>0$.

For bounded real $a$ on a nonzero $L^p$, its finite essential
endpoints $\alpha,\beta$ satisfy
$\mu(\{|a-\alpha|<\varepsilon\})>0$ and
$\mu(\{|a-\beta|<\varepsilon\})>0$ for every $\varepsilon>0$,
by the definition of essential infimum and supremum. Values
outside $[\alpha,\beta]$ have positive distance from $a$ almost
everywhere. Hence the endpoints and containment follow. If the
spectrum is $\{0\}$, both endpoints are zero, so $a=0$ almost
everywhere; the converse is the spectrum of zero on a nonzero
Banach space.

In the continuous example, $\lambda\notin\overline{a(\Omega)}$
is separated from all multiplier values. Conversely, if
$\lambda\in\overline{a(\Omega)}$, choose $x\in\Omega$ with
$|a(x)-\lambda|<\varepsilon/2$. Continuity supplies a nonempty
open neighborhood inside $\Omega$ on which $|a-\lambda|<\varepsilon$;
it has positive Lebesgue measure. This proves the last formula.
For example $a(x)=(1+e^{-x})^{-1}$ on $\mathbb R$ has range
$(0,1)$ but spectrum $[0,1]$, so writing the unclosed range
would be false.
\end{proof}

\begin{proposition}[The nonatomic compactness test and the measure boundary]
\label{fa2:p3:lp-multiplication-compact}
If $\mu$ is semifinite and nonatomic, multiplication is compact
as an everywhere-defined bounded operator on $L^p$ exactly when
$a=0$ almost everywhere. More strongly, for any nonzero multiplier
the image of the $L^p$ unit ball intersected with $D(M_a)$ is not
relatively compact. Atomlessness is necessary: on counting measure
on $\mathbb N$, the nonzero multiplier $a_n=1/n$ is compact.

Without semifiniteness, the first theorem remains valid, but the
ordinary essential-bound, essential-range and surjectivity tests
need not be valid. If $L^p=0$, the maximal multiplier is the
unique zero-space operator and its spectrum is empty.
\end{proposition}
\begin{proof}
Suppose $a\ne0$ on a set of positive measure. Some set
$\{\delta\le|a|\le N\}$ has positive measure, since their
countable union, for positive rational $\delta$ and integer $N$,
contains $\{a\ne0\}$. Semifiniteness gives a subset $F$ of finite
positive measure there. Nonatomicity successively splits the
remaining positive-measure subset of $F$ into two positive-measure
pieces. This produces pairwise disjoint $F_j\subset F$ with
$0<\mu(F_j)<\infty$. The normalized indicators
$f_j=1_{F_j}/\mu(F_j)^{1/p}$ belong to $D(M_a)$, have norm one,
and have disjoint supports. Thus, for $i\ne j$,
\[
 \|M_af_i-M_af_j\|_p^p
 =\|af_i\|_p^p+\|af_j\|_p^p\ge2\delta^p.
\]
No subsequence of their images is Cauchy, proving noncompactness.
The zero operator is compact. On $\ell^p$, truncating the sequence
$1/n$ after its first $m$ terms gives finite-rank multipliers
whose error norm is $1/(m+1)$; hence that nonzero multiplier is
compact and atomlessness cannot be dropped.

For the arbitrary-measure warning take a one-point measurable
space with mass $\infty$. Every finite-$p$ integrable function
vanishes there, so $L^p=0$. An arbitrarily large multiplier is
still the zero operator on this space. With $a=0$ it is surjective
although no positive ordinary almost-everywhere lower bound is
possible; with $a=1$ its ordinary essential range is $\{1\}$
although its operator spectrum is empty. Indeed on the zero space
$\lambda I-M_a$ is the unique bijective linear map for every
$\lambda$, with the unique bounded inverse. Such sets of infinite
measure containing no finite positive-measure subset are invisible
to finite-$p$ functions; semifiniteness is precisely what removes
this obstruction to the indicator tests above.
For an unbounded example use all subsets of $\mathbb N$ with
measure $\infty$ for every nonempty set. Again $L^p=0$,
but $a(n)=n$ is not ordinarily essentially bounded although its
maximal multiplication domain is all of $L^p$. Thus even the
full-domain test needs the stated measure qualification.
\end{proof}

% END INPUT


% BEGIN parts/part3-adjoint-sums.tex
\chapter{Perturbing Banach adjoint domains}
This is the general Banach-space assertion in
\cite[Problem 7]{B01}. The self-adjoint Hilbert perturbation
theorem alone does not prove it.

\begin{lemma}[One-sided gaps and annihilators]
\label{fa2:p3:subspace-gap}
For closed linear subspaces of a Banach space put
$\delta(M,N)=\sup_{x\in M,\|x\|\leq1}\operatorname{dist}(x,N)$.
Then $0\leq\delta(M,N)\leq1$,
$\delta(M,N)=\delta(N^\perp,M^\perp)$, and
\[
 \delta(M,L)\leq\delta(M,N)+\delta(N,L)
                              +\delta(M,N)\delta(N,L).
\]
If $N\subsetneqq M$, their one-sided gap $\delta(M,N)$ is one.
\end{lemma}
\begin{proof}
The upper bound follows by choosing the zero point of $N$.
Hahn--Banach gives
$\operatorname{dist}(x,N)=\sup_{f\in N^\perp,\|f\|\leq1}|f(x)|$.
For $f\in X'$, restriction and the norm-preserving extension
theorem give
$\operatorname{dist}(f,M^\perp)=\|f|_M\|$:
extend the restriction with that norm and subtract it from $f$
for the upper inequality; the reverse is immediate.
Taking the two suprema in either order proves the gap identity,
without identifying $X$ with its bidual.
For $\|x\|\leq1$ in $M$, choose $y\in N$ with
$\|x-y\|\leq\delta(M,N)+\varepsilon$. Then
$\|y\|\leq1+\delta(M,N)+\varepsilon$ and
$\operatorname{dist}(x,L)\leq\|x-y\|+
\delta(N,L)\|y\|$. Let $\varepsilon\downarrow0$ and take the
supremum. Finally Riesz's lemma on the proper closed subspace
$N$ of $M$ gives unit vectors with distance as close to one
as desired, proving the last assertion.
\end{proof}

\begin{theorem}[A linear maximal-monotonicity criterion]
\label{fa2:p3:banach-linear-maximal-monotone}
Let $X$ be real reflexive and suppose both $X$ and $X'$ are
strictly convex. Let $A:D(A)\subset X\to X'$ be closed,
densely defined, linear and monotone. Identify $X''$ with $X$.
Then $A$ is maximal monotone if and only if its Banach adjoint
is monotone on $D(A')$. In that case $A'$ is maximal monotone,
\[
 \ker A=\ker A',\qquad
 \overline{\operatorname{Ran}A}
                         =\overline{\operatorname{Ran}A'}.
\]
The range equality is an equality of norm closures, as in
\cite[Problem 16, B3]{B01}, not an unsupported equality of
unclosed ranges.
\end{theorem}
\begin{proof}
Here maximality means that any pair $(x,f)$ monotonically
related to the whole graph is already in the graph. Suppose
first that $A$ is maximal. For $v\in D(A')$ with
$\langle A'v,v\rangle<0$, the pair $(v,-A'v)$ is monotonically
related to the graph, since for $u\in D(A)$
\[
 \langle Au+A'v,u-v\rangle
       =\langle Au,u\rangle-\langle A'v,v\rangle>0.
\]
It would belong to that graph. Then $Av=-A'v$, whereas pairing
at the common domain point also gives
$\langle Av,v\rangle=\langle A'v,v\rangle$; these force zero,
a contradiction. Thus $A'$ is monotone.

Conversely suppose $A'$ monotone and fix a monotonically related
pair $(x,f)$. Let $J:X\to X'$ be the normalized duality map,
and $J^{-1}:X'\to X$ its inverse. Reflexivity and the two strict
convexities give unique norming vectors and functionals, so these
maps exist, and the derivatives of the squared norms are $J$
and $J^{-1}$. This is also proved in the norm-duality chapter:
the directional derivative is the maximum over norming
functionals, whose uniqueness makes it the displayed one.
Minimize on the graph domain
\[
 F(u)=\tfrac12\|Au-f\|^2+\tfrac12\|u-x\|^2
                            +\langle Au-f,u-x\rangle.
\]
The graph domain is reflexive, being a closed linear subspace of
$X\oplus X'$. The quadratic $u\mapsto\langle Au,u\rangle$
is convex by monotonicity. Thus $F$ is continuous convex in its
graph norm and weakly lower semicontinuous. It is coercive:
after expanding the last pairing, its only nonquadratic terms
are bounded linear functions of $u$ and $Au$, whereas
$\|Au-f\|^2+\|u-x\|^2$ dominates a positive multiple of
$\|u\|^2+\|Au\|^2$ minus a constant. The nonnegative
$\langle Au,u\rangle$ cannot spoil that bound. Reflexivity
therefore gives a minimizer $u$.

Write $v=u-x$, $r=Au-f$, $y=J^{-1}r$. Differentiating in
every direction $w\in D(A)$ gives
\[
 \langle Jv+r,w\rangle+\langle Aw,y+v\rangle=0.
\]
This says exactly that $y+v\in D(A')$ and
$A'(y+v)=-r-Jv$. Adjoint monotonicity yields
\[
 0\geq\langle r+Jv,y+v\rangle
 =\|r\|^2+\|v\|^2+\langle r,v\rangle+\langle Jv,y\rangle
 \geq(\|r\|-\|v\|)^2\geq0.
\]
Both Holder inequalities in this display are therefore
equalities. If either norm is zero both are zero. Otherwise
$\langle r,v\rangle=-\|r\|\|v\|$ and equality of the norms
force $r=-Jv$, by uniqueness of the norming functional.
But the original related-pair condition at this $u$ gives
$\langle r,v\rangle\geq0$. Hence $v=0$, $r=0$, so $x=u$
and $Ax=f$, proving maximality.

 The adjoint is closed and norm densely defined in the reflexive
 case. Indeed an element $y'\in X'$ annihilating $D(A')$ corresponds to
$(0,y')$ in the bidual graph of $A$, hence $(0,y')\in G(A)$,
so $y'=0$. This is the usual closability/adjoint-density graph
argument, with reflexivity on the target as well as the source.
It also gives $A''=A$. The criterion just proved, applied to
$A'$ and its monotone adjoint $A$, shows $A'$ maximal.

If $Au=0$, nonnegativity of
$\langle A(u+tw),u+tw\rangle$ for both signs of $t$ forces
$\langle Aw,u\rangle=0$ for every $w\in D(A)$. Hence
$u\in D(A')$ and $A'u=0$. Applying this to $A'$ gives the
reverse inclusion. Finally the closed-range annihilator
identities, with closures, identify both range closures as the
annihilator of this common kernel. This proves the last formula.
\end{proof}

\begin{theorem}[Adjoint of a relatively bounded sum]
\label{fa2:p3:banach-adjoint-sum}
Let $A:D(A)\subset X\to Y$ and $B:D(B)\subset X\to Y$ be
closed densely defined operators between Banach spaces, with
$D(A)\subset D(B)$ and $D(A')\subset D(B')$. Suppose
\begin{align*}
 \|Bx\|&\leq k\|Ax\|+C\|x\| &&(x\in D(A)),\quad 0\leq k<1,\ C\geq0,\\
 \|B'y'\|&\leq k'\|A'y'\|+C'\|y'\|
                &&(y'\in D(A')),\quad 0\leq k'<1,\ C'\geq0.
\end{align*}
Then $A+B$ is closed on $D(A)$ and
\[
 (A+B)'=A'+B'\quad\hbox{on exactly }D(A').
\]
Equality includes the adjoint domain. Neither reflexivity nor
symmetry nor a self-adjoint realization is assumed.
\end{theorem}
\begin{proof}
The first bound gives equivalence of the graph norms of $A$
and $A+tB$ for $0\leq t\leq1$:
\[
 (1-tk)\|Ax\|\leq\|(A+tB)x\|+tC\|x\|.
\]
The reverse bound is the triangle inequality. Completeness of
the graph domain therefore makes each $A+tB$ closed. The
dual bound similarly makes $A'+tB'$ closed on $D(A')$.
Always $A'+tB'\subset(A+tB)'$ by the defining pairing.

Here is a local equality argument. On $X\oplus Y$ use the
sum norm, so its dual has the maximum norm. The annihilator
of $G(A)$ is
$\{(-A'y',y'):y'\in D(A')\}$. Interchanging coordinates and
changing one sign is an isometry in the dual maximum norm.
For a unit graph vector $(x,Ax)$,
$(x,(A+tB)x)$ is at distance at most $t(k+C)$.
The gap lemma and the graph annihilator identity consequently give
\[
 \delta\big(G((A+tB)'),G(A')\big)\leq t(k+C).
\]
For a unit vector $(y',A'y')$ in the dual maximum graph norm,
its distance from $(y',(A'+tB')y')$ is at most $t(k'+C')$.
Thus $\delta(G(A'),G(A'+tB'))\leq t(k'+C')$.
If both bounds are at most $1/3$, the gap triangle inequality
gives a gap at most $7/9$ between $G((A+tB)')$ and its
closed subspace $G(A'+tB')$. A proper closed subspace would
give gap one. Therefore these two graphs are equal for all
sufficiently small $t$. This also justifies the domain assertion;
it is not a formal adjoint manipulation.

To reach $t=1$, suppose equality holds at some $t\in[0,1]$
and set $A_t=A+tB$. The graph estimate above gives the uniform
relative bounds
\[
 \|Bx\|\leq\frac{k}{1-k}\|A_tx\|+
                          \frac{C}{1-k}\|x\|,
\]
and, since $A_t'=A'+tB'$, the corresponding bound for $B'$
with $k',C'$. Choose $h_0>0$ so that $h_0$ times each of
the sums of these two coefficients is at most $1/3$; if a sum
is zero it imposes no restriction. These small relative
coefficients also make both perturbed graph domains complete,
by the same graph-norm argument. Apply the local proof with
base $A_t$ and increment $hB$, $0\leq h\leq h_0$.
It proves $A_{t+h}'=A_t'+hB'=A'+(t+h)B'$ as long as
$t+h\leq1$. Starting at zero, a finite subdivision of $[0,1]$
with mesh at most $h_0$ reaches one and proves the theorem.
\end{proof}

\begin{example}[The domain inclusion may otherwise be strict]
\label{fa2:p3:adjoint-sum-strict}
On $\ell^2$ let $Ax=(nx_n)$ on $D(A)=\{x:(nx_n)\in\ell^2\}$
and let $B=-A$ on that same domain. Both are self-adjoint and
closed, but $A+B$ is zero on $D(A)$ and its adjoint is zero
on all of $\ell^2$. In contrast $A'+B'$ remains defined only
on the corresponding proper dual domain. Thus a bare formal
sum does not identify the adjoint domain; the strict relative
bound in the theorem is substantive. Identifying the Hilbert
antilinear Riesz map with the scalar-linear dual transports this
example to the Banach-adjoint notation used there.
\end{example}
\begin{proof}
The maximal real multiplication operator is self-adjoint by the
proved spectral calculus. The domain is dense, since it contains
finitely supported sequences, and proper, for example excluding
$(1/n)_{n\geq1}$. Every vector pairs boundedly with the zero
operator, so the adjoint of the sum has the whole domain. The
two individual adjoint domains are the displayed weighted
square-integrability domain, proving strictness directly.
\end{proof}

% END INPUT


% BEGIN parts/part3-semigroups.tex
\chapter{Generation and approximation of semigroups}
In this chapter $X$ is an arbitrary real or complex Banach space.
A $C_0$ semigroup satisfies $T(0)=I$, $T(t+s)=T(t)T(s)$ and
$T(t)x\to x$ as $t\downarrow0$ for every $x$. Its generator is
$Ax=\lim_{h\downarrow0}(T(h)x-x)/h$ on the set where that norm
limit exists. We do not replace the general Banach generation theorem
by its Hilbert contraction special case. See \cite{B01,B04,B08}.

\begin{theorem}[Growth, generator, and Laplace resolvent]\label{fa2:p3:semigroup-basics}
Every $C_0$ semigroup has $\|T(t)\|\leq Me^{\omega t}$ for some
$M\geq1$, $\omega\in\mathbb R$. Its generator is closed and densely
defined, $T(t)D(A)\subset D(A)$, and
$\frac d{dt}T(t)x=AT(t)x=T(t)Ax$ on $D(A)$.
For $\operatorname{Re}\lambda>\omega$,
\[
 R(\lambda,A)x=\int_0^\infty e^{-\lambda t}T(t)x\,dt,
 \quad
 R(\lambda,A)^m x=\frac1{(m-1)!}
 \int_0^\infty t^{m-1}e^{-\lambda t}T(t)x\,dt.
\]
Thus $\|R(\lambda,A)^m\|\leq M/(\operatorname{Re}\lambda-\omega)^m$.
The generator uniquely determines the semigroup.
\end{theorem}
\begin{proof}
Uniform boundedness gives $\sup_{0\leq t\leq\delta}\|T(t)\|\leq C$
for some $\delta>0$; otherwise choose a sequence tending to zero with
unbounded norms, contradicting pointwise convergence. Writing
$t=n\delta+r$ bounds the norm by $C^{n+1}$, giving exponential growth.
The semigroup law and this bound extend strong continuity to every
$t\geq0$. For each $x$,
$x_h=h^{-1}\int_0^hT(s)x\,ds\to x$ and direct difference quotients
give $x_h\in D(A)$ and $Ax_h=(T(h)x-x)/h$.
This proves density. For $x\in D(A)$, commuting the difference
quotient with $T(t)$ gives domain invariance and the derivative.
Integrating it gives $T(t)x-x=\int_0^tT(s)Ax\,ds$.
If $x_n\to x$, $Ax_n\to y$, pass to this integral identity;
divide by $t$ and let $t\downarrow0$ to obtain $x\in D(A)$, $Ax=y$.
Thus $A$ is closed.
The Laplace integral converges absolutely by the growth bound.
Changing variables in $T(h)R(\lambda,A)x$ and taking its difference
quotient proves $AR(\lambda,A)x=\lambda R(\lambda,A)x-x$.
For $x\in D(A)$, integration by parts gives
$R(\lambda,A)(\lambda-A)x=x$. Hence the integral is the resolvent.
Fubini and the semigroup law turn its $m$th power into an integral
over $m$ positive times; the simplex of sum $t$ has volume
$t^{m-1}/(m-1)!$, giving the formula and estimate.
If $S,T$ have the same generator, for $x\in D(A)$ differentiate
$S(t-s)T(s)x$ on $0<s<t$; its derivative is zero. The endpoints
give $S(t)x=T(t)x$, and density proves equality everywhere.
\end{proof}

\begin{theorem}[Resolvent lines obstruct superexponential decay]\label{fa2:p3:resolvent-decay-obstruction}
Let $G$ generate a $C_0$ semigroup on a complex Hilbert space.
Suppose real $\xi_n\to-\infty$ satisfy
$\xi_n+i\mathbb R\subset\rho(G)$ and
$\sup_\tau\|R(\xi_n+i\tau,G)\|\leq d^{-1}$ with one $d>0$.
If $u\in C^1([0,\infty);H)\cap C([0,\infty);D(G))$ and
$\|u'(t)-Gu(t)\|\leq k\|u(t)\|$ for $k<d$, then either $u=0$
or $\int_0^\infty\|u(t)\|^2e^{bt}dt=\infty$ for some $b>0$.
This is a lower bound on possible decay, not a theorem asserting
exponential stability of arbitrary semigroups.
\end{theorem}
\begin{proof}
We first prove the weighted estimate
\[
 d^2\int_{\mathbb R}\|w(t)\|^2e^{-2\xi t}dt
 \leq\int_{\mathbb R}\|w'(t)-Gw(t)\|^2e^{-2\xi t}dt
\]
whenever the resolvent on $\xi+i\mathbb R$ has the stated bound.
Initially take compactly supported $w$ with local $W^{1,2}(H)$
regularity and $Gw\in L^2(H)$. For $v=e^{-\xi t}w$,
$e^{-\xi t}(w'-Gw)=v'-(G-\xi)v$.
Hilbert-valued Plancherel gives its Fourier transform
$(i\tau+\xi-G)\widehat v(\tau)$, and multiplication by the bounded
resolvent yields the inequality. The domain interchange is valid
because Fourier integrals of an integrable graph-valued function
lie in the closed graph domain; approximation gives the $L^2$
identity. Hilbert-valued Plancherel itself follows by expanding in
an orthonormal basis of the separable essential range and summing
the scalar identities by Tonelli.
For a locally graph-continuous $w$ with weighted $w$ and residual
in $L^2$, multiply by smooth cutoffs $\chi_R$, equal to one on
$[-R,R]$, supported in $[-2R,2R]$, with
$|\chi_R'|\leq C/R$. The residual of $\chi_Rw$ is
$\chi_R(w'-Gw)+\chi_R'w$; it converges in the weighted $L^2$
norm to the residual. The compact-support estimate and Fatou (or
ordinary $L^2$ convergence of $\chi_Rw$) give the same estimate.
Piecewise $C^1$ functions continuous at their junctions are allowed,
since their weak derivative has no point mass.
Suppose now all positive exponential weights of $u$ are square
integrable. Extend to $t<0$ by $u(0)\rho(t)$, with $\rho(0)=1$
and a smooth cutoff supported in $[-1,0]$. The extended function
has the local regularity just described. For $\xi=\xi_n<0$ the
weighted estimate, the residual bound on positive times, and
discarding the nonnegative negative-time left integral give
\[
 (d^2-k^2)\int_0^\infty\|u(t)\|^2e^{-2\xi_nt}dt
 \leq\int_{-1}^0\|Gu(0)\rho(t)-u(0)\rho'(t)\|^2
                                     e^{-2\xi_nt}dt.
\]
The right side tends to zero by dominated convergence, since
$\xi_n\to-\infty$ and $t<0$. The left side dominates
$(d^2-k^2)\int_0^T\|u(t)\|^2dt$ for every fixed $T$.
Thus $u=0$, contradicting nonzeroness. This proves the stated
alternative. The weighted Plancherel estimate and its cutoff
extension explicitly supply the integrability step needed by
\cite[Section 36.3]{B08}.
\end{proof}

\begin{corollary}[Self-adjoint spectral gaps]\label{fa2:p3:spectral-gap-decay-obstruction}
If $G$ is self-adjoint, bounded above, and its spectrum has gaps
$(\xi_n-d,\xi_n+d)$ with $\xi_n\to-\infty$, the preceding lower-decay
alternative applies to every trajectory with residual constant $k<d$.
\end{corollary}
\begin{proof}
The self-adjoint PVM gives
$\|R(\xi_n+i\tau,G)\|
=1/\operatorname{dist}(\xi_n+i\tau,\sigma(G))\leq1/d$.
Boundedness above makes $G$ a semigroup generator by its exponential
calculus. Thus all hypotheses of the preceding theorem hold.
\end{proof}

\begin{theorem}[Point spectrum and eventual compactness]\label{fa2:p3:semigroup-spectral-mapping}
For every complex Banach $C_0$ semigroup and every $t>0$,
$\sigma_p(T(t))\setminus\{0\}=e^{t\sigma_p(A)}$.
The full spectral mapping equality can fail: irrational rotations
on $L^2(\mathbb T)$ give a countable dense exponential image but
an entire unit-circle operator spectrum.
If $T(t_0)$ is compact for some $t_0>0$, full spectral mapping holds
for every $t>0$ after omitting zero. The generator spectrum then
consists of isolated eigenvalues of finite algebraic multiplicity,
with only finitely many in each right half-plane. Its growth bound
equals its spectral bound (both may be $-\infty$).
For any real $a$ whose vertical line misses the generator spectrum,
the part with real part greater than $a$ gives a finite exponential
polynomial expansion of $T(t)$, with remainder bounded by $C_a e^{at}$.
\end{theorem}
\begin{proof}
An $A$ eigenvector evolves as $e^{t\lambda}x$ by uniqueness of the
classical equation. Conversely if $T(t_0)x=zx$, $z\ne0$, choose
$\lambda$ with $e^{t_0\lambda}=z$ and define
$y_k=t_0^{-1}\int_0^{t_0}e^{-(\lambda+2\pi ik/t_0)s}T(s)x\,ds$.
The integrated generator identity and the endpoint equation give
$y_k\in D(A)$, $Ay_k=(\lambda+2\pi ik/t_0)y_k$.
At least one $y_k$ is nonzero: otherwise every scalar coefficient of
the continuous periodic function $e^{-\lambda s}T(s)x$ would have
all Fourier coefficients zero, forcing that function, hence $x$,
to vanish. This proves point spectral mapping.
For rotations, Fourier modes have generator eigenvalues $in$,
$n\in\mathbb Z$, and its resolvent is the diagonal inverse off this
set, so this is its whole spectrum. At an irrational multiple of
$2\pi$, the time operator has diagonal eigenvalues $e^{int}$ dense
in the circle and therefore spectrum the full circle; this proves
strictness of the general inclusion.
Now let $K=T(t_0)$ be compact. For each nonzero eigenvalue of $K$,
its finite-dimensional Riesz space reduces every $T(s)$ in the
algebraic direct-sum sense, since they commute with the resolvent.
On that space the continuous finite-dimensional semigroup is a
matrix exponential, so the whole space lies in $D(A)$ and its
generator is the matrix generator.
Cut off all $K$ eigenvalues of modulus greater than $r>0$.
Their total Riesz space is finite dimensional. On its complementary
space the spectral radius of $K$ is at most $r$, so the radius
formula and the locally bounded factors $T(s)$, $0\leq s\leq t_0$,
give a growth bound at most $\log r/t_0+\varepsilon$ there.
Its generator thus has no spectrum farther right by the Laplace
resolvent. Every right half-plane is therefore controlled by finitely
many matrix generators, proving the generator spectral assertions.
For any $t>0$, some power $B^n$ of $B=T(t)$ is compact. If
$\mu\ne0$ is in $\sigma(B)$, polynomial spectral mapping puts
$\mu^n$ in $\sigma(B^n)$. Its finite-dimensional Riesz space $E$
is invariant under $B$. On the complementary Riesz space,
$\mu^n-B^n$ is invertible; its commuting linear factorization
implies $\mu-B$ is invertible there. Thus $\mu$ is a spectral
value of the finite matrix $B|_E$, hence an eigenvalue of finite
algebraic multiplicity on the whole space. Point spectral mapping
now gives full mapping.
The growth bound is $\log r(T(t_0))/t_0$: the radius formula gives
the upper bound with an arbitrarily small excess, and the inequality
$r(T(t_0))\leq e^{\omega t_0}$ gives the reverse for every growth
bound $\omega$. Full spectral mapping identifies this with
$\sup\operatorname{Re}\sigma(A)$, including the quasinilpotent
case. For the last assertion take the finite Riesz sum above
$e^{at_0}$. Absence of eigenvalues on that circle leaves the
complementary radius strictly below it; powers there have the
stated $e^{at}$ bound. On the finite sum, Jordan's formula for the
matrix exponential is exactly the finite exponential polynomial.
Eventual compactness is not used to assert compact resolvent:
a nilpotent translation semigroup with infinite multiplicity is
eventually zero but can have noncompact resolvent.
\end{proof}

\begin{example}[Spectral bound need not be the growth bound]\label{fa2:p3:growth-spectral-gap}
On $H=\bigoplus_{n\geq1}\mathbb C^n$ let $S_n$ be the nilpotent
forward shift and
$A_n=(-1+in)I+S_n$.
The direct-sum generator has domain
$\{x:\sum n^2\|x_n\|^2<\infty\}$ and spectrum exactly
$\{-1+in:n\geq1\}$, but its contraction semigroup has
$\|T(t)\|=1$ for every $t\geq0$. Thus its spectral bound is $-1$
and its growth bound is zero. Also the adjoint of a $C_0$ semigroup
need not be $C_0$ on a nonreflexive space: translations on $L^1(\mathbb R)$
have adjoint translations on $L^\infty$, where a step function moves
distance one in the essential supremum norm for every nonzero shift.
\end{example}
\begin{proof}
The diagonal $inI$ generates unitaries; its bounded perturbation
$-I+\bigoplus S_n$ gives the displayed domain and component
semigroups $e^{(-1+in)t}e^{tS_n}$, each contractive.
Their direct sum is strongly continuous by finite-component
approximation and the common norm bound. For fixed $\lambda$ off
the displayed eigenvalues the block inverses exist; for large $n$
their norm is at most $(|\lambda+1-in|-1)^{-1}$ by the Neumann
series, which also bounds $n$ times their norm. The finitely many
remaining blocks cause no problem. Their direct sum is therefore a
bounded inverse into the exact generator domain, proving the spectrum.
For the unit vector with all $n$ coordinates $1/\sqrt n$,
$\|(S_n-I)x\|=1/\sqrt n$. Integrating the difference of the two
exponentials gives
$\|e^{tS_n}x-e^tx\|\leq te^t/\sqrt n$.
Thus $\sup_n\|e^{tS_n}\|=e^t$ and $\|T(t)\|=1$.
The adjoint translation example follows directly by testing an
indicator of a half-line; primal translation continuity on $L^1$
follows from the density of compact continuous functions.
\end{proof}

\begin{proposition}[Linear maximal monotone operators]\label{fa2:p3:linear-monotone}
If $L:D(L)\subset H\to H$ is linear, monotone, and $I+L$ is onto,
then it is closed, densely defined, and $I+\lambda L$ is onto for
every $\lambda>0$. Its resolvents are contractions and
$L_\lambda=L(I+\lambda L)^{-1}$ is bounded with
$\|L_\lambda x\|\leq\|Lx\|$ on $D(L)$,
$L_\lambda x\to Lx$ there. Moreover $D(L^2)$ is graph dense in
$D(L)$ and $-L$ generates the contraction semigroup above.
\end{proposition}
\begin{proof}
The lower bound $\|(I+L)x\|\geq\|x\|$ makes its inverse bounded
everywhere, hence $L$ closed. If $f\perp D(L)$ solve
$f=u+Lu$. Pairing with $u$ gives $0=\|u\|^2+\faip{u}{Lu}$,
so $u=0$ and $f=0$; thus the domain is dense.
Dissipativity of $-L$ and the one positive range condition give all
resolvents and generation by Lumer--Phillips.
For $J_\lambda=(I+\lambda L)^{-1}$, algebra on the exact domain gives
$LJ_\lambda x=J_\lambda Lx$ when $x\in D(L)$.
Also $J_\lambda x-x=-\lambda J_\lambda Lx\to0$ there, and
contractivity and density extend $J_\lambda\to I$ strongly.
The norm estimate and convergence of $L_\lambda x=J_\lambda Lx$
follow. Finally $J_\lambda x\in D(L^2)$ for $x\in D(L)$ because
$LJ_\lambda x=J_\lambda Lx\in D(L)$; the same two convergences
give graph density.
\end{proof}

\begin{theorem}[Semigroups of normal operators]\label{fa2:p3:normal-semigroup}
If every $T(t)$ in a $C_0$ semigroup on an arbitrary complex Hilbert space
is normal, its generator $A$ is normal, its spectrum lies in a
left half-plane, and $T(t)=e^{tA}$ in the unbounded measurable
calculus. Conversely every normal $A$ whose spectral real parts
are bounded above defines such a semigroup.
\end{theorem}
\begin{proof}
Fuglede makes all $T(s),T(t)^*$ commute. The unital C*-algebra they
generate is therefore commutative, with a PVM $E$ on its character
space. Let $A_h=(T(h)-I)/h$ and write $a_h$ for its continuous
scalar symbol. Since $T(2h)=T(h)^2$,
$a_{2h}-a_h=(h/2)a_h^2$.
For $x\in D(A)$ the norms $\|A_hx\|$ are bounded for small $h$.
Hence
\[
 \int|a_{2h}-a_h|\,d\mu_x
 \leq\frac h2\|A_hx\|^2.
\]
Summing over dyadic $h$ shows $a_{2^{-n}}$ converges almost everywhere
for each such $\mu_x$. Define the Borel function $b$ as that finite
limit where it exists and zero elsewhere. The bad set has zero
spectral projection, because it annihilates the dense set $D(A)$.
Fatou gives $D(A)\subset D(b(E))$.
The exponential series and the growth bound on $T(nh)$ give
$\|e^{tA_h}\|\leq M\exp(t(e^{\omega h}-1)/h)$, uniformly bounded
for $t$ in compact intervals and small $h$.
As these are normal, their scalar symbols satisfy a common bound;
at $t=1$ it gives a uniform upper bound on $\operatorname{Re}a_h$,
and therefore on $\operatorname{Re}b$.
For $x\in D(A)$, differentiate
$e^{(t-s)A_h}T(s)x$ and integrate. The common exponential bound
gives $\|e^{tA_h}x-T(t)x\|\leq C_Tt\|A_hx-Ax\|\to0$.
Density extends this to every vector. Along the dyadic sequence,
bounded convergence of the symbols also gives
$e^{tA_h}x\to e^{tb}(E)x$. Hence $T(t)=e^{tb}(E)$.
The generator of the latter has exactly domain $D(b(E))$:
dominated convergence gives the derivative there, using
$|(e^{hb}-1)/h|\leq|b|e^{h\max(\operatorname{Re}b,0)}$;
Fatou shows that existence of a norm derivative forces this square
integrability. Thus $A=b(E)$ and is normal. Its spectral support
lies in the half-plane bounding $\operatorname{Re}b$.
Conversely bounded real parts make $e^{tz}$ a bounded symbol for
$t\geq0$; the PVM product law and dominated convergence give the
semigroup and strong continuity, with the same derivative-domain
argument identifying its generator.
\end{proof}

\begin{theorem}[General Hille--Yosida theorem]\label{fa2:p3:hille-yosida}
A densely defined closed $A$ generates a $C_0$ semigroup with
$\|T(t)\|\leq Me^{\omega t}$ iff $(\omega,\infty)\subset\rho(A)$
and
\[
 \|R(\lambda,A)^m\|\leq\frac{M}{(\lambda-\omega)^m}
 \qquad(\lambda>\omega,\ m=1,2,\ldots).
\]
All resolvent powers are required when $M>1$.
For $\omega=0$, the bounded Yosida approximants
$A_\lambda=\lambda AR(\lambda,A)=\lambda^2R(\lambda,A)-\lambda I$
satisfy $e^{tA_\lambda}x\to T(t)x$ uniformly on bounded time intervals.
\end{theorem}
\begin{proof}
Necessity was proved above. Shift by $\omega I$ to reduce sufficiency
to $\omega=0$. The estimate on powers and the exponential series give
\[
 e^{tA_\lambda}=e^{-\lambda t}
 \sum_{k\geq0}\frac{(\lambda t)^k}{k!}(\lambda R(\lambda,A))^k,
 \qquad\|e^{tA_\lambda}\|\leq M.
\]
For $x\in D(A)$,
$\lambda R(\lambda,A)x-x=R(\lambda,A)Ax\to0$.
Uniform boundedness of $\lambda R(\lambda,A)$ and density extend
this convergence to every $x$. Therefore
$A_\lambda x=\lambda R(\lambda,A)Ax\to Ax$ on $D(A)$.
Resolvents commute, so all $A_\lambda$ commute. The fundamental
theorem of calculus applied to the product of their exponentials gives
\[
 e^{tA_\lambda}x-e^{tA_\mu}x
 =\int_0^t e^{(t-s)A_\lambda}e^{sA_\mu}
                   (A_\lambda-A_\mu)x\,ds.
\]
Its norm is at most $M^2t\|(A_\lambda-A_\mu)x\|$ for $x\in D(A)$.
Hence the exponentials are uniformly Cauchy there on bounded intervals.
Their uniform bounds extend convergence to all $x\in X$.
Call the limit $T(t)x$. Passing to limits proves the semigroup law,
the bound $\|T(t)\|\leq M$, and strong continuity, since convergence
is uniform down to $t=0$. For $x\in D(A)$, pass to the limit in
$e^{tA_\lambda}x-x=\int_0^te^{sA_\lambda}A_\lambda x\,ds$.
This gives $T(t)x-x=\int_0^tT(s)Ax\,ds$, so its generator $G$
extends $A$. For a positive $\lambda$, both $\lambda-A$ and
$\lambda-G$ are onto and the latter is injective by the Laplace
formula. For $x\in D(G)$ choose $y\in D(A)$ with
$(\lambda-A)y=(\lambda-G)x$; injectivity gives $x=y$.
Thus $G=A$. Undoing the shift proves the theorem and the stated
approximation, with shifted approximants if $\omega\ne0$.
\end{proof}

\begin{theorem}[Lumer--Phillips on Banach spaces]\label{fa2:p3:lumer-phillips}
A densely defined operator $A$ generates a contraction $C_0$ semigroup
iff it is dissipative,
$\|(\lambda-A)x\|\geq\lambda\|x\|$ for every $\lambda>0$, $x\in D(A)$,
and $\faRan(\lambda_0-A)=X$ for some $\lambda_0>0$.
On a Hilbert space dissipativity is equivalent to
$\operatorname{Re}\faip{x}{Ax}\leq0$ on $D(A)$.
\end{theorem}
\begin{proof}
Necessity follows from the Laplace resolvent bound of a contraction
semigroup. Conversely the range assumption and dissipativity give
a bounded everywhere inverse at $\lambda_0$, making $A$ closed.
Let $J$ be the set of positive $\lambda$ for which the range is all
of $X$. It is open by the geometric resolvent series, with inverse
norm at most $1/\lambda$. It is closed relative to $(0,\infty)$:
if $\lambda_n\to\lambda>0$, the resolvent identity and these bounds
make $R(\lambda_n,A)y$ Cauchy; closedness supplies a preimage of
$y$ under $\lambda-A$. Thus $J=(0,\infty)$.
The inverse norm bound implies every power bound with $M=1$;
Hille--Yosida gives a contraction semigroup.
On a Hilbert space the form inequality gives
$\|(\lambda-A)x\|^2\geq\lambda^2\|x\|^2+\|Ax\|^2$.
Conversely square the dissipativity estimate, expand, divide by
$2\lambda$ and let $\lambda\to\infty$.
\end{proof}

\begin{theorem}[Classical and inhomogeneous evolution]\label{fa2:p3:evolution}
If $A$ generates $T$ and $x\in D(A)$, $u(t)=T(t)x$ is the unique
classical solution of $u'=Au$, $u(0)=x$.
For $x\in D(A^k)$,
$u\in C^{k-j}([0,\infty);D(A^j))$ for $0\leq j\leq k$,
using graph norms. If $x\in D(A)$ and $f\in C^1([0,T];X)$, then
\[
 u(t)=T(t)x+\int_0^tT(t-s)f(s)\,ds
\]
is the unique classical solution of $u'=Au+f$.
For $x\in X$, $f\in L^1(0,T;X)$ the same formula defines its
continuous mild solution and depends continuously on the data.
\end{theorem}
\begin{proof}
The homogeneous assertions follow by iterating domain invariance and
$AT(t)x=T(t)Ax$. Uniqueness follows by differentiating
$T(t-s)u(s)$; on a bounded interval all derivatives are legitimate
on the graph domain. For the convolution, rewrite it as
$v(t)=\int_0^tT(s)f(t-s)\,ds$. Its norm derivative is
$v'(t)=T(t)f(0)+\int_0^tT(s)f'(t-s)\,ds$.
The integrated generator identity
$A\int_a^bT(s)y\,ds=T(b)y-T(a)y$ follows first for $y\in D(A)$
and then by closedness for every $y$. Integration by parts in the
smooth scalar time variable, or approximation of $f$ by piecewise
linear functions and this identity, yields
$v(t)\in D(A)$ and $Av(t)=v'(t)-f(t)$.
The same estimates show graph continuity. Uniqueness reduces to
the homogeneous case. For $L^1$ forcing, approximate by continuous
step-smoothed functions; the convolution norm is bounded by
$\sup_{s\leq T}\|T(s)\|\|f\|_{L^1}$, proving uniform convergence
and continuous dependence. This establishes the mild assertion.
\end{proof}

\begin{theorem}[Bounded and relatively bounded perturbations]\label{fa2:p3:semigroup-perturbation}
If $A$ generates a semigroup bounded by $Me^{\omega t}$ and
$B\in\faB(X)$, then $A+B$ on $D(A)$ generates a semigroup bounded
by $M e^{(\omega+M\|B\|)t}$.
On a Hilbert space, if $A$ generates contractions, $B$ is dissipative
on a domain containing $D(A)$, and
$\|Bx\|\leq a\|Ax\|+b\|x\|$ with $a<1$, then $A+B$ on $D(A)$
also generates contractions.
\end{theorem}
\begin{proof}
Define $S_0(t)=T(t)$ and
$S_{n+1}(t)x=\int_0^tT(t-s)BS_n(s)x\,ds$.
Induction gives
$\|S_n(t)\|\leq M^{n+1}\|B\|^nt^ne^{\omega t}/n!$.
Thus their sum $S(t)$ converges uniformly on bounded time intervals,
has the stated bound, and solves its Volterra equation. Splitting the
integral at a time $s$ shows $S(t+s)$ and $S(t)S(s)$ solve the same
Volterra equation with the same initial value; iterating its difference
and the factorial estimate proves uniqueness, hence the semigroup law.
Its Laplace transform for large $\lambda$ is the norm-convergent
Neumann sum
$R(\lambda,A)\sum_{n\geq0}(BR(\lambda,A))^n$,
which is the inverse of $\lambda-A-B$ on $D(A)$.
The generator has the same resolvent, so equals $A+B$ with that domain.
For the relative Hilbert assertion, dissipativity of $A$ gives
$\|AR(\lambda,A)\|\leq1$ and
$\|R(\lambda,A)\|\leq1/\lambda$ by the squared estimate in the
preceding theorem. Put $a_+=\max(a,0)$ and $b_+=\max(b,0)$.
Then $\|BR(\lambda,A)\|\leq a_++b_+/\lambda<1$ for large $\lambda$.
The relative bound makes the graph norms of $A$ and $A+B$ equivalent,
so the sum is closed; it is dissipative and has a positive resolvent
point by the factorization and Neumann series. Lumer--Phillips proves
the assertion. The strict bound $a<1$ is essential to this argument.
\end{proof}

\begin{theorem}[Strong resolvent approximation]\label{fa2:p3:trotter-kato}
Let $A_n,A$ generate contraction semigroups $T_n,T$ on $X$.
The following are equivalent: for some (then every) $\lambda>0$,
$R(\lambda,A_n)x\to R(\lambda,A)x$ for every $x$; and
$T_n(t)x\to T(t)x$ uniformly for $t$ in bounded intervals, for every $x$.
The same equivalence holds under a common growth bound $Me^{\omega t}$.
\end{theorem}
\begin{proof}
Semigroup convergence implies resolvent convergence by dominated
Laplace integration. A resolvent identity and finite chains of disks
of radius smaller than the reciprocal common resolvent bound extend
convergence at one real resolvent point to all sufficiently large
real points. For a fixed large $\mu$, the bounded Yosida generators
$A_{n,\mu}=\mu^2R(\mu,A_n)-\mu I$ converge strongly, with uniformly
bounded norms, so their exponentials converge strongly uniformly on
bounded time intervals by their power series.
We need uniformity when $\mu\to\infty$. In the contraction case the
Laplace formula gives $\mu R(\mu,A_n)x=\mathbb E T_n(E_\mu)x$,
where $E_\mu$ is exponentially distributed with mean $1/\mu$.
The exponential series then gives
$e^{tA_{n,\mu}}x=\mathbb E T_n(S)x$, where $S$ is a sum of a
Poisson number of independent such exponentials, the Poisson mean
being $\mu t$. Thus $\mathbb ES=t$ and
$\mathbb E(S-t)^2=2t/\mu$.
For $x\in D(A_n)$,
$\|T_n(s)x-T_n(t)x\|\leq|s-t|\|A_nx\|$.
Consequently
\[
 \|e^{tA_{n,\mu}}x-T_n(t)x\|
 \leq(2t/\mu)^{1/2}\|A_nx\|.
\]
Given $x\in D(A)$, write $x=R(\lambda,A)y$ and take
$x_n=R(\lambda,A_n)y$. Then $x_n\to x$ and
$A_nx_n=\lambda x_n-y$ is uniformly bounded. The preceding estimate,
the analogous estimate for $A$, and fixed-$\mu$ exponential
convergence prove semigroup convergence on the dense domain $D(A)$.
Uniform contraction bounds extend it to $X$.
For a common $Me^{\omega t}$ bound, shift by $\omega$ and use the
same random-time representation. The Lipschitz estimate acquires
$M e^{\omega\max(s,t)}$. Cauchy--Schwarz and the compound-Poisson
moment generating function
$\mathbb E e^{aS}=\exp(\mu t(\mu/(\mu-a)-1))$ show, for fixed
bounded $t$ and $\mu>2|a|+1$, that the error is bounded by
$C_T\mu^{-1/2}(\|A_nx\|+\|x\|)$, uniformly in $n$.
Indeed differentiating this formula twice bounds
$\mathbb E[(S-t)^2e^{2\omega S}]$ by $C_T/\mu$.
The rest of the argument is unchanged.
\end{proof}

\begin{theorem}[Chernoff products, implicit Euler, and Trotter]\label{fa2:p3:product-formula}
Let $F(h)$ be contractions with $F(0)=I$, and suppose on a core $D_0$
of a contraction generator $A$ that
$(F(h)x-x)/h\to Ax$. Then
$F(t/n)^nx\to T(t)x$ strongly, uniformly on bounded $t$ intervals.
In particular
$T(t)x=\lim_n(I-tA/n)^{-n}x$.
If $A,B$ generate contractions and the closure $C$ of $A+B$ on
$D(A)\cap D(B)$ generates contractions, then
\[
 (T_A(t/n)T_B(t/n))^nx\longrightarrow T_C(t)x.
\]
For bounded $A,B$, Strang's symmetric product
$(e^{tA/(2n)}e^{tB/n}e^{tA/(2n)})^n$ has operator-norm error
$O_T(n^{-2})$ on $0\leq t\leq T$.
\end{theorem}
\begin{proof}
Put $G_h=(F(h)-I)/h$. The exponential series shows $G_h$ generates
contractions. For $x\in D_0$,
$(\lambda-G_h)x\to(\lambda-A)x$; multiplying by the uniformly
bounded resolvents shows their strong convergence on
$(\lambda-A)D_0$, a dense set because $D_0$ is a graph core.
The bounds extend it to all $X$. Strong resolvent approximation gives
$e^{tG_h}\to T(t)$ uniformly on bounded intervals.
For $x\in D_0$ and integer $n$, the Poisson expansion of
$e^{n(F-I)}$ and the telescoping contraction estimate give
\[
 \|F^nx-e^{n(F-I)}x\|
 \leq\mathbb E|N-n|\|(F-I)x\|
 \leq\sqrt n\|(F-I)x\|,
\]
where $N$ is Poisson of mean $n$ and $F=F(h)$.
With $h=t/n$, consistency makes this error $O_T(n^{-1/2})$ for
fixed $x\in D_0$ (uniformity near zero follows directly from the
definition of its derivative). Density and contraction bounds give
the product convergence for all $x$.
For implicit Euler, $F(h)=(I-hA)^{-1}$ is contractive and
$(F(h)-I)x/h=F(h)Ax\to Ax$ on $D(A)$.
For the Trotter product, $F(h)=T_A(h)T_B(h)$ is contractive and
its derivative on the intersection is $A+B$; that domain is a core
for its closure by definition, so the preceding result applies.
For bounded operators, multiply their norm Taylor series through
degree two. The symmetric product agrees with $e^{h(A+B)}$ in
degrees zero, one and two, and their remainder norms are at most
$C_T h^3$. Telescoping $n$ products, whose partial norms are
uniformly bounded by $e^{T(\|A\|+\|B\|)}$, gives
$nC_T(t/n)^3=O_T(n^{-2})$. An unbounded second-order assertion needs
additional common-domain and commutator estimates; it is not inferred
from the first-order Trotter hypotheses alone.
\end{proof}

\begin{proposition}[Lax consistency and stability]\label{fa2:p3:lax-equivalence}
Suppose bounded time steps $C_h$ satisfy
$\sup_{nh\leq T}\|C_h^n\|<\infty$ and, on a graph core of a
generator $A$, $h^{-1}(C_h-I)x\to Ax$. If the stability estimate
and consistency extend to the graph-domain orbit uniformly on its
compact time intervals, then $C_h^nx\to T(t)x$ as $h\downarrow0$ with $nh\to t$,
uniformly for bounded times. Conversely, convergence for every $x$
in this same sense implies asymptotic stability: for every finite
$T\geq0$ there exist $h_T>0$ and $M_T<\infty$ such that
$\|C_h^n\|\leq M_T$ whenever $0<h<h_T$ and $nh\leq T$.
\end{proposition}
\begin{proof}
For $x$ in the graph domain use
\[
 C_h^nx-T(nh)x
 =\sum_{j=0}^{n-1}C_h^{n-1-j}(C_h-T(h))T(jh)x.
\]
Uniform consistency on the compact graph-domain orbit makes each
middle error $o(h)$ uniformly in $j$. Stability bounds the sum by
$C_T n o(h)\to0$. Density and the common bounds extend convergence
to every $x$. For necessity, fix finite $T>0$. If no small-step
uniform bound existed, choose $h_k\downarrow0$ and integers $n_k$
with $n_kh_k\leq T$ and $\|C_{h_k}^{n_k}\|>k$. Bounded-time
convergence makes $\{C_{h_k}^{n_k}x:k\geq1\}$ bounded for each
$x$, since $\{T(n_kh_k)x:k\geq1\}$ is bounded. Uniform boundedness
then contradicts the chosen operator norms. For $T=0$, only $n=0$
occurs and the bound is immediate.
This formulation makes the needed orbit-uniform consistency explicit;
it does not hide it in a bare pointwise consistency slogan.
\end{proof}

\begin{theorem}[Adjoint semigroups and spectral inclusion]\label{fa2:p3:adjoint-semigroup}
The adjoints $T(t)'$ form a weak-star continuous semigroup on $X'$.
If $X$ is reflexive, it is a $C_0$ semigroup with generator $A'$.
In general strong continuity on the entire dual can fail.
For every $C_0$ semigroup,
$e^{t\sigma(A)}\subset\sigma(T(t))\setminus\{0\}$;
equality need not hold without further hypotheses.
\end{theorem}
\begin{proof}
Weak-star continuity is the identity
$(T(t)'x')(x)=x'(T(t)x)$. On a reflexive space it is weak
continuity. To obtain strong continuity, the averages
$x'_h=h^{-1}\int_0^hT(s)'x'\,ds$ are norm-continuous vectors:
translation of the interval integral gives their norm continuity
using the local operator bound. Their span is norm dense, since its
annihilator in $X''=X$ is zero by weak convergence of these averages
to $x'$. Uniform boundedness then extends strong continuity to every
$x'$. Differentiating the pairing on $D(A)$ gives that its generator
is contained in $A'$. The Laplace resolvent is $R(\lambda,A)'$, the
inverse of $\lambda-A'$ on its exact domain, so equality follows.
For the spectral inclusion define
$B=\int_0^te^{\lambda(t-s)}T(s)\,ds$.
The integrated generator identity gives
$(\lambda-A)B=e^{\lambda t}I-T(t)$ and the corresponding identity
on $D(A)$ in reverse order. If the bounded right side were invertible,
its inverse commutes with $T(s)$ and the resolvent, so
$B(e^{\lambda t}I-T(t))^{-1}$ would invert $\lambda-A$.
Thus $\lambda\in\sigma(A)$ forces $e^{\lambda t}\in\sigma(T(t))$.
\end{proof}

\begin{theorem}[Stone]\label{fa2:p3:stone}
Every strongly continuous unitary group $U(t)$, $t\in\mathbb R$,
has a unique self-adjoint $H$ with $U(t)=e^{itH}$.
Its derivative generator is $iH$, with exactly $D(H)$ as domain.
Conversely every self-adjoint $H$ gives such a group by its PVM.
\end{theorem}
\begin{proof}
The positive-time generator $G$ is closed and dense by the semigroup
theorem; the negative group has generator $-G$ on the same domain,
as follows by changing variables in its difference quotients.
Differentiating $\faip{U(t)x}{U(t)y}$ on $D(G)$ gives
$\faip{Gx}{y}+\faip{x}{Gy}=0$.
Thus $H=-iG$ is symmetric. The Laplace formula for both time
directions makes $\lambda-G$ and $\lambda+G$ onto for $\lambda>0$.
Equivalently $H\pm i\lambda$ are onto, so $H$ is self-adjoint by
the deficiency criterion. Its PVM defines $V(t)=e^{itH}$; dominated
convergence gives strong continuity and, on $D(H)$, derivative
$iHV(t)x$. The semigroup uniqueness theorem gives $U(t)=V(t)$.
Conversely the same dominated convergence argument gives the group
and derivative from a self-adjoint $H$. If a vector has a norm
difference quotient, its norms are bounded; Fatou applied to
$|(e^{it\lambda}-1)/t|^2$ forces $\int\lambda^2d\mu_x<\infty$.
Thus the derivative domain is exactly $D(H)$ and uniqueness follows.
\end{proof}

\begin{theorem}[Self-adjoint smoothing]\label{fa2:p3:selfadjoint-smoothing}
For $L=L^*\geq0$, $T(t)=e^{-tL}$ is a self-adjoint contraction
semigroup. For every $t>0$, $T(t)H\subset\bigcap_{m\geq1}D(L^m)$ and
\[
 \|L^mT(t)\|\leq(m/(et))^m.
\]
Its trajectories are smooth away from zero and are the unique
continuous-at-zero classical solutions of $u'+Lu=0$ there.
If $L\geq cI$ with $c>0$, then $\|T(t)\|\leq e^{-ct}$.
\end{theorem}
\begin{proof}
All assertions follow from the measurable calculus with symbol
$e^{-t\lambda}$ and the elementary maximum
$\sup_{\lambda\geq0}\lambda^me^{-t\lambda}=(m/(et))^m$.
On a time interval bounded away from zero, every differentiated
symbol is uniformly bounded, so dominated convergence gives all
derivatives and graph-domain continuity. For uniqueness, the energy
derivative of a difference is
$\frac d{dt}\|u\|^2=-2\faip{u}{Lu}\leq0$; continuity at zero
forces a zero-initial difference to vanish. The lower spectral bound
gives the last norm estimate directly.
\end{proof}

% END INPUT


% BEGIN parts/part3-triples.tex
\chapter{Hilbert triples and the parabolic energy method}
Let $V\hookrightarrow H$ be a dense continuous injection of real
Hilbert spaces. Identifying $H$ with its dual gives
$V\hookrightarrow H\hookrightarrow V'$, with
$\langle h,v\rangle=\faip{h}{v}_H$.
Bochner spaces and their integrals have the meaning in
\ref{fa2:p4:bochner}. On a nonseparable $V$ we explicitly require
the strong measurability of the operator trajectories below; scalar
weak measurability alone is not automatically Bochner measurability.
This distinction is invisible in the usual separable PDE examples
but must not be hidden in the abstract theorem.

\begin{lemma}[The Hilbert-triple trace and energy identity]\label{fa2:p3:triple-trace}
Let $0<T<\infty$. If $u\in L^2(0,T;V)$ and its distributional derivative belongs to
$L^2(0,T;V')$, then it has a unique representative in $C([0,T];H)$,
and
\[
 \|u(t)\|_H^2-\|u(s)\|_H^2
 =2\int_s^t\langle u'(r),u(r)\rangle\,dr.
\]
The trace map into $C(H)$ is continuous for the sum of these two
$L^2$ norms.
\end{lemma}
\begin{proof}
For a smooth $V$-valued function the identity is differentiation of
its squared $H$ norm, since the triple pairing agrees with the
$H$ product. It gives, after choosing a time $s$ with squared norm
no greater than its time average,
\[
 \sup_t\|u(t)\|_H^2
 \leq T^{-1}\|u\|_{L^2(H)}^2
       +2\|u'\|_{L^2(V')}\|u\|_{L^2(V)}.
\]
To approximate a general $u$, first regard it as a
$W^{1,2}(0,T;V')$ function, which has an absolutely continuous
$V'$ representative (the fundamental theorem follows by integrating
its distributional derivative and testing functionals). Reflect it
across the two endpoints. Reflection preserves local $L^2(V)$ and
$L^2(V')$ derivative membership; no jump appears because the
$V'$ traces agree at the reflected endpoint.
Convolve with a smooth compactly supported time mollifier.
The resulting $u_\varepsilon$ is smooth with values in $V$, and
$u_\varepsilon\to u$ in $L^2(V)$,
$u_\varepsilon'\to u'$ in $L^2(V')$.
Apply the displayed estimate to their differences. It makes them
Cauchy in $C(H)$, defining the desired representative and proving
continuity of the trace map. The energy integrands converge in
$L^1$ by Cauchy--Schwarz, so passing to the limit gives the identity
at every pair of endpoint times. Two continuous representatives
agree almost everywhere and hence everywhere, proving uniqueness.
\end{proof}

\begin{theorem}[Nonautonomous parabolic forms]\label{fa2:p3:parabolic-triple}
Suppose $a(t;u,v)$ is bilinear, $t\mapsto a(t;u,v)$ measurable for
fixed $u,v$, and for almost every $t$, simultaneously for all $u,v$,
\[
 |a(t;u,v)|\leq M\|u\|_V\|v\|_V,\qquad
 a(t;v,v)\geq\alpha\|v\|_V^2-C\|v\|_H^2,\quad\alpha>0.
\]
Assume additionally $t\mapsto A(t)v$ is strongly measurable in $V'$
for each $v$, where $\langle A(t)u,v\rangle=a(t;u,v)$;
this is automatic if $V$ is separable. For every $f\in L^2(0,T;V')$
and $u_0\in H$, there is a unique
\[
 u\in L^2(0,T;V)\cap C([0,T];H),\quad u'\in L^2(0,T;V')
\]
with $u'+A(t)u=f$ in $V'$ almost everywhere and $u(0)=u_0$.
Its norm in these spaces is bounded by a constant depending only on
the displayed parameters, the embedding and $T$, times
$\|u_0\|_H+\|f\|_{L^2(V')}$.
\end{theorem}
\begin{proof}
The strong-measurability assumption and uniform bound make
$u\mapsto A(\cdot)u(\cdot)$ a bounded map $L^2(V)\to L^2(V')$:
prove it for simple functions and approximate a strongly measurable
$u$ in $L^2$. For a finite-dimensional subspace $E\subset V$ solve
the Galerkin system
$\faip{u_E'}{v}_H+a(t;u_E,v)=\langle f,v\rangle$ for $v\in E$,
with initial value the $H$-orthogonal projection of $u_0$ onto $E$.
In a basis of $E$ its positive definite mass matrix is invertible
and the coefficients are bounded measurable. Picard iteration of
the associated integral equation converges, with the usual
factorial estimate on bounded intervals, giving an absolutely
continuous finite-dimensional solution.
Test with $u_E$. Young's inequality and the Garding bound give
\[
 \tfrac12\frac d{dt}\|u_E\|_H^2
 +\tfrac\alpha2\|u_E\|_V^2
 \leq C\|u_E\|_H^2+(2\alpha)^{-1}\|f\|_{V'}^2.
\]
Gronwall and integration bound $u_E$ uniformly in $L^\infty(H)$
and $L^2(V)$. Direct finite-dimensional subspaces by inclusion.
Reflexivity of $L^2(V)$ gives a weakly convergent subnet to $u$.
For each fixed $v\in V$, the Galerkin equation holds eventually
in that subnet. Multiply by a smooth scalar time test, integrate
by parts, and pass to the weak limit. The bilinear term is a
continuous linear functional on $L^2(V)$, since $A(\cdot)'v$
is strongly measurable in the separable case; in the stated general
case the required scalar functional is bounded and measurable by
simple approximation and the uniform bound, so it still defines
an $L^2$ functional. The resulting equation is
$u'=f-Au\in L^2(V')$ in distributions. Tests not vanishing at
zero, but vanishing at $T$, retain the initial term
$\faip{u_0}{v}$ because the projections converge strongly to $u_0$.
The trace lemma now gives $u\in C(H)$ and identifies $u(0)=u_0$.
Testing the limit equation with $u$ is justified by that lemma and
gives the same energy estimate. The derivative estimate follows
from $u'=f-Au$. For two solutions the difference has zero forcing
and initial value; its energy inequality and Gronwall give zero.
Thus the solution is unique and all stated bounds follow.
No compact embedding and no countable Hilbert basis is used in the
net argument; where only scalar source measurability is imposed,
separability or the explicit strong-measurability bridge is needed.
\end{proof}

% END INPUT


% BEGIN parts/part3-hyperbolic.tex
\chapter{Time-dependent symmetric wave equations in a Hilbert triple}
\label{fa2:p3:hyperbolic-chapter}
This is the abstract hyperbolic theorem of \cite[Theorem 10.14]{B01}.
It is independent of a spatial differential realization and requires
no compact embedding. We include the weak-solution uniqueness argument
which recovers strong energy-space continuity; weak compactness by
itself does not supply that conclusion.

\begin{theorem}[Nonautonomous symmetric hyperbolic equation]
\label{fa2:p3:hyperbolic-triple}
Let $V\hookrightarrow H$ be a dense continuous embedding of real
Hilbert spaces and identify $H$ with a dense subspace of $V'$ by
its inner product. Suppose $a(t;u,v)$, $0\le t\le T$, is a family
of symmetric bounded bilinear forms on $V$, that each scalar function
$t\mapsto a(t;u,v)$ is $C^1$, and that uniformly in $t$,
\[
 a(t;v,v)\ge\alpha\|v\|_V^2-C\|v\|_H^2
 \qquad(\alpha>0, C\ge0).
\]
Write $\langle A(t)u,v\rangle=a(t;u,v)$.
For $u_0\in V$, $v_0\in H$ and $f\in L^2(0,T;H)$ there is a
unique solution
\[
 u\in C([0,T];V),\qquad u'\in C([0,T];H),\qquad
 u''\in L^2(0,T;V')
\]
of $u''+A(t)u=f$ a.e., with $u(0)=u_0$, $u'(0)=v_0$.
It satisfies the energy identity, for every $s,t$,
\begin{align*}
 &\|u'(t)\|_H^2+a(t;u(t),u(t))
   -\|u'(s)\|_H^2-a(s;u(s),u(s))\\
 &\hspace{1cm}=2\int_s^t(f(r),u'(r))_H\,dr
                +\int_s^t\partial_ra(r;u(r),u(r))\,dr.
\end{align*}
There is a constant $M_T$, depending only on the embedding and form
bounds, such that
\[
 \sup_t\big(\|u(t)\|_V^2+\|u'(t)\|_H^2\big)
 \le M_T\left(\|u_0\|_V^2+\|v_0\|_H^2+\|f\|_{L^2H}^2\right).
\]
The Hilbert spaces need not be separable.
\end{theorem}
\begin{proof}
\emph{Uniform bounds and a separable invariant carrier.}
Uniform boundedness, twice, applied to the scalar $C^1$ families
gives a constant $M$ with
$|a(t;u,v)|+|\partial_ta(t;u,v)|\le M\|u\|_V\|v\|_V$.
In particular the representing self-adjoint operators $B_t$ on
$V$, defined by $(B_tu,v)_V=a(t;u,v)$, are operator-norm Lipschitz
in $t$, by integrating the scalar derivative and taking the supremum
over the two unit balls. Let $J:V\to H$ be the embedding.
To avoid any nonseparable Bochner issue, take a countable set of
vectors spanning the essential range of the given strongly measurable
$f$ densely, and form the closed span $V_0$ of all finite words in
the bounded self-adjoint operators $J^*J$ and $B_r$ for rational
$r\in[0,T]$, applied to $u_0,J^*v_0$, and the $J^*$ images of
that countable set. This is a separable subspace reducing every
one of those operators. Norm continuity of $B_t$ makes it reducing
for every $t$. Put $H_0=\overline{JV_0}^{H}$.
For $w\in V_0^\perp$, $(Jw,Jv)_H=(w,J^*Jv)_V=0$ if $v\in V_0$.
Thus $JV_0^\perp\perp H_0$. Since $JV$ is dense in $H$, the
conditions $J^*v_0,J^*f(t)\in V_0$ imply $v_0,f(t)\in H_0$,
the latter a.e. Also $a(t;V_0,V_0^\perp)=0$.
A solution in the separable triple $V_0\subset H_0\subset V_0'$
therefore solves the full equation, because all complementary tests
vanish. We may perform the construction in that separable triple.
This reduction concerns the fixed data, not a separability assumption
on the theorem.

\emph{Galerkin existence in the weak energy class.}
Choose increasing finite-dimensional subspaces $V_n$ with dense
union in $V_0$. Let $u_{0n}\to u_0$ in $V_0$ and let $v_{0n}$ be
the $H_0$-orthogonal projections of $v_0$ onto $V_n$; then
$v_{0n}\to v_0$ in $H_0$. Solve the finite mass-matrix system
\[
 (u_n''(t),w)_H+a(t;u_n(t),w)=(f(t),w)_H\quad(w\in V_n)
\]
with these data. Its mass matrix is positive definite. Written as
a first-order integral equation its coefficient matrix is continuous
and the forcing integrable; Picard iteration with factorial bounds
gives a unique absolutely continuous state. Thus its velocity is
continuous and its acceleration is $L^2$ on each finite system.
Choose $\kappa\ge C$ and write
$\widetilde a_t=a_t+\kappa(\cdot,\cdot)_H$, which is coercive on $V$.
Testing with $u_n'$ gives, for
$E_n=\|u_n'\|_H^2+\widetilde a_t(u_n,u_n)$,
\[
 E_n'=2(f,u_n')_H+\partial_ta_t(u_n,u_n)
                         +2\kappa(u_n,u_n')_H
       \le c E_n+\|f\|_H^2.
\]
Coercivity and Gronwall bound $u_n$ in $L^\infty V_0$ and $u_n'$
in $L^\infty H_0$. Weak-star compactness in these spaces, with
separable preduals, gives a subsequence converging to $u,v$.
The integrated identity $u_n(t)=u_{0n}+\int_0^tu_n'$ gives $u'=v$
in distributions. Testing the Galerkin equations against a fixed
member of the dense union, multiplied by a scalar time test function,
then passing to the limit gives $v'=f-Au\in L^2V_0'$.
Density extends the identity to all tests in $V_0$.
In particular $u\in W^{1,\infty}(0,T;H_0)$ and
$v\in W^{1,2}(0,T;V_0')\cap L^\infty H_0$.
Integrating the tested equations up to the initial endpoint gives
the traces $u(0)=u_0$, $v(0)=v_0$.

We spell out the time representatives used here. The first Sobolev
inclusion makes $u$ strongly Lipschitz into $H_0$; its uniform $V_0$
bound then makes it weakly continuous into $V_0$. Indeed approximate
a functional in $V_0'$ by functionals from $H_0$, which are dense
because the embedding is dense and injective, and use the uniform
$V_0$ bound. Similarly the equation makes $(v(t),w)_H$ continuous
for each $w\in V_0$, and density plus the uniform $H_0$ bound makes
$v$ weakly continuous into $H_0$. These representatives are defined
at every time. Weak lower semicontinuity, applied to the Galerkin
energy estimate (first a.e., then at every time by weak continuity),
gives
\[
 E(t)\le e^{ct}\left(E(0)+\int_0^t\|f(r)\|_H^2\,dr\right),
 \quad E(t)=\|v(t)\|_H^2+\widetilde a_t(u(t),u(t)).
\]
For this limit it suffices to use the integrated weak lower
semicontinuity on time intervals; the continuous scalar bound above
then gives the a.e. pointwise inequality by differentiation.
At the initial time its right side tends to $E(0)$. Form Lipschitz
continuity and weak continuity therefore imply convergence of the
norm in the Hilbert space $V_0\oplus H_0$ equipped with the energy
inner product at time zero. Weak convergence plus convergence of
that squared norm is strong convergence. Thus $u,v$ have the required
strong energy trace at the starting time. The same construction and
estimate apply with any starting time and also with time reversed.

\emph{Uniqueness already in the weak energy class.}
Subtract two solutions with the same weak traces and data. Denote
their difference again by $u$, so $u''+Au=0$ and both initial traces
are zero. Set $U(t)=\int_0^tu(r)\,dr$ as a $V$-valued integral.
For fixed $s$, test on $[0,s]$ with
$w(t)=U(s)-U(t)=\int_t^su(r)\,dr$. It is Lipschitz into $V$,
$w'=-u$, and $w(s)=0$. The equation and Hilbert product rule give
\[
 \|u(s)\|_H^2+\widetilde a_0(U(s),U(s))
   =2\kappa\int_0^s(u,w)_H\,dt
                       -\int_0^s\partial_ta_t(w,w)\,dt.
\]
Integration by parts is legitimate since $u'\in W^{1,2}V'$
and $w\in W^{1,\infty}V$, while
$\int(u',u)_H=\|u(s)\|_H^2/2$ uses $u\in W^{1,\infty}H$.
The scalar form product rule follows by difference quotients and
the bounded derivative estimate. Since
$\|w(t)\|_V^2\le2\|U(s)\|_V^2+2\|U(t)\|_V^2$, the right side
in absolute value is at most
\[
 \kappa\int_0^s\|u(t)\|_H^2dt
   +c_1s\|U(s)\|_V^2+c_1\int_0^s\|U(t)\|_V^2dt.
\]
Choose $s$ in a fixed sufficiently short interval so $c_1s\le\alpha/2$.
Absorbing that term and applying Gronwall to
$\|u(s)\|_H^2+(\alpha/2)\|U(s)\|_V^2$ gives zero throughout
the interval. Its velocity is zero there as a distribution and
therefore in its weakly continuous representative. Restart the
same argument on consecutive intervals to obtain uniqueness on
$[0,T]$. No strong endpoint energy identity was assumed in this step.

\emph{Recovery of strong continuity everywhere.}
At an arbitrary time $t_0$ the weak representatives give data
$u(t_0)\in V$, $v(t_0)\in H$. Perform the preceding existence
construction forward from that time with these data. It has strong
energy continuity at its starting time. Weak-class uniqueness makes
it equal to the original solution on the forward interval. Apply
the same reasoning backward. Thus the original state is strongly
continuous at every $t_0$ in $V\oplus H$, including the one-sided
endpoint assertions. The uniform estimate follows from the displayed
Gronwall bound and equivalence of the shifted form norm with the
$V$ norm. Complementary tests give the full-space equation as
explained in the first paragraph; uniqueness in the full space
also follows from the same integrated-test argument.

\emph{Exact energy identity.}
For completeness we justify the time smoothing needed to test with
the velocity, which need not itself lie in $V$. On an interior time
interval choose an even smooth mollifier $\rho_\varepsilon$ and set
$u_\varepsilon=\rho_\varepsilon*u$. Let
$c_\varepsilon(t)=A(t)u_\varepsilon(t)-(Au)_\varepsilon(t)$.
Then
\[
 \|c_\varepsilon\|_{L^2V'}=o(\varepsilon),\qquad
 \|u_\varepsilon'\|_{L^2V}=O(\varepsilon^{-1}).
\]
Here is the required commutator estimate. Split its defining integral
into
\[
 \int\rho_\varepsilon(s)(A(t)-A(t-s))(u(t-s)-u(t))\,ds
 +\left[\int\rho_\varepsilon(s)(A(t)-A(t-s))\,ds\right]u(t).
\]
Operator Lipschitz continuity bounds the first term by $o(\varepsilon)$
in $L^2$, by continuity of translations of the $L^2(V)$ function $u$.
For a fixed vector $v$, the $V'$-valued function $A(t)v$ is Lipschitz
and has a strong a.e. derivative in $L^\infty(V')$. To verify the
latter without a general vector-measure theorem, its range spans a
separable Hilbert space; differentiate its countably many orthonormal
coordinates, use the Lipschitz bound and Fatou to build their
square-summable derivative, and integrate coordinatewise. This gives
the strong vector fundamental theorem. Evenness of the mollifier,
$\int s\rho_\varepsilon(s)ds=0$, and translation continuity of this
derivative in $L^2$ therefore make the bracket divided by
$\varepsilon$ tend strongly to zero on each fixed $v$.
Its operator norm divided by $\varepsilon$ is uniformly bounded.
Approximate $u(t)$ in $L^2(V)$ by finite-valued step functions to
obtain the same $o(\varepsilon)$ bound for the second term. This
proves the commutator estimate. The derivative estimate follows
directly from $\|\rho_\varepsilon'\|_1=O(\varepsilon^{-1})$.

Now the smoothed equation is
$u_\varepsilon''+A(t)u_\varepsilon=f_\varepsilon+c_\varepsilon$.
Testing with $u_\varepsilon'\in V$ and differentiating the scalar
form is legitimate. The commutator contribution tends to zero in
integral by the two estimates. Moreover $u_\varepsilon\to u$ in
$L^2V$ and $u_\varepsilon'\to u'$ in $L^2H$; the bounded form
derivative and Cauchy--Schwarz pass all remaining terms to the limit.
Thus in distributions
$\frac d{dt}(\|u'\|_H^2+a_t(u,u))
=2(f,u')_H+\partial_ta_t(u,u)$.
The right side is integrable. The energy is already continuous by
the preceding step, so integrating its distributional derivative
gives the claimed identity for every pair $s,t$, and by continuity
also at the endpoints. This completes all assertions.
\end{proof}

% END INPUT


% BEGIN parts/part3-models.tex
\chapter{Moment problems, unitary models, and monotone evolution}
This chapter separates abstract functional analysis from further
boundary-value and dynamical applications. It supplies the current
subcases of mixed source families rather than dropping them merely
because their principal curricular home is later. See
\cite{B01,B02,B03,B06,B08,B09}.

\begin{theorem}[Hamburger moment problem]\label{fa2:p3:hamburger-moment}
A real sequence $(m_n)_{n\geq0}$ is the sequence of moments of a finite
positive Borel measure on $\mathbb R$, all absolute moments finite,
iff every Hankel form
$\sum_{j,k=0}^Nm_{j+k}\overline{c_j}c_k$ is nonnegative.
It is equivalent to the existence of a self-adjoint $A$ and a vector
$e\in\bigcap_nD(A^n)$ with $m_n=\faip{e}{A^ne}$.
\end{theorem}
\begin{proof}
A representing measure makes the Hankel form the integral of
$|\sum c_jt^j|^2$, proving necessity. Conversely define on complex
polynomials
$\faip{p}{q}=\sum_{j,k}m_{j+k}\overline{p_j}q_k$.
The null space is stable under multiplication by $t$: Schwarz gives
$\faip{p}{r}=0$ for null $p$ and all $r$, hence
$\|tp\|^2=\faip{p}{t^2p}=0$.
Quotient and complete to a Hilbert space $H$. Multiplication by $t$
on polynomial classes is therefore a densely defined symmetric
operator $S$. Coefficient conjugation is an antiunitary involution
$J$ with $JS=SJ$, and the adjoint definition gives
$JS^*=S^*J$ with conjugated scalars. Thus it carries
$\ker(S^*-i)$ antiunitarily onto $\ker(S^*+i)$.
The closure of $S$ has equal deficiency indices and hence has a
self-adjoint extension $A$. With $e=[1]$, repeated multiplication
gives $A^ne=[t^n]$, so $e$ is in all power domains and the moments
have the stated values. The scalar PVM measure
$\mu(B)=\faip{e}{E_A(B)e}$ is finite, and
$\int|t|^{2n}d\mu=\|A^ne\|^2<\infty$.
Cauchy--Schwarz supplies every intermediate absolute moment;
the spectral integral gives $\int t^nd\mu=m_n$.
This also proves the equivalence with the operator formulation.
\end{proof}

\begin{proposition}[Fourier transform as a bounded spectral example]\label{fa2:p3:fourier-four-points}
The unitary Fourier transform on $L^2(\mathbb R)$ with phase
$e^{-2\pi ix\xi}$ has spectrum exactly $\{1,-1,i,-i\}$ and the
orthogonal decomposition into these four eigenspaces. A bounded
diagonal operator on $\ell^p$, $1\leq p\leq\infty$, has spectrum
the closure of its diagonal entries.
\end{proposition}
\begin{proof}
Fourier inversion gives $\mathcal F^2f(x)=f(-x)$, first on Schwartz
functions and then by density, so $\mathcal F^4=I$.
For each fourth root $\lambda$, the polynomial
$P_\lambda=\frac14\sum_{k=0}^3\lambda^{-k}\mathcal F^k$
is its spectral projection; the geometric-sum identities give
$\sum P_\lambda=I$ and pairwise orthogonality. All four are
nonzero: put $g=e^{-\pi x^2}$. The functions
$g$, $xg$, $(x^2-1/(4\pi))g$ and
$(x^3-3x/(4\pi))g$ have respective Fourier eigenvalues
$1,-i,-1,i$. This follows from the Gaussian transform and
$\mathcal F(xf)=i(\mathcal Ff)'/(2\pi)$, with derivatives computed
directly. Thus all four roots occur and the polynomial equation
excludes every other spectral value.
For a diagonal sequence $(a_n)$, if $\lambda$ stays a positive
distance from its closure, multiplication by $(\lambda-a_n)^{-1}$
is a bounded inverse on every stated $\ell^p$ space. If it belongs
to the closure, unit coordinate vectors along a converging
subsequence give approximate eigenvectors (or an exact eigenvector
if the value occurs). Hence no bounded inverse exists there.
\end{proof}

\begin{example}[A genuinely indeterminate Stieltjes sequence]\label{fa2:p3:moment-nonuniqueness}
For $-1\leq a\leq1$, the distinct positive measures
\[
 d\mu_a(x)=\frac{e^{-(\log x)^2}}{\sqrt\pi x}
            (1+a\sin(2\pi\log x))\,dx,\qquad x>0,
\]
have the same moments $m_n=e^{n^2/4}$ for every nonnegative integer
$n$. Thus positivity of both Hankel forms does not imply uniqueness.
\end{example}
\begin{proof}
The density is nonnegative and integrable. Substitute $y=\log x$.
The Gaussian integral gives
$\pi^{-1/2}\int e^{ny-y^2}dy=e^{n^2/4}$.
The perturbation integral is the imaginary part of
$\pi^{-1/2}\int e^{(n+2\pi i)y-y^2}dy
=e^{(n+2\pi i)^2/4}$, namely
$e^{n^2/4-\pi^2}\sin(n\pi)=0$.
The complex Gaussian identity follows by completing the square
first for real parameters and then using the identity theorem for
its entire parameter integral, justified by Gaussian domination on
compact parameter sets. Every absolute moment is finite by the
same Gaussian bound. Different $a$ give different densities on a
set of positive measure. This proves genuine Stieltjes and hence
Hamburger nonuniqueness.
\end{proof}

\begin{theorem}[A full one-dimensional compact-resolvent realization]\label{fa2:p3:sturm-liouville}
Let $a<b$, $p\in C^1[a,b]$ real with $p\geq p_0>0$, and
$q\in C[a,b]$ real. On $H=L^2(a,b)$ the Dirichlet realization
\[
 Lu=-(pu')'+qu,\qquad D(L)=H^2(a,b)\cap H_0^1(a,b)
\]
is self-adjoint and bounded below, and has compact resolvent.
Its eigenfunctions form an orthonormal basis, its eigenvalues have
finite multiplicity and tend to $+\infty$, and its heat semigroup
is compact at every positive time. The periodic realization has
domain
$\{u\in H^2:u(a)=u(b),\ p(a)u'(a)=p(b)u'(b)\}$ and has the
same conclusions. These are operator domains and boundary conditions,
not merely formal differential expressions.
\end{theorem}
\begin{proof}
For Dirichlet take $V=H_0^1(a,b)$ and the closed form
$l(v,u)=\int_a^bp\overline{v'}u'+q\bar v u$.
Adding $c\faip{v}{u}$ with $c>1+\|q\|_\infty$ makes it coercive
and gives a norm equivalent to the $H^1$ norm. The form construction
therefore gives a self-adjoint operator bounded below.
Its exact domain is as displayed. If $u\in V$ represents a form
functional by $f\in L^2$, testing compactly supported smooth $v$
gives $-(pu')'=f-qu\in L^2$. Thus $pu'\in H^1$;
$p\in C^1$ bounded away from zero implies $u'\in H^1$, so
$u\in H^2\cap H_0^1$. Conversely integration by parts shows every
such $u$ belongs to the operator domain with value $Lu$.
The inverse of $L+c$ maps $L^2$ boundedly into $V$ by the form
estimate. The embedding $H^1(a,b)\hookrightarrow L^2(a,b)$ is compact:
a bounded $H^1$ family has uniformly bounded continuous
representatives and $|u(x)-u(y)|\leq\|u'\|_2|x-y|^{1/2}$;
Arzela--Ascoli gives a uniformly convergent subsequence, hence an
$L^2$ one. Thus $(L+c)^{-1}$ is compact positive self-adjoint,
injective with dense range. Its compact spectral expansion gives a
complete eigenbasis, with positive eigenvalues tending to zero;
inverting and subtracting $c$ yields the claimed $L$ eigenvalues.
The heat multipliers $e^{-t\lambda_n}$ tend to zero for $t>0$,
so finite spectral truncations converge in norm, proving compactness.
For periodic conditions use $V=\{u\in H^1:u(a)=u(b)\}$.
The same distribution argument gives $H^2$ regularity; the boundary
term against arbitrary common endpoint values forces precisely
$p(a)u'(a)=p(b)u'(b)$. All other form and compactness arguments are
unchanged. This proves a full standard one-dimensional realization;
higher-dimensional boundary regularity is a separate application.
\end{proof}

\begin{proposition}[Periodic convolution as a compact spectral model]\label{fa2:p3:periodic-convolution}
For $k\in L^1(\mathbb T)$ and normalized Haar measure,
$Kf=k*f$ is bounded, compact, normal, and diagonal in the Fourier basis
$e_n(x)=e^{inx}$, with eigenvalues $\widehat k(n)$.
It is compact, its spectrum is the closure of these eigenvalues,
and it is self-adjoint iff they are real, positive iff they are
nonnegative. It is Hilbert--Schmidt exactly when $k\in L^2$,
with $\|K\|_2=\|k\|_2$. Trace class requires exactly the extra
condition $\sum_n|\widehat k(n)|<\infty$, even for continuous $k$.
\end{proposition}
\begin{proof}
Young's inequality bounds its operator norm by $\|k\|_1$.
Approximate $k$ in $L^1$ by trigonometric polynomials, whose
convolutions have finite rank. Their operators converge in norm,
proving compactness without an unjustified $L^2$ assumption.
Fubini gives
$Ke_n(x)=e_n(x)\int k(t)e^{-int}dt$.
Fourier completeness therefore identifies $K$ with that diagonal
operator. Its eigenvalues tend to zero by the preceding polynomial
approximation, and its singular values are their absolute values.
For $k\in L^2$ the product kernel has squared integral $\|k\|_2^2$,
giving the Hilbert--Schmidt norm. Conversely a square-summable
coefficient sequence has an $L^2$ function with those Fourier
coefficients; uniqueness of Fourier coefficients in $L^1$ makes
it equal to $k$. This proves the stated equivalence. The compact
spectral, positivity and trace-class assertions follow from the
diagonal model.
This computation uses spectral Fourier coefficients rather than
the pointwise sign of the kernel.
\end{proof}

\begin{theorem}[Stieltjes existence]\label{fa2:p3:stieltjes-moment}
The sequence has a representing positive measure on $[0,\infty)$ iff
both $(m_{j+k})$ and $(m_{j+k+1})$ define nonnegative finite Hankel
forms. The representing operator can be chosen nonnegative.
\end{theorem}
\begin{proof}
The two forms are the integrals of $|p|^2$ and $t|p|^2$ for a
representing measure, proving necessity. In the preceding polynomial
construction the second form says
$\faip{p}{Sp}\geq0$. Friedrichs gives a nonnegative self-adjoint
extension $A$ of $S$. Its spectral measure is supported on
$[0,\infty)$, and its vector $e$ has the required moments by the
same argument. This proves sufficiency without assuming a bounded
support or a bounded multiplication operator.
\end{proof}

\begin{proposition}[The Jacobi realization]\label{fa2:p3:jacobi-moment}
If all nonzero polynomial Hankel forms are strictly positive,
Gram--Schmidt gives real orthonormal polynomials $p_n$ with positive
leading coefficients and a three-term recurrence
\[
 tp_n=a_{n+1}p_{n+1}+b_np_n+a_np_{n-1},\qquad
 a_{n+1}>0,\ b_n\in\mathbb R,\ a_0p_{-1}=0.
\]
Under $p_n\leftrightarrow e_n$ the minimal multiplication operator
is the symmetric Jacobi matrix on finitely supported sequences.
Its deficiency indices are either $(0,0)$ or $(1,1)$.
If there is a nonzero null polynomial instead, the polynomial
Hilbert space is finite dimensional and the moment measure is a
uniquely determined finite atomic measure.
\end{proposition}
\begin{proof}
All moments are real, so Gram--Schmidt can use real coefficients.
Multiplication by $t$ has degree one higher. For $j<n-1$,
$\faip{p_j}{tp_n}=\faip{tp_j}{p_n}=0$ by degree orthogonality.
Symmetry gives equal neighboring coefficients, and comparison of
leading terms gives $a_{n+1}=\kappa_n/\kappa_{n+1}>0$.
The diagonal coefficient is real by symmetry. Completeness in the
polynomial Hilbert space gives the stated unitary identification.
The adjoint deficiency equation for a nonreal $z$ is the recurrence
$a_{n+1}x_{n+1}+b_nx_n+a_nx_{n-1}=zx_n$.
Since $a_{n+1}>0$ and the boundary has $a_0x_{-1}=0$, the full
sequence is determined by $x_0$. The square-summable solution
space therefore has dimension at most one. Conjugation interchanges
the two deficiencies, proving the alternatives.
If a null polynomial $p$ has degree $N$, all its multiples are
null by stability under multiplication. Polynomial division by $p$
therefore reduces every class to degree less than $N$, so the
Hilbert quotient is finite dimensional. Its symmetric multiplication
operator is everywhere defined and self-adjoint; hence it is
essentially self-adjoint and the determinacy theorem gives uniqueness.
Its finite spectral measure at $[1]$ is the required finite atomic
measure. Zero weights or repeated roots are automatically removed
by the quotient and the spectral projections.
\end{proof}

\begin{theorem}[Hamburger determinacy and essential self-adjointness]\label{fa2:p3:moment-determinacy}
For the polynomial Hilbert space and symmetric multiplication operator
$S$ just constructed, the representing Hamburger measure is unique
iff $S$ is essentially self-adjoint. In particular bounded even
moments, or more generally a common positive exponential moment
determined by the moment series, give uniqueness.
\end{theorem}
\begin{proof}
If $S$ is not essentially self-adjoint, its equal nonzero deficiency
indices permit two distinct self-adjoint extensions on the polynomial
Hilbert space. Their scalar measures at $e$ cannot be equal.
Indeed if equal to $\mu$, for polynomials $p,q$ their resolvent
coefficients would both be
$\int\overline{p(t)}q(t)/(z-t)\,d\mu(t)$.
The polynomial vectors are dense, so equality of these coefficients
would give equality of the bounded resolvents, hence of the two
operators. This contradiction gives distinct representing measures.
Conversely suppose $\overline S$ is self-adjoint and $\mu$ is any
representing measure. Sending polynomial classes to their functions
is an isometry into $L^2(\mu)$, with closed image $M$.
The full self-adjoint multiplication operator $M_t$ extends $S$;
by closedness it extends $\overline S$ on this embedded $M$.
For nonreal $z$ and $x\in M$, solve
$(z-\overline S)y=x$ in $M$. This also solves
$(z-M_t)y=x$, so uniqueness of the full resolvent puts
$(z-M_t)^{-1}M\subset M$. Both signs make $M$ reducing for those
resolvents. Thus the scalar Cauchy transform at $e=1$ is the one
fixed by $\overline S$, independent of $\mu$.
Finite measures are determined by their Cauchy transforms. For the
normalized kernel $P_\varepsilon(x)=\varepsilon/(\pi(x^2+\varepsilon^2))$,
the imaginary part at $a+i\varepsilon$ is $-\pi(P_\varepsilon*\mu)(a)$.
Integrating $P_\varepsilon*\mu$ against a compactly supported continuous
test and letting $\varepsilon\downarrow0$ recovers its integral against $\mu$.
For the last sufficient condition, if $m_{2n}\leq C$ then monotone
convergence of the positive series gives
$\int\cosh(rt)\,d\mu\leq C\cosh r$ for every $r>0$.
Thus the Fourier transform has its absolutely convergent moment
Taylor series, identical for every representing measure.
More generally the same argument works if the even moment series
converges for some $r>0$; analytic continuation of the Fourier
transform on the resulting strip and Fourier uniqueness prove
determinacy. A full Carleman quasi-analytic extension is a distinct
enhancement; it is not stated in the Lax locator verified here and
is not silently counted as proved by this exponential-moment argument.
\end{proof}

\begin{theorem}[Position, momentum, and uncertainty]\label{fa2:p3:position-momentum}
On $L^2(\mathbb R)$ set
$Qf(x)=xf(x)$ with $D(Q)=\{f:xf\in L^2\}$ and
$Pf=-if'$ with $D(P)=H^1(\mathbb R)$, using distributional derivatives.
Both are self-adjoint and the Schwartz space is a core for each and
a common invariant core. For the Fourier convention
$\widehat f(\xi)=\int e^{-2\pi ix\xi}f(x)\,dx$,
$\mathcal FP\mathcal F^{-1}=M_{2\pi\xi}$.
Their unitary groups are
$U(a)f(x)=e^{iax}f(x)$ and $V(b)f(x)=f(x+b)$, with
$U(a)V(b)=e^{-iab}V(b)U(a)$.
For a unit $f\in D(Q)\cap D(P)$,
\[
 \|(Q-q)f\|\,\|(P-p)f\|\geq\tfrac12,
 \quad q=\faip{f}{Qf},\ p=\faip{f}{Pf}.
\]
\end{theorem}
\begin{proof}
Real multiplication with its maximal square-integrability domain is
self-adjoint by the unbounded calculus. Plancherel and integration
by parts on Schwartz functions give
$\widehat{-if'}=2\pi\xi\widehat f$.
The distributional testing identity shows that the maximal domain
of this Fourier multiplication is exactly $H^1$, so it proves
self-adjointness of $P$ and the stated diagonalization.
Truncate and mollify in the weighted $L^2$ norms to approximate by
Schwartz functions in both graph norms. For the joint graph norm
one first uses smooth spatial cutoffs: the extra derivative term is
bounded by $C\|f\|_2/R$ and the weighted tails tend to zero.
Mollification then preserves convergence of derivatives, and
$Q(f*\rho_\varepsilon)=(Qf)*\rho_\varepsilon
                 +f*(x\rho_\varepsilon)$,
whose extra term tends to zero. This proves the joint core assertion.
The unitary formulas follow directly by Fourier inversion and real
multiplication calculus; substitution proves their Weyl relation.
On Schwartz functions integration by parts gives
$\faip{Qf}{Pf}-\faip{Pf}{Qf}=i\|f\|^2$.
Joint graph approximation extends this weak identity to the common
domain. Subtracting the real means does not change its imaginary
part, so Cauchy--Schwarz gives
$\|(Q-q)f\|\|(P-p)f\|\geq
|\operatorname{Im}\faip{(Q-q)f}{(P-p)f}|=1/2$.
This argument uses the weak form on the common domain, not an
undefined product $QP$ on every such vector.
\end{proof}

\begin{proposition}[No bounded canonical commutator]\label{fa2:p3:no-bounded-ccr}
Bounded operators on a nonzero Banach space cannot satisfy $AB-BA=I$.
\end{proposition}
\begin{proof}
Induction gives $AB^n-B^nA=nB^{n-1}$, so
$n\|B^{n-1}\|\leq2\|A\|\|B^n\|\leq2\|A\|\|B\|\|B^{n-1}\|$.
For $n>2\|A\|\|B\|$ this forces $B^{n-1}=0$.
The same identity then forces successively $B^{n-2}=0,\ldots,I=0$,
a contradiction. The zero-space exception is explicit.
\end{proof}

\begin{remark}[The uniqueness theorem for Weyl representations]
The Stone--von Neumann theorem concerns strongly continuous unitary
groups satisfying the Weyl relation with a fixed nonzero central
character. Irreducibility yields the scalar Schrödinger representation;
without it there is an arbitrary multiplicity Hilbert space.
Formal commutator identities on an unspecified common domain are
insufficient. The full Gaussian-corner proof, including arbitrary
multiplicity, is given in \ref{fa2:p3:stone-von-neumann}.
\end{remark}

\begin{proposition}[Koopman and abstract waves]\label{fa2:p3:kooper-wave}
If a measure-preserving invertible flow $\Phi_t$ has strongly
continuous pullbacks on $L^2$, then $U(t)f=f\circ\Phi_t$ is a
unitary group and Stone supplies its self-adjoint generator.
Strong continuity is an assumption to verify, not a consequence of
the group law alone. It holds for every continuous measure-preserving
flow on a compact Hausdorff space with finite regular measure,
including the smooth compact-manifold flows in the source.
On a probability space its fixed vectors are only constants iff
the flow is ergodic, and then its time averages tend in $L^2$ to
the integral of the initial function.
If $L=L^*\geq0$, the abstract wave equation has
the energy solution
\[
 u(t)=\cos(t\sqrt L)u_0+
      \frac{\sin(t\sqrt L)}{\sqrt L}v_0,
 \quad u_0\in D(\sqrt L),\ v_0\in H,
\]
where the second multiplier has value $t$ at zero. It belongs to
$C(\mathbb R;D(\sqrt L))\cap C^1(\mathbb R;H)$ and preserves
$\|\sqrt L u(t)\|^2+\|u'(t)\|^2$.
For $L\geq cI>0$ it is a unitary group on the Hilbert energy space
$D(\sqrt L)\oplus H$ with that energy norm.
\end{proposition}
\begin{proof}
Measure preservation gives equality of $L^2$ norms; invertibility
gives surjectivity and the flow law gives the group. The stated
continuity is exactly what Stone requires. For the compact continuous
flow, continuity on the compact product of a small time interval
and the space makes $f\circ\Phi_t\to f$ uniformly for continuous
$f$. Density of continuous functions in $L^2$ and the common
unitary bound extend convergence to every $L^2$ vector.
For ergodicity, a nonconstant fixed real function has a nontrivial
level set invariant modulo null sets under every time map; conversely
the indicator of an invariant set is fixed. Applying this also to
real and imaginary parts proves the fixed-space characterization.
The mean ergodic projection is therefore the projection onto the
constant function, namely $f\mapsto\int f$ for a probability measure.
For the wave equation, all multipliers in the formula are bounded
on the stated data: $|\sin(t\sqrt\lambda)|/\sqrt\lambda\leq|t|$.
Differentiate under the PVM using dominated convergence on bounded
time intervals. The pair
$(\sqrt\lambda u(t),u'(t))$ is the rotation of
$(\sqrt\lambda u_0,v_0)$, so its squared norm is preserved for each
spectral value; integration gives the energy identity and graph
continuity. The equation holds in $D(\sqrt L)'$, and classically in
$H$ for $u_0\in D(L),v_0\in D(\sqrt L)$.
Density and the energy identity give uniqueness of the energy
solution. When $L\geq cI$, the energy is a complete norm equivalent
to the product graph norm, and the inverse time gives a unitary group.
At zero spectral modes $u(t)=u_0+tv_0$; these are retained rather
than falsely treating the homogeneous energy seminorm as a norm.
\end{proof}

\begin{proposition}[The completed homogeneous wave energy space]\label{fa2:p3:homogeneous-wave-group}
For $L=L^*\geq0$ put $H_0=(\ker L)^\perp$ and complete
$D(\sqrt L)/\ker L$ in the norm $\|\sqrt L u\|$.
On the product of this displacement completion with $H_0$ the
abstract wave evolution is a strongly continuous unitary group,
without a positive spectral gap. Under the isometry
$(u,v)\mapsto(\sqrt L u,v)$ onto $H_0\oplus H_0$, it is
\[
 \begin{pmatrix}\cos(t\sqrt L)&\sin(t\sqrt L)\\
 -\sin(t\sqrt L)&\cos(t\sqrt L)\end{pmatrix}.
\]
Zero-frequency velocities may be included as a separate constant
component; in the original displacement space they give the
explicit drift $u_{\ker}(t)=u_{\ker}(0)+tv_{\ker}(0)$.
\end{proposition}
\begin{proof}
The map $\sqrt L$ has dense range in $H_0$ by the adjoint
kernel-range identity, so it identifies the displacement completion
isometrically with $H_0$. The displayed PVM matrix is unitary at
each spectral value, has the addition group law, and is strongly
continuous by bounded convergence. Hence it defines a unitary
$C_0$ group on the completed energy space. On the dense original
graph-domain data it agrees with the previous wave formula after
applying $\sqrt L$, so it is exactly the energy completion of that
evolution. Kernel displacements have zero homogeneous norm and are
quotiented out; keeping kernel velocities separately gives the
stated original zero-mode solution, rather than pretending a
degenerate energy seminorm was already a Hilbert norm.
\end{proof}

\begin{theorem}[Heat and Poisson semigroups]\label{fa2:p3:heat-poisson}
On $L^2(\mathbb R^d)$ the multipliers
$e^{-4\pi^2t|\xi|^2}$ and $e^{-2\pi t|\xi|}$ give contraction
$C_0$ semigroups with generators respectively
$\Delta$ on $H^2(\mathbb R^d)$ and
$-\sqrt{-\Delta}$ on $H^1(\mathbb R^d)$.
They smooth instantly into every power domain.
The heat kernels
$g_t(x)=(4\pi t)^{-d/2}e^{-|x|^2/(4t)}$ give contraction
$C_0$ semigroups by convolution also on $L^p$, $1\leq p<\infty$.
Their generator domain there is exactly
$\{f\in L^p:\Delta f\in L^p\text{ distributionally}\}$.
\end{theorem}
\begin{proof}
Plancherel and the real measurable calculus identify the two maximal
domains by square integrability of $|\xi|^2\widehat f$ and
$|\xi|\widehat f$, exactly the displayed Sobolev spaces.
The semigroup law is multiplication of symbols; dominated convergence
gives strong continuity and the precise derivative domains.
The smoothing bounds follow from
$\sup_{r\geq0}r^me^{-tr}<\infty$.
The Gaussian Fourier transform, obtained by the one-dimensional
Gaussian integral and tensor products, identifies the heat kernel.
It has integral one, is nonnegative, and its mass outside each fixed
ball tends to zero. Young's inequality gives contraction; translation
continuity and the approximate-identity estimate give $C_0$ on
$L^p$, $p<\infty$. Its generator agrees with $\Delta$ on test
functions by Taylor expansion and Gaussian moments.
If $f$ is in its generator domain, testing against compact smooth
functions and differentiating the dual heat flow shows $Af=\Delta f$
distributionally. Conversely if $f,\Delta f\in L^p$, choose
$u=R(\lambda,A)(\lambda f-\Delta f)$ using the Laplace resolvent.
The preceding distribution identity gives
$(\lambda-\Delta)(u-f)=0$. Every $L^p$ function defines a tempered
distribution. Fourier transformation and multiplication by
$(\lambda+4\pi^2|\xi|^2)^{-1}$, which preserves Schwartz tests,
show $u=f$. Thus $f\in D(A)$, proving the exact maximal domain.
Strong continuity on $L^\infty$ is not asserted.
\end{proof}

\begin{proposition}[Legendre as a complete classical example]\label{fa2:p3:legendre}
The polynomials
$P_n(x)=(2^nn!)^{-1}(d/dx)^n(x^2-1)^n$ are orthogonal on $[-1,1]$,
have squared $L^2$ norm $2/(2n+1)$, and
$\sqrt{(2n+1)/2}\,P_n$ form an orthonormal basis of $L^2(-1,1)$.
\end{proposition}
\begin{proof}
The first $n-1$ derivatives of $(x^2-1)^n$ vanish at both endpoints.
Integrating by parts $n$ times gives $\int P_nq=0$ for every
polynomial $q$ of degree less than $n$. The leading coefficient is
$c_n=(2n)!/(2^n(n!)^2)$, so
$\int P_n^2=c_n\int P_n x^n$.
The same integrations by parts give
$\int_{-1}^1P_nx^n=2^{-n}\int_{-1}^1(1-x^2)^n\,dx$.
If $I_n$ denotes the latter integral, integration of
$(x(1-x^2)^n)'$ gives
$I_n=2n(I_{n-1}-I_n)$, hence
$I_n=2^{2n+1}(n!)^2/(2n+1)!$ from $I_0=2$.
Substitution yields $2/(2n+1)$.
Their nonzero leading coefficients show that their finite spans
are exactly all polynomials. Weierstrass approximation makes
polynomials uniformly dense in continuous functions, and regular
measure approximation makes continuous functions dense in $L^2$.
Thus the normalized orthogonal family is complete. This supplies one
full classical-polynomial example; it does not claim the whole
orthogonal-polynomial curriculum of a later course.
\end{proof}

\begin{theorem}[Nonlinear monotone resolvents and evolution]\label{fa2:p3:nonlinear-monotone}
Let $A\subset H\times H$ be a possibly multivalued monotone graph
on a real Hilbert space, and assume $I+\lambda A$ is onto for every
$\lambda>0$. Its resolvent $J_\lambda=(I+\lambda A)^{-1}$ is
single valued, firmly nonexpansive, and
$A_\lambda=(I-J_\lambda)/\lambda$ is everywhere Lipschitz and
monotone, with $(J_\lambda x,A_\lambda x)\in A$.
For every $x\in\overline{D(A)}$ there is a contraction semigroup
trajectory $u(t)=S(t)x$ obtained from $u_\lambda'=-A_\lambda u_\lambda$.
For $x\in D(A)$ it is Lipschitz, satisfies
$u'(t)+Au(t)\ni0$ almost everywhere, and is unique in that class.
With $J_0=I$ on $\overline{D(A)}$, the implicit products for $n\ge1$ satisfy
$S(t)x=\lim_{n\to\infty}J_{t/n}^{\,n}x$, uniformly on bounded time intervals.
\end{theorem}
\begin{proof}
If $x=u+\lambda a$, $y=v+\lambda b$ with $(u,a),(v,b)\in A$,
monotonicity gives
$\|u-v\|^2\leq\faip{u-v}{x-y}$.
This proves uniqueness and firm nonexpansion of $J_\lambda$.
The same identity gives Lipschitz constant at most $1/\lambda$
and monotonicity for $A_\lambda$ (the weaker $2/\lambda$ bound
would suffice). The graph assumption also implies maximality:
if $(u,a)$ is monotonically related to all graph points, solve
$u+\lambda a=v+\lambda b$ with $(v,b)\in A$; then
$0\leq\faip{u-v}{a-b}=-\|u-v\|^2/\lambda$, so $(u,a)=(v,b)$.
The Lipschitz ODE has a global solution; monotonicity makes its
flow contractive by differentiating the squared distance.
For $x\in D(A)$ choose $a_0\in Ax$. Monotonicity against $(x,a_0)$
gives $\|A_\lambda x\|\leq\|a_0\|$. Time translation and flow
contraction give
$\|u_\lambda'(t)\|\leq\|A_\lambda x\|\leq C=\|a_0\|$.
Write $v_\lambda=J_\lambda u_\lambda$,
$a_\lambda=A_\lambda u_\lambda$. Then
$u_\lambda=v_\lambda+\lambda a_\lambda$ and
$(v_\lambda,a_\lambda)\in A$.
Comparing two parameters and using monotonicity yields
$\frac12(d/dt)\|u_\lambda-u_\mu\|^2
\leq2(\lambda+\mu)C^2$.
Thus $u_\lambda$ converges uniformly on bounded intervals to a
Lipschitz $u$, with derivative $-a$ obtained as a weak $L^2$ limit
of $a_\lambda$. To verify the graph relation without an uncountable
exceptional-set argument, put
$v=J_1(u+a)$ and $b=u+a-v$. These are strongly measurable and
$(v,b)\in A$. Integrated monotonicity gives
$\int\faip{a_\lambda-b}{v_\lambda-v}\geq0$.
Pass to the weak--strong limit to obtain
$0\leq\int\faip{a-b}{u-v}=-\int\|a-b\|^2$.
Hence $u=v$, $a=b$ almost everywhere and $a\in Au$.
Any two such solutions contract by the same squared-distance
calculation. Uniqueness gives the semigroup law and contraction;
extension from $D(A)$ to its closure uses uniform contraction bounds.
For implicit steps $x_k=J_hx_{k-1}$ with $h>0$, let
$a_k=(x_{k-1}-x_k)/h\in Ax_k$. Monotonicity successively gives
$\|a_k\|\leq C$ for initial $x\in D(A)$ and the chosen $a_0$.
Interpolate linearly with derivative $-a_k$, and let $\bar x_h=x_k$
on each step. Then $\|x_h-\bar x_h\|\leq hC$.
Comparison with $u$ gives
$\frac12(d/dt)\|x_h-u\|^2\leq2hC^2$.
At $t=nh>0$ this bounds the error by $2C\sqrt{th}$.
Density and contractions extend convergence to $\overline{D(A)}$,
uniformly for bounded positive times. At $t=0$, the semigroup and
every implicit product fix $x$, so the same convergence holds.
\end{proof}

\begin{proposition}[Convex subgradients supply the hypothesis]\label{fa2:p3:subgradient}
For a proper norm-lower-semicontinuous convex
$\Phi:H\to(-\infty,+\infty]$, its subgradient
$a\in\partial\Phi(u)$ iff
$\Phi(v)\geq\Phi(u)+\faip{a}{v-u}$ for all $v$, is monotone and
has onto $I+\lambda\partial\Phi$ for every $\lambda>0$.
Its resolvent uniquely minimizes
$\Phi(v)+\|v-f\|^2/(2\lambda)$.
\end{proposition}
\begin{proof}
Adding two subgradient inequalities proves monotonicity.
Hahn--Banach separation of the epigraph gives an affine lower bound
for $\Phi$, so the displayed functional is coercive and bounded
below. Convex norm-closed sublevel sets are weakly closed by
Hahn--Banach; reflexivity and coercivity give a minimizer.
Strict convexity of the squared norm gives uniqueness.
The one-sided derivative inequality along $u+t(v-u)$ at a minimum,
combined with convexity of $\Phi$, is exactly
$\Phi(v)\geq\Phi(u)+\faip{(f-u)/\lambda}{v-u}$.
Conversely this inequality makes the energy difference at $v$
at least $\|v-u\|^2/(2\lambda)$. Thus
$f\in u+\lambda\partial\Phi(u)$, proving surjectivity and the
resolvent description.
\end{proof}

% END INPUT


% BEGIN parts/part3-oscillator.tex
% Confining-potential and Hermite supplement, authored for the expanded notes.
\section{Confining potentials and the Hermite oscillator}

\begin{theorem}[Weighted energy space and a confining compact resolvent]
\label{fa2:p3:confining-potential}
Let $p\in C(\mathbb R)$, $p\ge0$, and set
\[
 V=\{u\in H^1(\mathbb R):\sqrt p\,u\in L^2\},\qquad
 \|u\|_V^2=\|u'\|_2^2+\|u\|_2^2+\|\sqrt p\,u\|_2^2.
\]
This is a separable Hilbert space with $C_c^\infty$ dense.
Suppose $p(x)\ge\delta>0$ for $|x|\ge R$.
The form $a(v,u)=\int(\overline{v'}u'+p\overline vu)$
is coercive on $V$. For every $f\in L^2$ it gives a unique
$u=Tf\in V$, with $a(v,u)=\langle v,f\rangle$.
The operator $T$ on $L^2$ is bounded, positive, self-adjoint,
injective and has dense range. Its inverse $A=T^{-1}$ is the
positive self-adjoint realization of $-u''+pu$, with
\[
 D(A)=\{u\in V:-u''+pu\in L^2\text{ in distributions}\}.
\]
If $p(x)\to\infty$ as $|x|\to\infty$, $T$ is compact.
There is then an orthonormal eigenbasis $e_n$ and numbers
$0<\lambda_n\to\infty$ with $Ae_n=\lambda_ne_n$,
and $D(A)=\{\sum c_ne_n:\sum\lambda_n^2|c_n|^2<\infty\}$.
\end{theorem}
\begin{proof}
Completeness follows by viewing $u\mapsto(u,u',\sqrt p\,u)$
as a closed subspace of the product of three $L^2$ spaces:
an almost-everywhere convergent subsequence identifies the
weighted component, and distribution testing identifies $u'$.
The product is separable. Cutoffs tend to $u$ in all three
norms; the extra derivative of the cutoff tends to zero.
For a fixed cutoff, convolution converges in $H^1$ and in the
weighted norm because $p$ is bounded on its enlarged compact
support. Thus $C_c^\infty$ is dense.

Choose $y\in[R,R+1]$ with
$|u(y)|^2\le\int_R^{R+1}|u|^2\le\delta^{-1}\int p|u|^2$.
For $|x|\le R$, the AC formula and Cauchy--Schwarz give
\[
 |u(x)|^2\le2\delta^{-1}\int p|u|^2
                         +2(2R+1)\int|u'|^2.
\]
Outside $[-R,R]$ the potential directly controls $\|u\|_2^2$.
Integration of this estimate proves
$\|u\|_V^2\le C a(u,u)$, hence coercivity.
The proved Lax--Milgram theorem constructs $T:L^2\to V$.
Its form symmetry gives
$\langle f,Tg\rangle=a(Tf,Tg)=\langle Tf,g\rangle$,
and $\langle f,Tf\rangle=a(Tf,Tf)\ge0$.
If $Tf=0$, testing against $C_c^\infty$ gives $f=0$;
self-adjointness then makes the range dense.
Testing the variational equation proves the distributional
domain formula in one direction; in the other, the equality
against compact tests extends by $V$-density, proving it for
all $v\in V$ and hence $u=Tf$. A bounded positive injective
self-adjoint operator has self-adjoint inverse on its dense
range by the measurable spectral calculus in this part;
its inverse domain is precisely the square-integrability
domain of $1/t$. This proves the self-adjoint realization.

If $p$ confines, the unit ball of $V$ has uniformly small
$L^2$ tails, bounded by $1/\inf_{|x|>R}p(x)$.
Rellich on each bounded interval and a diagonal subsequence
give global $L^2$ precompactness. Thus $T:L^2\to V\to L^2$
is compact. The positive compact spectral theorem, with zero
kernel, gives an orthonormal basis and eigenvalues
$\mu_n>0$ tending to zero. There are infinitely many because
$T$ is injective on infinite-dimensional $L^2$.
Set $\lambda_n=1/\mu_n$; the displayed inverse domain
follows directly from Parseval. This proves the general bridge
in \cite{B01}, Problem 51, before specializing to a model.
\end{proof}

\begin{theorem}[Hermite basis and the harmonic oscillator]
\label{fa2:p3:hermite-oscillator}
For $p(x)=x^2$, define
\[
 H_n(x)=(-1)^ne^{x^2}\frac{d^n}{dx^n}e^{-x^2},\qquad
 h_n(x)=\frac{e^{-x^2/2}H_n(x)}{(2^nn!\sqrt\pi)^{1/2}}.
\]
Then $(h_n)_{n\ge0}$ is an orthonormal basis of $L^2(\mathbb R)$,
$Ah_n=(2n+1)h_n$, and
\[
 D(A)=\left\{\sum c_nh_n:\sum(2n+1)^2|c_n|^2<\infty\right\}.
\]
The restriction of $-d^2/dx^2+x^2$ to $\mathcal S$ is
essentially self-adjoint, with this closure. Its resolvent is
compact. The sharp ground-state inequality is
$\|u'\|_2^2+\|xu\|_2^2\ge\|u\|_2^2$ on $V$;
equality holds exactly for Gaussian multiples.
\end{theorem}
\begin{proof}
Taylor's formula for $e^{-(x-t)^2}$ gives the generating identity
\[
 e^{2xt-t^2}=\sum_{n=0}^\infty H_n(x)t^n/n!.
\]
It shows that $H_n$ has degree $n$ and leading coefficient
$2^n$. Differentiating the generating function gives
$H_n''-2xH_n'+2nH_n=0$, by equating coefficients in
$(\partial_x^2-2x\partial_x+2t\partial_t)e^{2xt-t^2}=0$.
Consequently direct differentiation proves
$(-d^2/dx^2+x^2)h_n=(2n+1)h_n$.
Rodrigues' formula and $n$ integrations by parts give
\[
 \int H_nH_me^{-x^2}dx=0\ (m<n),\qquad
 \int H_n^2e^{-x^2}dx=2^nn!\sqrt\pi.
\]
All boundary terms vanish by Gaussian decay. This proves
orthonormality including the normalization constants.

To prove completeness, suppose $g\in L^2$ is orthogonal to
every $h_n$. Since their polynomial factors span all polynomials,
the integrable function $v=g e^{-x^2/2}$ has every moment zero
(replace $g$ by its conjugate if using the other pairing slot).
For every $a>0$, Cauchy--Schwarz gives
$\int|v(x)|e^{a|x|}dx<\infty$. Thus
$F(z)=\int v(x)e^{-2\pi izx}dx$ is entire, by locally
uniform dominated differentiation; all its derivatives at zero
vanish. Its Taylor series makes $F\equiv0$. Fourier uniqueness
for $L^1$ gives $v=0$ and hence $g=0$. This proves completeness,
rather than inferring a basis from the existence of eigenvectors.

Each $h_n$ is Schwartz and belongs to the operator domain of
the preceding theorem, so the spectral domain formula follows.
Finite Hermite sums are graph-norm dense in that domain by
truncation of the weighted coefficient square sum. They lie in
$\mathcal S$, and the Schwartz differential operator is a
restriction of $A$. Its closure therefore equals $A$, proving
essential self-adjointness. The eigenvalues of $(A+1)^{-1}$
are $(2n+2)^{-1}\to0$, giving compactness directly as well.
Finally integration by parts for compact smooth $u$ gives
\[
 \|u'+xu\|_2^2=\|u'\|_2^2+\|xu\|_2^2-\|u\|_2^2.
\]
$V$-density passes it to all $u\in V$. Equality means
$u'+xu=0$ distributionally; multiplying by $e^{x^2/2}$
on compact intervals gives a constant, hence precisely
$u=ce^{-x^2/2}$. This also proves the exact lowest eigenvalue.
\end{proof}

% END INPUT


% BEGIN parts/part3-krein-extensions.tex
\chapter{Extremal nonnegative extensions and Stieltjes determinacy}
\label{fa2:p3:krein-extensions-chapter}
The nonnegative moment problem has a different uniqueness criterion
from the unrestricted Hamburger problem. A symmetric operator can
have a unique nonnegative self-adjoint extension without having a
unique self-adjoint extension of every sign. The order argument below
supplies the precise bridge used in \cite[Section 33.5]{B08}.

\begin{theorem}[Friedrichs--Krein resolvent bounds]
\label{fa2:p3:krein-resolvent-order}
Let $S$ be densely defined, closed, symmetric and nonnegative on a
complex Hilbert space $H$. It has nonnegative self-adjoint extensions
$S_F,S_K$ such that for every $a>0$ and every nonnegative self-adjoint
extension $T$ on a Hilbert space $\mathcal H\supset H$, isometrically
containing $H$ as a closed subspace,
\[
 (S_F+a)^{-1}\ \le\ P_H(T+a)^{-1}|_H\ \le\ (S_K+a)^{-1}.
\]
Order is quadratic-form order on bounded self-adjoint operators.
The inclusion $T\supset S$ refers to both the action and its original
domain in $H$. The construction does not require $S\ge cI$ with $c>0$.
\end{theorem}
\begin{proof}
Put $q(x,y)=\langle x,Sy\rangle$ on $D(S)$ and let $W$ be its
Hilbert completion after quotienting its null vectors. If $q(y,y)=0$,
form Cauchy--Schwarz gives $\langle x,Sy\rangle=0$ for every
$x\in D(S)$, whence $Sy=0$ by density. Thus the functional below
is well defined on the quotient.
For $x\in H$ whose map $[y]\mapsto\langle x,Sy\rangle$ is
bounded on this energy space, let $Rx\in W$ be its representing
vector:
\[
 \langle Rx,[y]\rangle_W=\langle x,Sy\rangle_H
 \quad(y\in D(S)).
\]
Its domain is dense because it contains $D(S)$; for $x$ there,
$Rx=[x]$ by symmetry. The operator $R:H\to W$ is closed. Indeed
if $x_n\to x$ in $H$ and $Rx_n\to w$ in $W$, testing against each
$[y]$ gives $\langle w,[y]\rangle=\langle x,Sy\rangle$, proving
$x\in D(R)$ and $Rx=w$. Therefore
$q_K(x,y)=\langle Rx,Ry\rangle_W$ is a densely defined closed
nonnegative form on $D(R)$, with its graph norm. The closed-form
representation theorem gives a nonnegative self-adjoint $S_K$.
For $y\in D(S)$ and $x\in D(R)$,
$q_K(x,y)=\langle Rx,[y]\rangle=\langle x,Sy\rangle$;
thus $S_K$ extends $S$. Let $S_F$ be the Friedrichs extension
associated with the closure $q_F$ of $q$; its existence and extension
property were proved in \ref{fa2:p3:friedrichs}.

For any closed nonnegative form $b$ with associated operator $B$,
completion of the square at $u=(B+a)^{-1}f$ gives the variational
identity
\[
 \langle f,(B+a)^{-1}f\rangle
   =\sup_{x\in D(b)}\{2\operatorname{Re}\langle f,x\rangle
                           -b(x,x)-a\|x\|^2\}.
\]
The expression is real; replacing $x$ by $u+w$ makes its difference
from the maximum equal to $-b(w,w)-a\|w\|^2$. For $S_F$ the same
supremum can be taken over $D(S)$, its form core.

Let $b=q_T$ be the nonnegative form of $T$. For $y\in D(S)$ it
agrees with $q$ and $Ty=Sy$. Its variational supremum, for $f\in H$,
is therefore at least the supremum on $D(S)$, proving the lower
resolvent bound. For any $x\in D(q_T)$, form Cauchy--Schwarz gives
\[
 |\langle P_Hx,Sy\rangle|
   =|q_T(x,y)|\le q_T(x,x)^{1/2}q(y,y)^{1/2}.
\]
Consequently $P_Hx\in D(R)$ and
$q_K(P_Hx,P_Hx)\le q_T(x,x)$. Moreover
$\|P_Hx\|\le\|x\|$ and
$\langle f,x\rangle=\langle f,P_Hx\rangle$.
Every term in the variational supremum for the compressed resolvent
of $T$ is thus no greater than the corresponding term for $q_K$
at $P_Hx$. Taking suprema proves the upper bound.
This proves the assertions also for extensions in a larger space,
which is important when polynomials are not dense in a representing
measure's full $L^2$ space.
\end{proof}

\begin{theorem}[Stieltjes existence and exact uniqueness criterion]
\label{fa2:p3:stieltjes-determinacy}
Let $s_n$ be a real sequence and define $L(t^n)=s_n$ on the complex
polynomials. There is a finite positive measure $\mu$ on $[0,\infty)$
with all moments finite and $s_n=\int t^n\,d\mu$ if and only if
\[
 L(\overline p p)\ge0,\qquad L(t\overline p p)\ge0
 \quad\text{for every polynomial }p.
\]
In its polynomial Hilbert space, let $S$ be the closure of
multiplication by $t$ on the polynomial classes. The representing
measure is unique if and only if $S$ has exactly one nonnegative
self-adjoint extension in that Hilbert space, equivalently
$S_F=S_K$. This condition is not replaced by essential
self-adjointness, which corresponds to the stronger unrestricted
Hamburger uniqueness condition.
\end{theorem}
\begin{proof}
Necessity follows by integration. For sufficiency quotient the
polynomials by their zero seminorm for
$\langle p,q\rangle=L(\overline p q)$ and complete to $H$.
Multiplication is well defined: if $L(|p|^2)=0$, positivity and
polynomial Cauchy--Schwarz give
$L(\overline p t^2p)=0$, so $\|tp\|^2=0$.
It is densely defined symmetric, nonnegative by the second
condition, and hence closable. Its Friedrichs extension is
nonnegative self-adjoint. The scalar spectral measure of the
constant polynomial vector $\xi=[1]$ represents the sequence:
every polynomial vector remains in the domains of all the needed
powers of the extension, and $S_F^n\xi=[t^n]$. Thus
$\langle\xi,S_F^n\xi\rangle=s_n$ and the spectral support
is contained in $[0,\infty)$. Finiteness of all absolute moments follows
from the even moments and Cauchy--Schwarz. If $s_0=0$, polynomial
Cauchy--Schwarz makes every $s_n=0$; the zero measure is the
unique case and may be disposed of separately.

Suppose first that $S_F=S_K$. Any representing measure gives an
isometry $J:H\to L^2(\mu)$ by the polynomial classes. Its image
is closed, but need not be the whole $L^2(\mu)$. The multiplication
operator $T=M_t$ there is nonnegative self-adjoint and extends $S$
on that isometrically included subspace. The preceding bounds
force all the compressed resolvents to coincide. In particular
\[
 \int_0^\infty\frac{d\mu(t)}{t+a}
   =\langle\xi,(S_F+a)^{-1}\xi\rangle\quad(a>0)
\]
is independent of the representing measure. These values determine
a finite measure on $[0,\infty)$. Here is one elementary uniqueness
argument: differentiating at $a=1$, justified by the bounded kernels
and finite mass, gives equality of the integrals of
$(1+t)^{-n}$ for every $n\ge1$; the total mass $s_0$ is fixed too.
Push forward by $x=(1+t)^{-1}$ into $[0,1]$, assigning no mass at
zero. Polynomial density on $[0,1]$ and finite-measure uniqueness
make the pushforwards, hence the original measures, equal.

Conversely suppose there is only one representing measure. Each
nonnegative self-adjoint extension $T$ of $S$ in $H$ gives such a
measure $\mu_T$ at $\xi$. If two extensions give the same measure,
then for polynomials $p,q$ and $a>0$,
\[
 \langle[p],(T+a)^{-1}[q]\rangle
      =\int\frac{\overline{p(t)}q(t)}{t+a}\,d\mu_T(t).
\]
The integrals are finite by the available polynomial moments.
Thus their bounded resolvents agree on the dense polynomial
subspace, and hence on all of $H$; one such resolvent determines
the operator and its domain. The extensions are the same.
In particular uniqueness of the measure forces $S_F=S_K$.
Finally if those two extensions coincide, the order bounds applied
to any extension in $H$ force its resolvent to equal theirs.
This establishes all three equivalent uniqueness conditions and
retains the distinction from unrestricted self-adjoint extension.
\end{proof}

% END INPUT


% BEGIN parts/part3-weyl.tex
\chapter{Uniqueness of the regular finite-dimensional Weyl representation}
\label{fa2:p3:weyl-chapter}
An identity $[Q,P]=iI$ on an unspecified common domain is not the
hypothesis of the uniqueness theorem. We require the strongly
continuous unitary Weyl relations, which fix both the domains through
Stone's theorem and the nonzero central character. The proof below
uses a Gaussian integral to construct a rank-one corner, and also
gives the multiplicity version.
The Gaussian integrals below use angular Fourier frequency $p$;
it equals $2\pi$ times the frequency in the global Fourier convention.

\begin{theorem}[Stone--von Neumann, with multiplicity]
\label{fa2:p3:stone-von-neumann}
Let $d\ge1$, let $H\ne\{0\}$ be a complex Hilbert space, and let
$W:\mathbb R^{2d}\to\mathcal U(H)$ be strongly continuous with
\[
 W(z)W(w)=e^{i\omega(z,w)/2}W(z+w),\qquad W(0)=I,
 \quad \omega((q,p),(q',p'))=p\cdot q'-p'\cdot q.
\]
There is a nonzero Hilbert space $M$ and a unitary
$H\simeq L^2(\mathbb R^d)\otimes M$ under which
\[
 W(q,p)=W_0(q,p)\otimes I_M,
 \qquad (W_0(q,p)f)(t)=e^{ip\cdot(t-q/2)}f(t-q).
\]
In particular an irreducible such representation is unitarily
equivalent to $W_0$, and $W_0$ is irreducible. A specified nonzero
constant in place of $1$ in the phase is reduced to this form by a
linear rescaling of one set of coordinates.
\end{theorem}
\begin{proof}
Put $n=2d$, $c=(2\pi)^{-d}$ and $a(z)=e^{-|z|^2/4}$. Define
\[
 P=c\int_{\mathbb R^{2d}}a(z)W(z)\,dz.
\]
This is a strong integral: for each $x\in H$ the vector function
$a(z)W(z)x$ is Bochner integrable and its integral is bounded by
$\|a\|_1\|x\|$. We do not assert norm continuity or Bochner
measurability of $W$ as a function into $\mathcal B(H)$.
Adjunction, $W(z)^*=W(-z)$ and the even real Gaussian give $P^*=P$.
Vector-valued Fubini applies to all the following integrals, since
their norm bounds are products of integrable Gaussians.

The Weyl law and the change of variables $r=z+v$ give
\[
 P^2=c^2\int\left[\int a(z)a(r-z)e^{i\omega(z,r)/2}\,dz\right]W(r)\,dr=P.
\]
For clarity, complete the square: the product of Gaussians is
$e^{-|z-r/2|^2/2}e^{-|r|^2/8}$, and the Fourier transform of the
first factor at the symplectic frequency $r/2$ is
$(2\pi)^d e^{-|r|^2/8}$. The bracket is therefore
$(2\pi)^d a(r)$. This proves that $P$ is an orthogonal projection.
The same computation with one intervening Weyl operator gives
\begin{equation}\label{fa2:p3:gaussian-corner}
 P W(w) P=a(w)P.
\end{equation}
Here are the phases, to fix all normalizations. In the double integral
put $r=z+w+v$. Its phase becomes
$[\omega(z,r+w)+\omega(w,r)]/2$. Centering the product
$a(z)a(r-w-z)$ at $(r-w)/2$ cancels the constant phase. The Fourier
Gaussian integral leaves
$(2\pi)^d\exp(-|r-w|^2/8-|r+w|^2/8)
=(2\pi)^d a(r)a(w)$, proving the identity.

The projection is nonzero in every nonzero representation.
If $P=0$, conjugation by $W(w)$ would give, for every $x,y$,
\[
 \int a(z)e^{i\omega(w,z)}\langle x,W(z)y\rangle\,dz=0
 \quad(w\in\mathbb R^{2d}).
\]
The integrand without its phase belongs to $L^1$. Uniqueness of the
Fourier transform makes it zero almost everywhere. It is continuous,
so it is zero everywhere. At $z=0$ this gives $\langle x,y\rangle=0$
for every $x,y$, contradicting $H\ne0$. The required Fourier
uniqueness follows, for example, by convolving with the Gaussian
approximate identity and applying its integrable Fourier inversion,
then passing to the $L^1$ limit. It does not require a representation
theorem about $W$.

In the concrete representation let
$h(t)=\pi^{-d/4}e^{-|t|^2/2}$. Direct integration gives
$\|h\|_2=1$ and $\langle h,W_0(z)h\rangle=a(z)$.
One can also compute the kernel of its Gaussian projection:
integrating first in $p$ yields
\[
 (P_0f)(t)=\pi^{-d/2}e^{-|t|^2/2}
                \int e^{-|s|^2/2}f(s)\,ds
          =h(t)\langle h,f\rangle.
\]
First do this for Schwartz $f$, where Fubini is immediate, and then
extend by the bounded integral estimate. Thus $P_0$ is the projection
onto $\mathbb Ch$.
The span of $\{W_0(z)h:z\in\mathbb R^{2d}\}$ is dense in $L^2$.
Indeed if $f$ is orthogonal to it, use only $z=(0,p)$: the Fourier
transform of $\overline f h\in L^1$, by Cauchy--Schwarz, vanishes.
Fourier uniqueness and strict positivity of $h$ give $f=0$.

Set $M=PH$. On elementary vectors define
\[
 V\big(W_0(z)h\otimes\xi\big)=W(z)\xi\qquad(\xi\in M).
\]
This is an isometry on finite linear combinations. To check that it
is well defined, compute the Gram matrix: the Weyl law expresses
$W(z)^*W(w)$ as the same scalar phase times $W(w-z)$ in both
representations. Equation \eqref{fa2:p3:gaussian-corner} gives
\[
 \langle \xi,W(z)^*W(w)\eta\rangle
   =\langle W_0(z)h,W_0(w)h\rangle\langle\xi,\eta\rangle.
\]
Density therefore extends $V$ to an isometry on $L^2\otimes M$.
Its range is the closed span of $W(z)M$ and reduces every $W(w)$,
by the Weyl law. If the orthogonal complement were nonzero, restricting
the representation to it would be strongly continuous and would have
Gaussian projection zero, because it is orthogonal to $M$. This
contradicts the nonvanishing argument. Thus $V$ is unitary, and the
same Weyl law verifies that it intertwines all $W_0(w)\otimes I$
with $W(w)$.

If $\dim M>1$, a nontrivial orthogonal projection on $M$ supplies
a reducing subspace, so irreducibility forces $\dim M=1$.
Conversely an orthogonal projection commuting with all $W_0(z)$
commutes with their strong Gaussian integral $P_0$ and therefore
sends $h$ to $0$ or $h$. Its commutation and the density of the Weyl
orbit of $h$ make it $0$ or $I$. This proves irreducibility of $W_0$.
The asserted uniqueness and the rescaling conclusion follow.
\end{proof}

\begin{remark}[The precise scope of uniqueness]
The theorem concerns a finite number of position and momentum
coordinates, a nonzero fixed central character, unitarity, and strong
continuity of the Weyl relations. Dropping regularity, changing the
central character to zero, or passing to infinitely many canonical
coordinates changes the problem; none of those inequivalent
representations is ruled out by the theorem just proved.
\end{remark}

% END INPUT


% BEGIN parts/part3-source-refinements.tex
\chapter{Further source variants for semigroups and translation models}
\label{fa2:p3:source-refinements}

\begin{theorem}[Weak continuity at zero already gives a $C_0$ semigroup]
\label{fa2:p3:weak-semigroup-continuity}
Let $T(t)$, $t\ge0$, be bounded linear operators on a Banach space,
with $T(0)=I$ and $T(t+s)=T(t)T(s)$. If
$\ell(T(t)x)\to\ell(x)$ as $t\downarrow0$ for every $x,\ell$,
then $T$ is strongly continuous. This is the weak-continuity
variant in \cite[Section 34.1]{B08}.
\end{theorem}
\begin{proof}
Uniform boundedness, first on the dual tests and then on the
vectors, gives $\sup_{0\le t\le\delta}\|T(t)\|<\infty$ for
some $\delta>0$; otherwise the weakly convergent families at zero
would be unbounded. More explicitly, failure would give
$t_n\downarrow0$ with $\|T(t_n)\|>n$; weak convergence of
$T(t_n)x$ for every $x$, and twice uniform boundedness, contradict
that inequality. The semigroup property extends this bound to
each finite interval. For fixed $x$, $T(t+h)x\to T(t)x$ weakly
as $h\downarrow0$. Thus the scalar orbit functions are right
continuous and Borel measurable. Moreover every orbit value is
in the norm-closed linear span of the values at nonnegative
rational times: rational times decreasing to $t$ give weak
convergence, and a norm-closed linear subspace is weakly closed
by Hahn--Banach. This span is separable. On a separable Banach
space choose a countable norming family of continuous linear
functionals by taking supporting functionals at a dense subset
of its unit sphere. The norm distance to a fixed vector is the
supremum of their absolute values, hence is Borel on this orbit.
Approximation by the first point in a countable dense set within
distance $1/n$ now supplies measurable simple approximants.
The orbit is strongly measurable, bounded on compact intervals,
and Bochner integrable there.

Put $J_a x=\int_0^aT(s)x\,ds$. For $h>0$ small,
\[
 T(h)J_a x-J_a x=\int_a^{a+h}T(s)x\,ds-\int_0^hT(s)x\,ds,
\]
whose norm tends to zero. Thus every $J_a x$ is a strong
continuity vector. Their linear span is dense: a functional
annihilating all these vectors has
$\ell(J_a x)/a\to\ell(x)$ by weak continuity at zero, and
therefore annihilates every $x$. Hahn--Banach proves the density.
The locally uniform operator bound extends strong continuity
at zero from that span to all vectors. The semigroup property
gives right continuity at every time and, for $0<h<t$,
$\|T(t)x-T(t-h)x\|\le\|T(t-h)\|\|T(h)x-x\|$, giving left
continuity too.
\end{proof}

\begin{theorem}[Stable dual-consistent schemes produce weak solutions]
\label{fa2:p3:weak-scheme-existence}
Let $G$ be a closed densely defined operator on a complex Hilbert
space $H$, and suppose $D(G^*)$ is dense. Let $C_h\in\mathcal B(H)$
for $h\downarrow0$, and write $G_h=(C_h-I)/h$. Suppose that for
each $T<\infty$ the powers $C_h^n$, $nh\le T$, are uniformly
bounded, and
$G_h^*w\to G^*w$ for every $w\in D(G^*)$.
For each $x\in H$, the step functions
$u_h(t)=C_h^{\lfloor t/h\rfloor}x$ have weakly convergent
subsequences in $L^2(0,T;H)$. Every such weak limit belongs to
$L^\infty(0,T;H)$ and satisfies
\[
 \int_0^T\big(\langle w'(t),u(t)\rangle+
                    \langle G^*w(t),u(t)\rangle\big)dt
       +\langle w(0),x\rangle=0
\]
for $w\in C^1([0,T];H)\cap C([0,T];D(G^*)_{\rm graph})$
with $w(T)=0$. No uniqueness or strong convergence is asserted
without additional hypotheses. This is the weak-existence
variant of \cite[Section 34.3, Theorem 9]{B08}.
\end{theorem}
\begin{proof}
The power bound makes $u_h$ bounded in $L^2H$ and $L^\infty H$.
A bounded sequence in a Hilbert space has a weakly convergent
subsequence: its countable closed linear span is separable,
and a diagonal argument on a countable orthonormal basis,
followed by the norm bound, constructs its weak limit.
If the bound is $B$, weak lower semicontinuity on any measurable
set $E$ gives $\int_E\|u\|^2\le B^2|E|$. Thus $\|u\|\le B$
a.e., proving the $L^\infty$ assertion.

The maps $G_h^*:D(G^*)_{\rm graph}\to H$ are bounded and converge
pointwise to $G^*$. The graph domain is Banach, so uniform
boundedness gives a common norm bound. Their convergence is
therefore uniform on compact subsets of that graph domain,
in particular on the compact range of $w$. Take $N$ with
$Nh\le T<(N+1)h$, and set $u_n=C_h^nx$ and $w_n=w(nh)$.
The exact summation-by-parts identity is
\[
 h\sum_{n=0}^{N-1}\langle G_h^*w_{n+1},u_n\rangle
 =\langle w_N,u_N\rangle-\langle w_0,x\rangle
       -\sum_{n=0}^{N-1}\langle w_{n+1}-w_n,u_n\rangle.
\]
It follows by pairing $u_{n+1}-u_n=hG_hu_n$ with $w_{n+1}$
and telescoping. The endpoint term tends to zero because
$w(T)=0$, $w$ is continuous, and $u_N$ is bounded.
The difference quotients of $w$ converge uniformly to $w'$
in the corresponding piecewise-average sense. The preceding
compact-range convergence and continuity of $G^*w$ make
$G_h^*w_{n+1}$ uniformly close to $G^*w(t)$ on its cell.
The remaining interval of length at most $h$ contributes zero
in the limit by the uniform bounds. Pass to the weak $L^2$
limit in the two fixed continuous test functions. The displayed
weak identity follows, including its initial term and its signs.
\end{proof}

\begin{theorem}[Spectral mapping under eventual norm continuity]
\label{fa2:p3:eventual-norm-spectral-mapping}
Let $T$ be a $C_0$ semigroup on a complex Banach space with generator
$A$. If $t\mapsto T(t)$ is operator-norm continuous for all
$t>t_0$ for some $t_0\ge0$, then for every $t>0$,
\[
 \sigma(T(t))\setminus\{0\}=\exp(t\sigma(A)).
\]
This is the norm-continuity variant of
\cite[Section 34.5, Theorem 12$'$]{B08}; eventual compactness is
a sufficient stronger hypothesis, not the hypothesis used here.
\end{theorem}
\begin{proof}
The inclusion from right to left was proved in
\ref{fa2:p3:semigroup-spectral-mapping}. Put all $T(s)$ and
one real resolvent $R(\alpha,A)$ in a maximal commutative unital
subalgebra $\mathcal A$ of $\mathcal B(X)$, taking its norm closure.
Such an algebra exists by the maximal principle. It is inverse
closed: if an element is invertible in $\mathcal B(X)$, its inverse
commutes with every element of $\mathcal A$, and maximality puts
the inverse in $\mathcal A$. Thus spectra in the two algebras agree.
For $0\ne z\in\sigma(T(t))$ choose a character $\chi$ of
$\mathcal A$ with $\chi(T(t))=z$ and put $f(s)=\chi(T(s))$.
Then $f(s+r)=f(s)f(r)$. It never vanishes: if $f(s)=0$,
then $f(r)=0$ for every $r\ge s$, contradicting
$f(nt)=z^n\ne0$ for large integers $n$.
It is continuous for $s>t_0$ by the norm-continuity assumption,
and $f(s+h)=f(s)f(h)$ with a fixed $s>t_0$ gives continuity
at zero. It extends to a continuous nonzero multiplicative
function on $\mathbb R$ by $f(-s)=f(s)^{-1}$.
Such functions are $e^{\lambda s}$: near zero a continuous
logarithm exists, and additivity of its small increments follows
because their difference lies in $2\pi i\mathbb Z$ and is
continuous and zero at zero. A continuous additive function is
linear, first on rational and then on real parameters.

Take $b>t_0$ and $\alpha$ beyond the semigroup growth bound.
The Laplace formula multiplied by $T(b)$ is now a norm integral,
since $T(b+s)$ is norm continuous and exponentially bounded:
\[
 R(\alpha,A)T(b)=\int_0^\infty e^{-\alpha s}T(b+s)\,ds.
\]
The bound $|f(s)|\le\|T(s)\|$ gives $\operatorname{Re}\lambda$
no larger than that growth bound. Applying $\chi$ to the norm
integral and dividing by $f(b)$ yields
$\chi(R(\alpha,A))=(\alpha-\lambda)^{-1}\ne0$.
Resolvent spectral mapping, or directly the inverse identity
$(\mu-R)^{-1}$ versus $(\alpha-\mu^{-1}-A)^{-1}$, implies
$\lambda\in\sigma(A)$. Finally $z=f(t)=e^{\lambda t}$.
\end{proof}

\begin{proposition}[The full translation model and its exact multiplicity condition]
\label{fa2:p3:translation-model}
Let $M$ be any complex Hilbert space. On $L^2(\mathbb R;M)$,
$V(t)f(s)=f(s+t)$ is a strongly continuous unitary group with
self-adjoint Stone operator $-i\,d/ds$, on the domain of functions
whose distributional derivative belongs to $L^2(\mathbb R;M)$.
Its angular Fourier model is multiplication by $e^{it\lambda}$,
with Stone operator multiplication by $\lambda$ on
$\{g:\lambda g\in L^2\}$.
A self-adjoint operator, or its unitary group, has this full
translation representation exactly when it has a spectral model
with Lebesgue measure and constant multiplicity $M$.
Absolute continuity and full topological support alone do not
assert that constant multiplicity.
\end{proposition}
\begin{proof}
Finite sums $f(s)=\sum_{j=1}^n f_j(s)e_j$, with $f_j$ smooth
and compactly supported, are dense in $L^2(\mathbb R;M)$.
Indeed Bochner functions have essentially separable range;
simple approximation and scalar smooth approximation prove this
even if $M$ is nonseparable. Translation is isometric on such
sums and strongly continuous there, hence on the completion.
Scalar Plancherel applied to the finitely many coordinates
extends their angular Fourier transform to a unitary on that
completion; Parseval and density show that it is onto.
The scalar change of variables gives
$\mathcal F(V(t)f)(\lambda)=e^{it\lambda}\mathcal Ff(\lambda)$.
Dominated convergence shows that the strong difference quotient
exists when $\lambda\mathcal Ff\in L^2$, with value
$i\lambda\mathcal Ff$. Conversely existence of the quotient
and Fatou's lemma force that weighted square integrability,
and pointwise convergence identifies its limit. This proves
the exact generator domain. The identity between weak
derivatives and Fourier multiplication follows first by testing
finite smooth coordinate functions and then by distributional
duality and density. Thus the stated spatial domain is exact.
Conjugating this model by a specified spectral unitary proves
one implication of the representation criterion; Fourier
transformation of a translation representation proves the other.

For the warning take the scalar multiplication operator on
$L^2(E,d\lambda)$, where $E$ has positive measure and its
complement also has positive measure in every nonempty interval.
Such an $E$ is constructed by placing disjoint positive-measure
compact nowhere-dense subsets of each rational interval,
alternately into $E$ and its complement. At each finite stage
the unused open part of the next interval is nonempty, so the
construction continues; assign the remaining points arbitrarily.
Its scalar spectral measure is absolutely continuous and has
full topological support, but its multiplicity is zero on
$E^c$ and one on $E$. The uniqueness of measurable spectral
multiplicity from \ref{fa2:p2:spectral-multiplicity} prevents
equivalence with any nonzero constant-multiplicity Lebesgue
model on the whole line. This pinpoints the measure-class
condition needed in reading \cite[Section 35.5]{B08}.
\end{proof}

% END INPUT


% END INPUT

\part{Function spaces, distributions and analytic tools}

% BEGIN parts/part4.tex
% Part IV. New exposition; no source book or earlier manuscript is modified.
\providecommand{\pfnorm}[1]{\left\lVert#1\right\rVert}
\providecommand{\pvnorm}[1]{\left\lVert#1\right\rVert}
\providecommand{\pfabs}[1]{\left|#1\right|}
\providecommand{\pfsupp}{\operatorname{supp}}
\providecommand{\pvsupp}{\operatorname{supp}}
\providecommand{\pfD}{\mathcal D}
\providecommand{\pfS}{\mathcal S}
\providecommand{\pfF}{\mathcal F}
\chapter{Function Spaces, Integral Representation, and Vector Integration}
\label{fa2:p4:function-spaces}

Scalar measure theory is used with the following precise conventions.
Measures are positive and countably additive, functions are identified almost
everywhere, and all completed-measure assertions are understood on the
completion. Tonelli applies to nonnegative measurable functions on
$\sigma$-finite products, and Fubini applies when the absolute integral is finite.
Monotone convergence, Fatou, dominated convergence, and the construction of
Lebesgue measure are the scalar integration foundations. These statements,
including their proofs, occur respectively at D11 labels
\texttt{d11:05:monotone-convergence}, \texttt{d11:05:fatou},
\texttt{d11:06:dominated-convergence}, \texttt{d11:08:tonelli}, and
\texttt{d11:08:fubini} in the accepted measure-theory manuscript.
No vector integration or distribution theorem is being included in this
foundation convention. Unless stated otherwise, the scalar field is
$\mathbb K=\mathbb R$ or $\mathbb C$, and a functional is scalar-linear.
The integral duality pairing is $\int fg$, with no conjugation.
Hilbert products elsewhere in this book are conjugate-linear first.

\section{Scalar duals and representation}

\begin{theorem}[Sequence duals]\label{fa2:p4:sequence-duals}
Integration against a sequence gives surjective linear isometries
$\ell^1\simeq(c_0)'$, $\ell^\infty\simeq(\ell^1)'$, and
$\ell^q\simeq(\ell^p)'$ for $1<p<\infty$, $p^{-1}+q^{-1}=1$.
Moreover $c'=\ell^1\oplus\mathbb K\lambda$, where
$\lambda(x)=\lim x_n$; this is an algebraic decomposition with
$L(x)=\sum a_nx_n+b\lambda(x)$ and
$\pfnorm L=\sum|a_n|+|b|$.
\end{theorem}
\begin{proof}
Put $a_n=L(e_n)$. Testing finite sums with phases matching $a_n$ proves
$\sum|a_n|\le\pfnorm L$ for $L$ on $c_0$; finite-support density then
gives $L(x)=\sum a_nx_n$, and the same tests give norm equality.
For $L$ on $\ell^1$, $|a_n|\le\pfnorm L$, and truncation gives the
representation and equality. For $\ell^p$, test on finite truncations of
$x_n=\overline{a_n}|a_n|^{q-2}$, defining this expression as zero if
$a_n=0$. The inequality
$\sum_F|a_n|^q\le\pfnorm L(\sum_F|a_n|^q)^{1/p}$ shows that
$a\in\ell^q$; truncation and H\"older give the assertion.
Restrict a functional on $c$ to $c_0$ and subtract its $\ell^1$
representative. The difference vanishes on $c_0$ and hence is $b\lambda$.
For norm equality choose the first finitely many coordinates to match the
phases of $a_n$, and all subsequent coordinates equal to a phase matching
$b$. The tail of $a$ tends to zero, so these unit-ball tests approach
$\sum|a_n|+|b|$. The opposite inequality is the triangle inequality.
\end{proof}

% Further source families are proved explicitly below; these sections can be
% moved within this part without changing their stable labels.
% Further Part IV chapters are appended below in dependency order; forward
% references to interpolation do not use the full Young inequality in its proof.

\begin{theorem}[Radon--Nikodym and the Hilbert route]
\label{fa2:p4:radon-nikodym}
If $\mu$ is $\sigma$-finite and $\nu$ is a finite signed or complex measure
with $\nu\ll\mu$, there is a unique $g\in L^1(\mu)$ such that
$\nu(A)=\int_Ag\,d\mu$ for every measurable $A$.
More generally a $\sigma$-finite positive $\nu\ll\mu$ has a nonnegative
measurable density, possibly nonintegrable on the whole space.
\end{theorem}
\begin{proof}
First let both positive measures be finite and put $\rho=\mu+\nu$.
The functional $h\mapsto\int h\,d\nu$ is bounded on $L^2(\rho)$:
$|\int h\,d\nu|\le\nu(X)^{1/2}\pfnorm h_{L^2(\rho)}$.
Hilbert representation gives a real density $v\in L^2(\rho)$.
Indicator tests show $0\le v\le1$ and $d\nu=v\,d\rho$;
similarly $d\mu=(1-v)d\rho$. Absolute continuity implies
$\nu\{v=1\}=0$, because $\mu\{v=1\}=0$.
Thus $g=v/(1-v)$ off this null set satisfies
$\int_Ag\,d\mu=\int_Av\,d\rho=\nu(A)$.
Decompose finite signed measures into positive and negative parts, and
complex measures into real and imaginary parts. These decompositions and
finite variation are scalar measure results; they also follow by taking
suprema over finite measurable partitions. Apply the positive argument to
each part. For $\sigma$-finite spaces choose disjoint sets of finite
$\mu$- and $\nu$-measure and glue their densities. Uniqueness follows
by testing the real and imaginary positive-level sets of the difference.
For finite $\nu$, variation gives $\int|g|=|\nu|(X)<\infty$.
This route uses Hilbert geometry but does not use $L^p$ duality, so the next
proof is not circular. Compare \cite{B08}, Section 7.1.
\end{proof}

\begin{theorem}[$L^p$ duality, including the measure boundary]
\label{fa2:p4:lp-duality}
For any positive measure $\mu$ and $1<p<\infty$, integration is an
isometric linear bijection $L^q(\mu)\to L^p(\mu)'$.
For $p=1$ the same assertion with ordinary $L^\infty$ holds when $\mu$
is $\sigma$-finite. These assertions require neither separability nor
finite total mass. The ordinary $p=1$ statement is not asserted for an
arbitrary measure space.
\end{theorem}
\begin{proof}
On a finite-measure space, set $\nu(A)=L(1_A)$. Countable additivity
follows from convergence of disjoint-union indicators in $L^p$.
For a finite partition, choose phases $z_j$ with
$z_j\nu(A_j)=|\nu(A_j)|$. Then
$\sum|\nu(A_j)|\le\pfnorm L\mu(X)^{1/p}$, proving finite variation.
The preceding theorem supplies a density $g$. Truncated phase tests
$f=\overline g|g|^{q-2}1_{\{1/n\le|g|\le n\}}$ show
$\pfnorm g_q\le\pfnorm L$ if $p>1$; phase indicators of
$\{|g|>\pfnorm L+\varepsilon\}$ give the endpoint bound if $p=1$.
Simple-function density proves the representation for all $f$.
H\"older and the same phase tests give norm equality and uniqueness.
A countable finite-measure exhaustion proves the $\sigma$-finite version
by uniqueness on overlaps and monotone convergence of the $q$-norm.

For arbitrary $\mu$ and $p>1$, let $a(E)$ be the norm of the restriction
of $L$ to functions supported in a $\sigma$-finite measurable set $E$.
Let $M=\sup_Ea(E)$, choose $E_n$ with $a(E_n)\to M$, and put
$F=\bigcup E_n$. On $F$ the preceding construction gives a density of
norm $M$. If $A\supset F$ is $\sigma$-finite, its representing density
$g_A$ agrees on $F$ with this density. Since $q<\infty$ and
$\pfnorm {g_A}_q\le M$, its $q$-integral off $F$ is zero.
Every $L^p$ function has $\sigma$-finite nonzero support, since
$\{|f|>1/n\}$ has finite measure. Applying the preceding observation to
$A=F\cup\{f\ne0\}$ shows $L(f)=L(f1_F)$. Extend the density by
zero. Its norm is detected by its own $\sigma$-finite support, and the
same test proves uniqueness. This last support argument fails at
$q=\infty$; no endpoint assumption has been silently removed.
\end{proof}

\begin{theorem}[Riesz--Markov, precise scalar-linear version]
\label{fa2:p4:riesz-markov}
Let $X$ be locally compact Hausdorff. Every positive functional on
$C_c(X,\mathbb R)$ is integration against a unique positive Radon measure.
Every bounded functional on $C_0(X,\mathbb K)$ is integration against a
unique finite signed or complex Radon measure $\nu$, and
$\pfnorm L=|\nu|(X)$. Here Radon means compact-finite, outer regular
on Borel sets, and compact-inner regular on open sets. Finite Radon
measures are compact-inner regular on all Borel sets. For compact $X$,
$C_0(X)=C(X)$.
\end{theorem}
\begin{proof}
For a positive functional $I$ define
\[
 c(U)=\sup\{I(f):f\in C_c,\ 0\le f\le1,\ \pfsupp f\subset U\},
 \qquad \mu^*(A)=\inf_{U\supset A}c(U),
\]
where $U$ is open. Compact cutoffs exist by local compactness and the
compact Hausdorff separation lemma. A finite partition of unity on the
support of $f$ proves $c(U)\le\sum c(U_j)$ whenever
$U\subset\bigcup U_j$. Consequently $\mu^*$ is an outer measure and
$\mu^*(U)=c(U)$. For an open $G$ and an outer-finite $A$, choose
$U\supset A$ with $c(U)<\mu^*(A)+\varepsilon$. If
$0\le f\le1$ has compact support $K\subset U\cap G$, then adding
tests supported in $U\setminus K$ gives
$c(U)\ge I(f)+\mu^*(A\setminus G)$. Taking suprema proves the
Carath\'eodory inequality for $G$. Thus all Borel sets are measurable.
Cutoffs bound compact masses and prove open inner regularity; outer
regularity is built in. If $h\in C_c$ is globally nonnegative and
$h\ge1$ on $A$, then
$\mu^*(A)\le I(h)$, by testing on $\{h>1-\eta\}$ and letting
$\eta\downarrow0$.

For $f\ge0$ supported in compact $K$, split its range into intervals of
length $\eta$ and write $f=\sum_j\min(\eta,(f-(j-1)\eta)_+)$.
Both $I(f)$ and $\int f\,d\mu$ lie between the same lower and upper
level sums, whose difference is at most $\eta\mu(K)$.
Let $\eta\downarrow0$. This proves representation, and the formula for
$c(U)$ proves uniqueness. Positive functionals on $C_0$ correspond to
finite mass, with norm equal to that mass.

For a real bounded $L$ put
$L^+(f)=\sup_{0\le g\le f}L(g)$ for $f\ge0$.
The decomposition $g=\min(g,f_1)+(g-f_1)_+$ proves additivity on the
positive cone; hence $L^+$ extends to a positive bounded functional,
and so does $L^+-L$. This gives signed representation. Real and imaginary
parts give complex representation. To prove the exact norm, approximate
a finite measurable partition which nearly realizes $|\nu|(X)$ by
disjoint compact subsets, and choose continuous cutoffs with disjoint
supports and the appropriate constant phases. Their sum has norm at
most one and its integral is within any prescribed error of $|\nu|(X)$.
The upper bound is the variation inequality. These tests also prove
uniqueness for signed and complex measures. This is the same positive
content construction as D11.19, not an appeal to an unnamed representation
principle; see also \cite{B08,B09}.
\end{proof}

\begin{proposition}[The bounded-variation bridge on an interval]
\label{fa2:p4:stieltjes-bridge}
For a finite scalar Borel measure $\nu$ on $[a,b]$, define on $\mathbb R$
$F(t)=\nu([a,b]\cap(-\infty,t])$. Then $F$ is right-continuous,
has bounded variation, is zero for $t<a$, is constant for $t\ge b$,
and $\nu((s,t])=F(t)-F(s)$ after extending $\nu$ by zero.
Its total variation on $\mathbb R$ is $|\nu|([a,b])$.
Conversely every such normalized right-continuous BV function defines
a unique finite measure. For continuous $f$ the functional
$\int f\,d\nu$ is the corresponding Riemann--Stieltjes functional,
with both endpoint atoms retained by this extension convention.
\end{proposition}
\begin{proof}
Continuity from above of finite measures gives right continuity.
Interval partitions show the variation bound. Conversely finite unions
of half-open intervals have an additive increment function; right
continuity gives continuity at decreasing empty sequences, hence a
premeasure. The scalar extension theorem supplies the measure.
The variation identity follows first for real monotone differences,
then for complex increments using phase partitions and regularity.
Uniform continuity of $f$ bounds the difference of two tagged sums by
its modulus of continuity times $|\nu|([a,b])$. Their limit agrees with
the simple step-function integral. Normalization excludes additive
constants, and zero extension retains an atom at $a$; a convention on
$(a,b]$ alone would lose it. Left-continuous versions use $[s,t)$
instead and must be converted before comparing source formulas.
\end{proof}

\begin{corollary}[The literal normalized Riemann--Stieltjes form]
\label{fa2:p4:riemann-stieltjes-endpoint}
Let $a<b$ and let $\nu$ be a finite scalar measure on $[a,b]$.
For standard tagged Riemann--Stieltjes sums on this interval use
\[
 w(a)=0,\qquad w(t)=\nu([a,t])\quad(a<t\le b).
\]
Then $w$ has bounded variation, $\operatorname{Var}_{[a,b]}w
=|\nu|([a,b])$, and for every $f\in C[a,b]$,
\[
 \int_a^bf\,dw=\int_{[a,b]}f\,d\nu.
\]
Thus a continuous functional $L$ on $C[a,b]$ has this
normalized representation with $\operatorname{Var}w=\|L\|$.
An atom at $a$ produces the right jump
$w(a+)-w(a)=\nu(\{a\})$. Right continuity is imposed inside
$(a,b)$, and is not required at the normalized left endpoint.
\end{corollary}
\begin{proof}
For a partition $a=t_0<\cdots<t_n=b$, its first increment is
$w(t_1)-w(a)=\nu([a,t_1])$ and the later ones are
$\nu((t_{j-1},t_j])$. These disjoint cells cover $[a,b]$,
so their absolute increments sum to at most $|\nu|([a,b])$.
For tags $\xi_j\in[t_{j-1},t_j]$, the sum is the integral
of the corresponding cellwise step function. Its difference
from $\int f\,d\nu$ is at most
$\omega_f(\max_j(t_j-t_{j-1}))|\nu|([a,b])\to0$.
This proves existence and the standard Riemann--Stieltjes identity.
Every sum for $\|f\|_\infty\le1$ is bounded by
$\operatorname{Var}w$. Take limits and then the supremum in
the proved Riesz--Markov norm identity to obtain the reverse bound
$|\nu|([a,b])\le\operatorname{Var}w$.
Continuity of finite measures gives the endpoint limits.
The first cell thus includes the atom at $a$; using the global
right-continuous $F$ unchanged at $a$ in standard interval sums
would lose it. This gives the literal normalized form of
\cite{B03}, Theorem 4.4-1.
\end{proof}

\section{Weak compactness in $L^1$}

\begin{lemma}[A weakly null $L^1$ sequence is uniformly integrable]
\label{fa2:p4:l1-gliding-hump}
On a $\sigma$-finite space, if $f_n\rightharpoonup0$ in $L^1$, then
the integrals of $|f_n|$ are uniformly absolutely continuous on
small-measure sets and uniformly small outside a suitable finite-measure
set. The latter assertion is essential on an infinite-measure space.
\end{lemma}
\begin{proof}
If uniform absolute continuity fails, choose a subsequence, still
denoted $f_n$, and sets $A_n$ with $\int_{A_n}|f_n|>\varepsilon$.
They can be chosen with summable measures and with each future union
$\bigcup_{j>n}A_j$ having $|f_n|$-integral below $\varepsilon/8$:
at each step impose the finitely many absolute-continuity bounds of
previous functions, with geometrically summable error allocations.
Inductively construct measurable $|g_n|\le1$ by leaving $g_{n-1}$
unchanged off $A_n$ and putting on $A_n$ a constant phase times
$\overline{f_n}/|f_n|$. Choose that phase so that the new integral
has modulus at least $\int_{A_n}|f_n|$ (rotate the contribution on
$A_n$ to align with the outside contribution). Almost every point
belongs to only finitely many $A_n$, so $g_n\to g$ almost everywhere,
$|g|\le1$. Subsequent changes alter $\int f_ng_n$ by at most
$2\int_{\cup_{j>n}A_j}|f_n|<\varepsilon/4$. Hence
$|\int f_ng|>3\varepsilon/4$, contradicting weak convergence.
To force indices strictly increasing, observe that a finite family of
integrable functions is uniformly absolutely continuous.

For tails take an increasing finite-measure exhaustion $E_k$.
If tails fail uniformly, repeat the construction with $A_n=X\setminus
E_{k_n}$ and choose $k_n$ so large that all previous functions have
arbitrarily small integrals on $A_n$. Choose the selected index beyond
all previous ones. These sets decrease to the empty set modulo a null
set, so $g_n$ again stabilizes almost everywhere. The same contradiction
proves tail control. For a weakly convergent sequence subtract its limit;
a single integrable function has both properties.
\end{proof}

\begin{theorem}[Dunford--Pettis in scalar $L^1$]
\label{fa2:p4:dunford-pettis-l1}
Let $\mu$ be $\sigma$-finite. A set $\mathcal A\subset L^1(\mu)$
has weakly compact closure if and only if it is bounded and
\[
 \sup_{f\in\mathcal A}\int_E|f|\longrightarrow0\quad(\mu(E)\to0),
 \qquad
 \inf_{\mu(F)<\infty}\sup_{f\in\mathcal A}\int_{X\setminus F}|f|=0.
\]
\end{theorem}
\begin{proof}
If the closure is weakly compact, boundedness follows from uniform
boundedness. Failure of either displayed condition gives a sequence
violating the preceding lemma for every subsequence chosen by that
failure. Eberlein--\v Smulian in Part I supplies a weakly convergent
subsequence, a contradiction.

Conversely embed $L^1$ in $(L^\infty)'$ by integration and take its
weak-star compact closed hull inside a dual ball. Every member $L$ of
this closure satisfies $|L(1_E)|\le\sup_{f\in\mathcal A}\int_E|f|$.
If $E_n\downarrow\varnothing$, choose finite-measure $F$ with uniformly
small tails. Since $\mu(E_n\cap F)\downarrow0$, the first condition
then makes $L(1_{E_n})\to0$. Finite additivity plus this continuity
gives a countably additive finite scalar measure $\nu(E)=L(1_E)$.
It vanishes on null sets; variation is bounded by the dual-ball radius,
using finite phase partitions. Radon--Nikodym gives $\nu=g\mu$ with
$g\in L^1$. Uniform approximation by bounded simple functions proves
$L(h)=\int gh$ for every $h\in L^\infty$. Thus the entire compact
closure lies in the embedded $L^1$, and its inherited weak-star
topology is precisely $\sigma(L^1,L^\infty)$. This proves relative
weak compactness with both required tightness conditions.
\end{proof}

\section{Bochner integration}

\begin{definition}[Strong measurability and Bochner integrability]
\label{fa2:p4:bochner-definition}
A Banach-valued function is strongly measurable if it is an almost
everywhere pointwise norm limit of measurable finite-valued simple
functions. It is Bochner integrable if it is an $L^1$-norm limit of
integrable simple functions. For $s=\sum_j1_{A_j}x_j$, on disjoint sets
of finite measure, put $\int s=\sum_j\mu(A_j)x_j$.
\end{definition}

\begin{theorem}[Construction and convergence of the Bochner integral]
\label{fa2:p4:bochner}\label{fa2:p4:dominated-convergence}
\label{fa2:p4:lp-completeness}
A strongly measurable $f$ is Bochner integrable exactly when
$\int\pfnorm f<\infty$. Its integral is the unique continuous
extension of simple integration, with
$\pfnorm{\int f}\le\int\pfnorm f$.
Bounded linear maps commute with the integral. If $f_n\to f$ almost
everywhere in norm and $\pfnorm{f_n}\le g\in L^1$, then
$\int\pfnorm{f_n-f}\to0$ and $\int f_n\to\int f$.
For $1\le p\le\infty$ the strongly measurable $L^p(X;E)$ is Banach.
For $E=\mathbb K$ this proves scalar $L^p$ completeness on every
positive measure space: each strongly measurable scalar function
has an ordinary measurable representative and conversely, so the
almost-everywhere quotient is the usual scalar one.
\end{theorem}
\begin{proof}
Simple integration is independent of the chosen partition, and its
norm is bounded by its scalar $L^1$ norm. Completeness of $E$ therefore
defines a limit independent of approximants. A strongly measurable
function is essentially separably valued: all values of the approximating
simple functions lie in one separable closed subspace. Within that
subspace choose dense centers $a_1=0,a_2,\ldots$ and let $s_n(x)$
be a nearest center among $a_1,\ldots,a_n$, breaking ties by the first
index. Distances are measurable, so these functions are finite-valued
and measurable. Density gives $s_n(x)\to f(x)$, while the zero center
gives $\pfnorm{s_n-f}\le\pfnorm f$ and
$\pfnorm{s_n}\le2\pfnorm f$. Thus they are integrable when
$\int\pfnorm f<\infty$, and dominated convergence gives their
$L^1$ convergence. This is not a truncation of an independently
varying countable partition.
Conversely an $L^1$-Cauchy sequence of simple functions has a subsequence
with summable norms of successive differences. Scalar Tonelli makes
the differences norm-summable almost everywhere; completeness gives
the pointwise limit and $L^1$ convergence. This proves the criterion.
The same argument proves $L^p$ completeness for finite $p$, choosing
successive $L^p$ errors summable and using Minkowski on their norms;
for $p=\infty$ remove the countable union of exceptional null sets
and use pointwise uniform convergence. Strong measurability survives
these limits by a diagonal choice of simple approximants.
Bounded maps commute with simple integration and then with limits.
The asserted convergence is scalar dominated convergence applied to
$\pfnorm{f_n-f}\le2g$, followed by the integral norm inequality.
\end{proof}

\begin{proposition}[Pettis measurability criterion and weak integrals]
\label{fa2:p4:pettis-measurability}
On a complete measure space, $f:X\to E$ is strongly measurable if and
only if it is weakly measurable and essentially takes values in a
separable subspace. A Bochner integral is also a weak integral:
$\ell(\int f)=\int\ell(f)$ for every $\ell\in E'$.
Weak integrability alone does not imply Bochner integrability.
\end{proposition}
\begin{proof}
Necessity follows from simple approximants. On a separable subspace
choose a countable norming set of functionals: norm each vector in a
countable dense subset by Hahn--Banach. Then
$\pfnorm{x-y}=\sup_j|\ell_j(x-y)|$, so inverse images of norm balls
are measurable. Choosing a nearest center among the first $n$ dense
centers, with deterministic tie breaking, gives finite-valued measurable
approximants converging pointwise in norm, as in the preceding proof.
Completeness makes changes on the exceptional null set legitimate.
The weak integral identity was proved for bounded linear maps.
For the last assertion let $X=\mathbb N$ with counting measure and
$f(n)=e_n/n$ in $\ell^2$. Every $\ell\in(\ell^2)'$ gives an absolutely
summable scalar sequence by Cauchy--Schwarz, and
$\sum e_n/n\in\ell^2$ represents the weak integral. But
$\sum\pfnorm{f(n)}=\sum1/n=\infty$, so this function is not Bochner
integrable. This also identifies exactly which missing condition is
responsible.
\end{proof}

% The general reflexive vector-valued duality is deliberately not declared
% proved here until the vector Radon--Nikodym bridge has been written and audited.

\chapter{Test Functions and Distributions}
\label{fa2:p4:distributions}
\section{The topology before the dual}

\begin{definition}[The test-function LF topology]
\label{fa2:p4:test-lf}
Let $\Omega\subset\mathbb R^d$ be open. For compact $K\subset\Omega$,
let $\pfD_K$ be the smooth functions supported in $K$, with seminorms
$p_m(\phi)=\max_{|\alpha|\le m}\sup_x|\partial^\alpha\phi(x)|$.
Choose compact $K_j\subset\operatorname{int}K_{j+1}$ exhausting
$\Omega$. On $\pfD(\Omega)=\bigcup_j\pfD_{K_j}$ take the finest
locally convex topology making all inclusions continuous. Equivalently,
its continuous seminorms are exactly the seminorms whose restrictions
to every $\pfD_{K_j}$ are continuous. This is the strict LF topology.
It is a topology defined by seminorms, not a definition by convergent
sequences. A distribution is a continuous scalar-linear functional on
this space, denoted $u\in\pfD'(\Omega)$.
\end{definition}

\begin{theorem}[Bounded sets, continuity, and distribution order]
\label{fa2:p4:test-bounded}\label{fa2:p4:finite-order}
The subspace topology on $\pfD_K$ is its displayed Fr\'echet topology.
A set in $\pfD(\Omega)$ is bounded if and only if it is contained in
one $\pfD_K$ and all its $p_m$ are bounded. A linear functional $u$
is a distribution if and only if for every compact $K\subset\Omega$
there are $C_K<\infty$ and $m_K\in\mathbb N$ such that
\[
 |u(\phi)|\le C_Kp_{m_K}(\phi)\quad(\phi\in\pfD_K).
\]
A sequence tends to zero in $\pfD$ exactly when its supports lie in
one compact set and every derivative tends uniformly to zero there.
\end{theorem}
\begin{proof}
Each global seminorm $p_m$ is continuous on every $\pfD_{K_j}$ and
hence in the LF topology. This and continuity of the inclusion prove
the topology assertion. Completeness of $\pfD_K$ follows by integrating
uniform limits of successive derivatives along line segments in local
balls; the limits remain compatible derivatives, and their supports
stay in $K$.
If a bounded set $B$ had no common compact support, choose
$\phi_j\in B$ and $x_j\notin K_j$ with $\phi_j(x_j)\ne0$,
passing to a subsequence so that the points escape every compact set.
Choose $c_j>0$ with $c_j|\phi_j(x_j)|>j$. The seminorm
$p(\phi)=\sup_jc_j|\phi(x_j)|$ is finite and continuous on each
$\pfD_K$, since only finitely many $x_j$ meet $K$.
It is therefore LF-continuous but unbounded on $B$, a contradiction.
The $p_m$ then give the derivative bounds. Conversely, a set bounded
in one $\pfD_K$ is bounded under every continuous linear inclusion.

The seminorm criterion for a linear functional follows directly from
the final locally convex topology; a continuous seminorm on a
Fr\'echet step is bounded by finitely many of its increasing $p_m$.
For completeness of this use of the final topology, if $u$ is
continuous on every step, $\phi\mapsto|u(\phi)|$ is one of its
defining seminorms, so $u$ is continuous globally.
Every convergent sequence together with its limit is bounded in a TVS;
the bounded-set result and the subspace-topology result now give the
sequential statement. It is a consequence of the topology, and does
not replace its definition. Compare \cite{B04}, Sections 6.2--6.8.
\end{proof}

\begin{proposition}[Limits and calculus]
\label{fa2:p4:distribution-calculus}
Locally integrable functions and locally finite scalar Radon measures
define distributions by integration. Set
$\partial^\alpha u(\phi)=(-1)^{|\alpha|}u(\partial^\alpha\phi)$ and
$(au)(\phi)=u(a\phi)$ for $a\in C^\infty(\Omega)$.
These are distributions; the Leibniz rule holds. If $u_j(\phi)$ has
a limit for every test function, the limit is a distribution.
If $u_j\to u$ in this weak distribution sense and $a_j\to a$ in
$C^\infty_{\rm loc}$, then $a_ju_j\to au$.
\end{proposition}
\begin{proof}
Local integrability bounds the integral by a local $L^1$ norm times
$p_0$. Differentiation and multiplication satisfy finite-order
seminorm estimates by the ordinary finite Leibniz formula; its dual
gives the distribution Leibniz formula. On each Fr\'echet step the
family $u_j$ is pointwise bounded. Uniform boundedness for Fr\'echet
spaces gives $|u_j(\phi)|\le Cp_m(\phi)$, uniformly in $j$.
Pass to the limit and use the preceding theorem. Finally, on a fixed
compact support, $a_j\phi\to a\phi$ in every seminorm. The uniform
bound just obtained makes
$u_j((a_j-a)\phi)\to0$, while $u_j(a\phi)\to u(a\phi)$.
All pairings remain linear, so no extra conjugation enters these rules.
\end{proof}

\begin{lemma}[Smooth partitions of unity]
\label{fa2:p4:partition-unity}
Every open cover of an open subset of $\mathbb R^d$ admits a countable
locally finite subordinate smooth partition of unity, with compactly
supported nonnegative members.
\end{lemma}
\begin{proof}
Take compact exhaustion $K_j$ as above. For each compact shell
$K_j\setminus\operatorname{int}K_{j-1}$ choose finitely many balls
whose doubled closed balls lie in a member of the cover and in
$\operatorname{int}K_{j+1}\setminus K_{j-2}$.
These balls form a locally finite refinement. On each doubled ball
choose a nonnegative smooth bump positive on the smaller ball,
using $\exp(-1/(1-|x|^2))$ and rescaling. Their locally finite sum
$h$ is smooth and positive. Dividing each bump by $h$ gives the
partition. Compactness of each doubled closed ball gives compact
support and local finiteness justifies every derivative.
\end{proof}

\begin{theorem}[Localization, support, and gluing]
\label{fa2:p4:distribution-gluing}
Distributions on overlapping open sets that agree on all intersections
glue uniquely to a distribution on their union. The support of $u$
is the complement of the union of open sets where it vanishes; $u$
vanishes on that union. A compactly supported distribution extends
uniquely to a continuous functional on $C^\infty(\Omega)$ by
$u(f)=u(\chi f)$ with $\chi=1$ near its support.
Conversely every continuous functional on $C^\infty(\Omega)$ has
compact support and defines such a distribution.
\end{theorem}
\begin{proof}
For a subordinate partition $\eta_j$ define
$u(\phi)=\sum_j u_j(\eta_j\phi)$. Only finitely many terms meet
$\pfsupp\phi$; on each compact this finite sum obeys a finite-order
bound. A second partition gives the same result after inserting their
products and using compatibility. This also proves uniqueness and
the assertion about the union of vanishing opens. Cutoffs equal to
one near the support give the same pairing because their difference
vanishes near that support; use a finite partition on the remaining
test support to justify this statement. The finite-order estimate
for $\chi f$ proves continuity on $C^\infty$.
Conversely continuity of $L$ on $C^\infty$ gives a bound by finitely
many derivative seminorms on a compact $K$. Every test supported off
$K$ is annihilated. Thus its distribution support lies in $K$.
This also shows why merely vanishing pointwise on the support is
insufficient for multiplication: $x\delta'=-\delta$ although
$x=0$ on $\pfsupp\delta'=\{0\}$.
\end{proof}

\begin{theorem}[Point support and local continuous primitives]
\label{fa2:p4:point-support}\label{fa2:p4:distribution-structure}
A distribution supported at $a$ is a finite linear combination of
$\partial^\alpha\delta_a$. Every distribution is locally a finite
sum of derivatives of continuous functions. If its support is compact
and contained in an open $V\Subset\Omega$, these continuous functions
can be chosen with compact support in $V$. Globally one obtains a
locally finite sum of such derivatives.
\end{theorem}
\begin{proof}
For point support choose an order bound $m$ on a compact neighborhood
of $a$. If the derivatives of $\phi$ through order $m$ vanish at $a$,
Taylor's formula and a shrinking cutoff $\chi((x-a)/\varepsilon)$
give $p_m(\chi_\varepsilon\phi)\to0$.
Since $u(\phi)=u(\chi_\varepsilon\phi)$, it is zero. Thus $u$ depends
only on the finite Taylor jet and is a linear combination of jet
evaluations, equivalently delta derivatives.

For compact support, a fixed cutoff first gives
$|u(\phi)|\le C\max_{|\alpha|\le m}\pfnorm{\partial^\alpha\phi}_\infty$
with derivatives restricted to a fixed compact neighborhood.
Map $\phi$ into the finite product of continuous derivative functions
on that compact. Hahn--Banach extends $u$ to this product, and
Riesz--Markov represents its coordinates by finite compactly supported
measures $\mu_\alpha$. Thus
$u=\sum_{|\alpha|\le m}\partial^\alpha\mu_\alpha$, absorbing signs
in the measures. Define
$G(x)=\prod_{j=1}^d(x_j)_+$; then
$\partial_1^2\cdots\partial_d^2G=\delta_0$.
Each $F_\alpha=G*\mu_\alpha$ is continuous: the measure is supported
in a compact set and $G(x-\cdot)$ varies uniformly there when $x$
varies in a compact set. Consequently
$u=\sum\partial^{\alpha+(2,\ldots,2)}F_\alpha$.
Multiply by a cutoff equal to one near $\pfsupp u$ and supported in
$V$. The identity
\[
 \chi\partial^\gamma F
 =\sum_{\beta\le\gamma}(-1)^{|\gamma-\beta|}
 {\gamma\choose\beta}\partial^\beta
       (F\partial^{\gamma-\beta}\chi)
\]
follows by the Leibniz formula and proves the compact-support
assertion. A locally finite partition of unity reduces a general
distribution to compactly supported pieces, and the same argument
in locally finite enlarged patches proves the final assertion.
These are actual continuous primitives, rather than a formal claim
that distributions may be differentiated arbitrarily.
\end{proof}

\section{Convolution and regularization}

\begin{theorem}[Young inequalities and approximate identities]
\label{fa2:p4:young}\label{fa2:p4:approximate-identity}
If $1\le p,q,r\le\infty$ and $1+1/r=1/p+1/q$, then
$\pfnorm{f*g}_r\le\pfnorm f_p\pfnorm g_q$.
For $1\le p<\infty$, translations are norm-continuous on $L^p$.
If $\rho\in C_c^\infty$ is nonnegative with integral one and
$\rho_\varepsilon(x)=\varepsilon^{-d}\rho(x/\varepsilon)$, then
$\rho_\varepsilon*f\to f$ in $L^p$. In particular $C_c^\infty$
is dense in $L^p$. For continuous $f$ the convergence is locally
uniform, and is uniform globally when $f$ is uniformly continuous.
\end{theorem}
\begin{proof}
The $L^1*L^p$ inequality follows from Minkowski's integral inequality,
which itself follows by scalar duality from H\"older (or directly for
$p=1,\infty$). The $L^p*L^{p'}\to L^\infty$ bound is H\"older.
For fixed $g\in L^q$, interpolate these two endpoint estimates for
the map $f\mapsto f*g$ by the Riesz--Thorin theorem proved below;
the reciprocal exponents give the asserted Young relation. This
proof dependency is forward only to the independent interpolation
proof; that proof uses scalar duality and the three-lines lemma,
not Young's full inequality.
Translations are continuous first for compactly supported continuous
functions by uniform continuity and a common bounded support.
Such functions are dense in $L^p$: truncate the function and its
support, approximate by measurable simple functions, approximate
their finite-measure sets by finite unions of boxes using scalar
Lebesgue regularity, and round the box boundaries by continuous
cutoffs. The translation isometry extends the conclusion to all
$L^p$. Then
\[
 \pfnorm{\rho_\varepsilon*f-f}_p
 \le\int\rho_\varepsilon(y)\pfnorm{f(\cdot-y)-f}_p\,dy\to0.
\]
Truncate $f$ first and convolve to obtain compactly supported smooth
approximants; differentiating the compact convolution is justified by
dominated convergence. The continuous versions use the same estimate
with a modulus of continuity.
\end{proof}

\begin{theorem}[Distribution convolution]
\label{fa2:p4:distribution-convolution}
For $u\in\pfD'(\mathbb R^d)$ and $\phi\in\pfD$ put
$(u*\phi)(x)=u_y(\phi(x-y))$. This is smooth and
$\partial^\alpha(u*\phi)=(\partial^\alpha u)*\phi
=u*(\partial^\alpha\phi)$.
For distributions $u,v$ convolution is defined whenever the addition
map is proper on $\pfsupp u\times\pfsupp v$; in particular one compact
support suffices. It is commutative, derivatives commute with it, and
$\pfsupp(u*v)\subset\pfsupp u+\pfsupp v$.
Associativity holds when both inner convolutions exist by the pair-support condition
and triple addition is proper, in particular when two factors have compact support.
Finally $u*\rho_\varepsilon\to u$ in distributions.
\end{theorem}
\begin{proof}
On a compact set of $x$ the translated test functions have common
compact support in $y$. Taylor's formula with remainder in every
test seminorm permits differentiation through $u$, giving smoothness
and the derivative identities. On a compact support of an outer test
$\psi$, properness lets us insert a compact cutoff equal to one on
the part of $\pfsupp u\times\pfsupp v$ whose sum meets that support.
Define
\[
 (u*v)(\psi)=u_xv_y\bigl(\chi(x,y)\psi(x+y)\bigr).
\]
The iterated functional on compactly supported smooth functions is
independent of order: approximate such a function, together with all
derivatives required by the two finite-order bounds, by finite sums
$a_j(x)b_j(y)$. A tensor-product partition into small cubes followed
by Taylor polynomials provides these approximants, and on finite
products the assertion is algebraic. The same approximation proves
continuity of the tensor functional. Changing the cutoff adds a test
vanishing near the product support and contributes zero.
The support assertion follows because tests supported off the sum
vanish near that product support. Commutativity, derivatives, and
associativity now follow from the tensor definition and the identical
identities for finite products; properness is what permits all cutoffs.
Lastly, for each test $\psi$,
$(u*\rho_\varepsilon)(\psi)=u(\widetilde\rho_\varepsilon*\psi)$,
and the tests on the right tend to $\psi$ in one $\pfD_K$.
There is no assertion that two arbitrary distributions can be convolved.
\end{proof}

\chapter{Fourier Transform, Exponential Type, and Fundamental Solutions}
\label{fa2:p4:fourier-chapter}

\section{Schwartz space and the $L^2$ extension}
\begin{definition}[Fourier normalization]
\label{fa2:p4:fourier-normalization}
For $f\in L^1(\mathbb R^d)$ use
$\widehat f(\xi)=\int e^{-2\pi i x\cdot\xi}f(x)\,dx$.
The Schwartz space $\pfS$ has seminorms
$\sup_x|x^\alpha\partial^\beta f(x)|$.
Its continuous scalar-linear dual is $\pfS'$, the tempered
distributions. Set $\widehat u(\phi)=u(\widehat\phi)$.
Thus $\widehat{\partial_j u}=2\pi i\xi_j\widehat u$ and
$\widehat{x_ju}=-(2\pi i)^{-1}\partial_j\widehat u$.
\end{definition}

\begin{theorem}[Fourier inversion and Plancherel]
\label{fa2:p4:fourier-inversion}\label{fa2:p4:plancherel}
The Fourier transform is a continuous automorphism of $\pfS$,
$\widehat{\widehat f}(x)=f(-x)$, and
$f(x)=\int e^{2\pi i x\cdot\xi}\widehat f(\xi)\,d\xi$.
It extends uniquely to a unitary map on $L^2$ and to an automorphism
of $\pfS'$. For $f\in L^1\cap L^2$ the $L^2$ extension agrees with
the integral transform. Multiplication and differentiation obey the
displayed normalization; convolution of Schwartz functions transforms
to pointwise multiplication.
\end{theorem}
\begin{proof}
Integration by parts gives
$\xi^\alpha\partial_\xi^\beta\widehat f$
as a fixed constant times the transform of
$\partial_x^\alpha(x^\beta f)$. Its supremum is bounded by the
$L^1$ norm of that Schwartz function, itself bounded by finitely many
Schwartz seminorms. Hence the transform is continuous into $\pfS$.
In one dimension direct completion of squares, or differentiating
the Gaussian integral and solving the resulting differential equation,
gives $\widehat{e^{-\pi x^2}}=e^{-\pi\xi^2}$; products give dimension
$d$. For $\varepsilon>0$, Fubini therefore gives
\[
 \int e^{2\pi i x\cdot\xi}\widehat f(\xi)
        e^{-\pi\varepsilon|\xi|^2}\,d\xi
 =\int f(y)\varepsilon^{-d/2}e^{-\pi|x-y|^2/\varepsilon}\,dy.
\]
The right side tends to $f(x)$ by the approximate-identity argument;
the left side tends to its undamped integral by dominated convergence.
This proves inversion and the fourth-power identity. Apply inversion
to $g$ and Fubini to obtain
$\int f\overline g=\int\widehat f\overline{\widehat g}$ for
Schwartz functions. Smooth compact-support density in $L^2$ extends
this isometry to $L^2$; the inverse extends by the same construction,
so it is onto. If $f\in L^1\cap L^2$, approximate in both norms by
smooth compact-support functions. Their transforms converge uniformly
to the integral transform and in $L^2$ to the extension, which proves
agreement. The continuous inverse on $\pfS$ transposes to an inverse
on $\pfS'$. All differential identities follow first by integration
by parts and then by the definition of distribution derivatives.
Fubini proves the convolution identity. No separability of a Hilbert
space was used to replace this function-space calculation.
\end{proof}

\begin{proposition}[Tempered examples and compact support]
\label{fa2:p4:tempered-examples}
Every compactly supported distribution is tempered. A locally
integrable function of at most polynomial growth in the integral
sense defines a tempered distribution. In particular
$\widehat1=\delta_0$, $\widehat{\delta_0}=1$, and polynomials
transform to the corresponding delta derivatives. Not every
distribution is tempered.
\end{proposition}
\begin{proof}
For compact support insert a fixed compact cutoff in the finite-order
bound; the resulting seminorm is bounded by finitely many Schwartz
seminorms. If $\int(1+|x|)^{-N}|f(x)|\,dx<\infty$, bound its
pairing by that integral times $\sup(1+|x|)^N|\phi(x)|$.
Fourier inversion at zero proves the two basic identities, and the
calculus gives the polynomial formulas. The function $e^{x^2}$ on
$\mathbb R$ is locally integrable but not tempered: translate a fixed
nonnegative bump to $n$ and multiply it by $e^{-n^2/2}$.
These tests tend to zero in every Schwartz seminorm, whereas their
pairings with $e^{x^2}$ tend to infinity.
\end{proof}

\section{Paley--Wiener and entire divisibility}

\begin{theorem}[Paley--Wiener, smooth and distribution forms]
\label{fa2:p4:paley-wiener}
For a nonempty compact convex $K\subset\mathbb R^d$ put
$H_K(\eta)=\sup_{x\in K}x\cdot\eta$.
The Fourier--Laplace transform of $f\in C_c^\infty$ supported in
$K$ is entire and satisfies, for every $N$,
\[
 |\widehat f(z)|\le C_N(1+|z|)^{-N}e^{2\pi H_K(\operatorname{Im}z)}.
\]
Conversely every entire function with these bounds is the transform
of such an $f$. The transforms of distributions supported in $K$
are exactly the entire functions with a bound
\[
 |F(z)|\le C(1+|z|)^M e^{2\pi H_K(\operatorname{Im}z)}
\]
for some $C,M$. The exponent $M$ need not equal the least order of
the distribution.
\end{theorem}
\begin{proof}
Compact support permits termwise complex differentiation of the
integral. Integration by parts in a coordinate with
$|z_j|\ge|z|/\sqrt d$ gives arbitrary inverse powers of $|z|$ and
the displayed exponential factor. For a compactly supported
distribution, apply its finite-order bound to a cutoff around $K$
times $e^{-2\pi i xz}$. Choose cutoff width
$(1+|z|)^{-1}$; its derivatives cost only polynomial powers and
$H_{K+\delta B}(\operatorname{Im}z)
\le H_K(\operatorname{Im}z)+\delta|z|$.
This proves the polynomial bound and entire differentiability.

For the smooth converse, define $f$ by the inverse Fourier integral
on $\mathbb R^d$. The rapid bounds and differentiation under that
integral make $f$ smooth. Cauchy's theorem allows the contour shift
$\mathbb R^d\mapsto\mathbb R^d+it\eta$: shift each coordinate on
a rectangle, and use a decay power exceeding $d+1$ to make its side
integrals tend to zero. If $x\notin K$, finite-dimensional separation
gives $\eta$ with $x\cdot\eta>H_K(\eta)$. On the shifted contour
the inverse integral is bounded by a polynomial in $t$ times
$e^{-2\pi t(x\cdot\eta-H_K(\eta))}$, so $f(x)=0$.
This proves the support assertion without assuming an inverse theorem
for entire functions.

For a polynomial-growth entire $F$, its restriction to the real axes
defines a tempered distribution; let $u=\pfF^{-1}F$.
Multiply $F$ by $\widehat\rho_\varepsilon$, with
$\rho_\varepsilon$ a compactly supported smooth approximate identity.
The product has rapid bounds of type $K+\varepsilon B$, so the smooth
case gives $\pfsupp(u*\rho_\varepsilon)\subset K+\varepsilon B$.
Passing to distributions yields $\pfsupp u\subset K$.
Finally its entire transform and $F$ agree on the real axes and
hence everywhere: apply the one-variable identity theorem successively
in each variable. This also proves uniqueness. For nonconvex compact
support the exponential bound records its convex hull; it must not
be interpreted as recovering holes in the support.
\end{proof}

\begin{lemma}[Polynomial division preserves exponential type]
\label{fa2:p4:polynomial-division}
If $P$ is a nonzero polynomial and $G$ is entire, a Paley--Wiener
polynomial-growth bound for $PG$ with convex type $K$ implies one for
$G$ with the same $K$.
\end{lemma}
\begin{proof}
Choose a real direction $v$ in which the highest homogeneous part
of $P$ is nonzero. Then $t\mapsto P(z+tv)$ has fixed degree $m$
and fixed nonzero leading coefficient $a$, independent of $z$.
The mean-value inequality for the subharmonic function $\log|G|$
and Jensen's formula for this one-variable polynomial give
\[
 |G(z)|\le |a|^{-1}r^{-m}
      \max_{|t|=r}|P(z+tv)G(z+tv)|.
\]
Indeed the circular mean of $\log|P(z+tv)|$ is
$\log|a|+\sum_j\log\max(r,|t_j|)\ge\log|a|+m\log r$.
Zeros of $G$ at the center cause no difficulty, by a limiting
argument. Take $r=1$ and use subadditivity of $H_K$ and boundedness
of $\operatorname{Im}(tv)$ on that circle. Polynomial growth and
the same exponential type follow. Jensen's polynomial identity here
also follows directly by averaging
$\log|re^{i\theta}-t_j|$; its average is
$\log\max(r,|t_j|)$ by the harmonic mean-value property outside
the circle and the expansion of $\log(1-w)$ inside.
\end{proof}

\begin{theorem}[Compact-support solvability]
\label{fa2:p4:compact-pde-solvability}
Write $D_j=(2\pi i)^{-1}\partial_j$ and let $P\ne0$ be a polynomial.
For compactly supported $v\in\pfD'(\mathbb R^d)$, the equation
$P(D)u=v$ has a compactly supported distributional solution if and
only if $\widehat v=PG$ for an entire $G$.
Such a compactly supported solution is unique, and its support is
contained in the convex hull of $\pfsupp v$.
\end{theorem}
\begin{proof}
For $v=0$ take $u=0$; uniqueness below also covers that case without
assigning a support function to the empty set.
Necessity is the Fourier calculus. For $v\ne0$ apply the division
lemma and Paley--Wiener to $G$ with
$K=\operatorname{conv}(\pfsupp v)$. Their inverse transform gives
the solution and the support statement. If two compact solutions
differ, their entire transforms satisfy $PH=0$. On the nonempty
open set $P\ne0$ they vanish, and the entire identity theorem gives
$H=0$. Real-axis pointwise division alone would not provide the
entire $G$ and is insufficient.
\end{proof}

\begin{theorem}[Malgrange--Ehrenpreis for constant coefficients]
\label{fa2:p4:fundamental-solution}
Every nonzero constant-coefficient scalar linear differential
operator on $\mathbb R^d$ has a distributional fundamental solution.
If $P(D)E=\delta_0$ and $v$ has compact support, $E*v$ is a
distributional solution of $P(D)u=v$.
\end{theorem}
\begin{proof}
Normalize $D_j=(2\pi i)^{-1}\partial_j$ and rotate real coordinates
so that $P(\zeta_1,\xi')$, as a polynomial in $\zeta_1$, has a
constant nonzero leading coefficient. Divide it out. Its degree is
$m$. For each real $\xi'$, define $n(\xi')$ as the first height
in the finite list $-m-1,-m+1,\ldots,m+1$ satisfying
$\inf_{t\in\mathbb Q}|P(t+in,\xi')|\ge1$.
This countable Borel condition makes the finite-valued choice
measurable. Existence follows from a stronger sufficient condition:
some candidate has distance at least one from the imaginary part
of every root, since each root excludes at most one candidate
with strict distance less than one. Monic factorization then
gives the required denominator bound. The two conditions are
not claimed to be equivalent.
Then $|P(t+in(\xi'),\xi')|\ge1$ for every real $t$.
Define
\[
 E(\phi)=\int_{\mathbb R^{d-1}}
    \int_{\mathbb R+in(\xi')}
       \frac{\widehat\phi(-\zeta_1,-\xi')}
            {P(\zeta_1,\xi')}\,d\zeta_1\,d\xi'.
\]
For tests supported in a fixed compact set, integration by parts
in the Fourier integral bounds the numerator, uniformly over this
finite set of heights, by $C_Kp_N(\phi)(1+|t|+|\xi'|)^{-N}$.
Choose $N>d$. Thus the integral is absolutely convergent and defines
a distribution. In $E(P(-D)\phi)$ the polynomial cancels.
For fixed $\xi'$ the remaining entire rapidly decreasing integrand
can be shifted to the real line, since the rectangle side integrals
vanish. Fubini and inversion now give
$E(P(-D)\phi)=\int\widehat\phi(-\xi)\,d\xi=\phi(0)$.
Restore the divided leading coefficient. This proves the theorem
including continuity of the constructed functional, rather than only
formal division by a possibly vanishing real symbol.
The convolution assertion follows from the support-valid convolution
calculus. This is the contour route of \cite{B09}, Theorem 2.11;
the polynomial-division route above corresponds to \cite{B04}, 8.4.
\end{proof}

\chapter{Sobolev Spaces as Function Spaces}
\label{fa2:p4:sobolev-chapter}
\section{Weak derivatives, density, and representatives}

\begin{definition}[Sobolev norm and negative spaces]
\label{fa2:p4:sobolev-definition}
For open $\Omega\subset\mathbb R^d$, integer $m\ge0$, and
$1\le p\le\infty$, let $W^{m,p}(\Omega)$ consist of $L^p$ functions
whose distribution derivatives through order $m$ lie in $L^p$.
Use the norm $\sum_{|\alpha|\le m}\pfnorm{\partial^\alpha u}_p$;
the equivalent Euclidean product norm is used for its Hilbert structure
when $p=2$. For $p<\infty$, $W^{m,p}_0$ is the norm closure of
$C_c^\infty$. No corresponding closure characterization is silently
asserted at $p=\infty$.
\end{definition}

\begin{theorem}[Completeness, local smoothing, and calculus]
\label{fa2:p4:sobolev-calculus}
$W^{m,p}$ is Banach, reflexive for $1<p<\infty$, and separable for
$p<\infty$. Smooth functions in $\Omega$ are dense for finite $p$.
On every compact interior set, convolution smooths and converges
with all the weak derivatives of order at most $m$ in $L^p$ for
finite $p$. At $p=\infty$ it gives uniform derivative bounds,
local $L^r$ convergence for finite $r$, and weak-star derivative
convergence; $L^\infty$ norm convergence is not claimed.
For $m=1$, the weak derivative criterion is equivalent to the existence
of an almost-everywhere representative absolutely continuous on
almost every coordinate line, with the classical line derivatives in
$L^p$. Products of bounded $W^{1,p}$ functions obey the product rule.
If $F\in C^1(\mathbb R)$ has bounded derivative, then
$\nabla F(u)=F'(u)\nabla u$ whenever $F(u)\in L^p$; in particular
$F(0)=0$ suffices. Real truncations and absolute values obey their
almost-everywhere chain rules as well.
\end{theorem}
\begin{proof}
Embed $u$ by its derivative tuple in the finite product of $L^p$
spaces. Its image is closed: $L^p$ limits may be passed through
compact test integrals by H\"older, preserving all derivative identities.
Completeness, reflexivity, and separability follow from the corresponding
finite-product and closed-subspace results; scalar $L^p$ reflexivity
follows by applying duality twice. For local smoothing choose a cutoff
one on a neighborhood of the compact set. On that set
$\partial^\alpha(\rho_\varepsilon*u)
=\rho_\varepsilon*(\partial^\alpha u)$, and the approximate-identity
theorem gives convergence. For global smooth density on arbitrary
$\Omega$, take a locally finite compactly supported partition
$\eta_j$, convolve each $\eta_ju$ at a radius so small that its support
stays in its patch, and its $W^{m,p}$ error is below
$\varepsilon2^{-j}$. The locally finite sum is smooth; the triangle
inequality bounds its total error by $\varepsilon$.

On a box, test against products of a one-variable test and transverse
tests. Fubini and a countable dense family of these tests show that
almost every line has the indicated one-dimensional distribution
derivative. The one-dimensional primitive theorem below then gives
absolute continuity. Conversely integrate one-dimensional integration
by parts over transverse variables. An exhaustion by boxes handles
general $\Omega$.
The product rule is now the ordinary absolutely continuous product
rule on those lines; boundedness ensures both resulting derivative
terms are in $L^p$. For $C^1$ functions the same line rule gives the
chain formula, with bounded derivative providing the $L^p$ estimate.
For $|u|$ and truncations use the elementary fact that an absolutely
continuous function has derivative zero almost everywhere on each
level set. To prove that fact, restrict to density points of the level
set where the derivative exists and take difference quotients along
that set. Apply the ordinary derivative rule off the level set.
This proves also the precise a.e. values of the truncated gradients.
\end{proof}

\begin{theorem}[One-dimensional absolute continuity]
\label{fa2:p4:one-dimensional-sobolev}
On an interval $I$, each $u\in W^{1,p}(I)$ has a unique continuous
representative, locally absolutely continuous, with
$u(x)-u(y)=\int_y^xu'(t)\,dt$. It extends continuously to finite
endpoints. For finite $I$,
\[
 \pfnorm u_\infty\le |I|^{-1/p}\pfnorm u_p
                         +|I|^{1-1/p}\pfnorm{u'}_p.
\]
For $p>1$, $|u(x)-u(y)|\le\pfnorm{u'}_p|x-y|^{1-1/p}$.
$W^{1,\infty}$ is exactly the essentially bounded functions with
a Lipschitz representative, whose least Lipschitz constant is
$\pfnorm{u'}_\infty$. The global $L^\infty$ condition matters on an
unbounded interval.
\end{theorem}
\begin{proof}
Take a local primitive $v(x)=\int_{x_0}^xu'$. Then $(u-v)'=0$
as a distribution. A distribution on an interval with zero derivative
is constant: if $\phi$ has integral zero, its primitive is a compactly
supported test, so the distribution annihilates $\phi$.
Subtracting a fixed integral-one test proves the constant assertion.
Thus $u=v+c$ locally, and uniqueness patches the representatives.
H\"older gives the increment estimates and endpoint limits.
Choose a point where $|u|\le |I|^{-1/p}\pfnorm u_p$ (or an
arbitrarily close bound), and use the increment estimate to obtain
the supremum bound. A Lipschitz function is absolutely continuous
on compact intervals, has an essentially bounded derivative, and
its fundamental theorem of calculus proves the converse. Uniqueness
of continuous representatives follows because an open interval of
nonzero differences cannot have measure zero.
\end{proof}

\begin{theorem}[Extension under the actual boundary hypotheses]
\label{fa2:p4:sobolev-extension}
On a half-space, or on a $C^1$ domain with compact boundary, there
is a bounded linear extension
$E:W^{1,p}(\Omega)\to W^{1,p}(\mathbb R^d)$ simultaneously for
$1\le p\le\infty$. It may be chosen to map bounded domains to
functions supported in one fixed bounded neighborhood.
Consequently smooth functions up to the boundary are dense for
finite $p$. The assertion here does not replace $C^1$ by Lipschitz
without proving the additional chart theorem.
\end{theorem}
\begin{proof}
For $x_d>0$ reflect evenly: $Eu(x',x_d)=u(x',|x_d|)$.
The line characterization proves the weak derivative rule, with
a sign change for the normal derivative, and the $p$th-power
integrals at most double. Near a $C^1$ boundary write the domain
as $x_d>g(x')$ and flatten it by
$\Phi(y',y_d)=(y',y_d+g(y'))$. On a sufficiently small relatively
compact chart both derivative matrices and Jacobians are bounded,
with inverse Jacobian bounded. Change of variables and the chain
rule first for smooth functions, then by interior smooth density,
show bounded pullback and pushforward in $W^{1,p}$ for finite $p$.
For $p=\infty$, smooth locally inside the chart instead: the
functions and derivatives have uniform supremum bounds, so
their pullbacks have those bounds multiplied by the chart
constants and converge locally in $L^1$. Passing against compact
tests identifies the limit weak derivatives; the supremum bound
passes by the scalar integral bound or weak-star compactness.
Thus no $W^{1,\infty}$ norm-density assertion is used.
Multiply by a partition of unity on the finitely many boundary
charts and an interior patch, reflect the flattened boundary pieces,
and extend the interior piece by zero. Supports of the partition
pieces are contained strictly inside their charts, so no artificial
jump occurs at a chart edge. Finite summation gives one linear
extension and its norm bound. A final cutoff gives the bounded-support
claim. Convolve the extension and restrict to obtain global density.
For exterior domains the interior piece can be unbounded, but zero
extension is still legitimate since it vanishes in a neighborhood of
the boundary. This is the source generality of \cite{B01}, 9.7.
\end{proof}

\begin{proposition}[Translation characterization]
\label{fa2:p4:sobolev-translations}
On $\mathbb R^d$, $u\in W^{1,p}$ implies
$\pfnorm{u(\cdot+h)-u}_p\le |h|\pfnorm{\nabla u}_p$,
using a fixed Euclidean gradient norm. Conversely, for
$1<p<\infty$, $u\in L^p$ with these differences bounded by
$C|h|$ belongs to $W^{1,p}$, with each derivative norm at most $C$.
\end{proposition}
\begin{proof}
For smooth $u$ integrate $h\cdot\nabla u(x+th)$ from $0$ to $1$
and use Minkowski; finite-$p$ density proves the result there.
For $p=\infty$ convolve globally, preserving the derivative bound,
and pass a locally $L^1$-convergent a.e. subsequence for the fixed
$h$ and its translate. The pointwise bound passes to the limit
and gives the same essential-supremum estimate. For the converse,
difference quotients in each coordinate are bounded in reflexive
$L^p$. A weakly convergent subsequence as $h\to0$ gives $g_j$.
Moving the difference quotient to a compact test and passing to the
limit yields $\int g_j\phi=-\int u\partial_j\phi$.
Weak lower semicontinuity gives the bound. At $p=1$ bounded
differences yield bounded variation in general, not necessarily an
$L^1$ weak derivative; the characteristic function of an interval
is the elementary obstruction.
\end{proof}

\section{Embedding, compactness, and trace}

\begin{lemma}[The $p=1$ Sobolev inequality]
\label{fa2:p4:sobolev-one}
If $d>1$ and $u\in C_c^1(\mathbb R^d)$, then
$\pfnorm u_{d/(d-1)}\le C_d\sum_j\pfnorm{\partial_ju}_1$.
\end{lemma}
\begin{proof}
For each $j$ the fundamental theorem of calculus gives
$|u(x)|\le F_j(\widehat x_j)$, where
$F_j(\widehat x_j)=\int_{\mathbb R}|\partial_ju(x)|\,dx_j$.
Multiply these $d$ inequalities to obtain
$|u|^{d/(d-1)}\le\prod_jF_j(\widehat x_j)^{1/(d-1)}$.
The repeated H\"older inequality
\[
 \int_{\mathbb R^d}\prod_j a_j(\widehat x_j)^{1/(d-1)}
 \le\prod_j\left(\int_{\mathbb R^{d-1}}a_j\right)^{1/(d-1)}
\]
follows by induction on $d$: integrate the last coordinate using
H\"older on its $d-1$ factors, then apply the induction hypothesis
to the resulting $d-1$ variables. For $d=2$ it is Fubini.
Taking the power $(d-1)/d$ and the arithmetic--geometric mean
inequality proves the claim.
\end{proof}

\begin{theorem}[Sobolev and Morrey embeddings]
\label{fa2:p4:sobolev-embedding}
For $1\le p<d$, on $\mathbb R^d$,
$\pfnorm u_{p^*}\le C\pfnorm{\nabla u}_p$ for
$u\in C_c^\infty$, $p^*=dp/(d-p)$, and hence for
$u\in W^{1,p}(\mathbb R^d)$. On bounded $C^1$ domains the same
embedding holds with the full $W^{1,p}$ norm. For $p=d$ these
domains admit all finite target exponents. For $p>d$ there is a
continuous representative satisfying
$|u(x)-u(y)|\le C\pfnorm{\nabla u}_p|x-y|^{1-d/p}$ on
$\mathbb R^d$, and the analogous full-norm estimate on an extension
domain. Iteration gives the higher-order embeddings: if $mp<d$,
$W^{m,p}\hookrightarrow L^{dp/(d-mp)}$; if $mp=d$, all finite
exponents are allowed on bounded extension domains; if $mp>d$,
the derivatives have the H\"older regularity obtained when their
remaining Sobolev order exceeds $d/p$. In particular nonintegral
$m-d/p=k+\alpha$ gives $C^{k,\alpha}$, $0<\alpha<1$;
at a positive integer gap $k=m-d/p$, the conclusion is
$C^{k-1,\alpha}$ for every $0<\alpha<1$, without claiming a
Lipschitz endpoint for that last derivative.
\end{theorem}
\begin{proof}
For $p>1$ apply the previous lemma to $|u|^\gamma$, by a smooth
approximation at zero, with $\gamma=p(d-1)/(d-p)$.
H\"older on $\gamma|u|^{\gamma-1}|\nabla u|$ and
$(\gamma-1)p'=p^*$ gives the assertion after cancelling
$\pfnorm u_{p^*}^{\gamma-1}$. A zero norm is handled separately.
Cutoff and smooth density extend it to $W^{1,p}(\mathbb R^d)$;
the cutoff error tends to zero. Extension proves the bounded-domain
claim. At $p=d$ choose $p_0<d$ so that $p_0^*$ exceeds the desired
finite exponent, and use finite-measure H\"older inclusions.

For Morrey, averaging the identity
$u(x)-u(z)=-\int_0^1\nabla u(x+t(z-x))\cdot(z-x)\,dt$
over a ball of radius $r$ gives
\[
 |u(x)-u_{B(x,r)}|
 \le C\int_{B(x,r)}|\nabla u(z)|\,|x-z|^{1-d}\,dz
 \le C r^{1-d/p}\pfnorm{\nabla u}_{L^p(B(x,r))}.
\]
The last integral is finite exactly when $p>d$.
For $r=2|x-y|$, compare both averages to the average over an
enclosing ball to obtain the difference estimate. Initially smooth
functions satisfy it; for finite $p>d$ the estimate together with
$W^{1,p}$ norm convergence makes smooth approximants uniformly
Cauchy on compact sets. For $p=\infty$, locally mollified
functions have uniform Lipschitz and supremum bounds, so Ascoli
gives a locally uniform subsequential limit. Local $L^1$
convergence identifies it with $u$ a.e., and the Lipschitz bound
passes to this representative. Higher-order statements follow by
applying the first-order embedding to each derivative in succession.
At each step the reciprocal target exponent drops by $1/d$ until
the Morrey regime is reached; this gives precisely the stated
critical and noncritical exponents. No assertion of boundedness at
the critical $W^{1,d}$ endpoint for $d>1$ is made: near zero,
$(\log(1/|x|))^\alpha$ with $0<\alpha<1-1/d$ is a counterexample.
\end{proof}

\begin{theorem}[Kolmogorov--Riesz and Rellich]
\label{fa2:p4:kolmogorov-riesz}\label{fa2:p4:rellich}
For $1\le p<\infty$, a bounded family $\mathcal A\subset L^p(\mathbb R^d)$
is relatively norm compact exactly when it is uniformly tight in
$L^p$ and
$\sup_{f\in\mathcal A}\pfnorm{f(\cdot+h)-f}_p\to0$ as $h\to0$.
The translation condition alone gives compactness after restriction
to any bounded measurable region. On bounded $C^1$ domains,
$W^{1,p}\hookrightarrow L^q$ is compact for $q<p^*$ if $p<d$,
for every finite $q$ if $p=d$, and into $C(\overline\Omega)$ if
$p>d$. For $W^{1,p}_0$ the subcritical $L^q$ conclusions hold on
arbitrary bounded open sets.
\end{theorem}
\begin{proof}
Necessity follows from finite norm nets: each individual function
has small tails and continuous translations, and the norm of a
translation is one. For sufficiency truncate to a large cube, losing
uniformly little norm. Partition it into small cubes $Q$ of side
$\delta$ and replace $f$ on each by its average $f_Q$.
Jensen and Fubini give
\[
 \sum_Q\int_Q|f-f_Q|^p
 \le\delta^{-d}\int_{|h|\le\sqrt d\delta}
             \int|f(x+h)-f(x)|^p\,dx\,dh.
\]
The right side is bounded by a dimension constant times the uniform
translation modulus. The step averages form a bounded set in a
fixed finite-dimensional space, hence are totally bounded.
Completeness gives relative compactness. If no tails are assumed,
perform this argument only on a cube containing the target region.

For Rellich extend a bounded Sobolev family with common bounded
support. Translation differences in $L^p$ are $O(|h|)$.
Interpolation of this bound with the bounded critical Sobolev norm
gives a modulus in every $p\le q<p^*$; finite-measure H\"older
handles $q<p$. At $p=d$ first bound an exponent larger than $q$.
For $p>d$, Morrey gives uniform boundedness and equicontinuity,
and the finite-net proof of Arzel\`a--Ascoli gives compactness in
$C(\overline\Omega)$. For zero-boundary spaces extend compact
tests by zero and pass to their $W^{1,p}$ closure; no boundary
regularity is needed for that extension. Critical compactness fails:
rescaled bumps concentrating at an interior point have constant
critical norm and no convergent critical subsequence.
\end{proof}

\begin{theorem}[Trace, its kernel, and Poincar\'e]
\label{fa2:p4:sobolev-trace}\label{fa2:p4:poincare}
On a bounded $C^1$ domain and $1\le p<\infty$, restriction of smooth
functions extends to a bounded trace
$\operatorname{Tr}:W^{1,p}(\Omega)\to L^p(\partial\Omega)$.
Its kernel is $W^{1,p}_0(\Omega)$, and a function is in that kernel
exactly when its zero extension belongs to $W^{1,p}(\mathbb R^d)$.
On arbitrary bounded open $\Omega\subset\mathbb R^d$ with $d\ge1$,
$\pfnorm u_p\le C\pfnorm{\nabla u}_p$ for $u\in W^{1,p}_0$.
The same inequality holds trivially on $\Omega=\varnothing$ in every dimension.
On an interval these trace values are its two endpoint values.
\end{theorem}
\begin{proof}
On a flat half-space, for compactly supported smooth $u$,
integrate $-\partial_d|u|^p$ on each normal half-line. H\"older gives
$\pfnorm{u(\cdot,0)}_p^p
\le p\pfnorm u_p^{p-1}\pfnorm{\partial_du}_p$; for $p=1$ use
the absolute continuity inequality directly. Flattening charts and
a finite partition yield the trace bound, and smooth-up-to-boundary
density defines the unique trace.
Smooth Green's formula, passed through this approximation, shows
that the derivative of the zero extension equals the zero extension
of the derivative exactly when the boundary trace is zero.
More explicitly the boundary term against a compact test is
$\int_{\partial\Omega}\operatorname{Tr}(u)\phi\nu_j$; if all
these terms vanish, a local chart with nonzero normal component
and density of smooth boundary tests make the trace zero.
If the trace is zero, in each flattened chart translate the zero
extension slightly into the domain, then mollify at a smaller radius.
These are compactly supported interior smooth functions converging
in $W^{1,p}$, by translation continuity. Patch them with the
interior piece. Thus the trace kernel is precisely $W_0$.
For Poincar\'e with $d\ge1$, extend a compact test by zero into a box containing
$\Omega$. On every line parallel to one box axis the fundamental
theorem and H\"older bound its $L^p$ norm by the box length times
its derivative norm. Integrate transverse variables and pass to
the closure. This proof works on arbitrary bounded domains.
It also works when the domain has bounded projection on one axis.
\end{proof}

\begin{theorem}[Negative Sobolev representation]
\label{fa2:p4:negative-sobolev}
For $1\le p<\infty$ and $q=p'$, every functional
$F\in(W^{1,p}_0(\Omega))'$ has the representation
\[
 F(u)=\int f_0u+\sum_{j=1}^d\int f_j\partial_ju,
 \qquad f_j\in L^q,
 \qquad \pfnorm F=\max_{0\le j\le d}\pfnorm{f_j}_q,
\]
for the sum-component Sobolev norm used here and on the present
Lebesgue measure space. With the distribution convention this is
$F=f_0-\sum\partial_jf_j$. On bounded domains with $d\ge1$,
and trivially on $\Omega=\varnothing$ in every dimension,
$f_0$ can be omitted with a constant depending on Poincar\'e;
norm equality then refers to the derivative norm,
not automatically the original full norm.
\end{theorem}
\begin{proof}
Map $u$ isometrically into $(L^p)^{d+1}$ with its sum norm by
$u\mapsto(u,\partial_1u,\ldots,\partial_du)$.
Hahn--Banach extends $F$ to this product with the same norm.
The scalar dual theorem on Lebesgue space identifies its coordinates
with $f_j\in L^q$, and elementary finite-product duality identifies
the norm with the maximum of their norms. Restrict back to the
closed Sobolev image. On a bounded domain with $d\ge1$, the gradient-only map is
an isomorphism onto its closed image by Poincar\'e; applying the
same extension argument there omits $f_0$ and gives the asserted
norm interpretation. On the empty domain all spaces and functionals
are zero, so the derivative-only assertion is immediate in every dimension. For complex anti-dual data conjugate the
test variable throughout, exactly as in the form convention of
Part II; the displayed representation itself is linear in $u$.
\end{proof}

\section{Fourier Hilbert scales and regularity}

\begin{theorem}[The Fourier Sobolev scale]
\label{fa2:p4:sobolev-hilbert}
For real $s$ define $H^s(\mathbb R^d)$ by
$\pfnorm u_{H^s}^2=\int(1+4\pi^2|\xi|^2)^s|\widehat u(\xi)|^2d\xi$.
It is a Hilbert space of tempered distributions. Its linear dual is
$H^{-s}$ under the bilinear distribution pairing
\[
 (u,v)\longmapsto\int\widehat u(\xi)\widehat v(-\xi)d\xi,
\]
which extends $\int uv$ on smooth functions, and
$(I-\Delta)^{t/2}:H^s\to H^{s-t}$ is an isometric isomorphism.
For nonnegative integer $m$, $H^m=W^{m,2}$ with equivalent norms.
Compactly supported distributions belong to $H^{-N}$ for some $N$.
Multiplication by a compactly supported smooth function is bounded
on every $H^s$; derivatives lower the index by their order.
If $s>k+d/2$, $H^s\subset C^k$ with bounded continuous derivatives.
\end{theorem}
\begin{proof}
Multiply the Fourier transform by the positive weight
$(1+4\pi^2|\xi|^2)^{s/2}$ to identify the space with $L^2$.
Polynomial weights define tempered distributions, so the inverse
transform exists. Cauchy--Schwarz proves the dual pairing and its
norm; Riesz representation in this weighted $L^2$ proves surjectivity.
The multiplier assertion is direct. Plancherel and comparison of
$\sum_{|\alpha|\le m}|\xi^\alpha|^2$ with $(1+|\xi|^2)^m$
prove the integer assertion. A compact distribution has polynomially
bounded real Fourier transform by Paley--Wiener, which belongs to
weighted $L^2$ for sufficiently negative order.
For multiplication, use
$\widehat{au}=\widehat a*\widehat u$ and
\[
 (1+|\xi|^2)^{s/2}
 \le C_s(1+|\eta|^2)^{|s|/2}(1+|\xi-\eta|^2)^{s/2}.
\]
The weighted $L^1*L^2$ Young inequality bounds it because
$\widehat a$ is Schwartz. Derivatives are the corresponding
polynomial multipliers. Finally Cauchy--Schwarz makes
$\xi^\alpha\widehat u\in L^1$ for $|\alpha|\le k$ when
$s>k+d/2$. Fourier inversion and dominated convergence give the
bounded continuous derivatives and their norms.
\end{proof}

\begin{proposition}[Periodic and smooth-curve Hilbert scales]
\label{fa2:p4:periodic-sobolev}
On the circle,
$H^s=\{u:\sum_{n\in\mathbb Z}(1+4\pi^2n^2)^s|\widehat u(n)|^2<\infty\}$.
Its dual is $H^{-s}$, derivatives lower order, and
$H^{s+\varepsilon}\hookrightarrow H^s$ is compact for
$\varepsilon>0$. On a compact smooth curve, a finite smooth atlas
and partition of unity define equivalent $H^s$ norms independently
of atlas. A $C^\infty$ reparametrization is an isomorphism of these
spaces; the scalar-density Jacobian is included in the dual pairing.
\end{proposition}
\begin{proof}
Weighted sequence duality proves the first assertions; truncation
to $|n|\le N$ has error at most
$(1+4\pi^2N^2)^{-\varepsilon/2}$ between the two unit balls,
so the embedding is compact. For chart invariance at integer
nonnegative order use the chain rule, bounded derivatives of the
transition maps, and upper and lower Jacobian bounds on the
compact overlaps. For $0<\sigma<1$, Plancherel and Fubini give
\[
 \iint\frac{|v(x+h)-v(x)|^2}{|h|^{1+2\sigma}}\,dx\,dh
 =c_\sigma\int|\xi|^{2\sigma}|\widehat v(\xi)|^2\,d\xi,
\]
where $0<c_\sigma<\infty$, by scaling the inner integral.
Bi-Lipschitz change of variables preserves the corresponding
difference-quotient norm up to constants. Combine this with the
integer derivative description for $s=m+\sigma$.
Negative indices follow by duality; changing variables in the
integral pairing introduces its smooth nonvanishing Jacobian,
whose multiplication is bounded by the same estimates. A finite
atlas makes all constants finite. This establishes the general
function-space bridge used in \cite{B06}, Chapter 8, without
claiming any boundary layer-potential mapping theorem.
\end{proof}

\begin{theorem}[Constant-coefficient elliptic parametrix and local regularity]
\label{fa2:p4:elliptic-parametrix}
Let $L=P(D)$ have order $m$ and elliptic principal symbol.
There is a tempered distribution $Q$, smooth off zero, with
$LQ=\delta+r$, $r\in\pfS$. A parametrix supported in any prescribed
ball about zero can be chosen with a smooth remainder.
If $Lu$ is smooth on an open set, then $u$ is smooth there.
\end{theorem}
\begin{proof}
Ellipticity gives $|P(\xi)|\ge c|\xi|^m$ for large real $\xi$.
Choose $\chi\in C_c^\infty$ equal to one near the finite exceptional
region and set $\widehat Q=(1-\chi)/P$.
Repeated differentiation of the quotient gives
$|\partial^\alpha\widehat Q(\xi)|\le C_\alpha(1+|\xi|)^{-m-|\alpha|}$.
Integration by parts in the inverse transform, or the Hilbert-scale
Sobolev lemma after multiplying by powers of $x$, shows $Q$ is
smooth away from zero. Its symbol gives $LQ=\delta-\check\chi$.
Multiply $Q$ by a cutoff equal to one near zero and supported in
the prescribed ball; commutator terms are supported away from
zero and are smooth, proving the compact parametrix claim.
To localize, choose nested balls $B_1\Subset B_2\Subset\Omega$,
a cutoff $a=1$ on $B_2$, and a compact parametrix with support
smaller than the distance from $B_1$ to $\pfsupp\nabla a$.
Then on $B_1$,
$u=Q*L(au)-r*(au)$.
The part $Q*(aLu)$ is smooth because convolution of a compact
distribution with a smooth function is smooth. The commutator
$[L,a]u$ has support away from $B_1$ and its convolution contributes
zero there by the chosen support. The last term is smooth by
compact support of $au$. This proves local smoothness.
\end{proof}

\chapter{Interpolation and Standard Singular Integrals}
\label{fa2:p4:interpolation-chapter}

\begin{lemma}[Three lines]\label{fa2:p4:three-lines}
If $F$ is bounded and continuous on $0\le\operatorname{Re}z\le1$
and holomorphic inside, with $|F(j+it)|\le M_j$, then
$|F(\theta)|\le M_0^{1-\theta}M_1^\theta$, $0<\theta<1$.
\end{lemma}
\begin{proof}
For positive $M_j$ apply the maximum principle on large rectangles to
\[
 F(z)M_0^{z-1}M_1^{-z}\exp(\varepsilon(z-\theta)^2).
\]
On the vertical sides its modulus is at most $e^\varepsilon$.
On the horizontal sides boundedness of $F$ and the factor
$e^{-\varepsilon(\operatorname{Im}z)^2}$ give the same bound
when the rectangle is large. At $\theta$ let the height tend to
infinity and then $\varepsilon\downarrow0$. Replace a zero
$M_j$ by $M_j+\delta$ and let $\delta\downarrow0$.
Only the ordinary one-variable maximum principle is used here.
\end{proof}

\begin{theorem}[Riesz--Thorin, with common-domain compatibility]
\label{fa2:p4:riesz-thorin}
On $\sigma$-finite measure spaces let one complex-linear operator on
the complex space $L^{p_0}+L^{p_1}$ have compatible bounds
$L^{p_j}\to L^{q_j}$ with norms $M_j$, $1\le p_j,q_j\le\infty$.
For $0<\theta<1$ set
$p^{-1}=(1-\theta)p_0^{-1}+\theta p_1^{-1}$ and similarly $q$.
Then $\pvnorm{Tf}_q\le M_0^{1-\theta}M_1^\theta\pvnorm f_p$.
In particular the finite-measure $L^1/L^\infty$ self-map form has
$M_p\le M_1^{1/p}M_\infty^{1/p'}$; that special self-map
bound also holds for real operators, by the norm-preserving
endpoint complexification proved in
Proposition~\ref{fa2:p2:endpoint-interpolation-bridge}.
\end{theorem}
\begin{proof}
First assume $p,q'<\infty$. Take finite-support simple $f,h$
normalized in $L^p,L^{q'}$.
Where $f\ne0$ define
\[
 f_z=\frac f{|f|}|f|^{p((1-z)/p_0+z/p_1)},
 \qquad
 h_z=\frac h{|h|}|h|^{q'((1-z)/q'_0+z/q'_1)},
\]
and zero elsewhere, using reciprocal infinity equal to zero.
The scalar $F(z)=\int (Tf_z)h_z$ is a finite linear combination
of analytic exponentials. It is bounded on the strip: their
moduli depend only on $\operatorname{Re}z$, and only finitely
many fixed pairings occur. On either boundary the appropriate
norms of $f_z,h_z$ are one, including the supremum-norm endpoints.
Thus $|F(j+it)|\le M_j$. Three lines at $\theta$ and scalar
$L^q$ norm detection by simple $L^{q'}$ tests give the desired
bound for simple $f$. Density extends it to intermediate finite
$p$ and compatibility preserves the original operator on its
common domain. If $p=\infty$, both domain endpoints are infinity;
then scalar norm interpolation
$\pvnorm g_q\le\pvnorm g_{q_0}^{1-\theta}\pvnorm g_{q_1}^\theta$
is just H\"older and gives the result directly for $g=Tf$.
If $q'=\infty$ and $p<\infty$, both target endpoints are one.
Take a bounded simple test $h$ of supremum norm one and put
$h_z=h$, constant in $z$. The same analytic $f_z$ and three-lines
argument then apply without the undefined product $\infty\cdot0$.
Supremizing these tests gives the $L^1$ output bound, and finite-$p$
density finishes this case. This proof uses no convolution inequality,
so its use in the proof of full Young above is not circular.
\end{proof}

\begin{theorem}[Marcinkiewicz interpolation]
\label{fa2:p4:marcinkiewicz}
A sublinear operator with weak $(p_0,p_0)$ and weak $(p_1,p_1)$
bounds, $1\le p_0<p_1\le\infty$, is strongly bounded on $L^p$
for $p_0<p<p_1$. The infinity endpoint here means a strong
supremum bound. The result applies to the modulus of a linear
operator defined first on a dense intersection.
\end{theorem}
\begin{proof}
For finite endpoints split $f=f_{>c\lambda}+f_{\le c\lambda}$.
Sublinearity and the two weak estimates give
\[
 \mu\{|Tf|>\lambda\}
 \le (2C_0/\lambda)^{p_0}
           \int_{|f|>c\lambda}|f|^{p_0}
 +(2C_1/\lambda)^{p_1}
           \int_{|f|\le c\lambda}|f|^{p_1}.
\]
Multiply by $p\lambda^{p-1}$ and integrate. Tonelli evaluates
the two inner integrals in $\lambda$ as respectively constants
times $|f|^{p-p_0}/(p-p_0)$ and
$|f|^{p-p_1}/(p_1-p)$. The result is a finite constant times
$\pvnorm f_p^p$. At the infinity endpoint choose
$c=(2C_\infty)^{-1}$; the low part has image bounded by
$\lambda/2$ and only the first term remains. Truncation and
completion extend the bound from the initial intersection.
The integrals diverge at an endpoint, explaining why no strong
endpoint assertion follows from this argument.
\end{proof}

\begin{theorem}[Hardy--Littlewood maximal bounds]
\label{fa2:p4:maximal}
Define
\[
 Mf(x)=\sup_{r>0}\frac{1}{|B(x,r)|}\int_{B(x,r)}|f|.
\]
Then $\mu\{Mf>\lambda\}\le C_d\pvnorm f_1/\lambda$.
It is bounded on $L^p$ for $1<p\le\infty$. If $d\ge1$, it is not bounded on $L^1$.
Every locally integrable function has almost every point as a
Lebesgue point.
\end{theorem}
\begin{proof}
For a compact subset of a level set choose finitely many balls
with average exceeding $\lambda$. Select disjoint balls greedily
in decreasing radius; their fivefold dilations cover the original
family. Thus the covered measure is at most
$5^d\sum|B_j|\le5^d\pvnorm f_1/\lambda$.
Compact exhaustion and Lebesgue regularity give the full weak bound.
The supremum bound is one; Marcinkiewicz gives the interior exponents.
If $d\ge1$, for $f=1_{B(0,1)}$ we have $Mf(x)\ge c|x|^{-d}$
at infinity, and the radial integral diverges. In dimension zero the
space is a singleton of volume one, so $Mf=|f|$, the strong $L^1$
bound holds with norm one, and differentiation is immediate.
For differentiation approximate a local $L^1$ function by continuous
functions on a slightly larger bounded region. The limsup mean
oscillation of its difference from an approximant is bounded by
twice its maximal function plus its pointwise absolute value.
The weak bound and Markov's inequality make the set where this
exceeds a fixed threshold arbitrarily small. Continuous functions
have zero limiting oscillation. A countable exhaustion proves the
local differentiation conclusion.
\end{proof}

\begin{lemma}[Calder\'on--Zygmund decomposition]
\label{fa2:p4:cz-decomposition}
For $d\ge1$, $f\in L^1(\mathbb R^d)$ and $\lambda>0$, there are disjoint
dyadic cubes $Q_j$, $f=g+\sum b_j$, with
\[
\begin{aligned}
 &\sum|Q_j|\le\pvnorm f_1/\lambda,\qquad
 \pvnorm g_\infty\le2^d\lambda,\qquad
 \pvnorm g_1\le\pvnorm f_1,\\
 &\int b_j=0,\qquad
 \sum\pvnorm{b_j}_1\le2\pvnorm f_1,\qquad
 \pfsupp b_j\subset Q_j.
\end{aligned}
\]
If $f\in L^2$, also $g,b_j\in L^2$ and the sum converges in
$L^2$ as well as in $L^1$.
\end{lemma}
\begin{proof}
Take maximal dyadic cubes whose average of $|f|$ exceeds
$\lambda$. Large ancestors have small average by integrability,
so maximal cubes exist; dyadic nesting makes them disjoint.
Their parents give averages at most $2^d\lambda$.
Outside their union $|f|\le\lambda$ a.e. by differentiation.
Set $g=f$ there and $g=f_{Q_j}$ on each cube, and
$b_j=(f-f_{Q_j})1_{Q_j}$. All bounds follow by integration and
the disjoint supports. If $f\in L^2$, Jensen bounds the squared
averages by the original squared integrals; disjointness then
gives convergence of the bad part in $L^2$.
\end{proof}

\begin{theorem}[Standard convolution Calder\'on--Zygmund theorem]
\label{fa2:p4:calderon-zygmund}
Let $T$ be an $L^2$-bounded translation-invariant operator whose
distribution kernel agrees off zero with a $C^1$ function $k$,
where $|k(x)|\le C|x|^{-d}$ and
$|\nabla k(x)|\le C|x|^{-d-1}$.
Then $T$ is weak $(1,1)$ and bounded on every $L^p$,
$1<p<\infty$, with compatible extensions. This is a theorem
about standard convolution kernels; a general variable-kernel
theory is not being claimed from it.
\end{theorem}
\begin{proof}
Apply the decomposition to $f\in L^1\cap L^2$ at height
$\lambda$. The good part satisfies
$\pvnorm g_2^2\le2^d\lambda\pvnorm f_1$, so Chebyshev and the
$L^2$ bound control $\{|Tg|>\lambda/2\}$ by
$C\pvnorm f_1/\lambda$.
Enlarge each bad cube by a fixed dimension-dependent factor.
Their total measure has the same bound. Off the enlarged cube,
zero integral and the mean-value theorem give
\[
 Tb_j(x)=\int_{Q_j}[k(x-y)-k(x-c_j)]b_j(y)dy,
 \qquad
 \int_{(Q_j^*)^c}|Tb_j(x)|dx\le C\pvnorm{b_j}_1.
\]
Indeed $|x-c_j|>2|y-c_j|$ there, and the integral of
$C|y-c_j||x-c_j|^{-d-1}$ over this region is bounded uniformly.
The representation off the support follows first for smooth
approximants and then by the $L^2$ kernel distribution identity;
its right side is absolutely integrable there.
Sum the estimates and use Chebyshev to control the bad part
off all enlarged cubes. This proves weak $(1,1)$.
An $L^1$ approximating sequence has images Cauchy in measure by
that estimate; choose a subsequence with summable exceptional
sets to obtain the unique measurable extension and its weak bound.
Marcinkiewicz between one and two proves strong bounds for
$1<p<2$. The adjoint has kernel $\overline{k(-x)}$ and the same
bounds. Duality gives $p>2$, and compatibility on common dense
classes gives a single operator. This proves every estimate
claimed, including the endpoint weak estimate, rather than only
a formal multiplier assertion.
\end{proof}

\begin{proposition}[The standard smooth-kernel multiplier equivalence]
\label{fa2:p4:cz-multiplier}
For a distribution $K$ smooth off zero with
$|\partial^\alpha k(x)|\le C_\alpha|x|^{-d-|\alpha|}$, the
following equivalent conditions specify a Calder\'on--Zygmund
distribution: uniformly bounded pairings with all dilations of
normalized compactly supported $C^N$ bumps, for one sufficiently
large fixed $N$; or $\widehat K=m$ bounded, smooth off zero,
with $|\partial^\alpha m(\xi)|\le C_\alpha|\xi|^{-|\alpha|}$.
Equivalently, in the presence of all the displayed kernel bounds,
it suffices that $\widehat K$ be a bounded function.
\end{proposition}
\begin{proof}
An annular kernel of radius $r$ with these derivative bounds has
Fourier derivatives bounded by
$C_{\alpha,L}r^{|\alpha|}(1+r|\xi|)^{-L}$:
rescale to a fixed annulus and integrate by parts $L$ times.
Fix $\xi_0\ne0$ and split $K$ by a cutoff at
$r=|\xi_0|^{-1}$. The inner pairing against the exponential
and its $\xi$ derivatives is bounded by $Cr^{|\alpha|}$ by
the normalized-bump assumption. The outer dyadic annuli sum
absolutely after integration by parts with $L>|\alpha|$,
giving $C|\xi_0|^{-|\alpha|}$ uniformly in a neighborhood of
$\xi_0$. This constructs the smooth multiplier away from zero.

Conversely a dyadic frequency partition of a bounded $m$ with
these derivative bounds gives inverse kernels
$2^{jd}a_j(2^jx)$ with uniformly bounded Schwartz seminorms.
Their derivatives sum uniformly off zero, splitting at
$2^j|x|=1$, to give
$|\partial^\alpha k(x)|\le C_\alpha|x|^{-d-|\alpha|}$.
They converge in tempered distributions by dominated convergence
against Schwartz transforms. The possible discrepancy at frequency
zero in the first construction is a delta-derivative combination;
its inverse transform is a polynomial. Both physical kernels
decay off zero, so this polynomial must be zero.
Finally if $m\in L^\infty$, bump pairings are bounded by
$\pvnorm m_\infty\pvnorm{\widehat\phi}_1$.
For bumps supported in the unit ball, a uniform $C^N$ bound,
$N>d$, bounds this last norm by integration by parts; dilation
leaves that norm unchanged. These estimates also show that the
initial distribution is tempered: a cutoff near zero is compactly
supported and its remaining decaying kernel pairs with Schwartz
functions absolutely. This proves the equivalence in the source
standard convolution setting of \cite{B09}, Proposition 3.1.
\end{proof}

\begin{theorem}[Principal value, boundary values, and Hilbert transform]
\label{fa2:p4:hilbert-transform}\label{fa2:p4:principal-value}
On $\mathbb R$, $\operatorname{pv}(1/x)$ is tempered,
$(\log|x|)'=\operatorname{pv}(1/x)$, and
\[
 (x\pm i0)^{-1}=\operatorname{pv}(1/x)\mp i\pi\delta_0,
 \qquad \widehat{\operatorname{pv}(1/x)}=-i\pi\operatorname{sgn}\xi.
\]
The Hilbert transform has multiplier $-i\operatorname{sgn}\xi$,
is an isometry of $L^2$ with $H^2=-I$, and is bounded on $L^p$
for $1<p<\infty$ and weak $(1,1)$. It has neither a strong
$L^1$ nor a strong $L^\infty$ bound. Its symmetric truncated
integrals converge in $L^2$ to the multiplier definition.
\end{theorem}
\begin{proof}
In the symmetric integral near zero subtract $\phi(0)$;
the remaining quotient is locally bounded. Schwartz decay
handles infinity, proving existence and continuity as a tempered
distribution. Integration by parts off $(-\varepsilon,\varepsilon)$
leaves a boundary term
$\log\varepsilon[\phi(\varepsilon)-\phi(-\varepsilon)]\to0$,
proving the derivative identity.
Write $(x\pm i\varepsilon)^{-1}
=x/(x^2+\varepsilon^2)\mp i\varepsilon/(x^2+\varepsilon^2)$.
The first term tends to the principal value by subtracting the
test value at zero; the second is a mass-$\pi$ approximate
identity. This proves the boundary formulas.
For the transform insert damping $e^{-\varepsilon|x|}$ in the
odd integral. The identity
$\int_0^\infty e^{-\varepsilon x}\sin(bx)/x\,dx
=\arctan(b/\varepsilon)$ follows by differentiating in $b$
and integrating back from zero. Let $\varepsilon\downarrow0$
to obtain the multiplier with the stated sign and $2\pi$
normalization. Plancherel gives the isometry and its square.
The kernel $1/(\pi x)$ satisfies the standard CZ estimates,
so that theorem gives all the asserted $L^p$ and weak bounds.
For the truncated integral its multiplier is
$-(2i/\pi)\int_\varepsilon^\infty\sin(2\pi\xi t)/t\,dt$.
Dirichlet integration by parts bounds it uniformly, and it tends
pointwise a.e. to $-i\operatorname{sgn}\xi$. Dominated convergence
in the Plancherel identity proves $L^2$ convergence.
Finally
$H1_{[0,1]}(x)=\pi^{-1}\log|x/(x-1)|$ off the endpoints.
It has a nonintegrable $1/x$ tail and logarithmic unboundedness
near the endpoints, disproving both strong endpoint bounds.
\end{proof}

\begin{proposition}[Atomic $H^1$, its completeness, and the endpoint bound]
\label{fa2:p4:hardy-atomic}
An atom is supported in a ball $B$, has integral zero, and has
supremum norm at most $|B|^{-1}$. The sums
$f=\sum\lambda_ja_j$, $\sum|\lambda_j|<\infty$, form a Banach
space $H^1_{\rm at}$ with norm the infimum of coefficient sums.
It embeds continuously in $L^1$, all its members have integral
zero, and every standard convolution CZ operator is bounded
$H^1_{\rm at}\to L^1$.
\end{proposition}
\begin{proof}
Map $\ell^1$ of the set of atoms into $L^1$ by summation;
each element uses at most countably many atoms. Its kernel is
closed, and the quotient is Banach with precisely the infimum
norm. This also proves definiteness and the $L^1$ embedding.
Integral zero passes through the absolute $L^1$ sum.
For an atom, on a fixed enlargement of its ball use the $L^2$
bound and $\pvnorm a_2\le|B|^{-1/2}$. Outside that enlargement
use its zero integral and the kernel difference estimate from
the CZ proof. Both pieces have uniformly bounded $L^1$ norm.
Absolute coefficient summation gives a bounded map on the
quotient. It agrees with the previously constructed weak $(1,1)$
extension because partial atomic sums converge in $L^1$, whose
images converge in measure by the weak bound and in $L^1$ by
this estimate. Thus it is an extension of the same operator,
not a second unrelated endpoint operator.
\end{proof}

\begin{definition}[BMO]
\label{fa2:p4:bmo}
$\operatorname{BMO}(\mathbb R^d)$ consists of locally integrable
functions with seminorm
$\sup_B|B|^{-1}\int_B|b-b_B|<\infty$, modulo constants.
It is complete in this quotient seminorm. Bounded functions lie
in BMO. If $d\ge1$, the function $\log|x|$, assigned value zero
at the origin, is an unbounded BMO function.
\end{definition}
\begin{proof}
For completeness normalize means on one fixed ball to zero.
A Cauchy sequence in BMO is Cauchy in $L^1$ on that ball;
on any larger fixed ball its mean difference is controlled by
a finite chain of nested doubled balls, since
$|b_B-b_{2B}|\le2^d\pvnorm b_{\rm BMO}$.
Thus it is locally $L^1$-Cauchy. The local limit has the
limiting BMO bound by integration, and the same argument applied
to the difference proves norm convergence modulo constants.
In dimension zero the BMO quotient is the zero space.
For the logarithm in positive dimension, balls away from zero have bounded oscillation
by the gradient estimate after scaling; balls meeting a fixed
multiple of the origin reduce by scaling to a bounded ball,
where $\log|x|$ is integrable. Subtracting its scaling constant
does not change oscillation. This proves the example.
No $H^1$--BMO duality theorem is being inferred from these
definitions alone.
\end{proof}

\begin{theorem}[$p$-atoms and compactly supported mean-zero functions]
\label{fa2:p4:hardy-p-atoms}
For every $1<p<\infty$, a function supported in a ball $B$,
with integral zero and $\pfnorm a_p\le|B|^{1/p-1}$, belongs
to $H^1_{\rm at}$ with norm at most a constant depending only
on $p,d$. Thus $L^\infty$ atoms can be replaced by these
$p$-atoms with an equivalent norm. Every compactly supported
$L^p$ function of integral zero belongs to $H^1$.
\end{theorem}
\begin{proof}
Enclose $B$ in a cube of comparable measure and normalize its
side and the function. On this fixed cube apply the dyadic
decomposition at successive levels $\lambda_k=2^k\lambda_0$,
where $\lambda_0$ exceeds the original cube's average absolute
value. Let $g_k$ equal $a$ outside the maximal bad cubes and
its cube average inside. It is bounded by $2^d\lambda_k$,
has support in the original cube, and preserves its zero
integral. The bad cubes at level $k+1$ are contained in those
at level $k$. Hence on each level-$k$ bad cube $Q$, the
restriction of $g_{k+1}-g_k$ has integral zero, support $Q$,
and bound $C\lambda_k$. It is a multiple of an infinity-atom
on the enclosing ball with coefficient $C\lambda_k|Q|$.
The initial $g_0$ is also a multiple of one such atom.
The total coefficient bound is
\[
 C\lambda_0|B|+C\sum_{k\ge0}\lambda_k
                    |\{M_{\rm dyad}|a|>\lambda_k\}\cap Q_0|
 \le C\int_{Q_0}M_{\rm dyad}|a|
 \le C_{p,d}|Q_0|^{1-1/p}\pfnorm a_p.
\]
Here the distribution sum is bounded by the integral up to
the initial-level term, and the final inequality is the strong
$L^p$ maximal estimate. Also $g_k\to a$ in $L^1$:
the original $L^1$ integral on the shrinking bad set tends to
zero, and
$\lambda_k|\{M_{\rm dyad}|a|>\lambda_k\}|\to0$ because the
maximal function is integrable on this cube. The telescoping
series is therefore an actual atomic expansion of $a$.
Scaling back proves the uniform $p$-atom bound. Infinite
coefficient series then give the equivalent norm in both
directions, since infinity-atoms are $p$-atoms.
\end{proof}

\begin{theorem}[Smooth maximal bound and truncated Hilbert convergence]
\label{fa2:p4:hardy-maximal-forward}
For $\Phi\in C_c^1(\mathbb R^d)$,
$M_\Phi f=\sup_{t>0}|\Phi_t*f|$ satisfies
$\pfnorm{M_\Phi f}_1\le C_\Phi\pfnorm f_{H^1}$.
The bound is uniform over bumps supported in the unit ball
with uniformly bounded $C^1$ norm. On the real line symmetric
Hilbert truncations converge in $L^1$ for every $H^1$ function,
and their norms $H^1\to L^1$ are bounded independently of the
truncation parameter.
\end{theorem}
\begin{proof}
For a unit-ball atom the maximal function is bounded locally
in $L^1$ by its $L^2$ norm and the $L^2$ maximal bound
$M_\Phi a\le C M|a|$. Away from twice its support, subtract
$\Phi_t(x)$ using zero integral. The gradient bound gives
$|\Phi_t(x-y)-\Phi_t(x)|\le C|y|t^{-d-1}$,
and a nonzero term requires $t\ge c|x|$ because the bump
has compact support. Thus $M_\Phi a(x)\le C|x|^{-d-1}$,
an integrable tail. Translation and dilation give a uniform
bound on every atom; summing absolute coefficients proves
the maximal estimate.

For Hilbert truncations the local estimate again follows from
their uniformly bounded $L^2$ multipliers. Off a doubled atom
interval of radius $R$, the kernel-difference contribution has
integral bounded by $C R\int_{|x|>2R}|x|^{-2}dx$.
The truncation may cut only part of that interval; the extra
$x^{-1}$ times its partial integral is supported where
$||x|-\varepsilon|\le R$. Its absolute integral is at most
$C R/\max(R,\varepsilon)$, hence uniformly bounded as well.
For a fixed atom, once $\varepsilon<R$ its far-field formula
is untruncated; the common integrable tail bound and local
$L^2$ convergence imply $L^1$ convergence. Approximate a
general atomic sum by a finite sum and use the uniform operator
bound on the remainder. This proves convergence for the entire
$H^1$, including precisely the truncated-integral assertion
of \cite{B09}, Theorem 5.4.
\end{proof}

\begin{theorem}[The dual of atomic $H^1$]
\label{fa2:p4:hardy-bmo-dual}
Every continuous linear functional on $H^1_{\rm at}(\mathbb R^d)$
is pairing with a BMO function, unique modulo constants.
Its norm and the BMO seminorm are equivalent with dimension
constants. Pairing is first the integral on finite atom sums
and then its continuous extension; an absolutely convergent
Lebesgue integral $\int fg$ for every pair is not asserted.
\end{theorem}
\begin{proof}
For a bounded functional $L$, its restriction to mean-zero
$L^2$ functions supported in a ball $B$ has norm at most
$C|B|^{1/2}\pfnorm L$, by the $2$-atom theorem.
Hilbert representation gives $h_B\in L^2_0(B)$ with
$L(f)=\langle h_B,f\rangle$. Put $g_B=\overline{h_B}$;
then $L(f)=\int_Bfg_B$, with the same norm bound.
On nested balls the
representatives differ by constants, since they have the
same integrals against every mean-zero $L^2$ test.
Adjust these constants along an exhaustion to obtain one
$g\in L^2_{\rm loc}$. For any other ball its difference
from the local representative is likewise constant.
Consequently
$|B|^{-1}\int_B|g-g_B^{\rm mean}|
\le |B|^{-1/2}\pfnorm{g-g_B^{\rm mean}}_2
\le C\pfnorm L$, proving BMO and the representation on
finite atom sums, which are dense in $H^1$.

Conversely, for bounded BMO $g$ and an atom $a$,
$|\int ag|=|\int a(g-g_B)|\le\pfnorm g_{\rm BMO}$.
This gives the bound for finite sums and hence their completion.
For general real BMO $g$, truncate its values by a real
one-Lipschitz clipping map. The BMO seminorm of the clipped
function is at most twice that of $g$, since comparison with
the clipped mean bounds its mean oscillation. These bounded
functions converge locally in $L^1$ to $g$.
Their functionals therefore converge on every finite atom sum
and are uniformly bounded; density extends the limit to all
$H^1$. Treat real and imaginary parts for complex $g$.
This proves existence without presuming convergence of the
unbounded product integral. If two representatives agree on
all atoms, their integrals against all bounded compactly
supported mean-zero functions vanish; testing within each
ball makes their difference a constant there, hence globally.
This is the actual endpoint duality argument, rather than a
consequence claimed merely from the two space definitions.
\end{proof}

\chapter{Positive Representations, Tauberian Density, and Hardy Subspaces}
\label{fa2:p4:representation-chapter}

\begin{theorem}[Herglotz--Riesz and positive harmonic representation]
\label{fa2:p4:herglotz}
If $F$ is holomorphic in the disk and $\operatorname{Re}F\ge0$,
there are a unique positive finite measure $\mu$ on the circle
and a unique real $c$ such that
\[
 F(z)=ic+\int_{\mathbb T}\frac{\zeta+z}{\zeta-z}\,d\mu(\zeta).
\]
Here $c=\operatorname{Im}F(0)$ and $\mu(\mathbb T)=\operatorname{Re}F(0)$.
Every positive harmonic function is the Poisson integral of a
unique positive finite measure.
\end{theorem}
\begin{proof}
Put $u=\operatorname{Re}F$. The measures
$u(r\zeta)\,dm(\zeta)$ have constant mass $u(0)$, by the
harmonic mean-value property. Take a weak-star convergent subnet
as $r\uparrow1$. For $|z|<r$, the Poisson formula on the
radius-$r$ disk expresses $u(z)$ with kernel $P_{z/r}$.
It follows either by the Fourier series of a harmonic function
on a slightly larger disk, or by uniqueness of the harmonic
Dirichlet problem from the maximum principle. Pass to the limit:
$P_{z/r}\to P_z$ uniformly on the circle, giving
$u(z)=\int P_z\,d\mu$.
The integral of $(\zeta+z)/(\zeta-z)$ is holomorphic with this
real part, so its difference from $F$ is a purely imaginary
constant; evaluation at zero identifies it. Conversely the
displayed kernel has positive real part $P_z$ and defines such
a function by differentiation under the finite integral.
Uniqueness follows by its circle Fourier coefficients: Poisson
integration multiplies the $n$th coefficient by $r^{|n|}$,
so all moments of the difference measure vanish. Trigonometric
polynomials are dense in continuous periodic functions, by the
positive Fej\'er approximate identities, and Riesz--Markov then
gives zero measure. A positive harmonic function has a harmonic
conjugate on the disk by integration of its closed differential,
and hence the same argument applies.
\end{proof}

\begin{corollary}[Carath\'eodory's circle moment theorem]
\label{fa2:p4:circle-moment}
A two-sided sequence with $c_{-n}=\overline{c_n}$ is the moment
sequence $c_n=\int\zeta^n d\mu$ of a finite positive circle
measure exactly when all Toeplitz quadratic forms
$\sum_{j,k}\overline a_ja_kc_{k-j}$ are nonnegative.
\end{corollary}
\begin{proof}
Necessity is the integral of a squared polynomial modulus.
For sufficiency the Fej\'er-weighted polynomials
$\sum_{|n|\le N}(1-|n|/(N+1))c_ne^{-in\theta}$ are nonnegative:
they are $1/(N+1)$ times the Toeplitz quadratic form of the
vector $(1,e^{-i\theta},\ldots,e^{-iN\theta})$, with its sign
adjusted to the displayed moment convention. Their integrals
are $c_0$. Weak-star compactness gives a positive measure limit,
and for each fixed $n$ its moment tends to $c_n$.
If $c_0=0$, the $2\times2$ minors already force every $c_n=0$.
Uniqueness follows from trigonometric density. This circle theorem
is distinct from the real-line Hamburger/Stieltjes problems in
Part III and from Krein's weaker even-test positivity condition.
\end{proof}

\begin{theorem}[Bochner's positive-definite representation]
\label{fa2:p4:bochner-positive}
A continuous positive-definite function $\varphi$ on $\mathbb R^d$
is the Fourier transform, with the sign changed if desired, of
a unique finite positive measure; its mass is $\varphi(0)$.
Positive definiteness means
$\sum\overline a_ja_k\varphi(x_k-x_j)\ge0$ for all finite
choices. Continuity at zero alone suffices and implies uniform
continuity everywhere.
\end{theorem}
\begin{proof}
The finite Gram construction and Cauchy--Schwarz give
$|\varphi(x)|\le\varphi(0)$ and
$|\varphi(x+h)-\varphi(x)|^2
\le2\varphi(0)(\varphi(0)-\operatorname{Re}\varphi(h))$.
Thus continuity at zero implies the asserted uniform continuity.
For $\varepsilon>0$ put
$\varphi_\varepsilon(x)=\varphi(x)e^{-\pi\varepsilon|x|^2}$.
It is integrable and positive definite: the Gaussian is positive
definite by its positive Gaussian Fourier density, and products
of finite Gram matrices are Gram matrices of tensor vectors.
Riemann sums of its quadratic forms show that
\[
 \iint\varphi_\varepsilon(x-y)g(y)\overline{g(x)}dxdy\ge0
 \quad(g\in C_c^\infty).
\]
Plancherel gives nonnegativity of
$\widehat\varphi_\varepsilon$: if its continuous real function
were negative near a point, modulated increasingly broad smooth
test functions would concentrate $|\widehat g|^2$ there and
contradict this inequality. Gaussian damping and Fubini give
$\int\widehat\varphi_\varepsilon(\xi)e^{-\pi\delta|\xi|^2}d\xi
\to\varphi(0)$ as $\delta\downarrow0$.
Monotone convergence shows that this nonnegative transform is
integrable, of total mass $\varphi(0)$. Inversion gives its
positive-density representing measure $\mu_\varepsilon$.

These measures are uniformly tight as $\varepsilon\downarrow0$.
Indeed
\[
 \int(1-e^{-\pi r^2|\xi|^2})d\mu_\varepsilon
 =\varphi(0)-\int r^{-d}e^{-\pi|x|^2/r^2}
                                      \varphi_\varepsilon(x)dx
\]
tends uniformly to zero as $r\downarrow0$, using their uniform
continuity at zero and boundedness. Outside $|\xi|\ge1/r$ the
factor on the left is bounded below by $1-e^{-\pi}$.
Take a weak-star measure limit along $\varepsilon\downarrow0$.
Tightness retains its full mass and permits testing bounded
continuous characters by cutting them off on a large compact set.
Their integrals tend to $\varphi(x)$, proving representation.
Conversely a positive measure gives positive definiteness by
integrating $|\sum a_je^{2\pi ix_j\cdot\xi}|^2$, and continuity
by dominated convergence. If two finite measures have the same
transform, their convolutions with every Gaussian are zero by
inversion. Approximation against compact continuous tests then
proves equality. All limiting measures are scalar Radon measures
on Euclidean space, where finite Borel regularity is available.
\end{proof}

\begin{proposition}[Poisson summation]
\label{fa2:p4:poisson-summation}
For $f\in\mathcal S(\mathbb R^d)$,
$\sum_{n\in\mathbb Z^d}f(n)=\sum_{n\in\mathbb Z^d}\widehat f(n)$,
with absolute convergence. For an invertible real matrix $A$,
\[
 \sum_{n\in\mathbb Z^d}f(An)
 =|\det A|^{-1}\sum_{n\in\mathbb Z^d}
                                  \widehat f(A^{-\mathsf T}n).
\]
\end{proposition}
\begin{proof}
The periodization $\sum_nf(x+n)$ and each of its derivatives
converge uniformly on a fundamental cube by Schwartz decay.
Its $k$th Fourier coefficient is $\widehat f(k)$, by absolute
Fubini. These coefficients decay faster than any power, so their
Fourier series converges absolutely and uniformly to that
periodization. Evaluation at zero proves the identity.
Apply it to $f(Ax)$ and change variables in the Fourier integral
for the lattice form. This proves the mixed current component
from the trace/number-theory source families without claiming
their separate arithmetic applications.
\end{proof}

\begin{lemma}[Local Fourier flattening]
\label{fa2:p4:wiener-local}
For $f\in L^1$, a frequency $\xi_0$, and $\varepsilon>0$,
there is $h\in L^1$ with $\pvnorm h_1<\varepsilon$ and
$\widehat h(\xi)=\widehat f(\xi_0)-\widehat f(\xi)$ on a
neighborhood of $\xi_0$.
\end{lemma}
\begin{proof}
Choose a Schwartz function $g$ whose transform equals one near
zero. Put $g_R(x)=R^{-d}e^{2\pi ix\cdot\xi_0}g(x/R)$ and
$h_R=\widehat f(\xi_0)g_R-f*g_R$.
Its transform has the required equality near $\xi_0$.
Moreover its $L^1$ norm is bounded by
\[
 \int|f(y)|\pvnorm{g(\cdot)-g(\cdot-y/R)}_1dy\longrightarrow0,
\]
by translation continuity and dominated convergence. This proves
the small-norm correction needed for local inversion.
\end{proof}

\begin{theorem}[Wiener's translation-density and Tauberian theorem]
\label{fa2:p4:wiener-tauberian}
A closed translation-invariant subspace $Y\subset L^1(\mathbb R^d)$
equals $L^1$ if its members' transforms have no common zero.
In particular translates of $K\in L^1$ span densely exactly when
$\widehat K$ vanishes nowhere. If $\phi\in L^\infty$ and
$K*\phi(x)\to a\int K$ as $|x|\to\infty$, with such a $K$,
then $f*\phi(x)\to a\int f$ for every $f\in L^1$.
If in addition $\phi$ is slowly oscillating at infinity, it
itself tends to $a$. The analogous one-sided statement on
$\mathbb R$ holds with $x\to+\infty$.
Here slow oscillation means that for every $\varepsilon>0$
there are $\delta,R>0$ with
$|\phi(x-h)-\phi(x)|<\varepsilon$ whenever
$|h|<\delta$ and $|x|>R$ (respectively $x>R$).
\end{theorem}
\begin{proof}
A closed translation-invariant space is a convolution ideal:
$f*g$ is the Bochner integral of translated $f$ against $g$,
and its simple integral sums lie in $Y$.
At a frequency where $\widehat f\ne0$, normalize it to one.
The local lemma gives $\pvnorm h_1<1$ with
$\widehat h=1-\widehat f$ on a neighborhood $U$.
For any Schwartz $\psi$ with $\pvsupp\widehat\psi\subset U$,
Fourier inversion yields
$\psi-h*\psi=f*\psi$. The geometric convolution inverse
$\delta+\sum_{n\ge1}h^{*n}$ belongs to the unitized $L^1$
algebra, so $\psi\in Y$.
If there is no common zero, a finite frequency partition of
unity shows that every Schwartz function with compact Fourier
support lies in $Y$. These functions are dense in $L^1$:
convolve a compact smooth function with a Schwartz approximate
identity whose Fourier transform is compactly supported, and
then use compact smooth density. Hence $Y=L^1$.
If a common zero exists, Fourier evaluation there is a nonzero
bounded functional annihilating $Y$, proving the converse.

For the Tauberian assertion, the functions $f$ satisfying the
claimed convolution limit form a closed translation-invariant
subspace of $L^1$, by the bound
$\pvnorm{f*\phi}_\infty\le\pvnorm f_1\pvnorm\phi_\infty$.
It contains $K$, so is all of $L^1$. For slow oscillation,
choose a fixed nonnegative integral-one kernel supported so
close to zero that $\phi$ differs little from its convolution
outside a sufficiently large compact set. The convolution tends
to $a$ by the first assertion, and the oscillation error is
arbitrary. The one-sided proof uses the identical argument.
Prime-number and renewal applications require their own extra
arithmetic or probability hypotheses and are not identified
with this density theorem.
\end{proof}

\section{The disk Hardy space and invariant subspaces}

\begin{definition}[The analytic Hilbert space]
\label{fa2:p4:hardy-disk}
$H^2(\mathbb D)$ consists of $f(z)=\sum_{n\ge0}a_nz^n$ with
$\sum|a_n|^2<\infty$. Its boundary function is the $L^2$ Fourier
sum of those coefficients, and radial values are its Poisson
extension. An inner function is bounded holomorphic with
boundary modulus one almost everywhere. An outer function is
a cyclic vector for multiplication by $z$ in $H^2$.
\end{definition}

\begin{theorem}[Beurling's theorem]
\label{fa2:p4:beurling}
Every nonzero closed subspace of $H^2$ invariant under multiplication
by $z$ is $\vartheta H^2$ for an inner $\vartheta$, unique up
to a unimodular constant. Invariance under all bounded analytic
multipliers is equivalent to this shift invariance.
\end{theorem}
\begin{proof}
Let $M$ be such a subspace. The isometry $f\mapsto zf$ has
closed range $zM$, and $zM\ne M$: otherwise every member of
$M$ would have a zero of every order at zero.
Choose a unit vector $\vartheta\in M\ominus zM$.
For every $n\ge1$, $\langle\vartheta,z^n\vartheta\rangle=0$.
Thus all nonconstant Fourier coefficients of $|\vartheta|^2$
vanish and its integral is one, so its boundary modulus is one
a.e. Poisson majorization of the subharmonic function
$|\vartheta|^2$ gives $|\vartheta(z)|\le1$, proving innerness.
For any other $h\in M\ominus zM$, the analogous two
orthogonalities show that $\overline\vartheta h$ has only its
constant Fourier coefficient. Hence $h$ is a scalar multiple
of $\vartheta$, so this wandering subspace has dimension one.
Iterate $M=(M\ominus zM)\oplus zM$.
The remainder after removing all shifted copies belongs to
$\bigcap_nz^nM\subset\bigcap_nz^nH^2=\{0\}$.
Therefore $M=\bigoplus_{n\ge0}\mathbb Cz^n\vartheta
=\vartheta H^2$. Multiplication by $\vartheta$ is an isometry
by boundary modulus, so its range is closed.
If two inner generators give the same subspace, each divides
the other in $H^2$; the quotient and reciprocal are inner,
and the maximum principle makes them unimodular constants.
The representation immediately gives invariance under bounded
analytic multipliers; the converse includes the multiplier $z$.
\end{proof}

\begin{theorem}[Inner--outer factorization and the singular measure]
\label{fa2:p4:hardy-factorization}
Every nonzero $f\in H^2$ factors into an inner and an outer
function, uniquely up to opposite unimodular constants.
Its outer factor can be written
\[
 O(z)=c\exp\!\left(\int_{\mathbb T}
       \frac{\zeta+z}{\zeta-z}\log|f^*(\zeta)|dm(\zeta)\right).
\]
An inner function is a Blaschke product times
$c\exp(-\int(\zeta+z)/(\zeta-z)d\sigma(\zeta))$,
where $\sigma$ is a finite positive singular circle measure.
Its zeros and this singular measure are unique.
\end{theorem}
\begin{proof}
Jensen's formula implies $\log|f^*|\in L^1$.
Here is the needed integrability justification: remove a finite
zero at zero, use the lower bound on the circular mean of
$\log|f(r\zeta)|$, and use
$\log^+|f(r\zeta)|\le C+|f(r\zeta)|^2$.
The positive parts are uniformly integrable because the radial
values converge in $L^2$; Fatou then bounds the negative part
of the boundary logarithm. Poisson--Jensen and the same
positive-part convergence/negative-part Fatou argument yield
$\log|f(z)|\le P_z\log|f^*|$.
Thus the displayed zero-free $O$ has $|f/O|\le1$.
Jensen's inequality gives $|O(z)|^2\le P_z|f^*|^2$, so
$O\in H^2$. Its radial boundary modulus is $|f^*|$, by
Poisson differentiation of the integrable logarithm, and hence
$f/O$ is inner. If $O$ had a nonconstant inner divisor $J$,
apply the same logarithmic inequality to $O/J$; equality of
its boundary modulus with that of $O$ would give
$-\log|J(z)|\le0$. This contradicts inner boundedness unless
$J$ is constant. Beurling applied to the cyclic subspace of
$O$ now proves that $O$ is cyclic. It also proves uniqueness.

For an inner function, Jensen gives the Blaschke sum
$\sum(1-|z_j|)<\infty$. The canonical factors therefore
converge uniformly on compact subsets. For the resulting
product $B$, its circular mean of $-\log|B(r\zeta)|$ tends
to zero by summing $-\log\max(r,|z_j|)$ and dominated
convergence. Since bounded analytic functions have $L^2$
boundary values and radial Poisson limits a.e., Fatou makes
$|B^*|=1$. Dividing successively by finite Blaschke products
preserves the bound one, by the maximum principle; the limit
$S=I/B$ is thus zero-free inner. Choose a holomorphic logarithm
and apply Herglotz to the positive-real-part function
$-\log S$. Its real part has boundary value zero a.e., so
its representing measure is singular: the Poisson boundary
limit of a finite measure is its absolutely continuous density,
by Lebesgue differentiation. This proves the representation.
Uniqueness follows from the zeros, then from the unique Herglotz
measure and the imaginary constant of the logarithm.
\end{proof}

\begin{lemma}[Half-plane Hardy functions and a causal support]
\label{fa2:p4:hardy-halfplane}
An analytic function $F$ in the upper half-plane satisfies
$\sup_{y>0}\int_{\mathbb R}|F(x+iy)|^2dx<\infty$ exactly when
\[
 F(z)=\int_0^\infty e^{2\pi izt}f(t)dt
 \quad\text{for some }f\in L^2(0,\infty).
\]
The integral is absolutely convergent at $y>0$ and has boundary
values in $L^2$. The Cayley transform
$z=i(1+w)/(1-w)$ gives the isometry to disk $H^2$ by
$F\mapsto2\sqrt\pi\,F(i(1+w)/(1-w))/(1-w)$.
For $F\not\equiv0$, the corresponding inner factor has a distinguished
factor $e^{iaz}$ with uniquely determined $a\ge0$, the singular measure's
mass at $w=1$. If $F_1,F_2$ are nonzero and $F_1,F_2,F_1F_2$ belong
to half-plane $H^2$, these coefficients add under multiplication.
\end{lemma}
\begin{proof}
The forward integral has horizontal $L^2$ norm
$\pfnorm{e^{-2\pi yt}f(t)}_2$ by Plancherel, and converges
in $L^2$ to its boundary as $y\downarrow0$.
Conversely take inverse Fourier transforms of horizontal slices,
using the positive-exponential convention in the display.
The Cauchy--Riemann equation implies in distributions
$\partial_y f_y(t)=-2\pi t f_y(t)$.
To justify solving this identity, restrict to bounded $t$
intervals and multiply by $e^{2\pi yt}$; its distributional
$y$ derivative is zero, hence it is independent of $y$.
Thus $f_y=e^{-2\pi yt}f_0$ for a locally $L^2$ function.
The uniform slice bound as $y\to\infty$ forces $f_0=0$ for
$t<0$, and Fatou as $y\downarrow0$ gives $f_0\in L^2$.
This proves representation and strong boundary convergence.

For the Cayley statement, on the boundary the substitution
$x=i(1+e^{i\theta})/(1-e^{i\theta})$ has
$dx/d\theta=2/|1-e^{i\theta}|^2$.
This proves the norm identity including its constant.
It also maps the analytic subspaces onto one another: the disk
Cauchy integral formula, after this substitution and cancellation
of the displayed linear factors, becomes
$F(z)=(2\pi i)^{-1}\int_\mathbb R F(x)/(x-z)dx$.
The integral is defined by $L^2$ pairing with $(x-z)^{-1}$.
Its Fourier transform in the boundary variable is the projection
onto the causal half-line, as follows by closing the exponential
contour in the upper or lower half-plane. Equivalently this is
the preceding Cauchy--Riemann slice computation. It follows
that both transformed analytic spaces have precisely these
boundary functions, proving onto isometry rather than merely
a change-of-variable identity for norms.

For nonzero $F$, the singular mass $a$ at $1$ in the disk gives
$\exp(-a(1+w)/(1-w))=e^{iaz}$.
The Cayley prefactor $(1-w)^{-1}$ is zero-free analytic with
integrable boundary logarithm and exact Poisson-log modulus;
its inverse $1-w$ is bounded outer. The prefactor itself is
not being called a member of disk $H^2$. These identities show
that it contributes no singular atom.
For products write the disk correspondence as $U$. Then
$U(F_1F_2)=(1-w)U(F_1)U(F_2)/(2\sqrt\pi)$.
If $U(F_j)=I_jO_j$, the zero-free analytic product
$H=(1-w)O_1O_2/(2\sqrt\pi)$ has exact Poisson-log
modulus, and $|H^*|=|U(F_1F_2)^*|\in L^2$.
Jensen gives $|H(w)|^2\le P_w|H^*|^2$, hence $H\in H^2$.
The same exact-log argument used in the outer-factor theorem
shows $H$ has no nonconstant inner divisor, so it is outer.
Thus the product's inner factor is exactly $I_1I_2$.
Uniqueness of the Herglotz measure for their
zero-free parts makes the singular atoms, hence the coefficients
$a$, additive. This argument avoids an invalid multiplication
rule for arbitrary pointwise logarithmic limsups near zeros.
\end{proof}

\begin{theorem}[Titchmarsh convolution support theorem]
\label{fa2:p4:titchmarsh}
For nonzero compactly supported distributions $u,v$ on the real
line, let $a_u,b_u$ be the left and right endpoints of their
supports. Then
$a_{u*v}=a_u+a_v$ and $b_{u*v}=b_u+b_v$.
In particular a convolution cannot acquire an initial open
interval of zeros by cancellation beyond the sum of the two
initial support intervals. This is stronger than support inclusion.
\end{theorem}
\begin{proof}
First let $0\ne f\in L^2$ have support bounded below. Translate so it
lies in $[0,\infty)$ and use the preceding Fourier--Laplace
representation $F$, which is nonzero by injectivity. The earliest support
point of $f$ is $a/(2\pi)$, where $e^{iaz}$ is the distinguished factor of $F$.
Indeed support in $[c,\infty)$ is equivalent to
$e^{-2\pi icz}F\in H^2$ by the representation theorem and
translation. Removing any part of the factor $e^{iaz}$ keeps
the remaining inner times outer factor in $H^2$, so support
begins no earlier than $a/(2\pi)$. Conversely if the shifted
function is in $H^2$, multiplication by its exponential inner
factor and uniqueness force $a\ge2\pi c$.
For nonzero $f,g\in L^1\cap L^2$ with supports bounded below, translate
them into $[0,\infty)$. By Young, $f*g\in L^2$. Its Laplace transform
is the product of their nonzero analytic transforms, hence is nonzero.
The preceding lemma gives left-endpoint addition, which respects translation.

To reach compact distributions choose a sufficiently large integer
$m$ and the causal kernel
$h_m(t)=e^{-t}t_+^{m-1}/(m-1)!$.
It satisfies $(\partial_t+1)^mh_m=\delta_0$.
For $m$ larger than the finite order of a compact distribution
plus two, $U=h_m*u$ is a continuous $L^1\cap L^2$ function:
the defining finite-order estimate applies to the continuous
derivatives of $h_m$, and outside the fixed support interval
they are exponentially decreasing polynomials.
The convolution is legitimate because the supports are bounded
below, making addition proper on them. Causality gives
$a_U\ge a_u$, and recovery by $(\partial+1)^m$ gives
$a_u\ge a_U$; hence they agree. Do the same for $v$, giving
$V=h_n*v$. Both $U$ and $V$ are nonzero because applying their
respective differential operators recovers $u$ and $v$.
The function case gives a nonzero $U*V$ with left endpoint $a_U+a_V$.
Since $U*V=h_{m+n}*(u*v)$, the distribution $u*v$ is nonzero.
Causality and recovery identify its left endpoint with that of
$U*V$, which is $a_U+a_V=a_u+a_v$.
Reflect the two compact distributions in the origin and apply
the left-endpoint result to obtain the right-endpoint identity.
This proof supplies the analytic factorization bridge behind
the theorem, rather than deriving equality from the weaker
elementary support inclusion.
\end{proof}

\begin{corollary}[The source half-line $L^1$ version]
\label{fa2:p4:titchmarsh-halfline}
For nonzero $f,g\in L^1(0,\infty)$ the lower support endpoints
add under convolution. If $F(z)=\int_0^\infty e^{2\pi izt}f(t)dt$,
then $f=0$ on $(0,c)$ if and only if $e^{-2\pi icz}F(z)$ is
bounded on the upper half-plane.
\end{corollary}
\begin{proof}
For endpoint addition truncate both functions at a height larger
than the sum of their lower endpoints. The truncated convolutions
agree with the original convolution below that height by causality,
so the compact-support theorem proves the assertion locally at
its lower endpoint. For the bounded-transform criterion the forward
direction is the Laplace transform of $f(t+c)$.
Conversely test against $\phi\in C_c^\infty((0,c-\delta))$.
Fourier inversion and absolute Fubini give
$\int f\phi=\int_{\mathbb R}F(x)\widehat\phi(x)dx$.
Shift to $x+iy$ by Cauchy's theorem. The assumed bound on
$e^{-2\pi icz}F$ contributes $e^{-2\pi cy}$, while the test
transform contributes at most a polynomial in $y$ times
$e^{2\pi(c-\delta)y}(1+|x|)^{-N}$.
The rectangle sides vanish for $N>2$, and the shifted integral
tends to zero as $y\to\infty$. Thus all such tests are
annihilated, proving the vanishing interval. This is the precise
$L^1/H^\infty$ version in \cite{B08}, 38.3, separately from
the $L^2/H^2$ support bridge used above.
\end{proof}

\begin{proposition}[Hardy multipliers and divisibility]
\label{fa2:p4:hardy-divisibility}
Every bounded operator on disk $H^2$ commuting with the shift
is multiplication by a unique $b\in H^\infty$, with operator
norm $\pfnorm b_\infty$. For bounded analytic functions, an
inner $\vartheta$ divides $b$ in $H^\infty$ exactly when
$bH^2\subset\vartheta H^2$.
An outer $O\in H^\infty$ divides a bounded analytic $b$
exactly when the boundary quotient $b^*/O^*$ is essentially
bounded. If $b,c$ are not both zero, their greatest common inner
divisor, up to a unimodular constant, is the Beurling generator of
$\overline{bH^2+cH^2}$. If $b=c=0$, this closed sum is $\{0\}$
and has no inner Beurling generator; there is no greatest common
inner divisor within the class of inner functions.
\end{proposition}
\begin{proof}
For a commuting operator put $b=T1$. On polynomials $Tp=bp$.
The reproducing vector $k_w(z)=(1-\overline wz)^{-1}$ then
satisfies $T^*k_w=\overline{b(w)}k_w$, by bounded point
evaluation and polynomial density. Hence $|b(w)|\le\pfnorm T$.
Multiplication by bounded $b$ has norm at most $\pfnorm b_\infty$
by boundary integration, and the same kernel vectors give the
reverse bound. This proves the commutant assertion.
For the inner assertion the inclusion gives $b=\vartheta h$
with $h\in H^2$; boundary modulus makes $h$ essentially bounded,
so Poisson majorization puts it in $H^\infty$.
For the outer assertion its logarithmic formula and the
Poisson--Jensen inequality for $b$ give
$\log|b/O|\le P\log|b^*/O^*|\le\log C$.
Outer functions have no zeros, so the quotient is analytic and
bounded. The converse is immediate from boundary values.
For $b,c$ not both zero, the closed sum is a nonzero shift-invariant
subspace. Inclusion of generated subspaces translates inner
divisibility into reverse inclusion; Beurling and taking the
closed sum prove both divisibility and the greatest property.
If $b=c=0$, every inner function divides both. No inner function
can generate the zero subspace, because it belongs to its own
generated subspace and is nonzero. Nor can an inner $\vartheta$
be a greatest common divisor: $z\vartheta$ is another common
divisor, and it cannot divide $\vartheta$ since $1/z$ is not
analytic at zero.
These conclusions include the source result-bearing commutant
exercise and its algebraic factorization consequences, without
claiming arbitrary ideals of $H^\infty$ are principal.
\end{proof}

\section{The sharp trace spaces}

\begin{definition}[Fractional difference norm]
\label{fa2:p4:fractional-sobolev}
For $0<s<1$ and $1\le p<\infty$, set
\[
 [g]_{s,p}^p=\int_{\mathbb R^n}
       \frac{\pfnorm{g(\cdot+h)-g}_p^p}{|h|^{n+sp}}dh,
 \qquad W^{s,p}=\{g\in L^p:[g]_{s,p}<\infty\}.
\]
The difference integrand is assigned value zero at $h=0$.
For $1<s<2$ require first derivatives in $W^{s-1,p}$ as well.
For $0<s<1$, bi-Lipschitz changes of variables with their usual
almost-everywhere Jacobian bounds give equivalent order-$s$ norms.
Multiplication by a bounded globally Lipschitz scalar function is
bounded on $W^{s,p}(\mathbb R^n)$. In particular, smooth bounded
multipliers with bounded first derivatives have this property;
smoothness and boundedness alone are insufficient on the whole space.
For $0<s<1$ these coordinate and multiplier statements define the
corresponding spaces on a compact $C^1$ boundary by a finite atlas
and bounded Lipschitz cutoffs, with equivalent norms on overlaps.
\end{definition}
\begin{proof}
Translation and Tonelli identify the displayed seminorm with
$\iint |g(x)-g(y)|^p/|x-y|^{n+sp}\,dx\,dy$. For a bi-Lipschitz
coordinate map $\Phi$, change both variables; the distance
comparisons and the almost-everywhere upper Jacobian bounds for
$\Phi$ and $\Phi^{-1}$ give the estimates in both directions.

Let $M=\|a\|_\infty$ and let $L$ be a global Lipschitz constant.
Splitting a product difference and using
$|a(x+h)-a(x)|\le\min\{L|h|,2M\}$ gives
\[
 [ag]_{s,p}^p\le 2^{p-1}\left(M^p[g]_{s,p}^p+
 \|g\|_p^p\int_{\mathbb R^n}
 \frac{\min\{L|h|,2M\}^p}{|h|^{n+sp}}\,dh\right).
\]
For $n\ge1$, splitting the last integral at $|h|=1$ bounds it by
\[
 |\mathbb S^{n-1}|\left(\frac{L^p}{p(1-s)}+
                    \frac{(2M)^p}{sp}\right)<\infty.
\]
Also $\|ag\|_p\le M\|g\|_p$, proving bounded multiplication,
including $p=1$. The zero-dimensional case has zero difference
seminorm and needs only the $L^p$ bound. For smooth $a$ with bounded
first derivative, the segment integral gives a global Lipschitz
constant. On compact chart supports, $C^1$ transition maps have
the required local bi-Lipschitz and Jacobian bounds after shrinking
charts, and a finite subordinate cutoff family has bounded first
derivatives. Applying the two estimates on chart overlaps proves
atlas/cutoff independence. This boundary statement concerns
$0<s<1$; no higher-order invariance under merely $C^1$ charts is used.
\end{proof}

\begin{theorem}[Sharp first trace, including its endpoint boundary]
\label{fa2:p4:sharp-trace}
For $1<p<\infty$ and $s=1-1/p$, the trace on a half-space is
a bounded surjection $W^{1,p}\to W^{s,p}(\mathbb R^{d-1})$
with a bounded linear right inverse. The same holds on bounded
$C^1$ domains. At $p=1$ its range is all $L^1$ of the boundary;
surjectivity here is proved without asserting a bounded linear
right inverse. Its kernel in all these cases is the $W^{1,p}_0$
already identified above.
\end{theorem}
\begin{proof}
For a smooth compactly supported $u$ in the half-space write
$g=u(\cdot,0)$. At height $r=|h|$,
\[
 \pfnorm{g(\cdot+h)-g}_p
 \le2\int_0^r\pfnorm{\partial_tu(\cdot,t)}_pdt
       +r\pfnorm{\nabla'u(\cdot,r)}_p.
\]
Integrate its $p$th power in $h/|h|^{d-1+sp}$.
Radial integration gives weight $r^{-p}$ because $sp=p-1$.
The gradient term integrates to its $L^p$ norm, and Hardy's
inequality
\[
 \int_0^\infty\left(r^{-1}\int_0^ra(t)dt\right)^pdr
 \le(p/(p-1))^p\int_0^\infty a(t)^pdt
\]
controls the normal term. To verify this inequality integrate
by parts the left side truncated away from zero and infinity,
drop its favorable boundary terms, and apply H\"older to
$(\int_0^ra)^{p-1}a(r)r^{1-p}$; truncation then gives the
stated bound. Together with the earlier $L^p$ trace estimate
this proves continuity into $W^{s,p}$.

For surjectivity choose a smooth integral-one tangential bump
$\rho$ and set $Eg(x,t)=\chi(t)\rho_t*g(x)$, with $\chi=1$
near zero and compactly supported in $t\ge0$.
Both tangential and $t$ derivatives of $\rho_t$ have integral
zero and size at most $Ct^{-n-1}$ on $|h|\le Ct$.
Jensen therefore gives, for either derivative,
\[
 \pfnorm{\partial Eg(\cdot,t)}_p^p
 \le Ct^{-p-n}\int_{|h|\le Ct}
                    \pfnorm{g(\cdot-h)-g}_p^pdh
\]
apart from the cutoff term, bounded by $C\pfnorm g_p^p$.
Integrating in $t$ and interchanging integrals gives the weight
$|h|^{-n-p+1}$, exactly the $W^{s,p}$ seminorm.
The undifferentiated $L^p$ norm is bounded by Young.
Thus $E$ is bounded linear, and $\rho_t*g\to g$ in $L^p$
shows its trace is $g$. Smooth approximation in the fractional
norm is legitimate: translation is continuous in its double
integral by density in the underlying product $L^p$ space;
cutoff and convolution then give smooth density.

For $p=1$ let $g_j$ be smooth approximants with
$\pfnorm{g_j-g}_1\le2^{-j}\pfnorm g_1$.
Choose $t_j\downarrow0$, $t_{j+1}<t_j/2$, so that
$t_j(\pfnorm{g_j}_{W^{1,1}}+\pfnorm{g_{j+1}}_{W^{1,1}})
\le2^{-j}\pfnorm g_1$.
Between $t_{j+1}$ and $t_j$ interpolate linearly between the
two functions; above $t_1$ cut off smoothly. The integrated
normal derivative is at most
$\sum\pfnorm{g_{j+1}-g_j}_1<\infty$, and the tangential
derivative integral is bounded by the displayed summable
quantities. The slices tend to $g$ in $L^1$, so this $W^{1,1}$
function has trace $g$. The zero datum is handled separately.
This construction depends on the approximants and is not
claimed linear. Finally flattening finitely many $C^1$ charts,
using the stated fractional change-of-variable bounds and a
partition, transfers the result to the boundary of the domain.
\end{proof}

\begin{theorem}[Two boundary jets and Green's formula]
\label{fa2:p4:second-trace}
On a bounded $C^2$ domain and $1<p<\infty$, the map
$u\mapsto(\operatorname{Tr}u,\operatorname{Tr}\partial_\nu u)$
is bounded and onto
$W^{2-1/p,p}(\partial\Omega)\times W^{1-1/p,p}(\partial\Omega)$.
For $u,v\in H^1$ the first-order integration-by-parts formula
holds with their $L^2$ traces:
\[
 \int_\Omega(u\partial_j\overline v+
                 \partial_ju\,\overline v)
 =\int_{\partial\Omega}\operatorname{Tr}u\,
              \overline{\operatorname{Tr}v}\,\nu_j.
\]
For $u\in H^2$, $v\in H^1$,
\[
 -\int_\Omega\Delta u\,\overline v
 =\int_\Omega\nabla u\cdot\overline{\nabla v}
       -\int_{\partial\Omega}\partial_\nu u\,\overline v.
\]
The stated second-jet surjectivity does not include the integer
fractional-order endpoint $p=1$ without a separate trace-space
analysis.
\end{theorem}
\begin{proof}
In flat coordinates the first derivative trace theorem applied
to all derivatives gives the bound. For the right inverse take
an even tangential mollifier with integral one. For boundary
data $g\in W^{1+s,p}$, $h\in W^{s,p}$, use near zero
$u(x,t)=\rho_t*g(x)+t\rho_t*h(x)$, with a cutoff for large $t$.
Its boundary value is $g$ and its normal derivative tends to
$h$: evenness cancels first moments, so
$\partial_t\rho_t*g\to0$ for $g\in W^{1,p}$, while
$t\partial_t\rho_t*h\to0$ and $\rho_t*h\to h$ in $L^p$.
Here $h$ is the inward flat derivative $\partial_tu$;
for the outward derivative the flat coefficient below is $a=-1$.
For completeness of the $W^{2,p}$ estimate, each second
derivative of the first term is a scale-$t$ mean-zero first
derivative kernel applied to $\nabla g$, after integration by
parts. Its bound is the same
$t^{-p-n}\int_{|r|\le Ct}\pfnorm{\nabla g(\cdot-r)-\nabla g}_p^pdr$
as in the preceding right-inverse proof. Each second derivative
of $t\rho_t*h$ has the identical bound with $h$, using its
mean-zero scale-$t$ kernel of order $t^{-1}$.
Integrating proves the required fractional-norm bound.
For curved charts the prescribed normal derivative equals
$a(x')\partial_tu+b(x')\cdot\nabla'g$, with $a$ bounded
away from zero. Replace $h$ by $(h-b\cdot\nabla'g)/a$;
the $C^1$ coefficients preserve $W^{s,p}$. $C^2$ flattening
preserves $W^{2,p}$, so a finite patch construction gives
surjectivity. The complex conjugate in Green's formula is
written explicitly to match Hilbert convention.
The first-order identity holds for smooth functions by ordinary
integration by parts, and $W^{1,2}$ smooth-up-to-boundary density
passes it to $H^1$ pairs. Apply that already proved identity
to $\partial_ju\in H^1$ and $v\in H^1$ and sum in $j$
to obtain the displayed second formula. Thus no unproved
$W^{2,2}$ boundary density assertion is required here.
\end{proof}

\section{A compact Besicovitch set}

\begin{theorem}[Zero area with a unit segment in every direction]
\label{fa2:p4:besicovitch}
There is a compact planar set of Lebesgue measure zero containing
a line segment of length one in every unoriented direction.
\end{theorem}
\begin{proof}
Choose a sequence $c_k\in[0,1]$ as follows. At stage $r$ list
all dyadics $j/2^r$, $0\le j\le2^r$, repeating each one
consecutively $r$ times, and concatenate the stages.
For $a=(a_k)\in\{0,1\}^{\mathbb N}$ define
\[
 K=\left\{\left(x,\sum_{k\ge1}a_k2^{-k}(x-c_k)\right):
             0\le x\le1,\ a\in\{0,1\}^{\mathbb N}\right\}.
\]
The series converges uniformly on the compact product of the
interval and the binary product space. Thus $K$ is compact.
Every slope $m\in[0,1]$ has a binary expansion
$m=\sum a_k2^{-k}$. For these fixed digits the full segment
$y=mx-\sum a_k2^{-k}c_k$, $0\le x\le1$, lies in $K$;
its length is at least one.

Fix $x$. For every $N$, its vertical fiber is covered by
$2^N$ intervals, one per prefix, each of length at most
$\sum_{k>N}2^{-k}|x-c_k|$. At stage $r$ choose a listed
dyadic $q$ within $2^{-r}$ of $x$ and let $N$ be the index
immediately before its block of $r$ repetitions. Then the
total covering length is at most
\[
 2^N\sum_{k>N}2^{-k}|x-c_k|
 \le |x-q|+2^{-r}\le2^{1-r}.
\]
Hence every vertical fiber has one-dimensional measure zero.
Fubini for the compact Borel set $K$ gives planar measure zero.
Its segment directions cover angles $[0,\pi/4]$.
Take the union of its four rotations through $j\pi/4$,
$j=0,1,2,3$. This remains compact and null and covers all
unoriented angles in $[0,\pi)$. Each original segment contains
a unit subsegment. This construction needs no unstated
small-area triangle or limiting-area lemma.
\end{proof}

\section{Homogeneous extensions and explicit fundamental solutions}

\begin{theorem}[Homogeneous extension, residues, and the critical obstruction]
\label{fa2:p4:homogeneous-extension}
Let $k\in C^\infty(\mathbb R^d\setminus\{0\})$ satisfy
$k(rx)=r^\lambda k(x)$ for $r>0$, where $\lambda\in\mathbb C$.
If $\lambda\notin\{-d,-d-1,\ldots\}$, there is a unique
homogeneous distribution of degree $\lambda$ agreeing with $k$ off zero.
For $\lambda=-d-m$, $m\in\mathbb N_0$, such an extension exists
if and only if
\[
 \int_{S^{d-1}}k(\theta)\theta^\alpha\,d\sigma(\theta)=0
 \quad (|\alpha|=m).
\]
When it exists, the ambiguity is exactly
$\sum_{|\alpha|=m}c_\alpha\partial^\alpha\delta_0$.
In general there is a tempered finite-part extension $F$ with
\[
 F_a=a^{-d-m}(F+(\log a)R_m),\qquad
 R_m=\sum_{|\alpha|=m}\frac{(-1)^m}{\alpha!}
       \left(\int_{S^{d-1}}k(\theta)\theta^\alpha d\sigma\right)
                                      \partial^\alpha\delta_0.
\]
Here $u_a(\phi)=a^{-d}u(\phi(\cdot/a))$ is the dilation
whose action on functions is $u_a(x)=u(ax)$.
\end{theorem}
\begin{proof}
For $\operatorname{Re}z>-d$ define the degree-$z$ family
$I(z)(\phi)=\int k(x)|x|^{z-\lambda}\phi(x)dx$.
Write the integral as $|x|\le1$ and $|x|>1$. The latter is
entire in $z$ for each Schwartz test, locally uniformly by rapid
decay. In the former subtract the Taylor polynomial of degree
$N-1$ at zero. Polar coordinates give the precise continuation
\[
 I(z)(\phi)=\sum_{|\alpha|<N}
 \frac{\partial^\alpha\phi(0)}{\alpha!}
 \frac{\int_{S^{d-1}}k(\theta)\theta^\alpha d\sigma}
      {z+d+|\alpha|}+E_N(z)(\phi),
\]
where the remainder is holomorphic for $\operatorname{Re}z>-d-N$.
Indeed the Taylor remainder is bounded by $C|x|^N$ times a
finite derivative seminorm. The same estimates, and their versions
with powers of $\log|x|$, prove locally uniform continuity on
$\mathcal S$ and holomorphy. Thus this is a meromorphic
distribution family with at most the stated simple poles.
For tests supported off zero its value at $z=\lambda$ is the
integral against $k$. The change of variables proving
$I(z)_a=a^zI(z)$ in the initial half-plane extends meromorphically.
At a regular value it proves homogeneity. At $z=-d-m$, its
Laurent expansion has residue $R_m$ and constant term $F$;
comparing constant terms proves the displayed logarithmic law.
The residue is zero exactly under the stated moment conditions.

It remains to prove necessity and uniqueness, rather than only
produce a preferred extension. The difference of two extensions
is supported at zero and hence is a finite sum of delta derivatives,
by Theorem~\ref{fa2:p4:point-support}. These derivatives have
respective degrees $-d-|\alpha|$. Distinct powers of $a>0$
are linearly independent (put $a=e^t$ and differentiate, or use
a Vandermonde matrix). Consequently a point-supported homogeneous
distribution of degree $\lambda$ is zero except at the listed
critical degrees, where precisely the derivatives of order $m$
are allowed. If a homogeneous extension $U$ existed with
$R_m\ne0$, write $U-F$ as such a finite sum. Differentiate
the dilation identity at $a=1$. The dilation generator minus
$(-d-m)$ kills order-$m$ delta derivatives and is diagonal with
nonzero eigenvalues on all other orders; it therefore cannot
cancel the nonzero order-$m$ vector $R_m$. This contradiction
proves necessity and completes every alternative.
\end{proof}

\begin{proposition}[Fourier degrees and the Riesz-kernel constant]
\label{fa2:p4:homogeneous-fourier}
Every homogeneous distribution is tempered. Fourier transform
changes degree $\lambda$ to degree $-d-\lambda$.
If the original distribution is smooth away from zero, so is its
Fourier transform. In particular, for $0<a<d$,
\[
 \mathcal F(|x|^{-a})(\xi)=
 \pi^{a-d/2}\frac{\Gamma((d-a)/2)}{\Gamma(a/2)}|\xi|^{a-d}.
\]
\end{proposition}
\begin{proof}
A cutoff near zero is a compact distribution and hence tempered.
For the remaining dyadic annuli, rescale each to one fixed annulus.
The finite-order distribution bound there and homogeneity bound
the $j$th annular pairing by a fixed Schwartz seminorm times
$2^{j(\operatorname{Re}\lambda+d-N)}$, after taking a decay
order $N$ larger than the exponent and the finitely many derivative
orders involved. Summing proves a continuous extension to
$\mathcal S$. The Fourier dilation identity, first on tests by
change of variables and then by duality, proves the degree formula.

For smoothness let $u=u_0+u_1$, with $u_0$ supported near zero
and $u_1$ a smooth function vanishing near zero and equal to $k$
at infinity. The transform of $u_0$ is smooth. For every multiindex
$\alpha$, sufficiently many Laplacians make
$\Delta^N(x^\alpha u_1)$ integrable, since its large-radius
degree is $\operatorname{Re}\lambda+|\alpha|-2N<-d$.
Its Fourier transform is continuous. Dividing by
$(-4\pi^2|\xi|^2)^N$ off zero shows that every distributional
Fourier derivative is continuous there, hence the transform is
smooth there. This also proves the converse using inversion.

Finally use
\[
 |x|^{-a}=\frac{\pi^{a/2}}{\Gamma(a/2)}
       \int_0^\infty t^{a/2-1}e^{-\pi t|x|^2}dt,
 \qquad
 \widehat{e^{-\pi t|x|^2}}=t^{-d/2}e^{-\pi|\xi|^2/t}.
\]
Pair with a Schwartz test and interchange using absolute integrals;
the conditions $0<a<d$ give integrability at both ends (or first
use nonnegative Gaussian tests and then dominate the absolute test).
The substitution $v=\pi|\xi|^2/t$ gives exactly the stated
Gamma factor. This verifies both its sign and its $2\pi$
normalization, not just the degree of the answer.
\end{proof}

\begin{example}[Laplacian and causal heat fundamental solutions]
\label{fa2:p4:laplace-heat-fundamental}
Write $\omega_{d-1}=|S^{d-1}|$. Fundamental solutions for
$\Delta$ on $\mathbb R^d$ are
\[
 \Phi_d(x)=\frac{|x|^{2-d}}{(2-d)\omega_{d-1}}\ (d\ge3),
 \quad \Phi_2(x)=\frac{\log|x|}{2\pi},\quad
 \Phi_1(x)=\frac{|x|}{2}.
\]
A fundamental solution of $\partial_t-\Delta_x$ on space-time is
\[
 E(t,x)=1_{t>0}(4\pi t)^{-d/2}e^{-|x|^2/(4t)}.
\]
All formulas are distribution identities including the delta source.
\end{example}
\begin{proof}
The radial Laplacian vanishes off zero in each displayed case.
Integrate $\Phi_d\Delta\phi$ outside a radius-$\varepsilon$
ball twice by parts. Its inner boundary term containing
$\partial_r\Phi_d$ is the spherical average of $\phi$ and
tends to $\phi(0)$; the other term tends to zero, bounded by
$C\varepsilon$ for $d\ge3$ or
$C\varepsilon|\log\varepsilon|$ for $d=2$.
For $d=1$, integrating separately on the two half-lines gives
the jump of $\Phi_1'$ equal to one. Thus $\Delta\Phi_d=\delta$.

For heat, the Gaussian has mass one at each $t>0$, so $E$ is
locally integrable even through $t=0$. Direct differentiation
gives $\partial_tE=\Delta_xE$ for positive time. Pair with a
compact space-time test and integrate on $t>\varepsilon$;
integration by parts leaves
$\int E(\varepsilon,x)\phi(\varepsilon,x)dx$.
Changing variables $x=\sqrt\varepsilon y$ and dominated
convergence make it tend to $\phi(0,0)$.
Consequently $(\partial_t-\Delta)E=\delta_{(0,0)}$.
The earlier convolution theorem now solves the corresponding
equations for compactly supported distribution data; causality
follows from support addition. Neither a formal reciprocal
multiplier nor the positive-time equation alone would prove
the initial delta term.
\end{proof}

\section{Fourier failures and a finite-dimensionality theorem}

\begin{theorem}[Fourier divergence and nonsurjectivity of coefficients]
\label{fa2:p4:fourier-failures}
There is a continuous function on the circle whose symmetric Fourier
partial sums are unbounded at a prescribed point. The injective bounded
coefficient map $L^1(\mathbb T)\to c_0(\mathbb Z)$ is not onto.
There are also $L^1$ functions whose symmetric Fourier partial sums
do not converge in $L^1$. Normalized Haar measure is used here.
\end{theorem}
\begin{proof}
The Dirichlet kernel is
\[
 D_N(t)=\sum_{|n|\le N}e^{int}
       =\frac{\sin((N+1/2)t)}{\sin(t/2)}.
\]
Its $L^1$ norm tends to infinity, in fact is comparable with
$\log(N+2)$. For the upper bound split at $1/N$, use
$|D_N|\le2N+1$ there and $|D_N(t)|\le C/|t|$ for
$1/N\le|t|\le\pi$. For the lower bound use, in each period
of the numerator between $1/N$ and $1$, its subinterval of
fixed relative length where the numerator has modulus at least
$1/2$. The denominator is at most $|t|/2$; summing these
subinterval integrals gives $c\sum_{j\le cN}1/j$.

The functional $f\mapsto S_Nf(0)=\int D_N(-t)f(t)dm(t)$
on $C(\mathbb T)$ has norm exactly $\pfnorm{D_N}_1$:
the upper bound is immediate, and continuous approximations to
its measurable phase, outside arbitrarily small neighborhoods of
the finitely many zeros, give the reverse bound. Uniform boundedness
therefore gives one continuous $f$ for which these scalar values
are unbounded. Translation gives any prescribed point.

The Riemann--Lebesgue assertion follows by approximating an $L^1$
function by trigonometric polynomials (continuous approximation
followed by the positive Fej\'er approximate identities suffices).
The same identities prove uniqueness: all zero coefficients imply
all Fej\'er convolutions zero and their $L^1$ limit is the function.
If the coefficient map were onto, its bounded inverse would give
$\pfnorm f_1\le C\sup_n|\widehat f(n)|$. But $D_N$ has
coefficient supremum one and unbounded $L^1$ norm. This is a
concrete contradiction to surjectivity, using exactly the Banach
space inverse theorem rather than treating coefficient decay as
a sufficient condition. Finally convolution by $D_N$ on $L^1$
has norm $\pfnorm{D_N}_1$: the lower bound follows by applying
it to positive mass-one approximate identities, and the upper
bound is Young's inequality. Uniform boundedness again gives an
$L^1$ function with unbounded partial-sum norms, which precludes
norm convergence.
\end{proof}

\begin{theorem}[Grothendieck's bounded-function subspace theorem]
\label{fa2:p4:grothendieck-lp}
Let $\mu(X)<\infty$, $0<p<\infty$, and let $E$ be a closed
linear subspace of $L^p(\mu)$ whose every element belongs to
$L^\infty(\mu)$. Then $E$ is finite dimensional.
\end{theorem}
\begin{proof}
The zero measure case is immediate. For $p<1$, $L^p$ is an
$F$-space with complete metric $d(f,g)=\int|f-g|^p$.
Indeed take a Cauchy subsequence with summable $p$th-power
increments. Tonelli makes their pointwise $p$th powers summable,
hence their absolute values summable; the pointwise series has
a limit, and Fatou together with
$|a+b|^p\le|a|^p+|b|^p$ gives metric convergence.
Thus its closed subspace $E$ is an $F$-space. The inclusion
$E\to L^\infty$ has closed graph: $L^p$ convergence admits
an almost-everywhere convergent subsequence, and an $L^\infty$
limit agrees with that limit. The $F$-space closed graph theorem
(Theorem~\ref{fa2:p1:fspace-open-graph}), followed by scalar
homogeneity, gives
$\pfnorm f_\infty\le C\pfnorm f_p$.
We claim $\pfnorm f_\infty\le A\pfnorm f_2$ on $E$.
For $p\le2$ this follows from finite-measure H\"older, with
exponent $2/p>1$ when $p<2$, including $p<1$. For
$p>2$, integrate $|f|^p\le\pfnorm f_\infty^{p-2}|f|^2$
and substitute the closed-graph bound; after cancelling a nonzero
$\pfnorm f_p^{p-2}$ it gives
$\pfnorm f_p\le C^{(p-2)/2}\pfnorm f_2$ and the claim.

Choose any $L^2$-orthonormal $f_1,\ldots,f_n\in E$.
For every rational complex coefficient vector $a$,
$|\sum a_jf_j(x)|\le A(\sum|a_j|^2)^{1/2}$ outside
a null set. One common null set suffices for the countable
set of such vectors; continuity in $a$ extends the inequality
to all coefficients there. Taking the Euclidean dual norm gives
$\sum_{j=1}^n|f_j(x)|^2\le A^2$ a.e. Integration yields
$n\le A^2\mu(X)$. Any finite-dimensional subspace of $E$
admits this Gram--Schmidt process in $L^2$, since equality
there is the same almost-everywhere equality. Its dimensions
are therefore uniformly bounded, forcing $E$ itself to be
finite dimensional. Both hypotheses matter: without closedness
the polynomials on an interval give a counterexample; on an
infinite measure space, functions constant on disjoint unit
intervals give a closed $\ell^p\subset L^p\cap L^\infty$.
\end{proof}

\section{Interpolation and classical noncompact integral transforms}

\begin{corollary}[Hausdorff--Young]
\label{fa2:p4:hausdorff-young}
For $1\le p\le2$ and $p'=p/(p-1)$,
$\pfnorm{\widehat f}_{p'}\le\pfnorm f_p$ on $\mathbb R^d$,
with the compatible extended Fourier transform. On the circle,
$\pfnorm{(\widehat f(n))}_{\ell^{p'}}\le\pfnorm f_p$.
\end{corollary}
\begin{proof}
The $L^1\to L^\infty$ bound is one by the integral formula,
and the $L^2$ bound is one by Plancherel, including its circle
orthonormal-series form. Theorem~\ref{fa2:p4:riesz-thorin}
interpolates them. On the dense intersection with $L^1$ the
operators agree, which identifies the completion and proves
compatibility, not merely a bound for smooth inputs.
\end{proof}

\begin{theorem}[Laplace and Hilbert--Hankel transforms]
\label{fa2:p4:laplace-hankel}
On the positive half-line, $Lf(s)=\int_0^\infty e^{-st}f(t)dt$
extends to a bounded self-adjoint operator on $L^2$ with norm
$\sqrt\pi$. Its square is the Hilbert--Hankel (Carleman) operator
\[
 Hf(s)=\int_0^\infty\frac{f(t)}{s+t}dt.
\]
For $1<p<\infty$, $H$ has exact $L^p$ norm
$\pi/\sin(\pi/p)$. It has neither strong endpoint bound.
The Laplace transform $L:L^p\to L^p$ is unbounded for
$p\ne2$, $1\le p\le\infty$.
\end{theorem}
\begin{proof}
Weighted Schur with weight $t^{-1/2}$ gives
$\int_0^\infty e^{-st}t^{-1/2}dt=\sqrt\pi s^{-1/2}$
in both variables, proving $\pfnorm L_2\le\sqrt\pi$.
The symmetric real kernel proves self-adjointness by Fubini on
compactly supported bounded inputs and continuity. Positive Fubini
first, followed by linearity, computes $L^2=H$ on those inputs;
bounded extension gives the identity on $L^2$.

The isometry $U_pf(x)=e^{x/p}f(e^x)$ from $L^p(0,\infty)$
to $L^p(\mathbb R)$ converts $H$ into convolution with
\[
 h_p(x)=\frac{e^{x/p}}{1+e^x},\qquad
 \int_{\mathbb R}h_p(x)dx
 =\int_0^\infty\frac{r^{1/p-1}}{1+r}dr
 =\frac{\pi}{\sin(\pi/p)}.
\]
The last Euler integral follows, for example, by integrating
$z^{1/p-1}/(1+z)$ around a keyhole contour with argument
$0<\arg z<2\pi$: the small and large circular integrals tend
to zero because $0<1/p<1$, and the residue at $-1$ gives
the displayed value. Young gives the upper norm bound.
For the reverse bound use $g_R=1_{[-R,R]}$.
On $[-R+M,R-M]$, $(h_p*g_R)(x)\ge\int_{-M}^Mh_p$.
Divide by $\pfnorm{g_R}_p$, let $R\to\infty$, and then
$M\to\infty$. This proves exactness, including $p=2$.
Thus $\pfnorm{L^2}_2=\pi$; self-adjointness gives
$\pfnorm L_2^2=\pfnorm{L^2}_2$, proving the sharp Laplace norm.
At the endpoints $H1_{[1,2]}$ has a nonintegrable $1/s$ tail,
while $H1_{(0,1)}(s)=\log((s+1)/s)$ is unbounded.
Finally for $f_a(t)=f(at)$,
$Lf_a(s)=a^{-1}Lf(s/a)$. A bounded $L^p\to L^p$ norm
would bound the ratio $a^{2/p-1}\pfnorm{Lf}_p/\pfnorm f_p$
for all $a>0$. Choose a nonzero nonnegative compactly supported
input away from zero; its Laplace image belongs to every finite
$L^p$ and to $L^\infty$. Unless $p=2$, the ratio is unbounded
in one scaling direction. This proves the source exercise as well
as the main $L^2$ theorem.
\end{proof}

\section{Quantitative endpoint consequences}

\begin{theorem}[John--Nirenberg integrability and the BMO endpoint]
\label{fa2:p4:john-nirenberg}
For every BMO function $g$ and $1\le q<\infty$,
\[
 \left(\frac1{|B|}\int_B|g-g_B|^q\right)^{1/q}
 \le C_dq\pfnorm g_{\rm BMO}.
\]
There are dimension constants $c,C>0$ such that
\[
 \frac1{|B|}\int_B
 \exp\left(\frac{c|g-g_B|}{\pfnorm g_{\rm BMO}}\right)\le C
\]
when the denominator is nonzero. Standard convolution CZ operators
extend by transposition to $L^\infty\to\mathrm{BMO}/\mathbb C$;
in particular the Hilbert transform has this endpoint interpretation.
\end{theorem}
\begin{proof}
We retain the quantitative constant in the $p$-atom proof.
The maximal estimate can be taken as
$\pfnorm M_p\le C_dp/(p-1)$: in the weak estimate replace
$f$ by $f1_{|f|>\lambda/2}$, whose removed part has maximal
function at most $\lambda/2$, and integrate
$p\lambda^{p-1}\mu\{Mf>\lambda\}$ by Tonelli.
The resulting factor is $C_d2^pp/(p-1)$ before taking the
$p$th root, and is bounded by the displayed simpler constant.
The telescoping $p$-atom proof therefore has coefficient bound
$C_dp/(p-1)$, including constants uniform for $1<p\le2$.
For mean-zero $f\in L^p(B)$ this says
\[
 \pfnorm f_{H^1}\le C_d\frac p{p-1}
                         |B|^{1-1/p}\pfnorm f_p.
\]
Take $q\ge2$ and $p=q/(q-1)$. The BMO functional hence
has norm at most $C_dq|B|^{1/q}\pfnorm g_{\rm BMO}$
on this mean-zero $L^p$ space. Hahn--Banach and scalar $L^p$
duality represent its extension by an $L^q(B)$ function of that
norm. On bounded mean-zero tests its difference from $g$ is a
constant: testing $1_A-|A|1_B/|B|$ proves this directly.
Thus $g-g_B$ has $L^q$ norm at most twice that bound.
For $1\le q<2$ use the $q=2$ inequality and normalized
H\"older. This proves the asserted estimate for complex functions
as well, by applying the real argument to both components.

Expand the exponential as a nonnegative series. The integer
moment estimate bounds its $n$th term by
$(cC_dn)^n/n!$. Since $n!\ge(n/e)^n$, a sufficiently small
$c$ makes their sum a convergent geometric series, uniformly in
$B$. Zero BMO seminorm means a locally constant representative
and causes no exceptional case in the moment assertion.
Finally the bilinear transpose $T^{\mathsf t}=\mathcal C T^*\mathcal C$,
where $\mathcal C f=\overline f$, is another standard CZ
operator and is bounded $H^1\to L^1$ by the atomic endpoint
theorem. For $h\in L^\infty$ the functional
$f\mapsto\int hT^{\mathsf t}f$ on $H^1$ is bounded. The dual theorem
identifies it with a BMO class of norm at most $C\pfnorm h_\infty$.
On test inputs for which the original operator is also defined,
Fubini or the $L^2$ adjoint identity identifies this class with
that operator, up to the unavoidable constant. This proves a
consistent endpoint extension rather than assigning a divergent
uncorrected integral to arbitrary bounded data.
\end{proof}

\begin{proposition}[An atomic Fourier improvement over $L^1$]
\label{fa2:p4:hardy-fourier}
For $f\in H^1(\mathbb R^d)$,
\[
 \int_{\mathbb R^d}\frac{|\widehat f(\xi)|}{|\xi|^d}d\xi
 \le C_d\pfnorm f_{H^1}.
\]
For $d\ge1$ the cancellation condition on atoms is essential.
\end{proposition}
\begin{proof}
For $d=0$ the atomic space is zero because cancellation forces every
atom on the singleton to vanish; the bound is then immediate.
For $d\ge1$, for an atom on a ball of radius $R$, subtract its center
character to use zero integral and obtain
$|\widehat a(\xi)|\le CR|\xi|$. Its contribution from
$|\xi|\le R^{-1}$ is bounded by
$CR\int_0^{R^{-1}}dr=C$. On the complementary region,
Cauchy--Schwarz and Plancherel bound the integral by
\[
 \pfnorm a_2
 \left(\int_{|\xi|>R^{-1}}|\xi|^{-2d}d\xi\right)^{1/2}
 \le C R^{-d/2}R^{d/2}=C.
\]
Sum the atom coefficients and use Tonelli. The $L^1$ convergence
of atomic partial sums gives uniform convergence of their Fourier
transforms, so Fatou identifies their limit with $\widehat f$
and proves the assertion for the completed space. A function
with nonzero integral has nonzero transform near zero and the
weighted integral diverges, which proves the last sentence.
\end{proof}

\begin{example}[A common lacunary Hilbert subspace in every finite $L^p$]
\label{fa2:p4:lacunary-subspace}
There is a fixed increasing integer sequence $(n_j)$ such that,
for every $1\le p<\infty$, the closed span of
$e^{in_jt}$ in $L^p(\mathbb T)$ is isomorphic to $\ell^2$,
with equivalent coefficient and $L^p$ norms. These infinite-dimensional
closed spaces contain unbounded functions, in agreement with
Theorem~\ref{fa2:p4:grothendieck-lp}.
\end{example}
\begin{proof}
Take $n_j=2^{2^j}$. For every fixed integer $m$ there is $J_m$
such that $n_j>m\sum_{k<j}n_k$ whenever $j\ge J_m$.
Consider a tail polynomial $P=\sum_{j\ge J_m}a_je^{in_jt}$.
In the expansion of $\int|P|^{2m}$ an exponent can vanish only
when the positive and negative multisets of indices agree:
otherwise their largest unequal multiplicity has absolute
contribution at least $n_j$, exceeding the sum of all smaller
contributions. Hence
\[
 \pfnorm P_{2m}^{2m}
 =\sum_{\sum r_j=m}\left(\frac{m!}{\prod r_j!}\right)^2
                                  \prod|a_j|^{2r_j}
 \le m!\left(\sum|a_j|^2\right)^m.
\]
The finitely many head terms have supremum norm bounded by
$\sqrt{J_m}\,(\sum|a_j|^2)^{1/2}$. Combining head and tail
proves $\pfnorm P_p\le C_p\pfnorm a_2$ for every finite
$p$, by choosing $2m\ge p$ and using normalized measure.
For the reverse estimate use the $p=4$ upper bound and
$\pfnorm P_2\le\pfnorm P_1^{1/3}\pfnorm P_4^{2/3}$,
which is H\"older interpolation. Orthonormality gives
$\pfnorm P_2=\pfnorm a_2$, so
$\pfnorm P_1\ge c\pfnorm a_2$; all $p\ge1$ dominate
that norm. Completion therefore identifies each closed span
with precisely $\ell^2$, with both norm bounds proved.

For a concrete unbounded member use $a_j=1/j$. Its $L^2$
Fourier sum belongs to all the finite-$p$ spaces just obtained.
If it were essentially bounded, its Poisson extensions would be
uniformly bounded by that supremum norm. At angle zero they are
$\sum_j r^{n_j}/j$, which increase without bound as $r\uparrow1$
by monotone convergence. This contradiction proves unboundedness.
Thus the example retains the contrasting lacunary geometry in
the source exercise family, with no unproved general lacunary
multiplier assertion.
\end{proof}

\section{Further measure-space consequences from result-bearing problems}

\begin{theorem}[Setwise convergence and Hahn--Vitali--Saks]
\label{fa2:p4:hahn-vitali-saks}
On a $\sigma$-finite measure space, let $f_n\in L^1$ and
suppose $\int_Af_n$ converges to a finite scalar for every
measurable set $A$. Then $(f_n)$ is bounded and uniformly
integrable with finite-measure tail control, and converges
weakly in $L^1$ to some $f\in L^1$. In particular $L^1$ is
weakly sequentially complete. Weak convergence together with
almost-everywhere convergence implies $L^1$ norm convergence.
\end{theorem}
\begin{proof}
Choose $X_j\uparrow X$, $\mu(X_j)<\infty$, and put on the
measure algebra of measurable sets the complete metric
\[
 d(A,B)=\sum_j\frac{2^{-j}}{1+\mu(X_j)}
                                    \mu((A\triangle B)\cap X_j).
\]
Completeness follows by taking a subsequence with summable local
symmetric differences, and patching the almost-everywhere limiting
indicators; the original Cauchy sequence has that same metric limit.
For each $f\in L^1$, $A\mapsto\int_Af$ is continuous in this
metric: discard its small integral outside some $X_j$, and use
absolute continuity of its integral inside $X_j$.
For fixed $\varepsilon>0$ the closed sets
\[
 E_N=\{A:|\int_A(f_n-f_m)|\le\varepsilon
                                      \text{ for all }m,n\ge N\}
\]
cover this complete metric space. Baire gives a ball about $A_0$
contained in one $E_N$. If $d(B,\varnothing)$ is smaller than
the ball radius, both $A_0\setminus B$ and $A_0\cup B$ lie
in that ball. Subtract their integrals from that over $A_0$.
It follows that, for $g=f_n-f_m$ and $n,m\ge N$,
$|\int_Bg|\le4\varepsilon$. The same holds for every subset
of $B$. Splitting the real and imaginary parts into their positive
and negative sets consequently gives
$\int_B|g|\le16\varepsilon$.

All sufficiently small-$\mu$ sets $B$ have such small metric,
and so do all subsets of $X\setminus X_J$ for large $J$,
since the first $J$ summands vanish. Fix $m=N$ and use absolute
continuity and tails of the single $f_N$; the finitely many
earlier functions can be handled as well. This proves both
uniform small-set control and finite-measure tails for $(f_n)$.
To see boundedness, partition a fixed finite-measure $X_J$
into finitely many sets of measure below the small-set threshold
and finitely many large atoms. Such a partition exists: separate
the finitely many atoms of mass above half the threshold and
group the remaining small atoms and nonatomic pieces into sets
of mass between half the threshold and the threshold, with one
remainder. On the small pieces the preceding uniform integral
bound applies. On an atom an integrable function is constant
almost everywhere, so $\int_A|f_n|=|\int_Af_n|$, uniformly
bounded by the setwise convergence hypothesis. Tail control
finishes the uniform norm bound.

Dunford--Pettis now gives a weakly convergent subsequence.
Its limit $f$ has exactly the prescribed setwise integrals.
Uniform boundedness of $\pfnorm{f_n}_1$ permits uniform
simple approximation of every $L^\infty$ test, so the whole
sequence, not just the subsequence, converges weakly to $f$.
For weak sequential completeness, setwise convergence is automatic
for a weak Cauchy sequence. Finally weak convergence gives
uniform integrability and tails by that same criterion.
On a finite-measure part, almost-everywhere convergence plus
uniform integrability gives norm convergence: Egorov leaves a
small-measure exceptional set, controlled by uniform integrability,
and uniform convergence controls its complement. The common tails
then give global norm convergence. This proves the stated Vitali
consequence without silently assuming total measure finite.
\end{proof}

\begin{theorem}[De la Vall\'ee--Poussin's criterion]
\label{fa2:p4:de-la-vallee-poussin}
If $\mu(X)<\infty$, a family $F\subset L^1$ is bounded and
uniformly integrable if and only if there is a continuous increasing
convex $G:[0,\infty)\to[0,\infty)$ with $G(t)/t\to\infty$
and $\sup_{f\in F}\int G(|f|)<\infty$.
\end{theorem}
\begin{proof}
Given such $G$, for every $\varepsilon>0$ choose $R$ such
that $t\le\varepsilon G(t)$ for $t>R$. Then
$\int_A|f|\le R\mu(A)+\varepsilon\int G(|f|)$, proving
boundedness and uniform small-set control. Conversely boundedness
and uniform integrability imply
$\sup_f\int_{|f|>R}|f|\to0$: Markov makes these sets
uniformly small. Choose $R_j\uparrow\infty$ so that this tail
bound is at most $2^{-j}$, and set
$G(t)=\sum_{j\ge1}(t-R_j)_+$. The sum is locally finite,
continuous, increasing and convex, and Tonelli gives
$\sup_f\int G(|f|)\le\sum_j2^{-j}$. For every fixed $N$,
$G(t)/t\ge\sum_{j=1}^N(1-R_j/t)$ for large $t$, proving
superlinear growth. On infinite measure spaces this criterion
alone would not supply the missing spatial tail condition; no
such extension is being asserted.
\end{proof}

\begin{proposition}[Norm-preserving weak-star approximation of measures]
\label{fa2:p4:measure-approximation}
On a compact metric space every finite scalar Radon measure is
a sequential weak-star limit of finitely supported measures of
exactly the same variation norm. Probabilities may be approximated
by finitely supported probabilities. If $K=\overline\Omega$ for
a bounded open Euclidean set, the approximants may instead be
$u_n\,dx$, with $u_n\in C_c^\infty(\Omega)$ and
$\pfnorm{u_n}_1=\pfnorm\mu$; for a positive measure they
can be nonnegative. The same assertion holds for $C_0(\Omega)'$.
\end{proposition}
\begin{proof}
Take finite Borel partitions of mesh tending to zero and put
the mass of each cell at one point of that cell. Uniform continuity
shows weak-star convergence, with variations at most $\pfnorm\mu$.
Conversely for every continuous test $|h|\le1$, the limit of
its pairings is bounded by the liminf of these variations;
taking the supremum and using the measure-dual norm proves that
their variations tend to $\pfnorm\mu$. If $\mu\ne0$,
normalize the sufficiently late nonzero approximants to this
exact norm; the normalization tends to one. For positive measures
the unnormalized cell construction already has the correct mass.
The zero measure requires only zero approximants.

If $\Omega$ is empty use zero. For nonempty $\Omega$ in dimension
zero, the constant density $m$ represents $m\delta_0$ exactly,
including its variation norm and positivity. In positive dimension,
for $\overline\Omega$, replace each of the finitely many
nodes by a nearby distinct interior node, which is possible
because $\Omega$ is dense in its closure. Approximate each
node by a nonnegative integral-one smooth bump in a small
interior ball, choosing the finitely many balls disjoint and
with radii tending to zero. The scalar coefficients then give
a smooth density whose $L^1$ norm is the sum of their moduli,
exactly the atomic variation. Uniform continuity gives the
same weak-star limit; positivity is preserved. For $C_0(\Omega)$,
extend the finite measure by zero to $\overline\Omega$ and
apply this construction, or discard a small boundary tail first.
Finally $L^1(\Omega)\hookrightarrow M(\overline\Omega)$
is isometric by the full variation formula. For $d\ge1$,
$\pfnorm{u\,dx+c\delta_a}=\pfnorm u_1+|c|$ because
Lebesgue measure and the point mass are mutually singular. On the
mass-one singleton the corresponding norm is $|u(0)+c|$.
These bridges retain the actual approximation content of
\cite{B01}, Problem 24, beyond mere measure representation.
\end{proof}

\begin{proposition}[Sharp interval Poincar\'e--Wirtinger constants]
\label{fa2:p4:sharp-interval-poincare}
On $I=(0,1)$, writing $\bar u=\int_Iu$, one has
\[
 \pfnorm{u-\bar u}_\infty\le\pfnorm{u'}_1,\qquad
 \pfnorm u_\infty\le\tfrac12\pfnorm{u'}_1
                         \quad(u\in W^{1,1}_0),
\]
and $\pfnorm{u-\bar u}_1\le\tfrac12\pfnorm{u'}_1$.
All constants are sharp. The first and third are not attained by
nonconstant functions; the second is attained by any nonzero
nonnegative unimodal absolutely continuous tent with zero endpoints.
In $H^1$, $\pfnorm{u-\bar u}_2\le\pi^{-1}\pfnorm{u'}_2$,
with equality precisely for a constant plus a multiple of
$\cos\pi x$. In $H^1_0$, the sharp constant is also $1/\pi$,
with optimizers the multiples of $\sin\pi x$.
For $1<p\le\infty$ and $1\le q\le\infty$, the best
constants in the corresponding mean-zero and zero-trace
$\pfnorm u_q\le C\pfnorm{u'}_p$ inequalities are attained.
\end{proposition}
\begin{proof}
The AC formula gives
\[
 u(x)-\bar u=\int_0^1
       [1_{t<x}-(1-t)]u'(t)dt.
\]
For each $x$, the kernel has modulus strictly less than one
for almost every $t$, so the first bound is strict for a
nonzero derivative. A continuous representative attains its
supremum; approximations to a step of height one concentrated
near the endpoint show sharpness. Integrating the absolute
kernel in $x$ gives $2t(1-t)$, hence
$\pfnorm{u-\bar u}_1\le2\int t(1-t)|u'|\le\pfnorm{u'}_1/2$.
The second inequality is strict for a nonzero derivative, since
$t(1-t)<1/4$ a.e.; a smoothed unit step at $1/2$ approaches
the constant $1/2$. For zero endpoints, the two estimates
$|u(x)|\le\int_0^x|u'|$ and
$|u(x)|\le\int_x^1|u'|$ add to prove the supremum bound.
Unimodal tents give equality.

For the $L^2$ statements reflect evenly or oddly onto the
period-two circle. Parseval for its cosine or sine basis gives
\[
 \pfnorm{u-\bar u}_2^2=\sum_{n\ge1}|a_n|^2,
 \quad\pfnorm{u'}_2^2=\sum_{n\ge1}\pi^2n^2|a_n|^2,
\]
with the same coefficient normalization in both sums.
Smooth density proves the formula for all $H^1$ representatives;
the odd zero-trace extension gives the sine case. The smallest
nonzero frequency is $\pi$, and equality kills every higher
coefficient, proving constants and all optimizers.
For the final assertion normalize $\pfnorm u_q=1$, impose
zero mean or zero trace, and minimize $\pfnorm{u'}_p$.
The preceding bounds make a minimizing sequence bounded and
equicontinuous by $|u(x)-u(y)|\le C|x-y|^{1-1/p}$.
Ascoli gives a uniform limit preserving the constraint. For
finite $p>1$, weak compactness of the derivatives identifies
the derivative of the limit and lowers its norm. For $p=\infty$
the uniform Lipschitz bound passes directly to the limit.
Thus the minimum is attained in both cases, including $q=\infty$.
\end{proof}


% END INPUT


% BEGIN parts/part4-hardy-converse.tex
\chapter{A fixed smooth maximal function and atomic Hardy space}
\label{fa2:p4:hardy-converse-chapter}

We now prove the converse to the atomic maximal estimate. A uniform
estimate for each test function separately would not estimate their
supremum, so we introduce and control that supremum explicitly. This
also supplies the fixed-kernel characterization asked for in
\cite[Chapter 2, Section 8, Problem 6]{B09}, including its slower-decay variant.
Throughout $d\ge1$ and $\Phi_t(x)=t^{-d}\Phi(x/t)$.

\begin{theorem}[Fixed radial maximal characterization]
\label{fa2:p4:hardy-fixed-maximal}
Suppose $\Phi\in C^\infty(\mathbb R^d)$,
$c=\int\Phi\ne0$, and for every multi-index $\alpha$,
\[
 |\partial^\alpha\Phi(x)|\le C_\alpha
                  (1+|x|)^{-d-1-|\alpha|}.
\]
For $f\in L^1(\mathbb R^d)$ put
$F_\Phi f(x)=\sup_{t>0}|\Phi_t*f(x)|$.
Then $f\in H^1_{\rm at}$ if and only if $F_\Phi f\in L^1$, and
\[
 C^{-1}\|f\|_{H^1_{\rm at}}\le\|F_\Phi f\|_1
                  \le C\|f\|_{H^1_{\rm at}},
\]
with constants depending on $\Phi,d$.
In particular this holds for a Schwartz kernel of nonzero integral,
for the Gaussian heat kernel, and for the Poisson kernel
$C_d(1+|x|^2)^{-(d+1)/2}$ of integral one.
If $\Phi$ is Schwartz, the hypothesis $f\in L^1$ can be replaced
by $f\in\mathcal S'$ together with $F_\Phi f\in L^1$;
that latter condition itself recovers an $L^1$ representative.
\end{theorem}

\section{Elementary kernel estimates and a low-frequency inverse}

Fix $\epsilon=1/2$ and write
$\|u\|_w=\sup_x(1+|x|)^{d+\epsilon}|u(x)|$.
For integrable $u,v$ with finite weighted norms,
\begin{equation}\label{fa2:p4:hardy-weighted-convolution}
 \|u*v\|_w\le C_w(\|u\|_w\|v\|_1+\|u\|_1\|v\|_w).
\end{equation}
Indeed split the convolution integral into $|y|\ge|x|/2$ and
$|x-y|\ge|x|/2$; on each part put the corresponding factor in
its weighted supremum and integrate the other. This proves the
inequality with, for example, $C_w=2^{d+\epsilon}$.

\begin{lemma}[An integrable inverse on a small frequency ball]
\label{fa2:p4:hardy-lowpass-inverse}
There is a Schwartz function $\eta$ with $\widehat\eta=1$ near
zero and compactly supported smooth Fourier transform, and a
kernel $\Psi$ with
\[
 \eta=\Psi*\Phi,\qquad
 |\Psi(x)|\le C(1+|x|)^{-d-1/2}.
\]
The Fourier cutoff can be chosen with
$\operatorname{supp}\widehat\eta\subset B(0,2/R)$ and
$\widehat\eta=1$ on $B(0,1/R)$ for some fixed $R$.
\end{lemma}
\begin{proof}
Choose smooth Fourier cutoffs $\chi_0,\chi_1$ with
$\chi_0=1$ on $B(0,1)$, supported in $B(0,2)$, and
$\chi_1=1$ on $B(0,2)$, supported in $B(0,4)$.
Let $\eta$ and $\zeta$ have Fourier transforms
$\chi_0(R\xi)$ and $\chi_1(R\xi)$. Put
$h=\zeta-c^{-1}\Phi*\zeta$. Its $L^1$ norm tends to zero as
$R\to\infty$: since $\int\Phi=c$,
\[
 \|h\|_1\le |c|^{-1}\int|\Phi(y)|
       \|\zeta_1(\cdot)-\zeta_1(\cdot-y/R)\|_1\,dy\longrightarrow0,
\]
where $\zeta_1$ is the unscaled Schwartz cutoff. Translation
continuity and the majorant $2\|\zeta_1\|_1|\Phi(y)|$ justify
the limit. For each fixed $R$, $h$ has finite weighted norm
by \eqref{fa2:p4:hardy-weighted-convolution}, since $\Phi$
has the stated decay.
Choose $R$ so $a=\|h\|_1<(2C_w)^{-1}$. Induction in that
inequality gives
\[
 \|h^{*n}\|_1\le a^n,\qquad
 \|h^{*n}\|_w\le nC_w^{n-1}a^{n-1}\|h\|_w\quad(n\ge1).
\]
Thus $\Psi=c^{-1}(\eta+\sum_{n\ge1}h^{*n}*\eta)$
converges both in $L^1$ and in the weighted supremum norm.
Its Fourier transform is
$\widehat\eta/[c(1-\widehat h)]$. On the support of
$\widehat\eta$, $\widehat\zeta=1$ and
$1-\widehat h=\widehat\Phi/c$; off that support the numerator
is zero. Hence $\widehat\Psi\widehat\Phi=\widehat\eta$.
Injectivity of the Fourier transform on $L^1$ proves the
convolution identity. This construction does not assume that
$\widehat\Phi$ is smooth at zero, which would be false for
the Poisson example.
\end{proof}

We shall use a shifted averaging estimate. If
$|K(x)|\le C(1+|x|)^{-d-\tau}$, $\tau>0$, $H\ge0$ is locally
integrable and $K_s=s^{-d}K(\cdot/s)$, then
\begin{equation}\label{fa2:p4:hardy-shifted-average}
 \int |K_s(y-z)|H(z)\,dz
 \le C'(1+|x-y|/s)^d MH(x).
\end{equation}
To prove it, split about $y$ into the ball of radius $s$ and
annuli of radii $2^ls$. Each integral of $H$ is at most
$C(s+|x-y|+2^ls)^d MH(x)$. Multiplication by the kernel
bound and summation gives
$C MH(x)\sum_{l\ge0}2^{-l(d+\tau)}(1+|x-y|/s+2^l)^d$,
which is bounded by the asserted expression.

\section{From a radial maximal function to a grand maximal function}

Set $A=d+1$, $r=1-[4(d+1)]^{-1}$, so
$Ar>d$ and $A(1-r)=1/4<\epsilon$. Choose an integer
$N>2A+d+4$. Let $\mathcal B_N$ be the smooth functions supported
in the unit ball with $\|\partial^\alpha\psi\|_\infty\le1$
for $|\alpha|\le N$, and define
\[
 \mathcal G_N f(x)=\sup_{t>0,\,\psi\in\mathcal B_N}|\psi_t*f(x)|,
 \qquad
 P(x)=\sup_{j\in\mathbb Z,\,y\in\mathbb R^d}
       \frac{|\eta_{2^{-j}}*f(y)|}{(1+2^j|x-y|)^A}.
\]
Both suprema are measurable: each fixed convolution is continuous
in $x$, and a supremum of continuous functions is lower
semicontinuous. In the definition of $P$ one can also restrict
$y$ to a countable dense set.

\begin{lemma}[The compact-test transfer estimate]
\label{fa2:p4:hardy-compact-transfer}
For $f\in L^1$ and $P(x)<\infty$,
\[
 |\psi_t*f(z)|\le C P(x)(1+|x-z|/t)^A
 \quad(\psi\in\mathcal B_N,t>0).
\]
The same bound holds for any fixed uniformly bounded family of
compactly supported smooth tests with bounded $C^N$ norm and
one common support radius. In particular $\mathcal G_Nf\le CP$.
\end{lemma}
\begin{proof}
Write $\eta_j=\eta_{2^{-j}}$ and
$\delta_l=\eta_l-\eta_{l-1}$. The Fourier support choices imply
$\delta_l*f=\delta_l*(\eta_{l+1}*f)$.
All moments of $\delta_l$ vanish, because its Fourier transform
vanishes near zero. Approximate identity convergence gives
\[
 f=\eta_j*f+\sum_{l>j}\delta_l*(\eta_{l+1}*f)
\]
in $L^1$, with the right side understood as its telescoping
partial sums. Choose $j$ with $2^{-j}\le t<2^{-j+1}$.
For $l=j+k$, $k\ge1$, Taylor subtraction through degree $N-1$
in the convolution with $\delta_l$ gives the scaled estimate
\[
 |(\psi_t*\delta_l)(w)|
 \le C_M2^{jd}2^{-kN}(1+2^j|w|)^{-M}
\]
for any fixed $M$. Here is why it includes the spatial tail:
after scaling $2^jt$ into $[1,2]$, on a fixed ball use the
$N$th derivative remainder and
$\int|\delta_l(v)||v|^N\,dv=O(2^{-kN})$.
Outside a larger fixed ball the compact support of $\psi$
forces its convolution variable to be comparable to $w$;
arbitrary Schwartz decay of $\delta_l$ gives the same
$2^{-kN}$ with any prescribed additional power of $|w|$.
Consequently, taking $M>A+d$,
\[
 \int|(\psi_t*\delta_l)(w)|
                (1+2^{l+1}|w|)^A\,dw\le C2^{-k(N-A)}.
\]
The Peetre bound on $\eta_{l+1}*f(z-w)$ introduces in addition
$(1+2^{l+1}|x-z|)^A\le
2^{(k+1)A}(1+2^j|x-z|)^A$.
Thus the $l$th term is at most
$CP(x)(1+2^j|x-z|)^A2^{-k(N-2A)}$.
The low term $\psi_t*(\eta_j*f)$ has the same bound without
the geometric factor, by its compact-support weighted integral.
Sum the geometric series. Finally convolution of the telescoping
partial sums with fixed bounded $\psi_t$ converges uniformly
because those partial sums tend to $f$ in $L^1$.
This proves the estimate with no appeal to a grand-maximal theorem.
\end{proof}

The derivative decay of $\Phi$ yields the useful expansion
\begin{equation}\label{fa2:p4:hardy-tail-bumps}
 \Phi=\psi_0+\sum_{m\ge1}2^{-m}(\psi_m)_{2^m},
\end{equation}
where the $\psi_m$ have common compact support and uniformly
bounded $C^N$ norms. To construct it take a smooth annular
partition at radii $2^m$ and put
$\psi_m(y)=2^{m(d+1)}\rho_m(2^my)\Phi(2^my)$.
The displayed derivative assumptions bound every derivative
through $N$ uniformly; the inner piece gives $\psi_0$.
The series converges in $L^1$ and uniformly. The transfer estimate
at the coarser scales $2^m2^{-j}$ therefore gives
\begin{equation}\label{fa2:p4:hardy-radial-transfer}
 |\Phi_{2^{-j}}*f(z)|\le CP(x)(1+2^j|x-z|)^A.
\end{equation}
The geometric coefficient $2^{-m}$ is important: no nonexistent
high moment of the Poisson kernel is being used here.

\begin{lemma}[The nonlinear maximal estimate]
\label{fa2:p4:hardy-peetre}
If $F=F_\Phi f\in L^1$, then
\[
 P(x)\le C\{M(F^r)(x)\}^{1/r}\quad\hbox{a.e.},\qquad
 \|\mathcal G_Nf\|_1\le C\|F\|_1.
\]
\end{lemma}
\begin{proof}
The low-frequency inverse says
$\eta_j*f=\Psi_{2^{-j}}*(\Phi_{2^{-j}}*f)$.
Since the second factor is pointwise bounded by $F$, the shifted
averaging estimate with $H=F$ gives
$P(x)\le C MF(x)$, using $A>d$. Thus $P$ is finite almost
everywhere by the maximal weak $(1,1)$ bound. This justifies
the division in the next step; an infinite supremum is not
silently divided by itself.

Put $U_j=\Phi_{2^{-j}}*f$. Its bound by $F$ and
\eqref{fa2:p4:hardy-radial-transfer} give
\[
 |U_j(z)|\le C P(x)^{1-r}(1+2^j|x-z|)^{A(1-r)}F(z)^r.
\]
Insert this into $\eta_j*f(y)=\Psi_{2^{-j}}*U_j(y)$ and divide
by $(1+2^j|x-y|)^A$.
Using
$1+2^j|x-z|\le(1+2^j|x-y|)(1+2^j|y-z|)$,
the remaining kernel has decay exponent
$d+\epsilon-A(1-r)>d$. The shifted averaging estimate therefore
bounds the result by
\[
 C P(x)^{1-r}(1+2^j|x-y|)^{d-Ar}M(F^r)(x)
 \le C P(x)^{1-r}M(F^r)(x).
\]
Take the supremum over $j,y$ and divide by $P(x)^{1-r}$ when
$P(x)>0$; the zero case is immediate. This proves the pointwise
inequality. Since $1/r>1$, the strong maximal bound gives
\[
 \int P\le C\|M(F^r)\|_{1/r}^{1/r}
          \le C\|F^r\|_{1/r}^{1/r}=C\|F\|_1.
\]
The compact-test transfer then proves the grand-maximal estimate.
\end{proof}

\section{A full Whitney atomic decomposition}

\begin{lemma}[Smooth Whitney partitions, including nested levels]
\label{fa2:p4:hardy-whitney}
If $\Omega$ is a nonempty open set of finite measure, there are
dyadic cubes $Q_i$ with disjoint interiors covering $\Omega$,
and smooth $\eta_i\ge0$, such that
$\sum_i\eta_i=1$ on $\Omega$, each $\eta_i$ is supported in
$Q_i^*=(9/8)Q_i\Subset\Omega$,
$\int\eta_i\ge c|Q_i|$, and
$|\partial^\alpha\eta_i|\le C_\alpha\ell_i^{-|\alpha|}$.
The enlarged cubes have bounded overlap and neighboring sides
are comparable. If $\Omega'\subset\Omega$ has its analogous
cubes $Q_j'$, then intersection of $Q_j'^*$ and $Q_i^*$ implies
$\ell_j'\le C\ell_i$ and $Q_j'^*\subset C Q_i$.
Each fixed $Q_j'^*$ meets at most a fixed number of the $Q_i^*$.
\end{lemma}
\begin{proof}
Use the supremum-norm distance for this cubical geometry; its
conversion to Euclidean balls changes only dimension constants.
Take maximal dyadic cubes whose twentyfold enlargements lie in
$\Omega$. Every point belongs to a small admissible cube, and
finite measure prevents an infinite admissible ancestor chain.
Maximal dyadic cubes have disjoint interiors and cover the set,
with boundary ambiguities irrelevant and then covered by closures.
Their side lengths are comparable to their distances from
$\Omega^c$: the twentyfold inclusion gives the lower bound,
and failure for the parent gives the upper bound.
The distance function is one-Lipschitz, so overlapping fixed
small enlargements have comparable sides. Disjoint interiors
then imply bounded overlap by a volume comparison.
Choose bumps equal to one on $Q_i$, supported in $Q_i^*$,
with the scaled derivative bounds, and divide them by their
locally finite sum. That sum lies between one and a fixed
constant, proving all partition estimates, including the
lower bound for its integral.
For the nested assertion, distance to $\Omega'^c$ is no greater
than distance to $\Omega^c$. At a point of intersection this
gives $\ell_j'\le C\ell_i$, and elementary diameter comparison
gives containment in $C Q_i$. Across a fine cube the latter
distance varies by at most its diameter and is bounded below
by a fixed large multiple of its side (from the twentyfold
inclusion). All coarse cubes it meets consequently have
comparable sides. Disjoint interiors and a volume comparison
bound their number. These arguments also show local finiteness
inside each open set.
\end{proof}

\begin{theorem}[Grand maximal function gives atoms]
\label{fa2:p4:hardy-grand-atomic}
If $f\in L^1$ and $G=\mathcal G_N f\in L^1$, then
$f\in H^1_{\rm at}$ and $\|f\|_{H^1_{\rm at}}\le C\|G\|_1$.
\end{theorem}
\begin{proof}
A fixed nonnegative test in $\mathcal B_N$ with nonzero integral
is an approximate identity up to that integral, so
$|f(x)|\le CG(x)$ almost everywhere. For each integer $k$ put
$\Omega_k=\{G>2^k\}$. It is open and of finite measure; apply
the preceding lemma, using empty sums when a level is empty
(so $g_k=f$ there). Write its partition
as $\eta_i^k$ and its cubes as $Q_i^k$.
Set
\[
 c_i^k=\frac{\int f\eta_i^k}{\int\eta_i^k},\quad
 b_i^k=(f-c_i^k)\eta_i^k,\quad
 g_k=f-\sum_i b_i^k.
\]
Choose a point of $\Omega_k^c$ at distance at most $C\ell_i^k$
from $Q_i^k$, available from Whitney's upper distance bound.
The rescaled function $\eta_i^k$ is a fixed constant times an
admissible grand test centered at that point, at scale
$C\ell_i^k$. Thus $|\int f\eta_i^k|\le C2^k|Q_i^k|$ and
$|c_i^k|\le C2^k$. Consequently $\|g_k\|_\infty\le C2^k$:
outside $\Omega_k$ use $|f|\le CG$, and inside use
$g_k=\sum_i c_i^k\eta_i^k$.
All these sums converge in $L^1$, since
$\sum_i\|b_i^k\|_1\le\int_{\Omega_k}|f|+C2^k|\Omega_k|$.
Every $b_i^k$ has integral zero. Further,
$g_k\to f$ in $L^1$ as $k\to\infty$, because both terms in
that last bound tend to zero. Also $g_k\to0$ in $\mathcal S'$
as $k\to-\infty$, by its uniform bound times the $L^1$ norm
of a Schwartz test.

Localizing $g_{k+1}-g_k$ requires a cancellation correction;
arbitrary pieces of a mean-zero function need not be atoms.
For a fixed level abbreviate its quantities by $\eta_i,b_i$
and the next level by $\eta_j',b_j'$. Put
\[
 d_{ij}=\frac{\int b_j'\eta_i}{\int\eta_j'},\qquad
 h_i=b_i-\sum_j(b_j'\eta_i-d_{ij}\eta_j'),
\]
using only intersecting supports. Each bracket has integral
zero. For a fixed $j$, only a bounded number of $i$ occur,
and $\sum_i d_{ij}=0$, since $\sum_i\eta_i=1$ on the
support of $b_j'$ and $\int b_j'=0$. Therefore
\[
 \sum_i h_i=\sum_i b_i-\sum_j b_j'=g_{k+1}-g_k.
\]
The support of $h_i$ lies in a fixed dilation of $Q_i$ by the
nested Whitney geometry.

We verify its uniform bound. The cancellation of the $f$ terms
gives
\[
 b_i-\sum_jb_j'\eta_i
 =\eta_i\left[f1_{\Omega_k\setminus\Omega_{k+1}}
                      -c_i+\sum_jc_j'\eta_j'\right].
\]
Here $|f|\le C2^{k+1}$ on the displayed level strip,
$|c_i|\le C2^k$, and $|c_j'|\le C2^{k+1}$, so this part
is bounded by $C2^k$. To estimate $d_{ij}$, the product
$\eta_j'\eta_i$ is an admissible test up to a fixed constant
at the finer scale: its scaled derivatives through $N$ are
bounded because $\ell_j'\le C\ell_i$.
Use the outside point for $Q_j'$ at level $k+1$ to obtain
$|\int f\eta_j'\eta_i|\le C2^{k+1}|Q_j'|$.
The term $c_j'\int\eta_j'\eta_i$ has the same bound, and
$\int\eta_j'\ge c|Q_j'|$. Hence $|d_{ij}|\le C2^k$.
The fine partition has bounded overlap, so
$\|\sum_jd_{ij}\eta_j'\|_\infty\le C2^k$ as well.
Thus $\|h_i\|_\infty\le C2^k$, and $\int h_i=0$.
All rearrangements are legitimate in $L^1$: the fine cubes
meeting a fixed coarse cube lie in its fixed dilation and
have disjoint interiors, the above constants control their
integral sums, and each fine cube meets only boundedly many
coarse ones.

Enclose $C Q_i$ in a ball $B_i$ of comparable volume.
The function $h_i$ is a multiple of an atom with coefficient
at most $C2^k|Q_i|$. Summing over levels gives
\[
 \sum_{k,i}|\lambda_{ki}|
 \le C\sum_{k\in\mathbb Z}2^k|\Omega_k|
 \le2C\int G.
\]
The resulting atom series converges absolutely in $L^1$.
Its partial level sums telescope to $g_{n+1}-g_{-m}$,
which tend to $f$ in distributions by the limits already
proved. Thus its $L^1$ sum is exactly $f$, proving the
theorem. The two-sided level argument also proves, rather
than assumes, the vanishing integral of $f$.
\end{proof}

\section{Completion of the fixed-kernel theorem}

\begin{proof}[Proof of Theorem~\ref{fa2:p4:hardy-fixed-maximal}]
If $F_\Phi f\in L^1$, the nonlinear maximal estimate and the
Whitney decomposition give
$\|f\|_{H^1_{\rm at}}\le C\|\mathcal G_N f\|_1
\le C\|F_\Phi f\|_1$.
For the other direction first prove
$\|\mathcal G_Na\|_1\le C$ uniformly for every atom.
Near its ball, all compact normalized tests satisfy
$\mathcal G_Na\le C M|a|$; use the $L^2$ maximal bound,
$\|a\|_2\le|B|^{-1/2}$ and Cauchy--Schwarz on a fixed
enlargement. Away from the ball, subtract the test value
at its center. Its gradient bound and support condition
give $\mathcal G_Na(x)\le C R|x-x_B|^{-d-1}$, where
$R$ is the ball radius; the tail integral is bounded.
Summing absolute atomic coefficients yields
$\|\mathcal G_N f\|_1\le C\|f\|_{H^1_{\rm at}}$.
This estimates the supremum over all tests, not just each
individual test separately.
The annular expansion \eqref{fa2:p4:hardy-tail-bumps}
gives $F_\Phi f\le C\mathcal G_N f$, because all its
compact pieces have uniform $C^N$ bounds and summable
coefficients. This proves the reverse norm inequality
also for the polynomially decaying kernels.

Finally let $f\in\mathcal S'$ and $\Phi$ be Schwartz, with
$F_\Phi f\in L^1$. The functions $\Phi_t*f$ are dominated
pointwise by this one integrable function. Their family
is bounded, uniformly integrable and has uniform finite-
measure tail control, so the scalar $L^1$ compactness
theorem gives, along a sequence $t\downarrow0$, a weak
$L^1$ limit $u$. In distributions the same sequence tends
to $cf$: for every Schwartz test, the rescaled Schwartz
kernel is an approximate identity and converges to its
integral times that test in all Schwartz seminorms.
Thus $f=u/c$ is an $L^1$ distribution, and the case
already proved applies. For the slower-decay kernels we
assert the original $L^1$ formulation; convolution of an
arbitrary tempered distribution with a non-Schwartz
Poisson kernel has not been silently defined.
All scales $t>0$ and all integer levels were used, so
no low-frequency polynomial quotient or hidden zero-mean
assumption has entered the theorem. The zero function
satisfies all statements directly.
\end{proof}

% END INPUT


% BEGIN parts/part4-final-source-bridges.tex
\chapter{Final distribution and singular-integral source bridges}
\label{fa2:p4:final-source-bridges}
The following statements close the basic variants in
\cite[Chapter 3, Exercises 19,25,26,33 and Problems 1,2;
Chapter 4, Problem 3]{B09}. Distributions and their Fourier transforms
retain the scalar-linear convention of Part IV.

\begin{proposition}[A primitive of every distribution on the line]
\label{fa2:p4:distribution-primitive}
Every $u\in\mathcal D'(\mathbb R)$ has a primitive
$v\in\mathcal D'(\mathbb R)$, and any two primitives differ by a
constant distribution. The same assertion holds on an open interval.
\end{proposition}
\begin{proof}
Choose $\rho\in\mathcal D(\mathbb R)$ with $\int\rho=1$, and put
\[
 A\phi(x)=\int_{-\infty}^x
       \left(\phi(t)-\rho(t)\int\phi\right)dt,
 \qquad v(\phi)=-u(A\phi).
\]
The integrand has integral zero. For tests supported in a fixed compact
set $K$, all the $A\phi$ are supported in one fixed compact interval
containing $K\cup\operatorname{supp}\rho$. Integration bounds their
zeroth seminorm by a constant times $\|\phi\|_\infty$; differentiation
gives analogous bounds for every higher seminorm. Thus $A$ is
continuous on each test-function step and $v$ is a distribution.
For $\psi\in\mathcal D$, $\int\psi'=0$ and $A\psi'=\psi$, so
$v'(\psi)=-v(\psi')=u(\psi)$. If $w'=0$, every zero-integral
test is the derivative of a compactly supported smooth function,
and hence is annihilated by $w$. Consequently
$w(\phi)=w(\rho)\int\phi$, which proves uniqueness modulo constants.
On an interval choose $\rho$ inside the interval and integrate from
any point to the left of the supports. The same compact-support
argument works on every test-function step, including unbounded
intervals; no endpoint value is prescribed.
\end{proof}

\begin{theorem}[Positive distributions are locally finite measures]
\label{fa2:p4:positive-distribution-measure}
Let $\Omega\subset\mathbb R^d$ be open. A scalar-linear distribution
$u\in\mathcal D'(\Omega)$ is nonnegative on nonnegative real tests
if and only if it is integration against a unique positive locally
finite Radon measure on $\Omega$. No finite total-mass assertion is
implicit.
\end{theorem}
\begin{proof}
For compact $K\subset\Omega$ choose a nonnegative
$\chi\in\mathcal D(\Omega)$ equal to one on a neighborhood of $K$.
If a real test $\phi$ is supported in $K$, positivity applied to
$\|\phi\|_\infty\chi\pm\phi$ gives
\[
 |u(\phi)|\le u(\chi)\|\phi\|_\infty.
\]
The values on real tests are real: any such test is the difference
of two nonnegative tests, by the same cutoff. For complex tests
apply the real bound to real and imaginary parts. Thus $u$ has
order zero on every compact set.

For $f\in C_c(\Omega)$ approximate $f$ uniformly by smooth functions
with supports in one fixed compact neighborhood of its support,
using zero extension and convolution. The preceding estimate makes
their $u$-values Cauchy and makes the limit independent of the
approximation or neighborhood. This defines a linear functional
$I$ on $C_c(\Omega)$. If $f\ge0$, use nonnegative mollifiers to
preserve positivity, and obtain $I(f)\ge0$.
The positive $C_c$ version of
Theorem~\ref{fa2:p4:riesz-markov} therefore gives a compact-finite
Radon measure $\mu$ representing $I$ and hence $u$.
Conversely such a measure defines a positive distribution, since
$|\int\phi\,d\mu|\le\mu(K)\|\phi\|_\infty$ on each compact step.
Two representing measures agree on smooth tests and hence, by the
same uniform approximations, on $C_c$; Riesz--Markov uniqueness
then proves uniqueness.
\end{proof}

\begin{proposition}[Convexity and the distribution Hessian]
\label{fa2:p4:convex-distribution-hessian}
Let $\Omega\subset\mathbb R^d$ be open and convex and let
$f:\Omega\to\mathbb R$ be continuous. It is convex if and only if
for every $a\in\mathbb R^d$ the distribution
$\partial_a^2f=\sum_{i,j}a_i a_j\partial_i\partial_j f$ is positive.
In one dimension this says exactly that $f''$ is a positive
locally finite measure. The matrix of distribution second
derivatives thus has the precise positive-semidefinite meaning
given by all of these scalar tests.
\end{proposition}
\begin{proof}
Take a nonnegative smooth mass-one mollifier $\rho_\varepsilon$
supported in $B(0,\varepsilon)$. It is defined on
$\Omega_\varepsilon=\{x:x+\overline B(0,\varepsilon)\subset\Omega\}$;
this set is convex because $\Omega$ is convex.
If $f$ is convex, averaging its defining segment inequalities shows
that $f_\varepsilon=f*\rho_\varepsilon$ is convex there. For a smooth
function, restriction to each line shows that convexity is equivalent
to $\partial_a^2f_\varepsilon\ge0$ for all $a$.
As $f_\varepsilon\to f$ uniformly on compact subsets, its distribution
derivatives converge to those of $f$. Their nonnegative pairings
with nonnegative tests pass to the limit.

Conversely positive $\partial_a^2 f$ convolved with
$\rho_\varepsilon$ is pointwise nonnegative, since its value is
the pairing with a nonnegative translated test. Hence the smooth
Hessian of $f_\varepsilon$ is positive semidefinite. Twice
integrating its nonnegative second derivative on a segment proves
the convex segment inequality. Each fixed segment in $\Omega$
has a compact neighborhood inside $\Omega$, so the inequality
holds for all sufficiently small $\varepsilon$ and passes to
$f$ by local uniform convergence. The measure assertion follows
from Theorem~\ref{fa2:p4:positive-distribution-measure}.
\end{proof}

\begin{theorem}[Periodic distributions and polynomial Fourier growth]
\label{fa2:p4:periodic-distributions}
Write $\mathbb T^d=\mathbb R^d/\mathbb Z^d$ with normalized Haar
measure. Distributions on $\mathbb R^d$ invariant under every
integer translation are canonically the continuous linear
functionals on $C^\infty(\mathbb T^d)$. Their Fourier coefficients
\[
 c_n=U(e^{-2\pi i n\cdot x}),\qquad n\in\mathbb Z^d,
\]
have polynomial growth. Conversely every sequence with
$|c_n|\le C(1+|n|)^N$ defines exactly one such distribution, by
\[
 U(\phi)=\sum_{n\in\mathbb Z^d}c_n\widehat\phi(-n).
\]
The series converges absolutely on every smooth test, and its
finite Fourier sums converge in the distribution sense.
\end{theorem}
\begin{proof}
Periodization $P\psi(x)=\sum_{k\in\mathbb Z^d}\psi(x+k)$ is a
continuous map from each compact test-function step into
$C^\infty(\mathbb T^d)$. Choose a nonnegative compactly supported
smooth $\eta$ positive on $[0,1]^d$, and set
$\chi=\eta/(P\eta)$, with the positive periodic denominator read
on $\mathbb R^d$. Then $\chi\in\mathcal D$ and $P\chi=1$.
Given periodic $u$, define $U(\phi)=u(\chi\phi)$, where $\phi$
is lifted periodically. Continuity follows from the fixed compact
support of $\chi$. For any $\psi\in\mathcal D$, periodicity gives
\[
 U(P\psi)
 =\sum_k u\bigl(\chi(x)\psi(x+k)\bigr)
 =\sum_k u\bigl(\chi(x-k)\psi(x)\bigr)=u(\psi).
\]
Only finitely many terms meet the relevant compact supports.
Conversely $u(\psi)=U(P\psi)$ defines a periodic distribution.
For this distribution $u(\chi\phi)=U(\phi P\chi)=U(\phi)$.
The two constructions are inverse and in particular independent
of the chosen cutoff.

Continuity on the smooth Fr\'echet space gives
$|U(\phi)|\le C\max_{|\alpha|\le N}\|\partial^\alpha\phi\|_\infty$
for some finite $N$. Applying this to exponentials proves polynomial
growth. Integration by parts in the torus Fourier integral shows,
for every $L$, that
\[
 |\widehat\phi(n)|\le C_{L,d}(1+|n|)^{-L}
       \max_{|\alpha|\le L}\|\partial^\alpha\phi\|_\infty.
\]
One can select a largest coordinate of $n$ for each nonzero $n$
to obtain this estimate using $L$ derivatives in that coordinate.
It also implies that the Fourier series of a smooth function
converges with every derivative, by taking $L$ arbitrarily large.
Applying $U$ to that series gives the displayed formula.
For a polynomial-growth sequence choose $L>N+d$ in the same
estimate. Absolute summability gives a continuous functional and
distribution convergence of the finite sums. Testing on each
exponential recovers $c_n$, proving existence and uniqueness.
\end{proof}

\begin{lemma}[An elementary summable molecule decomposition]
\label{fa2:p4:hardy-decaying-molecule}
Suppose $g\in L^2(\mathbb R^d)$ has integral zero,
$\|g\|_2\le A$, and
$|g(x)|\le A|x|^{-d-1}$ for $|x|\ge2$. Then
$g\in H^1_{\rm at}$ and $\|g\|_{H^1}\le C_d A$.
The translated and dilated version has the same bound with the
correspondingly normalized $L^2$ size and tail.
\end{lemma}
\begin{proof}
Let $B_j=B(0,2^{j+1})$, $E_0=B_0$, and
$E_j=B_j\setminus B_{j-1}$ for $j\ge1$. Put
$g_j=g1_{E_j}$, $m_j=\int g_j$, and
$\phi_j=1_{B_j}/|B_j|$. The assumed bounds give
\[
 \|g_j\|_2\le C_d A\,2^{-j}|B_j|^{-1/2},\qquad
 |m_j|\le C_d A\,2^{-j}\quad(j\ge0).
\]
For $j=0$ these follow from Cauchy--Schwarz; for $j\ge1$
they follow by integrating the pointwise tail. Thus
$b_j=g_j-m_j\phi_j$ has zero integral, support in $B_j$,
and $L^2$ size at most $C_d A2^{-j}|B_j|^{-1/2}$.
Theorem~\ref{fa2:p4:hardy-p-atoms} gives
$\sum_j\|b_j\|_{H^1}\le C_d A$.

Set $M_j=\sum_{k\ge j}m_k$. Absolute integrability follows
from the same estimates, and zero total integral says $M_0=0$.
Also $|M_j|\le C_d A2^{-j}$. Absolute $L^1$ convergence and
telescoping give
\[
 \sum_{j\ge0}m_j\phi_j
       =\sum_{j\ge1}M_j(\phi_j-\phi_{j-1}).
\]
Each difference has zero integral, support in $B_j$, and
$L^2$ norm at most $C_d|B_j|^{-1/2}$. It therefore has uniformly
bounded $H^1$ norm by the same $2$-atom theorem. The last series
is absolutely summable in $H^1$. Adding it to $\sum b_j$ gives
$g$ as an actual $L^1$ and atomic sum, with the asserted norm.
Changes of center and scale preserve the normalized atom bounds
and the coefficient sum, proving the final assertion.
\end{proof}

\begin{theorem}[Convolution singular integrals preserve $H^1$]
\label{fa2:p4:cz-hardy-invariance}
Every standard convolution Calder\'on--Zygmund operator of
Part IV is bounded $H^1_{\rm at}(\mathbb R^d)\to
H^1_{\rm at}(\mathbb R^d)$. In particular this is stronger than
the previously proved $H^1\to L^1$ endpoint estimate.
\end{theorem}
\begin{proof}
Let $T$ have bounded Fourier multiplier $m$ and kernel $k$ off
zero, with $|\nabla k(x)|\le C|x|^{-d-1}$.
First take an infinity-atom supported in the unit ball centered
at zero. The $L^2$ bound gives $\|Ta\|_2\le C_d\|m\|_\infty$.
For $|x|\ge2$, the kernel identity and cancellation give
\[
 |Ta(x)|
 =\left|\int\bigl(k(x-y)-k(x)\bigr)a(y)\,dy\right|
 \le C|x|^{-d-1}\int|y|\,|a(y)|\,dy
 \le C|x|^{-d-1}.
\]
Consequently $Ta\in L^1$. Its Fourier transform is continuous,
and agrees almost everywhere with $m\widehat a$. Since
$\widehat a(0)=0$ and
$|\widehat a(\xi)|\le C_d|\xi|$, the latter product tends to
zero essentially as $\xi\to0$. Continuity then implies
$\widehat{Ta}(0)=0$, that is, $\int Ta=0$.
The preceding molecule lemma proves a uniform $H^1$ bound.

For an atom centered at $x_0$ of radius $r$, change variables
$x=x_0+rz$. The conjugated operator has multiplier $m(\xi/r)$
and kernel $r^d k(rz)$, with precisely the same $L^2$ and
gradient bounds, so the same uniform atomic bound follows.
If $f=\sum_j\lambda_j a_j$ is an atomic representation, then
$\sum_j\lambda_j Ta_j$ converges absolutely in the Banach
space $H^1$, with norm at most $C\sum_j|\lambda_j|$.
Its $L^1$ limit agrees with the already defined endpoint
$Tf$, by Proposition~\ref{fa2:p4:hardy-atomic}. Therefore it
is independent of the representation. Taking the infimum
over all representations proves the assertion.
\end{proof}

\begin{theorem}[The annular cancellation criterion]
\label{fa2:p4:cz-annular-criterion}
Let $k\in C^\infty(\mathbb R^d\setminus\{0\})$ satisfy
$|\partial^\alpha k(x)|\le C_\alpha|x|^{-d-|\alpha|}$.
There exists a Calder\'on--Zygmund distribution $K$ agreeing
with $k$ off zero if and only if
\[
 \sup_{0<r<R<\infty}
       \left|\int_{r<|x|<R}k(x)\,dx\right|<\infty.
\]
Here the integrals are ordinary integrals on finite annuli.
Such an extension is unique modulo a scalar multiple of $\delta_0$.
\end{theorem}
\begin{proof}
Assume the annular bound and define
\[
 K(\phi)=\int_{|x|<1}k(x)(\phi(x)-\phi(0))\,dx
              +\int_{|x|\ge1}k(x)\phi(x)\,dx.
\]
The subtraction makes the first integrand integrable, since it
is bounded in absolute value by $C|x|^{1-d}\|\nabla\phi\|_\infty$.
The second integral converges for Schwartz tests and is bounded
by a Schwartz seminorm. This defines a tempered distribution
agreeing with $k$ away from zero.

Let $\phi$ be supported in the unit ball with
$\|\phi\|_{C^1}\le1$, and write $\phi_s(x)=\phi(x/s)$.
For $0<s\le1$ split the defining integral at radius $s$:
\[
 K(\phi_s)
 =\int_{|x|<s}k(x)(\phi_s(x)-\phi(0))\,dx
      -\phi(0)\int_{s<|x|<1}k(x)\,dx.
\]
The first term has modulus at most
$C\int_0^s(s^{-1}t)\,dt/t\le C$, and the second is
bounded by the annular hypothesis. For $s\ge1$ the analogous
decomposition is
\[
 K(\phi_s)
 =\int_{|x|<s}k(x)(\phi_s(x)-\phi(0))\,dx
       +\phi(0)\int_{1<|x|<s}k(x)\,dx,
\]
with the same uniform bound. Thus all normalized dilated
$C^N$ bump pairings, for any fixed $N\ge1$, are bounded.
Proposition~\ref{fa2:p4:cz-multiplier} proves that $K$ is a
standard CZ distribution, including the bounded-multiplier
and $L^2$ assertions.

Conversely choose a smooth radial $\eta$ equal to one for
$|x|\le1$ and zero for $|x|\ge2$. The difference
$\eta(x/R)-\eta(x/r)$ vanishes near zero, so
\[
 \int k(x)\bigl(\eta(x/R)-\eta(x/r)\bigr)\,dx
      =K(\eta(\cdot/R))-K(\eta(\cdot/r))
\]
is uniformly bounded by the bump condition. Replacing each
smooth cutoff by its corresponding sharp ball indicator
changes the integral by at most
$C\int_r^{2r}dt/t+C\int_R^{2R}dt/t$, uniformly even when
the two transition annuli overlap. This proves the sharp
annular condition. Finally two such extensions differ by
a distribution supported at zero, hence by finitely many
delta derivatives. Its value on $\phi_s$ is a finite
polynomial in $s^{-1}$; uniform boundedness as $s\downarrow0$
for every bump eliminates every positive-order derivative.
Only a multiple of $\delta_0$ remains, and it does satisfy
the required bump bound.
\end{proof}

\begin{theorem}[Fourier divergence on a residual set of points]
\label{fa2:p4:generic-fourier-divergence}
There is a residual subset of $C(\mathbb T)$ such that, for
each of its functions $f$, the set
\[
 \{x\in\mathbb T:\sup_N|S_Nf(x)|=\infty\}
\]
is residual in $\mathbb T$. Here $S_N$ denotes the symmetric
Fourier partial sum, and residual means the complement of
a set of first category.
\end{theorem}
\begin{proof}
For each fixed $x$, the functionals $L_{N,x}f=S_Nf(x)$ have
unbounded norms by the Dirichlet-kernel calculation in
Theorem~\ref{fa2:p4:fourier-failures}. For each positive integer
$M$, the set
\[
 F_{x,M}=\{f:\sup_N|L_{N,x}f|\le M\}
\]
is closed and has empty interior. Indeed if it contained a
ball of radius $r$ centered at $f_0$, then testing $f_0$ and
$f_0+h$, $\|h\|_\infty<r$, would give
$|L_{N,x}h|\le2M$ for all $N$. Scaling would bound all the
operator norms by $2M/r$, a contradiction.

Choose a countable dense set $(x_j)$ in the circle. Baire's
theorem says that
$\mathcal G=\bigcap_{j,M}(C(\mathbb T)\setminus F_{x_j,M})$
is a dense $G_\delta$. For $f\in\mathcal G$ and fixed $M$, the
open set $O_M=\bigcup_N\{x:|S_Nf(x)|>M\}$ contains every $x_j$,
so is dense in the circle. A second application of Baire
shows that $\bigcap_M O_M$ is residual; it is exactly the
displayed set of unbounded partial sums. This proves the
two separate genericity assertions, rather than merely
divergence at one prescribed point.
\end{proof}

% END INPUT


% BEGIN parts/part4-vector-duality.tex
\chapter{Reflexive vector measures and the full Bochner dual}
\label{fa2:p4:vector-duality-chapter}

The range space need not be separable.  This matters: applying a
countable family of scalar Radon--Nikodym theorems is justified only
after a separable range has actually been obtained.  We obtain it
from weak compactness and the already proved Dunford--Pettis property
of scalar $L^1$.  The scalar measure is sigma-finite throughout the
duality theorem; all almost-everywhere statements can be made on its
completion without changing the Bochner spaces.  These results
complete the vector-valued function-space variants in
\cite{B01,B05,B08,B09}.

\section{A reflexive Radon--Nikodym theorem}

A countably additive vector measure $m:\Sigma\to Y$ is countably
additive in norm.  Its variation is
\[
 |m|(A)=\sup\left\{\sum_{j=1}^r\|m(A_j)\|:
                  A=\coprod_{j=1}^r A_j,\ A_j\in\Sigma\right\}.
\]
When $|m|(X)<\infty$, this is a finite positive countably additive
measure.  Indeed for disjoint $B_n$, partitions of the first finitely
many $B_n$ show $|m|(\bigcup B_n)\ge\sum_n|m|(B_n)$.
Conversely for any finite partition $(A_j)$ of the union, norm
countable additivity and triangle inequality give
$\sum_j\|m(A_j)\|\le\sum_n\sum_j\|m(A_j\cap B_n)\|
\le\sum_n|m|(B_n)$; take its supremum.  Also $m\ll\mu$ implies
$|m|\ll\mu$, because every measurable subset of a null set is null.

\begin{lemma}[A compact range before any density is assumed]
\label{fa2:p4:reflexive-measure-range}
A finite-variation vector measure taking values in a reflexive Banach
space has relatively norm compact range.  Consequently the closed
linear span of its range is a separable reflexive Banach space.
\end{lemma}
\begin{proof}
Put $\lambda=|m|$.  The case $\lambda(X)=0$ is trivial.  For a
scalar simple function $s=\sum_j a_j1_{A_j}$ define
$Ts=\sum_j a_jm(A_j)$.  Refining partitions makes this well defined
and $\|Ts\|\le\int|s|\,d\lambda$; it extends to a bounded
operator $T:L^1(\lambda)\to Y$ by completeness.
Since $Y$ is reflexive, $T$ is weakly compact: its unit-ball image
is bounded and lies in a weakly compact ball of $Y$.
By Theorem~\ref{fa2:p2:dp-classical-examples}, $L^1(\lambda)$ has
the Dunford--Pettis property, so $T$ is completely continuous.
The set $H=\{1_A:A\in\Sigma\}$ is bounded and uniformly integrable
in $L^1(\lambda)$, because
$\int_B|1_A|\,d\lambda\le\lambda(B)$ and $\lambda(X)<\infty$.
There is no tail condition to add on this finite measure space.
The scalar $L^1$ weak compactness criterion
(Theorem~\ref{fa2:p4:dunford-pettis-l1}) makes its weak closure
weakly compact.  A completely continuous map sends this closure
to a relatively norm compact set, by
Proposition~\ref{fa2:p2:complete-continuity} and
Eberlein--\v Smulian.  Since $T1_A=m(A)$, this proves the range
assertion.
A relatively compact metric set is separable: take a finite
$1/n$-net for each $n$.  Its closed linear span is separable,
and it is reflexive as a closed subspace of $Y$.
\end{proof}

\begin{theorem}[Radon--Nikodym for reflexive range, including nonseparable range]
\label{fa2:p4:vector-radon-nikodym}
Let $(X,\Sigma,\mu)$ be a sigma-finite positive measure space and
$Y$ a reflexive Banach space.  If $m:\Sigma\to Y$ is norm
countably additive, has finite variation, and $m\ll\mu$, then
there is a unique Bochner integrable $g:X\to Y$, up to null sets,
such that
\[
 m(A)=\int_Ag\,d\mu\quad(A\in\Sigma),\qquad
 |m|(A)=\int_A\|g\|\,d\mu.
\]
In particular no separability assumption on $Y$ is needed.
\end{theorem}
\begin{proof}
Let $Y_0$ be the separable reflexive closed span of the range given
by the preceding lemma.  Its dual is norm separable.  To verify
this useful fact, Alaoglu and separability of $Y_0$ make
$B_{Y_0'}$ weakly star compact metrizable, hence it has a countable
weakly star dense subset.  Reflexivity makes this its weak
topology.  The convex hull of that countable subset is weakly
dense in the ball, so Hahn--Banach separation makes its norm
closure the ball.  Rational convex combinations are a countable
norm dense set.  Choose a countable dense vector subspace
$D\subset Y_0'$ over $\mathbb Q$ or $\mathbb Q+i\mathbb Q$.

Scalar Radon--Nikodym gives $h=d|m|/d\mu\in L^1_+(\mu)$ and,
for every $\ell\in D$, a density $h_\ell$ for the scalar measure
$\ell m$.  Since $|\ell m|\le\|\ell\|\,|m|$,
$|h_\ell|\le\|\ell\|h$ almost everywhere.  Uniqueness of scalar
densities gives all rational linear relations among the $h_\ell$.
There are only countably many such bounds and relations.  Outside
one null set (also remove $\{h=\infty\}$), therefore,
$\ell\mapsto h_\ell(t)$ is a bounded scalar-linear functional
on $D$ of norm at most $h(t)$.  It extends by density to an
element of $Y_0''$.  Reflexivity supplies a unique $g(t)\in Y_0$
with $\ell(g(t))=h_\ell(t)$ for $\ell\in D$ and
$\|g(t)\|\le h(t)$.  Set $g=0$ on the removed null set.

For every $\ell\in Y_0'$, approximation by $D$ makes
$\ell(g(t))$ a pointwise limit of measurable functions; thus
$g$ is weakly measurable.  Its range lies in separable $Y_0$,
so Pettis's measurability criterion makes it strongly measurable
(work on the completed measure space).  Its norm is bounded
by the integrable $h$, hence $g$ is Bochner integrable.
For every $A$ and $\ell\in D$,
$\ell(\int_Ag)=\int_Ah_\ell=\ell(m(A))$.
Density of $D$ separates points of $Y_0$ and proves the vector
identity.

If a Bochner integrable function has zero integral on all sets,
its scalar tests do also.  Choose a countable norming family on
the separable span of its essential range and extend it by
Hahn--Banach to the ambient space.  Each scalar test vanishes
almost everywhere, so, outside their common null set, the vector
vanishes.  This proves uniqueness.
Finally for a simple density $s$ the variation of its vector
measure is exactly $\int\|s\|$, by its disjoint defining partition.
For general $g$ approximate in $L^1$ by simple functions.
The variation norm of the difference of their measures is at
most the $L^1$ norm of the density difference, while
$|\int\|s\|-\int\|g\||\le\int\|s-g\|$.
Passing to the limit, and applying the same reasoning after
restriction to each $A$, gives the stated variation formula.
\end{proof}

\begin{example}[The range condition cannot be discarded]
\label{fa2:p4:vector-rn-failure}
On $[0,1]$ with Lebesgue measure let $m(A)=1_A$ with values in
$Y=L^1[0,1]$.  This is norm countably additive and
$|m|(A)=\mu(A)$.  It has no Bochner density.  Indeed a measure
with Bochner density has relatively norm compact range:
approximate that density in $L^1$ by finite-valued simple
functions; the ranges of their integrals over arbitrary sets
lie in compact finite dimensional boxes and approximate the
range uniformly.  But the indicators of independent dyadic
half-sets have pairwise $L^1$ distance $1/2$.  Thus the range
here is not relatively norm compact.  Reflexivity is a
sufficient range hypothesis in the theorem, not a property
of all Banach spaces.
\end{example}

\section{The duality and reflexivity theorems}

\begin{theorem}[Full Bochner $L^p$ duality for reflexive $E$]
\label{fa2:p4:vector-lp-dual}
Let $\mu$ be sigma-finite and $E$ a reflexive real or complex
Banach space, without a separability assumption.  If
$1\le p<\infty$ and $1/p+1/q=1$, with $q=\infty$ for $p=1$,
then
\[
 L^q(\mu;E')\longrightarrow L^p(\mu;E)',\qquad
 g\longmapsto\left[f\longmapsto\int_Xg(t)(f(t))\,d\mu(t)\right]
\]
is a surjective linear isometry.  Both vector-valued spaces
consist of strongly measurable functions, modulo almost-everywhere
equality.  The pairing is bilinear, with no Hilbert conjugation.
\end{theorem}
\begin{proof}
H\"older and the Bochner measurability approximation make the
displayed scalar integral well defined and bound its functional
norm by $\|g\|_q$.
For surjectivity let $L\in L^p(\mu;E)'$.  First suppose
$\mu(X)<\infty$, and define $m(A)\in E'$ by
$m(A)(x)=L(1_Ax)$.  Then
$\|m(A)\|\le\|L\|\mu(A)^{1/p}$.
For disjoint sets, the tail bound proves norm countable
additivity.  If $(A_j)$ is a finite partition, choose unit
$x_j\in E$, adjusting phases, so
$m(A_j)(x_j)$ is arbitrarily close to $\|m(A_j)\|$ and real
nonnegative.  The simple function $\sum_j1_{A_j}x_j$ has
$L^p$ norm at most $\mu(X)^{1/p}$.  Thus
\[
 |m|(X)\le\|L\|\mu(X)^{1/p}<\infty.
\]
It is absolutely continuous with respect to $\mu$.
The space $E'$ is reflexive by Part I.  The vector
Radon--Nikodym theorem yields a strongly measurable
$g\in L^1(\mu;E')$ representing $m$.
It follows initially for finite simple $f$ that
$L(f)=\int g(f)$.

For sigma-finite $\mu$, partition $X$ into disjoint
finite-measure pieces and apply this construction to the
restriction of $L$ on each piece.  The resulting densities
give one strongly measurable $g$: their countably many
essential separable ranges have separable closed span, and
finite-valued approximants on the first finitely many pieces,
zero elsewhere, converge pointwise.  On every finite union
of these pieces, $g$ is integrable and the simple-function
identity holds.  It also holds for bounded strongly measurable
$f$ supported there: nearest-center simple approximants can
be chosen uniformly bounded and converge pointwise, giving
$L^p$ convergence and convergence against the integrable
$\|g\|$ by dominated convergence.

We prove the necessary $L^q$ bound explicitly; no measurable
selection of an uncountable norming family is assumed.
On the separable essential range of $g$ in $E'$, choose a
countable family $(x_k)\subset B_E$ which norms it.  Such a
family is obtained by nearly norming a countable dense
family $(a_j)$ of functionals: for each $j,r$ choose a
unit-ball vector with $|a_j(x)|>\|a_j\|-1/r$, and include
a countable dense set of scalar phases of all these vectors.
Density and arbitrary $r$ make their real evaluations have
supremum exactly $\|a\|$ for every $a$ in that range.
Given $0<\eta<1$, choose at each
$t$ the first $x_k$ such that
\[
 \operatorname{Re}g(t)(x_k)\ge(1-\eta)\|g(t)\|;
\]
put zero when $g(t)=0$.  This choice is measurable and
countably valued, hence strongly measurable, and has norm
at most one.
For $1<p<\infty$ multiply it by
$\|g(t)\|^{q-1}1_A$, where $A$ lies in a finite union of the
measure pieces and $1/n\le\|g(t)\|\le n$ there.  This is
a bounded strongly measurable test, so
\[
 (1-\eta)\int_A\|g\|^q
 \le\operatorname{Re}L(f)
 \le\|L\|\left(\int_A\|g\|^q\right)^{1/p}.
\]
Let $\eta\downarrow0$ and exhaust by these sets to get
$\|g\|_q\le\|L\|$.
For $p=1$, if $\|g\|>M>\|L\|$ on a set of positive measure,
sigma-finiteness gives a positive finite-measure portion
in one finite union of pieces.  Use the same norming
choice times its indicator and take $\eta$ with
$(1-\eta)M>\|L\|$ to obtain a contradiction.  Hence
$\|g\|_\infty\le\|L\|$.
H\"older and simple-function density now extend the
representation to all $L^p$ and give equality of norms.
The same tests applied to any representing $g$ prove
injectivity, including the endpoint.  This proves the
surjective isometry.
\end{proof}

\begin{theorem}[Bochner reflexivity without a separability restriction]
\label{fa2:p4:vector-lp-reflexive}
If $1<p<\infty$, $\mu$ is sigma-finite, and $E$ is reflexive,
then $L^p(\mu;E)$ is reflexive.  Conversely, if $\mu$ has a
set of positive finite measure and $L^p(\mu;E)$ is reflexive,
then $E$ is reflexive.  No analogous unconditional
reflexivity assertion holds for $p=1$ or $p=\infty$.
\end{theorem}
\begin{proof}
The preceding theorem identifies the dual with
$L^q(\mu;E')$.  Apply it again, since $E'$ is reflexive,
to identify the bidual with $L^p(\mu;E'')$.
Under these bilinear pairings the canonical embedding is
exactly $f\mapsto J_E\circ f$, because evaluation on a
dual function $g$ gives $\int g(t)(f(t))$.
The bounded linear inverse of the onto isometry
$J_E:E\to E''$ preserves strong measurability and norms,
so every member of $L^p(\mu;E'')$ is obtained this way.
Thus the canonical embedding is onto.
For the converse, a set $A$ with $0<\mu(A)<\infty$ gives
the closed isometric copy $x\mapsto\mu(A)^{-1/p}1_Ax$
of $E$ in $L^p(\mu;E)$.  Closed subspaces of reflexive
spaces are reflexive.  Finally the scalar spaces
$\ell^1=L^1(\mathbb N)$ and $\ell^\infty=L^\infty(\mathbb N)$
are nonreflexive, already providing the endpoint failures
with reflexive scalar range.
\end{proof}

% END INPUT


% BEGIN parts/part4-krein.tex
\chapter{Even-test positivity: cosine and hyperbolic-cosine measures}
\label{fa2:p4:krein-chapter}
Ordinary positive definiteness tests against all functions. Testing
only against even functions is weaker: $\cosh(as)$ passes those tests
and is unbounded. This is the distinct Krein result discussed in
\cite[Section 14.5]{B08}. We prove the representation through a
symmetric second-derivative operator. Its self-adjoint extension need
not be unique, and neither need the representing measure. The word
``uniquely'' in the printed general representation must therefore not
be imported into the statement. We also give the compensated odd-test
form, which includes the quadratic case missing from an uncompensated
finite-measure formula.

For $\lambda\in\mathbb R$ define the real even and odd solutions
\[
 C_\lambda(t)=
 \begin{cases}\cos(\sqrt{-\lambda}\,t)&\lambda<0,\\
 1&\lambda=0,\\\cosh(\sqrt\lambda\,t)&\lambda>0,
 \end{cases}
 \qquad
 S_\lambda(t)=
 \begin{cases}\sin(\sqrt{-\lambda}\,t)/\sqrt{-\lambda}&\lambda<0,\\
 t&\lambda=0,\\\sinh(\sqrt\lambda\,t)/\sqrt\lambda&\lambda>0.
 \end{cases}
\]
Both solve $w''=\lambda w$, with initial data $(1,0)$ and $(0,1)$.
Their addition identities hold also at zero by limits, in particular
\[
 C_\lambda(s-t)+C_\lambda(s+t)=2C_\lambda(s)C_\lambda(t),
 \quad C_\lambda(s-t)-C_\lambda(s+t)
                    =-2\lambda S_\lambda(s)S_\lambda(t).
\]

\section{A spectral kernel lemma with its measurability bridge}
\begin{lemma}[A continuous kernel satisfying the wave identity]
\label{fa2:p4:krein-kernel}
Let $F$ be continuous, real and even on $\mathbb R$. If
$K_+(s,t)=(F(s-t)+F(s+t))/2$ is positive definite, there is a
positive finite Borel measure $\nu$ on $\mathbb R$ with
\[
 F(s)=\int C_\lambda(s)\,d\nu(\lambda),\qquad
 \nu(\mathbb R)=F(0),\qquad
 \int_{\lambda>0}\cosh(R\sqrt\lambda)\,d\nu<\infty\quad(R>0).
\]
If instead $K_-(s,t)=(F(s-t)-F(s+t))/2$ is positive definite, there
is a positive locally finite measure $\nu$ for which
\[
 F(s)=F(0)+\int\frac{1-C_\lambda(s)}{\lambda}\,d\nu(\lambda),
\]
where the integrand at $\lambda=0$ is $-s^2/2$. Its integrals are
absolutely convergent, uniformly dominated on compact $s$-intervals.
\end{lemma}
\begin{proof}
We give the two constructions together, writing $\sigma=+$ or $-$.
The positive kernel construction gives a complex Hilbert space $H$
spanned by vectors $v(t)$ with
$\langle v(s),v(t)\rangle=K_\sigma(s,t)$. Null vectors are
quotiented out before completion. The kernel identities give
$v(-t)=v(t)$ for $+$ and $v(-t)=-v(t)$ for $-$, and continuity of
the kernel makes $v$ norm continuous. Consequently $H$ is separable:
the vectors at rational parameters have dense span. There is a
conjugation on $H$ fixing the real generators $v(t)$.

For a compactly supported smooth function $f$ put
$Jf=\int f(t)v(t)\,dt$. One may restrict to the corresponding parity
of $f$, since the other parity is killed. These vectors span densely,
by smooth approximate point masses. Define $S_0Jf=Jf''$.
The kernel has the distribution identity
$\partial_s^2K_\sigma=\partial_t^2K_\sigma$; integration by parts
twice on the compact supports therefore gives
$\langle Jf'',Jg\rangle=\langle Jf,Jg''\rangle$.
This also proves well-definedness: if $Jf=0$, then $Jf''$ is
orthogonal to the dense set of all $Jg$ and hence is zero.
Thus $S_0$ is densely defined symmetric. It commutes with the
conjugation, which carries $\ker(S_0^*-i)$ antiunitarily onto
$\ker(S_0^*+i)$. The deficiency-index theorem of Part III supplies
a self-adjoint extension $S$ of its closure. We do not claim
essential self-adjointness or uniqueness of this extension.

Use the separable spectral multiplicity model for $S$ to write
$UH=\int^\oplus H_\lambda\,d\mu(\lambda)$ with $S$ multiplication
by the real coordinate and $\mu$ a finite scalar control measure.
For an unbounded $S$ this model is obtained by applying the bounded
normal model to its Cayley transform and the Borel inverse coordinate;
the fixed-point spectral projection is zero. The fibers have finite
or countable dimension and may be placed in one $\ell^2$.

There is a genuine pointwise-distribution bridge here. On each
interval $[-R,R]$, the map $J_R:L^2[-R,R]\to H$ is Hilbert--Schmidt,
with squared norm $\int_{-R}^R\|v(t)\|^2dt<\infty$. This follows
by Parseval and Tonelli applied to the vector-valued continuous
kernel $v(t)$ and an orthonormal basis of $H$ (or of the source).
The Hilbert--Schmidt kernel theorem therefore represents $UJ_R$ by
a square-integrable fiber-vector kernel $h_R(\lambda,t)$.
Explicitly embed the fibers in $\ell^2$, hence the direct integral
in $L^2(\mu\times\text{counting measure})$, and apply the scalar
Hilbert--Schmidt kernel theorem. Its coordinate kernels combine by
Tonelli into this vector kernel and remain in the permitted fiber
a.e., since the operator's range does so.
For integer $R$ these kernels agree on overlaps by uniqueness, so
after one null set is removed they define an even, respectively odd,
$H_\lambda$-valued $h_\lambda\in L^2_{\rm loc}(\mathbb R)$.
The identity $S Jf=Jf''$ implies
\[
 \int h_\lambda(t)f''(t)\,dt
                  =\lambda\int h_\lambda(t)f(t)\,dt.
\]
Choose a countable dense collection of smooth tests on each integer
interval, in the topology controlling their first two derivatives.
The identities hold simultaneously outside one $\mu$-null set and
extend to all smooth compactly supported tests by the local $L^2$
bound. Thus $h_\lambda''=\lambda h_\lambda$ as a distribution.
The elementary distributional ODE has only the classical two
solutions: multiply the first-order factorizations by an exponential,
use that a distribution of derivative zero is constant, and treat
$\lambda=0$ by integrating twice. The same conclusion coordinatewise
in the countable fiber applies to a vector-valued distribution.
Parity consequently gives
\[
 h_\lambda(t)=c(\lambda) C_\lambda(t)\quad(\sigma=+),\qquad
 h_\lambda(t)=c(\lambda) S_\lambda(t)\quad(\sigma=-).
\]
The vector $c(\lambda)$ is measurable: on $[-1,1]$ integrate
$h_\lambda$ against the corresponding real solution and divide
by its strictly positive squared $L^2$ norm. It is a fiber vector
of finite norm a.e., since that solution is nonzero. Define
$d\nu=\|c(\lambda)\|^2d\mu$. The squared solution norm on
$[-1,1]$ has a positive lower bound on every compact $\lambda$-set;
the Hilbert--Schmidt identity therefore makes $\nu$ locally finite.
For smooth compactly supported $f,g$ of the chosen parity we have
the exact identity
\[
 \langle Jf,Jg\rangle
   =\int\overline{\left(\int f(t)Z_\lambda(t)dt\right)}
                 \left(\int g(t)Z_\lambda(t)dt\right)d\nu,
\]
where $Z=C$ or $S$. In addition
$\int\int_{-R}^R|Z_\lambda(t)|^2dt\,d\nu<\infty$ for every $R$.

For the plus case let $\rho\ge0$ be a smooth even mollifier of
integral one, and put $b_\varepsilon(\lambda)=
\int\rho_\varepsilon(t)C_\lambda(t)dt$. It tends to one for
each $\lambda$, has absolute value at most one for $\lambda\le0$,
and for $\lambda>0$ it is at least one and decreases to one as
$\varepsilon\downarrow0$. Since $J\rho_\varepsilon\to v(0)$,
Fatou first gives $\nu(\mathbb R)\le\|v(0)\|^2=F(0)$.
Now dominated convergence applies: on the negative half-line use
this finite mass, and on the positive half-line dominate by the
integrable $b_{\varepsilon_0}^2$. It gives equality of the mass.
Apply the same argument to
$(\rho_\varepsilon(t-s)+\rho_\varepsilon(t+s))/2$, whose transform
is $C_\lambda(s)b_\varepsilon(\lambda)$, to get
\[
 \|v(s)\|^2=\int C_\lambda(s)^2d\nu.
\]
Fatou gives integrability first; domination then gives equality.
On the positive half-line the dominant is
$C_\lambda(s)^2b_{\varepsilon_0}^2$, integrable by the identity
for that fixed smooth test. The negative half-line is bounded by
the already finite mass. Since $\cosh(2r)=2\cosh^2r-1$, this also
gives all the claimed hyperbolic exponential integrals. The cross
inner products of the two mollified vectors converge, by
Cauchy--Schwarz or the sum of their squared majorants. They yield
$F(s)=\langle v(0),v(s)\rangle=\int C_\lambda(s)d\nu$.

For the minus case use odd approximate point masses
$(\rho_\varepsilon(t-s)-\rho_\varepsilon(t+s))/2$.
Their transforms are $S_\lambda(s)b_\varepsilon(\lambda)$.
Fatou followed by the same positive-half-line domination gives
$\|v(s)\|^2=\int S_\lambda(s)^2d\nu$; on the negative half-line
the integrable $S_\lambda(s)^2$ obtained by Fatou is itself a
dominant, so finiteness of the total mass is not needed.
Using $\|v(s/2)\|^2=(F(0)-F(s))/2$ and
$C_\lambda(s)-1=2\lambda S_\lambda(s/2)^2$ proves the desired
compensated formula, including its limit at zero.
Its compact-interval domination follows also from the local
Hilbert--Schmidt bound. Explicitly, for $\lambda\le0$ the function
$\sup_{|s|\le R}|S_\lambda(s)|^2$ is bounded by a constant times
$(1+|\lambda|)^{-1}$ away from a bounded $\lambda$-interval, and
$\int_{-1}^1S_\lambda(t)^2dt$ bounds that weight below there.
For $\lambda>0$, the integral of $S_\lambda(t)^2$ on the longer
interval $[-R-1,R+1]$ dominates its maximum on $[-R,R]$, up to a
constant depending on $R$. This follows directly from the sinh
formula for large $\lambda$, and by continuity on the remaining
compact interval. Thus these suprema are integrable against $\nu$.
The identity of the compensated integrand with
$-2S_\lambda(s/2)^2$ now gives uniform compact domination.
If the constructed Hilbert space is zero, the plus case has $F=0$
and the minus case has constant $F$, and the zero measure supplies
both conclusions directly.
\end{proof}

\begin{theorem}[Krein's even-test representation and extreme functions]
\label{fa2:p4:krein-even}
A continuous real even $p$ satisfies
\[
 \iint p(s-t)f(s)f(t)\,dsdt\ge0
 \quad\text{for every real even continuous compactly supported }f
\]
if and only if it has a positive measure representation
$p(s)=\int C_\lambda(s)d\nu(\lambda)$ with the finite mass and
exponential integrability stated in the preceding lemma.
Equivalently it is a positive cosine mixture plus a positive cosh
mixture; the zero atom may be assigned only to the cosine part to
avoid duplicating the constant function. The normalized convex set
$p(0)=1$ has extreme points exactly $\cos(as)$, $a\ge0$, and
$\cosh(bs)$, $b>0$ (the constant is already the case $a=0$).
\end{theorem}
\begin{proof}
Smooth even approximations of point masses show that the integral
positivity is equivalent to positive definiteness of $K_+$.
Indeed symmetrizing each point mass in the finite quadratic sum
gives exactly that kernel; conversely integrate its finite-sum
positivity by Riemann sums on the compact support. Continuous
tests follow from smooth tests by uniform approximation and the
boundedness of $p$ on the relevant compact difference interval.
The preceding lemma gives the representation. Conversely the
cosine addition identity makes $K_+$ an integral of the positive
rank-one kernels $C_\lambda(s)C_\lambda(t)$. Compact domination
justifies integrating the test functions, proving positivity.
Push $\nu|_{(-\infty,0]}$ forward by $\sqrt{-\lambda}$ and
$\nu|_{(0,\infty)}$ by $\sqrt\lambda$ for the two-mixture form.

If $p(0)=1$, $\nu$ is a probability. For any measurable set $A$
of probability strictly between zero and one, its two normalized
restrictions give a convex decomposition of $p$ into the same
class; continuity follows from compact domination. If $p$ is
extreme, this forces
$\int_A C_\lambda(s)d\nu=\nu(A)p(s)$ for every $A$ and $s$.
For countably many rational $s$ this means
$C_\lambda(s)=p(s)$ off one null set; continuity extends it to all
$s$. Different $\lambda$ give different functions, as their second
derivatives at zero are different. Thus the probability is
concentrated at one $\lambda$ and $p=C_\lambda$.
Conversely the kernel of $C_\lambda$ is rank one. If it is the
average of two normalized positive kernels, the squared norm of
the vector corresponding to
$\delta_s-C_\lambda(s)\delta_0$ is zero in each, since it is zero
in their sum and each squared norm is nonnegative. Each kernel
vector is therefore $C_\lambda(s)$ times its vector at zero.
Its value against the latter gives $p_j(s)=C_\lambda(s)$.
Hence no nontrivial decomposition exists, proving extremality.
This argument does not assume uniqueness of a representing measure.
\end{proof}

\begin{proposition}[The measure need not be unique]
\label{fa2:p4:krein-nonunique}
General continuous even-test-positive functions can have different
positive cosine--cosh representations, even after removing the
duplicate zero atom.
\end{proposition}
\begin{proof}
Let $h(s)=e^{-\cosh s}$. It is real, even and Schwartz: each derivative
is a polynomial in $\sinh s,\cosh s$ times the same rapidly decreasing
factor. Fourier inversion writes it as the cosine transform of a
finite real signed measure $\eta$ with a real even Schwartz density
on $[0,\infty)$ (fold the ordinary angular-frequency Fourier density).
On the other hand expand the two exponentials in
$e^{-(e^s+e^{-s})/2}$ and group equal exponents. Absolute convergence
for every bounded $s$ gives
\[
 h(s)=c_0+2\sum_{n\ge1}(-1)^n c_n\cosh(ns),\qquad
 c_n=\sum_{j\ge0}\frac1{2^{n+2j}j!(n+j)!}>0.
\]
The absolute series has every exponential moment, since the
ungrouped absolute double sum is $e^{\cosh s}$. Denote this finite
signed cosh measure by $\theta$. Taking the Jordan decompositions
of the two measures gives
\[
 \int\cos(as)d\eta^+(a)+\int\cosh(bs)d\theta^-(b)
 =\int\cos(as)d\eta^-(a)+\int\cosh(bs)d\theta^+(b).
\]
Both sides are even-test positive by their positive representations.
They give different measures, since the nonzero positive-axis
atoms of $\theta$ alternate in sign and the cosine part is on the
negative spectral half-axis. Reassigning its zero atom to the
cosine part cannot remove this difference. Their common value at
zero is positive, so division by it gives a normalized example
if desired. This is a concrete counterexample to general
representation uniqueness, independently of any bibliographic claim.
\end{proof}

\begin{theorem}[The correct compensated odd-test form]
\label{fa2:p4:krein-odd}
A continuous real even $q$ is nonnegative on all odd real compactly
supported tests in the analogous quadratic integral if and only if
it has the compensated representation of the kernel lemma, for a
positive locally finite measure with compact-interval domination.
Its constant $q(0)$ is unrestricted. In particular $q(s)=-s^2$
has this positivity, although it has no representation as a finite
positive cosine mixture minus a finite positive cosh mixture.
\end{theorem}
\begin{proof}
Odd approximate point masses identify test positivity with positive
definiteness of $K_-$, and the lemma proves necessity. Conversely
the addition formula makes its kernel
$\int S_\lambda(s)S_\lambda(t)d\nu$, which is positive by the
same compact domination and Cauchy--Schwarz. The constant cancels
from every odd test. For $q=-s^2$ its kernel is $2st$, positive
of rank one; equivalently the quadratic integral on an odd $f$ is
$2(\int s f(s)ds)^2$. If a finite positive cosh measure had positive
mass on $[\varepsilon,\infty)$ for some $\varepsilon>0$, its cosh
transform would grow at least exponentially on the real line.
A bounded finite cosine transform could not cancel it to $-s^2$.
Thus such a cosh measure would have to be concentrated at zero,
making the proposed difference bounded, again a contradiction.
The compensated form represents the quadratic by an atom of mass
two at the spectral parameter zero and has no such defect.
\end{proof}

% END INPUT


% BEGIN parts/part4-convex-integral-exercises.tex
\chapter{Support constraints, nonlinear integral limits, and weak $L^p$}
These are the functional-analytic assertions of \cite[Problems
17--19 and 21]{B01}. They are stated with their actual measure
and endpoint assumptions, rather than absorbed into an unspecified
``application'' heading. All functions in the convex statements
are real valued.

\begin{theorem}[Support inequalities and constrained extensions]
\label{fa2:p4:constrained-support}
Let $X$ be reflexive, $M\subset X$ a closed subspace, and
$C\subset X'$ nonempty and convex. Then
\[
 f(u)\leq\sup_{g\in C}g(u)\ (u\in M)
       \quad\Longleftrightarrow\quad
 f\in\overline{M^\perp+C}.
\]
If $C$ is also closed and bounded the sum is already closed.
In particular, for $1<p<\infty$, $j,k\in L^{p'}$, $k\geq0$,
and $M=\{u\in L^p:\int ju=0\}$, the inequality
$\int fu\leq\int k|u|$ on $M$ is equivalent to
$f=\lambda j+g$ for some real $\lambda$ and $|g|\leq k$ a.e.
For a sigma-finite measure space, $M$ any linear subspace of $L^1$
and $a,b\in L^\infty$ with $a\leq b$, there is
$h\in M^\perp\cap L^\infty$ with $a\leq h\leq b$ iff
\[
 \int(a u^+-b u^-)\leq0\quad(u\in M).
\]
On measures without a scalar $L^1$ dual representation the
Hahn--Banach conclusion is instead a dominated continuous dual
functional; a global $L^\infty$ representative is not presumed.
\end{theorem}
\begin{proof}
One implication of the first statement holds on $M^\perp+C$
and survives norm limits. If $f$ is outside its closed convex
closure, separation in $X'$ gives $u\in X''=X$ with
$f(u)>\sup_{g\in M^\perp+C}g(u)$. Finiteness of this
supremum forces $u\in M$, by scaling arbitrary members of
$M^\perp$ and using its preannihilator. This contradicts the
proposed inequality. If $C$ is closed and bounded it is weakly
compact. In a weakly convergent net $m_i+g_i$, take a subnet
$g_i\rightharpoonup g\in C$; the remainder converges into the
weakly closed space $M^\perp$. This makes the sum weakly closed.

For the $L^p$ assertion, $C=\{g:|g|\leq k\}$ is convex,
norm closed and bounded in the reflexive dual. Its support
function is exactly $\int k|u|$, attained by
$g=k\operatorname{sgn}u$. The annihilator of $\ker j$ is the
span of $j$: for $j\ne0$ choose $u_0$ with $j(u_0)=1$ and
subtract $j(u)u_0$; for $j=0$ the assertion is immediate.
Apply the proved closed-sum result and the arbitrary-measure
interior $L^p$ duality theorem.

For the last assertion necessity follows from
$a u^+-b u^-\leq hu$. Conversely write $s=a+b$ and
$w=b-a\geq0$. The given inequality and the same one for $-u$
say $|\int su|\leq\int w|u|$ on $M$.
Extend $u\mapsto\int su$ from $M$ by real Hahn--Banach,
dominated by the sublinear function $\int w|u|$ on $L^1$.
It is bounded, and sigma-finite duality represents it by $v$
with $|v|\leq w$ a.e. The pointwise bound follows by testing
finite-measure subsets where it would fail and choosing their
sign; exhaustion covers the whole space. Put $h=(s-v)/2$.
It annihilates $M$ and lies between $a$ and $b$, as required.
\end{proof}

\begin{theorem}[Weak convergence with convergence of convex energy]
\label{fa2:p4:nemytskii-energy}
Let $\mu(\Omega)<\infty$, $1<p<\infty$, and let $g:\mathbb R\to
\mathbb R$ be continuous nondecreasing with
$|g(t)|\leq C(1+|t|^{p-1})$. Put $G(t)=\int_0^t g(s)ds$.
If $u_n\rightharpoonup u$ in $L^p$ and
$\limsup\int G(u_n)\leq\int G(u)$, then
$g(u_n)\to g(u)$ in every $L^q$, $1\leq q<p'$.
If $g$ is strictly increasing, $u_n\to u$ in every $L^q$,
$1\leq q<p$. If also
$|g(t)|\geq\alpha|t|^{p-1}-C$ for some $\alpha>0$,
the respective conclusions hold in $L^{p'}$ and, in the
strict case, in $L^p$. These endpoint conclusions can fail
without that lower growth assumption, even for strictly increasing $g$.
\end{theorem}
\begin{proof}
The growth bound puts $g(u)$ in $L^{p'}$ and $G(u)$ in $L^1$.
Convexity gives the nonnegative Bregman remainder
$D(a,b)=G(a)-G(b)-g(b)(a-b)$.
Weak convergence makes $\int g(u)(u_n-u)\to0$; the assumed
limsup and nonnegativity imply $\int D(u_n,u)\to0$.

For fixed $b$, $D(a_n,b)\to0$ implies $g(a_n)\to g(b)$.
On bounded intervals this follows by continuity and the fact
that an integral of the nonnegative difference $g(s)-g(b)$
can vanish only where $g$ stays constant between the endpoints.
An unbounded subsequence with $g(a_n)$ separated from $g(b)$
would eventually lie past a fixed point at which that difference
has the same sign and a fixed positive magnitude. Integrating
beyond that point makes $D$ grow, a contradiction. If $g$ is
strict, the same argument yields $a_n\to b$.
Every subsequence of $D(u_n,u)$ has a further subsequence
converging to zero a.e., by its $L^1$ convergence. The scalar
argument therefore gives convergence in measure of $g(u_n)$,
and of $u_n$ when strict. The $L^{p'}$ bound on $g(u_n)$
and $L^p$ bound on $u_n$ make their lower powers uniformly
integrable on this finite measure space. Vitali then proves
the stated subendpoint convergences for the whole sequence.

Under the lower growth assumption, monotonicity fixes the
eventual signs of $g$ at the two ends. Integration gives
$G(t)\geq c|t|^p-C_1|t|-C_2$ with $c>0$.
Young's inequality and the upper growth at $b$ consequently give
\[
 |a|^p\leq C_3D(a,b)+C_4(1+|b|^p).
\]
The family $D(u_n,u)$ converges to zero in $L^1$, hence is
uniformly integrable. This inequality makes $|u_n|^p$ uniformly
integrable, and the upper growth does the same for
$|g(u_n)|^{p'}$. The already proved convergence in measure and
Vitali now prove the endpoint assertions.

Here is a counterexample to dropping lower growth. Choose
$a_n\uparrow\infty$ so fast that
$a_{n-1}^{p-1}/a_n^{p-1}\to0$. Define a continuous
nondecreasing $g_0$ to be zero on the negative half-line, to
increase linearly from $a_{n-1}^{p-1}$ to $a_n^{p-1}$ on
$[(1-1/n)a_n,a_n]$ for $n\geq2$, and to be constant between
these ramps; start it continuously near zero with $g_0(0)=0$. The sequence
may be enlarged to make the intervals disjoint. Then
$g_0(t)\leq C(1+|t|^{p-1})$ and
$G_0(a_n)/a_n^p\leq
a_{n-1}^{p-1}/a_n^{p-1}+1/n+o(1)\to0$.
On $(0,1)$ put $u_n=a_n1_{(0,a_n^{-p})}$.
These converge weakly to zero in $L^p$: Holder bounds any
dual pairing by the $L^{p'}$ norm of its test on the shrinking
interval. Their $p$th norms equal one, their energies tend to
zero, but $\|g_0(u_n)\|_{p'}=1$. Adding $\arctan t$ to
$g_0$ makes $g$ strictly increasing, preserves the upper growth,
and adds only $O(a_n)$ to $G(a_n)$. Its endpoint failures
therefore persist. This verifies both warnings constructively.
\end{proof}

\begin{example}[A minimizer in a nonreflexive intersection space]
\label{fa2:p4:singular-weight-minimizer}
The Banach space $E=L^1(\mathbb R)\cap L^2(\mathbb R)$ with
sum norm is not reflexive. For $0<\alpha<1/2$, however,
\[
 \inf_{u\geq0,\ \int u\leq1}
 \left(\|u\|_2^2-\int |x|^{-\alpha}u(x)dx\right)
\]
has a unique minimizer in $E$.
\end{example}
\begin{proof}
A sum-norm Cauchy sequence has $L^1$ and $L^2$ limits; a.e.
subsequences identify them, proving completeness. The weight
is in $L^2$ near zero and bounded outside, so its pairing is
continuous in this sum norm. Write the weight as $w_2+w_\infty$;
on the feasible set its pairing is at most
$\|w_2\|_2\|u\|_2+\|w_\infty\|_\infty$. The energy is
bounded below and minimizing sequences are bounded in $L^2$.
Take a weakly convergent subsequence there. Its limit is
nonnegative, and testing $1_{[-R,R]}$ shows its integral on
that interval at most one; monotone convergence gives total
mass at most one. Thus it is feasible and in $E$.
For each fixed $R$ the truncated singular weight is in $L^2$,
so its integral converges weakly. Outside $[-R,R]$ its integral
on any feasible function is at most $R^{-\alpha}$.
Taking $R\to\infty$ proves convergence of the full linear
term. Weak lower semicontinuity of the squared $L^2$ norm
therefore yields a minimizer. Strict convexity gives uniqueness.

Finally $u_n=1_{[n,n+1]}$ are bounded in $E$ and weakly
converge to zero in $L^2$ by the tail norm of every $L^2$ test.
If $E$ were reflexive a subsequence would have an $E$ weak
limit, necessarily zero under the continuous inclusion in $L^2$.
But the continuous functional $u\mapsto\int u$ equals one
on that subsequence. This contradiction proves nonreflexivity;
the existence argument did not pretend that $E$ was reflexive.
\end{proof}

\begin{theorem}[Distribution functions and a Banach norm for weak $L^p$]
\label{fa2:p4:marcinkiewicz-space}
On a finite measure space write $\alpha_f(t)=\mu\{|f|>t\}$.
It is nonincreasing and right continuous. If $f_n\to f$ a.e.,
then for $t>0$
\[
 \alpha_f(t)\leq\liminf\alpha_{f_n}(t)
 \leq\limsup\alpha_{f_n}(t)\leq\alpha_f(t-).
\]
For $1<p<\infty$ put
\[
 Q_p(f)=\sup_{t>0}t\alpha_f(t)^{1/p},\qquad
 N_p(f)=\sup_{\mu(A)>0}\mu(A)^{-1/p'}\int_A|f|.
\]
Then $Q_p\leq N_p\leq p'Q_p$, and the a.e. quotient of
functions with finite $N_p$ is Banach in that norm.
One has $N_p(f)\leq\|f\|_p$,
$N_q(f)\leq\mu(\Omega)^{1/q-1/p}N_p(f)$ for $1<q\leq p$,
and
\[
 \|f\|_q\leq (p/(p-q))^{1/q}
            \mu(\Omega)^{1/q-1/p}Q_p(f)\quad(1\leq q<p).
\]
If $1<q<p<r<\infty$ and $1/p=\theta/q+(1-\theta)/r$,
then $\|f\|_p\leq C_{p,q,r}N_q(f)^\theta N_r(f)^{1-\theta}$.
If $d\ge1$, the inclusion of $L^p$ is strict on a Euclidean unit ball:
$|x|^{-d/p}$ belongs to weak $L^p$ there but not to $L^p$.
\end{theorem}
\begin{proof}
Right continuity is continuity of measure on the increasing sets
$\{|f|>t_n\}$ for $t_n\downarrow t$. The a.e. limit
inequalities follow from Fatou on indicators for the lower bound
and its reverse for bounded indicators on this finite measure
space for the upper bound. The latter upper set is
$\{|f|\geq t\}$, whose measure is $\alpha_f(t-)$.
There are at most countably many jumps, so distribution functions
converge at almost every positive level.
Tonelli gives the general layer formula
\[
 \int H(|f|)=\int_0^\infty \alpha_f(t)h(t)dt,
 \quad H(s)=\int_0^s h(t)dt,\quad h\geq0,
\]
and the tail formula adds the term $\alpha_f(a)H(a)$
when integration is restricted to $\{|f|>a\}$ and the
right integral starts at $a$. These hold also with value
infinity. In particular $H(s)=s^q$ gives the power formula;
if $f\in L^q$, then $t^q\alpha_f(t)\leq
\int_{|f|>t}|f|^q\to0$.

Choosing $A=\{|f|>t\}$ shows $Q_p\leq N_p$.
Conversely layer integration gives
$\int_A|f|\leq\int_0^\infty
\min(\mu(A),Q_p^pt^{-p})dt=p'Q_p\mu(A)^{1/p'}$.
The triangle inequality for $N_p$ is immediate and its
definiteness follows by testing positive-measure level sets.
The zero measure space consists of one quotient element and
is handled separately. Otherwise a sequence Cauchy for $N_p$
is Cauchy in $L^1$ by testing $A=\Omega$. Let $f$ be its
$L^1$ limit and take an a.e. convergent subsequence. Fatou on
every set $A$, followed by the supremum, gives
$N_p(f_n-f)\leq\liminf_m N_p(f_n-f_m)$. Thus convergence
holds in $N_p$, proving completeness.
Holder gives $N_p\leq\|f\|_p$, and comparing the powers
of $\mu(A)$ gives the $N_q$ bound.
For the strong $q$ estimate integrate
$\alpha_f(t)\leq\min(\mu(\Omega),Q_p^pt^{-p})$
against $qt^{q-1}$ and split at
$t=Q_p\mu(\Omega)^{-1/p}$. This gives the displayed constant.
For interpolation use
$\alpha_f(t)\leq\min((Q_q/t)^q,(Q_r/t)^r)$ and split
at $t_0=(Q_r^r/Q_q^q)^{1/(r-q)}$; direct integration gives
the powers $Q_q^\theta Q_r^{1-\theta}$ and a finite
constant because $q<p<r$. Zero quasi-norms are trivial.
On the zero-dimensional unit ball $Q_p(f)=N_p(f)=|f(0)|$ and
weak and strong $L^p$ coincide. Finally, in positive dimension, on the
unit ball the example has distribution
$\min(v_d,v_dt^{-p})$ for $t>0$, whereas the integral of
its $p$th power contains $\int_0^1dr/r=\infty$.
\end{proof}

% END INPUT

\part{Selected applications and representation bridges}

% BEGIN parts/part5.tex
\providecommand{\pvnorm}[1]{\left\lVert#1\right\rVert}
\providecommand{\pvsupp}{\operatorname{supp}}
\chapter{Integral Operators and Fredholm Equations}
\label{fa2:p5:integral-chapter}

The integral-operator models in this part retain their spaces and measures.
A formula with the same kernel does not automatically identify a map on
$C(K)$, a map on $L^2$, and a map on a space of measures.
Banach transposes are linear and have no conjugation in their defining
pairing. Hilbert adjoints use the conjugate-first, linear-second product.

\section{Boundedness, compactness, and exact model bridges}

\begin{lemma}[Arzel\`a--Ascoli on a compact Hausdorff space]
\label{fa2:p5:ascoli}
A uniformly bounded equicontinuous subset of $C(X)$, $X$ compact
Hausdorff, is relatively compact in the uniform norm. Conversely,
a relatively compact family has these properties.
\end{lemma}
\begin{proof}
Given $\varepsilon>0$, equicontinuity gives a finite open cover
$U_1,\ldots,U_m$ and points $x_j\in U_j$ such that
$|f(x)-f(x_j)|<\varepsilon$ for all $x\in U_j$ and all members
of the family. A bounded subset of $\mathbb K^m$ has a finite
$\varepsilon$-net for the value tuples $(f(x_j))$.
Choosing one family member for each occupied net cell gives a
$3\varepsilon$-net in uniform norm. Completeness of $C(X)$
turns total boundedness into relative compactness. Conversely,
choose a finite uniform $\varepsilon$-net of continuous functions;
their finite intersection of continuity neighborhoods proves
equicontinuity, and their finitely many norms prove boundedness.
The proof does not require that $X$ be metrizable.
\end{proof}

\begin{theorem}[Continuous kernels and their transpose]
\label{fa2:p5:continuous-kernel}
Let $X,Y$ be compact Hausdorff, $\nu$ a finite positive Radon
measure on $Y$, and $k\in C(X\times Y)$.
Then $Ku(x)=\int_Yk(x,y)u(y)\,d\nu(y)$ is a bounded compact map
$C(Y)\to C(X)$, with
$\pvnorm K\le\nu(Y)\pvnorm k_\infty$.
Under the linear Riesz--Markov identification,
$K'\eta=h\nu$, where $h(y)=\int_Xk(x,y)\,d\eta(x)$.
No conjugation occurs in this formula.
\end{theorem}
\begin{proof}
Continuity of the map $x\mapsto k(x,\cdot)$ into $C(Y)$ follows
from continuity on the compact product: finitely many neighborhoods
cover $Y$ to make the oscillation uniform in $y$.
The integral estimate gives boundedness, and the same observation
makes the image of the unit ball equicontinuous. Apply Ascoli.
For the transpose, interchange the two continuous-kernel integrals.
On arbitrary compact Hausdorff spaces this can be justified without
assuming that the Borel algebra of the product is the product
$\sigma$-algebra: finite sums $a_j(x)b_j(y)$ uniformly approximate
$k$ by Stone--Weierstrass, and for these sums the interchange is
an algebraic identity. Both iterated integrals are uniform-norm
continuous, so it passes to $k$. The resulting identity against
every $u\in C(Y)$ and Riesz--Markov uniqueness give the formula.
\end{proof}

\begin{theorem}[$L^2$ kernels and the continuous-to-$L^2$ bridge]
\label{fa2:p5:l2-kernel}
On $\sigma$-finite measure spaces $(X,\mu),(Y,\nu)$, a kernel
$k\in L^2(\mu\otimes\nu)$ defines a.e. a bounded compact map
$K:L^2(\nu)\to L^2(\mu)$, with
$\pvnorm K\le\pvnorm k_2$ and
\[
 K^*v(y)=\int_X\overline{k(x,y)}v(x)\,d\mu(x).
\]
It is Hilbert--Schmidt with squared Hilbert--Schmidt norm
$\int|k|^2$. For compact $X,Y$ with finite Radon measures and
continuous $k$, the two definitions agree on continuous functions;
their $L^2$ extensions agree by density. A Hermitian kernel on
one space gives a compact self-adjoint map.
\end{theorem}
\begin{proof}
Fubini gives $k(x,\cdot)\in L^2$ for almost every $x$.
Cauchy--Schwarz and integration in $x$ give the norm bound.
Finite sums $\sum a_j(x)b_j(y)$ are dense in the product $L^2$:
first on finite-measure rectangles the sets whose indicators are
in their closed span form a monotone class containing measurable
rectangles, and then a countable exhaustion and simple density
give the general case. Their operators are finite rank, and the
norm estimate proves compactness. The adjoint identity follows
by absolutely integrable Fubini, since
$\iint|k(x,y)u(y)v(x)|\le\pvnorm k_2\pvnorm u_2\pvnorm v_2$.
Choose finite-product approximants to $k$ and the separable closed
span $H_0$ of the complex conjugates of their $Y$-factors.
A subsequence converges in section $L^2$ for a.e. $x$, so
$\overline{k(x,\cdot)}\in H_0$ for a.e. $x$.
Since $Ku(x)=\langle\overline{k(x,\cdot)},u\rangle$ in our
conjugate-first convention, for an orthonormal basis of $H_0$
Parseval and Tonelli give
$\sum_j\pvnorm{Ke_j}_2^2=\int|k|^2$; the operator vanishes
on the perpendicular complement. Basis independence is the
Hilbert--Schmidt theorem of Part II.
For the continuous bridge, regularity approximates indicators in
$L^2$ by continuous cutoffs, proving $C(X)$ dense. The integral
formulas agree pointwise on continuous functions and hence in
$L^2$; boundedness gives the extension identity. Hermitian symmetry
is exactly the adjoint formula, not complex symmetry without a bar.
\end{proof}

\begin{theorem}[Schur's weighted kernel test]
\label{fa2:p5:schur}
On $\sigma$-finite measure spaces, suppose a product-measurable kernel
and positive measurable weights $a(x),b(y)$
satisfy
\[
 \int|k(x,y)|b(y)\,d\nu(y)\le A a(x),\qquad
 \int|k(x,y)|a(x)\,d\mu(x)\le B b(y).
\]
Then $K:L^2(\nu)\to L^2(\mu)$ is bounded with
$\pvnorm K\le\sqrt{AB}$, whenever the weights are positive and
finite a.e. The unweighted row and column bounds also give
$\pvnorm K_{L^p\to L^p}\le A^{1-1/p}B^{1/p}$ for
$1\le p\le\infty$.
\end{theorem}
\begin{proof}
Weighted Cauchy--Schwarz gives
$|Ku(x)|^2\le Aa(x)\int|k(x,y)||u(y)|^2/b(y)\,d\nu(y)$.
Tonelli and the second bound give $AB\pvnorm u_2^2$.
Construct the operator first on bounded simple functions supported
where the weights are bounded above and below, then extend by
density; the inequality also proves a.e. existence of the intended
integral for the extension. Unweighted column and row bounds give
the $L^1$ and $L^\infty$ estimates by Tonelli and the triangle
inequality. Riesz--Thorin supplies the displayed intermediate bound.
Boundedness does not imply compactness: translation convolution
on $\mathbb R$ gives many noncompact Schur-bounded examples.
\end{proof}

\begin{theorem}[Weakly singular kernels]
\label{fa2:p5:weakly-singular-kernel}
Let $G\subset\mathbb R^d$ be compact and let $k$ be continuous off
the diagonal with $|k(x,y)|\le C|x-y|^{\alpha-d}$ for some
$\alpha>0$ (a bounded kernel continuous off the diagonal may be added).
The resulting integral map is compact on $C(G)$ and on every $L^p(G)$,
$1\le p\le\infty$. The same assertion holds on a compact $n$-dimensional
$C^1$ surface with surface measure and singularity $|x-y|^{\alpha-n}$.
\end{theorem}
\begin{proof}
For $d\ge1$ and $0<\delta<1$, polar integration bounds the
singular component's row and column integrals by $C'\delta^\alpha$.
An added bounded kernel, still continuous off the diagonal,
contributes at most $C''\delta^d$, so the full discarded strip
has row and column bounds $O(\delta^\alpha+\delta^d)$.
For $d=0$, a compact subset of $\mathbb R^0$ is empty or a singleton,
and its integral map is finite rank directly. Cut the kernel off continuously near
the diagonal. The remaining continuous kernel is compact on $C(G)$.
It is also compact on every $L^p$: uniform approximation by finite
tensor sums gives operator-norm approximation by finite-rank maps,
including at $p=1,\infty$, by the finite-measure integral estimate.
Schur's endpoint/interpolated estimate makes the discarded strip
small in every $L^p$ operator norm; its row bound does so on $C(G)$.
Norm limits are compact on these complete targets.
On a surface of positive dimension $n$, finitely many $C^1$
graphs give $\mu(B(x,r))\le C r^n$ uniformly. Dyadic annuli
bound the singular component by $C'\delta^\alpha$, while an
added bounded off-diagonal-continuous kernel contributes
$O(\delta^n)$. Thus the full strip tends to zero, and the rest
of the proof is identical. A compact zero-dimensional $C^1$
surface is finite, so its integral maps are finite rank directly.
This proves the geometric bridge used for the compactness assertion
in \cite{B06}, 2.29--2.30, while leaving actual boundary-value
jump relations to their separate PDE theory.
\end{proof}

\section{Fredholm reduction and ill-posedness}

\begin{proposition}[Exact finite-rank reduction]
\label{fa2:p5:finite-kernel}
For $Ku=\sum_{j=1}^ma_j\ell_j(u)$ on a Banach space,
$u-\lambda Ku=f$ is equivalent to
\[
 c_i-\lambda\sum_j\ell_i(a_j)c_j=\ell_i(f),\qquad
 u=f+\lambda\sum_ja_jc_j.
\]
No independence of the $a_j$ or $\ell_j$ is required.
\end{proposition}
\begin{proof}
Apply $\ell_i$ to the original equation and put $c_i=\ell_i(u)$.
Conversely the reconstructed $u$ satisfies $\ell_i(u)=c_i$ by
the matrix equations and hence the original equation. This proves
both directions and preserves the full affine solution sets, even
when the chosen representation is redundant.
\end{proof}

\begin{theorem}[Fredholm kernel equations]
\label{fa2:p5:fredholm-kernel}
For any of the compact kernel maps above acting on the same Banach
space, $u-\lambda Ku=f$, $\lambda\ne0$, is uniquely solvable for
all data if and only if $\lambda^{-1}$ is not an eigenvalue of $K$.
Otherwise it is solvable exactly when every member of
$\ker(I-\lambda K)'$ annihilates $f$; the solution set is a
translate of the finite-dimensional homogeneous kernel.
In Hilbert space the condition is
$f\perp\ker(I-\overline\lambda K^*)$.
\end{theorem}
\begin{proof}
Apply the Riesz--Schauder/Fredholm theorem of Part II to
$I-\lambda K$. Its nullspace is precisely the eigenspace at
$\lambda^{-1}$, its range is closed, its index is zero, and
its range equals the preannihilator of the transpose kernel.
All kernel hypotheses have been verified above; thus this is an
application with an established bridge rather than a substitution
of an unproved compactness claim. The complex Hilbert form follows
from $(I-\lambda K)^*=I-\overline\lambda K^*$.
The relevant nullspace is in general the adjoint nullspace, not
the original one.
\end{proof}

\begin{theorem}[Picard criterion and instability]
\label{fa2:p5:picard}
Let $K:H\to G$ be compact with positive singular system
$(\sigma_j,e_j,f_j)$. Then $Ku=g$ is solvable exactly when
\[
 g\perp\ker K^*,\qquad
 \sum_j\frac{|\langle f_j,g\rangle|^2}{\sigma_j^2}<\infty.
\]
The series $u^\dagger=\sum_j\sigma_j^{-1}\langle f_j,g\rangle e_j$
is the unique minimum-norm solution, and all solutions are
$u^\dagger+\ker K$. An infinite-rank compact operator cannot
have a bounded inverse on its range when that range carries the
ambient norm.
\end{theorem}
\begin{proof}
Pairing $Ku=g$ with $f_j$ gives the coefficients of $u$ along
$e_j$. Bessel proves necessity of summability, and the adjoint
kernel identity gives the orthogonality condition. Conversely the
displayed series converges in $H$; applying bounded $K$ gives
the expansion of $g$ on $\overline{\operatorname{ran}K}$,
which is all of $g$ by that orthogonality. Orthogonal decomposition
along $\ker K$ proves minimality and uniqueness.
If rank is infinite, $\sigma_j\to0$; unit vectors $e_j$ have
images of norm $\sigma_j\to0$, so no bounded inverse estimate
is possible. For the explicit diagonal map $K(x_j)=(x_j/j)$
on $\ell^2$, $(1/j)_j\in\ell^2$ is not in its range, whereas
$K(e_j)=e_j/j\to0$. This separates existence from continuous
dependence with actual data.
\end{proof}

\chapter{Volterra Operators: Regular, Singular, and Noncompact Models}
\label{fa2:p5:volterra-chapter}

\section{Continuous kernels and comparison}

\begin{theorem}[Factorial bounds and full-interval resolvent]
\label{fa2:p5:volterra-resolvent}
Let $I=[0,T]$, $E$ Banach, and $k(t,s)\in\mathcal L(E)$ be
operator-norm continuous on $0\le s\le t\le T$, with bound $M$.
For $Vu(t)=\int_0^tk(t,s)u(s)\,ds$,
\[
 \pvnorm{V^n}\le M^nT^n/n!,\qquad
 (I-\lambda V)^{-1}=\sum_{n=0}^\infty\lambda^nV^n,
 \qquad \pvnorm{(I-\lambda V)^{-1}}\le e^{|\lambda|MT}.
\]
For every parameter and continuous datum the equation has a unique
solution on the entire finite interval. In a nonzero complex
state space its spectrum is $\{0\}$ when $T>0$; the zero interval
is treated as the zero integral operator. Iterated kernels give a
continuous resolvent kernel. Picard iterates from any initial
continuous function converge uniformly to the solution.
\end{theorem}
\begin{proof}
The $n$fold integral is over
$0\le s_n\le\cdots\le s_1\le t$, with the ordered operator
product $k(t,s_1)\cdots k(s_{n-1},s_n)$.
Its norm is bounded by $M^n$ times the simplex volume $t^n/n!$.
The exponential majorant proves uniform operator-norm convergence;
multiplying partial sums on either side by $I-\lambda V$ leaves
a remainder tending to zero, proving a two-sided inverse.
For the kernel series, its $n$th term is bounded by
$M^n(t-s)^{n-1}/(n-1)!$; uniform convergence proves continuity
and the resolvent identity by Fubini. Iteration error after $n$
steps is $(\lambda V)^n$ applied to the initial error, hence tends
to zero by the same estimate. Since all nonzero scalars lie in
the resolvent, the complex spectrum is contained in $\{0\}$;
zero is present because $(Vu)(0)=0$ makes $V$ nonsurjective on
$C(I,E)$ when $E\ne0$. This argument needs no smallness of $MT$.
\end{proof}

\begin{lemma}[Gronwall and a comparison principle]
\label{fa2:p5:gronwall}
If $a,M\ge0$ and a nonnegative continuous $z$ satisfies
$z(t)\le a+M\int_0^tz(s)ds$, then $z(t)\le ae^{Mt}$.
For nonnegative continuous Volterra kernels and $\lambda\ge0$,
successive iteration preserves order and gives the corresponding
resolvent comparison for $I-\lambda V$.
\end{lemma}
\begin{proof}
Iterate the inequality $n$ times. The data terms sum to
$a\sum_{j=0}^{n-1}(Mt)^j/j!$, while the remaining term is at
most $\pvnorm z_\infty(Mt)^n/n!$. Pass to the limit.
For a nonnegative kernel all iterated kernels are nonnegative;
applying its convergent resolvent series to an inequality therefore
preserves that inequality. For systems use componentwise order
and nonnegative matrix entries; scalar/Hilbert quadratic positivity
is a different notion and is not substituted here.
\end{proof}

\begin{proposition}[Smoothness, first-kind equations, and ODEs]
\label{fa2:p5:volterra-regularity}\label{fa2:p5:first-kind-smooth}
For a scalar kernel $k\in C^m$ on a neighborhood of the triangle
and $f\in C^m(I)$, the second-kind solution is $C^m$.
In the $C^1$ case,
$u'=f'+\lambda k(t,t)u+\lambda\int_0^t\partial_tk(t,s)u(s)ds$.
If $T>0$, $k\in C^1$ and $k(t,t)\ne0$ everywhere, the first-kind
equation $\int_0^tk(t,s)u(s)ds=g(t)$ has a continuous solution
exactly when $g(0)=0$, for $g\in C^1$; that solution is unique.
If $T=0$, it has a solution exactly when $g(0)=0$, and then every
scalar-valued function on $\{0\}$ is a solution, so uniqueness fails.
For operator-norm continuous $A:I\to\mathcal L(E)$ and continuous
$b:I\to E$, the equation
$u(t)=u_0+\int_0^t(A(s)u(s)+b(s))ds$ is equivalent to
$u'=Au+b$, $u(0)=u_0$, and is uniquely solvable on $I$.
\end{proposition}
\begin{proof}
The difference quotient of the triangular integral splits into a
short endpoint integral and an integral of kernel differences.
Uniform continuity of $\partial_tk$ and the norm integral estimate
give the displayed formula. It first makes $u$ continuously
differentiable. Differentiating again produces only continuous
diagonal derivatives of $k$ and already obtained derivatives of
$u$; induction proves $C^m$. The same induction applied to the
kernel resolvent identities proves its $C^m$ regularity.
For $T>0$, differentiating the first-kind equation gives
$u(t)=g'(t)/k(t,t)-\int_0^t\partial_tk(t,s)u(s)/k(t,t)ds$,
a continuous second-kind equation. Its unique solution integrates
back to the original equation exactly when $g(0)=0$. For $T=0$
the integral is zero and the equation is simply $0=g(0)$; the
value $u(0)$ is unrestricted, proving the separate endpoint case.
For the ODE take $k(t,s)=A(s)$ and apply the Banach-valued
resolvent theorem. The fundamental theorem for continuous Bochner
integrals differentiates the equation; integrating a classical
solution proves the converse. This proof states the domain and
regularity rather than identifying a formal derivative with an
unbounded operator realization.
\end{proof}

\section{Fractional and logarithmic kernels}

\begin{definition}[Fractional integrals]
\label{fa2:p5:fractional-definition}
For $\alpha>0$ put
$J_\alpha f(t)=\Gamma(\alpha)^{-1}\int_0^t(t-s)^{\alpha-1}f(s)ds$.
This normalization is explicit. Brunner's power-kernel operator
has no Gamma factor, and his parameter $\alpha=0$ denotes a
separate logarithmic kernel, not the nonintegrable power $r^{-1}$.
\end{definition}

\begin{theorem}[Beta convolution and weakly singular Volterra]
\label{fa2:p5:weak-volterra}
$J_\alpha J_\beta=J_{\alpha+\beta}$ on continuous functions and
on $L^p(0,T)$, $1\le p\le\infty$.
If $|k(t,s)|\le M$ and $k$ is continuous on the triangle, then
$V_\alpha f=\int_0^t(t-s)^{\alpha-1}k(t,s)f(s)ds$ is compact
on $C(I)$ and on each $L^p(0,T)$, and
\[
 \pvnorm{V_\alpha^n}\le
 \frac{M^n\Gamma(\alpha)^nT^{n\alpha}}{\Gamma(n\alpha+1)}.
\]
It is quasinilpotent and $I-\lambda V_\alpha$ is invertible for
every $\lambda$. The resolvent kernel has the same leading weak
singularity. These conclusions also hold for
$\int_0^t\log(t-s)k(t,s)f(s)ds$.
For constant $k=1$ and constant datum the normalized equation
$u=1+\lambda J_\alpha u$ has solution
$E_\alpha(\lambda t^\alpha)=\sum_{n\ge0}\lambda^nt^{n\alpha}/\Gamma(n\alpha+1)$.
\end{theorem}
\begin{proof}
Fubini and $s=v+(t-v)r$ give
$\int_v^t(t-s)^{\alpha-1}(s-v)^{\beta-1}ds
=B(\alpha,\beta)(t-v)^{\alpha+\beta-1}$.
The Gamma--Beta identity follows by changing variables in the
positive double Gamma integral, so the normalizations cancel.
Absolute Fubini is justified by this positive majorant and Young's
$L^1*L^p$ bound. Repeating the positive majorant gives the stated
iterate estimate, both by row/column bounds for $L^p$ and pointwise
on $C$. The Gamma denominator grows faster than any geometric
sequence: integrating its defining integral over
$[n\alpha,n\alpha+1]$ already gives an $n$th root bounded below
by a positive multiple of $n^\alpha$.
Thus the Neumann series converges in operator norm for every
parameter and proves quasinilpotence. A continuous cutoff of the
strip $t-s<\delta$ leaves a continuous kernel, and both discarded
row and column integrals are $O(\delta^\alpha)$; Schur proves
compactness as in the earlier kernel theorem. The iterated kernel
bound is
$M^n\Gamma(\alpha)^n(t-s)^{n\alpha-1}/\Gamma(n\alpha)$.
After dividing its resolvent series by $(t-s)^{\alpha-1}$ the
series converges uniformly; Beta rescaling proves continuity of
each amplitude. This proves the precise singularity assertion.
For the logarithmic kernel choose $0<\varepsilon<1$;
$|\log r|\le C_{\varepsilon,T}r^{-\varepsilon}$ on $(0,T]$.
The same majorant with $\alpha=1-\varepsilon$ proves invertibility
and compactness; its strip integral also tends to zero.
For the precise logarithmic leading term choose $0<\varepsilon<1/2$.
The terms of order at least two in the resolvent kernel series
have the uniform bound $O(r^{1-2\varepsilon})$, $r=t-s\downarrow0$:
factor $r^{2\alpha-1}$ from the Gamma majorant series starting
at $n=2$, whose remaining series is bounded for $0<r\le1$.
Thus the kernel of $(I-\lambda V)^{-1}-I$ is
$\lambda k(t,s)\log r+O(r^{1-2\varepsilon})$ locally at the
diagonal. This proves, separately, the logarithmic leading-singularity
claim; a power majorant alone would not identify that term.
Finally apply the semigroup
identity to the constant function to obtain the displayed series.
\end{proof}

\begin{theorem}[Abel inversion and the weakly singular first kind]
\label{fa2:p5:abel}
Let $T>0$. For $0<\alpha<1$, $J_\alpha u=g$ on $(0,T]$ (equivalently a.e.)
with $g\in C^1[0,T]$ has a unique $L^1$ solution,
\[
 u(t)=\frac{g(0)t^{-\alpha}}{\Gamma(1-\alpha)}+J_{1-\alpha}g'(t).
\]
It is continuous at zero exactly when $g(0)=0$; in that case its
value there is zero and the equation holds also at zero. If
$g(0)\ne0$, the integral convention at zero gives zero and is
not claimed to equal the nonzero datum there.
For a $C^1$ amplitude $k$ with nonvanishing
diagonal, the equation
$\int_0^t(t-s)^{\alpha-1}k(t,s)u(s)ds=g(t)$ with
$g\in C^1$, $g(0)=0$, has a unique continuous solution.
For $T=0$, the equation on $(0,T]$ is vacuous and $L^1[0,T]$
has only its zero equivalence class. For continuous functions,
either first-kind equation imposed also at zero is solvable
exactly when $g(0)=0$; every singleton value is then a solution.
The positive-length continuity and endpoint-reconstruction
claims are not asserted for that degenerate interval.
\end{theorem}
\begin{proof}
Apply $J_{1-\alpha}$ and differentiate $J_1u$.
Writing $g=g(0)+J_1g'$ and using the semigroup identity yields
the formula. Applying $J_\alpha$ back proves sufficiency, including
the endpoint constant since
$J_\alpha(t^{-\alpha}/\Gamma(1-\alpha))=1$ for $t>0$.
Uniqueness follows from injectivity of $J_\beta$ on $L^1$:
compose with $J_{1-\beta}$ for $0<\beta<1$ and differentiate.
The stated continuity follows from $J_{1-\alpha}g'=O(t^{1-\alpha})$.

For the amplitude equation, applying $J_{1-\alpha}$ gives
$\int_0^tG(t,s)u(s)ds=J_{1-\alpha}g(t)$, where
\[
 G(t,s)=\frac1{\Gamma(1-\alpha)}
 \int_0^1(1-r)^{-\alpha}r^{\alpha-1}
                  k(s+(t-s)r,s)\,dr.
\]
Both $G$ and $\partial_tG$ are continuous by the integrable Beta
weight, and $G(t,t)=\Gamma(\alpha)k(t,t)$.
Differentiation therefore gives a continuous second-kind equation
with datum $J_{1-\alpha}g'$ and nonzero diagonal coefficient.
Its unique continuous solution solves the original equation,
because the differentiated identity has zero integration constant
and $J_{1-\alpha}$ is injective. This proves the general amplitude
version rather than silently using translation-invariant inversion.
When $T=0$, the integration interval has measure zero, proving
the asserted $L^1$ statement. On singleton continuous functions
both original integrals are zero and the value of the unknown
is unrestricted; this proves the separate endpoint statement.
\end{proof}

\begin{theorem}[A valid H\"older gain and the source counterexample]
\label{fa2:p5:holder-repair}
Let $0<\alpha,\beta<1$ and $\gamma=\alpha+\beta<1$.
If $f\in C^\beta[0,T]$ and $f(0)=0$, then
$J_\alpha f\in C^\gamma[0,T]$.
The same conclusion holds for an amplitude $k\in C^1$ on a
neighborhood of the closed triangle. Mere continuity of $k$ is
insufficient.
\end{theorem}
\begin{proof}
In the increment from $t-h$ to $t$, subtract $f(t)$ inside both
integrals. The short interval contributes at most $C h^\gamma$.
The remaining integral, after putting $r=t-s$ and scaling by $h$,
is bounded by
\[
 C h^\gamma\int_1^{t/h}
       |x^{\alpha-1}-(x-1)^{\alpha-1}|x^\beta\,dx.
\]
This integral is uniformly bounded: the singularity at one is
integrable, and at infinity the integrand is
$O(x^{\gamma-2})$. The subtracted constant contributes
$|f(t)|\,|t^\alpha-(t-h)^\alpha|/\alpha$.
Since $|f(t)|\le Ct^\beta$, this is $O(h^\gamma)$ when $t\le2h$;
when $t>2h$ use the mean-value bound
$Ch t^{\gamma-1}\le Ch^\gamma$. This proves the constant-amplitude
claim, including the endpoint zero.

For a $C^1$ amplitude write the unnormalized operator as
$k(t,t)\Gamma(\alpha)J_\alpha f+R(t)$, with
\[
 R(t)=\int_0^t(t-s)^{\alpha-1}[k(t,s)-k(t,t)]f(s)ds.
\]
The bracket is $O(t-s)$. Its differentiated integrand is
\[
 (\alpha-1)(t-s)^{\alpha-2}[k(t,s)-k(t,t)]f(s)
 +(t-s)^{\alpha-1}
 [\partial_1k(t,s)-(\partial_1+\partial_2)k(t,t)]f(s).
\]
It is bounded by $C(t-s)^{\alpha-1}\pvnorm f_\infty$.
Splitting off a strip near $s=t$ justifies differentiation and
continuity of the derivative by dominated convergence away from
that strip; its integral is uniformly $O(\delta^\alpha)$.
The endpoint term is zero because the undifferentiated integrand
is $O((t-s)^\alpha)$. Thus $R\in C^1$, and multiplication by the
$C^1$ diagonal function preserves $C^\gamma$. This is an actually
proved sufficient condition, not a guessed correction to the source.

For the counterexample to \cite{B07}, Theorem 8.1.5, take
$T=1$, $t_0=1/2$, $\alpha=\beta=1/4$,
$f(s)=s$, and $k(t,s)=|t-1/2|^{1/4}$.
The printed hypotheses $k\in C$, $f\in C^\beta$, $f(0)=0$ hold,
but the unnormalized source operator equals
\[
 \int_0^t(t-s)^{-3/4}k(t,s)s\,ds
 =\frac{16}{5}t^{5/4}|t-1/2|^{1/4}.
\]
At $t_0$ its order-$1/2$ H\"older quotient is
$(16/5)(t_0+h)^{5/4}|h|^{-1/4}\to\infty$.
The source's surrounding pages impose no stronger global amplitude
assumption. This flags that particular printed statement; the earlier
continuous-kernel $C^1$ theorem is unaffected.
\end{proof}

\begin{example}[Weak singularity does not preserve smoothness]
\label{fa2:p5:fractional-regularity-example}
For smooth constant data, $u=1+\lambda J_\alpha u$ has
$u(t)=1+\lambda t^\alpha/\Gamma(1+\alpha)+O(t^{2\alpha})$.
For $T>0$, $0<\alpha<1$ and $\lambda\ne0$ it is not $C^1$ at zero.
For $T=0$ the solution is the singleton value $u(0)=1$, which
has a smooth constant extension.
Its derivatives of order $j$ on $(0,T]$ satisfy the weighted bound
$|u^{(j)}(t)|\le C_jt^{\alpha-j}$ near zero, after subtracting the
constant term, by termwise differentiation of the convergent
Mittag--Leffler series. Thus smooth-data regularity for bounded
kernels cannot be transferred to singular kernels.
\end{example}

\section{Nonlinear, delayed, and noncompact examples}

\begin{theorem}[Local nonlinear existence and a continuation criterion]
\label{fa2:p5:nonlinear-volterra}
For Banach $E$, suppose $f$ is continuous and $G(t,s,z)$ is
continuous on a triangular time cylinder and a closed ball about
$f(0)$, bounded there by $B$ and uniformly Lipschitz in $z$ with
constant $L$. On an interval $[0,h]$ on which
$\pvnorm{f-f(0)}_\infty+Bh\le R$, the equation
$u(t)=f(t)+\int_0^tG(t,s,u(s))ds$ has a unique solution in that
ball. Picard iterates converge uniformly. For weakly singular
factor $(t-s)^{\alpha-1}$ replace $Bh$ by $Bh^\alpha/\alpha$;
the same conclusions hold.
If these hypotheses hold on every bounded $z$-ball, with uniform
time continuity there, a maximal solution that stops before a
finite available endpoint has unbounded norm. In finite-dimensional
$E$ the required uniform time continuity follows from continuity
on compact cylinders.
\end{theorem}
\begin{proof}
Picard iteration remains in the indicated ball by its a priori
bound. Successive differences satisfy
$d_{n+1}(t)\le L\int_0^td_n(s)ds$, hence a factorial majorant;
the first difference is bounded. Summability proves uniform
convergence, and dominated convergence passes to the equation.
Apply the same iterate estimate to the difference of two solutions
to prove uniqueness without a further smallness condition on $Lh$.
In the singular case Beta iteration replaces $n!$ by
$\Gamma(n\alpha+1)$ and gives the same proof.

For continuation suppose the norm stays bounded up to $t_*<T$.
The short part of an increment is bounded by $B|h|$, and the
history part by $T$ times the uniform time modulus of $G$ on
that ball; hence $u$ has a norm limit at $t_*$.
The history integral becomes part of a continuous new datum,
and the local theorem restarts beyond $t_*$. For the singular
case its integrable kernel has translation-continuous $L^1$
increments; the same split works. This contradicts maximality.
The uniformity assumption for a noncompact Banach ball has been
stated explicitly instead of being inferred from ordinary continuity.
\end{proof}

\begin{example}[Existence without uniqueness, blow-up, and a third kind]
\label{fa2:p5:nonunique-blowup-third}
The equation $u(t)=\int_0^t2\sqrt{|u(s)|}\,ds$ has the solutions
$u_c(t)=0$ for $t\le c$ and $u_c(t)=(t-c)^2$ for $t\ge c$,
for any $c\ge0$. Direct substitution proves this and displays the
failure of Lipschitz uniqueness. The equation
$u(t)=1+\int_0^tu(s)^2ds$ has $u(t)=(1-t)^{-1}$ on $t<1$;
the ODE equivalence and uniqueness show that no bounded continuation
through one exists. For the third-kind equation
$t u(t)=\int_0^tu(s)ds$, every constant is a continuous solution.
For $tu(t)=t+\int_0^tu(s)ds$ on $[0,T]$ with $T>0$, no
continuous solution exists: divide by $t$ and let $t\downarrow0$,
obtaining $u(0)=1+u(0)$. On $[0,0]$ every scalar singleton
function is a continuous solution of this equation.
The vanishing leading coefficient is therefore a substantive change,
not a second-kind equation with a harmless parameter value.
\end{example}

\begin{proposition}[One complete delay model]
\label{fa2:p5:delay-model}
Let $\tau>0$, prescribe continuous history $\eta$ on $[-\tau,0]$,
and let $A,B$ be continuous scalar or bounded-operator coefficients.
Then
\[
 u(t)=\eta(0)+\int_0^t[A(s)u(s)+B(s)u(s-\tau)]ds,
 \qquad u=\eta\text{ on }[-\tau,0],
\]
has a unique continuous solution on every finite time interval and
is $C^1$ for positive times. At zero its derivative agrees with a
differentiable history exactly when the initial compatibility
condition holds; higher derivatives can acquire step-junction
discontinuities if the corresponding history compatibilities fail.
\end{proposition}
\begin{proof}
On $[0,\tau]$ the delayed term is the known continuous function
$B(s)\eta(s-\tau)$; solve the ordinary linear Volterra equation.
Having solved through $n\tau$, the delayed term on the next
interval is known and continuous. The same theorem with the
already determined endpoint value solves that interval uniquely.
Induction proves existence and uniqueness on each finite interval;
the integral formula differentiates to the displayed delay equation.
At a positive step junction the right and left first-derivative
formulas use the same continuous values and therefore agree.
At zero the left derivative is instead prescribed by the history.
Differentiating further, when coefficients and history permit it,
propagates those higher compatibility requirements by the same
step argument. Thus this model explicitly verifies the functional-
Volterra bridge while leaving state-dependent delays to their
separate theory.
\end{proof}

\begin{theorem}[Continuous nonlinear existence and extremal solutions]
\label{fa2:p5:hammerstein-extremal}
Let $G$ be continuous on the scalar compact triangular cylinder
$0\le s\le t\le T$, $|z|\le R$, and $f\in C[0,T]$.
If $\pfnorm f_\infty+T\sup|G|\le R$, the nonlinear
Volterra map has a fixed point in that ball even without a
Lipschitz assumption. The same conclusion holds with factor
$(t-s)^{\alpha-1}$, $\alpha>0$, replacing $T$ by $T^\alpha/\alpha$.
For real data, suppose instead that $G(t,s,z)$ is nondecreasing
in $z$ and continuous between a continuous subsolution $\ell$
and supersolution $v$, with $\ell\le v$. Then there are least
and greatest solutions in this order interval; successive iteration
from $\ell$ and $v$ converges uniformly to them.
\end{theorem}
\begin{proof}
On the closed convex ball the integral map is continuous by uniform
continuity on the compact cylinder. Its image is uniformly bounded.
The difference at two nearby times splits into a short integral and
an integral of differences in $t$, which tend uniformly to zero.
In the weakly singular case first discard a strip of width $\delta$
next to the diagonal, whose uniform integral bound is
$C\delta^\alpha$, and use uniform continuity on the remaining
compact triangle. Thus the image is equicontinuous, and Ascoli
makes the map compact. The stated bound makes the ball invariant.
The proved locally convex Schauder theorem
(Theorem~\ref{fa2:p1:schauder}, or its compact-map consequence
obtained by taking the closed convex hull of the image) gives a
fixed point. The hull is compact here: finite nets for the
precompact image make its convex hull totally bounded, and the
Banach space is complete.

For the order assertion the map preserves order and the interval.
Iterates of $\ell$ increase and iterates of $v$ decrease. The
same compactness bounds hold throughout this interval, so every
sequence of iterates has a uniformly convergent subsequence.
Monotonicity makes all subsequential limits the pointwise supremum
or infimum respectively; precompactness therefore gives full
uniform convergence. Continuity passes the limit through the map.
Every fixed point in the interval bounds the lower iterates and
is bounded by the upper iterates, proving the extremal assertions.
For $u=\int_0^t2\sqrt{\max(u(s),0)}ds$ on $[0,T]$,
$\ell=0$, $v=t^2$ give the least and greatest solutions $0,t^2$;
the waiting-time family in the preceding example shows the
intermediate multiplicity. This supplies an actual non-Lipschitz
existence and comparison mechanism, not only a uniqueness failure.
\end{proof}

\begin{theorem}[Fractional comparison and weighted endpoint derivatives]
\label{fa2:p5:fractional-gronwall}
\label{fa2:p5:weighted-volterra-regularity}
If $a,b\ge0$ and $z\ge0$ is continuous with
$z\le a+bJ_\alpha z$, then
$z(t)\le aE_\alpha(bt^\alpha)$.
Let $0<\alpha<1$, $k\in C^m$ near the closed triangle and
$f\in C^m[0,T]$. The solution of
$u=f+\lambda\int_0^t(t-s)^{\alpha-1}k(t,s)u(s)ds$
is $C^m$ for positive time and satisfies, for $1\le j\le m$,
\[
 |(u-f)^{(j)}(t)|\le C_jt^{\alpha-j}\quad(0<t\le T).
\]
For $m\ge1$ its precise first endpoint term is
\[
 u(t)=f(t)+\frac{\lambda k(0,0)f(0)}{\alpha}t^\alpha
                                      +O(t^{2\alpha}).
\]
For $T>0$, a nonzero displayed coefficient obstructs $C^1$
regularity. For $T=0$ the solution is the singleton value
$u(0)=f(0)$ and has a smooth constant extension, irrespective
of that coefficient. If $f(0)=0$ and $m\ge1$, the solution is $C^1$
through zero and $u'(0)=f'(0)$.
\end{theorem}
\begin{proof}
Iterating the comparison inequality gives the first $n$ terms
of $a\sum_{j\ge0}b^jt^{j\alpha}/\Gamma(j\alpha+1)$,
plus a nonnegative remainder bounded by
$\pfnorm z_\infty b^nt^{n\alpha}/\Gamma(n\alpha+1)$.
The Gamma estimate already proved makes the remainder tend to
zero, giving the full fractional Gronwall bound.

For derivative regularity we justify differentiation of the entire
resolvent series, including its endpoint weights. Set $r_0=1$.
Changing all $n$ ordered variables to $s_j=tr_j$ gives
\[
 V_\alpha^nf(t)=t^{n\alpha}
 \int_{0<r_n<\cdots<r_1<1}
 \prod_{j=1}^n(r_{j-1}-r_j)^{\alpha-1}
 \prod_{j=1}^nk(tr_{j-1},tr_j)f(tr_n)\,dr.
\]
The integral of the first product is
$\Gamma(\alpha)^n/\Gamma(n\alpha+1)$, by repeated Beta
integration. The product rule for $j$ derivatives has at most
$C_j(1+n)^j$ terms, each bounded by
$L_j^{n+1}$ for one $L_j\ge1$ dominating the $C^j$ norms
of $k,f$. Differentiating the factor $t^{n\alpha}$ as well
therefore gives, for $0<t\le T$,
\[
 |(V_\alpha^nf)^{(j)}(t)|
 \le C_{j,T}(1+n)^{2j}L_j^{n+1}
       \frac{\Gamma(\alpha)^n}{\Gamma(n\alpha+1)}
                                     t^{n\alpha-j}.
\]
Differentiation under the fixed simplex integral is justified by
this integrable positive majorant. After multiplication by
$|\lambda|^n$, the derivative series converges uniformly on
each positive-time compact interval, and after multiplication
by $t^{j-\alpha}$ it is uniformly bounded on $(0,T]$.
Its zeroth-order sum is the already proved solution, so ordinary
termwise differentiation proves the claimed derivatives and weights.

For the leading term, Lipschitz continuity gives
$k(t,s)f(s)=k(0,0)f(0)+O(t)$ on $0\le s\le t$.
Its first integral is the displayed coefficient plus
$O(t^{\alpha+1})$; the sum of terms of order $n\ge2$
is $O(t^{2\alpha})$ by the same Gamma majorant.
If $f(0)=0$, write $f(tr_n)=tr_n\int_0^1f'(\theta tr_n)d\theta$.
The first-derivative estimate then improves to a summable bound
$C(1+n)^2L^{n+1}\Gamma(\alpha)^n
t^{n\alpha}/\Gamma(n\alpha+1)$ for every $n\ge1$.
Thus $(u-f)'(t)\to0$ and $(u-f)(t)=O(t^{\alpha+1})$,
which proves the asserted $C^1$ endpoint and compatibility.
These estimates explain exactly which endpoint singularity is
retained and which stated compatibility removes it.
\end{proof}

\begin{proposition}[Galerkin existence through weak compactness]
\label{fa2:p5:galerkin-existence}
Let $H$ be a Hilbert space, with conjugate-first inner product,
and let $a$ be a bounded form, conjugate-linear first and linear
second, satisfying $\operatorname{Re}a(v,v)\ge c\pfnorm v^2$,
$c>0$. For every continuous conjugate-linear $F$ there is a
unique $u$ with $a(v,u)=F(v)$ for all $v\in H$.
It is the weak limit of the finite-dimensional Galerkin solutions;
along nested subspaces with dense union the convergence is strong.
The finite-dimensional spaces may be directed by inclusion, so
the existence argument requires no separability assumption.
\end{proposition}
\begin{proof}
On each finite-dimensional $M$, the matrix of the form is
injective by coercivity, hence invertible. Its solution $u_M$
satisfies $c\pfnorm{u_M}^2\le\operatorname{Re}F(u_M)
\le\pfnorm F\pfnorm{u_M}$, giving the uniform bound
$\pfnorm{u_M}\le\pfnorm F/c$. The Hilbert ball is weakly
compact; the directed net has a weakly convergent subnet.
For a fixed $v$, all sufficiently large $M$ contain $v$,
so bounded linearity in the second slot passes
$a(v,u_M)=F(v)$ to the weak limit. Coercivity makes a solution
unique. Every cluster point of this bounded net is therefore
the same $u$, proving full weak convergence (otherwise a subnet
outside one weak neighborhood would have a different cluster point).
Galerkin orthogonality and coercivity give
$\pfnorm{u-u_M}\le(\pfnorm a/c)\inf_{w\in M}\pfnorm{u-w}$,
by exactly the slot computation in C\'ea's lemma. Directed density
makes the last quantity tend to zero; the same holds for a
nested dense sequence when available. This gives the separate
weak-compactness existence route in \cite{B08}, rather than
only applying an already existing solution's error bound.
\end{proof}

\begin{corollary}[Sobolev energy spaces and compact lower-order terms]
\label{fa2:p5:sobolev-form-bridge}
On a bounded $C^1$ domain $\Omega\subset\mathbb R^d$ with $d\ge1$,
or on $\Omega=\varnothing$ in any dimension, put
$H=W^{1,2}_0(\Omega)$ with gradient norm, and let $H^\#$ denote the Banach space of continuous
conjugate-linear functionals, with pointwise scalar operations.
A bounded coercive form on $H$ solves every $H^\#$ datum uniquely.
In particular
$a(v,u)=\int\overline{\nabla v}\cdot A\nabla u$ is such a
form when $A$ is a bounded measurable matrix with
$\operatorname{Re}(\overline\xi\cdot A(x)\xi)\ge c|\xi|^2$
almost everywhere. A bounded zeroth-order form
$b(v,u)=\int q\overline vu$, $q\in L^\infty$, is compact
as a map $H\to H^\#$; adding it gives the precise Fredholm
compatibility alternative against the adjoint nullspace.
\end{corollary}
\begin{proof}
Poincar\'e makes the gradient norm equivalent to the complete
$W^{1,2}_0$ norm. The negative-Sobolev representation in Part IV
identifies the data as continuous forms on this space, with the
stated conjugation change. The displayed matrix inequality is
exactly coercivity; Cauchy--Schwarz proves boundedness. Apply
the preceding theorem. Rellich makes the inclusion
$J:H\to L^2$ compact, and multiplication by $q$ is bounded
on $L^2$. The bounded anti-Riesz embedding followed by pullback
is the linear map $y\mapsto[v\mapsto\langle Jv,y\rangle]$
from $L^2$ into $H^\#$. Composing it with $M_qJ$ shows
$B:u\mapsto b(\cdot,u)$ is compact; this map is not being
identified with the bilinear Banach transpose.
If $A_0:H\to H^\#$ is the coercive isomorphism and $B$ this
compact map, then $A_0+B=A_0(I+A_0^{-1}B)$.
The compact Fredholm alternative of Part II gives a closed
range of finite codimension, finite nullspace and index zero;
its range is exactly the annihilator of the adjoint nullspace.
This proves the energy-space and lower-order bridges. Additional
boundary regularity of an unbounded PDE realization is a separate
assertion and has not been inferred from these bounded forms.
\end{proof}

\begin{theorem}[Opial--Baillon nonlinear iteration and averages]
\label{fa2:p5:nonlinear-ergodic}
Let $T:H\to H$ be nonexpansive on a Hilbert space and have a
fixed point. Put $u_n=T^nf$, $S=(I+T)/2$,
$\mu_n=S^nf$, and $\sigma_n=n^{-1}\sum_{j=0}^{n-1}u_j$.
Both $\mu_n$ and $\sigma_n$ converge weakly to fixed points.
If $T$ is odd, both convergences are strong.
\end{theorem}
\begin{proof}
Use the underlying real inner product. The fixed set $K$ is
closed and convex: the two nonexpansivity balls about fixed
points $p,q$, with radii $t\pfnorm{p-q}$ and
$(1-t)\pfnorm{p-q}$, meet only at $(1-t)p+tq$.
Distances of $u_n$ to any fixed point decrease, hence all
sequences under consideration are bounded.
If $z_j\rightharpoonup z$ and $z_j-Tz_j\to0$, expanding
$\pfnorm{Tz-Tz_j}^2\le\pfnorm{z-z_j}^2$ leaves
$\pfnorm{Tz-z}^2\le0$ in the limit. This is the required
demiclosedness assertion.
Summing $\pfnorm{T\sigma_n-u_{j+1}}^2
\le\pfnorm{\sigma_n-u_j}^2$ for $0\le j<n$ gives
\[
 n\pfnorm{T\sigma_n-\sigma_n}^2
 \le\pfnorm{f-T\sigma_n}^2-\pfnorm{u_n-T\sigma_n}^2.
\]
Thus every weak cluster point of $\sigma_n$ lies in $K$.
For $p,q\in K$, the two decreasing squared distances of $u_n$
have limits, so $\langle u_n,p-q\rangle$ has a limit, and
its Ces\`aro averages do too. If $p,q$ are two weak cluster
points of $\sigma_n$, their two subsequences therefore give
$\langle p,p-q\rangle=\langle q,p-q\rangle$, hence $p=q$.
Hilbert weak sequential compactness proves full weak convergence.
For $S$ the parallelogram identity gives
\[
 \pfnorm{Su-Sv}^2+
 \pfnorm{(I-S)u-(I-S)v}^2\le\pfnorm{u-v}^2.
\]
With $v\in K$, telescoping proves
$\sum_n\pfnorm{\mu_{n+1}-\mu_n}^2<\infty$.
Thus $\mu_n-T\mu_n\to0$. Its decreasing distances to each
fixed point, demiclosedness, and the same uniqueness-of-cluster
argument prove weak convergence of $\mu_n$.

Suppose $T$ is odd. Each iterate of $T$ and of $S$ is odd
and nonexpansive. Expanding its inequalities at $v$ and $-v$
gives, for either such map $R$,
\[
 2|\langle u,v\rangle-\langle Ru,Rv\rangle|
 \le\pfnorm u^2+\pfnorm v^2-\pfnorm{Ru}^2-\pfnorm{Rv}^2.
\]
For $w_n=S^nf$ put $L=\lim\pfnorm{w_n}^2$,
$\delta_n=\pfnorm{w_n}^2-L\downarrow0$.
The inequality with $R=S^k$ shows that
$\langle w_n,w_{n+i}\rangle$ has a limit $m(i)$, with
error at most $\delta_n$, uniformly in $i$. Asymptotic
regularity gives $w_{n+i}-w_n\to0$ for each fixed $i$,
so $m(i)=L$. For every $j\ge n$ this yields
$\langle w_n,w_j\rangle\ge L-\delta_n$, hence
$\pfnorm{w_j-w_n}^2\le4\delta_n$. Thus $\mu_n$ converges
strongly.

For $u_n=T^nf$, the same reasoning gives numbers $\ell(i)$
and $\delta_n=\pfnorm{u_n}^2-\lim\pfnorm{u_k}^2\downarrow0$
with $|\langle u_n,u_{n+i}\rangle-\ell(i)|\le\delta_n$
for all $n,i$. Let $X_p=p^{-1}\sum_{i=0}^{p-1}\ell(i)$
and let $\sigma_n\rightharpoonup\sigma$ be the proved weak
limit. Directly separating the first $n$ terms of the average gives
\[
 |\langle u_n,\sigma_{n+p}\rangle-X_p|
 \le\delta_n+\frac{2n}{p}\pfnorm f^2.
\]
For fixed $n$, let $p\to\infty$. The upper and lower limits
of $X_p$ lie within $\delta_n$ of $\langle u_n,\sigma\rangle$.
As $\delta_n\to0$, $X_p$ has a limit $X$ and
$|\langle u_n,\sigma\rangle-X|\le\delta_n$.
Averaging in $n$ gives $X=\pfnorm\sigma^2$.
Expanding $\pfnorm{\sigma_n}^2$ and grouping pairs by their gap gives
\[
 \pfnorm{\sigma_n}^2
 \le\frac2{n^2}\sum_{i=0}^{n-1}(n-i)\ell(i)
       +\frac2n\sum_{i=0}^{n-1}\delta_i.
\]
The second term tends to zero; the first equals
$2n^{-2}\sum_{j=1}^njX_j\to X$. Weak lower semicontinuity
gives the reverse lower bound. Weak convergence and convergence
of norms imply $\sigma_n\to\sigma$ strongly. This proves
the nonlinear conclusions of \cite{B01}, Problem 29, separately
from the linear mean ergodic theorem.
\end{proof}

\begin{theorem}[Hardy averaging: the spectrum depends on the space]
\label{fa2:p5:hardy-averaging}
Let $Qf(x)=x^{-1}\int_0^xf(t)dt$ for $0<x\le1$.
On $C[0,1]$ put $Qf(0)=f(0)$. On the complex spaces,
\[
 \sigma_{C}(Q)=\{\lambda:|\lambda-1/2|\le1/2\},\qquad
 \sigma_{L^p}(Q)=\{\lambda:|\lambda-p'/2|\le p'/2\}
 \quad(1<p<\infty).
\]
The respective eigenvalue sets are the open disks together
with $1$ in the continuous case, and the open disk alone in
the $L^p$ case. The norm is $1$ on $C$ and $p'$ on $L^p$.
On $C^1[0,1]$ the spectrum is
$\{1\}\cup\{\lambda:|\lambda-1/4|\le1/4\}$;
its eigenvalues are $1$, $1/2$, and the latter open disk.
None of these operators is compact. But
$Q:L^p\to L^q$ is compact when $1\le q<p<\infty$,
and $Q:C\to L^q$ is compact for every finite $q$.
The source's real spectral intervals are precisely the real
sections of these complex disks, rather than complex spectra.
\end{theorem}
\begin{proof}
Since $Qf(x)=\int_0^1f(rx)dr$, the continuous-space bound
and endpoint continuity follow by dominated convergence, with
equality on constants. Minkowski and the dilation bound give
\[
 \|Qf\|_p\le\int_0^1r^{-1/p}dr\,\|f\|_p=p'\|f\|_p.
\]
The functions $x^{-1/p+\varepsilon}$ are eigenfunctions
with eigenvalues $(1/p'+\varepsilon)^{-1}$, proving sharpness.
For $\lambda\ne0$, set $F(x)=\int_0^xf$ in the eigenvalue
equation. Then $F'=F/(\lambda x)$, so every nonzero eigenfunction
is a multiple of $x^{1/\lambda-1}$. It is in $L^p$ exactly
when $\operatorname{Re}(1/\lambda)>1/p'$, and is continuous
through zero exactly when $\operatorname{Re}(1/\lambda)>1$
or $\lambda=1$. These are the asserted open disks and the
exceptional continuous constant. Zero is not an eigenvalue.

To exclude all exterior points from the spectrum, put
$a=1/\lambda$ and solve $(Q-\lambda I)u=f$ by
\[
 u(x)=-\lambda^{-1}f(x)-\lambda^{-2}
                 x^{a-1}\int_0^xt^{-a}f(t)dt.
\]
The integral converges and defines a bounded operator on $L^p$
if $\operatorname{Re}a<1/p'$, because its second factor equals
$\int_0^1r^{-a}f(rx)dr$ and Minkowski bounds its norm by
$\|f\|_p/(1/p'-\operatorname{Re}a)$.
For $C$ use instead $\operatorname{Re}a<1$ and the same
integral; dominated convergence also gives its endpoint value.
Differentiation of its primitive verifies both the equation
and uniqueness. These inequalities characterize the exterior
of the two closed disks. Since spectra are closed and contain
all their open-disk eigenvalues, they contain the disk closures;
the exterior resolvent formula proves equality.

The isomorphism $f\mapsto(f(0),f')$ from $C^1$ onto
$\mathbb C\oplus C$ converts $Q$ into $1\oplus R$, where
$Rv(x)=\int_0^1r v(rx)dr$.
The substitution $y=x^2$ conjugates $R$ on $C$ to $Q/2$:
if $w(y)=v(\sqrt y)$, then
$Rv(x)=\tfrac12 Qw(x^2)$. Its spectrum is the smaller disk,
and the continuous-case eigenvalue description gives its open
disk plus $1/2$. This proves the $C^1$ assertions without
incorrectly transferring the $C$ spectrum across a stronger norm.
Each space has infinitely many nonzero eigenvalues accumulating
away from zero, excluding compactness by the compact spectral theorem.

For compactness into the weaker targets discard an output interval
$0<x<\delta$. In the $L^p\to L^q$ case its norm is at most
$p'\delta^{1/q-1/p}$; for $C\to L^q$ it is at most
$\delta^{1/q}$. On $[\delta,1]$ the primitives of the source
unit ball are uniformly bounded and equicontinuous, by H\"older
(or the supremum norm), so division by $x$ preserves these
bounds and Ascoli gives compactness. A continuous output cutoff
between $\delta/2$ and $\delta$ yields compact operators on the
whole target and the same vanishing norm error. Thus the limiting
operator is compact into exactly the stated weaker targets.
\end{proof}

\begin{proposition}[An interior point that separates an open convex set]
\label{fa2:p5:cone-geometry}
If $D$ is a nonempty open convex subset of a real Hilbert space
and $x_0\notin D$, some $w_0\in D$ satisfies
$\langle w_0-x_0,w-x_0\rangle>0$ for every $w\in D$.
\end{proposition}
\begin{proof}
Translate $x_0$ to zero and put $C=\bigcup_{t>0}tD$.
This is a proper open convex cone. Set
$\Sigma=\{u:\langle u,c\rangle\le0\text{ for all }c\in C\}$.
Separation of any point outside $C$ from this open cone gives
\[
 C=\{x:\langle u,x\rangle<0\text{ for every }u\in\Sigma\setminus0\}.
\]
Indeed a nonzero supporting functional is strictly negative on
the open cone and conic scaling makes its separating level zero;
this also shows $\Sigma\ne\{0\}$. Choose $\omega\in C$.
Then $K=\{u\in\Sigma:\langle u,\omega\rangle=-1\}$ is
nonempty, closed and convex. If $B(\omega,r)\subset C$,
the polar inequality gives $r\|u\|\le1$ on $K$, so it
is bounded. Let $a=P_K0\ne0$. The projection inequality
$\langle a,u\rangle\ge\|a\|^2$ for $u\in K$,
after normalizing any $v\in\Sigma\setminus0$ into $K$,
gives $\langle a,v\rangle>0$. The displayed cone description
then yields $-a\in C$. Since $a\in\Sigma\setminus0$,
$\langle-a,c\rangle>0$ for every $c\in C$.
Write $-a=t w_0$ with $t>0$, $w_0\in D$. Division by $t$
proves the assertion. This retains the extra geometric conclusion
of \cite{B01}, Problem 25, rather than confusing it with the
weaker existence of an external supporting functional.
\end{proof}

\begin{proposition}[Small perturbations of one-sided regularizers]
\label{fa2:p5:one-sided-regularizer-perturbation}
Let $K:X\to Y$ and $R:Y\to X$ be bounded between Banach spaces,
and let $C:X\to Y$ be bounded with $\|R\|\|C\|<1$.
If $RK-I_X$ is compact, then
$\widetilde R=(I_X+RC)^{-1}R$ is a left regularizer of $K+C$.
If $KR-I_Y$ is compact, then
$\widetilde R=R(I_Y+CR)^{-1}$ is a right regularizer of $K+C$.
\end{proposition}
\begin{proof}
The stated product bound gives both inverses by the Neumann
series. Write $RK=I_X-A$ with $A$ compact. Then
\[
 (I_X+RC)^{-1}R(K+C)=I_X-(I_X+RC)^{-1}A,
\]
whose defect is compact. If $KR=I_Y-B$, the right-sided
identity is
$(K+C)R(I_Y+CR)^{-1}=I_Y-B(I_Y+CR)^{-1}$.
These are the required exact formulas, also when $R=0$.
The threshold printed in \cite{B06}, Problem 5.4, is
$\|C\|<\|R\|$ and is insufficient: on infinite-dimensional
$H$, take $K=I/2$, $R=2I$, $C=-I/2$. Its printed inequality
holds, but $K+C=0$ has neither regularizer, since the identity
on $H$ is not compact. The product criterion states and proves
the corrected mathematical assertion.
\end{proof}

\begin{example}[Dense least-squares trials can fail without regularization]
\label{fa2:p5:least-squares-nonconvergence}
There is a compact injective dense-range map on $\ell^2$ and
nested finite-dimensional trial spaces with dense union such that
the unique exact-data least-squares approximants have residuals
tending to zero but an unbounded subsequence of solution errors.
\end{example}
\begin{proof}
Use the real Hilbert basis $e_0,e_1,\ldots$ (the example also
works over complex scalars). Put $g=e_0+\sum_{j\ge1}2^{-j}e_j$,
$Ae_0=0$, and $Ae_j=\lambda_je_j$, with
$\lambda_{2k}=2^{-k^2}$, $\lambda_{2k+1}=2^{-k}$
for $k\ge1$, and $\lambda_1=1$.
These positive numbers tend to zero but are intentionally not
ordered. Define $Sx=Ax+\langle e_0,x\rangle g$.
It is compact. Its $e_0$ component shows injectivity; $S^*y=0$
first gives every positive coordinate of $y$ zero and then its
$e_0$ coordinate zero, so its range is dense.
For $f=Sg$ and $X_n=\operatorname{span}(e_0,\ldots,e_n)$,
the normal equations give a unique minimizer
$x_n=\sum_{j=0}^n\gamma_je_j$, since $S$ is injective.
Writing $g_j=2^{-j}$, the equations in coordinates $j\le n$
and then in the $e_0$ trial direction give respectively
\[
 \lambda_j(\gamma_j-g_j)=(1-\gamma_0)g_j,\qquad
 (\gamma_0-1)\left(1+\sum_{j>n}g_j^2\right)
                            =\sum_{j>n}\lambda_jg_j^2.
\]
Therefore
\[
 \|x_n-g\|\ge
 \frac{\lambda_{n+1}g_{n+1}^2g_n}
           {(1+\|g\|^2)\lambda_n}.
\]
For $n=2k$ the right side is
$c2^{k^2-7k-2}\to\infty$. Yet the residual of a minimizer
is at most $\|S(P_{X_n}g-g)\|\to0$.
This explicitly closes the instability mechanism in
\cite{B06}, Problem 17.2: density and vanishing residual alone
do not replace the noise/discretization stability conditions
proved in the regularization theorems.
\end{proof}

\begin{theorem}[Cordial noncompact Volterra counterexample]
\label{fa2:p5:cordial}
Let $T>0$. On $C[0,T]$ define, for $t>0$,
$Vf(t)=\int_0^t(t^2-s^2)^{-1/2}f(s)ds$ and
$Vf(0)=(\pi/2)f(0)$. This is bounded with norm $\pi/2$.
For every $r\ge0$, $t^r$ is an eigenfunction with eigenvalue
$\lambda_r=\int_0^{\pi/2}\sin^r\theta\,d\theta
=\tfrac12 B((r+1)/2,1/2)$.
Consequently $V$ is neither compact nor quasinilpotent.
\end{theorem}
\begin{proof}
The substitution $s=t\sin\theta$ gives
$Vf(t)=\int_0^{\pi/2}f(t\sin\theta)d\theta$, also at zero
by dominated convergence. This proves continuity, the norm bound,
and equality by the constant function. Applying it to $t^r$
gives the formula. The eigenvalues depend continuously and strictly
decrease from $\pi/2$ to zero, by dominated convergence and strict
pointwise inequality. Thus there are uncountably many nonzero
spectral values; a compact operator has at most countably many,
with only zero as possible accumulation. The positive eigenvalue
at $r=0$ also directly excludes quasinilpotence.
The zero-endpoint value must be this extension; setting it to zero
would not define a map on $C[0,T]$.
\end{proof}

\begin{theorem}[The singular system and power norm of integration]
\label{fa2:p5:volterra-singular-system}
For $Vf(t)=\int_0^tf(s)ds$ on $L^2(0,T)$, $T>0$, a singular system is
\[
 e_j(t)=\sqrt{2/T}\cos((j-\tfrac12)\pi t/T),\quad
 f_j(t)=\sqrt{2/T}\sin((j-\tfrac12)\pi t/T),\quad
 \sigma_j=\frac{T}{(j-\tfrac12)\pi}.
\]
In particular $\pvnorm V=2T/\pi$ although its spectrum is $\{0\}$.
Moreover
$\pvnorm{V^n}\sim T^n/(2n!)$ as $n\to\infty$.
\end{theorem}
\begin{proof}
Direct integration gives $Ve_j=\sigma_jf_j$ and
$V^*f_j=\sigma_je_j$, where $V^*g(s)=\int_s^Tg(t)dt$.
The half-integer sine and cosine systems are complete: extend a
function evenly at zero and oddly at $T$ (or the reverse), and
then periodically with period $4T$; its ordinary Fourier series
has exactly these modes. Fourier completeness proves both claims.
Thus the singular system is complete and its largest value is the
operator norm. The earlier factorial estimate and injectivity of
$I-\lambda V$ give zero spectrum, consistently with this nonzero
singular spectrum.

Scale to $T=1$. The kernel of $V^n$ is
$(t-s)_+^{n-1}/(n-1)!$. Use the unitary input coordinate $b=ns$
and output coordinate $a=n(1-t)$, with the $n^{-1/2}$ factors
required by $L^2$ measure. Then $n!V^n$ has kernel
$q_n(a,b)=(1-(a+b)/n)_+^{n-1}$ on $[0,n]^2$.
Extend it by zero to the positive quadrant. It tends pointwise
to $e^{-a-b}$ and, for $n\ge2$, is bounded by
$e^{-(a+b)/2}$. Dominated convergence of the squared difference
gives Hilbert--Schmidt, hence operator-norm convergence to the
rank-one kernel $e^{-a}e^{-b}$, whose norm is $1/2$.
The nonzero part of each extended operator has the same norm
as $n!V^n$, proving the asymptotic. Rescaling restores $T^n$.
This is a sharper result than the elementary factorial upper bound.
\end{proof}

\begin{proposition}[A genuine long-time convolution bridge]
\label{fa2:p5:long-time-convolution}
Let $k\in L^1(0,\infty)$ with $a=\pvnorm k_1<1$.
Then $r=\sum_{n\ge1}k^{*n}\in L^1$,
$\pvnorm r_1\le a/(1-a)$, and $u=f+r*f$ uniquely solves
$u=f+k*u$ in every $L^p$, $1\le p\le\infty$.
Its Laplace transform satisfies
$\mathcal Lr(z)=\mathcal Lk(z)/(1-\mathcal Lk(z))$ for
$\operatorname{Re}z\ge0$. If bounded $f(t)\to\ell$, then
$u(t)\to\ell/(1-\int k)$.
Neither finite-interval solvability nor $k\in L^1$ alone implies
long-time stability.
\end{proposition}
\begin{proof}
Young's inequality makes convolution by $k$ have norm at most $a$.
The geometric series converges in $L^1$ and in every indicated
operator space. Associativity and multiplication of partial sums
prove the inverse. Absolute Fubini gives the Laplace identity,
since $|\mathcal Lk(z)|\le a<1$.
In $r*f(t)=\int_0^tr(s)f(t-s)ds$, bounded convergence with the
integrable majorant $\pvnorm f_\infty|r|$ gives the asserted limit;
integrating the $r$ series gives $\int r=(\int k)/(1-\int k)$.
For contrast $k(t)=2e^{-t}\in L^1$ and $f=1$ give
$u(t)=2e^t-1$ by direct substitution. Thus the extra stability
hypothesis is mathematical content, not an automatic consequence
of the local Volterra theorem.
\end{proof}

\chapter{Approximation, Projection, and Regularization}
\label{fa2:p5:approximation-chapter}

\section{Norm, strong, and collectively compact approximation}

\begin{lemma}[Uniform convergence on compact sets]
\label{fa2:p5:compact-uniform-convergence}
If $T_n:X\to Y$ are uniformly bounded linear maps and converge
pointwise to $T$, then convergence is uniform on every compact
subset of $X$. If $X$ is Banach and pointwise convergence is given,
uniform boundedness of the sequence follows automatically.
If $A$ is an invertible bounded map between Banach spaces and
$\pvnorm{A^{-1}(A_n-A)}<1$, then $A_n$ is invertible, with the
usual Neumann inverse and quantitative inverse-difference bound.
\end{lemma}
\begin{proof}
For compact $K$ choose a finite $\varepsilon$-net.
The common norm bound controls the error between a point and its
net point, and convergence on finitely many net points controls
the remaining error. The pointwise limit is bounded by the same
bound. The second assertion is uniform boundedness from Part I.
For the inverse assertion factor
$A_n=A(I+A^{-1}(A_n-A))$ and sum the norm-convergent geometric
series. Multiplication by $A_n^{-1}$ and $A^{-1}$ gives
$A_n^{-1}-A^{-1}=A_n^{-1}(A-A_n)A^{-1}$ and its norm estimate.
\end{proof}

\begin{theorem}[Collectively compact stability]
\label{fa2:p5:collective-compact}
Let $K_n:X\to X$ on a Banach space be collectively compact:
the union of their images of the unit ball has compact closure.
Suppose $K_nx\to Kx$ for every $x$, and $I-K$ is injective.
Then $K$ is compact, $I-K$ is invertible, and for large $n$,
$I-K_n$ is invertible with uniformly bounded inverses.
If $u_n-K_nu_n=f_n$, $u-Ku=f$, then
\[
 \pvnorm{u_n-u}\le C\bigl(\pvnorm{(K_n-K)u}+
                                      \pvnorm{f_n-f}\bigr).
\]
\end{theorem}
\begin{proof}
The collective compact closure contains $K$ of the unit ball,
so $K$ is compact. Pointwise boundedness and uniform boundedness
give $\sup_n\pvnorm{K_n}<\infty$. The previous lemma, on the
collective compact closure, gives
$\pvnorm{(K_n-K)K_n}\to0$ and $\pvnorm{(K_n-K)K}\to0$.
Fredholm gives $R=(I-K)^{-1}$.
Put $B_n=I+RK_n$ and $S_n=R(K_n-K)K_n$. Direct multiplication
using $R-I=RK$ gives
$B_n(I-K_n)=I-S_n$. For large $n$, $\pvnorm{S_n}<1$,
so $I-K_n$ is injective. Each $K_n$ is compact, hence Fredholm
makes it surjective. Its inverse is
$(I-S_n)^{-1}B_n$, uniformly bounded by
$(1+\pvnorm{RK_n})/(1-\pvnorm{S_n})$.
Subtract the two equations and apply this inverse to obtain the
displayed estimate. Strong convergence alone would not justify
this inverse conclusion; collective compactness is exactly the
additional uniform compactness used in the proof.
\end{proof}

\begin{proposition}[The spectral-radius criterion for iteration]
\label{fa2:p5:iteration-radius}
For bounded $T$ on a complex Banach space,
$r(T)<1$ is equivalent to $\pvnorm{T^n}\to0$ and to the existence
of $C<\infty$, $q<1$ with $\pvnorm{T^n}\le Cq^n$.
It gives uniform geometric convergence of $x_{n+1}=Tx_n+b$
to $(I-T)^{-1}b$. Pointwise convergence $T^nx\to0$ for every
$x$ is weaker and does not require $r(T)<1$.
\end{proposition}
\begin{proof}
Choose $r(T)<q<1$ in the spectral-radius formula; for large $n$,
$\pvnorm{T^n}^{1/n}<q$, and absorb finitely many initial powers
in $C$. Conversely if $\pvnorm{T^n}\to0$, some power has norm
less than one, so $r(T)\le\pvnorm{T^m}^{1/m}<1$.
The iteration error is $T^n$ of the initial error.
On $\ell^2$ the left shift has powers tending strongly to zero
by the tail norm, but every power has norm one and its spectrum
is the closed unit disk. Thus a strong convergence statement
must not be promoted to the uniform geometric criterion.
\end{proof}

\section{General projection and the coercive Hilbert case}

\begin{theorem}[Abstract projection criterion]
\label{fa2:p5:projection-method}
Let $A:X\to Y$ be bounded and injective between Banach spaces.
Let finite-dimensional $X_n\subset X$, $Y_n\subset Y$ have the
same dimension, let $P_n:Y\to Y_n$ be bounded projections, and
assume $\operatorname{dist}(x,X_n)\to0$ for every $x$.
Suppose $P_nA|_{X_n}$ is invertible for all sufficiently large $n$.
Then the projection method converges for every exact datum $Ax$
if and only if
\[
 Q_n=(P_nA|_{X_n})^{-1}P_nA:X\to X_n
 \quad\hbox{have uniformly bounded norms}.
\]
In that case
$\pvnorm{Q_nx-x}\le(1+\sup\pvnorm{Q_n})\operatorname{dist}(x,X_n)$.
This does not assert stability for arbitrary noisy data unless the
maps $(P_nA|_{X_n})^{-1}P_n$ themselves are uniformly bounded.
\end{theorem}
\begin{proof}
$Q_n$ is the identity on $X_n$, so for every $v_n\in X_n$,
$Q_nx-x=(Q_n-I)(x-v_n)$. This proves sufficiency and the error
bound by taking the infimum. Necessity follows from uniform
boundedness applied to pointwise convergence of the bounded maps
$Q_n$ on the Banach domain. The noise qualification follows by
applying the data-solution map to the actual data error, whose
norm is governed by that map, not by its composition with $A$.
\end{proof}

\begin{corollary}[Compact second-kind projection]
\label{fa2:p5:second-kind-projection}
If $K$ is compact, $I-K$ injective, and bounded projections
$P_n:X\to X_n$ converge strongly to $I$, then
$u_n-P_nKu_n=P_nf$ has a unique solution in $X_n$ for large $n$,
and $u_n\to u=(I-K)^{-1}f$.
\end{corollary}
\begin{proof}
Uniform boundedness and compact-uniform convergence give
$\pvnorm{P_nK-K}\to0$. The inverse perturbation lemma makes
$I-P_nK$ invertible with uniformly bounded inverses.
For a right side in $X_n$ its solution satisfies
$u_n=P_nf+P_nKu_n\in X_n$. Subtract the equations:
$u_n-u=(I-P_nK)^{-1}(P_n-I)u$, using $f=u-Ku$.
This tends to zero, and also gives a best-approximation bound
through $P_n-I$ and its uniform norm.
\end{proof}

\begin{theorem}[Galerkin, C\'ea, and Ritz]
\label{fa2:p5:galerkin}\label{fa2:p5:cea}\label{fa2:p5:ritz}
Let $a(v,u)$ on a Hilbert space be conjugate-linear in $v$,
linear in $u$, bounded by $M\pvnorm v\pvnorm u$, and satisfy
$\operatorname{Re}a(u,u)\ge\alpha\pvnorm u^2$, $\alpha>0$.
For continuous anti-linear data $F$, let $u$ be its Lax--Milgram
solution. On every closed trial subspace $H_n$ there is a unique
Galerkin solution, and
\[
 a(v_n,u-u_n)=0,
 \qquad \pvnorm{u-u_n}\le(M/\alpha)
                    \inf_{v_n\in H_n}\pvnorm{u-v_n}.
\]
If $a$ is Hermitian, $u_n$ is its energy orthogonal projection,
and the error is exactly the minimum energy distance. Approximation
of every vector by the trial spaces implies convergence; a dense
union alone is not sufficient without eventual approximation.
\end{theorem}
\begin{proof}
The restricted form has the same constants and its subspace is
complete; Part II's Lax--Milgram theorem gives existence and the
bound. Subtraction gives the displayed test-slot orthogonality.
With $e=u-u_n$ and $v_n\in H_n$,
$a(e,e)=a(u-v_n,e)$ because $v_n-u_n$ is a permitted test.
Coercivity, boundedness, and cancellation of $\pvnorm e$ give
the estimate, treating $e=0$ separately.
For Hermitian $a$, it is a complete inner product equivalent to
the original one; orthogonality and Pythagoras give exact energy
minimization. The error estimate proves convergence whenever
the infimum tends to zero. For failure of the weaker condition,
in $\ell^2$ take $H_{2n}=\operatorname{span}(e_1,\ldots,e_n)$
and $H_{2n-1}=\{0\}$. Their union is dense, but the projection
of $e_1$ is zero along every odd index. In finite coordinates the
matrix is $a(\phi_i,\phi_j)$, invertible by the same coercivity;
this is an actual matrix consequence, not a separate assumption.
\end{proof}

\section{Regularization and the minimum-norm target}

\begin{theorem}[Tikhonov for every bounded Hilbert operator]
\label{fa2:p5:tikhonov}
Let $A:H\to G$ be bounded, $\alpha>0$, and $f\in G$.
The functional
$\pvnorm{Ax-f}^2+\alpha\pvnorm x^2$ has a unique minimizer
\[
 x_\alpha=R_\alpha f,
 \quad R_\alpha=(\alpha I+A^*A)^{-1}A^*,
 \quad \pvnorm{R_\alpha}\le\frac1{2\sqrt\alpha}.
\]
For exact solvable data $f=Ax^\dagger$, with
$x^\dagger\perp\ker A$, $x_\alpha\to x^\dagger$.
If $\pvnorm{f^\delta-f}\le\delta$, convergence remains true when
$\alpha(\delta)\to0$ and $\delta/\sqrt{\alpha(\delta)}\to0$.
\end{theorem}
\begin{proof}
The bounded form associated with $\alpha I+A^*A$ is coercive,
so it is invertible by Lax--Milgram. Expanding the functional
at its normal-equation solution gives
$\pvnorm{A(x-x_\alpha)}^2+\alpha\pvnorm{x-x_\alpha}^2$
plus its minimum value; the cross term vanishes. This proves
existence and uniqueness by an exact identity.
Put $B=A^*A$. The positive continuous calculus gives
$\pvnorm{R_\alpha}^2
=\pvnorm{B(\alpha+B)^{-2}}\le1/(4\alpha)$.
Also $R_\alpha A=B/(\alpha+B)$ is a contraction, zero on
$\ker A$, and converges strongly to the identity on its
perpendicular complement. To see the last claim without assuming
compactness, for $x=Bz$ its error has norm at most
$\alpha\pvnorm z$; such vectors are dense in
$(\ker B)^\perp$, and the contraction bound extends the convergence.
Thus exact-data convergence holds. Add the noise estimate
$\pvnorm{R_\alpha(f^\delta-f)}\le\delta/(2\sqrt\alpha)$.
The minimum-norm condition is essential; no regularization can
recover an unobservable component in $\ker A$ from $f$.
\end{proof}

\begin{proposition}[A source condition and its converse]
\label{fa2:p5:tikhonov-rate}
For exact data and $x^\dagger\perp\ker A$,
$x^\dagger\in\operatorname{ran}(A^*A)$ implies
$\pvnorm{x_\alpha-x^\dagger}\le\alpha\pvnorm w$ when
$x^\dagger=A^*Aw$. Conversely an $O(\alpha)$ exact-data error
implies that range condition.
\end{proposition}
\begin{proof}
The bias is $-\alpha(\alpha+B)^{-1}x^\dagger$.
On $x^\dagger=Bw$ the contraction $B/(\alpha+B)$ gives the bound.
Conversely boundedness of
$z_\alpha=(\alpha+B)^{-1}x^\dagger$ gives a weakly convergent
subsequence $z_\alpha\rightharpoonup w$ as $\alpha\to0$.
Since $Bz_\alpha=x^\dagger-\alpha z_\alpha\to x^\dagger$
in norm, weak continuity gives $Bw=x^\dagger$.
This proves necessity, not just the easy sufficient rate estimate.
\end{proof}

\begin{theorem}[Singular filters and cutoff]
\label{fa2:p5:spectral-regularization}
For compact $A$ with singular system, let
$R_\eta f=\sum_j\psi_\eta(\sigma_j)\langle f_j,f\rangle e_j$.
Assume for every parameter
$C_\eta=\sup_j|\psi_\eta(\sigma_j)|<\infty$,
$\sup_{\eta,j}|\sigma_j\psi_\eta(\sigma_j)|<\infty$, and
$\sigma_j\psi_\eta(\sigma_j)\to1$ for every fixed $j$.
Then $R_\eta Ax^\dagger\to x^\dagger$ for the minimum-norm
solution, and noisy convergence holds if
$\eta\to0$ and $\delta C_\eta\to0$.
For spectral cutoff $\psi_\varepsilon(\sigma)=\sigma^{-1}
1_{\{\sigma\ge\varepsilon\}}$, $C_\varepsilon\le1/\varepsilon$.
\end{theorem}
\begin{proof}
Bessel gives boundedness with norm at most $C_\eta$.
On exact data the squared error is the sum of
$|\sigma_j\psi_\eta(\sigma_j)-1|^2
|\langle e_j,x^\dagger\rangle|^2$.
Dominated convergence for the summable coefficient sequence
proves convergence. The noise bound is $\delta C_\eta$.
For cutoff the multiplier is either zero or one on the exact
coefficients, so all hypotheses hold. The Tikhonov filter
$\psi_\alpha(\sigma)=\sigma/(\alpha+\sigma^2)$ is another
instance, consistent with the preceding general bounded-operator
proof. A noise parameter, a filter parameter, and a discretization
dimension have not been identified with one another.
\end{proof}

\begin{theorem}[Moment projection as regularization]
\label{fa2:p5:moment-regularization}
Let $A:X\to Y$ be bounded between Hilbert spaces. Let $Y_n$ be
nested finite-dimensional subspaces with dense union
in $(\ker A^*)^\perp$, and put $X_n=A^*Y_n$.
For solvable $f=Ax^\dagger$, with $x^\dagger\perp\ker A$, the equations
$P_{Y_n}Ax_n=P_{Y_n}f$, $x_n\in X_n$, have a unique solution,
and $x_n=P_{X_n}x^\dagger\to x^\dagger$.
For noisy data the error is at most
$\operatorname{dist}(x^\dagger,X_n)+\delta/\beta_n$, where
$\beta_n>0$ is the smallest singular value of
$P_{Y_n}A|_{X_n}$ when $Y_n\ne0$. In zero-dimensional stages
the solution and the noise term are both zero, so no singular
value is needed (equivalently set $1/\beta_n=0$ there).
A choice $n(\delta)\to\infty$ with
$\delta/\beta_{n(\delta)}\to0$ is a regularization.
\end{theorem}
\begin{proof}
$A^*$ is injective on $Y_n$. The matrix of $P_{Y_n}AA^*$
has quadratic form $\pvnorm{A^*y}^2>0$ for nonzero $y\in Y_n$,
so it is invertible. The constraint says precisely
$\langle A^*y,x_n-x^\dagger\rangle=0$ for every $y\in Y_n$;
hence $x_n$ is its orthogonal projection.
The union of $A^*Y_n$ is dense in
$\overline{\operatorname{ran}A^*}=(\ker A)^\perp$ by continuity
and the density assumption, so the projections converge.
Finite-dimensional inversion bounds the noise term by
$\delta/\beta_n$. An admissible slowly increasing choice exists:
choose decreasing thresholds $\delta_n$ with
$\delta_n/\beta_n\le1/n$ and use dimension $n$ between
successive thresholds. Thus discretization becomes a proved
regularization, not an assumption that every Galerkin scheme is
stable for a first-kind problem.
\end{proof}

\begin{proposition}[Quasi-solutions and discrepancy constraints]
\label{fa2:p5:quasi-solution}
The residual has a minimizer on every closed Hilbert ball.
It is unique if $A$ is injective. Every nonempty discrepancy set
$\{x:\pvnorm{Ax-f}\le\delta\}$ has a unique minimum-norm member.
If $\operatorname{dist}(f,\overline{\operatorname{ran}A})<\delta<\pvnorm f$,
that member is a Tikhonov solution with a parameter chosen so that
its residual is $\delta$.
\end{proposition}
\begin{proof}
A closed Hilbert ball is weakly compact, and its residual norm
is weakly lower semicontinuous, proving existence. Strict convexity
of the squared residual on an injective image proves uniqueness.
The discrepancy set is closed convex, so projection of zero gives
its unique minimum-norm member. For the final assertion, the
Tikhonov residual is continuous in $\alpha>0$ by the inverse
resolvent identity. It increases from
$\operatorname{dist}(f,\overline{\operatorname{ran}A})$ to
$\pvnorm f$ as $\alpha$ increases from zero to infinity;
the positive spectral calculus proves these limits and monotonicity
for the squared residual, including the component in $\ker A^*$.
The intermediate value theorem gives a parameter. Comparing its
Tikhonov minimum with any feasible $x$ shows
$\delta^2+\alpha\pvnorm{x_\alpha}^2
\le\pvnorm{Ax-f}^2+\alpha\pvnorm x^2
\le\delta^2+\alpha\pvnorm x^2$.
Thus it is exactly the minimum-norm feasible member.
The strict lower inequality excludes an unattained range-distance
endpoint, at which the discrepancy set could be empty.
\end{proof}

\begin{theorem}[Reproducing kernels and integral-operator ranges]
\label{fa2:p5:rkhs}
A Hilbert space of functions with bounded point evaluations has
unique vectors $k_x$ with $f(x)=\langle k_x,f\rangle$ and
kernel $K(x,y)=\langle k_x,k_y\rangle$.
Conversely every positive definite scalar kernel defines a unique
Hilbert space with this reproducing property.
If $Au(x)=\int k(x,z)u(z)dz$ is a bounded map from $L^2$ into
functions with $k(x,\cdot)\in L^2$, its range, with minimum-preimage
norm, is such a Hilbert space and
$K(x,y)=\int k(x,z)\overline{k(y,z)}dz$.
\end{theorem}
\begin{proof}
Hilbert Riesz gives $k_x$, and its Gram matrices are positive.
Conversely start from finite sums $\sum c_jK(\cdot,x_j)$ with
quadratic norm $\sum\overline c_ic_jK(x_i,x_j)$.
Positivity and the two-vector Cauchy--Schwarz inequality show that
zero-norm sums vanish at every point, so quotienting is well defined.
Point evaluation has bound $K(x,x)^{1/2}$ and extends to the
completion, which is therefore still a space of functions.
Kernel sections span densely, forcing uniqueness of the construction.
For the range model quotient $L^2$ by the closed kernel of $A$,
or equivalently use its orthogonal complement. The pullback norm
makes the range complete even if it is not closed in its ambient
function norm. The evaluation vector in the domain is
$\overline{k(x,\cdot)}$, already perpendicular to the kernel,
and its image under $A$ is the reproducing vector.
Evaluate $A(\overline{k(y,\cdot)})$ at $x$ to obtain the
displayed $K(x,y)$, with exactly the
conjugation prescribed by the inner-product convention.
\end{proof}

\chapter{Convex Representation, Invariance, and Classical Counterexamples}
\label{fa2:p5:representation-chapter}

\section{Positive extension and affine fixed points}

\begin{theorem}[Positive extensions with a majorizing subspace]
\label{fa2:p5:positive-extension}
Let $E$ be a real ordered vector space and $M\subset E$ a subspace
such that every $x\in E$ is bounded above by some member of $M$.
Every positive linear functional $\ell$ on $M$ has a positive
linear extension to $E$. In particular this holds when $M$ contains
an order unit $e$. In the Archimedean order-unit norm the extension
is bounded with norm $\ell(e)$.
\end{theorem}
\begin{proof}
Set $p(x)=\inf\{\ell(m):m\in M,\ m\ge x\}$.
This is finite: an upper majorant gives an upper bound, and if
$n\in M$ majorizes $-x$, then $m+n\ge0$ gives
$\ell(m)\ge-\ell(n)$ for every eligible $m$.
The infimum is sublinear, by adding majorants and rescaling, and
$p(m)=\ell(m)$ on $M$. Real Hahn--Banach extends $\ell$ to
$L\le p$. For $x\ge0$, $p(-x)\le0$, hence $L(x)\ge0$.
If $-te\le x\le te$, positivity gives $|L(x)|\le tL(e)$.
Taking the infimum over $t$ proves the norm bound, and the test
$e$ proves equality when $\pvnorm e=1$. The Archimedean assumption
is what makes this gauge a norm rather than only a seminorm.
Positive extension is therefore a consequence of explicit order
majorization, not of an arbitrary unqualified Hahn--Banach extension.
\end{proof}

\begin{theorem}[Kakutani's equicontinuous affine-group theorem]
\label{fa2:p5:kakutani-affine}
An equicontinuous group of affine homeomorphisms of a nonempty
compact convex subset $K$ of a Hausdorff locally convex space has
a common fixed point.
\end{theorem}
\begin{proof}
The nonempty compact convex invariant subsets contain a minimal
one $Q$: order by reverse inclusion and use compactness for
intersections of chains, then Zorn. Suppose $x\ne y$ lie in $Q$.
Choose a zero neighborhood $W$ excluding $x-y$. Equicontinuity
gives a neighborhood $V$ such that differences in $V$ are sent
into $W$ by every group element. Applying the inverse group
elements shows $Tx-Ty\notin V$ for every $T$.
Let $z=(x+y)/2$. Its orbit closure has closed convex hull $Q$
by minimality. The Milman converse
(\ref{fa2:p1:milman-converse}) implies that every extreme point
$q$ of $Q$ lies in this orbit closure. Choose a net $T_iz\to q$;
compactness supplies a subnet with $T_ix\to a$, $T_iy\to b$.
Affineness gives $q=(a+b)/2$, so extremality forces $a=b=q$.
But all differences lie in the closed complement of the open
neighborhood $V$, contradicting their limit zero. Hence $Q$ is a
singleton, which is fixed by the entire group. Both inverses and
equicontinuity were used; commuting arbitrary maps are a distinct
fixed-point theorem.
\end{proof}

\begin{theorem}[Commuting affine maps and invariant means]
\label{fa2:p5:markov-kakutani}\label{fa2:p5:invariant-mean}
A commuting family of continuous affine self-maps of a nonempty
compact convex subset of a Hausdorff locally convex space has a
common fixed point. Every abelian group $G$ has a positive unital
translation-invariant functional on $\ell^\infty(G)$ and therefore
a translation-invariant finitely additive probability on its power set.
\end{theorem}
\begin{proof}
For one affine map $T$, the averages
$a_n=n^{-1}\sum_{j=0}^{n-1}T^jx$ lie in $K$ and
$Ta_n-a_n=(T^nx-x)/n\to0$, since compact sets are bounded.
A convergent subnet gives a fixed point. Its fixed set is compact
convex and invariant under every commuting map. Induction proves
common fixed points for each finite subfamily; compactness and the
finite-intersection property give the full conclusion.

The positive functionals $L$ with $L(1)=1$ form a nonempty
weak-star compact convex subset of the dual unit ball: point
evaluation gives a member, and positivity/unitality are closed
scalar tests. Transposition of translations gives commuting
weak-star continuous affine maps when $G$ is abelian.
The common fixed functional is the desired mean.
Define $m(A)=L(1_A)$; positivity, disjoint finite additivity,
normalization, and translation invariance follow immediately.
Countable additivity is not asserted. Nor can abelianness simply
be dropped: in the free group on $a,b$, partition nonidentity
reduced words by initial letters $a,a^{-1},b,b^{-1}$.
Singletons have mass zero in an infinite group. Writing
$A,A^{-},B,B^{-}$ for the four classes, left cancellation gives
$a(G\setminus A^{-})=A$ and
$b(G\setminus B^{-})=B$ give
$m(A)+m(A^{-})=1=m(B)+m(B^{-})$, whereas their disjoint union
has mass one. This explicit computation supplies the contradiction.
\end{proof}

\section{Haar measure without an action-continuity shortcut}

\begin{theorem}[Normalized Haar measure on compact groups]
\label{fa2:p5:haar}
Every compact Hausdorff group has a unique normalized left-invariant
regular Borel measure. It is also right invariant and invariant
under inversion. Every nonempty open set has positive mass.
\end{theorem}
\begin{proof}
For $f\in C(G)$, the map $s\mapsto L_sf$, $L_sf(x)=f(sx)$,
is norm-continuous: continuity of $(s,x)\mapsto f(sx)$ is uniform
in $x$ by compactness and a finite cover. Thus its orbit is norm
compact. Its closed convex hull is norm compact as well: finite
$\varepsilon$-nets reduce convex combinations to the convex hull
of finitely many points, a compact finite-dimensional simplex;
completeness closes the total-boundedness argument.
The left translations are an equicontinuous group of isometries
on this hull. Kakutani gives a fixed function, which must be
constant because $L_sf(e)=f(s)$. Thus a constant $c$ is uniformly
approximable by convex combinations of left translates.
Similarly a constant $c'$ is available from right translates.
If $\sum_i\alpha_if(a_ix)$ approximates $c$ and
$\sum_j\beta_jf(xb_j)$ approximates $c'$, their double sum
$\sum_{ij}\alpha_i\beta_jf(a_ib_j)$ approximates both. Hence
$c=c'$, and each type of hull has just this one constant, denoted
$Mf$.

Positivity, $M1=1$, scalar homogeneity, and both translation
invariances follow from this uniqueness. To see additivity, choose
a left averaging operator $P$ making $Pf$ uniformly close to
$Mf$. Since $Pg$ lies in the translate convex hull of $g$, its
own translate convex hull has the same unique constant $Mg$.
Choose another averaging operator $Q$ making $QPg$ close to $Mg$.
Being a contraction fixing constants, $Q$ keeps $QPf$ close to
$Mf$. Thus $QP(f+g)$ is close to $Mf+Mg$, and uniqueness for
$f+g$ gives $M(f+g)=Mf+Mg$. Riesz--Markov supplies a regular
probability $m$ representing this positive functional, already
left and right invariant. This construction occurs entirely in
$C(G)$; no unproved equicontinuity assertion about translations
on the probability-measure space was used.

If $\nu$ is another left-invariant probability, for $f\in C(G)$
the double integral of $f(yx)$ with $d\nu(x)dm(y)$ is both
$\int f\,d\nu$ by left invariance of $\nu$ and $\int f\,dm$
by right invariance of $m$. Its interchange is legitimate by
uniform tensor approximation of the continuous function on the
compact product, exactly as for the kernel theorem. Hence
$\nu=m$. The pushforward by inversion interchanges left and
right invariance, proving inversion invariance by uniqueness.
If a nonempty open set had mass zero, its translates, finitely
many of which cover $G$, would force $m(G)=0$. This proves strict
positivity on opens. Compare \cite{B04}, 5.11--5.14.
\end{proof}

\section{Extreme measures and interpolation}

\begin{theorem}[Stone--Weierstrass and extreme scalar measures]
\label{fa2:p5:stone-weierstrass}\label{fa2:p5:extreme-measures}
A self-adjoint point-separating unital subalgebra of $C(K,\mathbb C)$,
$K$ nonempty compact Hausdorff, is uniformly dense. The real form requires
no self-adjointness. On a locally compact Hausdorff space the
$C_0$ version replaces unitality by nonvanishing at every point.
The extreme probability measures on $K$ are exactly the Dirac
measures. The extreme points of the complex measure unit ball
are exactly $z\delta_x$, $|z|=1$ (real signs in the real case).
\end{theorem}
\begin{proof}
For the real theorem, the uniform closure is a lattice:
$|f|$ is uniformly approximable by polynomials in $f$, because
$\sqrt t$ is uniformly polynomial-approximable on a compact
positive interval. For a direct scalar proof, rescale to $[0,1]$
and use $\sqrt{1-t}=1-\sum_{n\ge1}a_nt^n$, where $a_n>0$
and $\sum a_n=1$; the latter follows by letting $t\uparrow1$
in the binomial series. Hence the series is uniformly convergent
even at its endpoint. Thus maxima and minima belong to the
closure. Given continuous $g$ and $\varepsilon>0$, for every pair
$x,y$ separation and constants give an algebra function matching
$g$ at $x,y$. For fixed $x$ take a finite maximum of these
functions so that it is everywhere greater than $g-\varepsilon$
and equals $g$ at $x$. Near $x$ it is below $g+\varepsilon$.
A finite minimum over such neighborhoods gives an approximant
between $g-\varepsilon$ and $g+\varepsilon$ globally.
For the complex theorem apply the real result to real parts of
the self-adjoint algebra and then to imaginary parts.
For $C_0$ adjoin constants on the one-point compactification.
Nonvanishing separates each finite point from infinity; subtract
the value at infinity from an approximant to return to the
original algebra. This is exactly where nonvanishing is needed.

A non-Dirac regular probability has a Borel set of mass strictly
between zero and one: its support contains two distinct points
with disjoint positive-mass neighborhoods. Restriction and
normalization give a nontrivial convex decomposition.
Conversely a convex decomposition of $\delta_x$ forces both
positive measures to be supported at $x$, hence equal to it.
For a norm-one scalar measure write $d\nu=h\,d|\nu|$,
$|h|=1$, by scalar Radon--Nikodym. If its variation is not a
point mass, split it on a set of intermediate variation mass;
the two normalized restrictions have norm one and give a
nontrivial convex decomposition. Interior points of the ball
are not extreme, by a small perturbation in any nonzero measure
direction. Finally if $z\delta_x=(\nu_1+\nu_2)/2$ with
$\pvnorm{\nu_i}\le1$, evaluation of the singleton and the
triangle inequality force $\nu_i(\{x\})=z$ and zero remaining
variation. This proves their equality.
\end{proof}

\begin{theorem}[Lyapunov convexity by purification]
\label{fa2:p5:lyapunov}
The range of a finite-dimensional finite nonatomic real vector
measure is compact and convex. Here nonatomic means its total
variation has no atoms; equivalently one may dominate its
coordinate variations by a finite nonatomic positive measure.
The complex case is the corresponding real $2n$-dimensional case.
\end{theorem}
\begin{proof}
Let $\sigma=\sum_j|\mu_j|$ and write $d\mu_j=h_jd\sigma$.
The set $B=\{g\in L^\infty(\sigma):0\le g\le1\}$ is weak-star
compact convex, and $A(g)=(\int gh_jd\sigma)_j$ is continuous
into $\mathbb R^n$. Thus $A(B)$ is compact convex.
For $b\in A(B)$ its fiber $B_b$ has an extreme point by
Krein--Milman. If such a $g$ were not an indicator, some
$F=\{\varepsilon\le g\le1-\varepsilon\}$ would have positive
measure. Nonatomicity splits $F$ into $n+1$ disjoint sets of
positive measure (repeatedly split a positive non-atom).
Their $A$-images are linearly dependent. A nonzero linear
combination $v$ of their indicators with $A(v)=0$, scaled to
$\pvnorm v_\infty\le\varepsilon$, makes $g\pm v$ distinct
members of the same fiber, a contradiction. Hence $g=1_E$ a.e.,
so every $b\in A(B)$ is $\mu(E)$. The actual range therefore
equals $A(B)$ and is compact convex. Infinite-dimensional
targets are not covered by this finite-dependence argument.
\end{proof}

\begin{theorem}[Bishop's antisymmetric-set theorem]
\label{fa2:p5:bishop}
Let $\mathcal A$ be a closed subalgebra of $C(K,\mathbb C)$.
A subset is $\mathcal A$-antisymmetric when every algebra member
real-valued there is constant there. If $g\in C(K)$ agrees on
each maximal antisymmetric set with some member of $\mathcal A$,
then $g\in\mathcal A$.
\end{theorem}
\begin{proof}
Consider an extreme point $\mu$ of the weak-star compact unit
ball of $\mathcal A^\perp$. Unless this annihilator is zero,
every such extreme point has norm one. Let $S=\pvsupp\mu$.
If $f\in\mathcal A$ is real on $S$, shift and rescale it there
to a function $h$ with $\varepsilon<h<1-\varepsilon$.
Although this affine shift need not itself belong to $\mathcal A$,
both $h\mu$ and $(1-h)\mu$ annihilate $\mathcal A$ because
$\mu$ and $f\mu$ do. Put $a=\int h\,d|\mu|$.
Their normalized forms have norm one and decompose
$\mu=a(h\mu/a)+(1-a)((1-h)\mu/(1-a))$.
Extremality implies $h=a$ $|\mu|$-a.e. and hence on its support
by continuity. Thus $S$ is antisymmetric.
It lies in a maximal closed antisymmetric set: unions of chains
containing $S$ remain antisymmetric, and closure preserves this
property. By hypothesis $g$ agrees there with an algebra member,
so $\int g\,d\mu=0$.
Krein--Milman and weak-star continuity of this integral show the
same for every annihilating measure. Hahn--Banach separation of
the closed subspace $\mathcal A$ then gives $g\in\mathcal A$.
This proves the genuine non-self-adjoint generalization rather
than invoking the ordinary Stone--Weierstrass theorem for an
algebra that need not be self-adjoint.
\end{proof}

\begin{theorem}[Bishop restriction interpolation]
\label{fa2:p5:bishop-interpolation}
Let $Y$ be a closed subspace of $C(K)$ and $F\subset K$ compact.
If $|\mu|(F)=0$ for every $\mu\in Y^\perp$, then restriction
$R:Y\to C(F)$ maps its open unit ball onto the open unit ball.
In particular it is onto. The condition concerns total variation,
not merely the possibly cancelling scalar mass $\mu(F)$.
\end{theorem}
\begin{proof}
For $\nu\in C(F)'$, extend $R'\nu$ from $Y$ to $C(K)$ by
Hahn--Banach without increasing norm, and represent the extension
by $\eta$. Extend $\nu$ by zero off $F$. Then
$\eta-\nu\in Y^\perp$, hence its restriction to $F$ is zero.
Thus $\pvnorm\nu\le\pvnorm\eta=\pvnorm{R'\nu}\le\pvnorm\nu$.
So $R'$ is an isometry. Separation now shows that the norm closure
of $R(rB_Y)$ contains $rB_{C(F)}$ for every $r>0$; otherwise a
separating functional would contradict that isometry.
For $\pvnorm g<1$ choose radii $r_j>0$ with
$r_0>\pvnorm g$ and $\sum r_j<1$. Successively lift the current
residual with a vector of norm at most $r_j$, approximating it
closely enough that the next residual has norm less than
$r_{j+1}$. Completeness of $Y$ makes the series converge to a
vector of norm less than one whose image is exactly $g$.
This gives actual interpolation, not only density of restrictions.
\end{proof}

\begin{proposition}[Finite Blaschke interpolation]
\label{fa2:p5:blaschke-interpolation}
Let the finitely many interpolation nodes $z_j$ lie in the open
unit disk. If prescribed values at them admit a holomorphic
solution of $|f|\le1$, they admit a finite Blaschke-product
solution, including a unimodular constant as degree zero.
\end{proposition}
\begin{proof}
For empty data any unimodular constant suffices. If a prescribed
value has modulus one, the maximum principle makes the admissible
solution that constant, which again suffices. Otherwise disk
automorphisms reduce one datum to $z_0=w_0=0$.
The Herglotz transform $H=(1+f)/(1-f)$ has positive real part,
$H(0)=1$, and finitely many prescribed values. Its representing
probability on the circle is supplied by
Theorem~\ref{fa2:p4:herglotz}. The vector of these values is a
barycenter of the continuous curve
$\zeta\mapsto((\zeta+z_j)/(\zeta-z_j))_{j=1}^n$
in $\mathbb R^{2n}$. It is therefore a convex combination of
at most $2n+1$ points. For completeness, any larger finite
combination has affinely dependent vertices; subtract a scalar
multiple of that dependence until one coefficient is zero and
repeat. This bounds the required number, and compactness makes
this finite convex hull closed, so integral barycenters belong
to it as limits of simple barycenters.
Thus an atomic probability gives a rational positive-real-part
$H$ with the same data. The function $(H-1)/(H+1)$ is rational
and holomorphic in the disk, has modulus one on the circle,
and its apparent singularities at atoms cancel. Divide it by
the Blaschke factors of its finitely many disk zeros; the quotient
and reciprocal are holomorphic with boundary modulus one, so
the maximum principle makes it a unimodular constant. Undoing
the initial automorphisms preserves finite Blaschke products.
Degenerate constant unimodular data are already solved by a
degree-zero product.
\end{proof}

\section{Synthesis, summability, and weak-star quadrature}

\begin{theorem}[Quantitative absolute synthesis]
\label{fa2:p5:absolute-synthesis}
Let $(x_n)$ be a sequence in the unit ball of a Banach space.
If $\sup_n|\ell(x_n)|\ge\beta\pvnorm\ell$ for every $\ell\in X'$,
$\beta>0$, then for every $x$ and $\varepsilon>0$ there are
$a\in\ell^1$ with $x=\sum a_nx_n$ and
$\pvnorm a_1<\beta^{-1}\pvnorm x+\varepsilon$.
\end{theorem}
\begin{proof}
The synthesis map $T:\ell^1\to X$ is bounded, and its transpose
satisfies $\pvnorm{T'\ell}_\infty\ge\beta\pvnorm\ell$.
Hahn--Banach separation implies that the closure of $T(rB)$
contains the ball of radius $\beta r$: otherwise a separating
functional violates that inequality. The successive residual
lifting argument used in the preceding interpolation theorem,
with a summable norm budget arbitrarily close to
$\beta^{-1}\pvnorm x$, gives an actual $a\in\ell^1$ and the
stated norm bound. Completeness ensures the synthesized series
converges to $x$, not merely to a point in a larger completion.
\end{proof}

\begin{corollary}[Absolute Poisson-kernel synthesis]
\label{fa2:p5:poisson-synthesis}
Let $(z_n)$ in the disk be nontangentially dense at almost every
boundary point: for a.e. $\zeta$ some subsequence tends to
$\zeta$ in a fixed Stolz cone. Every $f\in L^1(\mathbb T)$ has
an absolutely convergent expansion $f=\sum a_nP_{z_n}$ with
$\sum|a_n|<\pvnorm f_1+\varepsilon$.
\end{corollary}
\begin{proof}
Normalized Poisson kernels have $L^1$ norm one. For
$g\in L^\infty$, its Poisson extension tends to $g(\zeta)$
at almost every Lebesgue point along each nontangential approach.
Indeed the kernel is bounded by
$C h/(h^2+\theta^2)$, $h=1-|z|$, in a Stolz cone.
Split its integral against $g-g(\zeta)$ into dyadic annuli of
radii $2^jh$. The Lebesgue-point mean oscillation is small on
small arcs; the geometric kernel weights sum, and the remaining
large arcs contribute $O(h)$. Thus the limit assertion follows
directly from scalar Lebesgue differentiation.
Nontangential density gives
$\sup_n|\int P_{z_n}g|\ge\pvnorm g_\infty$, while the reverse
inequality is the kernel norm. Apply the absolute synthesis
theorem to $X=L^1$ and $\beta=1$. This records the actual density
condition needed for the Poisson application.
\end{proof}

\begin{theorem}[Silverman--Toeplitz summability]
\label{fa2:p5:toeplitz-summability}
For rows $(a_{nk})_k\in\ell^1$, the matrix transformation
$y_n=\sum_ka_{nk}x_k$ sends every convergent sequence to a
sequence with the same limit if and only if
\[
 \sup_n\sum_k|a_{nk}|<\infty,\qquad
 a_{nk}\to0\ (k\text{ fixed}),\qquad
 \sum_ka_{nk}\to1.
\]
\end{theorem}
\begin{proof}
For necessity use coordinate vectors and the constant sequence.
Uniform boundedness on the Banach space $c$ gives a common row
functional norm; that norm is exactly the row $\ell^1$ norm by
finite phase tests. For sufficiency subtract the limiting constant
from $x$. Its tail is uniformly small, and the common row norm
controls that tail. The finite head tends to zero by the coordinate
condition; the row-sum condition handles the removed constant.
\end{proof}

\begin{proposition}[Quadrature is weak-star convergence]
\label{fa2:p5:quadrature}
On compact $K$, scalar finite measures $\mu_n$ converge against
every $C(K)$ function exactly when their variations are uniformly
bounded and their integrals converge on a uniformly dense linear
subspace, to the restriction of one bounded functional.
For compact metric $K$ and a finite measure $\mu$, partitions of
mesh tending to zero with weights $\mu(A_j)$ and nodes in each
cell give weak-star convergent quadrature measures.
\end{proposition}
\begin{proof}
Necessity is uniform boundedness and Riesz--Markov. For sufficiency
approximate a continuous test in uniform norm by a dense-subspace
test and use the common variation bound on the error; this also
extends the limiting functional uniquely. For the partition rule,
uniform continuity bounds the integral error by
$|\mu|(K)$ times the modulus of continuity at the mesh.
Boundedness cannot be omitted from the dense-test criterion:
on $[0,1]$, $\mu_n=n(\delta_{1/n^2}-\delta_0)$ tends to zero
on every polynomial, but its value on $\sqrt x$ is one and its
variation is $2n$. Thus testing a dense family alone is insufficient.
\end{proof}

\begin{proposition}[Positive Markov kernels and invariant probabilities]
\label{fa2:p5:markov-kernel}
Let $K$ be nonempty compact Hausdorff. A positive unital linear map
$P:C(K)\to C(K)$ has an invariant regular
probability measure. If it is given by a continuous kernel with
respect to a probability $\nu$ and $k(x,y)\ge\varepsilon>0$,
then that probability is unique and the iterates converge uniformly
to its integral, geometrically in oscillation.
\end{proposition}
\begin{proof}
The transpose maps the weak-star compact convex probability set
continuously into itself. The compact-convex fixed-point theorem
\ref{fa2:p1:schauder} gives an invariant probability.
For the strict-kernel case write
$P=\varepsilon J+(1-\varepsilon)Q$, with
$Jf=\int f\,d\nu$ and $Q$ positive unital (the case
$\varepsilon=1$ is already constant). Thus for real $f$,
$\operatorname{osc}(Pf)\le(1-\varepsilon)\operatorname{osc}(f)$.
If two invariant probabilities are tested against $P^nf$, their
difference is bounded by this vanishing oscillation, giving
uniqueness. The invariant integral lies between the minimum and
maximum of $P^nf$, so the same estimate gives uniform convergence.
Apply it to real and imaginary parts for complex functions.
This is an order-positive Banach-space model, without asserting
self-adjointness of the kernel or of $P$.
\end{proof}

\begin{example}[A continuous spectrum from an explicit multiplier]
\label{fa2:p5:continuous-spectrum-example}
On $L^2(\mathbb R)$ the inverse of $I-\partial_x^2$, with its
Sobolev domain $H^2$, is the Fourier multiplier
$m(\xi)=(1+4\pi^2\xi^2)^{-1}$. It is bounded self-adjoint with
norm one and spectrum $[0,1]$, has no eigenvalues, and is not compact.
\end{example}
\begin{proof}
Plancherel identifies the inverse with multiplication by $m$;
$(1+4\pi^2\xi^2)m\widehat f=\widehat f\in L^2$, so $\mathcal F^{-1}(m\widehat f)\in H^2$.
The multiplier identity verifies both inverse equations. The spectrum is the essential range:
outside it the reciprocal is essentially bounded, while near a
value in it normalized indicators give arbitrarily small residuals.
Every level set of $m$ has Lebesgue measure zero, excluding
eigenvectors. Its essential range is the closed interval $[0,1]$.
Infinitely many disjoint subintervals on which $m\ge1/2$ give
orthonormal indicator vectors whose images remain separated in
norm, so the operator is not compact. This realizes the exercise
counterexample of \cite{B01}, 8.38, with its domain actually checked.
\end{proof}


% END INPUT

\backmatter
\begin{thebibliography}{B09}
\bibitem{B01} H.~Brezis, \emph{Functional Analysis, Sobolev Spaces and Partial Differential Equations}, Springer, 2011.
\bibitem{B02} J.~B.~Conway, \emph{A Course in Functional Analysis}, second edition, Springer, supplied 2007 reprint.
\bibitem{B03} E.~Kreyszig, \emph{Introductory Functional Analysis with Applications}, Wiley, supplied 1989 reprint.
\bibitem{B04} W.~Rudin, \emph{Functional Analysis}, second edition, McGraw-Hill, 1991.
\bibitem{B05} R.~E.~Megginson, \emph{An Introduction to Banach Space Theory}, Springer, 1998.
\bibitem{B06} R.~Kress, \emph{Linear Integral Equations}, third edition, Springer, 2014.
\bibitem{B07} H.~Brunner, \emph{Volterra Integral Equations}, Cambridge University Press, 2017.
\bibitem{B08} P.~D.~Lax, \emph{Functional Analysis}, Wiley, 2002.
\bibitem{B09} E.~M.~Stein and R.~Shakarchi, \emph{Functional Analysis: Introduction to Further Topics in Analysis}, Princeton University Press, 2011.
\bibitem{Enflo1973} P.~Enflo, A counterexample to the approximation problem in Banach spaces, \emph{Acta Mathematica} 130 (1973), 309--317. \href{https://doi.org/10.1007/BF02392270}{doi:10.1007/BF02392270}.
\end{thebibliography}
\end{document}
